
\documentclass[aps,preprint]{revtex4}
%%%%%%%%%%%%%%%%%%%%%%%%%%%%%%%%%%%%%%%%%%%%%%%%%%%%%%%%%%%%%%%%%%%%%%%%%%%%%%%%%%%%%%%%%%%%%%%%%%%%%%%%%%%%%%%%%%%%%%%%%%%%%%%%%%%%%%%%%%%%%%%%%%%%%%%%%%%%%%%%%%%%%%%%%%%%%%%%%%%%%%%%%%%%%%%%%%%%%%%%%%%%%%%%%%%%%%%%%%%%%%%%%%%%%%%%%%%%%%%%%%%%%%%%%%%%
\usepackage{amssymb}
\usepackage{amsfonts}
\usepackage{amsmath}
\usepackage{graphicx}

\setcounter{MaxMatrixCols}{10}
%TCIDATA{OutputFilter=LATEX.DLL}
%TCIDATA{Version=5.50.0.2960}
%TCIDATA{<META NAME="SaveForMode" CONTENT="1">}
%TCIDATA{BibliographyScheme=Manual}
%TCIDATA{Created=Monday, October 20, 2014 14:28:15}
%TCIDATA{LastRevised=Monday, March 09, 2026 06:25:26}
%TCIDATA{<META NAME="GraphicsSave" CONTENT="32">}
%TCIDATA{<META NAME="DocumentShell" CONTENT="Articles\SW\REVTeX 4">}
%TCIDATA{Language=American English}
%TCIDATA{CSTFile=revtex4.cst}

\newtheorem{theorem}{Theorem}
\newtheorem{acknowledgement}[theorem]{Acknowledgement}
\newtheorem{algorithm}[theorem]{Algorithm}
\newtheorem{axiom}[theorem]{Axiom}
\newtheorem{claim}[theorem]{Claim}
\newtheorem{conclusion}[theorem]{Conclusion}
\newtheorem{condition}[theorem]{Condition}
\newtheorem{conjecture}[theorem]{Conjecture}
\newtheorem{corollary}[theorem]{Corollary}
\newtheorem{criterion}[theorem]{Criterion}
\newtheorem{definition}[theorem]{Definition}
\newtheorem{example}[theorem]{Example}
\newtheorem{exercise}[theorem]{Exercise}
\newtheorem{lemma}[theorem]{Lemma}
\newtheorem{notation}[theorem]{Notation}
\newtheorem{problem}[theorem]{Problem}
\newtheorem{proposition}[theorem]{Proposition}
\newtheorem{remark}[theorem]{Remark}
\newtheorem{solution}[theorem]{Solution}
\newtheorem{summary}[theorem]{Summary}
\newenvironment{proof}[1][Proof]{\noindent\textbf{#1.} }{\ \rule{0.5em}{0.5em}}
\input{tcilatex}
\begin{document}

\preprint{HEP/123-qed}
\title[Short title for running header]{REV\TeX4}
\author{First Author}
\affiliation{My Institution}
\author{Second Author}
\affiliation{My Institution}
\author{Third Author}
\affiliation{Other Institution}
\keywords{one two three}
\pacs{PACS number}

\begin{abstract}
Shell document for REV\TeX{} 4.
\end{abstract}

\volumeyear{year}
\volumenumber{number}
\issuenumber{number}
\eid{identifier}
\date[Date text]{date}
\received[Received text]{date}
\revised[Revised text]{date}
\accepted[Accepted text]{date}
\published[Published text]{date}
\startpage{101}
\endpage{102}
\maketitle
\tableofcontents

\part{Complex Clifford and CAR Algebras}

\begin{notation}
\end{notation}

\begin{enumerate}
\item $V$ is a\textbf{\ real} vector space of dimension $r$ i.e.%
\begin{equation}
\dim \left( V\right) =r
\end{equation}

\item 
\begin{equation}
V_{\mathbb{C}}=\mathbb{C}\otimes _{\mathbb{R}}V=V\oplus _{\mathbb{R}}iV\cong
V+iV
\end{equation}%
\begin{equation}
\dim \left( V_{\mathbb{C}}\right) =r
\end{equation}

\item 
\begin{equation}
V_{\mathbb{R}}=V\oplus _{\mathbb{R}}V,
\end{equation}%
\begin{equation}
\dim \left( V_{\mathbb{R}}\right) =2r
\end{equation}

\item 
\begin{equation}
\mathcal{V}=\left\{ \left( V,0\right) \right\} ;\,i\mathcal{V}=\left\{
\left( 0,iV\right) \right\}
\end{equation}

\item 
\begin{equation}
\mathcal{V}_{1}=\left\{ \left( V,0\right) \right\} ;\,\mathcal{V}%
_{2}=\left\{ \left( 0,V\right) \right\} ,
\end{equation}

\item 
\begin{equation}
V_{\mathbb{C}}=V\oplus _{\mathbb{R}}iV=\mathcal{V}+i\mathcal{V}\cong V+iV
\end{equation}%
and 
\begin{equation}
V_{\mathbb{R}}=V\oplus _{\mathbb{R}}V=\mathcal{V}_{1}+\mathcal{V}_{2}\ncong
V+V.
\end{equation}

\item One can define the complex isomorphism $J_{V}\in \mathcal{L}\left(
V\right) $ 
\begin{equation}
W=J_{V}V\cong \mathcal{J}_{V}\mathcal{V}=\mathcal{W}
\end{equation}
\end{enumerate}

\section{Transformation of Bases and Super Notation}

\begin{notation}
\textbf{Super vectors} are $n\times 1$ or $1\times n$ arrays of elements
that belong to sets that do not belong to mathematical fields but have a
linear operations defined.
\end{notation}

\begin{enumerate}
\item Transformation of basis $T=\left( \mathbf{f}\right) ^{t}\left[ \left( 
\mathbf{e}_{b}\right) \right] \in \mathcal{L}\left( V\right) $ of a vector
space $V$ from $\left( \mathbf{e}\right) ^{t}$ to $\left( \mathbf{f}\right)
^{t}$ 
\begin{eqnarray}
\left( \mathbf{f}\right) ^{t} &=&T\left( \mathbf{e}\right) ^{t}=\left( 
\mathbf{e}\right) ^{t}\mathcal{R}_{e}\left( T\right) \\
\left( \mathbf{e}\right) ^{t} &=&T^{-1}\left( \mathbf{f}\right) ^{t}=\left( 
\mathbf{f}\right) ^{t}\mathcal{R}_{f}\left( T\right) ^{-1}.
\end{eqnarray}

\item 
\begin{eqnarray}
\left( \mathbf{f}_{b}\right) ^{t} &=&T_{b}\left( \mathbf{e}_{b}\right)
^{t}=\left( \mathbf{e}_{b}\right) ^{t}\mathcal{R}_{e_{b}}\left( T_{b}\right)
=\left( \mathbf{e}_{b}\right) ^{t}\mathcal{R}_{e}\left( T\right) ^{-t} \\
\left( \mathbf{e}_{b}\right) ^{t} &=&T_{b}^{-1}\left( \mathbf{f}_{b}\right)
^{t}=\left( \mathbf{f}_{b}\right) ^{t}\mathcal{R}_{f_{b}}\left( T_{b}\right)
^{-1}=\left( \mathbf{f}_{b}\right) ^{t}\mathcal{R}_{e}\left( T_{b}\right)
^{t}.
\end{eqnarray}

\item The definition of $T_{b}$ is 
\begin{equation}
\left( \mathbf{f}_{b}\right) ^{t}=T_{b}\left( \mathbf{e}_{b}\right)
^{t}=\left( T\mathbf{e}\right) _{b}^{t}.
\end{equation}

\item 
\begin{equation}
\left( \mathbf{e}_{b}\right) ^{t}\left( \mathbf{e}\right) =\left[ \left( 
\begin{array}{c}
e_{b1} \\ 
\vdots \\ 
e_{bn}%
\end{array}%
\right) \right] \otimes \left( 
\begin{array}{ccc}
e_{1} & \cdots & e_{n}%
\end{array}%
\right) =\left[ \left( 
\begin{array}{c}
e_{b1} \\ 
\vdots \\ 
e_{bn}%
\end{array}%
\right) \right] \left( 
\begin{array}{ccc}
e_{1} & \cdots & e_{n}%
\end{array}%
\right) =\left( 
\begin{array}{ccc}
e_{b1}\cdot e_{1} & \cdots & e_{b1}\cdot e_{n} \\ 
\vdots & \ddots & \vdots \\ 
e_{bn}\cdot e_{1} & \cdots & e_{bn}\cdot e_{n}%
\end{array}%
\right) =\left( 
\begin{array}{ccc}
e_{b1}\left( e_{1}\right) & \cdots & e_{b1}\left( e_{n}\right) \\ 
\vdots & \ddots & \vdots \\ 
e_{bn}\left( e_{1}\right) & \cdots & e_{bn}\left( e_{n}\right)%
\end{array}%
\right)
\end{equation}

\item 
\begin{gather}
\left( \mathbf{f}_{b}\right) ^{t}\left( \mathbf{f}\right) =\left( 
\begin{array}{c}
f_{b1} \\ 
\vdots \\ 
f_{bn}%
\end{array}%
\right) \otimes \left( 
\begin{array}{ccc}
f_{1} & \cdots & f_{n}%
\end{array}%
\right) =\left( 
\begin{array}{c}
f_{b1} \\ 
\vdots \\ 
f_{bn}%
\end{array}%
\right) \left( 
\begin{array}{ccc}
f_{1} & \cdots & f_{n}%
\end{array}%
\right) =\left( 
\begin{array}{ccc}
f_{b1}\cdot f_{1} & \cdots & f_{b1}\cdot f_{n} \\ 
\vdots & \ddots & \vdots \\ 
f_{bn}\cdot f_{1} & \cdots & f_{bn}\cdot f_{n}%
\end{array}%
\right) =\left( 
\begin{array}{ccc}
f_{b1}\left( f_{1}\right) & \cdots & f_{b1}\left( f_{n}\right) \\ 
\vdots & \ddots & \vdots \\ 
f_{bn}\left( f_{1}\right) & \cdots & f_{bn}\left( f_{n}\right)%
\end{array}%
\right) \\
\parallel \\
\left( 
\begin{array}{c}
\left( Te\right) _{b1} \\ 
\vdots \\ 
\left( Te\right) _{bn}%
\end{array}%
\right) \otimes \left( 
\begin{array}{ccc}
\left( Te\right) _{1} & \cdots & \left( Te\right) _{n}%
\end{array}%
\right) =\left( 
\begin{array}{c}
\left( Te\right) _{1} \\ 
\vdots \\ 
\left( Te\right) _{n}%
\end{array}%
\right) \otimes \left( 
\begin{array}{ccc}
\left( Te\right) _{b1} & \cdots & \left( Te\right) _{bn}%
\end{array}%
\right) \\
=\left( 
\begin{array}{ccc}
\left( Te\right) _{b1}\cdot \left( Te\right) _{1} & \cdots & \left(
Te\right) _{b1}\cdot \left( Te\right) _{n} \\ 
\vdots & \ddots & \vdots \\ 
\left( Te\right) _{bn}\cdot \left( Te\right) _{1} & \cdots & \left(
Te\right) _{bn}\cdot \left( Te\right) _{n}%
\end{array}%
\right) =\left( 
\begin{array}{ccc}
\left( Te\right) _{b1}\left( \left( Te\right) _{1}\right) & \cdots & \left(
Te\right) _{b1}\left( \left( Te\right) _{n}\right) \\ 
\vdots & \ddots & \vdots \\ 
\left( Te\right) _{bn}\left( \left( Te\right) _{1}\right) & \cdots & \left(
Te\right) _{bn}\left( \left( Te\right) _{n}\right)%
\end{array}%
\right) \\
\parallel \\
\mathcal{R}_{e}\left( T\right) ^{-1}\left( 
\begin{array}{c}
e_{b1} \\ 
\vdots \\ 
e_{bn}%
\end{array}%
\right) \left( 
\begin{array}{ccc}
e_{1} & \cdots & e_{n}%
\end{array}%
\right) \mathcal{R}_{e}\left( T\right) =\mathcal{R}_{e}\left( T\right) ^{-1}%
\mathcal{R}_{e}\left( T\right) =\mathbb{I}_{n}.
\end{gather}

\item This gives 
\begin{gather}
T_{b}\left( 
\begin{array}{ccc}
e_{b1} & \cdots & e_{bn}%
\end{array}%
\right) =\left( 
\begin{array}{ccc}
e_{b1} & \cdots & e_{bn}%
\end{array}%
\right) \mathcal{R}_{e_{b}}\left( T_{b}\right) \\
\Updownarrow \\
\left( 
\begin{array}{c}
T_{b}e_{b1} \\ 
\vdots \\ 
T_{b}e_{bn}%
\end{array}%
\right) =\left( 
\begin{array}{c}
\left( Te_{1}\right) _{b} \\ 
\vdots \\ 
\left( Te_{n}\right) _{b}%
\end{array}%
\right) =\mathcal{R}_{e}\left( T\right) ^{-1}\left( 
\begin{array}{c}
e_{b1} \\ 
\vdots \\ 
e_{bn}%
\end{array}%
\right) =\mathcal{R}_{e_{b}}\left( T_{b}\right) ^{t}\left( 
\begin{array}{c}
e_{b1} \\ 
\vdots \\ 
e_{bn}%
\end{array}%
\right)
\end{gather}%
and thus 
\begin{equation}
\left( \overset{}{\mathcal{R}_{e_{b}}\left( T_{b}\right) ^{t}\equiv \mathcal{%
R}_{e}\left( T\right) ^{-1}}\right) \Leftrightarrow .\left( \overset{}{%
\mathcal{R}_{e_{b}}\left( T_{b}\right) \equiv \mathcal{R}_{e}\left( T\right)
^{-t}}\right) .
\end{equation}

\item Biorthogonality preserving transformations on the dual space are given
by 
\begin{eqnarray}
\left( 
\begin{array}{ccc}
f_{b1} & \cdots & f_{bn}%
\end{array}%
\right) &=&\left( T_{fe}\right) _{b}\left( 
\begin{array}{ccc}
e_{b1} & \cdots & e_{bn}%
\end{array}%
\right) =\left( 
\begin{array}{ccc}
e_{b1} & \cdots & e_{bn}%
\end{array}%
\right) \mathcal{R}_{f}\left( T_{fe}\right) ^{-t} \\
&\parallel & \\
\left( \mathbf{f}_{b}\right) ^{t} &=&\left( T_{fe}\right) _{b}\left( \mathbf{%
e}_{b}\right) ^{t}=\left( \mathbf{e}_{b}\right) ^{t}\mathcal{R}_{e}\left(
T_{fe}\right) ^{-t}.
\end{eqnarray}
\end{enumerate}

\begin{remark}
\begin{enumerate}
\item 
\begin{enumerate}
\item 
\begin{remark}
The tensor product of super vectors to give super matrices are given in
terms of\textbf{\ the algebraic products} bewteen the elements the of the
super vectors. By choosing different types of products this can lead to
different typses of super matrix elements.
\end{remark}
\end{enumerate}

\item Resolutions of the Identity are given by 
\begin{equation}
\left( 
\begin{array}{ccc}
e_{1} & \cdots & e_{n}%
\end{array}%
\right) \cdot \left( 
\begin{array}{c}
e_{b1} \\ 
\vdots \\ 
e_{bn}%
\end{array}%
\right) =\limfunc{Tr}\left( \left( 
\begin{array}{ccc}
e_{1} & \cdots & e_{n}%
\end{array}%
\right) \otimes \left( 
\begin{array}{c}
e_{b1} \\ 
\vdots \\ 
e_{bn}%
\end{array}%
\right) \right) =\left( 
\begin{array}{ccc}
e_{1} & \cdots & e_{n}%
\end{array}%
\right) \left[ \left( 
\begin{array}{c}
e_{b1} \\ 
\vdots \\ 
e_{bn}%
\end{array}%
\right) \right] =\sum_{1\leq j\leq n}e_{j}\left[ e_{bj}\right] =I_{n}.
\end{equation}

\item Notation: Dot product of super vectors 
\begin{equation}
\left( 
\begin{array}{ccc}
e_{1} & \cdots & e_{n}%
\end{array}%
\right) \cdot \left( 
\begin{array}{c}
e_{b1} \\ 
\vdots \\ 
e_{bn}%
\end{array}%
\right) =\left( 
\begin{array}{ccc}
e_{1} & \cdots & e_{n}%
\end{array}%
\right) \left[ \left( 
\begin{array}{c}
e_{b1} \\ 
\vdots \\ 
e_{bn}%
\end{array}%
\right) \right]
\end{equation}

\item Representations of operators $A\in \mathcal{L}\left( V\right) $ are
defined by 
\begin{eqnarray}
A\left( \mathbf{f}\right) ^{t} &=&\left( \mathbf{f}\right) ^{t}\mathcal{R}%
_{f}\left( A\right) \equiv \left( 
\begin{array}{ccc}
f_{1} & \cdots & f_{n}%
\end{array}%
\right) \mathcal{R}_{f}\left( A\right) \\
A\left( \mathbf{e}\right) ^{t} &=&\left( \mathbf{e}\right) ^{t}\mathcal{R}%
_{e}\left( A\right) \equiv \left( 
\begin{array}{ccc}
e_{1} & \cdots & e_{n}%
\end{array}%
\right) \mathcal{R}_{e}\left( A\right)
\end{eqnarray}

\item Dyadic expansions of operators are given by

\begin{enumerate}
\item 
\begin{equation}
I_{n}AI_{n}=\left( 
\begin{array}{ccc}
e_{1} & \cdots & e_{n}%
\end{array}%
\right) \left[ \left( 
\begin{array}{c}
e_{b1} \\ 
\vdots \\ 
e_{bn}%
\end{array}%
\right) \right] A\left( 
\begin{array}{ccc}
e_{1} & \cdots & e_{n}%
\end{array}%
\right) \left[ \left( 
\begin{array}{c}
e_{b1} \\ 
\vdots \\ 
e_{bn}%
\end{array}%
\right) \right] =\left( 
\begin{array}{ccc}
e_{1} & \cdots & e_{n}%
\end{array}%
\right) \left( 
\begin{array}{ccc}
\mathcal{R}_{e}\left( A\right) & \cdots & \mathcal{R}_{e}\left( A\right) \\ 
\vdots & \ddots & \vdots \\ 
\mathcal{R}_{e}\left( A\right) & \cdots & \mathcal{R}_{e}\left( A\right)%
\end{array}%
\right) \left[ \left( 
\begin{array}{c}
e_{b1} \\ 
\vdots \\ 
e_{bn}%
\end{array}%
\right) \right] .
\end{equation}

\item The change of basis 
\begin{equation}
T=\left( \boldsymbol{f}\right) \left[ \left( \boldsymbol{e}_{b}\right) %
\right]
\end{equation}%
gives 
\begin{eqnarray}
I_{n}AI_{n} &=&\left( 
\begin{array}{ccc}
f_{1} & \cdots & f_{n}%
\end{array}%
\right) \left[ \left( 
\begin{array}{c}
f_{b1} \\ 
\vdots \\ 
f_{bn}%
\end{array}%
\right) \right] A\left( 
\begin{array}{ccc}
f_{1} & \cdots & f_{n}%
\end{array}%
\right) \left[ \left( 
\begin{array}{c}
f_{b1} \\ 
\vdots \\ 
f_{bn}%
\end{array}%
\right) \right] \\
&=&\left( 
\begin{array}{ccc}
e_{1} & \cdots & e_{n}%
\end{array}%
\right) \mathcal{R}_{e}\left( T\right) \left[ \left( 
\begin{array}{c}
f_{b1} \\ 
\vdots \\ 
f_{bn}%
\end{array}%
\right) \right] A\left( 
\begin{array}{ccc}
f_{1} & \cdots & f_{n}%
\end{array}%
\right) \mathcal{R}_{e}\left( T\right) ^{-1}\left[ \left( 
\begin{array}{c}
e_{b1} \\ 
\vdots \\ 
e_{bn}%
\end{array}%
\right) \right] \\
&=&\left( 
\begin{array}{ccc}
e_{1} & \cdots & e_{n}%
\end{array}%
\right) \mathcal{R}_{e}\left( T\right) \left[ \left( 
\begin{array}{c}
f_{b1} \\ 
\vdots \\ 
f_{bn}%
\end{array}%
\right) \right] A\left( 
\begin{array}{ccc}
f_{1} & \cdots & f_{n}%
\end{array}%
\right) \mathcal{R}_{e}\left( T\right) ^{-1}\left[ \left( 
\begin{array}{c}
e_{b1} \\ 
\vdots \\ 
e_{bn}%
\end{array}%
\right) \right] .
\end{eqnarray}

\item Noting that $\mathcal{R}_{e}\left( T\right) =\mathcal{R}_{f}\left(
T\right) $ this gives 
\begin{equation}
\left( 
\begin{array}{ccc}
e_{1} & \cdots & e_{n}%
\end{array}%
\right) \left( 
\begin{array}{ccc}
\mathcal{R}_{e}\left( A\right) _{11} & \cdots & \mathcal{R}_{e}\left(
A\right) _{1n} \\ 
\vdots & \ddots & \vdots \\ 
\mathcal{R}_{e}\left( A\right) _{n1} & \cdots & \mathcal{R}_{e}\left(
A\right) _{nn}%
\end{array}%
\right) \left[ \left( 
\begin{array}{c}
e_{b1} \\ 
\vdots \\ 
e_{bn}%
\end{array}%
\right) \right] =\left( 
\begin{array}{ccc}
e_{1} & \cdots & e_{n}%
\end{array}%
\right) \mathcal{R}_{f}\left( T\right) \left( 
\begin{array}{ccc}
\mathcal{R}_{f}\left( A\right) _{11} & \cdots & \mathcal{R}_{f}\left(
A\right) _{1n} \\ 
\vdots & \ddots & \vdots \\ 
\mathcal{R}_{f}\left( A\right) _{n1} & \cdots & \mathcal{R}_{f}\left(
A\right) _{nn}%
\end{array}%
\right) \mathcal{R}_{f}\left( T\right) ^{-1}\left[ \left( 
\begin{array}{c}
e_{b1} \\ 
\vdots \\ 
e_{bn}%
\end{array}%
\right) \right] ,
\end{equation}

\item showing that 
\begin{equation}
\left( 
\begin{array}{ccc}
\mathcal{R}_{e}\left( A\right) _{11} & \cdots & \mathcal{R}_{e}\left(
A\right) _{1n} \\ 
\vdots & \ddots & \vdots \\ 
\mathcal{R}_{e}\left( A\right) _{n1} & \cdots & \mathcal{R}_{e}\left(
A\right) _{nn}%
\end{array}%
\right) =\mathcal{R}_{f}\left( T\right) \left( 
\begin{array}{ccc}
\mathcal{R}_{f}\left( A\right) _{11} & \cdots & \mathcal{R}_{f}\left(
A\right) _{1n} \\ 
\vdots & \ddots & \vdots \\ 
\mathcal{R}_{f}\left( A\right) _{n1} & \cdots & \mathcal{R}_{f}\left(
A\right) _{nn}%
\end{array}%
\right) \mathcal{R}_{f}\left( T\right) ^{-1}.
\end{equation}
\end{enumerate}
\end{enumerate}

\begin{lemma}
\begin{gather}
A=\left( \mathbf{f}\right) \mathcal{R}_{f}\left( A\right) \left( \mathbf{f}%
_{b}\right) ^{t}=\left( \mathbf{e}\right) \mathcal{R}_{e}\left( A\right)
\left( \mathbf{e}_{b}\right) ^{t}=\left( \mathbf{f}\right) \mathcal{R}%
_{e}\left( T\right) \mathcal{R}_{e}\left( A\right) \mathcal{R}_{e}\left(
T\right) ^{-t}\left( \mathbf{f}_{b}\right) ^{t} \\
\parallel \\
AI_{n}=\left( 
\begin{array}{c}
f_{1} \\ 
\vdots \\ 
f_{n}%
\end{array}%
\right) \mathcal{R}_{f}\left( A\right) \left( 
\begin{array}{ccc}
f_{b1} & \cdots & f_{bn}%
\end{array}%
\right) =\left( 
\begin{array}{c}
e_{1} \\ 
\vdots \\ 
e_{n}%
\end{array}%
\right) \mathcal{R}_{e}\left( A\right) \left( 
\begin{array}{ccc}
e_{b1} & \cdots & e_{bn}%
\end{array}%
\right) =\mathcal{R}_{e}\left( T\right) ^{t}\left( 
\begin{array}{c}
f_{1} \\ 
\vdots \\ 
f_{n}%
\end{array}%
\right) \mathcal{R}_{e}\left( A\right) \left( 
\begin{array}{ccc}
f_{b1} & \cdots & f_{bn}%
\end{array}%
\right) \mathcal{R}_{e}\left( T\right) ^{-t} \\
=\left( 
\begin{array}{c}
f_{1} \\ 
\vdots \\ 
f_{n}%
\end{array}%
\right) \mathcal{R}_{e}\left( T\right) ^{-t}\mathcal{R}_{e}\left( A\right) 
\mathcal{R}_{e}\left( T\right) ^{t}\left( 
\begin{array}{ccc}
f_{b1} & \cdots & f_{bn}%
\end{array}%
\right)
\end{gather}%
i.e. 
\begin{equation}
\mathcal{R}_{f}\left( A\right) =\mathcal{R}_{e}\left( T\right) \mathcal{R}%
_{e}\left( A\right) \mathcal{R}_{e}\left( T\right) ^{-1}.\ 
\end{equation}
\end{lemma}
\end{remark}

\section{CAR-Fock Algebra}

\begin{definition}
The CAR-Fock algebra is defined as 
\begin{equation}
\mathcal{F=}\dbigwedge \left( \overset{}{\mathcal{F}_{1}}\right)
=\dbigoplus_{N=0}^{N=2r}\mathcal{F}_{N}=\dbigoplus_{N=0}^{N=2r}\tbigwedge%
\nolimits^{N}\mathcal{F}_{1}=\mathcal{B}\left( \overset{}{\mathcal{H}_{%
\mathcal{F}}}\right) ,
\end{equation}%
where 
\begin{equation}
\mathcal{H}_{\mathcal{F}}=\dbigoplus_{N=0}^{N=r}\mathcal{H}%
^{N}=\dbigoplus_{N=0}^{N=r}\tbigwedge\nolimits^{N}\mathcal{H}^{1}=\dbigwedge
\left( \overset{}{\mathcal{H}^{1}}\right) .
\end{equation}
\end{definition}

The CAR structure of $\mathcal{F}$ is induced by the CAR 
\begin{gather}
\left\{ a_{j},a_{k}\right\} =0 \\
\left\{ a_{j}^{\dagger },a_{k}\right\} =\delta _{jk}I_{\mathcal{F}} \\
1\leq j,k\leq r,
\end{gather}%
where 
\begin{equation}
a_{j}^{\dagger }\in \mathcal{B}\left( \mathcal{H}_{\mathcal{F}}\right)
,1\leq j\leq r.
\end{equation}

The Clifford algebra structure of $\mathcal{F}$ is induced by 
\begin{gather}
\left\{ x_{j},x_{k}\right\} =0 \\
\ x_{j}^{2}=\frac{1}{2}I_{\mathcal{F}} \\
1\leq j,k\leq r,
\end{gather}%
where 
\begin{equation}
\left( \boldsymbol{x}\right) =\left( 
\begin{array}{cc}
\boldsymbol{x}_{1} & \boldsymbol{x}_{2}%
\end{array}%
\right) =K\left( 
\begin{array}{cc}
\boldsymbol{a}^{\dagger } & \boldsymbol{a}%
\end{array}%
\right) =\left( 
\begin{array}{cc}
\boldsymbol{a}^{\dagger } & \boldsymbol{a}%
\end{array}%
\right) \mathbb{K}=\left( 
\begin{array}{cc}
\boldsymbol{a}^{\dagger } & \boldsymbol{a}%
\end{array}%
\right) \frac{1}{\sqrt{2}}\left( 
\begin{array}{cc}
\mathbb{I}_{r} & i\mathbb{I}_{r} \\ 
\mathbb{I}_{r} & -i\mathbb{I}_{r}%
\end{array}%
\right) .
\end{equation}

\begin{remark}
\begin{equation}
x_{j}=x_{j}^{\dagger },1\leq j\leq r.
\end{equation}
\end{remark}

In particular 
\begin{eqnarray}
x_{1} &=&\frac{1}{\sqrt{2}}\left( a^{\dagger }+a\right) \\
x_{2} &=&\frac{i}{\sqrt{2}}\left( a^{\dagger }-a\right)
\end{eqnarray}%
showing that 
\begin{equation}
x^{2}=\frac{1}{\sqrt{2}}\left( a^{\dagger }+a\right) \frac{1}{\sqrt{2}}%
\left( a^{\dagger }+a\right) =\frac{1}{2}\left( a^{\dagger }a+aa^{\dagger
}\right) =\frac{1}{2}I_{\mathcal{F}}.
\end{equation}

\begin{definition}
\begin{enumerate}
\item Zero Grade filtering 
\begin{equation}
\left\langle {}\right\rangle _{0}:\mathcal{F\rightarrow F}_{1}
\end{equation}%
by 
\begin{equation}
\left\langle A\right\rangle _{0}=\left\langle A_{0}+\cdots
+A_{2^{2r}}\right\rangle _{0}=A_{0},\quad A_{j}\in \mathcal{F}_{j},\,1\leq
j\leq 2^{2r}.
\end{equation}

\item Scalar projection 
\begin{equation}
\left\langle \left\langle {}\right\rangle \right\rangle _{0}:\mathcal{F}%
_{1}\rightarrow \mathbb{C}
\end{equation}%
given by 
\begin{equation}
\left\langle \left\langle zI_{\mathcal{F}}\right\rangle \right\rangle _{0}=z.
\end{equation}
\end{enumerate}
\end{definition}

\begin{definition}
A symmetric complex bilinear form 
\begin{equation}
\left( \left. {}\right\vert \right) _{\mathcal{F}_{1}}:\mathcal{F}_{1}\times 
\mathcal{F}_{1}\rightarrow \mathbb{C}
\end{equation}%
is defined by 
\begin{equation}
\left( \left. A\right\vert B\right) _{\mathcal{F}_{1}}=\left\langle
\left\langle \left\{ A,B\right\} \right\rangle \right\rangle
_{0}=\left\langle \left\langle \left( AB+BA\right) \right\rangle
\right\rangle _{0},\quad A,B\in \mathcal{F}_{1},
\end{equation}%
where%
\begin{equation}
\left( \left. A\right\vert B\right) _{\mathcal{F}_{1}}I_{\mathcal{F}%
}=\left\{ A,B\right\} =AB+BA,\quad A,B\in \mathcal{F}_{1}.
\end{equation}
\end{definition}

\begin{definition}
Clifford Multiplication 
\begin{equation}
A_{1}\diamond B_{1}=\frac{1}{2}\left( \left( \left. A_{1}\right\vert
B_{1}\right) _{\mathcal{F}_{1}}I_{\mathcal{F}}+A_{1}{\large \wedge }%
B_{1}\right) .
\end{equation}
\end{definition}

\begin{remark}
Clifford Multiplication is the same as operator multiplication in $\mathcal{F%
}$ 
\begin{equation}
A_{1}\diamond B_{1}=\frac{1}{2}\left( \left\{ A_{1},B_{1}\right\} I_{%
\mathcal{F}}+A_{1}{\large \wedge }B_{1}\right) =\frac{1}{2}\left(
A_{1}B_{1}+B_{1}A_{1}+A_{1}B_{1}-B_{1}A_{1}\right) =A_{1}B_{1}.
\end{equation}
\end{remark}

\begin{remark}
In the CAR algebra the commutator $\left[ ,\right] $ acting on $\mathcal{F}%
_{1}\times \mathcal{F}_{1}$ defines an antisymmetric product $\equiv $ wedge
product "$\wedge $" and the anti commutator acting on $\mathcal{F}_{1}\times 
\mathcal{F}_{1}$ defines a Jordan product $\left( \text{which is a symmetric
product}\right) ,$ whose scalar projection defines a complex symmetric inner
product $\left( \left. {}\right\vert \right) _{\mathcal{F}_{1}}$ on $%
\mathcal{F}_{1}\times \mathcal{F}_{1}.$
\end{remark}

\subsection{ Operator Valued Representations}

\begin{definition}
\textbf{Super Vector Basis:} $\left( \boldsymbol{X}\right) $ is a
partitioned super vector with elements from $\mathcal{F}$ expressed as 
\begin{equation}
\left( \boldsymbol{X}\right) _{S}=\left( \boldsymbol{X}_{0},\cdots ,%
\boldsymbol{X}_{N},\cdots ,\boldsymbol{X}_{2r}\right) ,
\end{equation}%
where%
\begin{equation}
\left( \boldsymbol{X}_{N}\right) =\left( X_{N1,}\cdots ,X_{N\binom{2r}{N}%
}\right) ,\quad X_{Nj}\in \mathcal{F}_{N},\,1\leq j\leq \binom{2r}{N}.
\end{equation}
\end{definition}

\begin{definition}
Let $Z\in \mathcal{F}$ then it can

\begin{enumerate}
\item be represented by and identified with $\mathfrak{R}_{\mathcal{F}%
}\left( Z\right) ,$ which is an operator super vector that belongs to $%
\mathcal{F}^{2^{2r}}$ i.e. 
\begin{equation}
Z\leftrightarrow \mathfrak{R}_{\mathcal{F}}\left( Z\right) =\left( 
\begin{array}{c}
Z_{0} \\ 
\vdots \\ 
Z_{N} \\ 
\vdots \\ 
Z_{2r}%
\end{array}%
\right) ,\,Z_{N}\in \mathcal{F}_{N},\,0\leq N\leq 2r.
\end{equation}

\item In general super vectors belonging to $\mathcal{F}^{2^{2r}}$ are of
the form 
\begin{equation}
\left( 
\begin{array}{c}
\alpha \\ 
\vdots \\ 
\mu \\ 
\vdots \\ 
\varrho%
\end{array}%
\right)
\end{equation}%
where all super coordinates belong to $\mathcal{F}.$

\item The scalar vector representations $\mathcal{R}_{X_{N}}\left( Z\right) $
of $Z_{N}$ wrt the basis $\left( \boldsymbol{X}_{N}\right) $ are produced by 
\begin{equation}
Z_{N}=\left( X_{N1,}\cdots ,X_{N\binom{2r}{N}}\right) \mathcal{R}%
_{X_{N}}\left( Z\right) =\left( X_{N1,}\cdots ,X_{N\binom{2r}{N}}\right)
\left( 
\begin{array}{c}
\mathcal{R}_{X_{N}}\left( Z\right) _{0} \\ 
\vdots \\ 
\mathcal{R}_{X_{N}}\left( Z\right) _{\binom{2r}{N}}%
\end{array}%
\right) ,\,0\leq N\leq 2^{2r}
\end{equation}%
and 
\begin{equation}
\,\mathcal{R}_{X_{N}}\left( Z\right) _{j}\in \mathbb{C},\,\,1\leq j\leq 
\binom{2r}{N}.
\end{equation}
\end{enumerate}
\end{definition}

In summary

\begin{definition}
The scalar vector representation $\mathcal{R}_{\mathcal{F}}\left( Z\right) $
of elements $Z\in \mathcal{F}$ wrt a basis $\left( \boldsymbol{X}\right)
_{S} $ of $\mathcal{F}$ is produced by 
\begin{equation}
Z\leftrightarrow \mathcal{R}_{\mathcal{F}}\left( Z\right) =\left( 
\begin{array}{c}
\left( I_{\mathcal{F}}\right) \mathcal{R}_{X_{0}}\left( Z\right) _{1} \\ 
\vdots \\ 
\left( X_{N1,}\cdots ,X_{N\binom{2r}{N}}\right) \left( 
\begin{array}{c}
\mathcal{R}_{X_{N}}\left( Z\right) _{1} \\ 
\vdots \\ 
\mathcal{R}_{X_{N}}\left( Z\right) _{\binom{2r}{N}}%
\end{array}%
\right) \\ 
\vdots \\ 
\left( X_{2r}\right) \mathcal{R}_{X_{2r}}\left( Z\right) _{1}%
\end{array}%
\right) ,\,Z_{N}\in \mathcal{F}_{N},\,0\leq N\leq 2r.
\end{equation}
\end{definition}

\begin{remark}
Note that 
\begin{equation}
\dim \left\{ \mathcal{F}_{0}\right\} =\dim \left\{ \mathcal{F}_{2r}\right\}
=1
\end{equation}%
thus there is only one coordinate associated with the representations $%
\mathcal{R}_{X_{0}}$ and $\mathcal{R}_{X_{2r}}$ .
\end{remark}

\begin{enumerate}
\item Operator super vectors 
\begin{equation}
\mathfrak{R}_{\mathcal{F}}\left( Z\right) \in \mathcal{M}\left( 2r+1,1,%
\mathcal{F}\right) .
\end{equation}

\item Scalar vectors 
\begin{equation}
\mathcal{R}_{\mathcal{F}}\left( Z\right) \in \mathcal{M}\left( 2^{2r},1,%
\mathbb{C}\right) .
\end{equation}
\end{enumerate}

\begin{remark}
Expansion of operators $Z\in \mathcal{F}$ in terms of super vectors 
\begin{equation}
Z=\sum_{0\leq N\leq 2r}\mathbf{X}_{N}\mathfrak{R}_{\boldsymbol{X}_{S}}\left(
Z\right) _{N}.
\end{equation}
\end{remark}

\begin{remark}
Each operator $Z_{N}\in \mathcal{F}_{N}\subset \mathcal{F}$ can be expanded
as 
\begin{equation}
Z_{N}=Z_{N}\left( X_{N1,}\cdots ,X_{N\binom{2r}{N}}\right) =\left(
X_{N1,}\cdots ,X_{N\binom{2r}{N}}\right) \mathcal{R}_{X_{N}}\left( Z\right)
_{N}=\left( X_{N1,}\cdots ,X_{N\binom{2r}{N}}\right) \left( 
\begin{array}{c}
\mathcal{R}_{X_{N}}\left( Z\right) _{1} \\ 
\vdots \\ 
\mathcal{R}_{X_{N}}\left( Z\right) _{\binom{2r}{N}}%
\end{array}%
\right) ,\quad \mathcal{R}_{X_{N}}\left( Z\right) _{j}\in \mathbb{C},\,1\leq
j\leq \binom{2r}{N}.
\end{equation}
\end{remark}

\begin{remark}
\begin{enumerate}
\item 
\begin{enumerate}
\item The $N^{th}$-order exterior power vector space $\mathcal{F}_{N}$ is
defined by 
\begin{equation}
\mathcal{F}_{N}=\dbigwedge\nolimits^{N}\mathcal{F}_{1}\ 
\end{equation}

\item The exterior algebra form of the Clifford algebra $\mathcal{F}$ is
given by 
\begin{equation}
\mathcal{F}=\dbigwedge \left( \mathcal{F}_{1}\right) \
=\dbigoplus\limits_{0\leq N\leq 2r}\dbigwedge\nolimits^{N}\mathcal{F}_{1}
\end{equation}

\item The super vector representations $\mathfrak{R}_{\mathcal{F}}\left\{ 
\mathcal{F}\right\} $ are column super vectors 
\begin{equation}
\mathfrak{R}_{\mathcal{F}}\left\{ \mathcal{F}\right\} \subset \mathcal{M}%
\left( 2r+1,1,\mathcal{F}\right)
\end{equation}%
explicitly denoted by

\item 
\begin{equation}
\mathfrak{R}_{\mathcal{F}}\left\{ Z\right\} =\left( 
\begin{array}{c}
Z_{0} \\ 
\vdots \\ 
Z_{2r}%
\end{array}%
\right) ,\,Z_{N}\in \mathcal{F}_{N},\,0\leq N\leq 2r.
\end{equation}

\item 
\begin{equation}
\mathfrak{R}_{\mathcal{F}}\left\{ Z\right\} \subset \mathcal{F}^{2^{2r}}.
\end{equation}

\item Homogeneous $N^{th}$ order operators $Z_{N}$ can be expanded as 
\begin{equation}
Z_{N}=\left( \boldsymbol{X}_{N}\right) \mathcal{R}_{\mathcal{F}_{N}}\left(
Z_{N}\right) =\left( 
\begin{array}{ccc}
X_{N1} & \cdots & X_{N\binom{2r}{N}}%
\end{array}%
\right) \left( 
\begin{array}{c}
\mathcal{R}_{\mathcal{X}_{N}}\left( Z_{N}\right) _{1} \\ 
\vdots \\ 
\mathcal{R}_{\mathcal{X}_{N}}\left( Z_{N}\right) _{\binom{2r}{N}}%
\end{array}%
\right) ,\quad \mathcal{R}_{\mathcal{X}_{N}}\left( Z_{N}\right) _{j}\in 
\mathbb{C},\,0\leq N\leq 2r,\,1\leq j\leq \binom{2r}{N}.
\end{equation}

\item $\mathcal{F}_{1}$ super vectors are of the form 
\begin{equation}
\left( 
\begin{array}{c}
\boldsymbol{a}^{\dagger } \\ 
\boldsymbol{a}%
\end{array}%
\right) \in \mathcal{M}\left( 2r,1,\mathcal{F}_{2}\right) ,
\end{equation}%
which is the\textbf{\ Nambu Representation} of creation and annihilation
operators.

\item Super matrices associated the CAR are 
\begin{equation}
\left( 
\begin{array}{cc}
\boldsymbol{a}^{\dagger }\boldsymbol{a}^{\dagger } & \boldsymbol{a}^{\dagger
}\boldsymbol{a} \\ 
\boldsymbol{aa}^{\dagger } & \boldsymbol{aa}%
\end{array}%
\right) \in \mathcal{M}\left( 2r,2r,\mathcal{F}_{2}\right) \subset \mathcal{M%
}\left( 2^{2r},2^{2r},\mathcal{F}\right) ,
\end{equation}%
and 
\begin{equation}
\left( 
\begin{array}{cc}
\boldsymbol{a}^{\dagger }\boldsymbol{a}^{\dagger } & \boldsymbol{aa}%
^{\dagger } \\ 
\boldsymbol{a}^{\dagger }\boldsymbol{a} & \boldsymbol{aa}%
\end{array}%
\right) \in \mathcal{M}\left( 2r,2r,\mathcal{F}_{2}\right) \subset \mathcal{M%
}\left( 2^{2r},2^{2r},\mathcal{F}\right) ,
\end{equation}%
\textbf{but note}

\item 
\begin{equation}
\left( 
\begin{array}{cc}
\boldsymbol{a}^{\dagger }\boldsymbol{a}^{\dagger } & \boldsymbol{a}^{\dagger
}\boldsymbol{a} \\ 
\boldsymbol{aa}^{\dagger } & \boldsymbol{aa}%
\end{array}%
\right) \left( 
\begin{array}{c}
\boldsymbol{a}^{\dagger } \\ 
\boldsymbol{a}%
\end{array}%
\right) \notin \mathcal{M}\left( 2r,1,\mathcal{F}_{1}\right) .
\end{equation}
\end{enumerate}
\end{enumerate}
\end{remark}

\subsection{Representations of Bogoliubov Transformations \label%
{Representations of Bogoliubov Transformations}}

\begin{remark}
The Bogoliubov transformation $B\in \mathcal{L}\left( \mathcal{F}_{1}\right) 
$ 
\begin{equation}
B\left( 
\begin{array}{cc}
\boldsymbol{a}^{\dagger } & \boldsymbol{a}%
\end{array}%
\right) =\left( 
\begin{array}{cc}
\boldsymbol{b}^{\dagger } & \boldsymbol{b}%
\end{array}%
\right) =\left( 
\begin{array}{cc}
\boldsymbol{a}^{\dagger } & \boldsymbol{a}%
\end{array}%
\right) \mathcal{R}_{a}\left( B\right) =\left( 
\begin{array}{cc}
\boldsymbol{a}^{\dagger } & \boldsymbol{a}%
\end{array}%
\right) \left( 
\begin{array}{cc}
\mathcal{R}_{a}\left( B\right) _{11} & \mathcal{R}_{a}\left( B\right) _{12}
\\ 
\mathcal{R}_{a}\left( B\right) _{21} & \mathcal{R}_{a}\left( B\right) _{22}%
\end{array}%
\right) =\left( 
\begin{array}{cc}
\boldsymbol{a}^{\dagger } & \boldsymbol{a}%
\end{array}%
\right) \left( 
\begin{array}{cc}
\mathcal{R}_{a}\left( B\right) _{11} & \overline{\mathcal{R}_{a}\left(
B\right) _{21}} \\ 
\mathcal{R}_{a}\left( B\right) _{21} & \overline{\mathcal{R}_{a}\left(
B\right) _{11}}%
\end{array}%
\right) =\left( 
\begin{array}{cc}
\boldsymbol{a}^{\dagger } & \boldsymbol{a}%
\end{array}%
\right) \left( 
\begin{array}{cc}
\mathbb{U} & \overline{\mathbb{V}} \\ 
\mathbb{V} & \overline{\mathbb{U}}%
\end{array}%
\right) ,
\end{equation}%
where $\mathcal{R}_{a}\left( B\right) \in \mathcal{M}\left( 2r,2r,\mathbb{C}%
\right) .$
\end{remark}

\begin{remark}
Using 
\begin{equation}
\left( 
\begin{array}{cc}
\boldsymbol{x}_{1} & \boldsymbol{x}_{2}%
\end{array}%
\right) =K\left( 
\begin{array}{cc}
\boldsymbol{a}^{\dagger } & \boldsymbol{a}%
\end{array}%
\right) =\left( 
\begin{array}{cc}
\boldsymbol{a}^{\dagger } & \boldsymbol{a}%
\end{array}%
\right) \mathcal{R}_{a}\left( K\right) =\left( 
\begin{array}{cc}
\boldsymbol{a}^{\dagger } & \boldsymbol{a}%
\end{array}%
\right) \mathbb{K}=\left( 
\begin{array}{cc}
\boldsymbol{a}^{\dagger } & \boldsymbol{a}%
\end{array}%
\right) \frac{1}{\sqrt{2}}\left( 
\begin{array}{cc}
\mathbb{I}_{r} & i\mathbb{I}_{r} \\ 
\mathbb{I}_{r} & -i\mathbb{I}_{r}%
\end{array}%
\right)
\end{equation}%
and 
\begin{equation}
\left( 
\begin{array}{cc}
\boldsymbol{a}^{\dagger } & \boldsymbol{a}%
\end{array}%
\right) =K^{\dagger }\left( 
\begin{array}{cc}
\boldsymbol{x}_{1} & \boldsymbol{x}_{2}%
\end{array}%
\right) =\left( 
\begin{array}{cc}
\boldsymbol{x}_{1} & \boldsymbol{x}_{2}%
\end{array}%
\right) \mathcal{R}_{x}\left( K\right) =\left( 
\begin{array}{cc}
\boldsymbol{x}_{1} & \boldsymbol{x}_{2}%
\end{array}%
\right) \mathbb{K}^{\dagger }=\left( 
\begin{array}{cc}
\boldsymbol{x}_{1} & \boldsymbol{x}_{2}%
\end{array}%
\right) \frac{1}{\sqrt{2}}\left( 
\begin{array}{cc}
\mathbb{I}_{r} & \mathbb{I}_{r} \\ 
-i\mathbb{I}_{r} & i\mathbb{I}_{r}%
\end{array}%
\right) .
\end{equation}
\end{remark}

\begin{remark}
One obtains the $\mathcal{R}_{x}$ representations of the Bogoliubov
transformation $B$%
\begin{eqnarray}
\left( 
\begin{array}{cc}
\boldsymbol{y}_{1} & \boldsymbol{y}_{2}%
\end{array}%
\right) &=&B\left( 
\begin{array}{cc}
\boldsymbol{x}_{1} & \boldsymbol{x}_{2}%
\end{array}%
\right) =\left( 
\begin{array}{cc}
\boldsymbol{x}_{1} & \boldsymbol{x}_{2}%
\end{array}%
\right) \mathcal{R}_{a}\left( K\right) ^{\dagger }\mathcal{R}_{a}\left(
B\right) \mathcal{R}_{a}\left( K\right) =\left( 
\begin{array}{cc}
\boldsymbol{x}_{1} & \boldsymbol{x}_{2}%
\end{array}%
\right) \mathcal{R}_{x}\left( B\right) \\
&=&\left( 
\begin{array}{cc}
\boldsymbol{x}_{1} & \boldsymbol{x}_{2}%
\end{array}%
\right) \left( 
\begin{array}{cc}
\func{Re}\left( \mathcal{R}_{a}\left( B\right) _{11}+\mathcal{R}_{a}\left(
B\right) _{22}\right) & -\func{Im}\left( \mathcal{R}_{a}\left( B\right)
_{11}+\mathcal{R}_{a}\left( B\right) _{22}\right) \\ 
\func{Im}\left( \mathcal{R}_{a}\left( B\right) _{11}-\mathcal{R}_{a}\left(
B\right) _{22}\right) & \func{Re}\left( \mathcal{R}_{a}\left( B\right) _{11}-%
\mathcal{R}_{a}\left( B\right) _{22}\right)%
\end{array}%
\right) =\left( 
\begin{array}{cc}
\boldsymbol{x}_{1} & \boldsymbol{x}_{2}%
\end{array}%
\right) \left( 
\begin{array}{cc}
\func{Re}\left( \mathbb{U+V}\right) & -\func{Im}\left( \mathbb{U+V}\right)
\\ 
\func{Im}\left( \mathbb{U-V}\right) & \func{Re}\left( \mathbb{U-V}\right)%
\end{array}%
\right) ,
\end{eqnarray}%
where $\mathcal{R}_{a}\left( B\right) $ and $\mathcal{R}_{x}\left( B\right) $
$\in \mathcal{M}\left( 2r,2r,\mathbb{C}\right) $ is given by 
\begin{equation}
\left( 
\begin{array}{cc}
\boldsymbol{b}^{\dagger } & \boldsymbol{b}%
\end{array}%
\right) =B\left( 
\begin{array}{cc}
\boldsymbol{a}^{\dagger } & \boldsymbol{a}%
\end{array}%
\right) =\left( 
\begin{array}{cc}
\boldsymbol{a}^{\dagger } & \boldsymbol{a}%
\end{array}%
\right) \mathcal{R}_{a}\left( B\right) =\left( 
\begin{array}{cc}
\boldsymbol{a}^{\dagger } & \boldsymbol{a}%
\end{array}%
\right) \left( 
\begin{array}{cc}
\mathbb{U} & \overline{\mathbb{V}} \\ 
\mathbb{V} & \overline{\mathbb{U}}%
\end{array}%
\right) .
\end{equation}
\end{remark}

\begin{remark}
The map $\Omega \in \mathcal{L}\left( \mathcal{M}\left( 2r,2r,\mathbb{C}%
\right) ,\mathcal{M}\left( 2r,2r,\mathbb{R}\right) \right) $ embeds the
space of real matrices into the space of complex matrices 
\begin{equation}
\Omega :\mathcal{M}\left( 2r,2r,\mathbb{R}\right) \rightarrow \mathcal{M}%
\left( 2r,2r,\mathbb{C}\right)
\end{equation}%
is defined by 
\begin{equation}
\Omega \left( \left( 
\begin{array}{cc}
\mathbb{A} & \mathbb{B} \\ 
\mathbb{C} & \mathbb{D}%
\end{array}%
\right) \right) =\frac{1}{2}\left( 
\begin{array}{cc}
\mathbb{A+D}-i\left( \mathbb{B-C}\right) & \mathbb{A-D}+i\left( \mathbb{B+C}%
\right) \\ 
\mathbb{A-D}-i\left( \mathbb{B+C}\right) & \mathbb{A+D}+i\left( \mathbb{B-C}%
\right)%
\end{array}%
\right)
\end{equation}%
and is a linear isomorphism.
\end{remark}

\begin{remark}
The matrix space $\mathcal{M}\left( 2r,2r,\mathbb{C}\right) $ can be
embedded by the matrix isomorphism $\digamma $ 
\begin{equation}
\digamma :\mathcal{M}\left( 2r,2r,\mathbb{C}\right) \rightarrow \mathcal{M}%
\left( 2r,2r,\mathcal{F}_{0}\right)
\end{equation}%
into the space of Clifford matrices $\mathcal{M}\left( 2r,2r,\mathcal{F}%
\right) \subset \mathcal{M}\left( 2^{2r},2^{2r},\mathcal{F}\right) $ by
identifying $\mathcal{M}\left( 2r,2r,\mathbb{C}\right) $ with $\mathcal{M}%
\left( 2r,2r,\mathcal{F}_{0}\right) $ by%
\begin{equation}
\digamma \left( \left( 
\begin{array}{cc}
\mathbb{M}_{11} & \mathbb{M}_{12} \\ 
\mathbb{M}_{21} & \mathbb{M}_{22}%
\end{array}%
\right) \right) =\left( 
\begin{array}{cc}
\mathbb{M}_{11} & \mathbb{M}_{12} \\ 
\mathbb{M}_{21} & \mathbb{M}_{22}%
\end{array}%
\right) \left( 
\begin{array}{cc}
I_{\mathcal{F}_{0}} & O_{r} \\ 
O_{r} & I_{\mathcal{F}_{0}}%
\end{array}%
\right) ,
\end{equation}%
where 
\begin{equation}
I_{\mathcal{F}_{0}}=\left( 
\begin{array}{cccc}
I_{\mathcal{F}} & 0 & \cdots & 0 \\ 
0 & \ddots & \cdots & \vdots \\ 
\vdots & \cdots & \ddots & 0 \\ 
0 & \cdots & 0 & I_{\mathcal{F}}%
\end{array}%
\right) \in \mathcal{M}\left( r,r,\mathcal{F}_{0}\right) .
\end{equation}%
This leads to the $\mathcal{M}\left( 2r,2r,\mathcal{F}_{0}\right) $ valued
representations $\mathfrak{R}_{\mathcal{F}}=\digamma \circ \mathcal{R}_{a}$
of Clifford Bogoliubov Transformations 
\begin{equation}
\digamma \left( \mathcal{R}_{a}\left( B\right) \right) =\left( 
\begin{array}{cc}
\mathcal{R}_{a}\left( B\right) _{11} & \mathcal{R}_{a}\left( B\right) _{12}
\\ 
\mathcal{R}_{a}\left( B\right) _{21} & \mathcal{R}_{a}\left( B\right) _{22}%
\end{array}%
\right) \left( 
\begin{array}{cc}
I_{\mathcal{F}_{0}} & O_{r} \\ 
O_{r} & I_{\mathcal{F}_{0}}%
\end{array}%
\right) =\left( 
\begin{array}{cc}
\mathbb{U} & \overline{\mathbb{V}} \\ 
\mathbb{V} & \overline{\mathbb{U}}%
\end{array}%
\right) \left( 
\begin{array}{cc}
I_{\mathcal{F}_{0}} & O_{r} \\ 
O_{r} & I_{\mathcal{F}_{0}}%
\end{array}%
\right) =\left( 
\begin{array}{cc}
\mathbb{U}I_{\mathcal{F}_{0}} & \overline{\mathbb{V}}I_{\mathcal{F}_{0}} \\ 
\mathbb{V}I_{\mathcal{F}_{0}} & \overline{\mathbb{U}}I_{\mathcal{F}_{0}}%
\end{array}%
\right) .
\end{equation}
\end{remark}

\begin{remark}
In terms of Clifford matrices the Bogoliubov transformation $B\in \mathcal{L}%
\left( \mathcal{F}_{1}\right) $ can be expressed as 
\begin{equation}
B\left( 
\begin{array}{cc}
\boldsymbol{a}^{\dagger } & \boldsymbol{a}%
\end{array}%
\right) =\left( 
\begin{array}{cc}
\boldsymbol{b}^{\dagger } & \boldsymbol{b}%
\end{array}%
\right) =\left( 
\begin{array}{cc}
\boldsymbol{a}^{\dagger } & \boldsymbol{a}%
\end{array}%
\right) \digamma \left( \mathcal{R}_{a}\left( B\right) \right) =\left( 
\begin{array}{cc}
\boldsymbol{a}^{\dagger } & \boldsymbol{a}%
\end{array}%
\right) \left( 
\begin{array}{cc}
\mathcal{R}_{a}\left( B\right) _{11}I_{\mathcal{F}_{0}} & \mathcal{R}%
_{a}\left( B\right) _{12}I_{\mathcal{F}_{0}} \\ 
\mathcal{R}_{a}\left( B\right) _{21}I_{\mathcal{F}_{0}} & \mathcal{R}%
_{a}\left( B\right) _{22}I_{\mathcal{F}_{0}}%
\end{array}%
\right) =\left( 
\begin{array}{cc}
\boldsymbol{a}^{\dagger } & \boldsymbol{a}%
\end{array}%
\right) \left( 
\begin{array}{cc}
\mathbb{U}I_{\mathcal{F}_{0}} & \overline{\mathbb{V}}I_{\mathcal{F}_{0}} \\ 
\mathbb{V}I_{\mathcal{F}_{0}} & \overline{\mathbb{U}}I_{\mathcal{F}_{0}}%
\end{array}%
\right) ,
\end{equation}%
where $\mathcal{R}_{a}\left( B\right) \in \mathcal{M}\left( 2r,2r,\mathbb{C}%
\right) .$
\end{remark}

\begin{remark}
\begin{enumerate}
\item Super operator representations $\equiv $ Clifford representations $%
\mathfrak{R}_{\boldsymbol{X}_{S}}\left( \Upsilon \right) \in \mathcal{M}%
\left( 2r+1,2r+1,\mathcal{F}\right) $ are produced by 
\begin{equation}
\Upsilon \left( \boldsymbol{X}_{0},\cdots ,\boldsymbol{X}_{N},\cdots ,%
\boldsymbol{X}_{2r}\right) =\left( \boldsymbol{X}_{0},\cdots ,\boldsymbol{X}%
_{N},\cdots ,\boldsymbol{X}_{2r}\right) \mathfrak{R}_{\boldsymbol{X}%
_{S}}\left( \Upsilon \right) =\left( \boldsymbol{X}_{0},\cdots ,\boldsymbol{X%
}_{N},\cdots ,\boldsymbol{X}_{2r}\right) \left( 
\begin{array}{ccc}
\mathfrak{R}_{\boldsymbol{X}_{S}}\left( \Upsilon \right) _{00} & \cdots & 
\mathfrak{R}_{\boldsymbol{X}_{S}}\left( \Upsilon \right) _{02r} \\ 
\vdots & \ddots & \vdots \\ 
\mathfrak{R}_{\boldsymbol{X}_{S}}\left( \Upsilon \right) _{2r0} & \cdots & 
\mathfrak{R}_{\boldsymbol{X}_{S}}\left( \Upsilon \right) _{2r2r}%
\end{array}%
\right) .
\end{equation}

\item Complex representations $\mathcal{R}_{X}\left( \Upsilon \right) \in 
\mathcal{M}\left( 2^{2r},2^{2r},\mathbb{C}\right) $ are produced by 
\begin{eqnarray}
\Upsilon \left( I_{\mathcal{F}},\cdots ,\left( X_{N1,}\cdots ,X_{N\binom{2r}{%
N}}\right) ,\cdots ,X_{2r1,}\right) &=&\left( I_{\mathcal{F}},\cdots ,\left(
X_{N1,}\cdots ,X_{N\binom{2r}{N}}\right) ,\cdots ,X_{2r}\right) \mathcal{R}%
_{X}\left( \Upsilon \right) \\
&=&\left( I_{\mathcal{F}},\cdots ,\left( X_{N1,}\cdots ,X_{N\binom{2r}{N}%
}\right) ,\cdots ,X_{2r}\right) \left( 
\begin{array}{ccc}
\mathcal{R}_{X}\left( \Upsilon \right) _{00} & \cdots & \mathcal{R}%
_{X}\left( \Upsilon \right) _{02^{2r}} \\ 
\vdots & \ddots & \vdots \\ 
\mathcal{R}_{X}\left( \Upsilon \right) _{2^{2r}0} & \cdots & \mathcal{R}%
_{X}\left( \Upsilon \right) _{2^{2r}2^{2r}}%
\end{array}%
\right) .
\end{eqnarray}
\end{enumerate}
\end{remark}

\begin{remark}
The action of operators $\Upsilon \in \mathcal{L}\left( \mathcal{F}\right) $
on elements $Z\in \mathcal{F}$ is given by

\begin{enumerate}
\item The Clifford matrix $\mathcal{M}\left( 2r+1,2r+1,\mathcal{F}\right) $
representation of the transformation $\Upsilon \in \mathcal{L}\left( 
\mathcal{F}\right) $ 
\begin{gather}
\Upsilon Z\left( \boldsymbol{X}_{0},\cdots ,\boldsymbol{X}_{N},\cdots ,%
\boldsymbol{X}_{2r}\right) =\left( \boldsymbol{X}_{0},\cdots ,\boldsymbol{X}%
_{N},\cdots ,\boldsymbol{X}_{2r}\right) \mathfrak{R}_{\boldsymbol{X}%
_{S}}\left( \Upsilon \right) \mathfrak{R}_{\boldsymbol{X}_{S}}\left( Z\right)
\\
=\left( \boldsymbol{X}_{0},\cdots ,\boldsymbol{X}_{N},\cdots ,\boldsymbol{X}%
_{2r}\right) \left( 
\begin{array}{ccc}
\mathfrak{R}_{\boldsymbol{X}_{S}}\left( \Upsilon \right) _{00} & \cdots & 
\mathfrak{R}_{\boldsymbol{X}_{S}}\left( \Upsilon \right) _{02r} \\ 
\vdots & \ddots & \vdots \\ 
\mathfrak{R}_{\boldsymbol{X}_{S}}\left( \Upsilon \right) _{2r0} & \cdots & 
\mathfrak{R}_{\boldsymbol{X}_{S}}\left( \Upsilon \right) _{2r2r}%
\end{array}%
\right) \left( 
\begin{array}{c}
\mathfrak{R}_{X_{S}}\left( Z\right) _{0} \\ 
\vdots \\ 
\mathfrak{R}_{X_{S}}\left( Z\right) _{N} \\ 
\vdots \\ 
\mathfrak{R}_{X_{S}}\left( Z\right) _{2r}%
\end{array}%
\right) =\left( \boldsymbol{X}_{0},\cdots ,\boldsymbol{X}_{N},\cdots ,%
\boldsymbol{X}_{2r}\right) \left( 
\begin{array}{ccc}
\mathfrak{R}_{\boldsymbol{X}_{S}}\left( \Upsilon \right) _{00} & \cdots & 
\mathfrak{R}_{\boldsymbol{X}_{S}}\left( \Upsilon \right) _{02r} \\ 
\vdots & \ddots & \vdots \\ 
\mathfrak{R}_{\boldsymbol{X}_{S}}\left( \Upsilon \right) _{2r0} & \cdots & 
\mathfrak{R}_{\boldsymbol{X}_{S}}\left( \Upsilon \right) _{2r2r}%
\end{array}%
\right) \left( 
\begin{array}{c}
Z_{0} \\ 
\vdots \\ 
Z_{N} \\ 
\vdots \\ 
Z_{2r}%
\end{array}%
\right)
\end{gather}

\item The complex matrix $\mathcal{M}\left( 2^{2r},2^{2r},\mathbb{C}\right) $
representation of the transformation $\Upsilon \in \mathcal{L}\left( 
\mathcal{F}\right) $ is given by 
\begin{gather}
\Upsilon Z\left( I_{\mathcal{F}},\cdots ,\left( X_{N1,}\cdots ,X_{N\binom{2r%
}{N}}\right) ,\cdots ,X_{2r1,}\right) =\left( I_{\mathcal{F}},\cdots ,\left(
X_{N1,}\cdots ,X_{N\binom{2r}{N}}\right) ,\cdots ,X_{2r}\right) \mathcal{R}%
_{X}\left( \Upsilon \right) \mathcal{R}_{X}\left( Z\right) \\
=\left( I_{\mathcal{F}},\cdots ,\left( X_{N1,}\cdots ,X_{N\binom{2r}{N}%
}\right) ,\cdots ,X_{2r}\right) \left( 
\begin{array}{ccc}
\mathcal{R}_{X}\left( \Upsilon \right) _{00} & \cdots & \mathcal{R}%
_{X}\left( \Upsilon \right) _{02^{2r}} \\ 
\vdots & \ddots & \vdots \\ 
\mathcal{R}_{X}\left( \Upsilon \right) _{2^{2r}0} & \cdots & \mathcal{R}%
_{X}\left( \Upsilon \right) _{2^{2r}2^{2r}}%
\end{array}%
\right) \left( 
\begin{array}{c}
\mathcal{R}_{X_{0}}\left( Z_{0}\right) \\ 
\vdots \\ 
\left( 
\begin{array}{c}
\mathcal{R}_{X_{N}}\left( Z\right) _{0} \\ 
\vdots \\ 
\mathcal{R}_{X_{N}}\left( Z\right) _{\binom{2r}{N}}%
\end{array}%
\right) \\ 
\vdots \\ 
\mathcal{R}_{X_{2r}}\left( Z_{2r}\right)%
\end{array}%
\right) .
\end{gather}
\end{enumerate}
\end{remark}

\subsubsection{Extensions of Bogoliubov Transformations \label{Extensions of
Bogoliubov Transformations}}

The transformation $B\in \mathcal{L}\left( \mathcal{F}_{1}\right) $ induces
a transformation $\boxtimes \left( B\right) \in \mathcal{L}\left( \mathcal{F}%
\right) $ by 
\begin{gather}
\boxtimes \left( B\right) \left( {\Huge \wedge }\left( \overset{}{%
\boldsymbol{x}}\right) \right) =\boxtimes \left( B\right) \left( {\Huge %
\wedge }\left( 
\begin{array}{cc}
\boldsymbol{x}_{1} & \boldsymbol{x}_{2}%
\end{array}%
\right) \right) =\boxtimes \left( \mathcal{B}\right) \left( \overset{}{%
\boldsymbol{X}_{0},\cdots ,\boldsymbol{X}_{N},\cdots ,\boldsymbol{X}_{2r}}%
\right) =\boxtimes \left( \mathcal{B}\right) \left( \overset{}{I_{\mathcal{F}%
},\cdots ,x_{j_{1}}\wedge \cdots \wedge x_{j_{N}},\cdots ,x_{j_{1}}\wedge
\cdots \wedge x_{j_{2r}}}\right) \\
=\left( \overset{}{\boxtimes ^{0}\left( B\right) I_{\mathcal{F}},\cdots
,\boxtimes ^{N}\left( B\right) \left( x_{j_{1}}\wedge \cdots \wedge
x_{j_{N}}\right) ,\cdots ,\boxtimes ^{2r}\left( B\right) \left(
x_{j_{1}}\wedge \cdots \wedge x_{j_{2r}}\right) }\right) =\left( \overset{}{%
B_{0}I_{\mathcal{F}},\cdots ,Bx_{j_{1}}\wedge \cdots \wedge
Bx_{j_{N}},\cdots ,Bx_{j_{1}}\wedge \cdots \wedge Bx_{j_{2r}}}\right) \\
1\leq j_{1}<\cdots <j_{N}\leq 2r,0\leq N\leq 2r.
\end{gather}

\begin{remark}
$\boxtimes \left( B\right) \in $ $\mathcal{L}\left( \mathcal{F}\right) $ is
a block diagonal transformation i.e. 
\begin{equation}
\boxtimes \left( B\right) :{\Huge \wedge }^{N}\mathcal{F}_{1}\rightarrow 
{\Huge \wedge }^{N}\mathcal{F}_{1},
\end{equation}%
which means that%
\begin{equation}
\left[ \boxtimes \left( B\right) ,P_{\mathcal{F}_{N}}\right] =0
\end{equation}%
is\textbf{\ gauge invariant}.
\end{remark}

\begin{remark}
The super matrix representation of the induced transformation $\boxtimes
\left( B\right) $ are Clifford matrices $\mathfrak{R}_{\boldsymbol{X}%
_{S}}\left( \boxtimes \left( B\right) \right) $ 
\begin{equation}
\mathfrak{R}_{\boldsymbol{X}_{S}}\left( \boxtimes \left( B\right) \right)
\in \mathcal{M}\left( 2r+1,2r+1,\mathcal{F}_{0}\right) .
\end{equation}%
General Clifford matrices $\mathfrak{M}$ that belong to $\mathcal{M}\left(
2r+1,2r+1,\mathcal{F}\right) $ act on Clifford vectors $\mathfrak{z}$
belonging to $\mathcal{M}\left( 2r+1,1,\mathcal{F}\right) $ to produce
Clifford vectors $\mathfrak{w}$ belonging to $\mathcal{M}\left( 2r+1,1,%
\mathcal{F}\right) $ i.e. 
\begin{equation}
\mathfrak{Mz=w}
\end{equation}%
are not gauge invariant i.e. they mix grades.
\end{remark}

\section{Complex Structures}

\subsection{Concrete Example}

\subsubsection{Complex Hilbert Space}

Consider a \emph{complex} Hilbert space $\mathfrak{H}$ of dimension $r$, the
inner product $\left\langle \left. {}\right\vert \right\rangle $ in $%
\mathfrak{H}$ can be decomposed in terms of \emph{real} \emph{symmetric} and 
\emph{real} \emph{antisymmetric} \emph{real} bilinear forms $h$ and $g$ 
\begin{equation}
\left\langle \left. x\right\vert y\right\rangle =h\left( x,y\right)
+ig\left( x,y\right) =\func{Re}\left\langle \left. x\right\vert
y\right\rangle +i\func{Im}\left\langle \left. x\right\vert y\right\rangle ,
\end{equation}%
where 
\begin{align}
h,g& :\mathfrak{H}\rightarrow \mathbb{R} \\
h\left( x,y\right) & =\func{Re}\left\langle \left. x\right\vert
y\right\rangle =\frac{1}{2}\left\{ \left\langle \left. x\right\vert
y\right\rangle +\overline{\left\langle \left. x\right\vert y\right\rangle }%
\right\} =\frac{1}{2}\left\{ \left\langle \left. x\right\vert y\right\rangle
+\left\langle \left. y\right\vert x\right\rangle \right\} =h\left(
y,x\right) \in \mathbb{R} \\
g\left( x,y\right) & =\func{Im}\left\langle \left. x\right\vert
y\right\rangle =-\frac{i}{2}\left\{ \left\langle \left. x\right\vert
y\right\rangle -\overline{\left\langle \left. x\right\vert y\right\rangle }%
\right\} =-\frac{i}{2}\left\{ \left\langle \left. x\right\vert
y\right\rangle -\left\langle \left. y\right\vert x\right\rangle \right\}
=-g\left( y,x\right) \in \mathbb{R}.
\end{align}%
$\lceil $ If $h$ were a complex symmetric bilinear form then 
\begin{equation}
h\left( \alpha x,y\right) =\alpha h\left( x,y\right) =h\left( y,\alpha
x\right) =\alpha h\left( y,x\right)
\end{equation}%
however 
\begin{equation}
h\left( \alpha x,y\right) =\frac{1}{2}\left\{ \left\langle \left. \alpha
x\right\vert y\right\rangle +\left\langle \left. y\right\vert \alpha
x\right\rangle \right\} =\frac{1}{2}\left\{ \overline{\alpha }\left\langle
\left. x\right\vert y\right\rangle +\alpha \left\langle \left. y\right\vert
x\right\rangle \right\} \neq \frac{1}{2}\alpha \left\{ \left\langle \left.
x\right\vert y\right\rangle +\left\langle \left. y\right\vert x\right\rangle
\right\}
\end{equation}%
unless 
\begin{equation}
\alpha =\overline{\alpha },
\end{equation}%
thus it is real symmetric.$\rfloor $

These real symmetric and antisymmetric forms determine each other and have
the properties

\begin{enumerate}
\item 
\begin{align}
g\left( x,y\right) & =h\left( ix,y\right) =h\left( y,ix\right) =-g\left(
y,x\right) \\
& \Rightarrow g\left( x,y\right) =-h\left( x,iy\right) .
\end{align}

\item 
\begin{align}
h\left( x,y\right) & =g\left( x,iy\right) \\
h\left( x,y\right) & =h\left( ix,iy\right) \\
g\left( x,y\right) & =g\left( ix,iy\right) .
\end{align}

\item 
\begin{align}
\left\langle \left. x\right\vert y\right\rangle & =\left\langle \left.
ix\right\vert iy\right\rangle \\
\left\langle \left. x\right\vert y\right\rangle & =h\left( x,y\right)
+ig\left( x,y\right)
\end{align}%
and thus in particular 
\begin{equation}
\left\langle \left. x\right\vert y\right\rangle =h\left( x,y\right)
+ih\left( ix,y\right) =g\left( x,iy\right) +ig\left( ix,iy\right) =h\left(
x,y\right) +ig\left( x,y\right) .
\end{equation}
\end{enumerate}

On the basis $\left\{ \left. e_{j},ie_{j}\right\vert 1\leq j\leq r\right\} $
of the real space $\mathfrak{H}_{e\mathbb{R}}$, where 
\begin{equation}
\mathfrak{H}_{e\mathbb{R}}=\mathbb{R}-\limfunc{linspan}\left\{ \left.
e_{j},ie_{j}\right\vert 1\leq j\leq r\right\} ,
\end{equation}%
the values of $h$ and $g$ are 
\begin{align}
h\left( e_{j},e_{k}\right) & =\frac{1}{2}\left\{ \left\langle \left.
e_{j}\right\vert e_{k}\right\rangle +\left\langle \left. e_{k}\right\vert
e_{j}\right\rangle \right\} =\frac{1}{2}2\delta _{jk}=\delta _{jk} \\
h\left( e_{j},ie_{k}\right) & =\frac{1}{2}\left\{ \left\langle \left.
e_{j}\right\vert ie_{k}\right\rangle +\left\langle \left. ie_{k}\right\vert
e_{j}\right\rangle \right\} =\frac{1}{2}\left\{ i\delta _{jk}-i\delta
_{jk}\right\} =0 \\
h\left( ie_{j},ie_{k}\right) & =\frac{1}{2}\left\{ \left\langle \left.
ie_{j}\right\vert ie_{k}\right\rangle +\left\langle \left. ie_{k}\right\vert
ie_{j}\right\rangle \right\} =\frac{1}{2}\left\{ -i^{2}\delta
_{jk}-i^{2}\delta _{jk}\right\} =\delta _{jk}
\end{align}%
and 
\begin{align}
g\left( e_{j},e_{k}\right) & =h\left( ie_{j},e_{k}\right) =\frac{1}{2}%
\left\{ \left\langle \left. ie_{j}\right\vert e_{k}\right\rangle
+\left\langle \left. e_{k}\right\vert ie_{j}\right\rangle \right\} =0 \\
g\left( e_{j},ie_{k}\right) & =h\left( ie_{j},ie_{k}\right) =\frac{1}{2}%
\left\{ \left\langle \left. ie_{j}\right\vert ie_{k}\right\rangle
+\left\langle \left. ie_{k}\right\vert ie_{j}\right\rangle \right\} =\delta
_{jk} \\
g\left( ie_{j},e_{k}\right) & =h\left( -e_{j},e_{k}\right) =\frac{1}{2}%
\left\{ \left\langle \left. -e_{j}\right\vert e_{k}\right\rangle
+\left\langle \left. e_{k}\right\vert -e_{j}\right\rangle \right\} =-\delta
_{jk} \\
g\left( ie_{j},ie_{k}\right) & =h\left( -e_{j},ie_{k}\right) =\frac{1}{2}%
\left\{ \left\langle \left. -e_{j}\right\vert ie_{k}\right\rangle
+\left\langle \left. ie_{k}\right\vert -e_{j}\right\rangle \right\} =0.
\end{align}%
Hence the action of $h$ in $\mathfrak{H}$ induces a \emph{real positive
definite inner product} in $\mathfrak{H}_{e\mathbb{R}}$ that is given by the
values of $h$ in $\mathfrak{H}$ and the action of $g$ induces a \emph{real
antisymetric bilinear form} that is given by the values of $g$ in $\mathfrak{%
H}$, so both $h$ and $g$ are defined on \emph{both} $\mathfrak{H}$ and $%
\mathfrak{H}_{e\mathbb{R}}$ as \emph{real bilinear forms}.

\begin{notation}
We denote $\left. h\right\vert _{\mathfrak{H\times H}}$ and $\left.
h\right\vert _{\mathfrak{H}_{e\mathbb{R}}\times \mathfrak{H}_{e\mathbb{R}}}$
by $h$ and $\left. g\right\vert _{\mathfrak{H\times H}}$ and $\left.
g\right\vert _{\mathfrak{H}_{e\mathbb{R}}\times \mathfrak{H}_{e\mathbb{R}}}$
by $g.$
\end{notation}

\begin{example}
Example of positive definiteness of the real symmetric inner product $h$ in $%
\mathfrak{H}$ and thus $\mathfrak{H}_{e\mathbb{R}}$ is 
\begin{align}
& h\left( i\left( 
\begin{array}{c}
1 \\ 
0%
\end{array}%
\right) ,i\left( 
\begin{array}{c}
1 \\ 
0%
\end{array}%
\right) \right) \overset{def}{=}\frac{1}{2}\left\{ \left\langle \left.
i\left( 
\begin{array}{c}
1 \\ 
0%
\end{array}%
\right) \right\vert i\left( 
\begin{array}{c}
1 \\ 
0%
\end{array}%
\right) \right\rangle +\left\langle \left. i\left( 
\begin{array}{c}
1 \\ 
0%
\end{array}%
\right) \right\vert i\left( 
\begin{array}{c}
1 \\ 
0%
\end{array}%
\right) \right\rangle \right\} \\
& =\overline{i\left( 
\begin{array}{cc}
1 & 0%
\end{array}%
\right) }i\left( 
\begin{array}{c}
1 \\ 
0%
\end{array}%
\right) =-ii=1,
\end{align}%
but note that this is not equal to a non positive definite complex symmetrc
inner product 
\begin{equation}
\left( \left. {}\right\vert \right) _{\mathfrak{H}}:\mathfrak{H\times H}%
\rightarrow \mathbb{C}
\end{equation}%
on these elements.
\end{example}

The complex sesquilinear bilinear form $\left\langle \left. {}\right\vert
\right\rangle _{\mathfrak{H}}$ that is the Hilbert space inner product is
defined by 
\begin{eqnarray}
\left\langle \left. v\right\vert v\right\rangle &\geq &0 \\
\left\langle \left. v\right\vert w\right\rangle &=&\overline{\left\langle
\left. w\right\vert v\right\rangle } \\
\left\langle \left. v\right\vert \alpha _{1}w_{1}+\alpha
_{2}w_{2}\right\rangle &=&\alpha _{1}\left\langle \left. v\right\vert
w_{1}\right\rangle +\alpha _{2}\left\langle \left. v\right\vert
w_{2}\right\rangle \\
\left\langle \alpha _{1}v_{1}+\alpha _{2}v_{2}\left. w\right\vert
\right\rangle &=&\overline{\alpha }_{1}\left\langle \left. v_{1}\right\vert
w\right\rangle +\overline{\alpha }_{2}\left\langle \left. v_{2}\right\vert
w\right\rangle
\end{eqnarray}%
and the complex symmetric bilinear form $\left( \left. {}\right\vert \right)
_{\mathfrak{H}}$ by 
\begin{eqnarray}
\left( \left. v\right\vert w\right) _{\mathfrak{H}} &=&\left( \left.
w\right\vert v\right) _{\mathfrak{H}} \\
\left( \left. v\right\vert \alpha _{1}w_{1}+\alpha _{2}w_{2}\right) _{%
\mathfrak{H}} &=&\alpha _{1}\left( \left. v\right\vert w_{1}\right) _{%
\mathfrak{H}}+\alpha _{2}\left( \left. v\right\vert w_{2}\right) _{\mathfrak{%
H}} \\
\left( \alpha _{1}v_{1}+\alpha _{2}v_{2}\left. w\right\vert \right) _{%
\mathfrak{H}} &=&\alpha _{1}\left( \left. v_{1}\right\vert w\right) _{%
\mathfrak{H}}+\alpha _{2}\left( \left. v_{2}\right\vert w\right) _{\mathfrak{%
H}}.
\end{eqnarray}%
The inner product $\left( \left. {}\right\vert \right) _{\mathfrak{H}}$ is
non definite as%
\begin{equation}
\left( \left. v\right\vert v\right) _{\mathfrak{H}}=0\nRightarrow v=0.
\end{equation}%
$\lceil $%
\begin{equation}
\left( 
\begin{array}{c}
i \\ 
1%
\end{array}%
\right) \left( 
\begin{array}{cc}
i & 1%
\end{array}%
\right) =-1+1=0
\end{equation}%
$\rfloor $

\subsection{If $e_{1}$ and $e_{2}$ are linearly independent then the
solution to 
\protect\begin{equation}
\protect\alpha _{1}e_{1}+\protect\alpha _{2}e_{2}=0 
\protect\end{equation}%
is 
\protect\begin{equation}
\protect\alpha _{1}=\protect\alpha _{2}=0. 
\protect\end{equation}%
One can show that if two vectors are orthogonal wrt any positive definite
Hilbert space inner product, e.g. $\left\langle \left. {}\right\vert
\right\rangle ,$ (over the same fields)$,$ then they are linearly
independent, as orthogonality gives 
\protect\begin{eqnarray}
&&\left\langle \left. \protect\alpha _{1}w_{1}+\protect\alpha %
_{2}w_{2}\right\vert w_{1}\right\rangle =\overline{\protect\alpha }%
_{1}\left\langle \left. w_{1}\right\vert w_{1}\right\rangle \\
&&\left\langle \left. \protect\alpha _{1}w_{1}+\protect\alpha %
_{2}w_{2}\right\vert w_{2}\right\rangle =\overline{\protect\alpha }%
_{2}\left\langle \left. w_{2}\right\vert w_{2}\right\rangle \\
&\Rightarrow &\text{If }\protect\alpha _{1}w_{1}+\protect\alpha _{2}w_{2}=0%
\text{ then }\protect\alpha _{1}=\protect\alpha _{2}=0, 
\protect\end{eqnarray}%
but the converse is not true i.e. linear dependence does imply orthogonality
as 
\protect\begin{equation*}
\text{If }\protect\alpha _{1}w_{1}+\protect\alpha _{2}w_{2}=0\text{ then }%
\protect\alpha _{1}=\protect\alpha _{2}=0 
\protect\end{equation*}%
only gives 
\protect\begin{eqnarray}
&&\left\langle \left. 0\right\vert w_{1}\right\rangle =0 \\
&&\left\langle \left. 0\right\vert w_{2}\right\rangle =0 
\protect\end{eqnarray}%
no matter the value of 
\protect\begin{equation}
\left\langle \left. w_{1}\right\vert w_{2}\right\rangle . 
\protect\end{equation}%
Operator Realification \label{Operator Realification}}

\begin{definition}
A \emph{real linear isomorphism }$\mathfrak{R}_{e}$ called a \emph{Vector
Realification} map 
\begin{equation}
\mathfrak{R}_{e}:\mathfrak{H\rightarrow H}_{e\mathbb{R}}
\end{equation}%
is given by 
\begin{equation}
\mathfrak{R}_{e}\left( v\right) =\mathfrak{R}_{e}\left( \sum_{1\leq j\leq
r}\left\langle \left. e_{j}\right\vert v\right\rangle \left\vert
e_{j}\right\rangle \right) =\mathfrak{R}_{e}\left( \sum_{1\leq j\leq
r}v_{j}\left\vert e_{j}\right\rangle \right) =\sum_{1\leq j\leq r}\left\{ 
\overset{}{\func{Re}\left\{ v_{j}\right\} \left\vert e_{j}\right\rangle +%
\func{Im}\left\{ v_{j}\right\} i\left\vert e_{j}\right\rangle }\right\}
=\left( 
\begin{array}{cc}
\left\vert \boldsymbol{e}\right\rangle & i\left\vert \boldsymbol{e}%
\right\rangle%
\end{array}%
\right) \left( 
\begin{array}{c}
\func{Re}_{e}\left\{ v\right\} \\ 
\func{Im}_{e}\left\{ v\right\}%
\end{array}%
\right) =\left( 
\begin{array}{cc}
\left\vert \boldsymbol{e}\right\rangle & i\left\vert \boldsymbol{e}%
\right\rangle%
\end{array}%
\right) \mathfrak{R}_{e}\left( v\right) ,
\end{equation}%
which is clearly basis dependent.
\end{definition}

\begin{definition}
The map $\mathfrak{R}_{e}$ is defined by 
\begin{equation}
\mathfrak{R}_{e}\left( v\right) =\left( 
\begin{array}{c}
\func{Re}_{e}\left\{ v\right\} \\ 
\func{Im}_{e}\left\{ v\right\}%
\end{array}%
\right)
\end{equation}
\end{definition}

The map $\mathfrak{R}_{e}$ is a real isometry wrt $h$ i.e. 
\begin{equation}
h\left( v,w\right) =h\left( \mathfrak{R}_{e}\left( v\right) ,\mathfrak{R}%
_{e}\left( w\right) \right)
\end{equation}%
and real symplectic wrt $g$ i.e. 
\begin{equation}
g\left( v,w\right) =g\left( \mathfrak{R}_{e}\left( v\right) ,\mathfrak{R}%
_{e}\left( w\right) \right) .
\end{equation}

\begin{definition}
A \emph{real algebra homomorphism, }$\widehat{\mathfrak{R}_{e}},$ called an 
\emph{Operator} \emph{Realification} map \emph{\ }%
\begin{equation}
\widehat{\mathfrak{R}_{e}}:\mathcal{B}\left( \mathfrak{H}\right) \rightarrow 
\mathcal{B}\left( \mathfrak{H}_{e\mathbb{R}}\right)
\end{equation}%
\emph{is }induced by $\mathfrak{R}_{e}$ by 
\begin{equation}
\widehat{\mathfrak{R}_{e}}\left( A\right) =\mathfrak{R}_{e}A\mathfrak{R}%
_{e}^{-1};\,A\in \mathcal{B}\left( \mathfrak{H}\right) ,
\end{equation}%
\emph{which is injective i.e. }%
\begin{equation}
\ker \left\{ \widehat{\mathfrak{R}_{e}}\right\} =0,
\end{equation}%
\emph{but} \textbf{not}\emph{\ in general surjective. }
\end{definition}

\begin{definition}
The map \emph{\ }$\widehat{\mathfrak{R}_{e}}$ is defined by 
\begin{equation}
\widehat{\mathfrak{R}_{e}}\left( A\right) =\left( 
\begin{array}{cc}
\mathcal{R}_{e}\left( \func{Re}_{e}\left( A\right) \right) & \mathcal{R}%
_{e}\left( \func{Im}_{e}\left( A\right) \right) \\ 
-\mathcal{R}_{e}\left( \func{Im}_{e}\left( A\right) \right) & \mathcal{R}%
_{e}\left( \func{Re}_{e}\left( A\right) \right)%
\end{array}%
\right)
\end{equation}
\end{definition}

\begin{remark}
\emph{In particular one can observe that even though }$iI_{\mathfrak{H}}\in 
\mathcal{Z}\left( \mathcal{B}\left( \mathfrak{H}\right) \right) ,$ which is
the centre of $\mathcal{B}\left( \mathfrak{H}\right) ,$ \emph{it does not
belong to } the centre of $\mathcal{B}\left( \mathfrak{H}_{e\mathbb{R}%
}\right) $ $\mathcal{Z}\left( \mathcal{B}\left( \mathfrak{H}_{e\mathbb{R}%
}\right) \right) .$
\end{remark}

The inner product $h$ and basis $\left\{ e_{j},ie_{j}\right\} $ of $%
\mathfrak{H}_{e\mathbb{R}}$ can be utilized to obtain a representation, $%
\mathcal{R}_{\left\{ e,ie\right\} },$ of $"iI_{\mathfrak{H}}",$ considered
as a real linear map 
\begin{equation}
iI_{\mathfrak{H}}:\mathfrak{H}_{e\mathbb{R}}\rightarrow \mathfrak{H}_{e%
\mathbb{R}},
\end{equation}%
in the real hilbert space $\mathfrak{H}_{e\mathbb{R}}$ defined by 
\begin{equation}
\mathfrak{H}_{e\mathbb{R}}=\mathbb{R}-\limfunc{linspan}\left\{ \left.
e_{j},ie_{j}\right\vert 1\leq j\leq r\right\} ,
\end{equation}%
that is given by the matrix 
\begin{align}
h\left( e_{j},i\left( e_{k}\right) \right) & =0 \\
h\left( e_{j},i\left( ie_{k}\right) \right) & =\frac{1}{2}\left\{
\left\langle \left. e_{j}\right\vert iie_{k}\right\rangle +\left\langle
\left. iie_{k}\right\vert e_{j}\right\rangle \right\} =\frac{1}{2}\left\{
-\delta _{jk}-\delta _{jk}\right\} =-\delta _{jk} \\
h\left( \left( ie_{j}\right) ,i\left( e_{k}\right) \right) & =\frac{1}{2}%
\left\{ \left\langle \left. ie_{j}\right\vert ie_{k}\right\rangle
+\left\langle \left. ie_{k}\right\vert ie_{j}\right\rangle \right\} =\frac{1%
}{2}\left\{ \delta _{jk}+\delta _{jk}\right\} =\delta _{jk} \\
h\left( \left( ie_{j}\right) ,i\left( ie_{k}\right) \right) & =\frac{1}{2}%
\left\{ \left\langle \left. ie_{j}\right\vert iie_{k}\right\rangle
+\left\langle \left. iie_{k}\right\vert ie_{j}\right\rangle \right\} =\frac{1%
}{2}\left\{ i\delta _{jk}-i\delta _{jk}\right\} =0
\end{align}%
wrt the basis $\left\{ \left\{ \left. e_{j},ie_{j}\right\vert 1\leq j\leq
r\right\} \right\} ,$ i.e. 
\begin{equation}
\mathcal{R}_{\left\{ e,ie\right\} }\left( \widehat{\mathfrak{R}_{e}}\left(
iI_{\mathfrak{H}}\right) \right) =\left( 
\begin{array}{cc}
0 & -\mathbb{I}_{r} \\ 
\mathbb{I}_{r} & 0%
\end{array}%
\right) .
\end{equation}

\begin{theorem}
$"iI_{\mathfrak{H}}"$ is represented by $\left( 
\begin{array}{cc}
0 & -\mathbb{I}_{r} \\ 
\mathbb{I}_{r} & 0%
\end{array}%
\right) $ in \emph{any }complete orthonormal basis $\left\{ \left. \left(
\left\vert e_{j}\right\rangle ,\left\vert ie_{j}\right\rangle \right)
\right\vert 1\leq j\leq r\right\} $ of $\mathfrak{H}_{e\mathbb{R}},$ wrt $h,$
induced by \emph{any} conb $\left\{ \left. \left\vert e_{j}\right\rangle
\right\vert 1\leq j\leq r\right\} $ of $\mathfrak{H.}$
\end{theorem}

\begin{proof}
See the preceding.
\end{proof}

\begin{remark}
As 
\begin{equation}
\left[ iI_{\mathfrak{H}},A\right] =0\quad \forall A\in \mathcal{B}\left( 
\mathfrak{H}\right)
\end{equation}%
then 
\begin{equation}
\left[ \widehat{\mathfrak{R}_{e}}\left( iI_{\mathfrak{H}}\right) ,\widehat{%
\mathfrak{R}_{e}}\left( A\right) \right] =0\quad \forall A\in \mathcal{B}%
\left( \mathfrak{H}\right) .
\end{equation}
\end{remark}

\begin{example}
In particular for $r=1$ one has 
\begin{gather}
\left( 
\begin{array}{cc}
a & b \\ 
c & d%
\end{array}%
\right) \left( 
\begin{array}{cc}
0 & -1 \\ 
1 & 0%
\end{array}%
\right) =\left( 
\begin{array}{cc}
b & -a \\ 
d & -c%
\end{array}%
\right) \\
\left( 
\begin{array}{cc}
0 & -1 \\ 
1 & 0%
\end{array}%
\right) \left( 
\begin{array}{cc}
a & b \\ 
c & d%
\end{array}%
\right) =\left( 
\begin{array}{cc}
-c & -d \\ 
a & b%
\end{array}%
\right)
\end{gather}%
then%
\begin{gather}
\left[ \left( 
\begin{array}{cc}
a & b \\ 
c & d%
\end{array}%
\right) ,\left( 
\begin{array}{cc}
0 & -1 \\ 
1 & 0%
\end{array}%
\right) \right] =0 \\
\Rightarrow b=-c;a=d
\end{gather}%
thus%
\begin{equation}
\mathcal{R}_{\left\{ e,ie\right\} }\left( \widehat{\mathfrak{R}_{e}}\left(
A\right) \right) =\left( 
\begin{array}{cc}
a & b \\ 
-b & a%
\end{array}%
\right) .
\end{equation}
\end{example}

\begin{theorem}
The image of the realification map $\widehat{\mathfrak{R}_{e}}$ 
\begin{equation}
\widehat{\mathfrak{R}_{e}}:\mathcal{B}\left( \mathfrak{H}\right) \rightarrow 
\mathcal{B}\left( \mathfrak{H}_{\mathbb{R}}\right)
\end{equation}%
is the\textbf{\ commutatnt of the complex unit }$J_{e}\in \mathcal{B}\left( 
\mathfrak{H}_{\mathbb{R}}\right) $ i.e. 
\begin{equation}
\widehat{\mathfrak{R}_{e}}\left\{ \mathcal{B}\left( \mathfrak{H}\right)
\right\} =\func{com}\left\{ J_{e}\right\} \subset \mathcal{B}\left( 
\mathfrak{H}_{\mathbb{R}}\right) .
\end{equation}

\begin{proof}
From the preceding
\end{proof}
\end{theorem}

\begin{remark}
In general

\begin{enumerate}
\item 
\begin{equation*}
A=\left\vert \mathbf{e}\right\rangle \mathcal{R}_{e}\left( A\right)
\left\langle \mathbf{e}\right\vert
\end{equation*}

\item 
\begin{equation}
\widehat{\mathfrak{R}_{e}}\left( iI_{2r}\right) =J_{e}
\end{equation}%
\begin{equation}
\mathcal{R}_{\left\{ e,ie\right\} }\left( J_{e}\right) =\left( 
\begin{array}{cc}
\mathbb{O} & -\mathbb{I}_{r} \\ 
\mathbb{I}_{r} & \mathbb{O}%
\end{array}%
\right)
\end{equation}

\item 
\begin{eqnarray}
&&\widehat{\mathfrak{R}_{e}}\left( \overset{}{\left\vert \mathbf{e}%
\right\rangle \mathcal{R}_{e}\left( A\right) \left\langle \mathbf{e}%
\right\vert }\right) =\left\vert \mathbf{e,}i\mathbf{e}\right\rangle 
\mathcal{R}_{\left\{ e,ie\right\} }\left( \widehat{\mathfrak{R}_{e}}\left(
A\right) \right) \left\langle \mathbf{e,}i\mathbf{e}\right\vert \\
&=&\widehat{\mathfrak{R}_{e}}\left( u^{\dagger }u\overset{}{\left\vert 
\mathbf{e}\right\rangle \mathcal{R}_{e}\left( A\right) \left\langle \mathbf{e%
}\right\vert }u^{\dagger }u\right) =\widehat{\mathfrak{R}_{e}}\left( \overset%
{}{u^{\dagger }\left\vert \mathbf{e}\right\rangle \mathcal{R}_{e}\left(
u\right) \mathcal{R}_{e}\left( A\right) \mathcal{R}_{e}\left( u\right)
^{\dagger }\left\langle \mathbf{e}\right\vert u}\right) =\widehat{\mathfrak{R%
}_{e}}\left( \overset{}{\left\vert u^{\dagger }\mathbf{e}\right\rangle 
\mathcal{R}_{e}\left( u\right) \mathcal{R}_{e}\left( A\right) \mathcal{R}%
_{e}\left( u\right) ^{\dagger }\left\langle u\mathbf{e}\right\vert }\right)
\\
&=&\widehat{\mathfrak{R}_{e}}\left( \overset{}{\left\vert \mathbf{f}%
\right\rangle \mathcal{R}_{e}\left( u\right) \mathcal{R}_{e}\left( A\right) 
\mathcal{R}_{e}\left( u\right) ^{\dagger }\left\langle \mathbf{f}\right\vert 
}\right) =\widehat{\mathfrak{R}_{e}}\left( \overset{}{\left\vert \mathbf{f}%
\right\rangle \mathcal{R}_{f}\left( A\right) \left\langle \mathbf{f}%
\right\vert }\right) ,
\end{eqnarray}%
where 
\begin{eqnarray}
\left\vert \mathbf{f}\right\rangle &=&u^{\dagger }\left\vert \mathbf{e}%
\right\rangle =\left\vert \mathbf{e}\right\rangle \mathcal{R}_{e}\left(
u\right) ^{\dagger }, \\
\mathcal{R}_{f}\left( A\right) &=&\mathcal{R}_{e}\left( u\right) \mathcal{R}%
_{e}\left( A\right) \mathcal{R}_{e}\left( u\right) ^{\dagger }
\end{eqnarray}%
and

\item 
\begin{equation}
A=\left\vert \mathbf{e}\right\rangle \mathcal{R}_{e}\left( A\right)
\left\langle \mathbf{e}\right\vert =\left\vert \mathbf{f}\right\rangle 
\mathcal{R}_{f}\left( A\right) \left\langle \mathbf{f}\right\vert
=\left\vert \mathbf{e}\right\rangle \left( \overset{}{\mathcal{R}_{e}\left(
u\right) ^{\dagger }\mathcal{R}_{f}\left( A\right) \mathcal{R}_{e}\left(
u\right) }\right) \left\langle \mathbf{e}\right\vert .
\end{equation}

\item 
\begin{eqnarray}
\widehat{\mathfrak{R}_{e}}\left( \overset{}{\left\vert \mathbf{f}%
\right\rangle \mathcal{R}_{f}\left( A\right) \left\langle \mathbf{f}%
\right\vert }\right) &=&\widehat{\mathfrak{R}_{e}}\left( \overset{}{%
\left\vert \mathbf{e}\right\rangle \left( \mathcal{R}_{e}\left( u\right)
^{\dagger }\mathcal{R}_{f}\left( A\right) \mathcal{R}_{e}\left( u\right)
\right) \left\langle \mathbf{e}\right\vert }\right) =\left\vert \mathbf{e,}i%
\mathbf{e}\right\rangle \left( \mathcal{R}_{\left\{ e,ie\right\} }\left( 
\widehat{\mathfrak{R}_{e}}\left( u\right) \right) ^{\dagger }\mathcal{R}%
_{\left\{ f,if\right\} }\left( \widehat{\mathfrak{R}_{e}}\left( A\right)
\right) \mathcal{R}_{\left\{ e,ie\right\} }\left( \widehat{\mathfrak{R}_{e}}%
\left( u\right) \right) \right) \left\langle \mathbf{e,}i\mathbf{e}%
\right\vert \\
&=&\left\vert \mathbf{e,}i\mathbf{e}\right\rangle \left( \left( 
\begin{array}{cc}
\mathcal{R}_{\left\{ e,ie\right\} }\left( \func{Re}_{e}\left( u\right)
\right) & \mathcal{R}_{\left\{ e,ie\right\} }\left( \func{Im}_{e}\left(
u\right) \right) \\ 
-\mathcal{R}_{\left\{ e,ie\right\} }\left( \func{Im}_{e}\left( u\right)
\right) & \mathcal{R}_{\left\{ e,ie\right\} }\left( \func{Re}_{e}\left(
u\right) \right)%
\end{array}%
\right) ^{\dagger }\mathcal{R}_{f}\left( \widehat{\mathfrak{R}_{e}}\left(
A\right) \right) \left( 
\begin{array}{cc}
\mathcal{R}_{\left\{ e,ie\right\} }\left( \func{Re}_{e}\left( u\right)
\right) & \mathcal{R}_{\left\{ e,ie\right\} }\left( \func{Im}_{e}\left(
u\right) \right) \\ 
-\mathcal{R}_{\left\{ e,ie\right\} }\left( \func{Im}_{e}\left( u\right)
\right) & \mathcal{R}_{\left\{ e,ie\right\} }\left( \func{Re}_{e}\left(
u\right) \right)%
\end{array}%
\right) \right) \left\langle \mathbf{e,}i\mathbf{e}\right\vert ,
\end{eqnarray}%
where%
\begin{equation}
\mathcal{R}_{\left\{ e,ie\right\} }\left( \overset{}{\widehat{\mathfrak{R}%
_{e}}\left( u\right) }\right) =\left( 
\begin{array}{cc}
\mathcal{R}_{\left\{ e,ie\right\} }\left( \widehat{\mathfrak{R}_{e}}\left( 
\func{Re}_{e}\left( u\right) \right) \right) & \mathcal{R}_{\left\{
e,ie\right\} }\left( \widehat{\mathfrak{R}_{e}}\left( \func{Im}_{e}\left(
u\right) \right) \right) \\ 
-\mathcal{R}_{\left\{ e,ie\right\} }\left( \widehat{\mathfrak{R}_{e}}\left( 
\func{Im}_{e}\left( u\right) \right) \right) & \mathcal{R}_{\left\{
e,ie\right\} }\left( \widehat{\mathfrak{R}_{e}}\left( \func{Re}_{e}\left(
u\right) \right) \right)%
\end{array}%
\right) .
\end{equation}

\item 
\begin{equation}
\widehat{\mathfrak{R}_{e}}\left( \mathcal{R}_{e}\left( u\right) \right) =%
\mathcal{R}_{\left\{ e,ie\right\} }\left( \widehat{\mathfrak{R}_{e}}\left( 
\overset{}{\left( \overset{}{\func{Re}_{e}\left( u\right) }\right) +i\left( 
\overset{}{\func{Im}_{e}\left( u\right) }\right) }\right) \right) =\left( 
\begin{array}{cc}
\mathcal{R}_{e}\left( \func{Re}_{e}\left( u\right) \right) & \mathcal{R}%
_{e}\left( \func{Im}_{e}\left( u\right) \right) \\ 
-\mathcal{R}_{e}\left( \func{Im}_{e}\left( u\right) \right) & \mathcal{R}%
_{e}\left( \func{Re}_{e}\left( u\right) \right)%
\end{array}%
\right) .
\end{equation}

\item 
\begin{eqnarray}
\widehat{\mathfrak{R}_{e}}\left( \mathcal{R}_{e}\left( \func{Re}_{e}\left(
u\right) \right) \right) &=&\left( 
\begin{array}{cc}
\mathcal{R}_{e}\left( \func{Re}_{e}\left( u\right) \right) & \mathbb{O}_{r}
\\ 
\mathbb{O}_{r} & \mathcal{R}_{e}\left( \func{Re}_{e}\left( u\right) \right)%
\end{array}%
\right) \\
\widehat{\mathfrak{R}_{e}}\left( \mathcal{R}_{e}\left( \func{Im}_{e}\left(
u\right) \right) \right) &=&\left( 
\begin{array}{cc}
\mathbb{O}_{r} & \mathcal{R}_{e}\left( \func{Im}_{e}\left( u\right) \right)
\\ 
-\mathcal{R}_{e}\left( \func{Im}_{e}\left( u\right) \right) & \mathbb{O}_{r}%
\end{array}%
\right) .
\end{eqnarray}%
\begin{equation}
\mathcal{R}_{\left\{ e,ie\right\} }\left( \overset{}{\widehat{\mathfrak{R}%
_{e}}\left( u\right) }\right) =\left( 
\begin{array}{cc}
\mathbb{U} & \mathbb{V} \\ 
-\mathbb{V} & \mathbb{U}%
\end{array}%
\right) =\left( 
\begin{array}{cc}
\mathcal{R}_{e}\left( \func{Re}_{e}\left( u\right) \right) & \mathcal{R}%
_{e}\left( \func{Im}_{e}\left( u\right) \right) \\ 
-\mathcal{R}_{e}\left( \func{Im}_{e}\left( u\right) \right) & \mathcal{R}%
_{e}\left( \func{Re}_{e}\left( u\right) \right)%
\end{array}%
\right)
\end{equation}

\begin{enumerate}
\item The unitary condition gives%
\begin{equation}
u^{\dagger }u=I_{2r}\Leftrightarrow \mathcal{R}_{e}\left( u\right) ^{\dagger
}\mathcal{R}_{e}\left( u\right) =\mathbb{I}_{2r}\Leftrightarrow \widehat{%
\mathfrak{R}_{e}}\left( u\right) ^{t}\widehat{\mathfrak{R}_{e}}\left(
u\right) =I_{4r}\Leftrightarrow \widehat{\mathfrak{R}_{e}}\left( \mathcal{R}%
_{e}\left( u\right) \right) ^{t}\widehat{\mathfrak{R}_{e}}\left( \mathcal{R}%
_{e}\left( u\right) \left( u\right) \right) =\mathbb{I}_{4r}\Leftrightarrow 
\mathcal{R}_{\left\{ e,ie\right\} }\left( \widehat{\mathfrak{R}_{e}}\left(
u\right) \right) ^{t}\mathcal{R}_{\left\{ e,ie\right\} }\left( \widehat{%
\mathfrak{R}_{e}}\left( u\right) \right) =\mathbb{I}_{4r}.
\end{equation}

\item In more detail 
\begin{equation}
\left( \mathcal{R}_{\left\{ e,ie\right\} }\left( \widehat{\mathfrak{R}_{e}}%
\left( \cong u\right) \right) \right) ^{t}\mathcal{R}_{\left\{ e,ie\right\}
}\left( \widehat{\mathfrak{R}_{e}}\left( u\right) \right) =\left( 
\begin{array}{cc}
\mathcal{R}_{\left\{ e,ie\right\} }\left( \widehat{\mathfrak{R}_{e}}\left( 
\func{Re}_{e}\left( u\right) \right) \right) ^{t} & -\mathcal{R}_{\left\{
e,ie\right\} }\left( \widehat{\mathfrak{R}_{e}}\left( \func{Im}_{e}\left(
u\right) \right) \right) ^{t} \\ 
\mathcal{R}_{\left\{ e,ie\right\} }\left( \widehat{\mathfrak{R}_{e}}\left( 
\func{Im}_{e}\left( u\right) \right) \right) ^{t} & \mathcal{R}_{\left\{
e,ie\right\} }\left( \widehat{\mathfrak{R}_{e}}\left( \func{Re}_{e}\left(
u\right) \right) \right) ^{t}%
\end{array}%
\right) \left( 
\begin{array}{cc}
\mathcal{R}_{\left\{ e,ie\right\} }\left( \widehat{\mathfrak{R}_{e}}\left( 
\func{Re}_{e}\left( u\right) \right) \right) & \mathcal{R}_{\left\{
e,ie\right\} }\left( \widehat{\mathfrak{R}_{e}}\left( \func{Im}_{e}\left(
u\right) \right) \right) \\ 
-\mathcal{R}_{\left\{ e,ie\right\} }\left( \widehat{\mathfrak{R}_{e}}\left( 
\func{Im}_{e}\left( u\right) \right) \right) & \mathcal{R}_{\left\{
e,ie\right\} }\left( \widehat{\mathfrak{R}_{e}}\left( \func{Re}_{e}\left(
u\right) \right) \right)%
\end{array}%
\right) =\left( 
\begin{array}{cc}
\mathbb{I}_{r} & \mathbb{O} \\ 
\mathbb{O} & \mathbb{I}_{r}%
\end{array}%
\right) .
\end{equation}

\item The commutator condition 
\begin{equation}
\left[ \widehat{\mathfrak{R}_{e}}\left( iI_{\mathfrak{H}}\right) ,\widehat{%
\mathfrak{R}_{e}}\left( A\right) \right] =0\quad \forall A\in \mathcal{B}%
\left( \mathfrak{H}\right)
\end{equation}
gives%
\begin{equation}
\left[ \widehat{\mathfrak{R}_{e}}\left( iI_{\mathfrak{H}}\right) ,\widehat{%
\mathfrak{R}_{e}}\left( u\right) \right] =0\quad \forall u\in U\left( 
\mathfrak{H}\right)
\end{equation}%
i.e. 
\begin{equation}
\widehat{\mathfrak{R}_{e}}\left( iI_{\mathfrak{H}}\right) \widehat{\mathfrak{%
R}_{e}}\left( u\right) =\widehat{\mathfrak{R}_{e}}\left( u\right) \widehat{%
\mathfrak{R}_{e}}\left( iI_{\mathfrak{H}}\right) \Leftrightarrow \mathcal{R}%
_{\left\{ e,ie\right\} }\left( J_{e}\right) \mathcal{R}_{\left\{
e,ie\right\} }\left( \widehat{\mathfrak{R}_{e}}\left( u\right) \right) =%
\mathcal{R}_{\left\{ e,ie\right\} }\left( \widehat{\mathfrak{R}_{e}}\left(
u\right) \right) \mathcal{R}_{\left\{ e,ie\right\} }\left( J_{e}\right)
\end{equation}%
\begin{gather}
\left( 
\begin{array}{cc}
\mathbb{O} & -\mathbb{I}_{2r} \\ 
\mathbb{I}_{2r} & \mathbb{O}%
\end{array}%
\right) \mathcal{R}_{\left\{ e,ie\right\} }\left( \widehat{\mathfrak{R}_{e}}%
\left( u\right) \right) =\mathcal{R}_{\left\{ e,ie\right\} }\left( \widehat{%
\mathfrak{R}_{e}}\left( u\right) \right) \left( 
\begin{array}{cc}
\mathbb{O} & -\mathbb{I}_{2r} \\ 
\mathbb{I}_{2r} & \mathbb{O}%
\end{array}%
\right) \\
\Updownarrow \\
\mathcal{R}_{\left\{ e,ie\right\} }\left( \widehat{\mathfrak{R}_{e}}\left(
u\right) \right) ^{t}\left( 
\begin{array}{cc}
\mathbb{O} & -\mathbb{I}_{2r} \\ 
\mathbb{I}_{2r} & \mathbb{O}%
\end{array}%
\right) \mathcal{R}_{\left\{ e,ie\right\} }\left( \widehat{\mathfrak{R}_{e}}%
\left( u\right) \right) =\left( 
\begin{array}{cc}
\mathbb{O} & -\mathbb{I}_{2r} \\ 
\mathbb{I}_{2r} & \mathbb{O}%
\end{array}%
\right) ,
\end{gather}%
which shows that 
\begin{equation}
\mathcal{R}_{\left\{ e,ie\right\} }\left( \widehat{\mathfrak{R}_{e}}\left(
u\right) \right) \in Sp\left( 4r,\mathbb{R}\right) ,
\end{equation}

\item Hence 
\begin{equation}
\mathcal{R}_{\left\{ e,ie\right\} }\left( \widehat{\mathfrak{R}_{e}}\left(
u\right) \right) \in Sp\left( 4r,\mathbb{R}\right) \cap O\left( 4r,\mathbb{R}%
\right)
\end{equation}%
and 
\begin{equation}
Sp\left( 4r,\mathbb{R}\right) \cap O\left( 4r,\mathbb{R}\right) \cong
U\left( 2r\right) .
\end{equation}
\end{enumerate}
\end{enumerate}
\end{remark}

\subsubsection{Real Spaces Induced fro Complex Spaces}

\bigskip It is interesting and important to focus on the differences and
similaraties of the following real spaces

\begin{enumerate}
\item 
\begin{equation}
\mathfrak{H}_{e\mathbb{R}}=\mathbb{R}-\limfunc{linspan}\left\{ \left.
e_{j},ie_{j}\right\vert 1\leq j\leq r\right\}
\end{equation}

\item 
\begin{equation}
\widetilde{\mathfrak{H}_{e\mathbb{R}}}=\,\func{Re}_{e}\left( \mathfrak{H}_{%
\mathbb{C}}\right) \oplus _{\mathbb{R}}i\func{Im}_{e}\left( \mathfrak{H}_{%
\mathbb{C}}\right) =\mathbb{R}-\limfunc{linspan}\left\{ \left.
e_{j}\right\vert 1\leq j\leq r\right\} \oplus \mathbb{R}-\limfunc{linspan}%
\left\{ \left. ie_{j}\right\vert 1\leq j\leq r\right\}
\end{equation}

\item 
\begin{equation}
\widehat{\mathfrak{H}_{e\mathbb{R}}}=\func{Re}_{e}\left( \mathfrak{H}_{%
\mathbb{C}}\right) \oplus _{\mathbb{R}}\func{Im}_{e}\left( \mathfrak{H}_{%
\mathbb{C}}\right) =\mathbb{R}-\limfunc{linspan}\left\{ \left.
e_{j}\right\vert 1\leq j\leq r\right\} \oplus \mathbb{R}-\limfunc{linspan}%
\left\{ \left. e_{j}\right\vert 1\leq j\leq r\right\} .
\end{equation}
\end{enumerate}

\begin{notation}
$\mathfrak{H}_{\mathbb{C}}\equiv \mathfrak{H}$
\end{notation}

\begin{remark}
In more detail 
\begin{align}
\mathfrak{H}_{e\mathbb{R}}& =\func{Re}_{e}\left( \mathfrak{H}_{\mathbb{C}%
}\right) +i\func{Im}_{e}\left( \mathfrak{H}_{\mathbb{C}}\right) =\mathbb{R}-%
\limfunc{linspan}\left\{ \left. e_{j},ie_{j}\right\vert 1\leq j\leq r\right\}
\\
& =\mathbb{R}-\limfunc{linspan}\left\{ \left. e_{j}\right\vert 1\leq j\leq
r\right\} +\mathbb{R}-\limfunc{linspan}\left\{ \left. ie_{j}\right\vert
1\leq j\leq r\right\} ,
\end{align}%
\begin{equation}
\shortparallel
\end{equation}%
\begin{align}
\widetilde{\mathfrak{H}_{e\mathbb{R}}}& =\func{Re}_{e}\left( \mathfrak{H}_{%
\mathbb{C}}\right) \oplus _{\mathbb{R}}i\func{Im}_{e}\left( \mathfrak{H}_{%
\mathbb{C}}\right) =\mathbb{R}-\limfunc{linspan}\left\{ \left. \left(
e_{j},0\right) ,\left( 0,ie_{j}\right) \right\vert 1\leq j\leq r\right\} \\
& =\mathbb{R}-\limfunc{linspan}\left\{ \left. \left( e_{j},0\right)
\right\vert 1\leq j\leq r\right\} +\mathbb{R}-\limfunc{linspan}\left\{
\left. \left( 0,ie_{j}\right) \right\vert 1\leq j\leq r\right\} \\
& \cong \mathbb{R}-\limfunc{linspan}\left\{ \left. e_{j}\right\vert 1\leq
j\leq r\right\} +\mathbb{R}-\limfunc{linspan}\left\{ \left.
ie_{j}\right\vert 1\leq j\leq r\right\} ,
\end{align}%
\begin{equation}
\shortparallel
\end{equation}%
\begin{align}
\widehat{\mathfrak{H}_{e\mathbb{R}}}& =\func{Re}_{e}\left( \mathfrak{H}_{%
\mathbb{C}}\right) \oplus _{\mathbb{R}}\func{Im}_{e}\left( \mathfrak{H}_{%
\mathbb{C}}\right) =\mathbb{R}-\limfunc{linspan}\left\{ \left. \left(
e_{j},0\right) ,\left( 0,e_{j}\right) \right\vert 1\leq j\leq r\right\} \\
& =\mathbb{R}-\limfunc{linspan}\left\{ \left. \left( e_{j},0\right)
\right\vert 1\leq j\leq r\right\} +\mathbb{R}-\limfunc{linspan}\left\{
\left. \left( 0,e_{j}\right) \right\vert 1\leq j\leq r\right\} \\
& \ncong \mathbb{R}-\limfunc{linspan}\left\{ \left. e_{j}\right\vert 1\leq
j\leq r\right\} +\mathbb{R}-\limfunc{linspan}\left\{ \left. e_{j}\right\vert
1\leq j\leq r\right\} .
\end{align}
\end{remark}

\subsubsection{\textbf{The }$r=1$\textbf{\ Case}}

\begin{example}
The simplest example is for $r=1$, which corresponds to the complex numbers $%
\mathbb{C}$ 
\begin{equation}
\mathbb{C}=\left\{ \left. x+iy\right\vert x,y\in \mathbb{R}\right\} =\left\{
\left. x+iy\right\vert x,y\in \mathbb{R}\right\}
\end{equation}%
and the associated real space $\mathbb{R}^{2}$. Complexifying the real space 
$\mathbb{R}$ gives 
\begin{equation}
\mathbb{C\otimes R}=\mathbb{R}\oplus _{\mathbb{C}}i\mathbb{R\cong R}+i%
\mathbb{R\cong R\oplus _{\mathbb{R}}R\cong R}^{2}=\mathbb{R\times R},
\end{equation}%
where we have used the basis $\left\{ 1,i\right\} $ of $\mathbb{C},$ ($%
\mathbb{R\times R}$ has natural vector space structures induced by $\mathbb{%
R\oplus _{\mathbb{R}}R}$ or $\mathbb{R}\mathbb{\oplus }\mathbb{_{\mathbb{R}}}%
i\mathbb{R}$, the direct sum $\oplus _{\mathbb{C}}$is the complex direct sum
in the preceeding equation as complex multiplication is defined in $\mathbb{%
C\otimes R}).$

One can consider the following three isomorphic spaces 
\begin{align}
\mathbb{R}\oplus _{\mathbb{R}}i\mathbb{R}& =\mathbb{R}-\limfunc{linspan}%
\left\{ \left( 1,0\right) ,\left( 0,i\right) \right\} \\
\mathbb{R}+i\mathbb{R}& =\mathbb{R}-\limfunc{linspan}\left\{ 1,i1\right\} =%
\mathbb{R}-\limfunc{linspan}\left\{ 1,i\right\} \\
\mathbb{R\oplus _{\mathbb{R}}R}& =\mathbb{R}-\limfunc{linspan}\left\{ \left(
1,0\right) ,\left( 0,1\right) \right\} \,=\mathbb{R}^{2}
\end{align}%
as real spaces 
\begin{equation}
\mathbb{R\oplus _{\mathbb{R}}}i\mathbb{R\cong R}+i\mathbb{R\cong R\oplus _{%
\mathbb{R}}R}
\end{equation}%
but $\mathbb{R\oplus _{\mathbb{R}}R}$ is not a complex space so the real
isomorphism only extends to a complex isomorphism when one considers $%
\mathbb{R}\oplus _{\mathbb{C}}i\mathbb{R\ }$and $\mathbb{R}+i\mathbb{R}.$

Vectors in these spaces can be expanded as 
\begin{align}
v_{1}& =a\left( 1,0\right) +b\left( 0,i\right) =\left( \left( 1,0\right)
,\left( 0,i\right) \right) \left( 
\begin{array}{c}
a \\ 
b%
\end{array}%
\right) \in \mathbb{R}\oplus _{\mathbb{R}}i\mathbb{R} \\
v_{2}& =a\left( 1,0\right) +b\left( 0,1\right) =\left( \left( 1,0\right)
,\left( 0,1\right) \right) \left( 
\begin{array}{c}
a \\ 
b%
\end{array}%
\right) \in \mathbb{R\oplus _{\mathbb{R}}R} \\
v_{3}& =a1+bi=\left( 1,i\right) \left( 
\begin{array}{c}
a \\ 
b%
\end{array}%
\right) \in \mathbb{R}+i\mathbb{R}.
\end{align}

The action of $i$ on these basis of $\mathbb{R\oplus _{\mathbb{R}}}i\mathbb{R%
}$ and $\mathbb{R}+i\mathbb{R}$ is given by 
\begin{align}
i\left( \left( 1,0\right) ,\left( 0,i\right) \right) & =\left( i\left(
1,0\right) ,i\left( 0,i\right) \right) =\left( \left( 0,i\right) ,\left(
-1,0\right) \right) =\left( \left( 1,0\right) ,\left( 0,i\right) \right)
\left( 
\begin{array}{cc}
0 & -1 \\ 
1 & 0%
\end{array}%
\right) \\
i\left( 1,i\right) & =\left( i,-1\right) =\left( 1,i\right) \left( 
\begin{array}{cc}
0 & -1 \\ 
1 & 0%
\end{array}%
\right)
\end{align}%
giving identical representations. \emph{Note that} $i$ \emph{does not} have
a representation on $\mathbb{R\oplus _{\mathbb{R}}R}$ as it does not act in
it as a linear automorphism i.e. it does not map $\mathbb{R\oplus _{\mathbb{R%
}}R}$ to itself. One can define a linear map $\mathcal{J}$ in $\mathbb{%
R\oplus _{\mathbb{R}}R}$ by the matrix $\left( 
\begin{array}{cc}
0 & -1 \\ 
1 & 0%
\end{array}%
\right) $ that has the following effect 
\begin{align}
\mathcal{J}\left( \left( 1,0\right) ,\left( 0,1\right) \right) & =\left( 
\mathcal{J}\left( 1,0\right) ,\mathcal{J}\left( 0,1\right) \right) =\left(
\left( 0,1\right) ,\left( -1,0\right) \right) =\left( \left( 1,0\right)
,\left( 0,1\right) \right) \left( 
\begin{array}{cc}
0 & -1 \\ 
1 & 0%
\end{array}%
\right) \\
\mathcal{J}\left( 1,0\right) & =\left( 0,1\right) \\
\mathcal{J}\left( 0,1\right) & =\left( -1,0\right) ,
\end{align}%
that allows us to define multiplication by $"i"$ by%
\begin{align}
i\left( 1,0\right) & =\mathcal{J}\left( 1,0\right) =\left( 0,1\right) \\
i\left( 0,1\right) & =\mathcal{J}\left( 0,1\right) =\left( -1,0\right)
\end{align}%
giving%
\begin{align}
v_{2}& =a\left( 1,0\right) +b\left( 0,1\right) =\left( \left( 1,0\right)
,\left( 0,1\right) \right) \left( 
\begin{array}{c}
a \\ 
b%
\end{array}%
\right) \in \mathbb{R\oplus _{\mathbb{R}}R} \\
& =a\left( 1,0\right) +b\mathcal{J}\left( 1,0\right) =\left( \left(
1,0\right) ,\mathcal{J}\left( 1,0\right) \right) \left( 
\begin{array}{c}
a \\ 
b%
\end{array}%
\right) \in \mathbb{R}\oplus _{\mathbb{R}}J\mathbb{R} \\
& =a\left( 1,0\right) +bi\left( 1,0\right) =\left( \left( 1,0\right)
,i\left( 1,0\right) \right) \left( 
\begin{array}{c}
a \\ 
b%
\end{array}%
\right) \in \mathbb{R}\oplus _{\mathbb{R}}i\mathbb{R}.
\end{align}%
$\mathcal{J},J$ and $i$ can be defined in terms of each other as 
\begin{align}
\mathcal{J}\left( 1,0\right) & =\left( 0,J\left( 1\right) \right) =\left(
0,i\right) \\
\mathcal{J}\left( 0,1\right) & =\left( -1,0\right) =\left( J\left( i\right)
,0\right) ,
\end{align}%
where 
\begin{align}
J& :\mathbb{R}\rightarrow i\mathbb{R} \\
J& :i\mathbb{R}\rightarrow \mathbb{R}.
\end{align}
\end{example}

\begin{remark}
Note 
\begin{align}
& \left( 
\begin{array}{cc}
0 & -1 \\ 
1 & 0%
\end{array}%
\right) \left( 
\begin{array}{cc}
0 & -1 \\ 
1 & 0%
\end{array}%
\right) =\left( 
\begin{array}{cc}
-1 & 0 \\ 
0 & -1%
\end{array}%
\right) \\
& \left( 
\begin{array}{cc}
0 & -1 \\ 
1 & 0%
\end{array}%
\right) \left( 
\begin{array}{c}
0 \\ 
b%
\end{array}%
\right) =\left( 
\begin{array}{c}
-b \\ 
0%
\end{array}%
\right) \\
& \left( 
\begin{array}{cc}
0 & -1 \\ 
1 & 0%
\end{array}%
\right) \left( 
\begin{array}{c}
a \\ 
0%
\end{array}%
\right) =\left( 
\begin{array}{c}
0 \\ 
a%
\end{array}%
\right) .
\end{align}
\end{remark}

Expanding in a more detailed fashion :

1. The complexified space $\mathbb{C\otimes R}$ $=\mathbb{R}\oplus _{\mathbb{%
R}}i\mathbb{R}$ (external complexification) gives 
\begin{gather}
v_{1}=\left( \left( 1,0\right) ,\left( 0,i\right) \right) \left( 
\begin{array}{c}
a \\ 
b%
\end{array}%
\right) =\left( \left( 1,0\right) ,\left( 0,i\right) \right) \left( 
\begin{array}{c}
a \\ 
0%
\end{array}%
\right) +\left( \left( 1,0\right) ,\left( 0,i\right) \right) \left( 
\begin{array}{c}
0 \\ 
b%
\end{array}%
\right) \\
=\left( \left( 1,0\right) ,\left( 0,i\right) \right) \left( 
\begin{array}{c}
a \\ 
0%
\end{array}%
\right) +\left( \left( 1,0\right) ,\left( 0,i\right) \right) \left( 
\begin{array}{cc}
0 & -1 \\ 
1 & 0%
\end{array}%
\right) \left( 
\begin{array}{c}
b \\ 
0%
\end{array}%
\right) =\left( \left( 1,0\right) ,\left( 0,i\right) \right) \left( 
\begin{array}{c}
a \\ 
0%
\end{array}%
\right) +\left( \left( 1,0\right) ,\left( 0,i\right) \right) \mathcal{R}%
\left( \mathcal{J}\right) \left( 
\begin{array}{c}
b \\ 
0%
\end{array}%
\right) \\
=\left( \left( 1,0\right) ,\left( 0,i\right) \right) \left( 
\begin{array}{c}
a \\ 
0%
\end{array}%
\right) +\mathcal{J}\left( \left( 1,0\right) ,\left( 0,i\right) \right)
\left( 
\begin{array}{c}
b \\ 
0%
\end{array}%
\right) =\left( \left( a,0\right) ,\left( 0,i\right) \right) +\mathcal{J}%
\left( \left( b,0\right) ,\left( 0,i\right) \right) =\left( a,0\right) +%
\mathcal{J}\left( b,0\right) \\
=v_{1a}+\mathcal{J}v_{1b};v_{1},v_{1a},v_{1b}\in \mathbb{R\oplus }i\mathbb{R}
\\
\longleftrightarrow a+Jb;a,b\in \mathbb{R}.
\end{gather}%
2. The complexified space $\mathbb{C}$ $=\mathbb{R}\oplus _{\mathbb{R}}J%
\mathbb{R}$ (internal complexification, through factorization) gives%
\begin{gather}
v_{2}=\left( \left( 1,0\right) ,\left( 0,1\right) \right) \left( 
\begin{array}{c}
a \\ 
b%
\end{array}%
\right) =\left( \left( 1,0\right) ,\left( 0,1\right) \right) \left( 
\begin{array}{c}
a \\ 
0%
\end{array}%
\right) +\left( \left( 1,0\right) ,\left( 0,1\right) \right) \left( 
\begin{array}{c}
0 \\ 
b%
\end{array}%
\right) \\
=\left( \left( 1,0\right) ,\left( 0,1\right) \right) \left( 
\begin{array}{c}
a \\ 
0%
\end{array}%
\right) +\left( \left( 1,0\right) ,\left( 0,1\right) \right) \left( 
\begin{array}{cc}
0 & -1 \\ 
1 & 0%
\end{array}%
\right) \left( 
\begin{array}{c}
b \\ 
0%
\end{array}%
\right) =\left( \left( 1,0\right) ,\left( 0,1\right) \right) \left( 
\begin{array}{c}
a \\ 
0%
\end{array}%
\right) +\left( \left( 1,0\right) ,\left( 0,1\right) \right) \mathcal{R}%
\left( \mathcal{J}\right) \left( 
\begin{array}{c}
b \\ 
0%
\end{array}%
\right) \\
=\left( \left( 1,0\right) ,\left( 0,1\right) \right) \left( 
\begin{array}{c}
a \\ 
0%
\end{array}%
\right) +\mathcal{J}\left( \left( 1,0\right) ,\left( 0,1\right) \right)
\left( 
\begin{array}{c}
b \\ 
0%
\end{array}%
\right) =\left( \left( a,0\right) ,\left( 0,1\right) \right) +\mathcal{J}%
\left( \left( b,0\right) ,\left( 0,1\right) \right) =\left( a,0\right) +%
\mathcal{J}\left( b,0\right) \\
=v_{2a}+\mathcal{J}v_{2b};\,v_{2},v_{2a},v_{2b}\in \mathbb{R\oplus }J\mathbb{%
R} \\
\longleftrightarrow a+Jb;\,a,b\in \mathbb{R}.
\end{gather}%
3. The space $\mathbb{C=R\oplus }i\mathbb{R}$%
\begin{gather}
v_{3}=\left( 1,i\right) \left( 
\begin{array}{c}
a \\ 
b%
\end{array}%
\right) =\left( 1,i\right) \left( 
\begin{array}{c}
a \\ 
0%
\end{array}%
\right) +\left( 1,i\right) \left( 
\begin{array}{c}
0 \\ 
b%
\end{array}%
\right) =\left( 1,i\right) \left( 
\begin{array}{c}
a \\ 
0%
\end{array}%
\right) +\left( 1,i\right) \left( 
\begin{array}{cc}
0 & -1 \\ 
1 & 0%
\end{array}%
\right) \left( 
\begin{array}{c}
b \\ 
0%
\end{array}%
\right) \\
=\left( 1,i\right) \left( 
\begin{array}{c}
a \\ 
0%
\end{array}%
\right) +\left( 1,i\right) \mathcal{R}\left( J\right) \left( 
\begin{array}{c}
b \\ 
0%
\end{array}%
\right) =\left( 1,i\right) \left( 
\begin{array}{c}
a \\ 
0%
\end{array}%
\right) +J\left( 1,i\right) \left( 
\begin{array}{c}
b \\ 
0%
\end{array}%
\right) =a+Jb \\
=v_{3a}+Jv_{3b}=a+ib
\end{gather}%
and 
\begin{align}
x+iy& =\left( 
\begin{array}{c}
x \\ 
0%
\end{array}%
\right) +\left( 
\begin{array}{cc}
0 & -1 \\ 
1 & 0%
\end{array}%
\right) \left( 
\begin{array}{c}
y \\ 
0%
\end{array}%
\right) \\
\left( x,0\right) +i\left( y,0\right) & =\left( x,0\right) +\left( 
\begin{array}{cc}
0 & -1 \\ 
1 & 0%
\end{array}%
\right) \left( y,0\right) =\left( x,0\right) +\left( 
\begin{array}{cc}
0 & -1 \\ 
1 & 0%
\end{array}%
\right) \left( y,0\right) .
\end{align}

\subsubsection{Complexification of Real Hilbert Space I (External)}

Let $V$ be a \textbf{rea}l Hilbert space of dimension $r$ with inner product 
$\left( \left. a\right\vert b\right) =\left( \left. b\right\vert a\right) $
then its complexification is defined to be 
\begin{align}
& V_{\mathbb{C}}=\mathbb{C}\otimes _{\mathbb{R}}V=V\oplus _{\mathbb{R}%
}iV\cong V+iV \\
& v_{c}=v_{1}+iv_{2};\,v_{1},v_{2}\in V;v_{c}\in V_{\mathbb{C}} \\
& \left( \text{see later for how multiplication by }i\text{, direct sum and
tensor are defined}\right) \\
& \mathbb{C-}\dim \left\{ V_{\mathbb{C}}\right\} =r;\,\mathbb{R-}\dim
\left\{ V_{\mathbb{C}}\right\} =2r.
\end{align}%
The real direct sum $\oplus _{\mathbb{R}}$can be replaced by the complex
direct sum $\oplus _{\mathbb{C}}$as complex multiplication is defined in $%
\mathbb{C}\otimes _{\mathbb{R}}V$ and different basis correspond to
multiplication by different complex numbers. We can also define the direct
sum $V_{\mathbb{R}}$ 
\begin{equation}
V_{\mathbb{R}}=V\oplus _{\mathbb{R}}V,
\end{equation}%
noting that the major difference between $V_{\mathbb{C}}$ and $V_{\mathbb{R}%
} $ is that $V$ and $iV$ are disjoint both as direct sum subspaces $\left( 
\mathcal{V},i\mathcal{V}\right) ,$where 
\begin{equation}
\mathcal{V}=\left\{ \left( V,0\right) \right\} ;\,i\mathcal{V}=\left\{
\left( 0,iV\right) \right\}
\end{equation}%
and individual vector spaces $\left( V,iV\right) $, while $V$ generates
disjoint direct sum subspaces $\left( \mathcal{V}_{1},\mathcal{V}_{2}\right) 
$ 
\begin{equation}
\mathcal{V}_{1}=\left\{ \left( V,0\right) \right\} ;\,\mathcal{V}%
_{2}=\left\{ \left( 0,V\right) \right\} ,
\end{equation}%
but obviously $V$ is not disjoint to itself. As discussed later if vector
spaces $V_{1}$ and $V_{2}$ are disjoint then 
\begin{equation}
\mathcal{W}=W_{1}\oplus W_{2}=\mathcal{W}_{1}+\mathcal{W}_{2}\cong
W_{1}+W_{2},
\end{equation}%
giving 
\begin{equation}
V_{\mathbb{C}}=V\oplus _{\mathbb{R}}iV=\mathcal{V}+i\mathcal{V}\cong V+iV
\end{equation}%
and 
\begin{equation}
V_{\mathbb{R}}=V\oplus _{\mathbb{R}}V=\mathcal{V}_{1}+\mathcal{V}_{2}\ncong
V+V.
\end{equation}

\begin{definition}
The complex unit "$i"$ defines a real linear map $\mathfrak{R}_{V_{\mathbb{C}%
}}\left( i\right) \in \mathcal{L}_{\mathbb{R}}\left( V_{\mathbb{C}}\right)
\equiv \mathcal{L}_{\mathbb{R}}\left( V\oplus _{\mathbb{R}}iV\right) $ 
\begin{equation}
\mathfrak{R}_{V_{\mathbb{C}}}\left( i\right) :V\oplus _{\mathbb{R}%
}iV\rightarrow V\oplus _{\mathbb{R}}iV
\end{equation}%
by 
\begin{equation}
\mathfrak{R}_{V_{\mathbb{C}}}\left( i\right) \left( \left( \overset{}{%
v_{1},iv_{2}}\right) \right) =i\left( v_{1},iv_{2}\right) =\left(
-v_{2},iv_{1}\right)
\end{equation}%
belonging to $\mathcal{L}\left( V\oplus _{\mathbb{R}}iV\right) .$
\end{definition}

In order to distinguish between the complex unit $i$ and the linear map $%
\mathfrak{R}_{V_{\mathbb{C}}}\left( i\right) $, which is a representation of 
$i$ in the real form $V_{\mathbb{C}}$ generated by $V,$ we denote the linear
map representation of $i$ $\mathfrak{R}_{V_{\mathbb{C}}}\left( i\right) \in 
\mathcal{L}_{\mathbb{R}}\left( V_{\mathbb{C}}\right) $ by $\mathcal{J}_{V}$
i.e.%
\begin{equation}
\mathcal{J}_{V}=\mathfrak{R}_{V}\left( i\right)
\end{equation}%
giving

\begin{remark}
\begin{enumerate}
\item 
\begin{equation}
\mathcal{J}_{V}\left( v_{1},iv_{2}\right) =\left( -v_{2},iv_{1}\right) .
\end{equation}

\item 
\begin{align}
i\left\{ \alpha \left( v_{1},iv_{2}\right) +\beta \left( w_{1},iw_{2}\right)
\right\} & =i\left\{ \left( \alpha v_{1},i\alpha v_{2}\right) +\left( \beta
w_{1},i\beta w_{2}\right) \right\} =i\left\{ \left( \alpha v_{1}+\beta
w_{1},i\left\{ \alpha v_{2}+\beta w_{2}\right\} \right) \right\} \\
& =\left( -\left\{ \alpha v_{2}+\beta w_{2}\right\} ,i\left\{ \alpha
v_{1}+\beta w_{1}\right\} \right)
\end{align}

\item 
\begin{align}
i\left\{ \left( \alpha v_{1},iv_{2}\right) +\left( \beta w_{1},iw_{2}\right)
\right\} & =i\left( \alpha v_{1},i\alpha v_{2}\right) +i\left( \beta
w_{1},i\beta w_{2}\right) =\left( -\alpha v_{2},i\alpha v_{1}\right) +\left(
-\beta w_{2},i\beta w_{1}\right) \\
& =\left( -\left\{ \alpha v_{2}+\beta w_{2}\right\} ,i\left\{ \alpha
v_{1}+\beta w_{1}\right\} \right) \\
\alpha ,\beta & \in \mathbb{R}
\end{align}
\end{enumerate}
\end{remark}

The analogous real linear map in $V_{\mathbb{R}}$ 
\begin{equation}
\mathcal{J}_{V}:V_{\mathbb{R}}\rightarrow V_{\mathbb{R}},
\end{equation}%
which we denote by the same symbol $\mathcal{J}_{V},$ as its action is
essentially the same, is given by 
\begin{equation}
\mathcal{J}_{V}\left( v_{1},v_{2}\right) =\left( -v_{2},v_{1}\right)
\end{equation}%
and 
\begin{align}
v_{1}& =\func{Re}_{V}\left( v_{c}\right) \\
v_{2}& =\func{Im}_{V}\left( v_{c}\right) \\
v_{1},v_{2}& \in V,v_{c}\in V_{\mathbb{C}}.
\end{align}%
In particular its action on basis elements of $V\oplus _{\mathbb{R}}V$ is 
\begin{align}
\mathcal{J}_{V}\left( e_{j},0\right) & =\left( 0,e_{j}\right) \\
\mathcal{J}_{V}\left( 0,e_{j}\right) & =\left( -e_{j},0\right) .
\end{align}

\begin{definition}
\begin{equation}
\mathcal{V}=\left\{ \left. \left( e_{j},0\right) \right\vert 1\leq j\leq
r\right\} \cong V
\end{equation}%
and 
\begin{equation}
\mathcal{W}=\left\{ \left. \left( 0,e_{j}\right) \right\vert 1\leq j\leq
r\right\} =\mathcal{J}_{V}\mathcal{V}\cong W.
\end{equation}
\end{definition}

\begin{definition}
One can define the complex isomorphism $J_{V}\in \mathcal{L}\left( V\right) $
\begin{equation}
W=J_{V}V\cong \mathcal{J}_{V}\mathcal{V}=\mathcal{W}
\end{equation}%
i.e. $W$ is isomorphic to a specified factored form of $\mathcal{W}$ through
the map $\mathcal{J}_{V}$ and thus can be expressed as the factorized form $%
J_{V}V$ as the product of the factors $J_{V}V=\left\{ \left.
J_{V}v\right\vert v\in V\right\} .$ The complexification of $V$ can thus be
expressed as 
\begin{equation}
V_{\mathbb{C}}=V\oplus _{\mathbb{R}}iV=V\oplus _{\mathbb{R}}J_{V}V=\mathcal{V%
}+i\mathcal{V}=\mathcal{V}+\mathcal{J}_{V}\mathcal{V}\cong V+iV.
\end{equation}
\end{definition}

This should be compared to action of $\mathcal{J}_{V}$ in a real form of $V_{%
\mathbb{C}}$ i.e.%
\begin{align}
\mathcal{J}_{V}\left( e_{j},0\right) & =\left( 0,ie_{j}\right) =\left(
0,J_{V}e_{j}\right) \\
\mathcal{J}_{V}\left( 0,J_{V}e_{j}\right) & =\mathcal{J}_{V}\left(
0,ie_{j}\right) =\left( J_{V}J_{V}e_{j},0\right) =\mathcal{J}_{V}\mathcal{J}%
_{V}\left( e_{j},0\right) =\left( -e_{j},0\right) ,
\end{align}%
which allows one to directly define $J_{V}$ in $V$ and $iV\equiv J_{V}V$ by 
\begin{align}
J_{V}e_{j}& =ie_{j}\equiv f_{j} \\
J_{V}f_{j}& =J_{V}J_{V}e_{j}=-e_{j}.
\end{align}%
\emph{Note that }$V$ determines the real and imaginary parts of vectors
belonging to $V_{\mathbb{C}}.$ \emph{Different subspaces }$V$ \emph{%
determine different real and imaginary parts of a vector }$v_{c}\in V_{%
\mathbb{C}}$\emph{.}

\emph{Complexification can always be cast in the form }$V\oplus _{\mathbb{R}%
}J_{V}V,$ \emph{where }$V$\emph{\ and }$J_{V}V$\emph{\ are disjoint real
subspaces}, \emph{it does matter where }$J_{V}V$\emph{\ comes from i.e. it
can come from }$\mathcal{J}_{V}$ $\emph{defined}$ \emph{in} $V\oplus _{%
\mathbb{R}}V$\emph{\ or in general when it is defined in }$V\oplus _{\mathbb{%
R}}W.$ \qquad

The square of the linear map $\mathcal{J}_{V}$ is given by 
\begin{align}
i^{2}\left( v_{1},iv_{2}\right) & =i\left( -v_{2},iv_{1}\right) =\left(
-v_{1},-iv_{2}\right) =-\left( v_{1},iv_{2}\right) \\
& \equiv \mathcal{J}_{V}^{2}\left( v_{1},iv_{2}\right) =\mathcal{J}%
_{V}\left( -v_{2},iv_{1}\right) =\left( -v_{1},-iv_{2}\right) =-I_{V_{%
\mathbb{C}}}\left( v_{1},iv_{2}\right) ,
\end{align}%
so that 
\begin{equation}
\mathcal{J}_{V}^{2}=-I_{V_{\mathbb{C}}}.
\end{equation}

\begin{definition}
The complex unit "$i"$ also defines a real linear map belonging to $\mathcal{%
L}\left( V,iV\right) $ and $\mathcal{L}\left( iV,V\right) $ given by 
\begin{align}
i& :V\rightarrow iV \\
i\left( v_{1}\right) & =iv_{1}\in iV
\end{align}%
and%
\begin{align}
i& :iV\rightarrow V \\
i\left( iv_{1}\right) & =-v_{1}\in V,
\end{align}%
which we denote by $J_{V}$ in both cases, giving 
\begin{align}
J_{V}& :V\rightarrow J_{V}V \\
J_{V}\left( v_{1}\right) & =J_{V}v_{1}\in J_{V}V
\end{align}%
and%
\begin{align}
J_{V}& :J_{V}V\rightarrow V \\
J_{V}\left( J_{V}v_{1}\right) & =J_{V}^{2}v_{1}=-v_{1}\in V.
\end{align}
\end{definition}

Using the linear map notation one obtains 
\begin{equation}
\mathcal{J}_{V}\left( v_{1},J_{V}v_{2}\right) =\left(
-v_{2},J_{V}v_{1}\right) .
\end{equation}%
\emph{In particular this gives} 
\begin{align}
\mathcal{J}_{V}\left( v_{1},0\right) & =\left( 0,J_{V}v_{1}\right) \\
\mathcal{J}_{V}\left( 0,J_{V}v_{2}\right) & =\left( -v_{2},0\right) ,
\end{align}%
which is consistent with 
\begin{align}
\mathcal{J}_{V}^{2}\left( v_{1},0\right) & =\mathcal{J}_{V}\left(
0,J_{V}v_{1}\right) =\left( -v_{1},0\right) =-I_{V}\left( v_{1},0\right) \\
\mathcal{J}_{V}^{2}\left( 0,J_{V}v_{2}\right) & =\mathcal{J}_{V}\left(
-v_{2},0\right) =\left( 0,-J_{V}v_{2}\right) =-I_{V}\left(
0,J_{V}v_{2}\right) .
\end{align}

If the spaces $V$ and $W$ are real Hilbert space we can extend their inner
products to the direct sum $V\oplus W$ by 
\begin{equation}
\left( \left. \left( v_{1},v_{2}\right) \right\vert \left(
w_{1},w_{2}\right) \right) _{V\oplus W}=\left( \left. v_{1}\right\vert
w_{1}\right) _{V}+\left( \left. v_{2}\right\vert w_{2}\right) _{W}
\end{equation}
giving in the case $V\oplus_{\mathbb{R}}iV$ 
\begin{align}
\left( \left. \left( v_{1},iv_{2}\right) \right\vert \left(
w_{1},iw_{2}\right) \right) _{V\oplus_{\mathbb{R}}iV} & =\left( \left.
v_{1}\right\vert w_{1}\right) _{V}+\left( \left. iv_{2}\right\vert
iw_{2}\right) _{iV} \\
v_{1},v_{2},w_{1},w_{2} & \in V.
\end{align}
and equivalently%
\begin{align}
\left( \left. \left( v_{1},J_{V}v_{2}\right) \right\vert \left(
w_{1},J_{V}w_{2}\right) \right) _{V\oplus_{\mathbb{R}}J_{V}V} & =\left(
\left. v_{1}\right\vert w_{1}\right) _{V}+\left( \left.
J_{V}v_{2}\right\vert J_{V}w_{2}\right) _{J_{V}V} \\
v_{1},v_{2},w_{1},w_{2} & \in V.
\end{align}
This gives 
\begin{align}
\left( \left. \left( v_{1},0\right) \right\vert \mathcal{J}_{V}\left(
v_{1},0\right) \right) & =\left( \left. \left( v_{1},0\right) \right\vert
\left( 0,J_{V}v_{1}\right) \right) =0 \\
\left( \left. \left( 0,J_{V}v_{2}\right) \right\vert \mathcal{J}_{V}\left(
0,J_{V}v_{2}\right) \right) & =\left( \left. \left( 0,J_{V}v_{2}\right)
\right\vert \left( -v_{2},0\right) \right) =0 \\
\left( \left. \left( v_{1},0\right) \right\vert \mathcal{J}_{V}\left(
0,J_{V}v_{1}\right) \right) & =\left( \left. \left( v_{1},0\right)
\right\vert \left( -v_{1},0\right) \right) =-1 \\
\left( \left. \left( 0,J_{V}v_{1}\right) \right\vert \mathcal{J}_{V}\left(
v_{1},0\right) \right) & =\left( \left. \left( 0,J_{V}v_{1}\right)
\right\vert \left( 0,J_{V}v_{1}\right) \right) =1 \\
v_{1},v_{2},w_{1},w_{2} & \in V.
\end{align}
$\lceil$%
\begin{align}
\left( \left. \left( v_{1},0\right) \right\vert i\left( v_{1},0\right)
\right) & =\left( \left. \left( v_{1},0\right) \right\vert \left(
0,iv_{1}\right) \right) =0 \\
\left( \left. \left( 0,iv_{1}\right) \right\vert i\left( 0,iv_{1}\right)
\right) & =\left( \left. \left( 0,iv_{1}\right) \right\vert \left(
-v_{1},0\right) \right) =0 \\
\left( \left. \left( v_{1},0\right) \right\vert i\left( 0,iv_{1}\right)
\right) & =\left( \left. \left( v_{1},0\right) \right\vert \left(
-v_{1},0\right) \right) =-1 \\
\left( \left. \left( 0,iv_{1}\right) \right\vert i\left( v_{1},0\right)
\right) & =\left( \left. \left( 0,iv_{1}\right) \right\vert \left(
0,iv_{1}\right) \right) =1 \\
v_{1},v_{2},w_{1},w_{2} & \in V.
\end{align}
$\rfloor$

\begin{example}
The representation of $\mathcal{J}_{V}$ on the basis $\left\{ \left(
e_{1},0\right) ,\left( 0,J_{V}e_{1}\right) \right\} $ is given by 
\begin{align}
i\left( \left( e_{1},0\right) ,\left( 0,ie_{1}\right) \right) & =\left(
\left( 0,ie_{1}\right) ,\left( -e_{1},0\right) \right) =\left( \left(
e_{1},0\right) ,\left( 0,ie_{1}\right) \right) \left( 
\begin{array}{cc}
0 & 1 \\ 
-1 & 0%
\end{array}%
\right) \\
\mathcal{J}_{V}\left( \left( e_{1},0\right) ,\left( 0,J_{V}e_{1}\right)
\right) & =\left( \left( 0,J_{V}e_{1}\right) ,\left( -e_{1},0\right) \right)
=\left( \left( e_{1},0\right) ,\left( 0,J_{V}e_{1}\right) \right) \left( 
\begin{array}{cc}
0 & 1 \\ 
-1 & 0%
\end{array}%
\right) .
\end{align}%
The space $iV$ is orthogonal to $V$ i.e. 
\begin{equation}
\left( \left. a\right\vert ib\right) =\left( \left. ia\right\vert b\right)
=\left( \left. a\right\vert J_{V}b\right) =\left( \left. J_{V}a\right\vert
b\right) =0,
\end{equation}%
where 
\begin{align}
\left( \left. a\right\vert ib\right) & =\left( \left. \left( a,0\right)
\right\vert \left( 0,ib\right) \right) _{V\oplus _{\mathbb{R}}iV} \\
\left( \left. a\right\vert J_{V}b\right) & =\left( \left. \left( a,0\right)
\right\vert \left( 0,J_{V}b\right) \right) _{V\oplus _{\mathbb{R}}J_{V}V}
\end{align}%
and 
\begin{equation}
\left( \left. J_{V}^{t}a\right\vert b\right) \overset{def}{=}\left( \left.
a\right\vert J_{V}b\right) .
\end{equation}%
\emph{Note that }$J_{V}$ \emph{is NOT antisymmetric, while }$\mathcal{J}_{V}$
\emph{is}.
\end{example}

\begin{remark}
Also note in the direct sum 
\begin{equation}
V_{\oplus }=V_{1}\oplus V_{2}=\widetilde{V_{1}}+\widetilde{V_{2}}=\left\{
\left. \left( a,b\right) \right\vert a\in V_{1},b\in V_{2}\right\} ,
\end{equation}%
for spaces (real or complex) $\widetilde{V_{1}}$ and $\widetilde{V_{2}}$ are
always disjoint even if $V_{1}$ $=V_{2},$ where 
\begin{align}
V_{1}& \longleftrightarrow \left( V_{1},0\right) \in \widetilde{V_{1}}\text{
giving }V_{1}\cong \widetilde{V_{1}} \\
V_{2}& \longleftrightarrow \left( 0,V_{2}\right) \in \widetilde{V_{2}}\text{
giving }V_{2}\cong \widetilde{V_{2}}
\end{align}%
i.e. 
\begin{align}
\left( a,0\right) +\left( 0,b\right) & =\left( a,b\right) \in V_{1}\oplus
V_{2} \\
& \leftrightarrow a+b\text{ if }V_{1}\cap V_{2}=\left\{ 0\right\} .
\end{align}
\end{remark}

The space $V_{\mathbb{C}}$ has a Hermitian inner product defined by 
\begin{gather}
\left\langle \left. x\right\vert y\right\rangle =\left\langle \left.
a+ib\right\vert c+id\right\rangle =\left\langle \left. a\right\vert
c\right\rangle +\left\langle \left. b\right\vert d\right\rangle +i\left\{
\left\langle \left. a\right\vert d\right\rangle -\left\langle \left.
b\right\vert c\right\rangle \right\} \\
\overset{\triangle}{=}\left( \left. a\right\vert c\right) +\left( \left.
b\right\vert d\right) +i\left\{ \left( \left. a\right\vert d\right) -\left(
\left. b\right\vert c\right) \right\} =h\left( x,y\right) +ig\left(
x,y\right)
\end{gather}
where 
\begin{gather}
h\left( x,y\right) =\left( \left. a\right\vert c\right) +\left( \left.
b\right\vert d\right) \\
g\left( x,y\right) =\left( \left. a\right\vert d\right) -\left( \left.
b\right\vert c\right) \\
x,y\in V_{\mathbb{C}},\,a,b,c,d\in V,
\end{gather}
such that 
\begin{equation}
\left. \left\langle \left. a\right\vert c\right\rangle \right\vert _{V\times
V}=h\left( a,c\right) \equiv\left( \left. a\right\vert c\right) .
\end{equation}

\begin{remark}
$V$ is a \emph{real} \emph{maximal} \emph{isotropic} (\emph{Lagrangian})
subspace of $V_{\mathbb{C}}$ and $V_{\mathbb{R}}$ wrt $g$ as 
\begin{align}
g\left( x,y\right) & =\func{Im}\left\langle \left. x\right\vert
y\right\rangle =0\forall x,y\in V \\
\text{as }x& =\func{Re}_{V}x=a;\,y=\func{Re}_{V}y=c \\
\text{i.e.\thinspace }b& =0\,\&\,d=0.
\end{align}%
The subspace $V$ defines the real and imaginary parts $\left( a,b\right) $
of the vector $z.$
\end{remark}

The dimensions of these vector spaces are%
\begin{align}
\mathbb{R-}\dim\left( V\right) & =r \\
\mathbb{R-}\dim\left( V_{\mathbb{R}}\right) & =2r \\
\mathbb{C-}\dim\left( V_{\mathbb{C}}\right) & =r \\
\mathbb{R-}\dim\left( V_{\mathbb{C}}\right) & =2r.
\end{align}

One has the following 
\begin{equation}
V_{\mathbb{C}}=V\oplus W=V\oplus iV=V\oplus J_{V}V=\mathcal{V}_{1}+\mathcal{V%
}_{2}=\mathcal{V}_{1}+\mathcal{J}_{V_{1}}\mathcal{V}_{1}\cong V_{1}+iV_{1}
\end{equation}%
\begin{align}
V\oplus W & =\mathcal{V}+\mathcal{W}=\limfunc{linspan}\left\{ \left. \left(
v,w\right) \right\vert v_{1}\in V,w\in W\right\} =\limfunc{linspan}\left\{
\left. \left( v,0\right) ,\left( 0,w\right) \right\vert v\in V,w\in W\right\}
\\
V\oplus iV & =\limfunc{linspan}\left\{ \left. \left( v_{1},iv_{2}\right)
\right\vert v_{1},v_{2}\in V_{1}\right\} =\limfunc{linspan}\left\{ \left.
\left( v_{1},0\right) ,\left( 0,iv_{2}\right) \right\vert v_{1},v_{2}\in
V\right\} \\
V\oplus J_{V}V & =\limfunc{linspan}\left\{ \left. \left(
v_{1},J_{V}v_{2}\right) \right\vert v_{1},v_{2}\in V\right\} =\limfunc{%
linspan}\left\{ \left. \left( v_{1},0\right) ,\left( 0,J_{V}v_{2}\right)
\right\vert v_{1},v_{2}\in V\right\} \\
\mathcal{V}+\mathcal{J}_{V}\mathcal{W} & =\limfunc{linspan}\left\{ \left.
\left( v_{1},0\right) ,\mathcal{J}_{V}\left( 0,v_{2}\right) \right\vert
v_{1},v_{2}\in V\right\} .
\end{align}

External complexification of real vector space $V$ (usually equipped with
positive definite inner products and thus are real Hilbert spaces) is
equivalent to internal complexification of $V$ produced by a map 
\begin{equation}
\mathcal{J}_{V}:V\oplus V\rightarrow V\oplus V,
\end{equation}%
where 
\begin{equation}
\mathcal{J}_{V}^{2}=-I_{V\oplus V}
\end{equation}%
and if $V$ is a Hilbert space 
\begin{equation}
\mathcal{J}_{V}=-\mathcal{J}_{V}^{t}.
\end{equation}%
This allows one to define the map 
\begin{equation}
J_{V}:V\rightarrow J_{V}V
\end{equation}%
by 
\begin{align}
\left( 0,J_{V}e_{j}\right) & =\mathcal{J}_{V}\left( e_{j},0\right) =\left(
0,f_{j}\right) \\
J_{V}e_{j}& =f_{j} \\
\left( -e_{j},0\right) & =\mathcal{J}_{V}\left( 0,J_{V}e_{j}\right) =\left(
J_{V}^{2}e_{j},0\right) \\
J_{V}f_{j}& =-e_{j},
\end{align}%
such that 
\begin{equation}
J_{V}^{2}=-I_{V}.
\end{equation}%
The complexification is then given by 
\begin{equation}
V_{\mathbb{C}}=V\oplus J_{V}V.
\end{equation}

\subsubsection{Complexification of Real Hilbert Space II (Internal)}

Let $V_{\mathbb{R}}$ be a real Hilbert space of dimension $2r$ with inner
product $\left( \left. a\right\vert b\right) =\left( \left. b\right\vert
a\right) $ then internal complexifications can be defined by using the
orthogonal direct sum \emph{decompositions} 
\begin{equation}
V_{\mathbb{R}}=V\oplus V^{\bot}=\mathcal{V}+\mathcal{V^{\bot}},
\end{equation}
where 
\begin{equation}
\dim\left( V\right) =\dim\left( V^{\bot}\right) =r
\end{equation}
and orthogonality is defined in terms of the inner product in $V_{\mathbb{R}%
}.$

The vectors in $V^{\bot}$ and $V$ are disjoint can be of quite a different
nature, e.g. it is possible $V^{\bot}=V^{d},$ where $V^{d}$ is the dual
space of $V.$ As we display later in order to be rigorous and unambiguous
one must use the isomorphisms $\mathcal{V}\cong V$ and $\mathcal{V^{\bot}}%
\cong V^{\bot}$ giving 
\begin{equation}
V_{\mathbb{R}}=\mathcal{V}+\mathcal{V^{\bot}}=V\oplus V^{\bot}\cong
V+V^{\bot }.
\end{equation}

The complexification is accomplished by using the complex structure, which
can be viewed as belonging to three different spaces of mappings but denoted
by the same symbol, $J_{V}\in \mathcal{L}\left( V,V^{\bot }\right) ,\mathcal{%
L}\left( V^{\bot },V\right) $ and $\mathcal{L}\left( V+V^{\bot }\right) $ 
\begin{align}
J_{V}& :V\rightarrow V^{\bot } \\
J_{V}& :V^{\bot }\rightarrow V \\
J_{V}& :V+V^{\bot }\rightarrow V+V^{\bot }
\end{align}%
given by 
\begin{align}
J_{V}& =\dsum\limits_{1\leq j\leq r}\left\{ \left\vert e_{j}\right) \left(
e_{j+r}\right\vert -\left\vert e_{j+r}\right) \left( e_{j}\right\vert
\right\} =T_{VV^{\bot }}-T_{V^{\bot }V} \\
J_{V}^{2}& =-I_{2r}.
\end{align}%
Explicitly 
\begin{align}
J_{V}\left( a\right) & =a^{\bot } \\
J_{V}\left( b^{\bot }\right) & =-b
\end{align}%
The map $J_{V}$ induces the map $\mathcal{J}_{\mathcal{V}}$ in $V\oplus
V^{\bot }$ given by 
\begin{equation}
\mathcal{J}_{\mathcal{V}}\left( a,b^{\bot }\right) =\left( -b,a^{\bot
}\right) ;\quad \left( a,b^{\bot }\right) ,\left( -b,a^{\bot }\right) \in
V\oplus V^{\bot }=\mathcal{V}+\mathcal{V^{\bot }}
\end{equation}%
where we have used the isomorphisms 
\begin{align}
\mathcal{V}& =\left\{ \left. \left( a,0\right) \right\vert a\in V\right\}
\cong V \\
\mathcal{V^{\bot }}& =\left\{ \left. \left( 0,b^{\bot }\right) \right\vert
b^{\bot }\in V^{\bot }\right\} \cong V^{\bot }.
\end{align}%
The transition operators 
\begin{align}
T_{V^{\bot }V}& :V\rightarrow V^{\bot } \\
T_{VV^{\bot }}& :V^{\bot }\rightarrow V \\
T_{V^{\bot }V}& =-\dsum\limits_{1\leq j\leq r}\left\vert e_{j+r}\right)
\left( e_{j}\right\vert \\
T_{VV^{\bot }}& =\dsum\limits_{1\leq j\leq r}\left\vert e_{j}\right) \left(
e_{j+r}\right\vert
\end{align}%
are partial isometries%
\begin{align}
T_{V^{\bot }V}T_{VV^{\bot }}& =T_{V^{\bot }V}T_{VV^{\bot }}^{t}=P_{V^{\bot }}
\\
T_{VV^{\bot }}T_{V^{\bot }V}& =T_{VV^{\bot }}T_{VV^{\bot }}^{t}=P_{V} \\
T_{V^{\bot }V}& =T_{VV^{\bot }}^{t}.
\end{align}%
The complex structures, (associated to different subspaces $V$), have the
properties%
\begin{align}
J_{V}^{t}& =-\dsum\limits_{1\leq j\leq r}\left\{ \left\vert e_{j+r}\right)
\left( e_{j}\right\vert -\left\vert e_{j}\right) \left( e_{j+r}\right\vert
\right\} =-J_{V} \\
J_{V}^{t}J_{V}& =I_{2r}\Rightarrow J_{V}\in O\left( 2r,\mathbb{R}\right) \\
J_{V}^{2}& =-I_{2r},
\end{align}%
producing the complex space

\begin{equation}
V_{\mathbb{C}}=V\oplus J_{V}V=\mathcal{V}+\mathcal{J}_{\mathcal{V}}\mathcal{%
V\,}\,\,\mathcal{\cong \,}\,V+J_{V}V,
\end{equation}%
where complex multiplication is defined by 
\begin{equation}
i\left( a,b^{\bot }\right) \overset{def}{=}J_{V}\left( a,b^{\bot }\right)
=\left( -b^{\bot },a\right) .
\end{equation}%
giving 
\begin{equation}
V_{\mathbb{C}}\cong \left\{ a+J_{V}\left( b\right) \right\} =\left\{
a+ib\right\} ;a,b\in V.
\end{equation}%
Different $V$ 's produce the same vector space $V_{\mathbb{C}},$ \emph{but
not the same Hilbert space in general as they are not unitarily equivalent,
i.e. the induced inner products are not equivalent (isometric).}

\begin{definition}
Complex Conjugations $\sigma _{V}$%
\begin{equation}
\sigma _{V}=\dsum\limits_{1\leq j\leq r}\left\{ \left\vert e_{j}\right)
\left( e_{j}\right\vert -\left\vert e_{j+r}\right) \left( e_{j+r}\right\vert
\right\} =P_{V}-P_{V^{\bot }}=\sigma _{V}^{t}.
\end{equation}
\end{definition}

\begin{theorem}
Complexifications $J_{V}$ are in $1-1$ correspondence with complex
conjugations $\sigma _{V}$
\end{theorem}

\begin{proof}
\begin{align}
\sigma _{V}& =\dsum\limits_{1\leq j\leq r}\left\{ \left\vert e_{j}\right)
\left( e_{j}\right\vert -\left\vert e_{j+r}\right) \left( e_{j+r}\right\vert
\right\} =P_{V}-P_{V^{\bot }}=\sigma _{V}^{t} \\
\sigma _{V}^{2}& =\dsum\limits_{1\leq j\leq r}\left\{ \left\vert
e_{j}\right) \left( e_{j}\right\vert -\left\vert e_{j+r}\right) \left(
e_{j+r}\right\vert \right\} \dsum\limits_{1\leq k\leq r}\left\{ \left\vert
e_{k}\right) \left( e_{k}\right\vert -\left\vert e_{k+r}\right) \left(
e_{k+r}\right\vert \right\} \\
& =\dsum\limits_{1\leq j\leq r}\left\{ \left\vert e_{j}\right) \left(
e_{j}\right\vert +\left\vert e_{j+r}\right) \left( e_{j+r}\right\vert
\right\} =I_{2r},
\end{align}%
(thus $\sigma _{V}$ is an involution and belongs to $O\left( V_{\mathbb{R}},%
\mathbb{R}\right) $), while 
\begin{align}
J_{V}& =\dsum\limits_{1\leq j\leq r}\left\{ \left\vert e_{j}\right) \left(
e_{j+r}\right\vert -\left\vert e_{j+r}\right) \left( e_{j}\right\vert
\right\} =T_{VV^{\bot }}-T_{V^{\bot }V} \\
& =\dsum\limits_{1\leq j\leq r}\left\{ -S_{V}\left\vert e_{j+r}\right)
\left( e_{j+r}\right\vert +S_{V}\left\vert e_{j}\right) \left(
e_{j}\right\vert \right\} =S_{V}\left\{ \dsum\limits_{1\leq j\leq r}\left\{
\left\vert e_{j}\right) \left( e_{j}\right\vert -\left\vert e_{j+r}\right)
\left( e_{j+r}\right\vert \right\} \right\} =S_{V}\sigma _{V},
\end{align}%
where $S_{V}$ is invertible and has the properties 
\begin{align}
& S_{V}\left\vert e_{j}\right) =-\left\vert e_{j+r}\right) \\
& S_{V}\left\vert e_{j+r}\right) =-\left\vert e_{j}\right) \\
& \Rightarrow \left( e_{j+r}\left\vert S_{V}\right. e_{j}\right) =-1\text{
and }\left( e_{j}\left\vert S_{V}\right. e_{j}\right) =0 \\
& S_{V}^{2}\left\vert e_{j}\right) =-S_{V}\left\vert e_{j+r}\right)
=\left\vert e_{j}\right) \\
& S_{V}^{2}=I_{2r}.
\end{align}%
\emph{This shows that }$S_{V}$ 
\begin{equation}
S_{V}=-\dsum\limits_{1\leq j\leq r}\left\{ \left\vert e_{j}\right) \left(
e_{j+r}\right\vert +\left\vert e_{j+r}\right) \left( e_{j}\right\vert
\right\} =T_{VV^{\bot }}+T_{V^{\bot }V}=S_{V}^{t}
\end{equation}%
is also an involution, (but not a complex conjugation), and thus is its own
inverse, non singular and belongs to $O\left( V_{\mathbb{R}},\mathbb{R}%
\right) $.
\end{proof}

\begin{corollary}
The maps $\sigma _{V}$ ,$S_{V}$ and $J_{V}$ anti commute
\end{corollary}

\begin{proof}
First we observe that 
\begin{gather}
J_{V}=\dsum\limits_{1\leq j\leq r}\left\{ \left\vert e_{j}\right) \left(
e_{j+r}\right\vert -\left\vert e_{j+r}\right) \left( e_{j}\right\vert
\right\} =-J_{V}^{t} \\
J_{V}^{2}=-I_{V} \\
\sigma _{V}=\dsum\limits_{1\leq j\leq r}\left\{ \left\vert e_{j}\right)
\left( e_{j}\right\vert -\left\vert e_{j+r}\right) \left( e_{j+r}\right\vert
\right\} =\sigma _{V}^{t} \\
\sigma _{V}^{2}=I_{V} \\
S_{V}=-\dsum\limits_{1\leq j\leq r}\left\{ \left\vert e_{j}\right) \left(
e_{j+r}\right\vert +\left\vert e_{j+r}\right) \left( e_{j}\right\vert
\right\} =S_{V}^{t} \\
S_{V}^{2}=I_{V}
\end{gather}%
and show that $\left\{ \sigma _{V},S_{V},J_{V}\right\} $ anti commute by the
following: 
\begin{gather}
J_{V}\sigma _{V}=\dsum\limits_{1\leq j\leq r}\left\{ \left\vert e_{j}\right)
\left( e_{j+r}\right\vert -\left\vert e_{j+r}\right) \left( e_{j}\right\vert
\right\} \dsum\limits_{1\leq j\leq r}\left\{ \left\vert e_{j}\right) \left(
e_{j}\right\vert -\left\vert e_{j+r}\right) \left( e_{j+r}\right\vert
\right\} =-\dsum\limits_{1\leq j\leq r}\left\{ \left\vert e_{j}\right)
\left( e_{j+r}\right\vert +\left\vert e_{j+r}\right) \left( e_{j}\right\vert
\right\} \\
\sigma _{V}J_{V}=\dsum\limits_{1\leq j\leq r}\left\{ \left\vert e_{j}\right)
\left( e_{j}\right\vert -\left\vert e_{j+r}\right) \left( e_{j+r}\right\vert
\right\} \dsum\limits_{1\leq j\leq r}\left\{ \left\vert e_{j}\right) \left(
e_{j+r}\right\vert -\left\vert e_{j+r}\right) \left( e_{j}\right\vert
\right\} =\dsum\limits_{1\leq j\leq r}\left\{ \left\vert e_{j}\right) \left(
e_{j+r}\right\vert +\left\vert e_{j+r}\right) \left( e_{j}\right\vert
\right\} \\
J_{V}\sigma _{V}=-\sigma _{V}J_{V} \\
S_{V}\sigma _{V}=-\dsum\limits_{1\leq j\leq r}\left\{ \left\vert
e_{j}\right) \left( e_{j+r}\right\vert +\left\vert e_{j+r}\right) \left(
e_{j}\right\vert \right\} \dsum\limits_{1\leq j\leq r}\left\{ \left\vert
e_{j}\right) \left( e_{j}\right\vert -\left\vert e_{j+r}\right) \left(
e_{j+r}\right\vert \right\} =\dsum\limits_{1\leq j\leq r}\left\{ \left\vert
e_{j}\right) \left( e_{j+r}\right\vert -\left\vert e_{j+r}\right) \left(
e_{j}\right\vert \right\} \\
\sigma _{V}S_{V}=-\dsum\limits_{1\leq j\leq r}\left\{ \left\vert
e_{j}\right) \left( e_{j}\right\vert -\left\vert e_{j+r}\right) \left(
e_{j+r}\right\vert \right\} \dsum\limits_{1\leq j\leq r}\left\{ \left\vert
e_{j}\right) \left( e_{j+r}\right\vert +\left\vert e_{j+r}\right) \left(
e_{j}\right\vert \right\} =\dsum\limits_{1\leq j\leq r}\left\{ \left\vert
e_{j+r}\right) \left( e_{j}\right\vert -\left\vert e_{j}\right) \left(
e_{j+r}\right\vert \right\} \\
S_{V}\sigma _{V}=-\sigma _{V}S_{V} \\
J_{V}S_{V}=-\dsum\limits_{1\leq j\leq r}\left\{ \left\vert e_{j}\right)
\left( e_{j+r}\right\vert -\left\vert e_{j+r}\right) \left( e_{j}\right\vert
\right\} \dsum\limits_{1\leq j\leq r}\left\{ \left\vert e_{j}\right) \left(
e_{j+r}\right\vert +\left\vert e_{j+r}\right) \left( e_{j}\right\vert
\right\} =-\dsum\limits_{1\leq j\leq r}\left\{ \left\vert e_{j}\right)
\left( e_{j}\right\vert -\left\vert e_{j+r}\right) \left( e_{j+r}\right\vert
\right\} \\
S_{V}J_{V}=-\dsum\limits_{1\leq j\leq r}\left\{ \left\vert e_{j}\right)
\left( e_{j+r}\right\vert +\left\vert e_{j+r}\right) \left( e_{j}\right\vert
\right\} \dsum\limits_{1\leq j\leq r}\left\{ \left\vert e_{j}\right) \left(
e_{j+r}\right\vert -\left\vert e_{j+r}\right) \left( e_{j}\right\vert
\right\} =-\dsum\limits_{1\leq j\leq r}\left\{ -\left\vert e_{j}\right)
\left( e_{j}\right\vert +\left\vert e_{j+r}\right) \left( e_{j+r}\right\vert
\right\} \\
J_{V}S_{V}=-S_{V}J_{V}.
\end{gather}%
Also%
\begin{eqnarray}
S_{V}J_{V}S_{V} &=&-J_{V} \\
S_{V}\sigma _{V}S_{V} &=&-\sigma _{V} \\
\sigma _{V}S_{V}\sigma _{V} &=&-S_{V} \\
\sigma _{V}J_{V}\sigma _{V} &=&-J_{V} \\
J_{V}\sigma _{V}J_{V} &=&-\sigma _{V} \\
J_{V}S_{V}J_{V} &=&-S_{V}.
\end{eqnarray}
\end{proof}

\section{Cartan Decomposition}

The polar decomposition factors the general linear group $GL\left( V_{%
\mathbb{R}},\mathbb{R}\right) $ and corresponds to the \emph{Cartan
Decomposition }of\emph{\ the Lie algebra }$gl\left( V_{\mathbb{R}},\mathbb{R}%
\right) $ i.e. 
\begin{equation}
gl\left( V_{\mathbb{R}},\mathbb{R}\right) =\mathfrak{l}\oplus \mathfrak{p,}
\end{equation}%
where 
\begin{eqnarray}
\left[ \mathfrak{l,l}\right] &\subseteq &\mathfrak{l} \\
\left[ \mathfrak{l,p}\right] &\subseteq &\mathfrak{p} \\
\left[ \mathfrak{p,p}\right] &\subseteq &\mathfrak{l,}
\end{eqnarray}%
and $\mathfrak{l}$\emph{\ are antisymmetric} and $\mathfrak{p}$\emph{\ are
symmetric}.

\begin{itemize}
\item $A\in \mathfrak{p}$ and $B\in \mathfrak{p}$ symmetric then $\left[ A,B%
\right] \in \mathfrak{l}$ i.e. is antisymmetric%
\begin{equation}
\left[ A,B\right] ^{t}=-\left[ A,B\right]
\end{equation}

\item $A\in \mathfrak{l}$ and $B\in \mathfrak{l}$ antisymmetric then $\left[
A,B\right] \in \mathfrak{l}$ i.e. is antisymmetric%
\begin{equation}
\left[ A,B\right] ^{t}=-\left[ A,B\right]
\end{equation}

\item $A\in \mathfrak{p}$ symmetric and $B\in \mathfrak{l}$ antisymmetric
then $\left[ A,B\right] \in \mathfrak{p}$ i.e. is symmetric%
\begin{equation}
\left[ A,B\right] ^{t}=-BA+AB=\left[ A,B\right] .
\end{equation}
\end{itemize}

\begin{definition}
One can define the $\emph{Adjoint}$ $\emph{Representation}$ $ad$ of the Lie
algebra $gl\left( V_{\mathbb{R}},\mathbb{R}\right) $ and its \emph{Universal
Enveloping Algebra} $ugl\left( V_{\mathbb{R}},\mathbb{R}\right) $ by 
\begin{eqnarray}
ad &:&ugl\left( V_{\mathbb{R}},\mathbb{R}\right) \rightarrow \mathcal{B}%
\left( ugl\left( V_{\mathbb{R}},\mathbb{R}\right) \right) \\
ad\left( X\right) \left( A\right) &=&\left[ X,A\right] \,;\,X,A\in ugl\left(
V_{\mathbb{R}},\mathbb{R}\right) .
\end{eqnarray}
\end{definition}

This is a representation as the \emph{Jacobi Identity} 
\begin{equation}
\left[ x,\left[ y,z\right] \right] +\left[ y,\left[ z,x\right] \right] +%
\left[ z,\left[ x,y\right] \right] =0
\end{equation}%
gives%
\begin{equation}
\left( \overset{}{ad\left( x\right) ad\left( y\right) }\right) \left(
z\right) =ad\left( \left[ x,y\right] \right) \left( z\right) .
\end{equation}

In the present case $G=$ $GL\left( V_{\mathbb{R}},\mathbb{R}\right) $, its
Lie algebra $\mathcal{L}_{G}=gl\left( V_{\mathbb{R}},\mathbb{R}\right) $ and 
\begin{eqnarray}
gl\left( V_{\mathbb{R}},\mathbb{R}\right) &=&\mathcal{B}\left( V_{\mathbb{R}%
}\right) \text{ equipped with the Lie product }\left[ ,\right] \\
ugl\left( V_{\mathbb{R}},\mathbb{R}\right) &=&\mathcal{B}\left( V_{\mathbb{R}%
}\right) \text{ equipped with the usual operator product.}
\end{eqnarray}

\begin{definition}
One can define an inner product $\left\langle \left. {}\right\vert
\right\rangle _{K}\equiv B\left( ,\right) $ called the\textbf{\ Killing Form}%
\emph{\ }by\emph{\ }%
\begin{equation}
\left\langle \left. X\right\vert Y\right\rangle _{K}=\func{Tr}\left\{
ad\left( X\right) ad\left( Y\right) \right\} .
\end{equation}
\end{definition}

\begin{remark}
A Lie algebra is \emph{semisimple }$\emph{\limfunc{Iff}}$ the Killing form
on that Lie algebra is \emph{non degenerate, though it may be non definite}.
\end{remark}

\begin{definition}
An involution, $\Theta ,$ is a \emph{Cartan Involution} if 
\begin{equation}
-\left\langle \left. X\right\vert \Theta \left( X\right) \right\rangle
_{K}>0\,\forall X\in \mathcal{L}_{G},
\end{equation}%
so that
\end{definition}

\begin{definition}
\begin{equation}
B_{\Theta }\left( X,Y\right) =-\left\langle \left. X\right\vert \Theta
\left( Y\right) \right\rangle _{K}
\end{equation}%
is a positive definite bilinear form. The \emph{Cartan Involution} $\Theta $
defines
\end{definition}

\begin{definition}
the Cartan decomposition for semisimple Lie algebras by%
\begin{eqnarray}
\Theta \left( \mathfrak{l}\right) &=&\mathfrak{l} \\
\Theta \left( \mathfrak{p}\right) &=&-\mathfrak{p.}
\end{eqnarray}
\end{definition}

\begin{remark}
$\left\langle \left. {}\right\vert \right\rangle _{K}$ is negative definite
on $\mathfrak{l}$ and positive definite on $\mathfrak{p.}$If $\mathcal{L}%
_{G} $ is the Lie algebra of a compact semisimple group $G$ then $%
\left\langle \left. {}\right\vert \right\rangle _{K}$ is negative definite
on $\mathcal{L}_{G}$ and thus the identity operator and 
\begin{equation}
\mathcal{L}_{G}=\mathfrak{l}
\end{equation}%
and 
\begin{equation}
\mathfrak{p}=\left\{ \phi \right\} .
\end{equation}
\end{remark}

In the case of $gl\left( V_{\mathbb{R}},\mathbb{R}\right) $ the Cartan
involution $\Theta \in \mathcal{B}\left( \mathcal{B}\left( V_{\mathbb{R}%
}\right) \right) $ and is given by 
\begin{equation}
\Theta \left( X\right) =-X^{t}.
\end{equation}%
and 
\begin{equation}
\mathfrak{l}=so\left( V_{\mathbb{R}},\mathbb{R}\right)
\end{equation}
is the real Lie algebra of skew-symmetric operators giving 
\begin{equation}
SO\left( V_{\mathbb{R}},\mathbb{R}\right) =\exp \left\{ so\left( V_{\mathbb{R%
}},\mathbb{R}\right) \right\}
\end{equation}%
and%
\begin{equation}
\mathfrak{p}=so\left( V_{\mathbb{R}},\mathbb{R}\right) ^{\bot },
\end{equation}
(orthogonality is defined by the Killing form), is the subspace of symmetric
operators, (a Jordan algebra), that produces the cone of nonsingular
positive operators under the exponential map 
\begin{equation}
\mathcal{B}\left( V_{\mathbb{R}}\right) _{+}=\exp \left\{ so\left( V_{%
\mathbb{R}},\mathbb{R}\right) ^{\bot }\right\} .
\end{equation}

The subspace of symmetric operators contains a maximal\textbf{\ abelian
subspace }of real operators that map onto a maximal abelian group of
positive operators (maximal torus). In general the exponential map is a
diffeomorphism from $so\left( V_{\mathbb{R}},\mathbb{R}\right) ^{\bot }$
onto the space of nonsingular positive definite operators. Up to this
exponential map, the global \emph{Cartan decomposition} is the \emph{polar
decomposition} of an operator$.$

\subsection{Invariance Groups}

All conjugations and complex structures are generated by $t\in O\left( V_{%
\mathbb{R}},\mathbb{R}\right) $ acting on a reference subspace $V,$where 
\begin{equation}
V_{\mathbb{R}}=V\oplus V^{\bot },
\end{equation}%
i.e. 
\begin{align}
t\sigma _{V}t^{t}& =\dsum\limits_{1\leq j\leq r}\left\{ \left\vert
te_{j}\right) \left( te_{j}\right\vert -\left\vert te_{j+r}\right) \left(
te_{j+r}\right\vert \right\} \\
tJ_{V}t^{t}& =\dsum\limits_{1\leq j\leq r}\left\{ \left\vert te_{j}\right)
\left( te_{j+r}\right\vert -\left\vert te_{j+r}\right) \left(
te_{j}\right\vert \right\} ,
\end{align}%
as all possible conb's of $V_{\mathbb{R}}$ are produced by such
transformations. Thus the preceeding invariance groups must be subgroups of $%
O\left( V_{\mathbb{R}},\mathbb{R}\right) $ (assuming $r<\infty )$.

\begin{remark}
A real antisymmetric bilinear form $g$ is given by 
\begin{equation}
g\left( v_{1},v_{2}\right) =\left( \left. \overset{}{g}\right\vert
v_{1}\wedge v_{2}\right) ,
\end{equation}%
where 
\begin{equation}
g=\dsum\limits_{1\leq j\leq r}\alpha _{j}\left\vert e_{j}e_{j+r}\right)
;\alpha _{j}\geq 0;1\leq j\leq 2r.
\end{equation}
\end{remark}

\begin{remark}
The real antisymmetric bilinear form $g$ corresponds to a real antisymmetric
operator $G$ 
\begin{equation}
G=\dsum\limits_{1\leq j\leq r}\alpha _{j}\left( \left\vert e_{j}\right)
\left( e_{j+r}\right\vert -\left\vert e_{j+r}\right) \left( e_{j}\right\vert
\right) .
\end{equation}
$G$ has the properties 
\begin{equation*}
G=-G^{t}
\end{equation*}%
and 
\begin{align}
G^{2}& =\dsum\limits_{1\leq j,k\leq r}\alpha _{j}\left\{ \left\vert
e_{j}\right) \left( e_{j+r}\right\vert -\left\vert e_{j+r}\right) \left(
e_{j}\right\vert \right\} \alpha _{k}\left\{ \left\vert e_{k}\right) \left(
e_{k+r}\right\vert -\left\vert e_{k+r}\right) \left( e_{k}\right\vert
\right\} \\
& =-\dsum\limits_{1\leq j\leq r}\alpha _{j}^{2}\left\{ \left\vert
e_{j}\right) \left( e_{j}\right\vert +\left\vert e_{j+r}\right) \left(
e_{j+r}\right\vert \right\} .
\end{align}
\end{remark}

If $G$ is a complexification then 
\begin{equation}
G^{2}=-I_{2r}\Rightarrow \alpha _{j}^{2}=1;1\leq j\leq 2r,
\end{equation}%
which coupled with the $\alpha _{j}>0$ gives that 
\begin{equation}
\alpha _{j}=1;1\leq j\leq 2r.
\end{equation}%
Thus if $G$ is a complexification the corresponding $g$ it must have the
form 
\begin{equation}
g=\dsum\limits_{1\leq j\leq r}\left\vert e_{j}e_{j+r}\right) \leftrightarrow
\dsum\limits_{1\leq j\leq r}\left\{ \left\vert e_{j}\right) \left(
e_{j+r}\right\vert -\left\vert e_{j+r}\right) \left( e_{j}\right\vert
\right\} =G.
\end{equation}

\begin{remark}
All geminals can be generated by 
\begin{eqnarray}
g\left( \mathbb{\lambda },\mathbb{\zeta }\right) &=&\dsum\limits_{1\leq
j\leq r}\alpha _{e}{}_{j}\left\vert f_{j}f_{j+r}\right) =\dsum\limits_{1\leq
j\leq r}\left( {\LARGE \wedge }^{2}O\left( \mathbb{\lambda }\right) \right)
\left( {\LARGE \wedge }^{2}\alpha _{e}\right) \left( \mathbb{\zeta }\right)
\left\vert e_{j}e_{j+r}\right) \\
&=&\dsum\limits_{1\leq j\leq r}\left( {\huge \wedge }^{2}O\left( \mathbb{%
\lambda }\right) \right) \alpha _{e}\left( f\right) _{j}\left\vert
e_{j}e_{j+r}\right) =\dsum\limits_{1\leq j\leq r}\alpha _{e}{}_{j}\left\vert
O\left( \mathbb{\lambda }\right) e_{j}O\left( \mathbb{\lambda }\right)
e_{j+r}\right) .
\end{eqnarray}%
and all antisymmetric operarors by%
\begin{equation}
G\left( \mathbb{\lambda },\mathbb{\zeta }\right) =\dsum\limits_{1\leq j\leq
r}\left\{ \left\vert f_{j}\right) \alpha _{e}{}_{j}\left( f_{j+r}\right\vert
-\left\vert f_{j+r}\right) \alpha _{e}{}_{j}\left( f_{j}\right\vert \right\}
=\dsum\limits_{1\leq j\leq r}O\left( \mathbb{\lambda }\right) \alpha
_{e}\left( \mathbb{\zeta }\right) ^{\frac{1}{2}}\left\{ \left\vert
e_{j}\right) \left( e_{j+r}\right\vert -\left\vert e_{j+r}\right) \left(
e_{j}\right\vert \right\} \alpha _{e}\left( \mathbb{\zeta }\right) ^{\frac{1%
}{2}}O\left( \mathbb{\lambda }\right) ^{t}.
\end{equation}
\end{remark}

In general the transformations between complex units are 
\begin{align}
& \left( tJ_{V}t^{t}\right) ^{2}=\dsum\limits_{1\leq j\leq r}\left\{
\left\vert te_{j}\right) \left( te_{j+r}\right\vert -\left\vert
te_{j+r}\right) \left( te_{j}\right\vert \right\} \dsum\limits_{1\leq j\leq
r}\left\{ \left\vert te_{j}\right) \left( te_{j+r}\right\vert -\left\vert
te_{j+r}\right) \left( te_{j}\right\vert \right\} \\
& =\dsum\limits_{1\leq j,k\leq r}\left\{ \left\vert te_{j}\right) \left(
te_{j+r}\right\vert -\left\vert te_{j+r}\right) \left( te_{j}\right\vert
\right\} \left\{ \left\vert te_{k}\right) \left( te_{k+r}\right\vert
-\left\vert te_{k+r}\right) \left( te_{k}\right\vert \right\} \\
& =\dsum\limits_{1\leq j,k\leq r}\left\{ \left( \left\vert te_{j}\right)
\left( te_{j+r}\right\vert -\left\vert te_{j+r}\right) \left(
te_{j}\right\vert \right) \left( \left\vert te_{k}\right) \left(
te_{k+r}\right\vert -\left\vert te_{k+r}\right) \left( te_{k}\right\vert
\right) \right\} \\
& =\dsum\limits_{1\leq j,k\leq r}\left\{ \left( \left\vert te_{j}\right)
\left( e_{j+r}\right\vert \left\vert t^{t}te_{k}\right) \left(
te_{k+r}\right\vert -\left\vert te_{j+r}\right) \left( e_{j}\right\vert
\left\vert t^{t}te_{k}\right) \right) \right. \\
& \left. \times \left( te_{k+r}\right\vert -\left\vert te_{j}\right) \left(
e_{j+r}\right\vert \left\vert t^{t}te_{k+r}\right) \left( te_{k}\right\vert
+\left\vert te_{j+r}\right) \left( e_{j}\right\vert \left\vert
t^{t}te_{k+r}\right) \left( te_{k}\right\vert )\right\} \\
& =\dsum\limits_{1\leq j,k\leq r}\left\{ -\left\vert te_{j}\right) \left(
e_{j+r}\right\vert \left\vert t^{t}te_{k+r}\right) \left( te_{k}\right\vert
-\left\vert te_{j+r}\right) \left( e_{j}\right\vert \left\vert
t^{t}te_{k}\right) \left( te_{k+r}\right\vert +\left\vert te_{j}\right)
\left( e_{j+r}\right\vert \left\vert t^{t}te_{k}\right) \left(
te_{k+r}\right\vert +\left\vert te_{j+r}\right) \left( e_{j}\right\vert
\left\vert t^{t}te_{k+r}\right) \left( te_{k}\right\vert \right\} \\
& =-t\dsum\limits_{1\leq j,k\leq r}\left\{ \left\vert e_{j}\right) \left(
e_{k}\right\vert \left( e_{j+r}\right\vert \left\vert t^{t}te_{k+r}\right)
+\left\vert e_{j+r}\right) \left( e_{k+r}\right\vert \left( e_{j}\right\vert
\left\vert t^{t}te_{k}\right) -\left\vert e_{j}\right) \left(
e_{k+r}\right\vert \left( e_{j+r}\right\vert \left\vert t^{t}te_{k}\right)
-\left\vert e_{j+r}\right) \left( e_{k}\right\vert \left( e_{j}\right\vert
\left\vert t^{t}te_{k+r}\right) \right\} t^{t}.
\end{align}%
$\left( tJ_{V}t^{t}\right) ^{2}$ can be succinctly expressed as 
\begin{align}
tJ_{V}t^{t}=& \left\vert te_{1},\cdots ,te_{2r}\right) \left( te_{1},\cdots
,te_{2r}\right\vert \left( 
\begin{array}{cc}
0 & I_{r} \\ 
-I_{r} & 0%
\end{array}%
\right) \left\vert te_{1},\cdots ,te_{2r}\right) \left( te_{1},\cdots
,te_{2r}\right\vert \\
& \times \left\vert te_{1},\cdots ,te_{2r}\right) \left( te_{1},\cdots
,te_{2r}\right\vert \left( 
\begin{array}{cc}
0 & I_{r} \\ 
-I_{r} & 0%
\end{array}%
\right) \left\vert te_{1},\cdots ,te_{2r}\right) \left( te_{1},\cdots
,te_{2r}\right\vert .
\end{align}%
In order that $tJ_{V}t^{t}$ is a complexification one has the constraint 
\begin{equation}
\left( tJ_{V}t^{t}\right) ^{2}=-I_{2r},
\end{equation}%
thus $\left( tJ_{V}t^{t}\right) \in O\left( 2r,\mathbb{R}\right) $ i.e. 
\begin{equation}
\left( tJ_{V}t^{t}\right) \left( tJ_{V}t^{t}\right)
^{t}=tJ_{V}t^{t}tJ_{V}^{t}t^{t}=I_{2r}.
\end{equation}%
Hence \label{03/04/2026}%
\begin{align}
& t\dsum\limits_{1\leq j,k\leq r}\left\{ \left\vert e_{j}\right) \left(
e_{k}\right\vert \left( \left. e_{j+r}\right\vert t^{t}te_{k+r}\right)
+\left\vert e_{j+r}\right) \left( e_{k+r}\right\vert \left( \left.
e_{j}\right\vert t^{t}te_{k}\right) -\left\vert e_{j}\right) \left(
e_{k+r}\right\vert \left( \left. e_{j+r}\right\vert t^{t}te_{k}\right)
-\left\vert e_{j+r}\right) \left( e_{k}\right\vert \left( \left.
e_{j}\right\vert t^{t}te_{k+r}\right) \right\} t^{t} \\
& =\dsum\limits_{1\leq j\leq r}\left\{ \left\vert e_{j}\right) \left(
e_{j}\right\vert +\left\vert e_{j+r}\right) \left( e_{j+r}\right\vert
\right\} =I_{2r}.
\end{align}

\begin{remark}
In the following one should observe that 
\begin{equation}
O\left( r,r,\mathbb{R}\right) \supset U\left( r,r\right)
\end{equation}%
\textbf{NOT vice versa i.e. }%
\begin{equation}
U\left( r,r\right) \not\supset O\left( r,r,\mathbb{R}\right)
\end{equation}%
\textbf{unlike the compact case} 
\begin{equation}
U\left( 2r\right) \supset O\left( 2r,\mathbb{R}\right)
\end{equation}
\end{remark}

\begin{remark}
One should observe that
\end{remark}

\begin{enumerate}
\item 
\begin{gather}
t\in Sp\left( 2r,\mathbb{R}\right) \supset O\left( r,r,\mathbb{R}\right)
\supset U\left( r,r\right) \\
\Downarrow \\
Sp\left( 2r,\mathbb{R}\right) \cap O\left( r,r,\mathbb{R}\right) =O\left(
r,r,\mathbb{R}\right) .
\end{gather}%
and 
\begin{equation}
O\left( r,r,\mathbb{R}\right) \cap U\left( r,r\right) =U\left( r,r\right)
\end{equation}

\item 
\begin{gather}
t\in Sp\left( 2r,\mathbb{R}\right) \supset O\left( r,r,\mathbb{R}\right)
\supset U\left( r,r\right) \\
\Downarrow \\
Sp\left( 2r,\mathbb{R}\right) \cap U\left( r,r\right) =U\left( r,r\right) .
\end{gather}%
and 
\begin{equation}
O\left( r,r,\mathbb{R}\right) \cap U\left( r,r\right) =U\left( r,r\right)
\end{equation}

\item 
\begin{equation}
O\left( r,r,\mathbb{R}\right) \cap U\left( r,r\right) =U\left( r,r\right) .
\end{equation}

\item 
\begin{equation}
Sp\left( 2r,\mathbb{R}\right) \cap O\left( 2r,\mathbb{R}\right) =U\left(
r\right) .
\end{equation}
\end{enumerate}

\begin{remark}
What are 
\begin{equation}
\left. Sp\left( 2r,\mathbb{R}\right) \right/ O\left( r,r,\mathbb{R}\right)
\end{equation}%
and%
\begin{equation}
\left. O\left( r,r,\mathbb{R}\right) \right/ \left( O\left( r,\mathbb{R}%
\right) \times O\left( r,\mathbb{R}\right) \right) ?
\end{equation}
\end{remark}

\subparagraph{Example}

If $t\in Sp\left( 2,\mathbb{R}\right) $%
\begin{align}
\left( 
\begin{array}{cc}
a & b \\ 
c & d%
\end{array}%
\right) \left( 
\begin{array}{cc}
0 & 1 \\ 
-1 & 0%
\end{array}%
\right) \left( 
\begin{array}{cc}
a^{t} & c^{t} \\ 
b^{t} & d^{t}%
\end{array}%
\right) & =\left( 
\begin{array}{cc}
0 & 1 \\ 
-1 & 0%
\end{array}%
\right) \\
\left( 
\begin{array}{cc}
-b & a \\ 
-d & c%
\end{array}%
\right) \left( 
\begin{array}{cc}
a^{t} & c^{t} \\ 
b^{t} & d^{t}%
\end{array}%
\right) & =\left( 
\begin{array}{cc}
ab^{t}-ba^{t} & ad^{t}-bc^{t} \\ 
cb^{t}-da^{t} & cd^{t}-dc^{t}%
\end{array}%
\right) =\left( 
\begin{array}{cc}
0 & 1 \\ 
-1 & 0%
\end{array}%
\right)
\end{align}%
giving 
\begin{equation}
ad=1+bc.
\end{equation}%
For example 
\begin{align}
\left( 
\begin{array}{cc}
2 & 7 \\ 
1 & 4%
\end{array}%
\right) & \in Sp\left( 2,\mathbb{R}\right) , \\
\text{e.g. }\left( 2\right) \left( 4\right) & =1+\left( 7\right) \left(
1\right)
\end{align}%
i.e. 
\begin{equation}
\left( 
\begin{array}{cc}
2 & 7 \\ 
1 & 4%
\end{array}%
\right) \left( 
\begin{array}{cc}
0 & 1 \\ 
-1 & 0%
\end{array}%
\right) \left( 
\begin{array}{cc}
2 & 1 \\ 
7 & 4%
\end{array}%
\right) =\left( 
\begin{array}{cc}
-7 & 2 \\ 
-4 & 1%
\end{array}%
\right) \left( 
\begin{array}{cc}
2 & 1 \\ 
7 & 4%
\end{array}%
\right) =\left( 
\begin{array}{cc}
0 & 1 \\ 
-1 & 0%
\end{array}%
\right) ,
\end{equation}%
which corresponds to the symplectic condition 
\begin{equation}
tJ_{V}t^{t}=J_{V}.
\end{equation}%
This gives 
\begin{align}
& \left( 
\begin{array}{cc}
0 & -1 \\ 
1 & 0%
\end{array}%
\right) \left( 
\begin{array}{cc}
2 & 7 \\ 
1 & 4%
\end{array}%
\right) \left( 
\begin{array}{cc}
0 & 1 \\ 
-1 & 0%
\end{array}%
\right) \left( 
\begin{array}{cc}
2 & 1 \\ 
7 & 4%
\end{array}%
\right) =\left( 
\begin{array}{cc}
-1 & -4 \\ 
2 & 7%
\end{array}%
\right) \left( 
\begin{array}{cc}
7 & 4 \\ 
-2 & -1%
\end{array}%
\right) \\
& =\left( 
\begin{array}{cc}
0 & -1 \\ 
1 & 0%
\end{array}%
\right) \left( 
\begin{array}{cc}
0 & 1 \\ 
-1 & 0%
\end{array}%
\right) =\left( 
\begin{array}{cc}
1 & 0 \\ 
0 & 1%
\end{array}%
\right) ,
\end{align}%
which corresponds to%
\begin{equation}
\left( J_{V}t\right) \left( tJ_{V}\right) ^{t}=\left( J_{V}t\right) \left(
J_{V}^{t}t^{t}\right) =-\left( J_{V}t\right) \left( J_{V}t^{t}\right) =I_{V},
\end{equation}%
where%
\begin{align}
\left( J_{V}t\right) & =\left( 
\begin{array}{cc}
0 & -1 \\ 
1 & 0%
\end{array}%
\right) \left( 
\begin{array}{cc}
2 & 7 \\ 
1 & 4%
\end{array}%
\right) =\left( 
\begin{array}{cc}
-1 & 4 \\ 
2 & 7%
\end{array}%
\right) \\
\left( tJ_{V}\right) ^{t}& =J_{V}^{t}t^{t}=\left( 
\begin{array}{cc}
0 & 1 \\ 
-1 & 0%
\end{array}%
\right) \left( 
\begin{array}{cc}
2 & 1 \\ 
7 & 4%
\end{array}%
\right) =\left( 
\begin{array}{cc}
7 & 4 \\ 
-2 & -1%
\end{array}%
\right) ,
\end{align}%
i.e. $\left( J_{V}t\right) $ and $\left( tJ_{V}\right) ^{t}$ are \emph{%
biorthogonal to each other}. Note that $\left( J_{V}t\right) $ and $\left(
tJ_{V}\right) ^{t}\notin O\left( 2,\mathbb{R}\right) $ i.e. neither $\left(
J_{V}t\right) $ nor $\left( tJ_{V}\right) ^{t}$ belong to the orthogonal
group as 
\begin{align}
\left( 
\begin{array}{cc}
-1 & 4 \\ 
2 & 7%
\end{array}%
\right) \left( 
\begin{array}{cc}
-1 & 2 \\ 
4 & 7%
\end{array}%
\right) & =\left( 
\begin{array}{cc}
16 & 26 \\ 
26 & 53%
\end{array}%
\right) \\
\left( 
\begin{array}{cc}
7 & 4 \\ 
-2 & -1%
\end{array}%
\right) \left( 
\begin{array}{cc}
7 & -2 \\ 
4 & -1%
\end{array}%
\right) & =\left( 
\begin{array}{cc}
65 & -18 \\ 
-18 & 6%
\end{array}%
\right) ,
\end{align}%
but they produce biorthogonal vectors 
\begin{align}
\left( 
\begin{array}{cc}
-1 & -4%
\end{array}%
\right) \left( 
\begin{array}{c}
7 \\ 
-2%
\end{array}%
\right) & =1 \\
\left( 
\begin{array}{cc}
2 & 7%
\end{array}%
\right) \left( 
\begin{array}{c}
4 \\ 
-1%
\end{array}%
\right) & =1 \\
\left( 
\begin{array}{cc}
-1 & -4%
\end{array}%
\right) \left( 
\begin{array}{c}
4 \\ 
-1%
\end{array}%
\right) & =0 \\
\left( 
\begin{array}{cc}
2 & 7%
\end{array}%
\right) \left( 
\begin{array}{c}
7 \\ 
-2%
\end{array}%
\right) & =0,
\end{align}%
thus $\left\{ \left( 
\begin{array}{c}
-1 \\ 
-4%
\end{array}%
\right) ,\left( 
\begin{array}{c}
2 \\ 
7%
\end{array}%
\right) \right\} $ and $\left\{ \left( 
\begin{array}{c}
7 \\ 
-2%
\end{array}%
\right) ,\left( 
\begin{array}{c}
4 \\ 
-1%
\end{array}%
\right) \right\} $ are biorthogonal bases. This corresponds to the rows of $%
J_{V}t$ being biorthogonal to the columns of $J_{V}^{t}t^{t}=\left(
J_{V}t\right) ^{t}.$

Note that $t$ does not have to be an orthogonal matrix i.e. 
\begin{equation}
t=\left( 
\begin{array}{cc}
2 & 7 \\ 
1 & 4%
\end{array}
\right) \notin O\left( 2,\mathbb{R}\right) ,
\end{equation}
as 
\begin{align}
\left( 
\begin{array}{cc}
2 & 7 \\ 
1 & 4%
\end{array}
\right) \left( 
\begin{array}{cc}
2 & 7 \\ 
1 & 4%
\end{array}
\right) ^{t} & =\left( 
\begin{array}{cc}
2 & 7 \\ 
1 & 4%
\end{array}
\right) \left( 
\begin{array}{cc}
2 & 1 \\ 
7 & 4%
\end{array}
\right) =\left( 
\begin{array}{cc}
53 & 30 \\ 
30 & 17%
\end{array}
\right) \\
\left( 
\begin{array}{cc}
2 & 7 \\ 
1 & 4%
\end{array}
\right) ^{t}\left( 
\begin{array}{cc}
2 & 7 \\ 
1 & 4%
\end{array}
\right) & =\left( 
\begin{array}{cc}
2 & 1 \\ 
7 & 4%
\end{array}
\right) \left( 
\begin{array}{cc}
2 & 7 \\ 
1 & 4%
\end{array}
\right) =\left( 
\begin{array}{cc}
5 & 18 \\ 
18 & 65%
\end{array}
\right) .
\end{align}
The invariance condition 
\begin{align}
& t\dsum \limits_{1\leq j,k\leq r}\left\{ \left\vert e_{j}\right) \left(
e_{k}\right\vert \left( \left. e_{j+r}\right\vert t^{t}te_{k+r}\right)
+\left\vert e_{j+r}\right) \left( e_{k+r}\right\vert \left( \left.
e_{j}\right\vert t^{t}te_{k}\right) -\left\vert e_{j}\right) \left(
e_{k+r}\right\vert \left( \left. e_{j+r}\right\vert t^{t}te_{k}\right)
-\left\vert e_{j+r}\right) \left( e_{k}\right\vert \left( \left.
e_{j}\right\vert t^{t}te_{k+r}\right) \right\} t^{t} \\
& =\dsum \limits_{1\leq j\leq r}\left\{ \left\vert e_{j}\right) \left(
e_{j}\right\vert +\left\vert e_{j+r}\right) \left( e_{j+r}\right\vert
\right\} =I_{2r}
\end{align}
can be written as 
\begin{align}
& t\left\vert e_{1}\cdots e_{2r}\right) \left( e_{1}\cdots e_{2r}\right\vert
t^{t}t\left( 
\begin{array}{cc}
\left( \left. e_{j+r}\right\vert t^{t}te_{k+r}\right) & -\left( \left.
e_{j+r}\right\vert t^{t}te_{k}\right) \\ 
-\left( \left. e_{j}\right\vert t^{t}te_{k+r}\right) & \left( \left.
e_{j}\right\vert t^{t}te_{k}\right)%
\end{array}
\right) \left\vert e_{1}\cdots e_{2r}\right) \left( e_{1}\cdots
e_{2r}\right\vert t^{t} \\
& =t\left\vert e_{1}\cdots e_{2r}\right) \left( 
\begin{array}{cc}
\left( \left. e_{j+r}\right\vert t^{t}te_{k+r}\right) & -\left( \left.
e_{j+r}\right\vert t^{t}te_{k}\right) \\ 
-\left( \left. e_{j}\right\vert t^{t}te_{k+r}\right) & \left( \left.
e_{j}\right\vert t^{t}te_{k}\right)%
\end{array}
\right) \left( e_{1}\cdots e_{2r}\right\vert t^{t} \\
& =\left\vert e_{1}\cdots e_{2r}\right) \mathcal{R}_{e}\left( t\right)
\left( 
\begin{array}{cc}
\left( \left. e_{j+r}\right\vert t^{t}te_{k+r}\right) & -\left( \left.
e_{j+r}\right\vert t^{t}te_{k}\right) \\ 
-\left( \left. e_{j}\right\vert t^{t}te_{k+r}\right) & \left( \left.
e_{j}\right\vert t^{t}te_{k}\right)%
\end{array}
\right) \mathcal{R}_{e}\left( t\right) ^{t}\left( e_{1}\cdots
e_{2r}\right\vert ,
\end{align}
which corresponds to 
\begin{align}
& \left( 
\begin{array}{cc}
2 & 7 \\ 
1 & 4%
\end{array}
\right) \left( 
\begin{array}{cc}
0 & 1 \\ 
-1 & 0%
\end{array}
\right) \left( 
\begin{array}{cc}
2 & 1 \\ 
7 & 4%
\end{array}
\right) \left( \left( 
\begin{array}{cc}
2 & 7 \\ 
1 & 4%
\end{array}
\right) \left( 
\begin{array}{cc}
0 & 1 \\ 
-1 & 0%
\end{array}
\right) \left( 
\begin{array}{cc}
2 & 1 \\ 
7 & 4%
\end{array}
\right) \right) ^{t} \\
& =\left( 
\begin{array}{cc}
2 & 7 \\ 
1 & 4%
\end{array}
\right) \left( 
\begin{array}{cc}
0 & 1 \\ 
-1 & 0%
\end{array}
\right) \left( 
\begin{array}{cc}
2 & 1 \\ 
7 & 4%
\end{array}
\right) \left( 
\begin{array}{cc}
2 & 7 \\ 
1 & 4%
\end{array}
\right) \left( 
\begin{array}{cc}
0 & -1 \\ 
1 & 0%
\end{array}
\right) \left( 
\begin{array}{cc}
2 & 1 \\ 
7 & 4%
\end{array}
\right) \\
& =\left( 
\begin{array}{cc}
2 & 7 \\ 
1 & 4%
\end{array}
\right) \left( 
\begin{array}{cc}
0 & 1 \\ 
-1 & 0%
\end{array}
\right) \left( 
\begin{array}{cc}
5 & 18 \\ 
18 & 65%
\end{array}
\right) \left( 
\begin{array}{cc}
0 & -1 \\ 
1 & 0%
\end{array}
\right) \left( 
\begin{array}{cc}
2 & 1 \\ 
7 & 4%
\end{array}
\right) \\
& =\left( 
\begin{array}{cc}
-7 & 2 \\ 
-4 & 1%
\end{array}
\right) \left( 
\begin{array}{cc}
5 & 18 \\ 
18 & 65%
\end{array}
\right) \left( 
\begin{array}{cc}
-7 & -4 \\ 
2 & 1%
\end{array}
\right) =\left( 
\begin{array}{cc}
-35+36 & -126+130 \\ 
-2 & -72+65%
\end{array}
\right) \left( 
\begin{array}{cc}
-7 & -4 \\ 
2 & 1%
\end{array}
\right) \\
& =\left( 
\begin{array}{cc}
1 & 4 \\ 
-2 & -7%
\end{array}
\right) \left( 
\begin{array}{cc}
-7 & -4 \\ 
2 & 1%
\end{array}
\right) =\left( 
\begin{array}{cc}
1 & 0 \\ 
0 & 1%
\end{array}
\right) .
\end{align}

If $t\in O\left( 1,1,\mathbb{R}\right) $ then%
\begin{align}
\left( 
\begin{array}{cc}
a & b \\ 
c & d%
\end{array}%
\right) \left( 
\begin{array}{cc}
1 & 0 \\ 
0 & -1%
\end{array}%
\right) \left( 
\begin{array}{cc}
a^{t} & c^{t} \\ 
b^{t} & d^{t}%
\end{array}%
\right) & =\left( 
\begin{array}{cc}
1 & 0 \\ 
0 & -1%
\end{array}%
\right) \\
\left( 
\begin{array}{cc}
a & -b \\ 
c & -d%
\end{array}%
\right) \left( 
\begin{array}{cc}
a^{t} & c^{t} \\ 
b^{t} & d^{t}%
\end{array}%
\right) & =\left( 
\begin{array}{cc}
aa^{t}-bb^{t} & ac^{t}-bd^{t} \\ 
ca^{t}-bd^{t} & cc^{t}-dd^{t}%
\end{array}%
\right) =\left( 
\begin{array}{cc}
1 & 0 \\ 
0 & -1%
\end{array}%
\right) .
\end{align}%
For $t\in O\left( 1,1,\mathbb{R}\right) $ 
\begin{align}
a^{2}& =1+b^{2}\Longleftrightarrow a^{2}-b^{2}=1\Rightarrow a=\cosh \left(
\theta \right) ;b=\sinh \left( \theta \right) \\
d^{2}& =1+c^{2}\Longleftrightarrow d^{2}-c^{2}=1\Rightarrow d=\cosh \left(
\psi \right) ;c=\sinh \left( \psi \right) \\
ca& =bd\Leftrightarrow \sinh \left( \psi \right) \cosh \left( \theta \right)
-\sinh \left( \theta \right) \cosh \left( \psi \right) =0\Leftrightarrow
\sinh \left( \psi -\theta \right) =0\Leftrightarrow \psi =\theta
\end{align}%
thus%
\begin{equation}
\left( 
\begin{array}{cc}
a & b \\ 
c & d%
\end{array}%
\right) =\left( 
\begin{array}{cc}
\cosh \left( \theta \right) & \sinh \left( \theta \right) \\ 
\sinh \left( \theta \right) & \cosh \left( \theta \right)%
\end{array}%
\right)
\end{equation}%
and one hence can see that 
\begin{equation}
\left( 
\begin{array}{cc}
2 & 7 \\ 
1 & 4%
\end{array}%
\right) \notin O\left( 1,1,\mathbb{R}\right) ,
\end{equation}%
i.e. 
\begin{equation}
\left( 
\begin{array}{cc}
2 & 7 \\ 
1 & 4%
\end{array}%
\right) \left( 
\begin{array}{cc}
1 & 0 \\ 
0 & -1%
\end{array}%
\right) \left( 
\begin{array}{cc}
2 & 1 \\ 
7 & 4%
\end{array}%
\right) =\left( 
\begin{array}{cc}
2 & -7 \\ 
1 & -4%
\end{array}%
\right) \left( 
\begin{array}{cc}
2 & 1 \\ 
7 & 4%
\end{array}%
\right) =\left( 
\begin{array}{cc}
-45 & -26 \\ 
-26 & -15%
\end{array}%
\right) \neq \left( 
\begin{array}{cc}
1 & 0 \\ 
0 & -1%
\end{array}%
\right) .
\end{equation}%
However if 
\begin{equation}
\left( 
\begin{array}{cc}
a & b \\ 
c & d%
\end{array}%
\right) =\left( 
\begin{array}{cc}
\cosh \left( \theta \right) & \sinh \left( \theta \right) \\ 
\sinh \left( \theta \right) & \cosh \left( \theta \right)%
\end{array}%
\right)
\end{equation}%
then 
\begin{equation}
ad-bc=\cosh ^{2}\left( \theta \right) -\sinh ^{2}\left( \theta \right) =1
\end{equation}%
and $t\in Sp\left( 2,\mathbb{R}\right) .$ Hence 
\begin{equation}
Sp\left( 2,\mathbb{R}\right) \supset O\left( 1,1,\mathbb{R}\right) .
\end{equation}

If $t\in O\left( 2,\mathbb{R}\right) $%
\begin{equation}
\left( 
\begin{array}{cc}
a & b \\ 
c & d%
\end{array}%
\right) \left( 
\begin{array}{cc}
a^{t} & c^{t} \\ 
b^{t} & d^{t}%
\end{array}%
\right) =\left( 
\begin{array}{cc}
aa^{t}+bb^{t} & ac^{t}+bd^{t} \\ 
ca^{t}+bd^{t} & cc^{t}+dd^{t}%
\end{array}%
\right) =\left( 
\begin{array}{cc}
1 & 0 \\ 
0 & 1%
\end{array}%
\right)
\end{equation}%
\begin{align}
aa^{t}+bb^{t}& =1 \\
ac^{t}+bd^{t}& =0_{r},
\end{align}%
thus 
\begin{align}
a^{2}+b^{2}& =1\Leftrightarrow a=\cos \theta ;b=\sin \theta \\
c& =\sin \theta ;d=-\cos \theta
\end{align}%
giving%
\begin{equation}
\left( 
\begin{array}{cc}
\cos \theta & \sin \theta \\ 
\sin \theta & -\cos \theta%
\end{array}%
\right) \left( 
\begin{array}{cc}
\cos \theta & \sin \theta \\ 
\sin \theta & -\cos \theta%
\end{array}%
\right) =\left( 
\begin{array}{cc}
1 & 0 \\ 
0 & 1%
\end{array}%
\right) .
\end{equation}

The transformations $S_{W}$,$\,J_{W}$ and $\sigma _{W}$ have the properties 
\begin{align}
S_{W}^{2}& =\left( 
\begin{array}{cc}
0 & -1 \\ 
-1 & 0%
\end{array}%
\right) \left( 
\begin{array}{cc}
0 & -1 \\ 
-1 & 0%
\end{array}%
\right) =\left( 
\begin{array}{cc}
1 & 0 \\ 
0 & 1%
\end{array}%
\right) =I_{2r} \\
S_{W}J_{W}& =\left( 
\begin{array}{cc}
0 & -1 \\ 
-1 & 0%
\end{array}%
\right) \left( 
\begin{array}{cc}
0 & 1 \\ 
-1 & 0%
\end{array}%
\right) =\left( 
\begin{array}{cc}
1 & 0 \\ 
0 & -1%
\end{array}%
\right) =\sigma _{W} \\
J_{W}S_{W}& =\left( 
\begin{array}{cc}
0 & 1 \\ 
-1 & 0%
\end{array}%
\right) \left( 
\begin{array}{cc}
0 & -1 \\ 
-1 & 0%
\end{array}%
\right) =\left( 
\begin{array}{cc}
-1 & 0 \\ 
0 & 1%
\end{array}%
\right) =-\sigma _{W} \\
S_{W}\sigma _{W}& =\left( 
\begin{array}{cc}
0 & -1 \\ 
-1 & 0%
\end{array}%
\right) \left( 
\begin{array}{cc}
1 & 0 \\ 
0 & -1%
\end{array}%
\right) =\left( 
\begin{array}{cc}
0 & 1 \\ 
-1 & 0%
\end{array}%
\right) =J_{W} \\
\sigma _{W}S_{W}& =\left( 
\begin{array}{cc}
1 & 0 \\ 
0 & -1%
\end{array}%
\right) \left( 
\begin{array}{cc}
0 & -1 \\ 
-1 & 0%
\end{array}%
\right) =\left( 
\begin{array}{cc}
0 & -1 \\ 
1 & 0%
\end{array}%
\right) =-J_{W}
\end{align}%
and%
\begin{align}
S_{W}\sigma _{W}S_{W}^{t}& =\left( 
\begin{array}{cc}
0 & -1 \\ 
-1 & 0%
\end{array}%
\right) \left( 
\begin{array}{cc}
1 & 0 \\ 
0 & -1%
\end{array}%
\right) \left( 
\begin{array}{cc}
0 & -1 \\ 
-1 & 0%
\end{array}%
\right) =\left( 
\begin{array}{cc}
0 & 1 \\ 
-1 & 0%
\end{array}%
\right) \left( 
\begin{array}{cc}
0 & -1 \\ 
-1 & 0%
\end{array}%
\right) =\left( 
\begin{array}{cc}
-1 & 0 \\ 
0 & 1%
\end{array}%
\right) =-\sigma _{W} \\
S_{W}J_{W}S_{W}^{t}& =\left( 
\begin{array}{cc}
0 & -1 \\ 
-1 & 0%
\end{array}%
\right) \left( 
\begin{array}{cc}
0 & 1 \\ 
-1 & 0%
\end{array}%
\right) \left( 
\begin{array}{cc}
0 & -1 \\ 
-1 & 0%
\end{array}%
\right) =\left( 
\begin{array}{cc}
1 & 0 \\ 
0 & -1%
\end{array}%
\right) \left( 
\begin{array}{cc}
0 & -1 \\ 
-1 & 0%
\end{array}%
\right) =\left( 
\begin{array}{cc}
0 & -1 \\ 
1 & 0%
\end{array}%
\right) =-J_{W}.
\end{align}

\begin{align}
\left( 
\begin{array}{cc}
2 & 7 \\ 
1 & 4%
\end{array}%
\right) \left( 
\begin{array}{cc}
0 & 1 \\ 
-1 & 0%
\end{array}%
\right) \left( 
\begin{array}{cc}
2 & 7 \\ 
1 & 4%
\end{array}%
\right) ^{t}& =\left( 
\begin{array}{cc}
0 & 1 \\ 
-1 & 0%
\end{array}%
\right) \\
\left( 
\begin{array}{cc}
2 & 7 \\ 
1 & 4%
\end{array}%
\right) \left( 
\begin{array}{cc}
1 & 0 \\ 
0 & -1%
\end{array}%
\right) \left( 
\begin{array}{cc}
2 & 7 \\ 
1 & 4%
\end{array}%
\right) ^{t}& \neq \left( 
\begin{array}{cc}
1 & 0 \\ 
0 & -1%
\end{array}%
\right) .
\end{align}

The dimensions of the real and complexified spaces are 
\begin{align}
\mathbb{R-}\dim \left( V_{\mathbb{C}}\right) & =\mathbb{R-}\dim \left( V_{%
\mathbb{R}}\right) =2r \\
\mathbb{C-}\dim \left( V_{\mathbb{C}}\right) & =r.
\end{align}

The complexified space 
\begin{equation}
\mathcal{V}_{\mathbb{C}}=V\oplus _{\mathbb{R}}J_{V}V=V\oplus V^{\bot
}=\left\{ \left. \left( u,v\right) \right\vert u,v\in V\right\}
\leftrightarrow \left\{ \left. u+J_{V}v\right\vert u\in V,v\in V\right\} ,
\end{equation}%
\begin{equation}
J_{V}\left( v\right) =v^{\bot }\equiv J_{V}\left( v,0\right) =\left(
0,v^{\bot }\right) ;v\in V,
\end{equation}%
where 
\begin{equation}
V=\left\{ v\in V\right\} \cong \left\{ \left. \left( v,0\right) \right\vert
v\in V\right\}
\end{equation}%
and $\bot $ is wrt the real inner product $\left( \left. {}\right\vert
\right) $ and $v^{\bot }$ is the image of $v$ under $J_{V}.$ As the
subspaces $V$ and $J_{V}V$ are disjoint one has 
\begin{equation}
\mathcal{V}_{\mathbb{C}}=V\oplus _{\mathbb{R}}J_{V}V=\mathcal{V}+J_{V}%
\mathcal{V}\cong V+J_{V}V
\end{equation}%
showing that 
\begin{equation}
V_{\mathbb{C}}\cong \mathcal{V}_{\mathbb{C}}.
\end{equation}

\begin{remark}
Recall multiplication by the complex unit is defined by 
\begin{equation}
i\left( a,b^{\bot }\right) \overset{def}{=}\mathcal{J}_{V}\left( a,b^{\bot
}\right) =\left( b^{\bot },a\right) ;\mathcal{J}_{\mathcal{V}}\in \mathcal{L}%
\left( \mathcal{V}\right) .
\end{equation}%
giving 
\begin{equation}
V_{\mathbb{C}}=\left\{ a+b^{\bot }\right\} =\left\{ a+J_{V}\left( b\right)
\right\} =\left\{ a+ib\right\} ;a,b\in V,J_{V}\in \mathcal{L}\left( V\right)
.
\end{equation}
\end{remark}

The conjugation $\sigma _{\mathcal{V}}\in \mathcal{L}\left( \mathcal{V}%
\right) $ has the effect 
\begin{align}
\sigma _{V}\left( \left( v_{1},v_{2}^{\bot }\right) \right) & =\left(
v_{1},-v_{2}^{\bot }\right) \\
\sigma _{V}\left( v_{1}+J_{V}\left( v_{2}\right) \right) &
=v_{1}-J_{V}\left( v_{2}\right) \equiv \sigma _{V}\left( v_{1}+i\left(
v_{2}\right) \right) =v_{1}-i\left( v_{2}\right) ;J_{V}\in \mathcal{L}\left(
V\right) .
\end{align}

\begin{remark}
We can define the real and imaginary part real spaces of $V_{\mathbb{C}}$ as 
\begin{align}
\Sigma _{Vr}\left\{ V_{\mathbb{C}}\right\} & =\left\{ \frac{1}{2}\left(
v+\sigma _{V}\left( v\right) \right) \left\vert v=\left( v_{1},v_{2}^{\bot
}\right) \in V_{\mathbb{C}}\right. \right\} \\
& =\left\{ \frac{1}{2}\left( \left( v_{1},v_{2}^{\bot }\right) +\left(
v_{1},-v_{2}^{\bot }\right) \right) =\left( v_{1},0\right) \left\vert \left(
v_{1},v_{2}^{\bot }\right) \in V_{\mathbb{C}}\right. \right\} \\
& =\left\{ \left. \left( v_{1},0\right) \right\vert v_{1}\in V_{1}\right\}
\cong V_{1}=\left\{ \frac{1}{2}\left( v_{1}+J_{V}\left( v_{2}\right)
+v_{1}-J_{V}\left( v_{2}\right) \right) \right\} =\left\{ v_{1}\right\} \\
\Sigma _{Vi}\left\{ V_{\mathbb{C}}\right\} & =\left\{ v=\left(
v_{1},v_{2}^{\bot }\right) \in V_{\mathbb{C}}\left\vert \frac{1}{2}\left(
v-\sigma _{V}\left( v\right) \right) \right. \right\} \\
& =\left\{ \frac{1}{2}\left( \left( v_{1},v_{2}^{\bot }\right) +\left(
-v_{1},v_{2}^{\bot }\right) \right) =\left( 0,v_{2}^{\bot }\right)
\left\vert \left( v_{1},v_{2}^{\bot }\right) \in V_{\mathbb{C}}\right.
\right\} \\
& =\left\{ \left. \left( 0,v_{2}^{\bot }\right) \right\vert v_{2}^{\bot }\in
V_{1}^{\bot }\right\} \cong V_{1}^{\bot }=\left\{ \frac{1}{2}\left(
v_{1}+J_{V}\left( v_{2}\right) -v_{1}+J_{V}\left( v_{2}\right) \right)
\right\} =\left\{ J_{V}\left( v_{2}\right) \in V_{1}^{\bot },v_{2}\in
V\right\}
\end{align}%
to obtain 
\begin{equation}
v=\Sigma _{Vr}\left( v\right) +\Sigma _{Vi}\left( v\right) \in V_{\mathbb{C}%
},
\end{equation}%
so that%
\begin{align}
\mathcal{V}_{\mathbb{C}}& =\mathcal{V}+\mathcal{J}_{V}\mathcal{V}=\Sigma
_{Vr}\left( V_{\mathbb{C}}\right) \oplus \Sigma _{Vi}\left( V_{\mathbb{C}%
}\right) \equiv \left( v_{1},0\right) +\mathcal{J}_{\mathcal{V}}\left(
v_{2},0\right) =\left( v_{1},v_{2}^{\bot }\right) \\
& \cong \Sigma _{Vr}\left( V_{\mathbb{C}}\right) +\Sigma _{Vi}\left( V_{%
\mathbb{C}}\right) .
\end{align}
\end{remark}

\begin{remark}
\begin{gather}
-\mathcal{J}_{\mathcal{V}}\Sigma _{i}\left( V_{\mathbb{C}}\right) =\left\{ -%
\mathcal{J}_{\mathcal{V}}\left( 0,v_{2}^{\bot }\right) =\left(
v_{2},0\right) \right\} \\
\updownarrow \\
-i\Sigma _{i}\left( V_{\mathbb{C}}\right) =-iv_{2}^{\bot }=-iiv_{2}=v_{2}
\end{gather}
\end{remark}

\begin{remark}
\begin{equation}
V_{\mathbb{C}}=V_{1}\oplus _{\mathbb{R}}iV_{1}=V_{1}\oplus _{\mathbb{R}}%
\mathcal{J}_{\mathcal{V}}V_{1}
\end{equation}%
\begin{align}
\Sigma _{r}\left\{ V_{\mathbb{C}}\right\} & \leftrightarrow \left\{ \frac{1}{%
2}\left( v_{1}+iv_{2}+v_{1}-iv_{2}\right) =v_{1}\right\} \cong V_{1} \\
\Sigma _{i}\left\{ V_{\mathbb{C}}\right\} & \leftrightarrow \left\{ \frac{1}{%
2}\left( v_{1}+iv_{2}-v_{1}+iv_{2}\right) =iv_{2}\right\} \cong V_{1}^{\bot
}.
\end{align}%
$\rfloor $
\end{remark}

\subsection{Clifford Algebra Example}

\subsubsection{Complexification}

The generating Hilbert space $\mathcal{F}_{1}$ can be decomposed as the
complex sum of two disjoint \emph{non self adjoint} spaces (spaces that are
not closed under the adjoint operation) that is isomorphic to their complex
direct sum 
\begin{equation}
\mathcal{F}_{1}=\mathbf{b}^{\dagger }+_{\mathbb{C}}\mathbf{b\cong b}%
^{\dagger }\oplus _{\mathbb{C}}\mathbf{b}
\end{equation}%
or into the complex sum of two disjoint \emph{self adjoint} spaces (spaces
that are closed under the adjoint operation) of dimension $r$ i.e.%
\begin{equation}
\mathcal{F}_{1}=\mathbf{x}_{1}+_{\mathbb{C}}\mathbf{x}_{2}\,\mathbf{\cong x}%
_{1}\oplus _{\mathbb{C}}\mathbf{x}_{2}.
\end{equation}%
The spaces $\mathbf{b}^{\dagger }$and $\mathbf{b}$ are each complex spaces,
while $\mathbf{x}_{1}$ and $\mathbf{x}_{2}$ are real spaces that have
complexifications $\mathbf{x}_{1\mathbb{C}}=\mathbf{x}_{1}\otimes \mathbb{C}$
and $\mathbf{x}_{2\mathbb{C}}=\mathbf{x}_{2}\otimes \mathbb{C}.$

One has the following identities 
\begin{align}
\mathcal{F}_{1\mathbb{C}}& =\mathbf{x}_{1}\oplus _{\mathbb{C}}\mathbf{x}%
_{2}\equiv \mathbf{x}_{1\mathbb{C}}\oplus \mathbf{x}_{2\mathbb{C}} \\
\mathcal{F}_{1\mathbb{R}}& =\mathbf{x}_{1\mathbb{R}}\oplus _{\mathbb{R}}%
\mathbf{x}_{2\mathbb{R}} \\
\mathbf{x}_{1\mathbb{C}}& =\mathbb{C}\text{-}\limfunc{linspan}\left\{ \left.
x_{j}\right\vert 1\leq j\leq r\right\} \\
\mathbf{x}_{2\mathbb{C}}& =\mathbb{C}\text{-}\limfunc{linspan}\left\{ \left.
x_{j+r}\right\vert 1\leq j\leq r\right\} \\
\mathbf{x}_{1}& \equiv \mathbf{x}_{1\mathbb{R}}=\mathbb{R}\text{-}\limfunc{%
linspan}\left\{ \left. x_{j}\right\vert 1\leq j\leq r\right\} \\
\mathbf{x}_{2}& \equiv \mathbf{x}_{2\mathbb{R}}=\mathbb{R}\text{-}\limfunc{%
linspan}\left\{ \left. x_{j+r}\right\vert 1\leq j\leq r\right\} ,
\end{align}%
where the \emph{unnormalized} real basis is related to the complex one by%
\begin{equation}
\left( 
\begin{array}{cc}
x_{j} & x_{j+r}%
\end{array}%
\right) =K\left( 
\begin{array}{cc}
a_{j}^{\dagger } & a_{j}%
\end{array}%
\right) =\left( 
\begin{array}{cc}
a_{j}^{\dagger } & a_{j}%
\end{array}%
\right) \mathbb{K}_{j}=\left( 
\begin{array}{cc}
a_{j}^{\dagger } & a_{j}%
\end{array}%
\right) \left( 
\begin{array}{cc}
\frac{1}{\sqrt{2}} & \frac{i}{\sqrt{2}} \\ 
\frac{1}{\sqrt{2}} & -\frac{i}{\sqrt{2}}%
\end{array}%
\right) ,
\end{equation}%
giving%
\begin{align}
& \mathbf{x}_{1}=\Omega _{h}\left\{ \mathbf{a}^{\dagger }\right\} =\limfunc{%
linspan}\left\{ \left\vert \frac{1}{\sqrt{2}}\left( a_{j}^{\dagger
}+a_{j}\right) \right. 1\leq j\leq r\right\} \\
& \mathbf{x}_{2}=i\Omega _{ah}\left\{ \mathbf{a}^{\dagger }\right\} =%
\limfunc{linspan}\left\{ \left\vert \frac{i}{\sqrt{2}}\left( a_{j}^{\dagger
}-a_{j}\right) \right. 1\leq j\leq r\right\} \\
& \Omega _{ah}\left\{ \mathbf{a}^{\dagger }\right\} =\limfunc{linspan}%
\left\{ \left\vert \frac{1}{\sqrt{2}}\left( a_{j}^{\dagger }-a_{j}\right)
\right. 1\leq j\leq r\right\}
\end{align}%
and%
\begin{equation}
\mathcal{F}_{1\mathbb{C}}=\Omega _{h}\left( \mathbf{a}^{\dagger }\right)
\oplus i\Omega _{ah}\left( \mathbf{a}^{\dagger }\right) \cong \Omega
_{h}\left( \mathbf{a}^{\dagger }\right) +i\Omega _{ah}\left( \mathbf{a}%
^{\dagger }\right) =\mathbf{x}_{1}+\mathbf{x}_{2}=\mathcal{F}_{1}.
\end{equation}

\begin{remark}
\begin{align*}
x_{j}^{2}& =\left( \frac{1}{\sqrt{2}}\left( a_{j}^{\dagger }+a_{j}\right)
\right) ^{2}=\frac{1}{2}\left( a_{j}^{\dagger }a_{j}+a_{j}a_{j}^{\dagger
}\right) =\frac{1}{2}I_{\mathcal{F}} \\
x_{j+r}^{2}& =\left( \frac{i}{\sqrt{2}}\left( a_{j}^{\dagger }-a_{j}\right)
\right) ^{2}=\frac{1}{2}\left( a_{j}^{\dagger }a_{j}+a_{j}a_{j}^{\dagger
}\right) =\frac{1}{2}I_{\mathcal{F}}
\end{align*}
\end{remark}

One can introduce the complexifications $\widehat{J_{x}}$ 
\begin{align}
\widehat{J_{x}}\left( x_{k}\right) & =\left[ J_{x},x_{k}\right] =\left[
\dsum \limits_{1\leq j\leq r}x_{j}x_{j+r},x_{k}\right] \\
& =\dsum \limits_{1\leq j\leq r}\left[ x_{j}x_{j+r},x_{k}\right] =\left[
x_{k}x_{k+r},x_{k}\right] =x_{k}x_{k+r}x_{k}-x_{k}x_{k}x_{k+r}=-x_{k+r} \\
\widehat{J_{x}}\left( x_{k+r}\right) & =\left[ J_{x},x_{k+r}\right] \\
& =\left[ \dsum \limits_{1\leq j\leq r}x_{j}x_{j+r},x_{k+r}\right] =\dsum
\limits_{1\leq j\leq r}\left[ x_{j}x_{j+r},x_{k+r}\right] =\left[
x_{k}x_{k+r},x_{k+r}\right] =x_{k}x_{k+r}x_{k+r}-x_{k+r}x_{k}x_{k+r}=x_{k}
\end{align}
of the real spaces $\left\{ \mathbf{x}_{1},\mathbf{x}_{2}\right\} $ that
produce $\mathcal{F}_{1}$, that are in general not unitarily equivalent when 
$r=\infty$.

Using unnormalized $\left\{ \left. x_{j}\right\vert 1\leq j\leq 2r\right\} $
one obtains 
\begin{align}
\mathcal{F}_{1\mathbb{R}}& =\mathbf{x}_{1\mathbb{R}}+\mathbf{x}_{2\mathbb{R}}
\\
\mathbf{x}_{1\mathbb{C}}& =\mathbf{x}_{1}+\widehat{J_{x}}\mathbf{x}%
_{1}\equiv \mathbf{x}_{1}+i\mathbf{x}_{1} \\
\mathbf{x}_{2\mathbb{C}}& =\mathbf{x}_{2}+\widehat{J_{x}}\mathbf{x}%
_{2}\equiv \mathbf{x}_{2}+i\mathbf{x}_{2} \\
\mathcal{F}_{1x}& =\mathbf{x}_{1\mathbb{C}}+\mathbf{x}_{2\mathbb{C}}.
\end{align}%
In finite dimensions all bases $\mathbf{x}$ give the same space $\mathcal{F}%
_{1}$ (i.e. $\mathcal{F}_{1}=\mathcal{F}_{1x}$ $\forall x)$ up to unitary
equivalence $\left( \text{seems reasonable but maybe not true}\right) ,$ but
in infinite dimensions spaces unitarily inequivalence can occur.

Multiplication by the complex unit is defined by 
\begin{align}
& i\left( a,b\right) \overset{def}{=}\widehat{\mathcal{J}_{x}}\left(
a,b\right) ;a\in\mathbf{x}_{1};b\in\mathbf{x}_{2} \\
& \widehat{\mathcal{J}_{x}}\left( a,0\right) =\left( 0,\widehat{J_{x}}\left(
a\right) \right) ;a\in\mathbf{x}_{1} \\
& \widehat{\mathcal{J}_{x}}\left( 0,b\right) =\left( \widehat{J_{x}}\left(
b\right) ,0\right) ;b\in\mathbf{x}_{2},
\end{align}%
\begin{equation}
\left( a,b\right) \in\mathbf{x}_{1}\oplus_{\mathbb{R}}\mathbf{x}_{2}\cong%
\mathbf{x}_{1}+_{\mathbb{R}}\mathbf{x}_{2},
\end{equation}
and 
\begin{equation}
\mathbf{x}_{1}\bot\mathbf{x}_{2}
\end{equation}
i.e. 
\begin{gather}
a\in\mathbf{x}_{1}\subset\mathcal{F}_{1} \\
b\in\mathbf{x}_{2}\subset\mathcal{F}_{1} \\
\left( a,0\right) +\left( 0,b\right) \leftrightarrow a+b.
\end{gather}

\subsubsection{Groups}

The action of the group $O\left( 2r,\mathbb{R}\right) $ on $\mathcal{F}$ can
be realized as inner automorphisms 
\begin{align}
\exp\left\{ \sum_{1\leq j<k\leq2r}g_{jk}x_{j}x_{k}\right\} A\exp\left\{
-\sum_{1\leq j<k\leq2r}g_{jk}x_{j}x_{k}\right\} & =\exp\left\{ \sum_{1\leq
j<k\leq2r}g_{jk}x_{j}x_{k}\right\} A\exp\left\{ \sum_{1\leq
j<k\leq2r}g_{jk}x_{j}x_{k}\right\} ^{t} \\
g_{jk} & \in\mathbb{R},1\leq j<k\leq2r
\end{align}
based on the Lie algebra $\mathcal{F}_{2}.$ The general linear group $%
GL\left( 2^{2r},\mathbb{C}\right) $ and all of its subgroups can be
represented as inner automorphisms generated by the Lie algebra $\mathcal{F}$
i.e. 
\begin{align}
& \exp\left\{ \sum_{0\leq n\leq2r}\sum_{1\leq j_{1}<\cdots<j_{n}\leq
2r}g_{j_{1}\cdots j_{n}}x_{j_{1}}\cdots x_{j_{n}}\right\} A\exp\left\{
-\sum_{0\leq n\leq2r}\sum_{1\leq j_{1}<\cdots<j_{n}\leq2r}g_{j_{1}\cdots
j_{n}}x_{j_{1}}\cdots x_{j_{n}}\right\} , \\
& \left\{ \left. g_{j_{1}\cdots j_{n}}\in\mathbb{C}\right\vert 1\leq
j_{1}<\cdots<j_{n}\leq2r,0\leq n\leq2r\right\} ,
\end{align}
when $r<\infty$, if $r=\infty$ various topological details must be taken
into account and handled. In particular $GL\left( 2r,\mathbb{C}\right) $ and
thus $Sp\left( 2r,\mathbb{C}\right) $ are subgroups of $GL\left( 2^{2r},%
\mathbb{C}\right) $ that have invariant subspaces $\tbigwedge ^{n}\mathcal{F}%
_{1}$ as they map $\mathcal{F}_{1}$ to $\mathcal{F}_{1},$ though not
necessarily generated by transformations produced by the Lie algebra $%
\mathcal{F}_{2},$ as inner automorphisms, unless they also belong to the
subgroup $O\left( 2r,\mathbb{R}\right) $.

\begin{remark}
All algebra automorphisms in finite dimensions are inner, this does mean
that all vector space automorphisms on an algebra are inner. In particular 
\begin{equation}
GL\left( 2^{2r},\mathbb{C}\right) :\mathcal{F}\rightarrow \mathcal{F}
\end{equation}%
looks like it acts as vector space automorphisms, however $\mathcal{F}$ is
the Lie algebra $gl\left( 2^{2r},\mathbb{C}\right) ,$ which always carries
the adjoint representation$,$which is the same as the inner automorphism
representations of $GL\left( 2^{2r},\mathbb{C}\right) $ and $gl\left( 2^{2r},%
\mathbb{C}\right) $ on $\mathcal{F}.$ This does mean that $\tbigwedge^{n}%
\mathcal{F}_{1}$ are obviously invariant subspaces, but this is something
that must be proven.
\end{remark}

Symplectic forms on $\mathcal{F}_{1}$ are defined by 
\begin{align}
\left\{ \left. x_{k}\right\vert x_{l}\right\} & =\left( x_{k}\left\vert 
\widehat{J_{x}}\left( \overset{}{x_{l}}\right) \right. \right) =\left\langle
\left\langle x_{k}\widehat{\sum_{1\leq j\leq r}x_{j}x_{j+r}}\left( \overset{}%
{x_{l}}\right) \right\rangle \right\rangle _{0} \\
& =-\left\{ \left. x_{l}\right\vert x_{k}\right\} =-\left( x_{l}\left\vert 
\widehat{J_{x}}\left( \overset{}{x_{k}}\right) \right. \right)
=-\left\langle \left\langle x_{l}\widehat{\sum_{1\leq j\leq r}x_{j}x_{j+r}}%
\left( \overset{}{x_{k}}\right) \right\rangle \right\rangle _{0}=\left\{ 
\begin{array}{c}
1\text{ if }k=l+r \\ 
-1\text{ if }k-r=l \\ 
0\text{ otherwise}%
\end{array}%
\right. .
\end{align}

\begin{remark}
\begin{equation}
\left\langle \left\langle x_{k+r}\widehat{\sum_{1\leq j\leq r}x_{j}x_{j+r}}%
\left( x_{l}\right) \right\rangle \right\rangle _{0}=\left\langle
\left\langle x_{kk+r}\widehat{\sum_{1\leq j\leq r}x_{j}x_{j+r}}\left(
x_{l+r}\right) \right\rangle \right\rangle _{0}=0;1\leq k\neq l\leq r
\end{equation}%
\begin{equation}
\left\langle \left\langle x_{k}\widehat{\sum_{1\leq j\leq r}x_{j}x_{j+r}}%
\left( x_{l}\right) \right\rangle \right\rangle _{0}=\left\langle
\left\langle x_{k+rk+r}\widehat{\sum_{1\leq j\leq r}x_{j}x_{j+r}}\left(
x_{l+r}\right) \right\rangle \right\rangle _{0}=0;1\leq k\neq l\leq r
\end{equation}
\end{remark}

\begin{example}
Analysing the $2\times 2$ example where 
\begin{equation}
\mathcal{R}_{e}\left( t\right) =\left( 
\begin{array}{cc}
2 & 7 \\ 
1 & 4%
\end{array}%
\right) \in Sp\left( 2,\mathbb{R}\right) ,
\end{equation}%
gives%
\begin{align}
& t\left( 
\begin{array}{cc}
e_{1} & e_{2}%
\end{array}%
\right) =\left( 
\begin{array}{cc}
e_{1} & e_{2}%
\end{array}%
\right) \left( 
\begin{array}{cc}
2 & 7 \\ 
1 & 4%
\end{array}%
\right) \\
& =\left( 
\begin{array}{cc}
2e_{1}+e_{2} & 7e_{1}+4e_{2}%
\end{array}%
\right) =\left( 
\begin{array}{cc}
f_{1} & f_{2}%
\end{array}%
\right)
\end{align}%
Normalization of $\mathbf{f}$ gives 
\begin{align}
& \left\vert \mathbf{f}\right) \Delta _{f}^{-1}\left( \mathbf{f}\right\vert
\left\vert \mathbf{f}\right) \Delta _{f}^{-1}\left( \mathbf{f}\right\vert
=\left\vert \mathbf{f}\right) \Delta _{f}^{-1}\left( \mathbf{f}\right\vert
;\quad \Delta _{f}=\left( \left. \mathbf{f}\right\vert \mathbf{f}\right) \\
& v=\left\vert \mathbf{f}\right) \Delta _{f}^{-1}\left( \mathbf{f}%
\right\vert v=\left\vert \mathbf{f}\right) \Delta _{f}^{-1}\left( \mathbf{f}%
\right\vert \left\vert \mathbf{f}\right) \Delta _{f}^{-\frac{1}{2}}\mathbf{v}%
=\left\vert \mathbf{f}\right) \Delta _{f}^{-\frac{1}{2}}\mathbf{v}%
=v\Leftrightarrow \left\vert \mathbf{f}\right) \Delta _{f}^{-1}\left( 
\mathbf{f}\right\vert =I \\
& \left\vert \mathbf{f}_{n}\right) =\left\vert \mathbf{f}\right) \Delta
_{f}^{-\frac{1}{2}} \\
& \left( \mathbf{f}_{n}\right\vert \left\vert \mathbf{f}_{n}\right) =\Delta
_{f}^{-\frac{1}{2}}\Delta _{f}\Delta _{f}^{-\frac{1}{2}}=\mathbf{I.}
\end{align}%
Complexifications are given by 
\begin{align}
& J_{e}=\left\vert e_{1}\right) \left( e_{2}\right\vert -\left\vert
e_{2}\right) \left( e_{1}\right\vert \\
& tJ_{e}t^{t}=\left\vert te_{1}\right) \left( te_{2}\right\vert -\left\vert
te_{2}\right) \left( te_{1}\right\vert =\left\vert f_{1}\right) \left(
f_{2}\right\vert -\left\vert f_{2}\right) \left( f_{1}\right\vert =J_{f} \\
& \left( tJ_{e}t^{t}\right) ^{2}=\left( \left\vert f_{1}\right) \left(
f_{2}\right\vert -\left\vert f_{2}\right) \left( f_{1}\right\vert \right)
\left( \left\vert f_{1}\right) \left( f_{2}\right\vert -\left\vert
f_{2}\right) \left( f_{1}\right\vert \right) \\
& =\left\vert f_{1}\right) \left( f_{2}\right\vert \left\vert f_{1}\right)
\left( f_{2}\right\vert -\left\vert f_{1}\right) \left( f_{2}\right\vert
\left\vert f_{2}\right) \left( f_{1}\right\vert -\left\vert f_{2}\right)
\left( f_{1}\right\vert \left\vert f_{1}\right) \left( f_{2}\right\vert
+\left\vert f_{2}\right) \left( f_{1}\right\vert \left\vert f_{2}\right)
\left( f_{1}\right\vert .
\end{align}%
Representations of $J_{f}$ are given by 
\begin{align}
& tJ_{e}t^{t}\left( 
\begin{array}{cc}
f_{1} & f_{2}%
\end{array}%
\right) =\left( 
\begin{array}{cc}
f_{1} & f_{2}%
\end{array}%
\right) \mathcal{R}_{f}\left( tJ_{e}t^{t}\right) \\
& \left( 
\begin{array}{c}
f_{1} \\ 
f_{2}%
\end{array}%
\right) tJ_{e}t^{t}\left( 
\begin{array}{cc}
f_{1} & f_{2}%
\end{array}%
\right) =\mathcal{M}_{f}\left( tJ_{e}t^{t}\right) \\
& =\left( 
\begin{array}{c}
\mathbf{f}_{1} \\ 
\mathbf{f}_{2}%
\end{array}%
\right) \left( 
\begin{array}{cc}
\mathbf{f}_{1} & \mathbf{f}_{2}%
\end{array}%
\right) \mathcal{R}_{f}\left( tJ_{e}t^{t}\right) =\left( \left. \mathbf{f}%
\right\vert \mathbf{f}\right) \mathcal{R}_{f}\left( tJ_{e}t^{t}\right)
=\Delta _{f}\mathcal{R}_{f}\left( tJ_{e}t^{t}\right) \\
& \mathcal{R}_{f}\left( tJ_{e}t^{t}\right) =\Delta _{f}^{-1}\mathcal{M}%
_{f}\left( tJ_{e}t^{t}\right) \\
& \Delta _{f}\mathcal{R}_{f}\left( tJ_{e}t^{t}\right) =\mathcal{M}_{f}\left(
tJ_{e}t^{t}\right) ,
\end{align}%
giving%
\begin{align}
tJ_{e}t^{t}\left( 
\begin{array}{cc}
e_{1} & e_{2}%
\end{array}%
\right) & =\left( 
\begin{array}{cc}
e_{1} & e_{2}%
\end{array}%
\right) \mathcal{R}_{e}\left( tJ_{e}t^{t}\right) =\left( 
\begin{array}{cc}
e_{1} & e_{2}%
\end{array}%
\right) \left( 
\begin{array}{cc}
2 & 7 \\ 
1 & 4%
\end{array}%
\right) \left( 
\begin{array}{cc}
0 & 1 \\ 
-1 & 0%
\end{array}%
\right) \left( 
\begin{array}{cc}
2 & 7 \\ 
1 & 4%
\end{array}%
\right) ^{t} \\
& =\left( 
\begin{array}{cc}
e_{1} & e_{2}%
\end{array}%
\right) \left( 
\begin{array}{cc}
0 & 1 \\ 
-1 & 0%
\end{array}%
\right) .
\end{align}%
Then%
\begin{align}
& \mathcal{R}_{e}\left( tJ_{e}t^{t}\right) =\left( 
\begin{array}{cc}
2 & 7 \\ 
1 & 4%
\end{array}%
\right) \left( 
\begin{array}{cc}
0 & 1 \\ 
-1 & 0%
\end{array}%
\right) \left( 
\begin{array}{cc}
2 & 1 \\ 
7 & 4%
\end{array}%
\right) =\left( 
\begin{array}{cc}
-7 & 2 \\ 
-4 & 1%
\end{array}%
\right) \left( 
\begin{array}{cc}
2 & 1 \\ 
7 & 4%
\end{array}%
\right) =\left( 
\begin{array}{cc}
0 & 1 \\ 
-1 & 0%
\end{array}%
\right) \\
& tJ_{e}t^{t}=\left\vert \left( 
\begin{array}{cc}
2 & 7 \\ 
1 & 4%
\end{array}%
\right) \left( 
\begin{array}{c}
1 \\ 
0%
\end{array}%
\right) \right) \left( \left( 
\begin{array}{cc}
2 & 7 \\ 
1 & 4%
\end{array}%
\right) \left( 
\begin{array}{c}
0 \\ 
1%
\end{array}%
\right) \right\vert -\left\vert \left( 
\begin{array}{cc}
2 & 7 \\ 
1 & 4%
\end{array}%
\right) \left( 
\begin{array}{c}
0 \\ 
1%
\end{array}%
\right) \right) \left( \left( 
\begin{array}{cc}
2 & 7 \\ 
1 & 4%
\end{array}%
\right) \left( 
\begin{array}{c}
1 \\ 
0%
\end{array}%
\right) \right\vert \\
& =\left( 
\begin{array}{cc}
2 & 7 \\ 
1 & 4%
\end{array}%
\right) \left\vert \left( 
\begin{array}{c}
1 \\ 
0%
\end{array}%
\right) \right) \left( \left( 
\begin{array}{c}
0 \\ 
1%
\end{array}%
\right) \right\vert \left( 
\begin{array}{cc}
2 & 7 \\ 
1 & 4%
\end{array}%
\right) ^{t}-\left( 
\begin{array}{cc}
2 & 7 \\ 
1 & 4%
\end{array}%
\right) \left\vert \left( 
\begin{array}{c}
0 \\ 
1%
\end{array}%
\right) \right) \left( \left( 
\begin{array}{c}
1 \\ 
0%
\end{array}%
\right) \right\vert \left( 
\begin{array}{cc}
2 & 7 \\ 
1 & 4%
\end{array}%
\right) ^{t} \\
& =\left\vert \left( 
\begin{array}{c}
2 \\ 
1%
\end{array}%
\right) \right) \left( \left( 
\begin{array}{c}
7 \\ 
4%
\end{array}%
\right) \right\vert -\left\vert \left( 
\begin{array}{c}
7 \\ 
4%
\end{array}%
\right) \right) \left( \left( 
\begin{array}{c}
2 \\ 
1%
\end{array}%
\right) \right\vert \\
& \Leftrightarrow \mathcal{R}_{e}\left( tJ_{e}t^{t}\right) =\mathcal{M}%
_{e}\left( tJ_{e}t^{t}\right) =\left( 
\begin{array}{cc}
14 & 8 \\ 
7 & 4%
\end{array}%
\right) -\left( 
\begin{array}{cc}
0 & -7 \\ 
-8 & 0%
\end{array}%
\right) =\left( 
\begin{array}{cc}
0 & 1 \\ 
-1 & 0%
\end{array}%
\right) .
\end{align}%
Ket-Bra form of $tJ_{e}t^{t}$ 
\begin{equation}
tJ_{e}t^{t}=\left\vert \left( 
\begin{array}{c}
2 \\ 
1%
\end{array}%
\right) \right) \left( \left( 
\begin{array}{c}
7 \\ 
4%
\end{array}%
\right) \right\vert -\left\vert \left( 
\begin{array}{c}
7 \\ 
4%
\end{array}%
\right) \right) \left( \left( 
\begin{array}{c}
2 \\ 
1%
\end{array}%
\right) \right\vert
\end{equation}%
$\lceil $Thus 
\begin{align}
\left( 
\begin{array}{cc}
2 & 7 \\ 
1 & 4%
\end{array}%
\right) \left( 
\begin{array}{c}
1 \\ 
0%
\end{array}%
\right) & =\left( 
\begin{array}{c}
2 \\ 
1%
\end{array}%
\right) \\
\left( 
\begin{array}{cc}
2 & 7 \\ 
1 & 4%
\end{array}%
\right) \left( 
\begin{array}{c}
0 \\ 
1%
\end{array}%
\right) & =\left( 
\begin{array}{c}
7 \\ 
4%
\end{array}%
\right)
\end{align}%
$\rfloor $
\end{example}

has the following action on basis elements

\begin{enumerate}
\item 
\begin{align}
& \left\{ \left\vert \left( 
\begin{array}{c}
2 \\ 
1%
\end{array}%
\right) \right) \left( \left( 
\begin{array}{c}
7 \\ 
4%
\end{array}%
\right) \right\vert -\left\vert \left( 
\begin{array}{c}
7 \\ 
4%
\end{array}%
\right) \right) \left( \left( 
\begin{array}{c}
2 \\ 
1%
\end{array}%
\right) \right\vert \right\} \left( 
\begin{array}{c}
1 \\ 
0%
\end{array}%
\right) \\
& =\left\vert \left( 
\begin{array}{c}
2 \\ 
1%
\end{array}%
\right) \right) 7-\left\vert \left( 
\begin{array}{c}
7 \\ 
4%
\end{array}%
\right) \right) 2=\left\vert \left( 
\begin{array}{c}
14-14 \\ 
7-8%
\end{array}%
\right) \right) =\left\vert \left( 
\begin{array}{c}
0 \\ 
-1%
\end{array}%
\right) \right)
\end{align}

\item 
\begin{align}
& \left\{ \left\vert \left( 
\begin{array}{c}
2 \\ 
1%
\end{array}%
\right) \right) \left( \left( 
\begin{array}{c}
7 \\ 
4%
\end{array}%
\right) \right\vert -\left\vert \left( 
\begin{array}{c}
7 \\ 
4%
\end{array}%
\right) \right) \left( \left( 
\begin{array}{c}
2 \\ 
1%
\end{array}%
\right) \right\vert \right\} \left( 
\begin{array}{c}
0 \\ 
1%
\end{array}%
\right) \\
& =\left\vert \left( 
\begin{array}{c}
2 \\ 
1%
\end{array}%
\right) \right) 4-\left\vert \left( 
\begin{array}{c}
7 \\ 
4%
\end{array}%
\right) \right) 1=\left\vert \left( 
\begin{array}{c}
8-7 \\ 
4-4%
\end{array}%
\right) \right) =\left\vert \left( 
\begin{array}{c}
1 \\ 
0%
\end{array}%
\right) \right) .
\end{align}
\end{enumerate}

Ket-Bra form of $\left( tJ_{e}t^{t}\right) ^{2}$ 
\begin{equation}
\left( tJ_{e}t^{t}\right) ^{2}=\left( \left\vert \left( 
\begin{array}{c}
2 \\ 
1%
\end{array}%
\right) \right) \left( \left( 
\begin{array}{c}
7 \\ 
4%
\end{array}%
\right) \right\vert -\left\vert \left( 
\begin{array}{c}
7 \\ 
4%
\end{array}%
\right) \right) \left( \left( 
\begin{array}{c}
2 \\ 
1%
\end{array}%
\right) \right\vert \right) ^{2}
\end{equation}%
has the following action on basis elements

\begin{enumerate}
\item 
\begin{align}
& \left\{ \left\vert \left( 
\begin{array}{c}
2 \\ 
1%
\end{array}%
\right) \right) \left( \left( 
\begin{array}{c}
7 \\ 
4%
\end{array}%
\right) \right\vert -\left\vert \left( 
\begin{array}{c}
7 \\ 
4%
\end{array}%
\right) \right) \left( \left( 
\begin{array}{c}
2 \\ 
1%
\end{array}%
\right) \right\vert \right\} ^{2}\left( 
\begin{array}{c}
1 \\ 
0%
\end{array}%
\right) \\
& =\left\{ \left\vert \left( 
\begin{array}{c}
2 \\ 
1%
\end{array}%
\right) \right) \left( \left( 
\begin{array}{c}
7 \\ 
4%
\end{array}%
\right) \right\vert -\left\vert \left( 
\begin{array}{c}
7 \\ 
4%
\end{array}%
\right) \right) \left( \left( 
\begin{array}{c}
2 \\ 
1%
\end{array}%
\right) \right\vert \right\} \left( 
\begin{array}{c}
0 \\ 
-1%
\end{array}%
\right) \\
& =\left\vert \left( 
\begin{array}{c}
2 \\ 
1%
\end{array}%
\right) \right) \left( -4\right) -\left\vert \left( 
\begin{array}{c}
7 \\ 
4%
\end{array}%
\right) \right) \left( -1\right) \\
& =\left\vert \left( 
\begin{array}{c}
-8+7 \\ 
-4+4%
\end{array}%
\right) \right) =\left\vert \left( 
\begin{array}{c}
-1 \\ 
0%
\end{array}%
\right) \right)
\end{align}

\item 
\begin{align}
& \left\{ \left\vert \left( 
\begin{array}{c}
2 \\ 
1%
\end{array}%
\right) \right) \left( \left( 
\begin{array}{c}
7 \\ 
4%
\end{array}%
\right) \right\vert -\left\vert \left( 
\begin{array}{c}
7 \\ 
4%
\end{array}%
\right) \right) \left( \left( 
\begin{array}{c}
2 \\ 
1%
\end{array}%
\right) \right\vert \right\} ^{2}\left( 
\begin{array}{c}
0 \\ 
1%
\end{array}%
\right) \\
& =\left\{ \left\vert \left( 
\begin{array}{c}
2 \\ 
1%
\end{array}%
\right) \right) \left( \left( 
\begin{array}{c}
7 \\ 
4%
\end{array}%
\right) \right\vert -\left\vert \left( 
\begin{array}{c}
7 \\ 
4%
\end{array}%
\right) \right) \left( \left( 
\begin{array}{c}
2 \\ 
1%
\end{array}%
\right) \right\vert \right\} \left( 
\begin{array}{c}
1 \\ 
0%
\end{array}%
\right) =\left( 
\begin{array}{c}
0 \\ 
-1%
\end{array}%
\right)
\end{align}%
thus 
\begin{equation}
\left( tJ_{e}t^{t}\right) ^{2}=-I_{2}.
\end{equation}%
The Ket-Bra form of $\left( tJ_{e}t^{t}\right) ^{2}$ is 
\begin{align}
\left( tJ_{e}t^{t}\right) ^{2}& =\left( \left\vert \left( 
\begin{array}{c}
2 \\ 
1%
\end{array}%
\right) \right) \left( \left( 
\begin{array}{c}
7 \\ 
4%
\end{array}%
\right) \right\vert -\left\vert \left( 
\begin{array}{c}
7 \\ 
4%
\end{array}%
\right) \right) \left( \left( 
\begin{array}{c}
2 \\ 
1%
\end{array}%
\right) \right\vert \right) \\
& \times \left( \left\vert \left( 
\begin{array}{c}
2 \\ 
1%
\end{array}%
\right) \right) \left( \left( 
\begin{array}{c}
7 \\ 
4%
\end{array}%
\right) \right\vert -\left\vert \left( 
\begin{array}{c}
7 \\ 
4%
\end{array}%
\right) \right) \left( \left( 
\begin{array}{c}
2 \\ 
1%
\end{array}%
\right) \right\vert \right)
\end{align}%
\begin{align}
& =\left\vert \left( 
\begin{array}{c}
2 \\ 
1%
\end{array}%
\right) \right) \left( \left( 
\begin{array}{c}
7 \\ 
4%
\end{array}%
\right) \right\vert \left\vert \left( 
\begin{array}{c}
2 \\ 
1%
\end{array}%
\right) \right) \left( \left( 
\begin{array}{c}
7 \\ 
4%
\end{array}%
\right) \right\vert -\left\vert \left( 
\begin{array}{c}
2 \\ 
1%
\end{array}%
\right) \right) \left( \left( 
\begin{array}{c}
7 \\ 
4%
\end{array}%
\right) \right\vert \left\vert \left( 
\begin{array}{c}
7 \\ 
4%
\end{array}%
\right) \right) \left( \left( 
\begin{array}{c}
2 \\ 
1%
\end{array}%
\right) \right\vert \\
& -\left\vert \left( 
\begin{array}{c}
7 \\ 
4%
\end{array}%
\right) \right) \left( \left( 
\begin{array}{c}
2 \\ 
1%
\end{array}%
\right) \right\vert \left\vert \left( 
\begin{array}{c}
2 \\ 
1%
\end{array}%
\right) \right) \left( \left( 
\begin{array}{c}
7 \\ 
4%
\end{array}%
\right) \right\vert +\left\vert \left( 
\begin{array}{c}
7 \\ 
4%
\end{array}%
\right) \right) \left( \left( 
\begin{array}{c}
2 \\ 
1%
\end{array}%
\right) \right\vert \left\vert \left( 
\begin{array}{c}
7 \\ 
4%
\end{array}%
\right) \right) \left( \left( 
\begin{array}{c}
2 \\ 
1%
\end{array}%
\right) \right\vert
\end{align}%
\begin{align}
=& \left\vert \left( 
\begin{array}{c}
2 \\ 
1%
\end{array}%
\right) \right) 18\left( \left( 
\begin{array}{c}
7 \\ 
4%
\end{array}%
\right) \right\vert -\left\vert \left( 
\begin{array}{c}
2 \\ 
1%
\end{array}%
\right) \right) 65\left( \left( 
\begin{array}{c}
2 \\ 
1%
\end{array}%
\right) \right\vert \\
& -\left\vert \left( 
\begin{array}{c}
7 \\ 
4%
\end{array}%
\right) \right) 5\left( \left( 
\begin{array}{c}
7 \\ 
4%
\end{array}%
\right) \right\vert +\left\vert \left( 
\begin{array}{c}
7 \\ 
4%
\end{array}%
\right) \right) 18\left( \left( 
\begin{array}{c}
2 \\ 
1%
\end{array}%
\right) \right\vert
\end{align}%
\begin{equation}
=\left( \left\vert \left( 
\begin{array}{c}
2 \\ 
1%
\end{array}%
\right) \right) 18-\left\vert \left( 
\begin{array}{c}
7 \\ 
4%
\end{array}%
\right) \right) 5\right) \left( \left( 
\begin{array}{c}
7 \\ 
4%
\end{array}%
\right) \right\vert +\left( \left\vert \left( 
\begin{array}{c}
-2 \\ 
-1%
\end{array}%
\right) \right) 65+\left\vert \left( 
\begin{array}{c}
7 \\ 
4%
\end{array}%
\right) \right) 18\right) \left( \left( 
\begin{array}{c}
2 \\ 
1%
\end{array}%
\right) \right\vert
\end{equation}%
\begin{equation}
=\left( \left\vert \left( 
\begin{array}{c}
36 \\ 
18%
\end{array}%
\right) \right) +\left\vert \left( 
\begin{array}{c}
-35 \\ 
-20%
\end{array}%
\right) \right) \right) \left( \left( 
\begin{array}{c}
7 \\ 
4%
\end{array}%
\right) \right\vert +\left( \left\vert \left( 
\begin{array}{c}
-130 \\ 
-65%
\end{array}%
\right) \right) +\left\vert \left( 
\begin{array}{c}
126 \\ 
72%
\end{array}%
\right) \right) \right) \left( \left( 
\begin{array}{c}
2 \\ 
1%
\end{array}%
\right) \right\vert
\end{equation}%
\begin{equation}
=\left\vert \left( 
\begin{array}{c}
1 \\ 
-2%
\end{array}%
\right) \right) \left( \left( 
\begin{array}{c}
7 \\ 
4%
\end{array}%
\right) \right\vert +\left\vert \left( 
\begin{array}{c}
-4 \\ 
7%
\end{array}%
\right) \right) \left( \left( 
\begin{array}{c}
2 \\ 
1%
\end{array}%
\right) \right\vert .
\end{equation}
\end{enumerate}

The action of this form of $\left( tJ_{e}t^{t}\right) ^{2}$ on basis
elements is

\begin{enumerate}
\item 
\begin{align}
& \left\{ \left\vert \left( 
\begin{array}{c}
1 \\ 
-2%
\end{array}%
\right) \right) \left( \left( 
\begin{array}{c}
7 \\ 
4%
\end{array}%
\right) \right\vert +\left\vert \left( 
\begin{array}{c}
-4 \\ 
7%
\end{array}%
\right) \right) \left( \left( 
\begin{array}{c}
2 \\ 
1%
\end{array}%
\right) \right\vert \right\} \left( 
\begin{array}{c}
1 \\ 
0%
\end{array}%
\right) \\
& =\left\vert \left( 
\begin{array}{c}
1 \\ 
-2%
\end{array}%
\right) \right) 7+\left\vert \left( 
\begin{array}{c}
-4 \\ 
7%
\end{array}%
\right) \right) 2=\left\vert \left( 
\begin{array}{c}
-1 \\ 
0%
\end{array}%
\right) \right)
\end{align}

\item 
\begin{align}
& \left\{ \left\vert \left( 
\begin{array}{c}
1 \\ 
-2%
\end{array}%
\right) \right) \left( \left( 
\begin{array}{c}
7 \\ 
4%
\end{array}%
\right) \right\vert +\left\vert \left( 
\begin{array}{c}
-4 \\ 
7%
\end{array}%
\right) \right) \left( \left( 
\begin{array}{c}
2 \\ 
1%
\end{array}%
\right) \right\vert \right\} \left( 
\begin{array}{c}
0 \\ 
1%
\end{array}%
\right) \\
& =\left\vert \left( 
\begin{array}{c}
1 \\ 
-2%
\end{array}%
\right) \right) 4+\left\vert \left( 
\begin{array}{c}
-4 \\ 
7%
\end{array}%
\right) \right) 1=\left\vert \left( 
\begin{array}{c}
0 \\ 
-1%
\end{array}%
\right) \right) ,
\end{align}%
thus%
\begin{align}
& \left\vert \left( 
\begin{array}{c}
-1 \\ 
2%
\end{array}%
\right) \right) \left( \left( 
\begin{array}{c}
7 \\ 
4%
\end{array}%
\right) \right\vert +\left\vert \left( 
\begin{array}{c}
4 \\ 
-7%
\end{array}%
\right) \right) \left( \left( 
\begin{array}{c}
2 \\ 
1%
\end{array}%
\right) \right\vert \\
& =\left\vert \left( 
\begin{array}{c}
7 \\ 
4%
\end{array}%
\right) \right) \left( \left( 
\begin{array}{c}
-1 \\ 
2%
\end{array}%
\right) \right\vert +\left\vert \left( 
\begin{array}{c}
2 \\ 
1%
\end{array}%
\right) \right) \left( \left( 
\begin{array}{c}
4 \\ 
-7%
\end{array}%
\right) \right\vert =I_{2}
\end{align}%
as the identity is invariant under transpose.
\end{enumerate}

Denoting bases elements as 
\begin{align}
f_{1}^{d} & =\left( 
\begin{array}{c}
-1 \\ 
2%
\end{array}
\right) =-\left( 
\begin{array}{cc}
0 & 1 \\ 
-1 & 0%
\end{array}
\right) \left( 
\begin{array}{c}
2 \\ 
1%
\end{array}
\right) =-J_{e}f_{2}=-J_{e}te_{2} \\
f_{2}^{d} & =\left( 
\begin{array}{c}
4 \\ 
-7%
\end{array}
\right) =\left( 
\begin{array}{cc}
0 & 1 \\ 
-1 & 0%
\end{array}
\right) \left( 
\begin{array}{c}
7 \\ 
4%
\end{array}
\right) =J_{e}f_{1}=J_{e}te_{1} \\
f_{1} & =\left( 
\begin{array}{c}
7 \\ 
4%
\end{array}
\right) =t\left( 
\begin{array}{c}
1 \\ 
0%
\end{array}
\right) =te_{1} \\
f_{2} & =\left( 
\begin{array}{c}
2 \\ 
1%
\end{array}
\right) =t\left( 
\begin{array}{c}
0 \\ 
1%
\end{array}
\right) =te_{2}
\end{align}
one obtains \emph{for symplectic} $t$ 
\begin{align}
\left\vert f_{1}^{d}\right) \left( f_{1}\right\vert +\left\vert
f_{2}^{d}\right) \left( f_{2}\right\vert & =I_{2}=\left\vert
f_{1}^{d}\right) \left( te_{1}\right\vert +\left\vert f_{2}^{d}\right)
\left( te_{2}\right\vert =\left\vert te_{1}\right) \left(
f_{1}^{d}\right\vert +\left\vert te_{2}\right) \left( f_{2}^{d}\right\vert
=\left\vert f_{1}\right) \left( f_{1}^{d}\right\vert +\left\vert
f_{2}\right) \left( f_{2}^{d}\right\vert \\
& =-\left\vert te_{1}\right) \left( J_{e}te_{2}\right\vert +\left\vert
te_{2}\right) \left( J_{e}te_{1}\right\vert =\left\vert te_{2}\right) \left(
J_{e}te_{1}\right\vert -\left\vert te_{1}\right) \left(
J_{e}te_{2}\right\vert .
\end{align}

One should observe that $\left\{ f_{1},f_{2}\right\} ,$ $\left\{
f_{1}^{d},f_{2}^{d}\right\} $ are biorthogonal bases i.e. 
\begin{align}
& \left\vert f_{1}^{d}\right) \left( f_{1}\right\vert +\left\vert
f_{2}^{d}\right) \left( f_{2}\right\vert =I_{2} \\
& \left( \left. f_{1}\right\vert f_{2}^{d}\right) =\left( \left.
f_{1}^{d}\right\vert f_{2}\right) =\left( \left. f_{2}\right\vert
f_{1}^{d}\right) =\left( \left. f_{2}^{d}\right\vert f_{1}\right) =0 \\
& \left( \left. f_{1}\right\vert f_{1}^{d}\right) =\left( \left.
f_{1}^{d}\right\vert f_{1}\right) =\left( \left. f_{2}\right\vert
f_{2}^{d}\right) =\left( \left. f_{2}^{d}\right\vert f_{2}\right) =1.
\end{align}%
The resolution of the identity gives 
\begin{align}
& \left\{ \left\vert f_{1}^{d}\right) \left( f_{1}\right\vert +\left\vert
f_{2}^{d}\right) \left( f_{2}\right\vert \right\} \left\vert
f_{1}^{d}\right) =\left\vert f_{1}^{d}\right) \\
& \left\{ \left\vert f_{1}^{d}\right) \left( f_{1}\right\vert +\left\vert
f_{2}^{d}\right) \left( f_{2}\right\vert \right\} \left\vert
f_{2}^{d}\right) =\left\vert f_{2}^{d}\right)
\end{align}%
\begin{align}
& \left\{ \left\vert f_{1}\right) \left( f_{1}^{d}\right\vert +\left\vert
f_{2}\right) \left( f_{2}^{d}\right\vert \right\} \left\vert
f_{1}^{d}\right) =\left\vert f_{1}^{d}\right) \\
& \left\{ \left\vert f_{1}\right) \left( f_{1}^{d}\right\vert +\left\vert
f_{2}\right) \left( f_{2}^{d}\right\vert \right\} \left\vert
f_{2}^{d}\right) =\left\vert f_{2}^{d}\right)
\end{align}%
and 
\begin{align}
& \left\{ \left\vert f_{1}^{d}\right) \left( f_{1}\right\vert +\left\vert
f_{2}^{d}\right) \left( f_{2}\right\vert \right\} \left\vert f_{1}\right)
=\left\vert f_{1}\right) \\
& \left\{ \left\vert f_{1}^{d}\right) \left( f_{1}\right\vert +\left\vert
f_{2}^{d}\right) \left( f_{2}\right\vert \right\} \left\vert f_{2}\right)
=\left\vert f_{2}\right)
\end{align}%
\begin{align}
& \left\{ \left\vert f_{1}^{d}\right) \left( f_{1}\right\vert +\left\vert
f_{2}^{d}\right) \left( f_{2}\right\vert \right\} \left\vert f_{1}\right)
=\left\vert f_{1}\right) \\
& \left\{ \left\vert f_{1}^{d}\right) \left( f_{1}\right\vert +\left\vert
f_{2}^{d}\right) \left( f_{2}\right\vert \right\} \left\vert f_{2}\right)
=\left\vert f_{2}\right)
\end{align}

\begin{example}
\begin{equation}
f_{1}^{d}=\left( 
\begin{array}{c}
-1 \\ 
2%
\end{array}%
\right) ,f_{2}^{d}=\left( 
\begin{array}{c}
4 \\ 
-7%
\end{array}%
\right) ;f_{1}=\left( 
\begin{array}{c}
7 \\ 
4%
\end{array}%
\right) ^{\prime }f_{2}=\left( 
\begin{array}{c}
2 \\ 
1%
\end{array}%
\right)
\end{equation}%
\begin{align}
\left\vert f_{1}\right) \left( f_{1}^{d}\right\vert +\left\vert f_{2}\right)
\left( f_{2}^{d}\right\vert =& \left\{ \left\vert \left( 
\begin{array}{c}
7 \\ 
4%
\end{array}%
\right) \right) \left( \left( 
\begin{array}{c}
-1 \\ 
2%
\end{array}%
\right) \right\vert +\left\vert \left( 
\begin{array}{c}
2 \\ 
1%
\end{array}%
\right) \right) \left( \left( 
\begin{array}{c}
4 \\ 
-7%
\end{array}%
\right) \right\vert \right\} \\
& \leftrightarrow \left( 
\begin{array}{c}
7 \\ 
4%
\end{array}%
\right) \left( 
\begin{array}{cc}
-1 & 2%
\end{array}%
\right) +\left( 
\begin{array}{c}
2 \\ 
1%
\end{array}%
\right) \left( 
\begin{array}{cc}
4 & -7%
\end{array}%
\right) =\left( 
\begin{array}{cc}
1 & 0 \\ 
0 & 1%
\end{array}%
\right)
\end{align}%
and%
\begin{align}
\left\{ \left\vert f_{1}^{d}\right) \left( f_{1}\right\vert +\left\vert
f_{2}^{d}\right) \left( f_{2}\right\vert \right\} & =\left\{ \left\vert
\left( 
\begin{array}{c}
-1 \\ 
2%
\end{array}%
\right) \right) \left( \left( 
\begin{array}{c}
7 \\ 
4%
\end{array}%
\right) \right\vert +\left\vert \left( 
\begin{array}{c}
4 \\ 
-7%
\end{array}%
\right) \right) \left( \left( 
\begin{array}{c}
2 \\ 
1%
\end{array}%
\right) \right\vert \right\} \\
& \leftrightarrow \left( 
\begin{array}{c}
-1 \\ 
2%
\end{array}%
\right) \left( 
\begin{array}{cc}
7 & 4%
\end{array}%
\right) +\left( 
\begin{array}{c}
4 \\ 
-7%
\end{array}%
\right) \left( 
\begin{array}{cc}
2 & 1%
\end{array}%
\right) =\left( 
\begin{array}{cc}
1 & 0 \\ 
0 & 1%
\end{array}%
\right) .
\end{align}
\end{example}

These bases 
\begin{equation}
f_{1}^{d}=\left( 
\begin{array}{c}
-1 \\ 
2%
\end{array}%
\right) ,f_{2}^{d}=\left( 
\begin{array}{c}
4 \\ 
-7%
\end{array}%
\right) ;f_{1}=\left( 
\begin{array}{c}
7 \\ 
4%
\end{array}%
\right) ^{\prime }f_{2}=\left( 
\begin{array}{c}
2 \\ 
1%
\end{array}%
\right)
\end{equation}%
allows one to form the complexification 
\begin{align}
tJ_{e}t^{t}& =\left\vert te_{1}\right) \left( te_{2}\right\vert -\left\vert
te_{2}\right) \left( te_{1}\right\vert =\left\vert f_{2}\right) \left(
f_{1}\right\vert -\left\vert f_{1}\right) \left( f_{2}\right\vert =J_{f} \\
& =\left\vert \left( 
\begin{array}{c}
2 \\ 
1%
\end{array}%
\right) \right) \left( \left( 
\begin{array}{c}
7 \\ 
4%
\end{array}%
\right) \right\vert -\left\vert \left( 
\begin{array}{c}
7 \\ 
4%
\end{array}%
\right) \right) \left( \left( 
\begin{array}{c}
2 \\ 
1%
\end{array}%
\right) \right\vert
\end{align}%
corresponding to

\begin{enumerate}
\item 
\begin{align}
J_{e}f_{1}^{d}& =f_{2} \\
-J_{e}f_{2}^{d}& =f_{1}
\end{align}

\item 
\begin{equation}
J_{f}=\left\vert f_{2}\right) \left( f_{1}\right\vert -\left\vert
f_{1}\right) \left( f_{2}\right\vert =\left\vert J_{e}f_{1}^{d}\right)
\left( f_{2}\right\vert +\left\vert J_{e}f_{2}^{d}\right) \left(
f_{1}\right\vert =J_{e}\left( \left\vert f_{1}^{d}\right) \left(
f_{2}\right\vert +\left\vert f_{2}^{d}\right) \left( f_{1}\right\vert \right)
\end{equation}

\item 
\begin{align}
& \left( \left\vert f_{1}^{d}\right) \left( f_{2}\right\vert -\left\vert
f_{2}^{d}\right) \left( f_{1}\right\vert \right) \left( \left\vert
f_{1}^{d}\right) \left( f_{2}\right\vert -\left\vert f_{2}^{d}\right) \left(
f_{1}\right\vert \right) =\left\vert f_{1}^{d}\right) \left(
f_{2}\right\vert \left\vert f_{1}^{d}\right) \left( f_{2}\right\vert
-\left\vert f_{1}^{d}\right) \left( f_{2}\right\vert \left\vert
f_{2}^{d}\right) \left( f_{1}\right\vert -\left\vert f_{2}^{d}\right) \left(
f_{1}\right\vert \left\vert f_{1}^{d}\right) \left( f_{2}\right\vert
+\left\vert f_{2}^{d}\right) \left( f_{1}\right\vert \left\vert
f_{2}^{d}\right) \left( f_{1}\right\vert \\
& =-\left\vert f_{1}^{d}\right) \left( f_{1}\right\vert -\left\vert
f_{2}^{d}\right) \left( f_{2}\right\vert =-\left( \left\vert
f_{1}^{d}\right) \left( f_{1}\right\vert +\left\vert f_{2}^{d}\right) \left(
f_{2}\right\vert \right) =-I_{2}
\end{align}

\item 
\begin{equation}
\left( \left\vert f_{1}^{d}\right) \left( f_{2}\right\vert -\left\vert
f_{2}^{d}\right) \left( f_{1}\right\vert \right) ^{t}=\left\vert
f_{2}\right) \left( f_{1}^{d}\right\vert -\left\vert f_{1}\right) \left(
f_{2}^{d}\right\vert \neq -\left( \left\vert f_{1}^{d}\right) \left(
f_{2}\right\vert -\left\vert f_{2}^{d}\right) \left( f_{1}\right\vert \right)
\end{equation}
\end{enumerate}

\begin{example}
\begin{align}
& \left\vert f_{1}^{d}\right) \left( f_{2}\right\vert -\left\vert
f_{2}^{d}\right) \left( f_{1}\right\vert =\left\vert \left( 
\begin{array}{c}
-1 \\ 
2%
\end{array}%
\right) \right) \left( \left( 
\begin{array}{c}
2 \\ 
1%
\end{array}%
\right) \right\vert -\left\vert \left( 
\begin{array}{c}
4 \\ 
-7%
\end{array}%
\right) \right) \left( \left( 
\begin{array}{c}
7 \\ 
4%
\end{array}%
\right) \right\vert \\
& \leftrightarrow \left( 
\begin{array}{c}
-1 \\ 
2%
\end{array}%
\right) \left( 
\begin{array}{cc}
2 & 1%
\end{array}%
\right) -\left( 
\begin{array}{c}
4 \\ 
-7%
\end{array}%
\right) \left( 
\begin{array}{cc}
7 & 4%
\end{array}%
\right) \\
& \left( \left\vert \left( 
\begin{array}{c}
-1 \\ 
2%
\end{array}%
\right) \right) \left( \left( 
\begin{array}{c}
2 \\ 
1%
\end{array}%
\right) \right\vert -\left\vert \left( 
\begin{array}{c}
4 \\ 
-7%
\end{array}%
\right) \right) \left( \left( 
\begin{array}{c}
7 \\ 
4%
\end{array}%
\right) \right\vert \right) \left( \left\vert \left( 
\begin{array}{c}
-1 \\ 
2%
\end{array}%
\right) \right) \left( \left( 
\begin{array}{c}
2 \\ 
1%
\end{array}%
\right) \right\vert -\left\vert \left( 
\begin{array}{c}
4 \\ 
-7%
\end{array}%
\right) \right) \left( \left( 
\begin{array}{c}
7 \\ 
4%
\end{array}%
\right) \right\vert \right) \\
& =-\left( \left\vert \left( 
\begin{array}{c}
-1 \\ 
2%
\end{array}%
\right) \right) \left( \left( 
\begin{array}{c}
7 \\ 
4%
\end{array}%
\right) \right\vert +\left\vert \left( 
\begin{array}{c}
4 \\ 
-7%
\end{array}%
\right) \right) \left( \left( 
\begin{array}{c}
2 \\ 
1%
\end{array}%
\right) \right\vert \right)
\end{align}%
The basis 
\begin{equation}
\left\{ f_{1}^{d}=\left( 
\begin{array}{c}
-1 \\ 
2%
\end{array}%
\right) ,f_{2}^{d}=\left( 
\begin{array}{c}
4 \\ 
-7%
\end{array}%
\right) \right\} ,\left\{ f_{1}=\left( 
\begin{array}{c}
7 \\ 
4%
\end{array}%
\right) ,f_{2}=\left( 
\begin{array}{c}
2 \\ 
1%
\end{array}%
\right) \right\}
\end{equation}%
of $\mathbb{R}^{2}$ could be used to form the direct sum 
\begin{equation}
\mathcal{V}=V_{1}\oplus V_{2}=\mathbb{R}^{2}\oplus \mathbb{R}%
^{2}=V_{1}\oplus V_{1}^{d}=\mathcal{V}_{1}+\mathcal{V}_{1}^{d}\cong \mathbb{R%
}^{4},
\end{equation}%
but it is more interesting to consider them as two non orthogonal and
unnormalized basis of $V,$ that are biorthogonal to each other, where 
\begin{align}
V& =\mathbb{R}-\limfunc{linspan}\left\{ f_{1}^{d},f_{2}^{d}\right\} =\mathbb{%
R}-\limfunc{linspan}\left\{ \left( 
\begin{array}{c}
-1 \\ 
2%
\end{array}%
\right) ,\left( 
\begin{array}{c}
4 \\ 
-7%
\end{array}%
\right) \right\} =V_{1} \\
& =\mathbb{R}-\limfunc{linspan}\left\{ f_{1},f_{2}\right\} =\mathbb{R}-%
\limfunc{linspan}\left\{ \left( 
\begin{array}{c}
7 \\ 
4%
\end{array}%
\right) ,\left( 
\begin{array}{c}
2 \\ 
1%
\end{array}%
\right) \right\} =V_{2}=V_{1}^{d}.
\end{align}%
Of course 
\begin{equation}
V=V_{1}=V_{2}=V_{1}^{d}\cong \mathbb{R}^{2}.
\end{equation}%
One can consider the complexification 
\begin{equation}
J_{V_{1}}=\left\vert f_{1}\right) \left( f_{2}^{d}\right\vert -\left\vert
f_{2}\right) \left( f_{1}^{d}\right\vert
\end{equation}%
that satisfies the two complexifiction conditions 
\begin{equation}
\left( \left\vert f_{1}\right) \left( f_{2}^{d}\right\vert -\left\vert
f_{2}\right) \left( f_{1}^{d}\right\vert \right) \left( \left\vert
f_{1}\right) \left( f_{2}^{d}\right\vert -\left\vert f_{2}\right) \left(
f_{1}^{d}\right\vert \right) =-\left\vert f_{1}\right) \left(
f_{1}^{d}\right\vert -\left\vert f_{2}\right) \left( f_{2}^{d}\right\vert
=-I_{2r}
\end{equation}%
and%
\begin{equation}
\left\{ f_{1}^{d}=\left( 
\begin{array}{c}
-1 \\ 
2%
\end{array}%
\right) ,f_{2}^{d}=\left( 
\begin{array}{c}
4 \\ 
-7%
\end{array}%
\right) \right\} ,\left\{ f_{1}=\left( 
\begin{array}{c}
7 \\ 
4%
\end{array}%
\right) ,f_{2}=\left( 
\begin{array}{c}
2 \\ 
1%
\end{array}%
\right) \right\}
\end{equation}%
where%
\begin{align}
J_{V_{1}}& \in \mathcal{L}\left( V_{1}\right) ,\mathcal{L}\left(
V_{1},V_{1}^{d}\right) \\
J_{V_{1}}^{t}& \in \mathcal{L}\left( V_{1}\right) ,\mathcal{L}\left(
V_{1}^{d},V_{1}\right) .
\end{align}%
The transpose gives 
\begin{equation}
\left( \left\vert f_{1}\right) \left( f_{2}^{d}\right\vert -\left\vert
f_{2}\right) \left( f_{1}^{d}\right\vert \right) ^{t}=\left\vert
f_{2}^{d}\right) \left( f_{1}\right\vert -\left\vert f_{1}^{d}\right) \left(
f_{2}\right\vert =-\left( \left\vert f_{1}\right) \left(
f_{2}^{d}\right\vert -\left\vert f_{2}\right) \left( f_{1}^{d}\right\vert
\right) .
\end{equation}%
$i.e.$%
\begin{equation}
\left( \left\vert \left( 
\begin{array}{c}
7 \\ 
4%
\end{array}%
\right) \right) \left( \left( 
\begin{array}{c}
4 \\ 
-7%
\end{array}%
\right) \right\vert -\left\vert \left( 
\begin{array}{c}
2 \\ 
1%
\end{array}%
\right) \right) \left( \left( 
\begin{array}{c}
-1 \\ 
2%
\end{array}%
\right) \right\vert \right) ^{t}=\left\vert \left( 
\begin{array}{c}
4 \\ 
-7%
\end{array}%
\right) \right) \left( \left( 
\begin{array}{c}
7 \\ 
4%
\end{array}%
\right) \right\vert -\left\vert \left( 
\begin{array}{c}
-1 \\ 
2%
\end{array}%
\right) \right) \left( \left( 
\begin{array}{c}
2 \\ 
1%
\end{array}%
\right) \right\vert
\end{equation}
\end{example}

The complexification of $V$ to give $V_{\mathbb{C}},$given by $J_{V_{1}}$
(i.e. based on $V_{1})$ is based on 
\begin{equation}
V_{1}^{d}=J_{V_{1}}V_{1}
\end{equation}

e.g. 
\begin{align}
f_{1}^{d} & =\left( 
\begin{array}{c}
-1 \\ 
2%
\end{array}
\right) =-\left( 
\begin{array}{cc}
0 & 1 \\ 
-1 & 0%
\end{array}
\right) \left( 
\begin{array}{c}
2 \\ 
1%
\end{array}
\right) =-J_{V_{1}}f_{2}=-J_{V_{1}}te_{2} \\
f_{2}^{d} & =\left( 
\begin{array}{c}
4 \\ 
-7%
\end{array}
\right) =\left( 
\begin{array}{cc}
0 & 1 \\ 
-1 & 0%
\end{array}
\right) \left( 
\begin{array}{c}
7 \\ 
4%
\end{array}
\right) =J_{V_{1}}f_{1}=J_{V_{1}}te_{1}
\end{align}%
\begin{align}
f_{2} & =te_{2}=-J_{V_{1}}f_{1}^{d}=-\left( 
\begin{array}{cc}
0 & 1 \\ 
-1 & 0%
\end{array}
\right) \left( 
\begin{array}{c}
-1 \\ 
2%
\end{array}
\right) =\left( 
\begin{array}{c}
2 \\ 
1%
\end{array}
\right) \\
f_{1} & =te_{1}=-J_{V_{1}}f_{2}^{d}=\left( 
\begin{array}{c}
4 \\ 
-7%
\end{array}
\right) =\left( 
\begin{array}{c}
7 \\ 
4%
\end{array}
\right) .
\end{align}
The complexification is explicitly given by 
\begin{equation}
V_{\mathbb{C}}\cong V_{1}+J_{V_{1}}V_{1}
\end{equation}
that corresponds to $V_{1}$ and $J_{V_{1}}V_{1}$ being disjoint spaces and 
\begin{equation}
V_{\mathbb{C}}=V_{1}\oplus J_{V_{1}}V_{1}
\end{equation}
and making identiifications using the natural isomorphisms. $V_{1}$ and $%
J_{V_{1}}V_{1}$ are \emph{always disjoint subspaces }by construction and
definition although they are isomorphic.

\begin{summary}
We consider the complexification 
\begin{align}
J_{f}& =\left\vert \left( 
\begin{array}{c}
2 \\ 
1%
\end{array}%
\right) \right) \left( \left( 
\begin{array}{c}
7 \\ 
4%
\end{array}%
\right) \right\vert -\left\vert \left( 
\begin{array}{c}
7 \\ 
4%
\end{array}%
\right) \right) \left( \left( 
\begin{array}{c}
2 \\ 
1%
\end{array}%
\right) \right\vert \\
& =\left\vert \left( 
\begin{array}{cc}
0 & 1 \\ 
-1 & 0%
\end{array}%
\right) \left( 
\begin{array}{c}
-1 \\ 
2%
\end{array}%
\right) \right) \left( \left( 
\begin{array}{c}
7 \\ 
4%
\end{array}%
\right) \right\vert -\left\vert \left( 
\begin{array}{cc}
0 & 1 \\ 
-1 & 0%
\end{array}%
\right) \left( 
\begin{array}{c}
4 \\ 
-7%
\end{array}%
\right) \right) \left( \left( 
\begin{array}{c}
2 \\ 
1%
\end{array}%
\right) \right\vert \\
& =\left( 
\begin{array}{cc}
0 & 1 \\ 
-1 & 0%
\end{array}%
\right) \left( \left\vert \left( 
\begin{array}{c}
-1 \\ 
2%
\end{array}%
\right) \right) \left( \left( 
\begin{array}{c}
7 \\ 
4%
\end{array}%
\right) \right\vert -\left\vert \left( 
\begin{array}{c}
4 \\ 
-7%
\end{array}%
\right) \right) \left( \left( 
\begin{array}{c}
2 \\ 
1%
\end{array}%
\right) \right\vert \right) \\
& =J_{e}I_{2}
\end{align}%
formed from the biorthogonal bases 
\begin{equation}
\left\{ \left( 
\begin{array}{c}
-1 \\ 
2%
\end{array}%
\right) ,\left( 
\begin{array}{c}
4 \\ 
-7%
\end{array}%
\right) \right\} ,\left\{ \left( 
\begin{array}{c}
7 \\ 
4%
\end{array}%
\right) ,\left( 
\begin{array}{c}
2 \\ 
1%
\end{array}%
\right) \right\} ,
\end{equation}%
corresponding to the spaces%
\begin{align}
V& =\mathbb{R}-\limfunc{linspan}\left\{ f_{1}^{d},f_{2}^{d}\right\} =\mathbb{%
R}-\limfunc{linspan}\left\{ \left( 
\begin{array}{c}
-1 \\ 
2%
\end{array}%
\right) ,\left( 
\begin{array}{c}
4 \\ 
-7%
\end{array}%
\right) \right\} =V_{1} \\
& =\mathbb{R}-\limfunc{linspan}\left\{ f_{1},f_{2}\right\} =\mathbb{R}-%
\limfunc{linspan}\left\{ \left( 
\begin{array}{c}
7 \\ 
4%
\end{array}%
\right) ,\left( 
\begin{array}{c}
2 \\ 
1%
\end{array}%
\right) \right\} =V_{2}=V_{1}^{d}.
\end{align}%
$\lceil $%
\begin{align}
& \left( \left\vert \left( 
\begin{array}{c}
2 \\ 
1%
\end{array}%
\right) \right) \left( \left( 
\begin{array}{c}
7 \\ 
4%
\end{array}%
\right) \right\vert -\left\vert \left( 
\begin{array}{c}
7 \\ 
4%
\end{array}%
\right) \right) \left( \left( 
\begin{array}{c}
2 \\ 
1%
\end{array}%
\right) \right\vert \right) \left( \left\vert \left( 
\begin{array}{c}
2 \\ 
1%
\end{array}%
\right) \right) \left( \left( 
\begin{array}{c}
7 \\ 
4%
\end{array}%
\right) \right\vert -\left\vert \left( 
\begin{array}{c}
7 \\ 
4%
\end{array}%
\right) \right) \left( \left( 
\begin{array}{c}
2 \\ 
1%
\end{array}%
\right) \right\vert \right) \\
& =-\left( \left\vert \left( 
\begin{array}{c}
-1 \\ 
2%
\end{array}%
\right) \right) \left( \left( 
\begin{array}{c}
7 \\ 
4%
\end{array}%
\right) \right\vert +\left\vert \left( 
\begin{array}{c}
4 \\ 
-7%
\end{array}%
\right) \right) \left( \left( 
\begin{array}{c}
2 \\ 
1%
\end{array}%
\right) \right\vert \right) =-I_{2}
\end{align}%
\begin{align}
& \left( \left\vert \left( 
\begin{array}{c}
-1 \\ 
2%
\end{array}%
\right) \right) \left( \left( 
\begin{array}{c}
2 \\ 
1%
\end{array}%
\right) \right\vert -\left\vert \left( 
\begin{array}{c}
4 \\ 
-7%
\end{array}%
\right) \right) \left( \left( 
\begin{array}{c}
7 \\ 
4%
\end{array}%
\right) \right\vert \right) \left( \left\vert \left( 
\begin{array}{c}
-1 \\ 
2%
\end{array}%
\right) \right) \left( \left( 
\begin{array}{c}
2 \\ 
1%
\end{array}%
\right) \right\vert -\left\vert \left( 
\begin{array}{c}
4 \\ 
-7%
\end{array}%
\right) \right) \left( \left( 
\begin{array}{c}
7 \\ 
4%
\end{array}%
\right) \right\vert \right) \\
& =-\left( \left\vert \left( 
\begin{array}{c}
-1 \\ 
2%
\end{array}%
\right) \right) \left( \left( 
\begin{array}{c}
7 \\ 
4%
\end{array}%
\right) \right\vert +\left\vert \left( 
\begin{array}{c}
4 \\ 
-7%
\end{array}%
\right) \right) \left( \left( 
\begin{array}{c}
2 \\ 
1%
\end{array}%
\right) \right\vert \right) =-I_{2}
\end{align}%
\begin{align}
& \left( \left\vert \left( 
\begin{array}{c}
-1 \\ 
2%
\end{array}%
\right) \right) \left( \left( 
\begin{array}{c}
7 \\ 
4%
\end{array}%
\right) \right\vert +\left\vert \left( 
\begin{array}{c}
4 \\ 
-7%
\end{array}%
\right) \right) \left( \left( 
\begin{array}{c}
2 \\ 
1%
\end{array}%
\right) \right\vert \right) \left( \left\vert \left( 
\begin{array}{c}
-1 \\ 
2%
\end{array}%
\right) \right) \left( \left( 
\begin{array}{c}
7 \\ 
4%
\end{array}%
\right) \right\vert +\left\vert \left( 
\begin{array}{c}
4 \\ 
-7%
\end{array}%
\right) \right) \left( \left( 
\begin{array}{c}
2 \\ 
1%
\end{array}%
\right) \right\vert \right) \\
& =\left( \left\vert \left( 
\begin{array}{c}
-1 \\ 
2%
\end{array}%
\right) \right) \left( \left( 
\begin{array}{c}
7 \\ 
4%
\end{array}%
\right) \right\vert +\left\vert \left( 
\begin{array}{c}
4 \\ 
-7%
\end{array}%
\right) \right) \left( \left( 
\begin{array}{c}
2 \\ 
1%
\end{array}%
\right) \right\vert \right) =I_{2}
\end{align}%
\begin{equation}
\left( \left\vert \left( 
\begin{array}{c}
-1 \\ 
2%
\end{array}%
\right) \right) \left( \left( 
\begin{array}{c}
2 \\ 
1%
\end{array}%
\right) \right\vert -\left\vert \left( 
\begin{array}{c}
4 \\ 
-7%
\end{array}%
\right) \right) \left( \left( 
\begin{array}{c}
7 \\ 
4%
\end{array}%
\right) \right\vert \right) ^{t}=\left\vert \left( 
\begin{array}{c}
2 \\ 
1%
\end{array}%
\right) \right) \left( \left( 
\begin{array}{c}
-1 \\ 
2%
\end{array}%
\right) -\left\vert \left( 
\begin{array}{c}
7 \\ 
4%
\end{array}%
\right) \right) \left( \left( 
\begin{array}{c}
4 \\ 
-7%
\end{array}%
\right) \right) \right\vert
\end{equation}%
$\rfloor $
\end{summary}

The transpose 
\begin{equation}
J_{V_{1}}^{t}=J_{V_{1}^{d}}
\end{equation}
is also a complexification $V$ based on $V_{1}^{d}.$

\begin{example}
The latter identities give for example 
\begin{align}
\left\{ \left\vert \left( 
\begin{array}{c}
-1 \\ 
2%
\end{array}%
\right) \right) \left( \left( 
\begin{array}{c}
7 \\ 
4%
\end{array}%
\right) \right\vert +\left\vert \left( 
\begin{array}{c}
4 \\ 
-7%
\end{array}%
\right) \right) \left( \left( 
\begin{array}{c}
2 \\ 
1%
\end{array}%
\right) \right\vert \right\} \left\vert \left( 
\begin{array}{c}
7 \\ 
4%
\end{array}%
\right) \right) & =\left\vert \left( 
\begin{array}{c}
7 \\ 
4%
\end{array}%
\right) \right) \\
\left\{ \left\vert \left( 
\begin{array}{c}
-1 \\ 
2%
\end{array}%
\right) \right) 65+\left\vert \left( 
\begin{array}{c}
4 \\ 
-7%
\end{array}%
\right) \right) 18\right\} & =\left\vert \left( 
\begin{array}{c}
7 \\ 
4%
\end{array}%
\right) \right) .
\end{align}
\end{example}

\begin{remark}
dyadic expansions%
\begin{equation}
\left\vert \mathbf{e}\right) \left( \mathbf{e}\right\vert \mathcal{R}%
_{e}\left( \left( tJ_{e}t^{t}\right) ^{2}\right) \left\vert \mathbf{e}%
\right) \left( \mathbf{e}\right\vert =-\left\vert \mathbf{e}\right) \left( 
\mathbf{e}\right\vert
\end{equation}%
\begin{align}
\left\vert \mathbf{f}\right) \Delta _{f}^{-1}\left( \mathbf{f}\right\vert
A\left\vert \mathbf{f}\right) \Delta _{f}^{-1}\left( \mathbf{f}\right\vert &
=\left\vert \mathbf{f}\right) \Delta _{f}^{-1}\mathcal{M}_{f}\left( A\right)
\Delta _{f}^{-1}\left( \mathbf{f}\right\vert \\
& =\left\vert \mathbf{f}\right) \Delta _{f}^{-1}\Delta _{f}\mathcal{R}%
_{f}\left( A\right) \Delta _{f}^{-1}\left( \mathbf{f}\right\vert =\left\vert 
\mathbf{f}\right) \mathcal{R}_{f}\left( A\right) \Delta _{f}^{-1}\left( 
\mathbf{f}\right\vert
\end{align}%
Hence%
\begin{equation}
A=\left\vert \mathbf{f}\right) \mathcal{R}_{f}\left( A\right) \Delta
_{f}^{-1}\left( \mathbf{f}\right\vert =\left\vert \mathbf{f}\right) \Delta
_{f}^{-1}\mathcal{M}_{f}\left( A\right) \Delta _{f}^{-1}\left( \mathbf{f}%
\right\vert .
\end{equation}
\end{remark}

\begin{remark}
\begin{enumerate}
\item 
\begin{align}
\left( 
\begin{array}{c}
1 \\ 
-2%
\end{array}%
\right) & =\left( 
\begin{array}{cc}
0 & 1 \\ 
-1 & 0%
\end{array}%
\right) \left( 
\begin{array}{c}
2 \\ 
1%
\end{array}%
\right) ;\left( 
\begin{array}{c}
4 \\ 
-7%
\end{array}%
\right) =\left( 
\begin{array}{cc}
0 & 1 \\ 
-1 & 0%
\end{array}%
\right) \left( 
\begin{array}{c}
7 \\ 
4%
\end{array}%
\right) \\
-f_{1}^{d}& =J_{e}f_{2};f_{2}^{d}=J_{e}f_{1}
\end{align}

\item 
\begin{equation}
\left( tJ_{e}t^{t}\right) \left( tJ_{e}t^{t}\right) ^{t}=I
\end{equation}

\item 
\begin{equation}
tJ_{e}t^{t}=\left\vert \left( 
\begin{array}{c}
2 \\ 
1%
\end{array}%
\right) \right) \left( \left( 
\begin{array}{c}
7 \\ 
4%
\end{array}%
\right) \right\vert -\left\vert \left( 
\begin{array}{c}
7 \\ 
4%
\end{array}%
\right) \right) \left( \left( 
\begin{array}{c}
2 \\ 
1%
\end{array}%
\right) \right\vert
\end{equation}

\item 
\begin{align}
tJ_{e}e_{1}& =\left( 
\begin{array}{cc}
-7 & 2 \\ 
-4 & 1%
\end{array}%
\right) \left( 
\begin{array}{c}
1 \\ 
0%
\end{array}%
\right) =\left( 
\begin{array}{c}
-7 \\ 
-4%
\end{array}%
\right) =-f_{2} \\
tJ_{e}e_{2}& =\left( 
\begin{array}{cc}
-7 & 2 \\ 
-4 & 1%
\end{array}%
\right) \left( 
\begin{array}{c}
0 \\ 
1%
\end{array}%
\right) =\left( 
\begin{array}{c}
2 \\ 
1%
\end{array}%
\right) =f_{1}
\end{align}

\item 
\begin{align}
te_{1}& =\left( 
\begin{array}{cc}
2 & 7 \\ 
1 & 4%
\end{array}%
\right) \left( 
\begin{array}{c}
1 \\ 
0%
\end{array}%
\right) =\left( 
\begin{array}{c}
2 \\ 
1%
\end{array}%
\right) =f_{1} \\
te_{2}& =\left( 
\begin{array}{cc}
2 & 7 \\ 
1 & 4%
\end{array}%
\right) \left( 
\begin{array}{c}
0 \\ 
1%
\end{array}%
\right) =\left( 
\begin{array}{c}
7 \\ 
4%
\end{array}%
\right) =f_{2}
\end{align}

\item 
\begin{align}
& tJ_{e}t^{t}J_{e}=J_{e}J_{e}=-I=-tJ_{e}t^{t}J_{e}^{t} \\
& \left( 
\begin{array}{cc}
2 & 7 \\ 
1 & 4%
\end{array}%
\right) \left( 
\begin{array}{cc}
0 & 1 \\ 
-1 & 0%
\end{array}%
\right) \left( 
\begin{array}{cc}
2 & 1 \\ 
7 & 4%
\end{array}%
\right) \left( 
\begin{array}{cc}
0 & -1 \\ 
1 & 0%
\end{array}%
\right) =\left( 
\begin{array}{cc}
1 & 0 \\ 
0 & 1%
\end{array}%
\right) \\
& \left( 
\begin{array}{cc}
-7 & 2 \\ 
-4 & 1%
\end{array}%
\right) \left( 
\begin{array}{cc}
1 & -2 \\ 
4 & -7%
\end{array}%
\right) =\left( 
\begin{array}{cc}
1 & 0 \\ 
0 & 1%
\end{array}%
\right) \\
& \left( tJ_{e}\right) \left( t^{t}J_{e}^{t}\right) =\left( tJ_{e}\right)
\left( J_{e}t\right) ^{t}=I \\
& \left( J_{e}t\right) ^{t}=\left( J_{e}t\right)
^{-1}=t^{-1}J_{e}^{t}=-t^{-1}J_{e}. \\
& J_{e}t=J_{e}t^{-t}
\end{align}

\item 
\begin{align}
f_{1}^{d}& =\left( 
\begin{array}{c}
-1 \\ 
2%
\end{array}%
\right) =-\left( 
\begin{array}{cc}
0 & 1 \\ 
-1 & 0%
\end{array}%
\right) \left( 
\begin{array}{c}
2 \\ 
1%
\end{array}%
\right) =-J_{e}f_{2}=-J_{e}te_{2} \\
f_{2}^{d}& =\left( 
\begin{array}{c}
4 \\ 
-7%
\end{array}%
\right) =\left( 
\begin{array}{cc}
0 & 1 \\ 
-1 & 0%
\end{array}%
\right) \left( 
\begin{array}{c}
7 \\ 
4%
\end{array}%
\right) =J_{e}f_{1}=J_{e}te_{1} \\
f_{1}& =\left( 
\begin{array}{c}
7 \\ 
4%
\end{array}%
\right) =t\left( 
\begin{array}{c}
1 \\ 
0%
\end{array}%
\right) =te_{1} \\
f_{2}& =\left( 
\begin{array}{c}
2 \\ 
1%
\end{array}%
\right) =t\left( 
\begin{array}{c}
0 \\ 
1%
\end{array}%
\right) =te_{2}
\end{align}
\end{enumerate}
\end{remark}

\section{Free Vector Space}

\qquad The free vector space $\mathcal{V}$ generated by a set $\mathcal{S}$
generated by $\left\{ \left. s_{j}\right\vert 1\leq j\leq\kappa\right\} $ ($%
\kappa$ is the cardenality of the set if it has no structure) over a field $%
\mathbb{F}$ is given by

\begin{equation}
\mathcal{V}=\limfunc{lincom}\left\{ \mathcal{S}\right\} =\left\{
\sum\limits_{1\leq j\leq\kappa}c_{j}s_{j}\right\} ,
\end{equation}
where $c_{j}\in\mathbb{F}$ and $\sum$ denotes formal sums (i.e. convergence
is ignored) and multipliction and addition have the properties 
\begin{gather}
cs=c\sum\limits_{1\leq j\leq\kappa}c_{j}s_{j}=\sum\limits_{1\leq j\leq\kappa
}cc_{j}s_{j};\quad c\in\mathbb{F} \\
c_{1}s+c_{2}s=\left( c_{1}+c_{2}\right) s.
\end{gather}

\subsection{Direct Sum}

Consider two vector spaces $V_{1}$ and $V_{2}$ then their \emph{direct
(Cartesian) product} 
\begin{equation}
V_{1}\times V_{2}=\left\{ \left. \left( v_{1},v_{2}\right) \right\vert
v_{1}\in V_{1},v_{2}\in V_{2}\right\} ,
\end{equation}
can be endowered with a vector space structure by considering the \emph{free
vector space} over $V_{1}\times V_{2}$ (formal linear combinations of $%
\left( v_{1},v_{2}\right) )$ and defining the equivalences 
\begin{align}
\lambda\left( v_{1},v_{2}\right) & \sim\left( \lambda v_{1},\lambda
v_{2}\right) \\
\left( v_{1},v_{2}\right) +\left( v_{3},v_{4}\right) & \sim\left(
v_{1}+v_{3},v_{2}+v_{4}\right) \\
\lambda_{1}\left( v_{1},v_{2}\right) +\lambda_{2}\left( v_{3},v_{4}\right) &
\sim\left( \lambda_{1}v_{1}+\lambda_{2}v_{3},\lambda_{1}v_{2}+\lambda
_{2}v_{4}\right)
\end{align}
thus producing the direct sum vector space 
\begin{equation}
V=V_{1}\oplus_{\mathbb{F}}V_{2};\lambda\in\mathbb{F}.
\end{equation}

The scalars $\lambda$ in our case are real or complex, leading to $\oplus_{%
\mathbb{R}}$and $\oplus_{\mathbb{C}}.$ Other mathematical fields, $\mathbb{F}%
,$ could be used in the definition of vector spaces and direct sums (and
tensor products) and if $\lambda$ belongs to a ring then modules are formed.
One can consider real vector subspaces and complex dirct sums i.e. 
\begin{align}
\lambda_{1}\left( v_{1},v_{2}\right) +\lambda_{2}\left( v_{3},v_{4}\right) &
\sim\left( \lambda_{1}v_{1}+\lambda_{2}v_{3},\lambda_{1}v_{2}+\lambda
_{2}v_{4}\right) \\
\lambda_{1},\lambda_{2} & \in\mathbb{C},
\end{align}
which makes $V_{1}$ and $V_{2}$ complex spaces, so 
\begin{equation}
V_{1\mathbb{R}}\oplus_{\mathbb{C}}V_{2\mathbb{R}}=V_{1\mathbb{C}}\oplus_{%
\mathbb{C}}V_{2\mathbb{C}},
\end{equation}
one can also see that 
\begin{equation}
V_{1\mathbb{C}}\oplus_{\mathbb{R}}V_{2\mathbb{C}}=V_{1\mathbb{C}}\oplus_{%
\mathbb{C}}V_{2\mathbb{C}}.
\end{equation}

The subspaces $\widetilde{V_{1}}=\left\{ \left. \left( v_{1},0\right)
\right\vert v_{1}\in V_{1}\right\} $ and $\widetilde{V_{2}}=\left\{ \left.
\left( 0,v_{2}\right) \right\vert v_{2}\in V_{2}\right\} $ of $V$ are \emph{%
disjoint} and are isomorphic to $V_{1}$ and $V_{2}$ i.e. 
\begin{align}
& \widetilde{V_{1}}=\left\{ \left. \left( v_{1},0\right) \right\vert
v_{1}\in V_{1}\right\} \cong V_{1} \\
& \widetilde{V_{2}}=\left\{ \left. \left( 0,v_{2}\right) \right\vert
v_{2}\in V_{2}\right\} \cong V_{2} \\
& \left\{ \left. \left( v_{1},0\right) \right\vert v_{1}\in V_{1}\right\}
\cap\left\{ \left. \left( 0,v_{2}\right) \right\vert v_{2}\in V_{2}\right\}
=\left\{ \left( 0,0\right) \right\} .
\end{align}
\emph{Hence the explicit form of the direct sum is} 
\begin{equation}
V=\left\{ \left. \left( v_{1},v_{2}\right) \right\vert v_{1}\in
V_{1},v_{2}\in V_{2}\right\} =V_{1}\oplus_{\mathbb{F}}V_{2}.
\end{equation}
The spaces $V_{1}$ and $V_{2}$ could be subspaces (not neccesarily disjoint)
of a third vector space $W.$

In the particular case when $V_{2}=V_{1}$ we have%
\begin{equation}
V=V_{1}\oplus V_{1}=\left\{ \left. \left( v_{1},v_{2}\right) \right\vert
v_{1}\in V_{1},v_{2}\in V_{1}\right\}
\end{equation}
and%
\begin{align}
& \widetilde{V_{1}}=\left\{ \left. \left( v_{1},0\right) \right\vert
v_{1}\in V_{1}\right\} \cong V_{1} \\
& \widetilde{V_{2}}=\left\{ \left. \left( 0,v_{2}\right) \right\vert
v_{2}\in V_{1}\right\} \cong V_{1} \\
& \left\{ \left. \left( v_{1},0\right) \right\vert v_{1}\in V_{1}\right\}
\cap\left\{ \left. \left( 0,v_{2}\right) \right\vert v_{2}\in V_{1}\right\}
=\left\{ \left( 0,0\right) \right\} ,
\end{align}
\emph{BUT }in this case the spaces $V_{1}$ and $V_{2}$ are not disjoint
while $\widetilde{V_{1}}$ and $\widetilde{V_{2}}$ always are i.e. \emph{\ }%
\begin{align}
V_{1}\cap V_{1} & =V_{1} \\
\widetilde{V_{1}}\cap\widetilde{V_{2}} & =\left\{ \left( 0,0\right) \right\}
.
\end{align}
Note that 
\begin{align}
V & \neq\widetilde{V_{1}}\oplus\widetilde{V_{2}}=\left\{ \left. \overset{}{%
\left( \left( v_{1},0\right) ,\left( 0,v_{2}\right) \right) }\right\vert
\left( v_{1},0\right) \in\widetilde{V_{1}},\left( 0,v_{2}\right) \in%
\widetilde{V_{2}}\right\} , \\
& \text{but} \\
V & \cong\widetilde{V_{1}}\oplus\widetilde{V_{2}}\quad\text{i.e. }%
v_{1}\longleftrightarrow\left( v_{1},0\right) \text{ and }%
v_{2}\longleftrightarrow\left( 0,v_{2}\right) .
\end{align}
It is always true 
\begin{equation}
\widetilde{V_{1}}\oplus\widetilde{V_{2}}\cong\widetilde{V_{1}}+\widetilde{%
V_{2}}
\end{equation}
i.e. 
\begin{equation}
\left( \left( v_{1},v_{2}\right) ,\left( w_{1},w_{2}\right) \right)
\longleftrightarrow\left( v_{1},v_{2}\right) +\left( w_{1},w_{2}\right)
;v_{1},w_{1}\in V_{1},v_{2},w_{2}\in V_{2}
\end{equation}
and in particular when $V_{1}=V_{2}$ 
\begin{equation}
\left( \left( v_{1},v_{2}\right) ,\left( w_{1},w_{2}\right) \right)
\longleftrightarrow\left( v_{1},v_{2}\right) +\left( w_{1},w_{2}\right)
;v_{1},w_{1}\in V_{1},v_{2},w_{2}\in V_{1}.
\end{equation}

\subsection{Sum of Subspaces}

The \emph{sum of two subspaces} 
\begin{equation}
W=V_{1}+V_{2}=\left\{ \left. v_{1}+v_{2}\right\vert v_{1}\in V_{1},v_{2}\in
V_{2}\right\} =V_{1}\cup V_{2}
\end{equation}
is not isomorphic to the direct sum $V_{1}\oplus V_{2}$ unless $V_{1}$ and $%
V_{2}$ are disjoint. In particular 
\begin{align}
& V+V=V, \\
& \dim\left\{ V\right\} =\dim\left\{ V\right\} =r
\end{align}
but 
\begin{align}
& V\oplus V\supset\mathcal{V}\cong V \\
& \dim\left\{ V\oplus V\right\} =2r.
\end{align}

One can see this by observing the following special cases that can easily be
generalized

\begin{align}
V_{1} & =\limfunc{linspan}\left\{ e_{1},e_{2},e_{3},e_{4}\right\} \\
V_{2} & =\limfunc{linspan}\left\{ e_{5},e_{6},e_{3},e_{4}\right\} \\
V_{1}\oplus V_{2} & =\limfunc{linspan}\left\{ \left( e_{1},0\right) ,\left(
e_{2},0\right) ,\left( e_{3},0\right) ,\left( e_{4},0\right) ,\left(
0,e_{3}\right) ,\left( 0,e_{4}\right) ,\left( 0,e_{5}\right) ,\left(
0,e_{6}\right) \right\} \\
V_{1}+V_{2} & =\limfunc{linspan}\left\{
e_{1},e_{2},e_{3},e_{4},e_{5},e_{6}\right\} =V_{1}\cup V_{2}\cong\mathcal{V}%
\subset V_{1}\oplus V_{2}
\end{align}

and%
\begin{align}
V_{1} & =\limfunc{linspan}\left\{ e_{1},e_{2},e_{3},e_{4}\right\} \\
V_{2} & =\limfunc{linspan}\left\{ e_{5},e_{6},e_{7},e_{8}\right\} \\
V_{1}\oplus V_{2} & =\limfunc{linspan}\left\{ \left( e_{1},0\right) ,\left(
e_{2},0\right) ,\left( e_{3},0\right) ,\left( e_{4},0\right) ,\left(
0,e_{5}\right) ,\left( 0,e_{6}\right) ,\left( 0,e_{7}\right) ,\left(
0,e_{8}\right) \right\} \\
V_{1}+V_{2} & =\limfunc{linspan}\left\{
e_{1},e_{2},e_{3},e_{4},e_{5},e_{6},e_{7},e_{8}\right\} =V_{1}\cup
V_{2}\cong V_{1}\oplus V_{2}.
\end{align}

\subsubsection{Commutivity}

In what sense does 
\begin{equation}
V_{1}\oplus V_{2}\equiv\left\{ \left. \left( v,w\right) \right\vert v\in
V_{1},w\in V_{2}\right\} =V_{2}\oplus V_{1}\equiv\left\{ \left. \left(
w,v\right) \right\vert v\in V_{1},w\in V_{2}\right\} .
\end{equation}
Clearly 
\begin{equation}
V_{1}\oplus V_{2}\cong V_{2}\oplus V_{1}
\end{equation}
but they are \emph{not in general equal,} however if they are disjoint the
isomorphism is given by 
\begin{align}
\left( v,w\right) & \leftrightarrow v+w \\
\left( w,v\right) & \leftrightarrow v+w.
\end{align}

\subsection{Decomposition and Composition}

\subsubsection{Decomposition}

Given a vector space $V$ one can always find a \emph{complimentary}
subspace, $V_{1}^{c},$ to a subspace $V_{1}\subset V$ that is by
construction disjoint to $V_{1}.$

\begin{remark}
The subspace $V_{1}^{c}$ is not unique, if there is an inner product in $V$
the orthogonal compliment is unique and can be choosen amongst all
compliments as a canonical representative. This discussion is analogous to
the difference between orthogonal and non orthogonal projectors
\end{remark}

Then 
\begin{equation}
V=V_{1}\cup V_{1}^{c}=V_{1}+V_{1}^{c}\quad\cong\quad\mathcal{V}=V_{1}\oplus
V_{1}^{c}=\left\{ \left. \left( v_{1},v_{2}\right) \right\vert v_{1}\in
V_{1},v_{2}\in V_{1}^{c}\right\} =\widetilde{V_{1}}+\widetilde{V_{1}^{c}},
\end{equation}
where 
\begin{align}
V_{1} & \cong\widetilde{V_{1}}=\left\{ \left. \left( v_{1},0\right)
\right\vert v_{1}\in V_{1}\right\} \\
V_{1}^{c} & \cong\widetilde{V_{1}^{c}}=\left\{ \left. \left( 0,v_{2}\right)
\right\vert v_{2}\in V_{1}^{c}\right\} .
\end{align}
Taking into account the preceeding isomorphisms one says that $V$ is \emph{%
decomposed} into the direct sum of $V_{1}$ and $V_{1}^{c},$ and state that 
\begin{equation}
V=V_{1}\oplus V_{1}^{c},
\end{equation}
where we make the identification 
\begin{equation}
V\equiv\mathcal{V},
\end{equation}
although technically the direct sum of $V_{1}$ and $V_{1}^{c}$ is not
contained in $V.$ This decomposion can be paraphrased as any vector $v\in V$
can be decomposed into a sum of vectors each corresponding to a vector in
each subspace that occurs in the direct sum decomposition of $V.$

\subsubsection{Composition}

If one considers two vector spaces $V_{1}$ and $V_{2},$ (which could be
identical i.e. $V_{1}=$ $V_{2}),$ one can form their direct sum to \emph{%
compose }the vector space\emph{\ }$V$ 
\begin{equation}
V=V_{1}\oplus V_{2}=\left\{ \left. \left( v_{1},v_{2}\right) \right\vert
v\in V_{1},v_{2}\in V_{2}\right\} =\widetilde{V_{1}}+\widetilde{V_{2}}=%
\widetilde{V_{1}}+\widetilde{V_{1}^{c}}.
\end{equation}
$\widetilde{V_{1}}$ and $\widetilde{V_{2}}$ are not subspaces of $V,$ though
they are isomorphic to subspaces of $V$ i.e. 
\begin{align}
V_{1} & \cong\widetilde{V_{1}}=\left\{ \left. \left( v_{1},0\right)
\right\vert v_{1}\in V_{1}\right\} \\
V_{2} & \cong\widetilde{V_{2}}=\left\{ \left. \left( 0,v_{2}\right)
\right\vert v_{2}\in V_{2}\right\}
\end{align}
and 
\begin{equation}
\widetilde{V_{1}^{c}}=\widetilde{V_{2}}.
\end{equation}
The space $V$ \emph{composed} by the direct sum of $V_{1}$ and $V_{2}$ can
always be \emph{decomposed} as the direct sum 
\begin{equation}
V=V_{1}\oplus V_{2},
\end{equation}
which in general 
\begin{equation}
\ncong V_{1}+V_{2}
\end{equation}
\emph{unless} $V_{1}$ and $V_{2}$ are disjoint vector sub spaces of $V$.

\begin{remark}
For completeness we consider the \emph{diagonal} \emph{sum of two subspaces }%
given by\emph{\ }%
\begin{equation}
V_{1}+_{d}V_{2}=\left\{ \left. v_{1}+M\left( v_{1}\right) \right\vert
v_{1}\in V_{1}\right\} .
\end{equation}
\end{remark}

\begin{remark}
$1+1$
\end{remark}

\begin{align}
\lambda \left( \mathbf{z}_{1},\mathbf{z}_{2}\right) & \sim \left( \lambda 
\mathbf{z}_{1},\lambda \mathbf{z}_{2}\right) \\
\left( \mathbf{z}_{1},\mathbf{z}_{2}\right) +\left( \mathbf{z}_{3},\mathbf{z}%
_{4}\right) & \sim \left( \mathbf{z}_{1}+\mathbf{z}_{3},\mathbf{z}_{3}+%
\mathbf{z}_{4}\right) \\
\lambda _{1}\left( \mathbf{z}_{1},\mathbf{z}_{2}\right) +\lambda _{2}\left( 
\mathbf{z}_{3},\mathbf{z}_{4}\right) & \sim \left( \lambda _{1}\mathbf{z}%
_{1}+\lambda _{2}\mathbf{z}_{3},\lambda _{1}\mathbf{z}_{2}+\lambda _{2}%
\mathbf{z}_{4}\right)
\end{align}%
\begin{gather}
\left\langle \left. \mathbf{z}_{1}\right\vert \mathbf{z}_{2}\right\rangle =%
\overline{z}_{11}z_{21}+\overline{z}_{12}z_{22}\text{ } \\
\left\Vert \mathbf{z}_{1}\right\Vert ^{2}=\overline{z}_{11}z_{11}+\overline{z%
}_{12}z_{12}=\left\vert z_{11}\right\vert ^{2}+\left\vert z_{12}\right\vert
^{2} \\
\text{corresponding to } \\
\mathbb{C}^{2}=\mathbb{C\oplus C}
\end{gather}

\subsection{Group Representations}

Bounded linear group representation are homomorphism from the group to space
of bounded linear maps of a normed vector space i.e. 
\begin{equation}
\mathcal{R}_{V}:G\rightarrow\mathcal{B}\left( V\right) .
\end{equation}
Representations can be classified as reducible or irreducible. If they are
reducible then the vector space $V$ can be decomposed into a direct sum of $%
G $-invariant subspaces that produce irreducible representations of $G$ i.e. 
\begin{equation}
V=\tbigoplus _{1\leq p\leq d}V_{p},
\end{equation}
where the subspaces $\left\{ V_{p}\right\} $ can be the same or distinct
and\qquad\qquad\qquad\qquad\ 
\begin{equation}
\mathcal{R}_{V}\left( G\right) :V_{p}\rightarrow V_{p};1\leq p\leq d.
\end{equation}
Two representations $\mathcal{R}_{V}$ and $\mathcal{R}_{W}$ are equivalent
if 
\begin{equation}
T\mathcal{R}_{V}\left( G\right) T^{-1}=\mathcal{R}_{W}\left( G\right) ,
\end{equation}
where%
\begin{equation}
T:V\rightarrow W,
\end{equation}
i.e. the representations are intwined 
\begin{equation}
T\mathcal{R}_{V}\left( G\right) =\mathcal{R}_{W}\left( G\right) T.
\end{equation}

\subsubsection{Physical Interpretation}

If $G$ is the symmetry group of the "particles" under consideration then the
irreducible representations $\mathcal{R}_{Vp_{1}\cdots p_{\kappa }}$
contained in the one particle space $V$ describe the possible states of the
particle.

\subsubsection{Functions on a Group}

If $V$ is a Banach space then one construct complex values functions on $G$ 
\begin{equation}
f:G\rightarrow \mathbb{C}.
\end{equation}%
By using the representations one can define the functions as 
\begin{equation}
f_{lv}\left( g\right) =\left\langle \left. l\right\vert \mathcal{R}%
_{V}\left( g\right) v\right\rangle ;\,l\in V^{d},v\in V.
\end{equation}%
If the representation space and its dual have bases $\left\{ v_{k}\right\} $%
and $\left\{ l_{j}\right\} $ respectively then any function generated by
bounded linear representations can be expressed as combinations of 
\begin{equation}
f_{l_{j}v_{k}}\left( g\right) =\left\langle \left. l_{j}\right\vert \mathcal{%
R}_{V}\left( g\right) v_{k}\right\rangle .
\end{equation}%
The product of basis functions gives 
\begin{equation}
\sum_{1\leq k\leq r}f_{l_{j}v_{k}}\left( g_{1}\right) f_{l_{k}v_{m}}\left(
g_{2}\right) =f_{l_{j}v_{m}}\left( g_{1}g_{2}\right) .
\end{equation}%
In particular 
\begin{equation}
\sum_{1\leq k\leq r}f_{l_{j}v_{k}}\left( g^{-1}\right) f_{l_{k}v_{m}}\left(
g\right) =\sum_{1\leq k\leq r}f_{l_{j}v_{k}}\left( g\right)
f_{l_{k}v_{m}}\left( g^{-1}\right) =f_{l_{j}v_{m}}\left( id\right)
=\left\langle \left. l_{j}\right\vert I_{V}v_{m}\right\rangle =\delta _{jm}.
\end{equation}

\subsubsection{Left Regular Representation}

One can form a group vector space, $V\left( G\right) ,$ by considering
formal linear combinations 
\begin{equation}
\mathbf{\lambda}=\sum_{1\leq j\leq\left\vert G\right\vert }\lambda\left(
g_{j}\right) g_{j},\quad\mathbf{\lambda\in}V\left( G\right) .
\end{equation}
Then the action of $G$ on this vector space is given by 
\begin{equation}
L_{g}\sum_{1\leq j\leq\left\vert G\right\vert }\lambda\left( g_{j}\right)
g_{j}=\sum_{1\leq j\leq\left\vert G\right\vert }g\lambda\left( g_{j}\right)
g_{j}=\sum_{1\leq j\leq O\left( G\right) }\mu\left( g_{j}\right) g_{j},
\end{equation}
i.e. 
\begin{equation}
L_{g}\lambda\left( h\right) =\mu\left( h\right) .
\end{equation}
The left regular representation can also be defined by the action of the
group on the space of functions on the group as 
\begin{equation}
L_{g}f\left( h\right) =f\left( g^{-1}h\right) ,
\end{equation}
thus 
\begin{equation}
L_{g}\lambda\left( h\right) =\lambda\left( g^{-1}h\right) .
\end{equation}
This gives 
\begin{equation}
L_{g}\sum_{1\leq j\leq\left\vert G\right\vert }\lambda\left( g_{j}\right)
g_{j}=\sum_{1\leq j\leq\left\vert G\right\vert }\lambda\left(
g^{-1}g_{j}\right) g_{j}.
\end{equation}

\begin{remark}
In th definition one must use $g^{-1}$ otherwise $L$ is not a group
homomorphism as 
\begin{equation}
\left( \left( L_{g_{1}}L_{g_{2}}\right) \left( f\right) \right) \left(
h\right) =\left( L_{g_{1}}\left( L_{g_{2}}\left( f\right) \right) \right)
\left( h\right) =L_{g_{1}}\left( L_{g_{2}}\left( f\right) \right) \left(
h\right) =\left( L_{g_{2}}f\right) \left( g_{1}^{-1}h\right) =f\left(
g_{2}^{-1}g_{1}^{-1}h\right) =f\left( \left( g_{1}g_{2}\right) ^{-1}h\right)
=L_{g_{1}g_{2}}f\left( h\right) ,
\end{equation}%
thus 
\begin{equation}
\left( L_{g_{1}}L_{g_{2}}\right) \left( f\left( h\right) \right)
=L_{g_{1}g_{2}}\left( f\left( h\right) \right) \quad \forall f\in F\left( G,%
\mathbb{C}\right) ,\forall h\in G,\forall g_{1}g_{2}\in G
\end{equation}%
so 
\begin{equation}
L:G\rightarrow F\left( G,G\right)
\end{equation}%
is a group homomorphism as 
\begin{equation}
L\left( g_{1}\right) \circ L\left( g_{2}\right) =L\left( g_{1}g_{2}\right) ,
\end{equation}%
where 
\begin{equation}
L\left( g\right) \equiv L_{g}.
\end{equation}%
On the other hand 
\begin{equation}
\mathcal{L}_{g}f\left( h\right) =f\left( gh\right)
\end{equation}%
gives 
\begin{equation}
\left( \mathcal{L}_{g_{1}}\mathcal{L}_{g_{2}}\right) \left( f\right) \left(
h\right) =\mathcal{L}_{g_{1}}\left( \mathcal{L}_{g_{2}}\left( f\right)
\right) \left( h\right) =\left( \mathcal{L}_{g_{2}}\left( f\right) \right)
\left( g_{1}h\right) =f\left( g_{2}g_{1}h\right) =\left( \mathcal{L}%
_{g_{2}g_{1}}f\right) \left( h\right) ,
\end{equation}%
which shows that $\mathcal{L}$ is not a group homomorphism
\end{remark}

\subsubsection{Group Algebra}

\paragraph{Finite Dimensional Groups}

The vector space structure is given by 
\begin{equation}
\mathbf{\lambda}_{1}+\mathbf{\lambda}_{2}=\sum_{1\leq j\leq\left\vert
G\right\vert }\lambda_{1}\left( g_{j}\right) g_{j}+\sum_{1\leq
j\leq\left\vert G\right\vert }\lambda_{2}\left( g_{j}\right)
g_{j}=\sum_{1\leq j\leq\left\vert G\right\vert }\left( \lambda_{1}\left(
g_{j}\right) +\lambda_{2}\left( g_{j}\right) \right) g_{j}
\end{equation}
and the multiplication by 
\begin{align}
\mathbf{\lambda}_{1}\ast\mathbf{\lambda}_{2} & =\sum_{1\leq j\leq\left\vert
G\right\vert }\lambda_{1}\left( g_{j}\right) g_{j}\sum_{1\leq k\leq
\left\vert G\right\vert }\lambda_{2}\left( g_{k}\right) g_{k}=\sum_{1\leq
j,k\leq\left\vert G\right\vert }\lambda_{1}\left( g_{j}\right) \lambda
_{2}\left( g_{k}\right) g_{j}g_{k}=\sum_{1\leq j,k\leq\left\vert
G\right\vert }\lambda_{1}\left( g_{j}\right) L_{g_{j}}\lambda_{2}\left(
g_{k}\right) g_{k} \\
& =\sum_{1\leq j,k\leq\left\vert G\right\vert }\lambda\left( g_{j}\right)
\lambda_{2}\left( g_{j}^{-1}g_{k}\right) g_{k}=\sum_{1\leq k\leq\left\vert
G\right\vert }\lambda_{3}\left( g_{k}\right) g_{k}=\mathbf{\lambda}_{3}
\end{align}
where 
\begin{equation}
\lambda_{3}\left( g_{k}\right) =\sum_{1\leq j\leq\left\vert G\right\vert
}\lambda_{1}\left( g_{j}\right) \lambda_{2}\left( g_{j}^{-1}g_{k}\right) .
\end{equation}
which is called the \emph{convolution product }in the space $F\left(
G\right) $ i.e. 
\begin{equation}
\left( \lambda_{1}\ast\lambda_{2}\right) \left( g\right) =\sum_{\ h\in
G}\lambda_{1}\left( h\right) \lambda_{2}\left( h^{-1}g\right) .
\end{equation}
The vector space $V\left( G\right) $ equipped with the product "$\ast"$
becomes an algebra $A\left( G\right) .$

Linear combination in $F\left( G\right) $ is defined as 
\begin{equation}
\left( \lambda_{1}+\lambda_{2}\right) \left( g\right) =\lambda_{1}\left(
g\right) +\lambda_{2}\left( g\right) ,
\end{equation}
which shows that $F\left( G\right) $ equipped with the convolution product
is also an algebra that is in fact isomorhpic to $A\left( G\right) .$

\paragraph{Infinite Dimensional Groups}

In the infinite dimensional case a norms can be introducd in $V\left(
G\right) $ and $F\left( G\right) $ for \emph{locally compact groups} 
\begin{align}
\left\Vert \mathbf{\lambda}\right\Vert _{L_{1}} & =\left\Vert \mathbb{\lambda%
}\right\Vert _{L_{1}}=\int_{G}\left\vert \mathbb{\lambda }\left( g\right)
\right\vert d\mu\left( g\right) \\
\left\Vert \mathbf{\lambda}\right\Vert _{L_{2}} & =\left\Vert \mathbb{\lambda%
}\right\Vert _{L_{2}}=\int_{G}\left\vert \mathbb{\lambda }\left( g\right)
\right\vert ^{2}d\mu\left( g\right) ,
\end{align}
where $\mu$ is Haar measure for $G$ and one considers the completions of $%
F\left( G\right) $ in these norms to produce $L_{1}\left( G,\mu\right) $ and 
$L_{2}\left( G,\mu\right) .$

The convolution product in $L_{1}\left( G,\mu\right) $ is expressed as 
\begin{equation}
\left( \lambda_{1}\ast\lambda_{2}\right) \left( g\right)
=\int_{G}\lambda_{1}\left( h\right) \lambda_{2}\left( h^{-1}g\right)
d\mu\left( h\right)
\end{equation}
and it can be shown that $\left( \lambda_{1}\ast\lambda_{2}\right) \in
L_{1}\left( G,\mu\right) ,$ thus making $L_{1}\left( G,\mu\right) $ a Banach 
$^{\ast}$-algebra where the involution is defined as 
\begin{equation}
f^{\ast}\left( g\right) =\overline{f\left( g^{-1}\right) }\Delta\left(
g^{-1}\right)
\end{equation}
and $\Delta$ is the modular function for the group $G.$

\subsection{Vector Tensor Product}

One can endower the Cartesian product of two vector spaces 
\begin{equation}
V_{1}\times V_{2}=\left\{ \left. \left( v_{1},v_{2}\right) \right\vert
v_{1}\in V_{1},v_{2}\in V_{2}\right\} ,
\end{equation}
with a tensor product structure by generating the free vector space based on 
$V_{1}\times V_{2}$ and then assigning the equivalences 
\begin{align}
\left( v_{1},v_{3}\right) +\left( v_{2},v_{3}\right) & \sim\left(
v_{1}+v_{2},v_{3}\right) \\
\left( v_{1},v_{2}\right) +\left( v_{1},v_{3}\right) & \sim\left(
v_{1},v_{2}+v_{3}\right) \\
\lambda\left( v_{1},v_{2}\right) & \sim\left( \lambda v_{1},v_{2}\right)
\sim\left( v_{1},\lambda v_{2}\right)
\end{align}
to produce $V_{1}\otimes_{\mathbb{F}}V_{2},$ where $\lambda\in\mathbb{F=}$
usually $\mathbb{R}$ or $\mathbb{C}.$

\begin{example}
Simple example:
\end{example}

The cartesian direct product is 
\begin{equation}
\mathbb{R}^{2}=\mathbb{R\times R}=\left\{ \left. \left( x,y\right)
\right\vert x,y\in \mathbb{R}\right\}
\end{equation}%
The direct sum imposes a vector space structure 
\begin{equation}
\mathbb{R}^{2}=\mathbb{R\oplus R}
\end{equation}%
given by 
\begin{align}
a\left( x,y\right) & \sim \left( ax,ay\right) \\
\left( x_{1},y_{1}\right) +\left( x_{2},y_{2}\right) & \sim \left(
x_{1}+x_{2},y_{1}+y_{2}\right) \\
a\left( x_{1},y_{1}\right) +b\left( x_{2},y_{2}\right) & \sim \left(
ax_{1}+bx_{2},ay_{1}+by_{2}\right) .
\end{align}%
The tensor product $\mathbb{R\otimes R}$%
\begin{equation}
\mathbb{R\otimes R\neq R}^{2}
\end{equation}%
as it is defined by 
\begin{align}
\left( x,y_{1}\right) +\left( x,y_{2}\right) & \sim \left(
x,y_{1}+y_{2}\right) \\
\left( x_{1},y\right) +\left( x_{2},y\right) & \sim \left(
x_{1}+x_{2},y\right) \\
a\left( x,y\right) & \sim \left( ax,y\right) \sim \left( x,ay\right) .
\end{align}

\begin{remark}
$1\times 1$%
\begin{align}
\left( \mathbf{z}_{1},\mathbf{z}_{3}\right) +\left( \mathbf{z}_{2},\mathbf{z}%
_{3}\right) & \sim \left( \mathbf{z}_{1}+\mathbf{z}_{2},\mathbf{z}_{3}\right)
\\
\left( \mathbf{z}_{1},\mathbf{z}_{2}\right) +\left( \mathbf{z}_{1},\mathbf{z}%
_{3}\right) & \sim \left( \mathbf{z}_{1},\mathbf{z}_{2}+\mathbf{z}_{3}\right)
\\
\lambda \left( \mathbf{z}_{1},\mathbf{z}_{2}\right) & \sim \left( \lambda 
\mathbf{z}_{1},\mathbf{z}_{2}\right) \sim \left( \mathbf{z}_{1},\lambda 
\mathbf{z}_{2}\right)
\end{align}%
\begin{gather}
\left\langle \left. \mathbf{z}_{1}\right\vert \mathbf{z}_{2}\right\rangle
=\left( \overline{z}_{11}z_{21}\right) \left( \overline{z}_{12}z_{22}\right)
\\
\left\Vert \mathbf{z}_{1}\right\Vert ^{2}=\text{ }\left\vert
z_{11}\right\vert ^{2}\left\vert z_{12}\right\vert ^{2} \\
\text{corresponding to }\mathbb{C\otimes C}.
\end{gather}
\end{remark}

\subsubsection{Matric Tensor Product}

The set of matrices $M_{n}\left( \mathbb{C}\right) $ forms an algebra i.e. 
\emph{has both vector space and ring properties}. There is a natural
definition of tensor product 
\begin{equation}
A\otimes B=\left( 
\begin{array}{ccc}
a_{11} & \cdots & a_{1n} \\ 
\vdots & \ddots & \vdots \\ 
a_{n1} & \cdots & a_{nn}%
\end{array}%
\right) \otimes \left( 
\begin{array}{ccc}
b_{11} & \cdots & b_{1n} \\ 
\vdots & \ddots & \vdots \\ 
b_{n1} & \cdots & b_{nn}%
\end{array}%
\right) =\left( 
\begin{array}{ccc}
a_{11}B & \cdots & a_{1n}B \\ 
\vdots & \ddots & \vdots \\ 
a_{n1}B & \cdots & a_{nn}B%
\end{array}%
\right) ,
\end{equation}%
which shows that 
\begin{equation}
M_{m}\left( \mathbb{C}\right) \otimes M_{n}\left( \mathbb{C}\right)
=M_{m}\left( M_{n}\left( \mathbb{C}\right) \right) ,
\end{equation}
i.e. a general element of $M_{m}\left( \mathbb{C}\right) \otimes M_{n}\left( 
\mathbb{C}\right) $ can be expanded as 
\begin{equation}
\sum_{1\leq j\leq n^{2}m^{2}}\left( 
\begin{array}{ccc}
a_{j11}B_{j} & \cdots & a_{j1n}B_{j} \\ 
\vdots & \ddots & \vdots \\ 
a_{jn1}B_{j} & \cdots & a_{jnn}B_{j}%
\end{array}%
\right) =\left( 
\begin{array}{ccc}
B_{11} & \cdots & B_{1n} \\ 
\vdots & \ddots & \vdots \\ 
B_{n1} & \cdots & B_{nn}%
\end{array}%
\right) .
\end{equation}

One can form the tensor algebra 
\begin{equation}
\bigotimes\left( M_{n}\left( \mathbb{C}\right) \right) =\tbigoplus
\limits_{0\leq n\leq\infty}\tbigotimes \nolimits^{n}\left( M_{n}\left( 
\mathbb{C}\right) \right) .
\end{equation}

\subsection{Complexification \label{3/8/2026}}

\subsubsection{External Complexification}

The formal definition of the complexification $V_{\mathbb{C}}$ of the real
space $V$ is in terms of the real tensor product 
\begin{equation}
V_{\mathbb{C}}=V\otimes_{\mathbb{R}}\mathbb{C}\text{ or }\mathbb{C}\otimes_{%
\mathbb{R}}V.
\end{equation}
One should note that 
\begin{equation}
V\otimes_{\mathbb{R}}\mathbb{C\cong C}\otimes_{\mathbb{R}}V
\end{equation}
but 
\begin{equation}
V\otimes_{\mathbb{R}}\mathbb{C\neq C}\otimes_{\mathbb{R}}V.
\end{equation}
In the proceeding we set 
\begin{equation}
V_{\mathbb{C}}=\text{ }\mathbb{C}\otimes_{\mathbb{R}}V.
\end{equation}
$\mathbb{C}$ can be considered as a $2$-dimensional abelian $c^{\ast}$-real
algebra, (thus a vector space) with a basis $\left\{ 1,i\right\} ,$ i.e. 
\begin{equation}
\mathbb{C}=\left\{ \left. z=a1+bi=1a+ib\right\vert a,b\in\mathbb{R}\right\}
\end{equation}
and \emph{complex multiplication} in $V_{\mathbb{C}}$ is defined by%
\begin{equation}
\lambda\left( z\otimes_{\mathbb{R}}v\right) =\left( \lambda z\right)
\otimes_{\mathbb{R}}v;\,v\in V;\lambda,z\in\mathbb{C}.
\end{equation}
The tensor product can be decomposed as a direct sum 
\begin{align}
z\otimes_{\mathbb{R}}v & =\left( \func{Re}\left( z\right) +i\func{Im}\left(
z\right) \right) \otimes_{\mathbb{R}}v=\func{Re}\left( z\right) \otimes_{%
\mathbb{R}}v+i\func{Im}\left( z\right) \otimes_{\mathbb{R}}v \\
& =1\otimes_{\mathbb{R}}\func{Re}\left( z\right) v+i\otimes _{\mathbb{R}}%
\func{Im}\left( z\right) v
\end{align}
and multiplication by complex \emph{numbers} can be expressed in terms of
the direct sum as 
\begin{align}
\left( \lambda z\right) \otimes_{\mathbb{R}}v & =\left( \func{Re}\left(
\lambda z\right) +i\func{Im}\left( \lambda z\right) \right) \otimes_{\mathbb{%
R}}v=\func{Re}\left( \lambda z\right) \otimes_{\mathbb{R}}v+i\func{Im}\left(
\lambda z\right) \otimes_{\mathbb{R}}v \\
& =1\otimes_{\mathbb{R}}\func{Re}\left( \lambda z\right) v+i\otimes_{\mathbb{%
R}}\func{Im}\left( \lambda z\right) v.
\end{align}
In particular if $z=\alpha\in\mathbb{R}$ 
\begin{equation}
\alpha\otimes_{\mathbb{R}}v=\alpha1\otimes_{\mathbb{R}}v+0i\otimes _{\mathbb{%
R}}v=1\otimes_{\mathbb{R}}\alpha v=\left( \alpha v,0\right)
\end{equation}
and if $z=\alpha i,\alpha\in\mathbb{R}$%
\begin{equation}
\left( i\alpha\right) \otimes_{\mathbb{R}}v=i\alpha\otimes_{\mathbb{R}%
}v=\left( 0,i\alpha v\right) .
\end{equation}
Multiplication by the \emph{operator} $"i"$ is given by its action on the
real parts as 
\begin{equation}
i\left( \alpha\otimes_{\mathbb{R}}v\right) =i\left( \alpha v,0\right)
=\left( i\alpha\right) 1\otimes_{\mathbb{R}}v=0\left( 1\otimes_{\mathbb{R}%
}v\right) +\left( i\alpha\right) \otimes_{\mathbb{R}}v=\left( 0,i\alpha
v\right)
\end{equation}
and imaginary parts as 
\begin{equation}
i\left( \left( i\alpha\right) \otimes_{\mathbb{R}}v\right) =-\alpha
1\otimes_{\mathbb{R}}v+0i\otimes_{\mathbb{R}}v=\left( -\alpha v,0\right)
\end{equation}
of the direct sum.

It is possible to find other bases of $\mathbb{C}$ that are real linear
combinations of $\left\{ 1,i\right\} ,$ but it is conventional and practical
to use the fixed canonical basis $\left\{ 1,i\right\} .$

$\lceil$One could choose a different basis for $\mathbb{C}$ i.e. $\left\{
1+i,1-i\right\} $ so that 
\begin{equation}
\left( 1,i\right) \left( 
\begin{array}{cc}
1 & 1 \\ 
1 & -1%
\end{array}
\right) =\left( 1+i,1-i\right)
\end{equation}
and 
\begin{equation}
w=v_{3}\otimes\left( 1+i\right) +v_{4}\otimes\left( 1-i\right)
\end{equation}%
\begin{align}
v_{1} & =v_{3}+v_{4} \\
v_{2} & =v_{3}-v_{4}.
\end{align}

$\rfloor$

$\lceil z\in \mathbb{C}$ has real matric representations as two dimensional
abelian subalgebras of $\mathcal{B}\left( \mathbb{R}^{2r}\right) ,r\in 
\mathbb{N}$ 
\begin{align}
\mathfrak{R}_{r}\left( z\right) & =\left\{ \left. \left( 
\begin{array}{cc}
a\mathbb{I}_{r} & -b\mathbb{I}_{r} \\ 
b\mathbb{I}_{r} & a\mathbb{I}_{r}%
\end{array}%
\right) \right\vert a,b\in \mathbb{R}\right\} \\
\mathfrak{R}_{r}\left( 1\right) & =\left( 
\begin{array}{cc}
\mathbb{I}_{r} & 0 \\ 
0 & \mathbb{I}_{r}%
\end{array}%
\right) =\mathbb{I}_{2r} \\
\mathfrak{R}_{r}\left( i\right) & =\left( 
\begin{array}{cc}
0 & -\mathbb{I}_{r} \\ 
\mathbb{I}_{r} & 0%
\end{array}%
\right) =\mathbb{J}_{2r},
\end{align}%
giving%
\begin{align}
V_{\mathbb{C}}& =\mathfrak{R}_{r}\left( \mathbb{C}\right) \otimes _{\mathbb{R%
}}V=\mathfrak{R}_{r}\left( 1\right) \otimes _{\mathbb{R}}V+\mathfrak{R}%
_{r}\left( i\right) \otimes _{\mathbb{R}}V \\
& =I_{2r}\otimes _{\mathbb{R}}V+J_{2r}\otimes _{\mathbb{R}}V.
\end{align}

More abstractly $\mathbb{C}$ has operator representations as real two
dimensional abelian subalgebras of $\mathcal{L}\left( W_{1}\oplus
W_{2}\right) ,$ where $\dim\left\{ W_{1}\right\} =\dim\left\{ W_{2}\right\} $
and 
\begin{align}
J_{2r} & :W_{1}\rightarrow W_{2} \\
J_{2r} & :W_{2}\rightarrow W_{1}.
\end{align}
The vector spaces $\left\{ W_{1},W_{2}\right\} $ could both be choosen to be 
$V.$

$\rfloor$

\paragraph{Tensor product to Direct Sum}

From the properties of tensor product every vector $v_{c}$ in $V_{\mathbb{C}%
} $ can be expressed via the direct sum and the isomorphic sum%
\begin{equation}
V_{\mathbb{C}}=\mathbb{C}\otimes_{\mathbb{R}}V=1\otimes V+i\otimes
V\cong1\otimes V\oplus i\otimes V=1V\oplus iV\cong1V+iV
\end{equation}
i.e. 
\begin{align*}
V_{\mathbb{C}} & \ni v_{c}=1\otimes v_{1}+i\otimes v_{2}=\left( 1\otimes
v_{1},0\right) +\left( 0,i\otimes v_{2}\right) =\left( 1v_{1},0\right)
+\left( 0,iv_{2}\right) =\left( v_{1},iv_{2}\right) \\
& =1\otimes v_{1}+i\otimes v_{2}=1v_{1}+iv_{2},
\end{align*}
which is the results of $1$ and $i$ being scalars and the tensor product
being the standard pointwise product of the complex numbers i.e. 
\begin{align}
1\otimes v & =1v=v \\
i\otimes v & =iv \\
v & \in V.
\end{align}
$\lceil$Example when $r=4$%
\begin{align}
V_{\mathbb{C}} & =\mathbb{C}\otimes_{\mathbb{R}}V=\left\{ 1,i\right\}
\otimes_{\mathbb{R}}\left\{ e_{1},e_{2},e_{3},e_{4}\right\} =\left\{
1\otimes e_{1},1\otimes e_{2},1\otimes e_{3},1\otimes e_{4},i\otimes
e_{1},i\otimes e_{2},i\otimes e_{3},i\otimes e_{4}\right\} \\
& =\left\{ e_{1},e_{2},e_{3},e_{4},ie_{1},ie_{2},ie_{3},ie_{4}\right\}
\end{align}
giving%
\begin{align}
v_{c} & =a_{1}e_{1}+a_{2}e_{2}+a_{3}e_{3}+a_{4}e_{4}\in\mathbb{C}\otimes_{%
\mathbb{R}}V \\
& =\left( \func{Re}\left( a_{1}\right) +i\func{Im}\left( a_{1}\right)
\right) e_{1}+\left( \func{Re}\left( a_{2}\right) +i\func{Im}\left(
a_{2}\right) \right) e_{2}+\left( \func{Re}\left( a_{3}\right) +i\func{Im}%
\left( a_{3}\right) \right) e_{3}+\left( \func{Re}\left( a_{4}\right) +i%
\func{Im}\left( a_{4}\right) \right) e_{4} \\
& =\func{Re}\left( a_{1}\right) e_{1}+\func{Re}\left( a_{2}\right) e_{2}+%
\func{Re}\left( a_{3}\right) e_{3}+\func{Re}\left( a_{4}\right) e_{4}+i\func{%
Im}\left( a_{1}\right) e_{1}+i\func{Im}\left( a_{2}\right) e_{2}+i\func{Im}%
\left( a_{3}\right) e_{3}+i\func{Im}\left( a_{4}\right) e_{4} \\
& =\func{Re}\left( v_{c}\right) +i\func{Im}\left( v_{c}\right) \in V_{%
\mathbb{C}}\cong V\oplus iV=\left\{ \left( \func{Re}\left( v_{c}\right) ,i%
\func{Im}\left( v_{c}\right) \right) \right\} .
\end{align}
$\rfloor$

\emph{The real spaces }$\left\{ 1\otimes_{\mathbb{R}}V,i\otimes_{\mathbb{R}%
}V\right\} =$\emph{\ }$\left\{ V,iV\right\} $\emph{\ are disjoint (wrt real
linear combinations) thus} 
\begin{equation}
V\oplus iV=\widetilde{1V}+\widetilde{iV}\cong V+iV.
\end{equation}
To summerize the tensor product can be expanded as 
\begin{align}
V_{\mathbb{C}} & =\mathbb{C}\otimes_{\mathbb{R}}V=\left\{ \left. \left(
z,v\right) \right\vert z\in\mathbb{C},v\in V\right\} =\left\{ \func{Re}%
\left( z\right) 1\otimes_{\mathbb{R}}V\oplus \func{Im}\left( z\right)
i\otimes_{\mathbb{R}}V\right\} \\
& =\left\{ \left( \func{Re}\left( z\right) 1\otimes v,\func{Im}\left(
z\right) i\otimes v\right) \right\} =\left\{ \left( \func{Re}\left( z\right)
1\otimes v+\func{Im}\left( z\right) i\otimes v\right) \leftrightarrow\left( 
\func{Re}\left( z\right) v+\func{Im}\left( z\right) iv\right) \right\} \\
& =\left\{ \func{Re}\left( zv\right) +i\func{Im}\left( zv\right) \right\}
\leftrightarrow\left\{ \left( \func{Re}\left( zv\right) ,i\func{Im}\left(
zv\right) \right) \right\} ;z\in\mathbb{C},v\in V,
\end{align}
i.e. 
\begin{align}
V_{\mathbb{C}} & =\mathbb{C}\otimes_{\mathbb{R}}V=V\oplus iV\equiv\left\{
\left( z,v\right) =z\otimes v\right\} =\left\{ \left( \func{Re}\left(
z\right) v,\func{Im}\left( z\right) iv\right) \right\} \\
& =\left\{ \left( \func{Re}\left( v\right) ,i\func{Im}\left( v\right)
\right) \right\} \leftrightarrow\left\{ \left( \func{Re}\left( v\right) +%
\func{Im}\left( v\right) i\right) \right\} =\left\{ \left( \func{Re}\left(
v\right) +i\func{Im}\left( v\right) \right) \right\} ;z\in\mathbb{C},v\in V,
\end{align}
where we haved used 
\begin{equation}
\left\{ \left( \func{Re}\left( v\right) +i\func{Im}\left( v\right) \right)
\right\} =\left\{ \left( \func{Re}\left( zv\right) +i\func{Im}\left(
zv\right) \right) \right\} ;v\in V;z\in\mathbb{C}.
\end{equation}
One can observe that 
\begin{align}
V_{\mathbb{C}} & =\mathbb{C-}\limfunc{linspan}\left\{ \left. \func{Re}\left(
ze_{j}\right) +i\func{Im}\left( ze_{j}\right) \right\vert 1\leq j\leq
r\right\} =\mathbb{C-}\limfunc{linspan}\left\{ \left. \func{Re}\left(
z\right) e_{j}+i\func{Im}\left( z\right) e_{j}\right\vert 1\leq j\leq
r\right\} \\
& =\mathbb{C-}\limfunc{linspan}\left\{ \left. \alpha e_{j}+\beta
ie_{j}\right\vert 1\leq j\leq r\right\} =\mathbb{C-}\limfunc{linspan}\left\{
\left. \left( \alpha+i\beta\right) e_{j}\right\vert 1\leq j\leq r\right\} \\
& =\mathbb{R-}\limfunc{linspan}\left\{ \left. \alpha e_{j},\beta
ie_{j}\right\vert 1\leq j\leq r\right\} ;\alpha,\beta\in\mathbb{R} \\
& \mathbb{\leftrightarrow}\left\{ \left. \left( \left(
\alpha_{1},\beta_{1}\right) ,\cdots,\left( \alpha_{r},\beta_{r}\right)
\right) \right\vert \alpha_{j},\beta_{j}\in\mathbb{R}^{r};1\leq j\leq
r\right\} \leftrightarrow\left\{ \left. \left( z_{1},\cdots,z_{r}\right)
\right\vert z_{j},1\leq j\leq r\right\} .
\end{align}
The last isomorphisms are of course the same as produced by faithful
representations produced by a basis, i.e 
\begin{equation}
V_{\mathbb{C}}\cong\mathbb{R}^{2r}\cong\mathbb{C}^{r}.
\end{equation}

Complex multiplication is given by 
\begin{equation}
z_{1}\otimes _{\mathbb{R}}V=\left( x_{1}+iy_{1}\right) \otimes _{\mathbb{R}%
}V=\left\{ x_{1}1\otimes _{\mathbb{R}}V\oplus iy_{1}\otimes _{\mathbb{R}%
}V\right\} =\left\{ \left( x_{1}v_{1},iy_{1}v_{2}\right) \right\}
\end{equation}%
\begin{equation}
x\otimes _{\mathbb{R}}V=x\otimes _{\mathbb{R}}V=x1\otimes _{\mathbb{R}%
}V=\left( xv,0\right)
\end{equation}%
\begin{equation}
iy\otimes _{\mathbb{R}}V=iy\otimes _{\mathbb{R}}V=iy\otimes _{\mathbb{R}%
}V=\left( 0,iyv\right)
\end{equation}%
\begin{align}
z_{1}z_{2}\otimes _{\mathbb{R}}V& =\left( x_{1}+iy_{1}\right) \left(
x_{2}+iy_{2}\right) \otimes _{\mathbb{R}}V=\left( \left(
x_{1}x_{2}-y_{1}y_{2}\right) +i\left( x_{1}y_{2}+y_{1}x_{2}\right) \right)
\otimes _{\mathbb{R}}V \\
& =\left\{ \left( x_{1}x_{2}-y_{1}y_{2}\right) 1\otimes _{\mathbb{R}}V\oplus
i\left( x_{1}y_{2}+y_{1}x_{2}\right) \otimes _{\mathbb{R}}V\right\} \\
& =\left( \left( x_{1}x_{2}-y_{1}y_{2}\right) v_{1},\left(
x_{1}y_{2}+y_{1}x_{2}\right) iv_{2}\right)
\end{align}%
In particular 
\begin{align}
x\left( v_{1},iv_{2}\right) & =\left( xv_{1},xiv_{2}\right) \\
iy\left( v_{1},iv_{2}\right) & =\left( -yv_{2},iyv_{1}\right) \\
v_{1},v_{2}& \in V,x,y\in \mathbb{R}.
\end{align}%
\begin{align}
z\left( v_{1},iv_{2}\right) & =(x+iy)\left( v_{1},iv_{2}\right) =\left(
xv_{1}-yv_{2},iyv_{1}+xiv_{2}\right) \\
v_{1},v_{2}& \in V.
\end{align}

$\lceil$%
\begin{gather}
\mathcal{V}_{1}=\left\{ \left. \left( v,0\right) \right\vert v\in V\right\}
\cong V \\
i\mathcal{V}_{2}=\left\{ \left. \left( 0,iv\right) \right\vert v\in
V\right\} \cong iV \\
\left\{ \left. \left( v,0\right) \right\vert v\in V\right\} \cap\left\{
\left. \left( 0,iv\right) \right\vert v\in V\right\} =\left\{ \left(
0,0\right) \right\} \\
V\oplus iV=\mathcal{V}_{1}+i\mathcal{V}_{2}\cong V+iV.
\end{gather}
$\rfloor$

\paragraph{i (imaginary multiplication) induces a Complexification J$_{V}$
map in$\mathcal{\ }$ $\mathcal{V}_{\mathbb{R}}=V\oplus_{\mathbb{R}}V=%
\mathcal{V}_{1}+\mathcal{V}_{2}$}

Multiplication by the imaginary unit in $V_{\mathbb{C}}$ has the properties 
\begin{align}
i\left( v,0\right) & =\left( 0,iv\right) \\
i\left( 0,iv\right) & =\left( -v,0\right) \\
v & \in V.
\end{align}
This induces a real linear map (operator) $\mathcal{J}_{\mathcal{V}}\in%
\mathcal{L}\left( V\oplus_{R}iV\right) $ 
\begin{align}
\mathcal{J}_{\mathcal{V}} & :V\oplus_{R}iV\rightarrow V\oplus_{R}iV \\
\mathcal{J}_{\mathcal{V}} & :\widetilde{V_{1}}\equiv\mathcal{V}%
_{1}\rightarrow\widetilde{V_{2}}=i\mathcal{V}_{2} \\
\mathcal{J}_{\mathcal{V}} & :\widetilde{V_{2}}\rightarrow\widetilde{V_{1}},
\end{align}
given by 
\begin{align}
\mathcal{J}_{\mathcal{V}}\left( v,0\right) & =\left( 0,iv\right) =i\left(
v,0\right) \\
\mathcal{J}_{\mathcal{V}}\left( 0,iv\right) & =\left( -v,0\right) =i\left(
0,iv\right) \\
v & \in V.
\end{align}
This map then induces a linear map, which we name $J_{V}$ 
\begin{equation}
J_{V}:\mathcal{V}_{\mathbb{R}}\rightarrow\mathcal{V}_{\mathbb{R}},
\end{equation}
in the real space $\mathcal{V}_{\mathbb{R}}$ 
\begin{equation}
\mathcal{V}_{\mathbb{R}}=V\oplus_{\mathbb{R}}V=\mathcal{V}_{1}+\mathcal{V}%
_{2}\ncong V+V,
\end{equation}
where 
\begin{align}
\mathcal{V}_{1} & =\left\{ \left. \left( v,0\right) \right\vert v\in
V\right\} \\
\mathcal{V}_{2} & =\left\{ \left. \left( 0,v\right) \right\vert v\in
V\right\} ,
\end{align}
given by%
\begin{align}
J_{V}\left( v,0\right) & =\left( 0,v\right) \\
J_{V}\left( 0,v\right) & =\left( -v,0\right) \\
v,w & \in V.
\end{align}
\emph{Note that }$J_{V}$\emph{\ is not defined on }$V,$ but 
\begin{equation}
\mathcal{V}_{1}\cong\mathcal{V}_{2}\cong\widetilde{V_{1}}\cong\widetilde{%
V_{2}}\cong V.
\end{equation}

Reversing the order \emph{we can now define} multiplication by an imaginary
unit in $\mathcal{V}_{\mathbb{R}}$, to produce $\mathcal{V}_{\mathbb{C}}$ by 
\begin{align}
i\left( v_{1},v_{2}\right) & =J_{V}\left( v_{1},v_{2}\right) \\
i\left( v_{1},0\right) & =J_{V}\left( v_{1},0\right) =\left( 0,v_{1}\right)
\\
i\left( 0,v_{2}\right) & =J_{V}\left( 0,v_{2}\right) =\left( -v_{2},0\right)
\\
v_{1},v_{2} & \in V.
\end{align}
This gives 
\begin{equation}
\widetilde{V_{2}}=J_{V}\widetilde{V_{1}}=J_{V}\widetilde{V_{1}}=\widetilde{iV%
}_{1},
\end{equation}
so that 
\begin{equation}
\mathcal{V}_{1}+\mathcal{V}_{2}=\mathcal{V}_{1}+J_{V}\mathcal{V}_{1}\cong
V\oplus iV=V_{\mathbb{C}}.
\end{equation}

\paragraph{Complex and Real Linear Independence}

The real and complex linear spans of $V\oplus_{R}iV$ are equal i.e. 
\begin{equation}
V_{\mathbb{C}}=\mathbb{C}-\limfunc{linspan}\left\{ \left. \left( 1\otimes
v_{1},i\otimes v_{2}\right) \right\vert v_{1},v_{2}\in V\right\} =\mathbb{R}-%
\limfunc{linspan}\left\{ \left. \left( 1\otimes v_{1},i\otimes v_{2}\right)
\right\vert v_{1},v_{2}\in V\right\} .
\end{equation}
This can be seen by considering the following over complete basis for $V_{%
\mathbb{C}}$ $\left\{ \left( 1\otimes e_{1}\right) ,\cdots,\left( 1\otimes
e_{r}\right) ;\left( i\otimes e_{1}\right) ,\cdots,\left( i\otimes
e_{r}\right) \right\} $ and 
\begin{equation}
v=\sum_{1\leq j\leq r}\lambda_{j}\left( 1\otimes e_{j}\right) +\sum_{1\leq
j\leq r}\lambda_{j+r}\left( i\otimes e_{j}\right) ;\lambda_{j}\in \mathbb{R}%
,1\leq j\leq r.
\end{equation}
The vectors 
\begin{equation}
\left( 1\otimes v,0\right) \equiv\left( v,0\right) \text{ and }\left(
0,i\otimes v\right) \equiv\left( 0,iv\right)
\end{equation}
are $\emph{not}$ \emph{complex linearly independent }i.e. 
\begin{equation}
\lambda_{1}\left( 1\otimes v,0\right) \text{ }+\lambda_{2}\text{ }\left(
0,i\otimes v\right) =\left( 0,0\right) ;\lambda_{1},\lambda_{2}\in\mathbb{C}
\end{equation}
has non trivial solutions i.e. 
\begin{align}
& \left( a+ib\right) \left( 1\otimes v,0\right) \text{ }+\left( c+id\right) 
\text{ }\left( 0,i\otimes v\right) \\
& =a\left( 1\otimes v,0\right) +ib\left( 1\otimes v,0\right) +c\left(
0,i\otimes v\right) +id\left( 0,i\otimes v\right) \\
& =\left( a1\otimes v,0\right) +\left( 0,ib\otimes v\right) +\left(
0,ic\otimes v\right) -\left( d1\otimes v,0\right) \\
& =\left( \left( a-d\right) 1\otimes v,0\right) +\left( 0,i\left( b+c\right)
\otimes v\right) =\left( 0,0\right) \\
& \Rightarrow a=d\,\&\,b=-c.
\end{align}
They are however \emph{real linearly dependent.} Any $v\in V_{\mathbb{C}}$
can be expanded in terms of the complex linearly independent basis $\left\{
\left. e_{j}\right\vert 1\leq j\leq r\right\} $ as 
\begin{equation}
v=\sum_{1\leq j\leq r}\left( \lambda_{j}+i\lambda_{j+s}\right) \otimes
e_{j}=\sum_{1\leq j\leq r}\left( \lambda_{j}+i\lambda_{j+s},e_{j}\right)
=\sum_{1\leq j\leq r}\left( z_{j},e_{j}\right) ,
\end{equation}
where 
\begin{equation}
z_{j}\mathbb{\otimes}e_{j}=\left( z_{j},e_{j}\right) \in\mathbb{C\otimes}V.
\end{equation}
The vectors 
\begin{equation}
\left( 1\otimes v,0\right) \text{ and }\left( 0,i\otimes v\right)
\end{equation}
are \emph{real linearly independent. }

\subsubsection{Internal Complexification}

In the preceeding subsection we saw \emph{external} \emph{complexification,}
(complex numbers are added from the outside), of a real $r$-dimensional real
space $V$ given by the tensor product $\mathbb{C}\otimes_{\mathbb{R}}V=$ $%
V\otimes_{\mathbb{R}}iV$ $=\widetilde{V}+i\widetilde{V}$ $=\mathcal{V}_{%
\mathbb{C}}\cong$ $V+J_{V}V$ $=V_{\mathbb{C}}$ could be expressed in terms
of a complexification map $\mathcal{J}_{V}$ acting in the real space $%
\mathcal{V}_{\mathbb{R}}=V\oplus_{\mathbb{R}}V$ that induces the map 
\begin{equation}
J_{V}:V\rightarrow J_{V}V
\end{equation}
to give 
\begin{equation}
\mathbb{C}\otimes_{\mathbb{R}}V=V\oplus_{\mathbb{R}}J_{V}V=\widetilde{V}+%
\widetilde{J_{V}V}.
\end{equation}
This could alternativetly be described as an example of \emph{internal
complexification, }(nothing added\emph{\ }from the outside),\emph{\ }of the
real direct sum space $\mathcal{V}_{\mathbb{R}}$ 
\begin{equation}
\mathcal{V}_{\mathbb{R}}=V\oplus_{\mathbb{R}}V=\mathcal{V}_{1}+\mathcal{V}%
_{2}
\end{equation}
of dimension $2r.$

$\lceil$Notation 
\begin{equation}
\mathcal{V}_{1}=\widetilde{V_{1}}
\end{equation}

$\rfloor$

The essential properties needed to produce a complex space are not dependent
on $\mathcal{V}_{\mathbb{R}}$ being the direct sum of identical spaces but
are also possesed by spaces $\mathcal{V}_{\mathbb{R}}$ that have \emph{%
decompositions} into disjoint direct sums 
\begin{equation}
\mathcal{V}_{\mathbb{R}}=V\oplus_{\mathbb{R}}V^{c}=\widetilde{V}+\widetilde{V%
}^{c}=\widetilde{V}+\widetilde{V^{c}},
\end{equation}
where the compliment of $\widetilde{V}$ is taken in $\mathcal{V}_{\mathbb{R}%
} $ i.e. 
\begin{align}
\mathcal{V}_{1} & =\widetilde{V}=\left\{ \left. \left( v,0\right)
\right\vert v\in V\right\} \\
\mathcal{V}_{2} & =\widetilde{V^{c}}=\left\{ \left. \left( 0,v^{c}\right)
\right\vert v^{c}\in V^{c}\right\}
\end{align}
and one has used the isomorphisms between real spaces $V\cong\widetilde{V}$,$%
V^{c}\cong\widetilde{V^{c}}$ and $V+V^{c}\cong\widetilde{V}+\widetilde{V^{c}}
$ and one utilizes a complexification map $\mathcal{J}_{V}.$

The map $J_{V}\in\mathcal{L}\left( \widetilde{V},\widetilde{V^{c}}\right) $
and $\mathcal{L}\left( \widetilde{V^{c}},\widetilde{V}\right) $ is defined
by 
\begin{align}
\mathcal{J}_{V} & :V\oplus_{\mathbb{R}}V^{c}\rightarrow V\oplus_{\mathbb{R}%
}V^{c} \\
\mathcal{J}_{V} & :\widetilde{V}\rightarrow\widetilde{V^{c}} \\
\mathcal{J}_{V} & :\widetilde{V^{c}}\rightarrow\widetilde{V}
\end{align}
and%
\begin{align}
\mathcal{J}_{V}\left( v,0\right) & =\left( 0,w\right) \\
\mathcal{J}_{V}\left( 0,w\right) & =\left( -v,0\right) \\
& v\in V,w\in V^{c}.
\end{align}
Using $\mathcal{J}_{V}$ we can define multiplication by an imaginary unit $i$
by 
\begin{align}
i\left( v,0\right) \overset{def}{=}\mathcal{J}_{V}\left( v,0\right) &
=\left( 0,w\right) \\
i\left( 0,w\right) \overset{def}{=}\mathcal{J}_{V}\left( 0,w\right) &
=\left( -v,0\right)
\end{align}
to obtain%
\begin{align}
\mathcal{V}_{\mathbb{C}} & =\widetilde{V}+\widetilde{V^{c}}=\widetilde{V}+%
\mathcal{J}_{V}\widetilde{V}=V\oplus iV=\widetilde{V}+\widetilde{iV} \\
& =\left\{ \left. \left( v_{1},0\right) +i\left( v_{2},0\right) \right\vert
v_{1},v_{2}\in V\right\} \leftrightarrow\left\{ \left.
v_{1}+iv_{2}\right\vert v_{1},v_{2}\in V\right\} .
\end{align}

One can use the same complexification map $J_{V}$ to form the
complexification of the complimentary subspace 
\begin{align}
\mathcal{V}_{\mathbb{C}}^{c} & =\widetilde{V^{c}}+\widetilde{V}=\widetilde{%
V^{c}}+\mathcal{J}_{V}\widetilde{V^{c}}=V^{c}\oplus iV^{c}=\widetilde{V^{c}}+%
\widetilde{iV^{c}} \\
& =\left\{ \left. \left( 0,w_{1}\right) +i\left( 0,w_{2}\right) \right\vert
w_{1},w_{2}\in V^{c}\right\} \leftrightarrow\left\{ \left.
w_{1}+iw_{2}\right\vert w_{1},w_{2}\in V^{c}\right\} ,
\end{align}
which leads to an $r$-dimensional complex space $V_{\mathbb{C}}^{c}$ that is
distinct from $V_{\mathbb{C}}$, though it is isomorphic.

This leads to the following identities for complexifications $V_{\mathbb{C}%
},V_{\mathbb{C}}^{c}$ of an element $\left( v,w\right) \in V\oplus _{\mathbb{%
R}}V^{c}$ i.e. 
\begin{align}
\left( v,w\right) & \rightarrow v+iv^{\prime},v^{\prime}\in V \\
\left( v,w\right) & \rightarrow w-iw^{\prime},w^{\prime}\in V^{c}
\end{align}%
\begin{align}
\left( v,w\right) & =\left( v,J_{V}v^{\prime}\right) =\left(
v,iv^{\prime}\right) \leftrightarrow v+iv^{\prime}\in V_{\mathbb{C}%
};v,v^{\prime}\in V \\
w & =J_{V}v^{\prime}\Rightarrow J_{V}w=-v^{\prime}\equiv iw=-v^{\prime} \\
\left( v,w\right) & =\left( -J_{V}w^{\prime},w\right) =\left(
-iw^{\prime},w\right) \leftrightarrow-iw^{\prime}+w\in V_{\mathbb{C}%
}^{c};w,w^{\prime}\in V^{c} \\
v & =-J_{V}w^{\prime}\Rightarrow J_{V}v=w^{\prime}\equiv iv=w^{\prime}.
\end{align}
One could infer from this that for external complexifications 
\begin{equation}
V_{\mathbb{C}}^{c}\cong\overline{V_{\mathbb{C}}}.
\end{equation}

In fact complexification can be seen only to depend on the existance of a
complexification map $J_{VW}$ defined on 
\begin{equation}
\mathcal{V}_{\mathbb{R}}=V\oplus_{\mathbb{R}}W=\widetilde{V}+\widetilde{W}
\end{equation}
i.e.%
\begin{equation}
\mathcal{J}_{VW}:V\oplus_{\mathbb{R}}W\rightarrow V\oplus_{\mathbb{R}}W,
\end{equation}
such that 
\begin{align}
\mathcal{J}_{VW} & :\widetilde{V}\rightarrow\widetilde{W} \\
\mathcal{J}_{VW} & :\widetilde{W}\rightarrow\widetilde{V} \\
\mathcal{J}_{VW}^{2} & =-I_{2r}
\end{align}
and a direct sum decomposition that is not neccesarily disjoint. If an inner
product $\left( \left. {}\right\vert \right) _{\mathcal{V}_{\mathbb{R}}}$is
defined in $\mathcal{V}_{\mathbb{R}}$ that is defined from inner products $%
\left( \left. {}\right\vert \right) _{V}$ and $\left( \left. {}\right\vert
\right) _{W}$ defined in $V$ and $W$ then 
\begin{equation}
\mathcal{J}_{VW}=-\mathcal{J}_{VW}^{t}
\end{equation}
and 
\begin{equation}
\mathcal{J}_{VW}\in O\left( V\oplus_{\mathbb{R}}W,\mathbb{R}\right) .
\end{equation}

\paragraph{Extensions of Complexification}

One can consider the $2r$-dimensional complex space 
\begin{equation}
V_{\mathbb{C}}\oplus_{\mathbb{C}}V_{\mathbb{C}}^{c}
\end{equation}
which corresponds to the \emph{external} complexification of 
\begin{equation}
V\oplus_{\mathbb{R}}V^{c}.
\end{equation}
The same complexification map $J_{V}$ is used to form $V_{\mathbb{C}}$ and $%
V_{\mathbb{C}}^{c}$ and hence $V_{\mathbb{C}}\oplus_{\mathbb{C}}V_{\mathbb{C}%
}^{c}.$

\paragraph{External Complexification in Terms of Internal
Complexification---concrete examples}

A bases of $V\oplus_{\mathbb{R}}V^{c}$ is%
\begin{equation}
\left\{ \left. \left( e_{j},0\right) ,\left( 0,f_{j}\right) \right\vert
1\leq j\leq r\right\} .
\end{equation}
The complexification map $J_{V}$ 
\begin{equation}
\mathcal{J}_{V}:V\oplus_{\mathbb{R}}V^{c}\rightarrow V\oplus_{\mathbb{R}%
}V^{c},
\end{equation}
where $J_{V}\in\mathcal{L}\left( V,V^{c}\right) $ 
\begin{equation}
J_{V}:V\rightarrow V^{c}
\end{equation}
in the case of external complification one defines 
\begin{equation}
i\otimes e_{j}=ie_{j}=J_{V}\left( e_{j}\right) =f_{j},
\end{equation}
and $J_{V}\in\mathcal{L}\left( V^{c},V\right) $ 
\begin{equation}
J_{V}:V^{c}\rightarrow V
\end{equation}
gives%
\begin{equation}
i\otimes f_{j}=if_{j}=J_{V}\left( f_{j}\right) =-e_{j}.
\end{equation}
Summarizing 
\begin{align}
J_{V}\left( e_{j},0\right) & =\left( 0,f_{j}\right) \equiv i\left(
e_{j},0\right) =\left( 0,f_{j}\right) \leftrightarrow J_{V}e_{j}\equiv
ie_{j}=f_{j} \\
J_{V}\left( 0,f_{j}\right) & =\left( -e_{j},0\right) \equiv i\left(
0,f_{j}\right) =\left( -e_{j},0\right) \leftrightarrow J_{V}f_{j}\equiv
if_{j}=-e_{j}.
\end{align}

Multiplication by $i$ induces a real linear map $\mathcal{J}_{V}\in$ $%
\mathcal{L}\left( V\oplus_{\mathbb{R}}iV^{c}\right) =\mathcal{L}\left( 
\mathcal{V}_{\mathbb{C}}\right) ,$ that maintains the direct sum
decomposition and gives a real representation of $\mathbb{C}$ on $V_{\mathbb{%
C}}.$ Its properties are given by its action on the basis elements $\left\{
e_{j},ie_{j}\right\} $ 
\begin{align}
\mathcal{J}_{V}\left( e_{j},0\right) & =\left( 0,ie_{j}\right) \equiv
J_{V}e_{j}=ie_{j} \\
\mathcal{J}_{V}\left( 0,ie_{j}\right) & =\left( -e_{j},0\right) \equiv
J_{V}\left( ie_{j}\right) =-e_{j}.
\end{align}
Using the first identity 
\begin{equation}
ie_{j}=J_{V}e_{j}
\end{equation}
and substituting into the second equation one obtains 
\begin{equation}
J_{V}\left( J_{V}e_{j}\right) =J_{V}^{2}e_{j}=-e_{j},
\end{equation}
and 
\begin{equation}
J_{V}^{2}\left( ie_{j}\right) =J_{V}^{3}\left( e_{j}\right) =-J_{V}\left(
e_{j}\right) =-\left( ie\right) ,
\end{equation}
which shows that as expected 
\begin{equation}
J_{V}^{2}=-I_{V_{\mathbb{C}}}.
\end{equation}

\subsubsection{Invariance Transformations}

Orthogonal transformations $t\in O\left( V\right) $ generate unitary
transformations $u\in U\left( V_{\mathbb{C}}\right) $ as 
\begin{align}
\left\langle \left. ux\right\vert uy\right\rangle & =\left( \left.
ta\right\vert tc\right) +\left( \left. tb\right\vert td\right) +i\left\{
\left( \left. ta\right\vert td\right) -\left( \left. tb\right\vert tc\right)
\right\} \\
& =\left( \left. a\right\vert c\right) +\left( \left. b\right\vert d\right)
+i\left\{ \left( \left. a\right\vert d\right) -\left( \left. b\right\vert
c\right) \right\} =\left\langle \left. x\right\vert y\right\rangle ,
\end{align}
where 
\begin{align}
a,b,c,d & \in V \\
ia,ib,ic,id & \in V^{c}
\end{align}
and the inner product $\left\langle \left. {}\right\vert \right\rangle $ is
expressed in terms of the inner product $\left( \left. {}\right\vert \right) 
$ in $V\oplus_{\mathbb{R}}W,$ and hence in $V$ and $W,$ where 
\begin{equation}
W=\left\{ 
\begin{array}{c}
V^{c} \\ 
V%
\end{array}
\right. ,
\end{equation}
as 
\begin{align}
\left\langle \left. x\right\vert y\right\rangle & =h\left( \left. \left(
a,b\right) \right\vert \left( c,d\right) \right) +ig\left( \left. \left(
a,b\right) \right\vert \left( c,d\right) \right) \\
& =h\left( \left. \left( a,b\right) \right\vert \left( c,d\right) \right)
+ih\left( \left. i\left( a,b\right) \right\vert \left( c,d\right) \right) \\
& =g\left( \left( a,b\right) ,i\left( c,d\right) \right) +ig\left( \left.
\left( a,b\right) \right\vert \left( c,d\right) \right) ,
\end{align}
where 
\begin{align}
h\left( \left. \left( a,b\right) \right\vert \left( c,d\right) \right) &
=\left( \left. a\right\vert c\right) +\left( \left. b\right\vert d\right)
=\left( a,b\right) \left( 
\begin{array}{cc}
I & 0 \\ 
0 & I%
\end{array}
\right) \left( 
\begin{array}{c}
c \\ 
d%
\end{array}
\right) \\
g\left( \left. \left( a,b\right) \right\vert \left( c,d\right) \right) &
=\left( \left. a\right\vert d\right) -\left( \left. b\right\vert c\right)
=\left( a,b\right) \left( 
\begin{array}{cc}
0 & I \\ 
-I & 0%
\end{array}
\right) \left( 
\begin{array}{c}
c \\ 
d%
\end{array}
\right) .
\end{align}
Hence\emph{\ all }unitary transformations $u\in U\left( V_{\mathbb{C}%
}\right) $ are generated by $\mathcal{T}\in O\left( V\oplus_{\mathbb{R}%
}V^{c}\right) \cap Sp\left( V\oplus_{\mathbb{R}}V^{c}\right) =$ $O\left(
V,V^{c}\right) $ $\left( ?\text{One needs to examine this in detail}\right)
\supset O\left( V\right) ,O\left( V^{c}\right) .$

$\lceil$%
\begin{align}
\left\langle \left. ux\right\vert uy\right\rangle & =h\left( \left. \mathcal{%
T}\left( a,b\right) \right\vert \mathcal{T}\left( c,d\right) \right)
+ig\left( \left. \mathcal{T}\left( a,b\right) \right\vert \mathcal{T}\left(
c,d\right) \right) =g\left( \left. \mathcal{T}\left( a,b\right) \right\vert 
\mathcal{T}\left( c,d\right) \right) +ig\left( \left. \mathcal{T}\left(
a,b\right) \right\vert J^{t}\mathcal{T}\left( c,d\right) \right) \\
& =\left( a,b\right) \mathcal{T}^{t}\left( 
\begin{array}{cc}
I & 0 \\ 
0 & I%
\end{array}
\right) \mathcal{T}\left( 
\begin{array}{c}
c \\ 
d%
\end{array}
\right) +i\left( a,b\right) \mathcal{T}^{t}\left( 
\begin{array}{cc}
0 & I \\ 
-I & 0%
\end{array}
\right) \mathcal{T}\left( 
\begin{array}{c}
c \\ 
d%
\end{array}
\right) \,\forall\left( a,b\right) ,\left( c,d\right) ,
\end{align}
where 
\begin{equation}
J^{t}=\left( 
\begin{array}{cc}
0 & I \\ 
-I & 0%
\end{array}
\right) .
\end{equation}
$\rfloor$

Every complex Hilbert space $\mathfrak{H}$ of complex dimension $r$ can be
represented as the complexification of a real Hilbert space (Euclidean
space) $V_{e}\subset\mathfrak{H}$ of real dimension $r$ by choosing an conb $%
\left\{ \left. \left\vert e_{j}\right\rangle \right\vert 1\leq j\leq
r\right\} $ of the Hilbert space $\mathfrak{H}$ and setting 
\begin{equation}
V_{e}=\mathbb{R-}\limfunc{linspan}\left\{ \left. \left\vert
e_{j}\right\rangle \right\vert 1\leq j\leq r\right\} .
\end{equation}
It is easy to see that 
\begin{equation}
\mathfrak{H}=V_{e}\oplus_{\mathbb{R}}iV_{e},
\end{equation}
which carries a \emph{natural conjugation }%
\begin{align}
\sigma_{e} & :\mathfrak{H}\rightarrow\sigma_{e}\left( \mathfrak{H}\right) =%
\mathfrak{H} \\
\sigma_{e}\left( a+ib\right) & =a-ib\,\forall a,b\in V_{e},
\end{align}
where a \emph{conjugation} 
\begin{equation}
\sigma:\mathfrak{H}\rightarrow\sigma\left( \mathfrak{H}\right) =\mathfrak{H}
\end{equation}
is an \emph{antiunitary} $\left( \text{conjugate, complex antilinear and
real linear}\right) $ map with the properties 
\begin{align}
\sigma^{2} & =I_{\mathfrak{H}} \\
\left\langle \left. \sigma\left( x\right) \right\vert \sigma\left( y\right)
\right\rangle & =\overline{\left\langle \left. x\right\vert y\right\rangle }.
\end{align}

Hence a complexification of a complete Euclidean ($\equiv$ real Hilbert)
space $V$ $\left( \text{real dimension }r\right) $ leads to a natural
conjugation on the Hilbert space $V_{\mathbb{C}}$ $\left( \text{complex
dimension }r\right) .$ Conversely a given conjugation $\sigma$ in a complex
Hilbert space $\mathfrak{H}$ $\left( \text{complex dimension }r\right) $
determines a unique way to represent $\mathfrak{H}$ as $V_{\mathbb{C}}$ by
defining a "real" subspace 
\begin{equation}
V_{\sigma}=\left\{ \left. \frac{1}{2}\left( x+\sigma\left( x\right) \right)
\right\vert x\in\mathfrak{H}\right\} =\left\{ \left. \sigma\left(
\varsigma\right) =\varsigma\right\vert \varsigma\in\mathfrak{H}\right\}
\subset\mathfrak{H}
\end{equation}
leading to 
\begin{equation}
\mathfrak{H}\equiv V_{\mathbb{C}}=V_{\sigma}\oplus_{\mathbb{R}}iV_{\sigma}.
\end{equation}
One can also consider the complex direct sum 
\begin{equation}
\mathfrak{H}_{\mathbb{C}}=V_{\sigma}\oplus_{\mathbb{C}}iV_{\sigma }=%
\mathfrak{H}\oplus_{\mathbb{C}}\mathfrak{H,}
\end{equation}
where $\mathfrak{H}_{\mathbb{C}}$ has complex dimension $2r.$

\emph{Conjugations are in }$1-1$\emph{\ correspondence with
complexifications, which are in }$1-1$\emph{\ correspondence with
equivalence classes of bases --- two orthonormal basis being equivalent if
they are related by a real orthogonal transformation. Each equivalence class
corresponds to a real subspace of }$\mathfrak{H}$ \emph{of the same
dimension as} $\mathfrak{H}$\emph{.}

\section{Real and Complex Linear Operators on Complexified and Real Vector
Spaces}

First we collect together a few observations about the real and complex
spaces induced by the complex Hilbert space $\mathfrak{H}$ and the operators
acting on those spaces:

\begin{enumerate}
\item The\ complex Hilbert space $\mathfrak{H}$ is closed and generated by
real linear combinations i.e. 
\begin{equation}
\sum_{1\leq j\leq p}x_{j}v_{j}\in\mathfrak{H\,\forall}\left\{ \left. v_{j}\in%
\mathfrak{H}\right\vert 1\leq j\leq p\right\} ,\left\{ \left. x_{j}\in%
\mathbb{R}\right\vert 1\leq j\leq p\right\} ;p\in\mathbb{N}.
\end{equation}

\item The elements $v$, $iv$ belonging to $\mathfrak{H}$ are real linearly
independent i.e. the only solution to 
\begin{equation}
x_{1}v+x_{2}iv=0;\,x_{1},x_{2}\in\mathbb{R}
\end{equation}
is 
\begin{equation}
x_{1}=x_{2}=0.
\end{equation}

\item The expansion coefficients of $v$ wrt the over-complete basis $\left\{
\left\vert \mathbf{e}\right\rangle ,i\left\vert \mathbf{e}\right\rangle
\right\} $ of $\mathfrak{H}$ are 
\begin{align}
\left\langle \left. e_{j}\right\vert v\right\rangle _{\mathfrak{H}}&
=x_{j}+iy_{j} \\
\left\langle \left. ie_{j}\right\vert v\right\rangle _{\mathfrak{H}}&
=y_{j}-ix_{j}.
\end{align}%
In particular 
\begin{gather}
\left\langle \left. e_{j}\right\vert e_{k}\right\rangle _{\mathfrak{H}%
}=\delta _{jk} \\
\left\langle \left. ie_{j}\right\vert ie_{k}\right\rangle _{\mathfrak{H}%
}=\delta _{jk} \\
\left\langle \left. e_{j}\right\vert ie_{k}\right\rangle _{\mathfrak{H}%
}=i\delta _{jk} \\
1\leq j,k\leq r.
\end{gather}%
The resolution of the identity for this over-complete basis is 
\begin{equation}
\frac{1}{2}\sum \left\{ \left\vert e_{j}\right\rangle \left\langle
e_{j}\right\vert +\left\vert ie_{j}\right\rangle \left\langle
ie_{j}\right\vert \right\} =I_{r}.
\end{equation}

\item One can define the real linear maps 
\begin{align}
\func{Re}_{e}& :\mathfrak{H\rightarrow H} \\
\func{Im}_{e}& :\mathfrak{H\rightarrow H}
\end{align}%
by%
\begin{align}
\func{Re}_{e}\left( v\right) & =\sum_{1\leq j\leq r}\func{Re}\left(
v_{j}\right) e_{j}=\left\vert \mathbf{e}\right\rangle \mathcal{R}_{e}\left( 
\func{Re}_{e}\left( v\right) \right) =\left\vert \mathbf{e}\right\rangle 
\mathbf{x}=\sum_{1\leq j\leq r}\left\vert e_{j}\right\rangle x_{j}\equiv
\sum_{1\leq j\leq r}x_{j}e_{j} \\
\func{Im}_{e}\left( v\right) & =\sum_{1\leq j\leq r}\func{Im}\left(
v_{j}\right) e_{j}=\left\vert \mathbf{e}\right\rangle \mathcal{R}_{e}\left( 
\func{Im}_{e}\left( v\right) \right) =\left\vert \mathbf{e}\right\rangle 
\mathbf{y}=\sum_{1\leq j\leq r}\left\vert e_{j}\right\rangle y_{j}\equiv
\sum_{1\leq j\leq r}y_{j}e_{j},
\end{align}%
which gives 
\begin{equation}
v=\func{Re}_{e}\left( v\right) +i\func{Im}_{e}\left( v\right) .
\end{equation}%
Note that the components $\left\{ \func{Re}_{e}\left( v\right) ,i\func{Im}%
_{e}\left( v\right) \right\} $ are not in general orthogonal as%
\begin{equation}
\left\langle \left. \func{Re}_{e}\left( v\right) \right\vert i\func{Im}%
_{e}\left( v\right) \right\rangle _{\mathfrak{H}}=i\sum_{1\leq j\leq r}\func{%
Re}\left( v_{j}\right) \func{Im}\left( v_{j}\right) ,
\end{equation}%
which does not in general equal zero.

\item $\func{Re}_{e}$ and $\func{Im}_{e}$ define real subspaces $V_{e}$ and $%
iV_{e}$ each of real dimension $r$ that have real conb's $\left\vert \mathbf{%
e}\right\rangle $ and $\left\vert i\mathbf{e}\right\rangle $ that generate $%
\mathfrak{H}$ by a real direct sum i.e. 
\begin{equation}
\mathfrak{H}=\func{Re}_{e}\left( \mathfrak{H}\right) \oplus _{\mathbb{R}}i%
\func{Im}_{e}\left( \mathfrak{H}\right) .
\end{equation}

\item This direct sum is real isomorphic to the real Hilbert space $%
\mathfrak{H}_{\mathbb{R}},$where 
\begin{equation}
\mathfrak{H}_{\mathbb{R}}=\func{Re}_{e}\left( \mathfrak{H}\right) \oplus _{%
\mathbb{R}}\func{Im}_{e}\left( \mathfrak{H}\right)
\end{equation}%
and $\mathbb{R-}\dim \left\{ \mathfrak{H}_{\mathbb{R}}\right\} =2r$ i.e. 
\begin{eqnarray}
\mathfrak{R}_{e} &:&\mathfrak{H\rightarrow H}_{\mathbb{R}} \\
\mathfrak{R}_{e}\left( v\right) &=&\left( \func{Re}_{e}\left( v\right) ,%
\func{Im}_{e}\left( v\right) \right) .
\end{eqnarray}%
$\mathfrak{R}_{e}$ is injective and a real linear mapping i.e. 
\begin{equation}
\mathfrak{R}_{e}\left( v_{1}-v_{2}\right) =0\Leftrightarrow v_{1}=v_{2}
\end{equation}%
and 
\begin{equation}
\mathfrak{R}_{e}\left( \alpha _{1}v_{1}+\alpha _{2}v_{2}\right) =\alpha _{1}%
\mathfrak{R}_{e}\left( v_{1}\right) +\alpha _{2}\mathfrak{R}_{e}\left(
v_{2}\right) \quad \forall \alpha _{1},\alpha _{2}\in \mathbb{R},\,\forall
v_{1},v_{2}\in \mathfrak{H}.
\end{equation}

\item This direct sum decomposition is left invariant by $G=GL\left( r,%
\mathbb{R}\right) +GL\left( r,\mathbb{R}\right) .$ One can see from this
direct sum decomposition and this invariance group that the decomposition is
in $1-1$ correspondence with an equivalence class of basis $\left[ \mathbf{e}%
\right] _{G},$where 
\begin{equation}
\left[ \mathbf{e}\right] _{G}=\left\{ \left. T\left\vert \mathbf{e}%
\right\rangle \right\vert T\in GL\left( r,\mathbb{R}\right) +GL\left( r,%
\mathbb{R}\right) \right\} .
\end{equation}%
If one restricts attention to conb's of $\mathfrak{H}$ the equivalence class
is given by 
\begin{equation}
\left[ \mathbf{e}\right] =\left\{ \left. T\left\vert \mathbf{e}\right\rangle
\right\vert 0\in O\left( r,\mathbb{R}\right) +O\left( r,\mathbb{R}\right)
\right\} .
\end{equation}

\item The inner product in $\mathfrak{H}$ can be expressed as 
\begin{align}
& \left\langle \left. \func{Re}_{e}\left( v_{1}\right) +i\func{Im}_{e}\left(
v_{1}\right) \right\vert \func{Re}_{e}\left( v_{2}\right) +i\func{Im}%
_{e}\left( v_{2}\right) \right\rangle _{\mathfrak{H}} \\
& =\left\langle \left. \func{Re}_{e}\left( v_{1}\right) \right\vert \func{Re}%
_{e}\left( v_{2}\right) \right\rangle _{\mathfrak{H}}+\left\langle \left. 
\func{Re}_{e}\left( v_{1}\right) \right\vert i\func{Im}_{e}\left(
v_{2}\right) \right\rangle _{\mathfrak{H}}+\left\langle \left. i\func{Im}%
_{e}\left( v_{1}\right) \right\vert \func{Re}_{e}\left( v_{2}\right)
\right\rangle _{\mathfrak{H}}+\left\langle \left. i\func{Im}_{e}\left(
v_{1}\right) \right\vert i\func{Im}_{e}\left( v_{2}\right) \right\rangle _{%
\mathfrak{H}} \\
& =\left\langle \left. \func{Re}_{e}\left( v_{1}\right) \right\vert \func{Re}%
_{e}\left( v_{2}\right) \right\rangle _{\mathfrak{H}}+i\left\langle \left. 
\func{Re}_{e}\left( v_{1}\right) \right\vert \func{Im}_{e}\left(
v_{2}\right) \right\rangle _{\mathfrak{H}}-i\left\langle \left. \func{Im}%
_{e}\left( v_{1}\right) \right\vert \func{Re}_{e}\left( v_{2}\right)
\right\rangle _{\mathfrak{H}}+\left\langle \left. \func{Im}_{e}\left(
v_{1}\right) \right\vert \func{Im}_{e}\left( v_{2}\right) \right\rangle _{%
\mathfrak{H}} \\
& =\frame{$\left\langle \left. \func{Re}_{e}\left( v_{1}\right) \right\vert 
\func{Re}_{e}\left( v_{2}\right) \right\rangle _{\mathfrak{H}}$+$%
\left\langle \left. \func{Im}_{e}\left( v_{1}\right) \right\vert \func{Im}%
_{e}\left( v_{2}\right) \right\rangle _{\mathfrak{H}}$}+i\frame{$\left(
\left\langle \left. \func{Re}_{e}\left( v_{1}\right) \right\vert \func{Im}%
_{e}\left( v_{2}\right) \right\rangle _{\mathfrak{H}}-\left\langle \left. 
\func{Im}_{e}\left( v_{1}\right) \right\vert \func{Re}_{e}\left(
v_{2}\right) \right\rangle _{\mathfrak{H}}\right) $} \\
& =\left( \left. \left( \func{Re}_{e}\left( v_{1}\right) ,\func{Im}%
_{e}\left( v_{1}\right) \right) \right\vert \left( \func{Re}_{e}\left(
v_{2}\right) ,\func{Im}_{e}\left( v_{2}\right) \right) \right) _{\mathfrak{H}%
}+i\left\{ \left. \left( \func{Re}_{e}\left( v_{1}\right) ,\func{Im}%
_{e}\left( v_{1}\right) \right) \right\vert \left( \func{Re}_{e}\left(
v_{2}\right) ,\func{Im}_{e}\left( v_{2}\right) \right) \right\} _{\mathfrak{H%
}},
\end{align}%
which defines a \emph{positive definite inner product} $h\left( ,\right)
=\left( \left. {}\right\vert \right) _{\mathfrak{H}_{\mathbb{R}}}$ in $%
\mathfrak{H}_{\mathbb{R}}$ 
\begin{equation}
\left( \left. \mathfrak{R}_{e}\left( v_{1}\right) \right\vert \mathfrak{R}%
_{e}\left( v_{2}\right) \right) _{\mathfrak{H}_{\mathbb{R}}}=\left\langle
\left. \func{Re}_{e}\left( v_{1}\right) \right\vert \func{Re}_{e}\left(
v_{2}\right) \right\rangle _{\mathfrak{H}}+\left\langle \left. \func{Im}%
_{e}\left( v_{1}\right) \right\vert \func{Im}_{e}\left( v_{2}\right)
\right\rangle _{\mathfrak{H}}=\func{Re}\left\{ \left\langle \left.
v_{1}\right\vert v_{2}\right\rangle _{\mathfrak{H}}\right\}
\end{equation}%
and a symplectic form $g\left( ,\right) =\left\{ \left. {}\right\vert
\right\} _{\mathfrak{H}_{\mathbb{R}}}$ in $\mathfrak{H}_{\mathbb{R}}$%
\begin{equation}
\left\{ \left. \mathfrak{R}_{e}\left( v_{1}\right) \right\vert \mathfrak{R}%
_{e}\left( v_{2}\right) \right\} _{\mathfrak{H}_{\mathbb{R}}}=\left\langle
\left. \func{Re}_{e}\left( v_{1}\right) \right\vert \func{Im}_{e}\left(
v_{2}\right) \right\rangle _{\mathfrak{H}}-\left\langle \left. \func{Im}%
_{e}\left( v_{1}\right) \right\vert \func{Re}_{e}\left( v_{2}\right)
\right\rangle _{\mathfrak{H}}=\func{Im}\left\{ \left\langle \left.
v_{1}\right\vert v_{2}\right\rangle _{\mathfrak{H}}\right\} .
\end{equation}

\begin{remark}
\begin{equation}
\func{Re}\left\{ \left\langle \left. v_{1}\right\vert v_{2}\right\rangle _{%
\mathfrak{H}}\right\} =\left\langle \left. \func{Re}_{e}\left( v_{1}\right)
\right\vert \func{Re}_{e}\left( v_{2}\right) \right\rangle _{\mathfrak{H}%
}+\left\langle \left. \func{Im}_{e}\left( v_{1}\right) \right\vert \func{Im}%
_{e}\left( v_{2}\right) \right\rangle _{\mathfrak{H}}
\end{equation}

and 
\begin{equation}
\func{Im}\left\{ \left\langle \left. v_{1}\right\vert v_{2}\right\rangle _{%
\mathfrak{H}}\right\} =\left\langle \left. \func{Re}_{e}\left( v_{1}\right)
\right\vert \func{Im}_{e}\left( v_{2}\right) \right\rangle _{\mathfrak{H}%
}-\left\langle \left. \func{Im}_{e}\left( v_{1}\right) \right\vert \func{Re}%
_{e}\left( v_{2}\right) \right\rangle _{\mathfrak{H}}
\end{equation}

are \emph{basis} \emph{independent i.e. these equalities are true for all }$%
\left\vert \mathbf{e}\right\rangle $\emph{.}
\end{remark}

\item Hence the inner prouct in $\mathfrak{H}$ can be decomposed into a sum
of a real symmetric inner product 
\begin{equation}
\left( \left. \left( \func{Re}_{e}\left( v_{1}\right) ,\func{Im}_{e}\left(
v_{1}\right) \right) \right\vert \left( \func{Re}_{e}\left( v_{2}\right) ,%
\func{Im}_{e}\left( v_{2}\right) \right) \right) _{\mathfrak{H}}
\end{equation}%
and a real symplectic inner product 
\begin{equation}
\left\{ \left. \left( \func{Re}_{e}\left( v_{1}\right) ,\func{Im}_{e}\left(
v_{1}\right) \right) \right\vert \left( \func{Re}_{e}\left( v_{2}\right) ,%
\func{Im}_{e}\left( v_{2}\right) \right) \right\} _{\mathfrak{H}}
\end{equation}%
i.e.%
\begin{equation}
\left\langle \left. v_{1}\right\vert v_{2}\right\rangle _{\mathfrak{H}%
}=\left( \left. \left( \func{Re}_{e}\left( v_{1}\right) ,\func{Im}_{e}\left(
v_{1}\right) \right) \right\vert \left( \func{Re}_{e}\left( v_{2}\right) ,%
\func{Im}_{e}\left( v_{2}\right) \right) \right) _{\mathfrak{H}}+\left\{
\left. \left( \func{Re}_{e}\left( v_{1}\right) ,\func{Im}_{e}\left(
v_{1}\right) \right) \right\vert \left( \func{Re}_{e}\left( v_{2}\right) ,%
\func{Im}_{e}\left( v_{2}\right) \right) \right\} _{\mathfrak{H}}
\end{equation}%
for \emph{any} conb $\left\vert \mathbf{e}\right\rangle $ of $\mathfrak{H}$
i.e. is \emph{basis independent}$\mathfrak{.}$

\item The real linear mapping $\mathfrak{R}_{e}$ is an isometric isomorphism
i.e a real Hilbert space isomorphism where the inner product in $\mathfrak{H}
$ is $\left\langle \left. {}\right\vert \right\rangle _{\mathfrak{H}}$ and
the inner product in $\mathfrak{H}_{\mathbb{R}}$ is $\left( \left.
{}\right\vert \right) _{\mathfrak{H}_{\mathbb{R}}}$ as 
\begin{equation}
\left\langle \left. v\right\vert v\right\rangle _{\mathfrak{H}}=\left(
\left. \mathfrak{R}_{e}\left( v\right) \right\vert \mathfrak{R}_{e}\left(
v\right) \right) _{\mathfrak{H}_{\mathbb{R}}}.
\end{equation}

\item The invariance group of $\left( \left. {}\right\vert \right) _{%
\mathfrak{H}_{\mathbb{R}}}$ is $O\left( 2r,\mathbb{R}\right) $ and the
invariance group $\left\{ \left. \cdot \right\vert \cdot \right\} _{%
\mathfrak{H}_{\mathbb{R}}}$ is $Sp\left( 2r,\mathbb{R}\right) $, thus the
invariance group of, $\left\langle \left. {}\right\vert \right\rangle _{%
\mathfrak{H}},$ $U\left( r\right) =O\left( 2r,\mathbb{R}\right) \cap
Sp\left( 2r,\mathbb{R}\right) .$

\item A \emph{real symmetric} inner product in $\func{Re}_{e}\left( 
\mathfrak{H}\right) $ with invariance gr $O\left( r,\mathbb{R}\right) $ is
given by 
\begin{equation}
\left( \left. o\func{Re}_{e}\left( v_{1}\right) \right\vert o\func{Re}%
_{e}\left( v_{2}\right) \right) _{\func{Re}_{e}\left( \mathfrak{H}\right)
}=\left( \left. \func{Re}_{e}\left( v_{1}\right) \right\vert \func{Re}%
_{e}\left( v_{2}\right) \right) _{\func{Re}_{e}\left( \mathfrak{H}\right)
},\,o\in O\left( r,\mathbb{R}\right) .
\end{equation}

\item For any conb $\left\vert \mathbf{e}\right\rangle $ of $\mathfrak{H}$
the elements $\left\{ \left. e_{j},ie_{j}\right\vert 1\leq j\leq r\right\} $
are $\mathbb{R}$-linearly independent and form a conb basis for a real
Hilbert space $\mathfrak{H}_{\mathbb{R}}$ of real dimension $2r,$ producing
the \emph{real linear} \emph{isomorphism} 
\begin{equation}
\mathfrak{R}_{e}:\mathfrak{H\rightarrow H}_{\mathbb{R}}=\func{Re}_{e}\left( 
\mathfrak{H}\right) \oplus _{\mathbb{R}}\func{Im}_{e}\left( \mathfrak{H}%
\right)
\end{equation}%
given by 
\begin{equation}
\mathfrak{R}_{e}\left( v\right) =\mathfrak{R}_{e}\left( \sum_{1\leq j\leq
r}x_{j}e_{j}+y_{j}ie_{j}\right) =\sum_{1\leq j\leq r}\left\{ x_{j}\mathfrak{R%
}_{e}\left( e_{j}\right) +y_{j}\mathfrak{R}_{e}\left( ie_{j}\right) \right\}
=\left( \func{Re}_{e}\left( v\right) ,\func{Im}_{e}\left( v\right) \right) .
\end{equation}%
The real Hilbert space $\mathfrak{H}_{\mathbb{R}}$ does not depend on the
basis $\left\vert \mathbf{e}\right\rangle ,$but the mapping $\mathfrak{R}%
_{e} $ obviously \emph{does}. A \emph{real symmetric} inner product in $%
\mathfrak{H}_{\mathbb{R}}$ with invariance group $O\left( 2r,\mathbb{R}%
\right) $ is given by 
\begin{equation}
\left( \left. \left( \func{Re}_{e}\left( v_{1}\right) ,\func{Im}_{e}\left(
v_{1}\right) \right) \right\vert \left( \func{Re}_{e}\left( v_{2}\right) ,%
\func{Im}_{e}\left( v_{2}\right) \right) \right) _{\mathfrak{H}_{\mathbb{R}%
}}=\left( \left. \left( \func{Re}_{e}\left( v_{1}\right) ,\func{Im}%
_{e}\left( v_{1}\right) \right) \right\vert \left( \func{Re}_{e}\left(
v_{2}\right) ,\func{Im}_{e}\left( v_{2}\right) \right) \right) _{\mathfrak{H}%
}=\func{Re}\left\{ \left\langle \left. v_{1}\right\vert v_{2}\right\rangle _{%
\mathfrak{H}}\right\} .
\end{equation}%
This makes $\mathfrak{R}_{e}$ a real Hilbert space isometric isomorphism.

\item $\left( a\right) \quad \mathfrak{R}_{e}$ is $O(r,\mathbb{R)}$
equivalent to $\mathfrak{R}_{f}$ written as 
\begin{equation}
\mathfrak{R}_{e}\overset{O(r,\mathbb{R)}}{\backsim }\mathfrak{R}_{f}
\end{equation}%
if $\exists o\in O(r,\mathbb{R)}+O(r,\mathbb{R)\subset }O(2r,\mathbb{R)}$
and 
\begin{equation}
\mathfrak{R}_{e}=o\mathfrak{R}_{f}o^{t}.
\end{equation}

$\left( b\right) \quad \mathfrak{R}_{e}$ is $U(r,\mathbb{R)}$ equivalent to $%
\mathfrak{R}_{f}$ written as 
\begin{equation}
\mathfrak{R}_{e}\overset{U(r,\mathbb{R)}}{\backsim }\mathfrak{R}_{f}
\end{equation}%
if $\exists u\in U(r,\mathbb{R)}$ $=O\left( 2r,\mathbb{R}\right) \cap
Sp\left( 2r,\mathbb{R}\right) \subset O(2r,\mathbb{R)}$ and 
\begin{equation}
\mathfrak{R}_{e}=u\mathfrak{R}_{f}u^{t}.
\end{equation}

$\left( c\right) \quad \mathfrak{R}_{e}$ is $O(2r,\mathbb{R)}$ equivalent to 
$\mathfrak{R}_{f}$ written as 
\begin{equation}
\mathfrak{R}_{e}\overset{O(2r,\mathbb{R)}}{\backsim }\mathfrak{R}_{f}
\end{equation}%
if $\exists q\in O(2r,\mathbb{R)}$ and $\mathfrak{R}_{e}=q\mathfrak{R}%
_{f}q^{t}.$

These equivalences are related by 
\begin{equation}
\mathfrak{R}_{e}\overset{O(r,\mathbb{R)}}{\backsim }\mathfrak{R}%
_{f}\Rightarrow \mathfrak{R}_{e}\overset{U(r,\mathbb{R)}}{\backsim }%
\mathfrak{R}_{f}\Rightarrow \mathfrak{R}_{e}\overset{O(2r,\mathbb{R)}}{%
\backsim }\mathfrak{R}_{f}.
\end{equation}

\item The Hilbert space $\mathfrak{H}$ can be expressed as the direct sum 
\begin{gather}
\mathfrak{H}=\func{Re}_{e}\left( \mathfrak{H}\right) \oplus _{\mathbb{R}}i%
\func{Im}_{e}\left( \mathfrak{H}\right) \\
\mathbb{C}-\dim \left\{ \mathfrak{H}\right\} =r,
\end{gather}%
for \emph{any} conb $\left\vert \mathbf{e}\right\rangle $ of $\mathfrak{H.}$
The complex direct sum gives the complex Hilbert space $\mathfrak{H}_{%
\mathbb{C}}$ 
\begin{gather}
\mathfrak{H}_{\mathbb{C}}=\func{Re}_{e}\left( \mathfrak{H}\right) \oplus _{%
\mathbb{C}}i\func{Im}_{e}\left( \mathfrak{H}\right) =\func{Re}_{e}\left( 
\mathfrak{H}\right) \oplus _{\mathbb{C}}\func{Im}_{e}\left( \mathfrak{H}%
\right) \\
\mathbb{C}-\dim \left\{ \mathfrak{H}\right\} =2r.
\end{gather}%
The inner product in $\mathfrak{H}_{\mathbb{C}}$ is 
\begin{equation}
\left\langle \left. \left( v_{1},v_{2}\right) \right\vert \left(
w_{1},w_{2}\right) \right\rangle _{\mathfrak{H}_{\mathbb{C}}}=\left\langle
\left. v_{1}\right\vert w_{1}\right\rangle _{\mathfrak{H}}+\left\langle
\left. v_{2}\right\vert w_{2}\right\rangle _{\mathfrak{H}}.
\end{equation}%
$\lceil $Direct sum is defined by 
\begin{align}
\lambda \left( v_{1},v_{2}\right) & \sim \left( \lambda v_{1},\lambda
v_{2}\right) \\
\left( v_{1},v_{2}\right) +\left( v_{3},v_{4}\right) & \sim \left(
v_{1}+v_{3},v_{2}+v_{4}\right) \\
\lambda _{1}\left( v_{1},v_{2}\right) +\lambda _{2}\left( v_{3},v_{4}\right)
& \sim \left( \lambda _{1}v_{1}+\lambda _{2}v_{3},\lambda _{1}v_{2}+\lambda
_{2}v_{4}\right) .
\end{align}%
In the cases of $\mathfrak{H}=\mathbb{C}$ this becomes 
\begin{gather}
\mathfrak{H}=\mathbb{C}=\func{Re}_{e}\left( \mathbb{C}\right) \oplus _{%
\mathbb{R}}i\func{Im}_{e}\left( \mathbb{C}\right) =\mathbb{R}\oplus _{%
\mathbb{R}}i\mathbb{R}=\mathbb{R}-\limfunc{linspan}\left\{ \left( 1,0\right)
,\left( 0,i1\right) \right\} \\
=\left\{ \left. \left( x,iy\right) =x+iy\right\vert \quad \forall x,y\in 
\mathbb{R}\right\} .
\end{gather}%
In the case $\mathfrak{H}_{\mathbb{R}}=\mathbb{R}$ this becomes%
\begin{gather}
\mathfrak{H}_{\mathbb{R}}=\mathbb{R}^{2}=\mathbb{R}\oplus _{\mathbb{R}}%
\mathbb{R}=\mathbb{R}-\limfunc{linspan}\left\{ \left( 1,0\right) ,\left(
0,1\right) \right\} \\
=\left\{ \left. \left( x,y\right) \right\vert \forall x,y\in \mathbb{R}%
\right\} \cong \mathbb{R}\oplus _{\mathbb{R}}i\mathbb{R}=\mathbb{C} \\
\text{Note that }\left( x,y\right) \neq x+y.
\end{gather}%
In the case $\mathfrak{H}_{\mathbb{C}}=\mathbb{C}^{2}$ this becomes%
\begin{gather}
\mathfrak{H}_{\mathbb{C}}=\mathbb{C}^{2}=\mathbb{R}\oplus _{\mathbb{C}}i%
\mathbb{R}=\mathbb{C}-\limfunc{linspan}\left\{ \left( 1,0\right) ,\left(
0,i1\right) \right\} \\
=\left\{ \left( z,w\right) \forall z,w\in \mathbb{C}\right\} \cong \mathbb{C}%
\oplus _{\mathbb{C}}\mathbb{C} \\
\text{Note that }\left( z,w\right) \neq z+w.
\end{gather}%
$\rfloor $

$\lceil $The product of two complex numbers is 
\begin{equation}
z_{1}z_{2}=\left( x_{1}+iy_{1}\right) \left( x_{2}+iy_{2}\right)
=x_{1}x_{2}-y_{1}y_{2}+i\left( x_{1}y_{2}+y_{1}x_{2}\right) .
\end{equation}%
$\rfloor $

\item $T\in \mathcal{B}\left( \mathfrak{H}_{\mathbb{R}}\right) $%
\begin{equation}
T=\left( 
\begin{array}{cc}
\left( \left\vert \mathbf{e}\right\rangle ,\mathbf{0}\right) & \left( 
\mathbf{0,}\left\vert \mathbf{e}\right\rangle \right)%
\end{array}%
\right) \left( 
\begin{array}{cc}
\mathbb{T}_{11} & \mathbb{T}_{12} \\ 
\mathbb{T}_{21} & \mathbb{T}_{22}%
\end{array}%
\right) \left( 
\begin{array}{c}
\left( \left\langle \mathbf{e}\right\vert ,\mathbf{0}\right) \\ 
\left( \mathbf{0,}\left\langle \mathbf{e}\right\vert \right)%
\end{array}%
\right) \equiv \left( 
\begin{array}{cc}
\left\vert \mathbf{\varepsilon }_{1}\right\rangle & \left\vert \mathbf{%
\varepsilon }_{2}\right\rangle%
\end{array}%
\right) \left( 
\begin{array}{cc}
\mathbb{T}_{11} & \mathbb{T}_{12} \\ 
\mathbb{T}_{21} & \mathbb{T}_{22}%
\end{array}%
\right) \left( 
\begin{array}{c}
\left\langle \mathbf{\varepsilon }_{1}\right\vert \\ 
\left\langle \mathbf{\varepsilon }_{2}\right\vert%
\end{array}%
\right) ,\mathbb{T}_{jk}\in \mathbb{R},
\end{equation}%
$T\in \mathcal{B}\left( \mathfrak{H}_{\mathbb{C}}\right) $%
\begin{equation}
T=\left( 
\begin{array}{cc}
\left( \left\vert \mathbf{e}\right\rangle ,\mathbf{0}\right) & \left( 
\mathbf{0,}\left\vert \mathbf{e}\right\rangle \right)%
\end{array}%
\right) \left( 
\begin{array}{cc}
\mathbb{T}_{11} & \mathbb{T}_{12} \\ 
\mathbb{T}_{21} & \mathbb{T}_{22}%
\end{array}%
\right) \left( 
\begin{array}{c}
\left( \left\langle \mathbf{e}\right\vert ,\mathbf{0}\right) \\ 
\left( \mathbf{0,}\left\langle \mathbf{e}\right\vert \right)%
\end{array}%
\right) \equiv \left( 
\begin{array}{cc}
\left\vert \mathbf{\varepsilon }_{1}\right\rangle & \left\vert \mathbf{%
\varepsilon }_{2}\right\rangle%
\end{array}%
\right) \left( 
\begin{array}{cc}
\mathbb{T}_{11} & \mathbb{T}_{12} \\ 
\mathbb{T}_{21} & \mathbb{T}_{22}%
\end{array}%
\right) \left( 
\begin{array}{c}
\left\langle \mathbf{\varepsilon }_{1}\right\vert \\ 
\left\langle \mathbf{\varepsilon }_{2}\right\vert%
\end{array}%
\right) ,\mathbb{T}_{jk}\in \mathbb{C},
\end{equation}%
$T\in \mathcal{B}\left( \mathfrak{H}\right) $%
\begin{equation}
T=\left\vert \mathbf{e}\right\rangle \mathbb{T}_{11}\left\langle \mathbf{e}%
\right\vert ,\mathbb{T}_{jk}\in \mathbb{C}.
\end{equation}

\item The Hilbert spaces $\mathfrak{H}_{\mathbb{R}}$ and $\mathfrak{H}_{%
\mathbb{C}}$ carry representations of $U\left( r\right) ,O\left( 2r,\mathbb{R%
}\right) ,$ $O\left( r,r,\mathbb{R}\right) $ and $Sp\left( 2r,\mathbb{R}%
\right) $ wrt the complex conjugate basis $\left\{ \left( \left\vert \mathbf{%
e}\right\rangle ,\mathbf{0}\right) ,\left( \mathbf{0,}\sigma _{e}\left(
\left\vert \mathbf{e}\right\rangle \right) \right) \right\} $ as 
\begin{gather}
o\in O\left( 2r,\mathbb{R}\right) \subset U\left( 2r\right) \subset GL\left(
2r,\mathbb{C}\right) \\
o=\left( 
\begin{array}{cc}
\left( \left\vert \mathbf{e}\right\rangle ,\mathbf{0}\right) & \left( 
\mathbf{0,}\sigma _{e}\left( \left\vert \mathbf{e}\right\rangle \right)
\right)%
\end{array}%
\right) \left( 
\begin{array}{cc}
\mathbb{U} & \overline{\mathbb{V}} \\ 
\mathbb{V} & \overline{\mathbb{U}}%
\end{array}%
\right) \left( 
\begin{array}{c}
\left( \left\langle \mathbf{e}\right\vert ,\mathbf{0}\right) \\ 
\left( \mathbf{0,}\sigma _{e}\left( \left\langle \mathbf{e}\right\vert
\right) \right)%
\end{array}%
\right) =\left( 
\begin{array}{cc}
\left\vert \mathbf{\varepsilon }_{1}\right\rangle & \left\vert \mathbf{%
\varepsilon }_{2}\right\rangle%
\end{array}%
\right) \left( 
\begin{array}{cc}
\mathbb{U} & \overline{\mathbb{V}} \\ 
\mathbb{V} & \overline{\mathbb{U}}%
\end{array}%
\right) \left( 
\begin{array}{c}
\left\langle \mathbf{\varepsilon }_{1}\right\vert \\ 
\left\langle \mathbf{\varepsilon }_{2}\right\vert%
\end{array}%
\right) \\
=\left( 
\begin{array}{cc}
\mathbf{x}_{1} & \mathbf{x}_{2}%
\end{array}%
\right) \left( 
\begin{array}{cc}
\func{Re}\left( \mathbb{U+V}\right) & -\func{Im}\left( \mathbb{U+V}\right)
\\ 
\func{Im}\left( \mathbb{U-V}\right) & \func{Re}\left( \mathbb{U-V}\right)%
\end{array}%
\right) \left( 
\begin{array}{c}
\mathbf{x}_{1} \\ 
\mathbf{x}_{2}%
\end{array}%
\right)
\end{gather}%
and 
\begin{gather}
u\in U\left( r\right) \subset O\left( 2r,\mathbb{R}\right) \\
u=\left( 
\begin{array}{cc}
\left( \left\vert \mathbf{e}\right\rangle ,\mathbf{0}\right) & \left( 
\mathbf{0,}\sigma _{e}\left( \left\vert \mathbf{e}\right\rangle \right)
\right)%
\end{array}%
\right) \left( 
\begin{array}{cc}
\mathbb{U} & \mathbb{O} \\ 
\mathbb{O} & \overline{\mathbb{U}}%
\end{array}%
\right) \left( 
\begin{array}{c}
\left( \left\langle \mathbf{e}\right\vert ,\mathbf{0}\right) \\ 
\left( \mathbf{0,}\sigma _{e}\left( \left\langle \mathbf{e}\right\vert
\right) \right)%
\end{array}%
\right) =\left( 
\begin{array}{cc}
\left\vert \mathbf{\varepsilon }_{1}\right\rangle & \left\vert \mathbf{%
\varepsilon }_{2}\right\rangle%
\end{array}%
\right) \left( 
\begin{array}{cc}
\mathbb{U} & \mathbb{O} \\ 
\mathbb{O} & \overline{\mathbb{U}}%
\end{array}%
\right) \left( 
\begin{array}{c}
\left\langle \mathbf{\varepsilon }_{1}\right\vert \\ 
\left\langle \mathbf{\varepsilon }_{2}\right\vert%
\end{array}%
\right) \\
=\left( 
\begin{array}{cc}
\mathbf{x}_{1} & \mathbf{x}_{2}%
\end{array}%
\right) \left( 
\begin{array}{cc}
\func{Re}\left( \mathbb{U}\right) & -\func{Im}\left( \mathbb{U}\right) \\ 
\func{Im}\left( \mathbb{U}\right) & \func{Re}\left( \mathbb{U}\right)%
\end{array}%
\right) \left( 
\begin{array}{c}
\mathbf{x}_{1} \\ 
\mathbf{x}_{2}%
\end{array}%
\right) .
\end{gather}
\end{enumerate}

\subsection{Matrix Representation of the Complex Numbers}

\subsubsection{Complex Numbers $\mathbb{C}$}

Complex numbers can be expressed

\begin{enumerate}
\item Polar Form%
\begin{equation*}
z=r\left( \overset{}{\cos \left( \theta \right) +i\sin \left( \theta \right) 
}\right) ,\quad r\in \left\vert 0,\infty \right) ,\quad \theta \in
\left\vert 0,2\pi \right) .
\end{equation*}

\item Polar Form using Eulers Formula%
\begin{equation}
z=re^{i\theta },\quad r\in \left\vert 0,\infty \right) ,\quad \theta \in
\left\vert 0,2\pi \right) .
\end{equation}

\item Reectangular Form 
\begin{equation}
z=\func{Re}\left( z\right) +i\func{Im}\left( z\right) ,\quad \func{Re}\left(
z\right) ,\func{Im}\left( z\right) \in \left( -\infty ,\infty \right) .
\end{equation}
\end{enumerate}

\begin{remark}
\begin{eqnarray}
e^{2\pi i} &=&1 \\
e^{\pi i} &=&-1.
\end{eqnarray}
\end{remark}

Let $r=1:$ Two different basis of $\mathbb{C}^{1}$ are $\left\{ 1\right\} $
and $\left\{ i\right\} $ as 
\begin{equation}
\mathbb{C}^{1}=\left\{ \left. z1\right\vert z\in \mathbb{C}\right\} =\left\{
\left. zi\right\vert z\in \mathbb{C}\right\} .
\end{equation}%
They are also bases of $\mathbb{R}^{1}$ as 
\begin{equation}
\mathbb{R}^{1}=\left\{ \left. x1\right\vert x\in \mathbb{R}\right\} \cong
\left\{ \left. xi\right\vert x\in \mathbb{R}\right\} .
\end{equation}%
$\mathbb{C}^{1}$ can be expressed as the real direct sum 
\begin{equation}
\mathbb{C}^{1}=\mathbb{R}^{1}\oplus _{\mathbb{R}}i\mathbb{R}^{1}=\left\{
\left. \left( x1,yi\right) \right\vert x,y\in \mathbb{R}\right\} \overset{%
\text{as }\mathbb{R}^{1}\text{ and }i\mathbb{R}^{1}\text{ are disjoint}}{=}%
\left\{ \left. x1+iy\right\vert x,y\in \mathbb{R}\right\} ,
\end{equation}%
which shows that $\left\{ \left( 1,0\right) ,\left( 0,i\right) \right\} $ is
a basis for $\mathbb{R}^{1}\oplus _{\mathbb{R}}i\mathbb{R}^{1}$ and 
\begin{equation}
\mathbb{C}^{1}=\mathbb{R}^{1}\oplus _{\mathbb{R}}i\mathbb{R}^{1}=\mathbb{R}-%
\limfunc{linspan}\left\{ \left( 1,0\right) ,\left( 0,i\right) \right\} .
\end{equation}%
On the other hand%
\begin{equation}
\mathbb{C}^{2}=\mathbb{R}^{1}\oplus _{\mathbb{C}}i\mathbb{R}^{1}=\mathbb{C}-%
\limfunc{linspan}\left\{ \left( 1,0\right) ,\left( 0,i\right) \right\}
\end{equation}%
and in terms of another basis 
\begin{equation}
\mathbb{C}^{2}=\mathbb{R}^{1}\oplus _{\mathbb{C}}\mathbb{R}^{1}=\mathbb{C}-%
\limfunc{linspan}\left\{ \left( 1,0\right) ,\left( 0,1\right) \right\} ,
\end{equation}%
while%
\begin{equation}
\mathbb{R}^{2}=\mathbb{R}^{1}\oplus _{\mathbb{R}}\mathbb{R}^{1}=\mathbb{R}%
\text{-}\limfunc{linspan}\left\{ \left( 1,0\right) ,\left( 0,1\right)
\right\} \cong \mathbb{R}^{1}\oplus _{\mathbb{R}}i\mathbb{R}^{1}=\mathbb{R}%
\text{-}\limfunc{linspan}\left\{ \left( 1,0\right) ,\left( 0,i\right)
\right\} .
\end{equation}

\begin{definition}
One can consider a representation $\rho $ of $\mathbb{C}$ (real $c^{\ast }$%
-algebra homomorphism) given by the \emph{real} injective but not surjective
mapping 
\begin{equation}
\rho :\mathbb{C}\rightarrow \mathcal{B}\left( \mathbb{R}^{2}\right) =%
\mathcal{M}\left( 2,2,\mathbb{R}\right)
\end{equation}%
defined by 
\begin{equation}
\rho \left( z\right) =\left( 
\begin{array}{cc}
x & -y \\ 
y & x%
\end{array}%
\right) =\left( 
\begin{array}{cc}
\func{Re}\left\{ z\right\} & -\func{Im}\left\{ z\right\} \\ 
\func{Im}\left\{ z\right\} & \func{Re}\left\{ z\right\}%
\end{array}%
\right) =\frac{1}{2}\left( 
\begin{array}{cc}
z+\overline{z} & i\left( z-\overline{z}\right) \\ 
-i\left( z-\overline{z}\right) & z+\overline{z}%
\end{array}%
\right) =\left( 
\begin{array}{cc}
r\cos \theta & -r\sin \theta \\ 
r\sin \theta & r\cos \theta%
\end{array}%
\right) =r\left( 
\begin{array}{cc}
\cos \theta & -\sin \theta \\ 
\sin \theta & \cos \theta%
\end{array}%
\right) ,
\end{equation}%
where
\end{definition}

\begin{enumerate}
\item 
\begin{eqnarray}
z &=&x+iy=re^{i\theta } \\
z^{\ast } &\equiv &\overline{z}=x-iy=re^{-i\theta } \\
z^{-1} &=&r^{-1}e^{-i\theta } \\
\theta &\in &\left\vert 0,2\pi \right) .
\end{eqnarray}

\item 
\begin{eqnarray}
r &=&+\sqrt{x^{2}+y^{2}} \\
\theta &=&\tan ^{-1}\left\{ \frac{y}{x}\right\} ,\quad \theta \in \left\vert
0,2\pi \right) .
\end{eqnarray}

\item 
\begin{equation}
\rho \left( z^{-1}\right) =\rho \left( z\right) ^{-1}=r^{-1}\left( 
\begin{array}{cc}
\cos \theta & \sin \theta \\ 
-\sin \theta & \cos \theta%
\end{array}%
\right)
\end{equation}%
as 
\begin{equation}
\left( 
\begin{array}{cc}
\cos \theta & \sin \theta \\ 
-\sin \theta & \cos \theta%
\end{array}%
\right) \left( 
\begin{array}{cc}
\cos \theta & -\sin \theta \\ 
\sin \theta & \cos \theta%
\end{array}%
\right) =\left( 
\begin{array}{cc}
\cos ^{2}\theta +\sin ^{2}\theta & -\cos \theta \sin \theta +\sin \theta
\cos \theta \\ 
-\sin \theta \cos \theta +\cos \theta \sin \theta & \sin ^{2}\theta +\cos
^{2}\theta%
\end{array}%
\right) =\left( 
\begin{array}{cc}
1 & 0 \\ 
0 & 1%
\end{array}%
\right)
\end{equation}%
and 
\begin{equation}
\left( 
\begin{array}{cc}
\cos \theta & \sin \theta \\ 
-\sin \theta & \cos \theta%
\end{array}%
\right) ^{-1}=\left( 
\begin{array}{cc}
\cos \theta & \sin \theta \\ 
-\sin \theta & \cos \theta%
\end{array}%
\right) ^{t}
\end{equation}

\item 
\begin{equation}
\rho \left( z^{\ast }\right) =\rho \left( z\right) ^{t}
\end{equation}
\end{enumerate}

\begin{definition}
\begin{enumerate}
\item Conformal maps are transforms that preserve angles,

\item Linear Conformal = locally a similarity transformation (rotation +
uniform scaling).

\item Holomorphic maps with non-zero derivatives i.e. a smooth maps 
\begin{equation}
f:U\subset \mathbb{R}^{n}\rightarrow \mathbb{R}^{n}
\end{equation}%
are conformal at a point $p$ if its differential $df(p)$ (a linear map) is a
similarity transformation. That is:

\item 
\begin{equation}
df(p)=\U{3bb} (p)\cdot O(p),
\end{equation}

\item where%
\begin{equation}
\U{3bb} (p)>0
\end{equation}%
is a scale factor depending on $p$ and $O(p)$ is an orthogonal matrix
(rotation/reflection).
\end{enumerate}
\end{definition}

\begin{remark}
Linear conformal maps are real transforms that preserve angles, though not
all conformal linear transforms are orthogonal. In classical terms this is
the difference between \emph{congruence} and \emph{similarity}, as
exemplified by $SSS$ (side-side-side) \emph{congruence of triangles} and $%
AAA $ (angle-angle-angle) \emph{similarity of triangles}. The group of
conformal linear maps of $\mathbb{R}^{n}$ is denoted $CO(n,\mathbb{R})$ for
the conformal orthogonal group, and consists of the product of the
orthogonal group with the group of dilations. If $n$ is odd, these two
subgroups do not intersect, and they are a \emph{direct product}: 
\begin{equation}
CO(2k+1)=O(2k+1,\mathbb{R})\times \mathbb{R}_{\ast },
\end{equation}%
where $\mathbb{R}_{\ast }=\mathbb{R}\setminus \{0\}$ is the real
multiplicative group, while if $n$ is even, these subgroups intersect in $%
\pm I$, so this is not a direct product, but it is a direct product with the
subgroup of dilation by a positive scalar: 
\begin{equation}
CO(2k,\mathbb{R})=O(2k,\mathbb{R})\times \mathbb{R}_{+}.
\end{equation}%
Similarly one can define $CSO(n,\mathbb{R})$; note that this is always: 
\begin{equation}
CSO(n,\mathbb{R})=CO(n,\mathbb{R})\cap GL_{+}(n,\mathbb{R})=SO(n,\mathbb{R}%
)\times \mathbb{R}_{+}.
\end{equation}
\end{remark}

\begin{theorem}
The image $\rho \left\{ \mathbb{C}\right\} $ of the representation $\rho ,$
which is a $c^{\ast }$-algebra isomorphism from $\mathbb{C}$ to $\rho
\left\{ \mathbb{C}\right\} \subset \mathcal{B}\left( \mathbb{R}^{2}\right) $
and is expressed by 
\begin{equation}
\rho \left\{ \mathbb{C}\right\} =CO^{+}\left( 2,\mathbb{R}\right) =\mathbb{R}%
_{+}\mathbb{\times }SO\left( 2,\mathbb{R}\right) \subset \mathcal{B}\left( 
\mathbb{R}^{2}\right) \cong \mathcal{M}\left( 2,2,\mathbb{R}\right) ,
\end{equation}%
where $CO^{+}\left( 2,\mathbb{R}\right) $ stands for the \emph{connected
component of the identity of the linear} \emph{conformal orthogonal group}
in dimension $2$.

\begin{proof}
\begin{equation}
\det \left( \rho \left( z\right) \right) =r>0
\end{equation}%
as 
\begin{equation}
\det \left\{ \left( 
\begin{array}{cc}
\cos \theta & -\sin \theta \\ 
\sin \theta & \cos \theta%
\end{array}%
\right) \right\} =\cos ^{2}\theta +\sin ^{2}\theta =+1.
\end{equation}
\end{proof}
\end{theorem}

\begin{corollary}
If $\left\Vert z\right\Vert =1$ $\Leftrightarrow r=1$ then%
\begin{equation}
\rho \left( z\right) \in SO\left( 2,\mathbb{R}\right)
\end{equation}%
and 
\begin{equation}
\rho \left\{ \mathbb{C}\right\} =SO\left( 2,\mathbb{R}\right) .
\end{equation}

\begin{proof}
If $\left\Vert z\right\Vert =1$ $\Leftrightarrow r=1$ then 
\begin{eqnarray}
\rho \left( z\right) &=&\left( 
\begin{array}{cc}
\cos \theta & -\sin \theta \\ 
\sin \theta & \cos \theta%
\end{array}%
\right) \\
\rho \left( z^{-1}\right) &=&\left( 
\begin{array}{cc}
\cos \theta & \sin \theta \\ 
-\sin \theta & \cos \theta%
\end{array}%
\right) =\rho \left( z\right) ^{t}=\rho \left( \overline{z}\right)
\end{eqnarray}%
so 
\begin{equation}
\rho \left( zz^{-1}\right) =\rho \left( z\overline{z}\right) =\mathbb{I}%
_{2}=\rho \left( z\right) \rho \left( z\right) ^{t}
\end{equation}
\end{proof}
\end{corollary}

hence 
\begin{equation}
\rho \left( z\right) \in SO\left( 2,\mathbb{R}\right) .\blacksquare
\end{equation}

\section{Bogoliubov Transformations}

\subsection{Matrices $\mathcal{M}\left( m,n,\mathcal{F}\right) $ with
Clifford Algebra $\mathcal{F}$ Elements}

\begin{definition}
\begin{equation}
\mathcal{M}\left( m,n,\mathcal{F}\right) =\left\{ \left. \left( 
\begin{array}{ccc}
C_{11} & \cdots & C_{1n} \\ 
\vdots & \ddots & \vdots \\ 
C_{m1} & \cdots & C_{mn}%
\end{array}%
\right) \right\vert C_{jk}\in \mathcal{F},\quad 1\leq j\leq m,\,1\leq k\leq
n\right\}
\end{equation}
\end{definition}

\begin{definition}
Clifford Algebra "$t$" Transpose $\equiv $ Reversion 
\begin{equation}
\left( x_{1}x_{2}\cdots x_{2r-1}x_{2r}\right) ^{t}=x_{2r}x_{2r-1}\cdots
x_{2}x_{1},
\end{equation}%
which gives 
\begin{equation}
\left( 
\begin{array}{ccc}
C_{11} & \cdots & C_{1n} \\ 
\vdots & \ddots & \vdots \\ 
C_{m1} & \cdots & C_{mn}%
\end{array}%
\right) ^{t}=\left( 
\begin{array}{ccc}
C_{11}^{t} & \cdots & C_{m1}^{t} \\ 
\vdots & \ddots & \vdots \\ 
C_{1n}^{t} & \cdots & C_{mn}^{t}%
\end{array}%
\right) .
\end{equation}
\end{definition}

\subsection{Generators of $\mathcal{F}_{1}$}

\begin{remark}
The action of the transpose on generators $\left\{ \left.
x_{j},a_{j},a_{j}^{\dagger },c_{j}\right\vert 1\leq j\leq 2r\right\} $ of $%
\mathcal{F}_{1}$

\begin{enumerate}
\item Self adjoint generators 
\begin{equation}
x_{j}^{t}=x_{j},\,1\leq j\leq 2r.
\end{equation}

\item Non self adjoint generators 
\begin{equation}
a_{j}^{t}=a_{j},\,\left( a_{j}^{\dagger }\right) ^{t}=a_{j}^{\dagger
},\,1\leq j\leq r.
\end{equation}

\item General generators 
\begin{equation}
c_{j}^{t}=c_{j},\,1\leq j\leq 2r.
\end{equation}
\end{enumerate}
\end{remark}

\begin{remark}
If $c_{j}\in \mathcal{F}_{1},\,1\leq j\leq 2r$ then 
\begin{eqnarray}
\left( 
\begin{array}{c}
c_{1} \\ 
\vdots \\ 
c_{2r}%
\end{array}%
\right) ^{t} &=&\left( 
\begin{array}{ccc}
c_{1} & \cdots & c_{2r}%
\end{array}%
\right) \\
\left( 
\begin{array}{ccc}
c_{1} & \cdots & c_{2r}%
\end{array}%
\right) ^{t} &=&\left( 
\begin{array}{c}
c_{1} \\ 
\vdots \\ 
c_{2r}%
\end{array}%
\right) .
\end{eqnarray}
\end{remark}

\begin{remark}
\begin{enumerate}
\item In partitioned form 
\begin{equation}
\left( \boldsymbol{c}\right) ^{t}=\left( 
\begin{array}{cc}
\left( \boldsymbol{c}_{1}\right) ^{t} & \left( \boldsymbol{c}_{2}\right) ^{t}%
\end{array}%
\right) ,
\end{equation}%
where 
\begin{eqnarray}
\left( \boldsymbol{c}_{1}\right) ^{t} &=&\left( 
\begin{array}{ccc}
c_{1} & \cdots & c_{r}%
\end{array}%
\right) \\
\left( \boldsymbol{c}_{2}\right) ^{t} &=&\left( 
\begin{array}{ccc}
c_{r+1} & \cdots & c_{2r}%
\end{array}%
\right) .
\end{eqnarray}

\item 
\begin{equation}
\left( \boldsymbol{c}\right) \left( \boldsymbol{c}\right) ^{t}=\left( 
\begin{array}{c}
c_{1} \\ 
\vdots \\ 
c_{2r}%
\end{array}%
\right) \left( 
\begin{array}{ccc}
c_{1} & \cdots & c_{2r}%
\end{array}%
\right) =\left( 
\begin{array}{ccc}
c_{1}c_{1} & \cdots & c_{1}c_{2r} \\ 
\vdots & \ddots & \vdots \\ 
c_{2r}c_{1} & \cdots & c_{2r}c_{2r}%
\end{array}%
\right) =\left( 
\begin{array}{c}
\left( \boldsymbol{c}_{1}\right) ^{t} \\ 
\left( \boldsymbol{c}_{2}\right) ^{t}%
\end{array}%
\right) \left( 
\begin{array}{cc}
\boldsymbol{c}_{1} & \boldsymbol{c}_{2}%
\end{array}%
\right) =\left( 
\begin{array}{cc}
\left( \boldsymbol{c}_{1}\right) ^{t}\left( \boldsymbol{c}_{1}\right) & 
\left( \boldsymbol{c}_{1}\right) ^{t}\left( \boldsymbol{c}_{2}\right) \\ 
\left( \boldsymbol{c}_{2}\right) ^{t}\left( \boldsymbol{c}_{1}\right) & 
\left( \boldsymbol{c}_{2}\right) ^{t}\left( \boldsymbol{c}_{2}\right)%
\end{array}%
\right) .
\end{equation}

\item 
\begin{equation}
\left( \boldsymbol{c}\right) ^{t}\left( \boldsymbol{c}\right) =\left( 
\begin{array}{ccc}
c_{1} & \cdots & c_{2r}%
\end{array}%
\right) \left( 
\begin{array}{c}
c_{1} \\ 
\vdots \\ 
c_{2r}%
\end{array}%
\right) =\left( \boldsymbol{c}_{1}\right) ^{t}\left( \boldsymbol{c}%
_{1}\right) +\left( \boldsymbol{c}_{2}\right) ^{t}\left( \boldsymbol{c}%
_{2}\right) =\sum_{1\leq j\leq 2r}c_{j}c_{j}.
\end{equation}
\end{enumerate}
\end{remark}

\begin{definition}
Representations $\mathcal{R}_{c}\left( T\right) $ of $T\in GL(2r,\mathcal{F}%
_{1}\mathbb{)}\cong GL(2r,\mathbb{C)}$ then 
\begin{equation}
T\left( 
\begin{array}{ccc}
c_{1} & \cdots & c_{2r}%
\end{array}%
\right) =\left( \boldsymbol{c}\right) ^{t}\mathcal{R}_{c}\left( T\right)
=\left( 
\begin{array}{ccc}
c_{1} & \cdots & c_{2r}%
\end{array}%
\right) \mathcal{R}_{c}\left( T\right) =\left( 
\begin{array}{cc}
\boldsymbol{c}_{1} & \boldsymbol{c}_{2}%
\end{array}%
\right) \mathcal{R}_{c}\left( T\right) =\left( 
\begin{array}{cccccc}
c_{1} & \cdots & c_{r} & c_{r+1} & \cdots & c_{2r}%
\end{array}%
\right) \left( 
\begin{array}{cc}
\mathcal{R}_{c}\left( T\right) _{11} & \mathcal{R}_{c}\left( T\right) _{12}
\\ 
\mathcal{R}_{c}\left( T\right) _{21} & \mathcal{R}_{c}\left( T\right) _{22}%
\end{array}%
\right) .
\end{equation}
\end{definition}

\subsection{Self Adjoint Generators of $\mathcal{F}_{1}$}

\begin{remark}
Any operator $A$ can be expressed as a real sum of self adjoint operators 
\begin{equation}
\func{Self}\left( A\right) =\frac{1}{2}\left( A+A^{\dagger }\right)
\end{equation}%
and anti self adjoint operators%
\begin{equation}
i\func{ASelf}\left( A\right) =\frac{1}{2}\left( A-A^{\dagger }\right) ,
\end{equation}%
i.e. 
\begin{equation*}
A=\frac{1}{\sqrt{2}}\left( A+A^{\dagger }\right) +\frac{1}{\sqrt{2}}\left(
A-A^{\dagger }\right) =\func{Self}\left( A\right) +i\func{ASelf}\left(
A\right) ,
\end{equation*}%
which can be displayed by the transformation 
\begin{equation}
\left( 
\begin{array}{cc}
A^{\dagger } & A%
\end{array}%
\right) =\left( 
\begin{array}{cc}
\frac{1}{\sqrt{2}}\left( A^{\dagger }+A\right) & \frac{i}{\sqrt{2}}\left(
A^{\dagger }-A\right)%
\end{array}%
\right) \frac{1}{\sqrt{2}}\left( 
\begin{array}{cc}
1 & 1 \\ 
-i & i%
\end{array}%
\right) =\left( 
\begin{array}{cc}
\frac{1}{2}\left( A^{\dagger }+A\right) +\frac{1}{2}\left( A^{\dagger
}-A\right) & \frac{1}{2}\left( A+A^{\dagger }\right) -\frac{1}{2}\left(
A^{\dagger }-A\right)%
\end{array}%
\right) .
\end{equation}%
The inverse relationship is given by 
\begin{equation}
\left( 
\begin{array}{cc}
\frac{1}{\sqrt{2}}\left( A^{\dagger }+A\right) & \frac{i}{\sqrt{2}}\left(
A^{\dagger }-A\right)%
\end{array}%
\right) =\left( 
\begin{array}{cc}
A^{\dagger } & A%
\end{array}%
\right) \frac{1}{\sqrt{2}}\left( 
\begin{array}{cc}
1 & i \\ 
1 & -i%
\end{array}%
\right) .
\end{equation}%
If we define 
\begin{eqnarray}
X_{1} &=&\frac{1}{\sqrt{2}}\left( A^{\dagger }+A\right) \\
X_{2} &=&\frac{i}{\sqrt{2}}\left( A^{\dagger }-A\right)
\end{eqnarray}%
the transformation to self adjoint parts can be written as%
\begin{equation}
\left( 
\begin{array}{cc}
X_{1} & X_{2}%
\end{array}%
\right) =\left( 
\begin{array}{cc}
A^{\dagger } & A%
\end{array}%
\right) \frac{1}{\sqrt{2}}\left( 
\begin{array}{cc}
1 & i \\ 
1 & -i%
\end{array}%
\right)
\end{equation}%
and the inverse relationship of self adjoint parts back to the non self
adjoint operator as 
\begin{equation}
\left( 
\begin{array}{cc}
A^{\dagger } & A%
\end{array}%
\right) =\left( 
\begin{array}{cc}
X_{1} & X_{2}%
\end{array}%
\right) \frac{1}{\sqrt{2}}\left( 
\begin{array}{cc}
1 & i \\ 
i & i%
\end{array}%
\right) .
\end{equation}
\end{remark}

\begin{lemma}
The order one subspace $\mathcal{F}_{1}$ of the Clifford algebra $\mathcal{F}
$ can be generated by a non self adjoint basis 
\begin{equation}
\left( 
\begin{array}{cc}
\left( \boldsymbol{a}^{\dagger }\right) ^{t} & \left( \boldsymbol{a}\right)
^{t}%
\end{array}%
\right) =\left( 
\begin{array}{cccccc}
a_{1}^{\dagger } & \cdots & a_{r}^{\dagger } & a_{1} & \cdots & a_{r}%
\end{array}%
\right)
\end{equation}%
and its associated self adjoint basis 
\begin{equation}
\left( \boldsymbol{x}\right) ^{t}=\left( 
\begin{array}{cc}
\left( \boldsymbol{x}_{1}\right) ^{t} & \left( \boldsymbol{x}_{2}\right) ^{t}%
\end{array}%
\right) =\left( 
\begin{array}{cccccc}
x_{1} & \cdots & x_{r} & x_{r+1} & \cdots & x_{2r}%
\end{array}%
\right) ,
\end{equation}%
where these bases are related by 
\begin{gather}
\left( 
\begin{array}{cc}
\left( \boldsymbol{x}_{1}\right) ^{t} & \left( \boldsymbol{x}_{2}\right) ^{t}%
\end{array}%
\right) \\
\parallel \\
u_{xa}\left( 
\begin{array}{cc}
\left( \boldsymbol{a}^{\dagger }\right) ^{t} & \left( \boldsymbol{a}\right)
^{t}%
\end{array}%
\right) \\
\parallel \\
\left( 
\begin{array}{cc}
\left( \boldsymbol{a}^{\dagger }\right) ^{t} & \left( \boldsymbol{a}\right)
^{t}%
\end{array}%
\right) \frac{1}{\sqrt{2}}\left( 
\begin{array}{cc}
\mathbb{I}_{r} & i\mathbb{I}_{r} \\ 
\mathbb{I}_{r} & -i\mathbb{I}_{r}%
\end{array}%
\right) \\
\parallel \\
\frac{1}{\sqrt{2}}\left( 
\begin{array}{cc}
\left( \overset{}{\left( \boldsymbol{a}^{\dagger }\right) +\left( 
\boldsymbol{a}\right) }\right) ^{t} & i\left( \overset{}{\left( \boldsymbol{a%
}^{\dagger }\right) -\left( \boldsymbol{a}\right) }\right) ^{t}%
\end{array}%
\right) .
\end{gather}
\end{lemma}

\begin{remark}
The transformations $u_{xa}$ from a non sel adjoint basis to a self adjoint
basis 
\begin{equation}
u_{xa}:\left( \overset{}{\left( \boldsymbol{a}^{\dagger }\right) ^{t},\left( 
\boldsymbol{a}\right) ^{t}}\right) \rightarrow \left( \overset{}{\left( 
\boldsymbol{x}_{1}\right) ^{t},\left( \boldsymbol{x}_{2}\right) ^{t}}\right)
\end{equation}%
are given by 
\begin{equation}
\frac{1}{\sqrt{2}}\left( 
\begin{array}{cc}
\mathbb{I}_{r} & i\mathbb{I}_{r} \\ 
\mathbb{I}_{r} & -i\mathbb{I}_{r}%
\end{array}%
\right) =\mathcal{R}_{\left( a^{\dagger },a\right) }\left( u_{xa}\right) \in
U\left( 2r\right) ,
\end{equation}%
where 
\begin{equation}
u_{xa}\left( 
\begin{array}{cc}
\left( \boldsymbol{a}^{\dagger }\right) ^{t} & \left( \boldsymbol{a}\right)
^{t}%
\end{array}%
\right) =\left( 
\begin{array}{cc}
\left( \boldsymbol{a}^{\dagger }\right) ^{t} & \left( \boldsymbol{a}\right)
^{t}%
\end{array}%
\right) \mathcal{R}_{\left( a^{\dagger },a\right) }\left( u_{xa}\right) .
\end{equation}
\end{remark}

\begin{remark}
Let $T\in GL\left( 2r,\mathcal{F}_{1}\right) $ then

\begin{enumerate}
\item 
\begin{equation}
T\left( 
\begin{array}{cc}
\left( \boldsymbol{a}^{\dagger }\right) ^{t} & \left( \boldsymbol{a}\right)
^{t}%
\end{array}%
\right) =\left( 
\begin{array}{cc}
\left( \boldsymbol{a}^{\dagger }\right) ^{t} & \left( \boldsymbol{a}\right)
^{t}%
\end{array}%
\right) \mathcal{R}_{\left( a^{\dagger },a\right) }\left( T\right) =\left( 
\begin{array}{cc}
\left( \boldsymbol{a}^{\dagger }\right) ^{t} & \left( \boldsymbol{a}\right)
^{t}%
\end{array}%
\right) \left( 
\begin{array}{cc}
\mathcal{R}_{\left( a^{\dagger },a\right) }\left( T\right) _{11} & \mathcal{R%
}_{\left( a^{\dagger },a\right) }\left( T\right) _{12} \\ 
\mathcal{R}_{\left( a^{\dagger },a\right) }\left( T\right) _{21} & \mathcal{R%
}_{\left( a^{\dagger },a\right) }\left( T\right) _{22}%
\end{array}%
\right)
\end{equation}

\item 
\begin{equation}
T\left( 
\begin{array}{cc}
\left( \boldsymbol{x}_{1}\right) ^{t} & \left( \boldsymbol{x}_{2}\right) ^{t}%
\end{array}%
\right) =\left( 
\begin{array}{cc}
\left( \boldsymbol{x}_{1}\right) ^{t} & \left( \boldsymbol{x}_{2}\right) ^{t}%
\end{array}%
\right) \mathcal{R}_{\left( x\right) }\left( T\right) =\left( 
\begin{array}{cc}
\left( \boldsymbol{x}_{1}\right) ^{t} & \left( \boldsymbol{x}_{2}\right) ^{t}%
\end{array}%
\right) \left( 
\begin{array}{cc}
\mathcal{R}_{\left( x\right) }\left( T\right) _{11} & \mathcal{R}_{\left(
x\right) }\left( T\right) _{12} \\ 
\mathcal{R}_{\left( x\right) }\left( T\right) _{21} & \mathcal{R}_{\left(
x\right) }\left( T\right) _{22}%
\end{array}%
\right)
\end{equation}

\begin{enumerate}
\item The self adjoint basis and non sel adjoint are given by 
\begin{eqnarray}
\left( 
\begin{array}{cc}
\left( \boldsymbol{x}_{1}\right) ^{t} & \left( \boldsymbol{x}_{2}\right) ^{t}%
\end{array}%
\right) &=&\left( 
\begin{array}{cc}
\left( \boldsymbol{a}^{\dagger }\right) ^{t} & \left( \boldsymbol{a}\right)
^{t}%
\end{array}%
\right) \frac{1}{\sqrt{2}}\left( 
\begin{array}{cc}
\mathbb{I}_{r} & i\mathbb{I}_{r} \\ 
\mathbb{I}_{r} & -i\mathbb{I}_{r}%
\end{array}%
\right) \\
\left( 
\begin{array}{cc}
\left( \boldsymbol{x}_{1}\right) ^{t} & \left( \boldsymbol{x}_{2}\right) ^{t}%
\end{array}%
\right) \frac{1}{\sqrt{2}}\left( 
\begin{array}{cc}
\mathbb{I}_{r} & i\mathbb{I}_{r} \\ 
\mathbb{I}_{r} & -i\mathbb{I}_{r}%
\end{array}%
\right) ^{\dagger } &=&\left( 
\begin{array}{cc}
\left( \boldsymbol{a}^{\dagger }\right) ^{t} & \left( \boldsymbol{a}\right)
^{t}%
\end{array}%
\right)
\end{eqnarray}

\item Representations of operators $T\in \mathcal{L}\left( \mathcal{F}%
_{1}\right) $ are given by 
\begin{eqnarray}
T\left( 
\begin{array}{cc}
\left( \boldsymbol{x}_{1}\right) ^{t} & \left( \boldsymbol{x}_{2}\right) ^{t}%
\end{array}%
\right) &=&T\left( 
\begin{array}{cc}
\left( \boldsymbol{a}^{\dagger }\right) ^{t} & \left( \boldsymbol{a}\right)
^{t}%
\end{array}%
\right) \frac{1}{\sqrt{2}}\left( 
\begin{array}{cc}
\mathbb{I}_{r} & i\mathbb{I}_{r} \\ 
\mathbb{I}_{r} & -i\mathbb{I}_{r}%
\end{array}%
\right) \\
&=&\left( 
\begin{array}{cc}
\left( \boldsymbol{a}^{\dagger }\right) ^{t} & \left( \boldsymbol{a}\right)
^{t}%
\end{array}%
\right) \frac{1}{\sqrt{2}}\left( 
\begin{array}{cc}
\mathbb{I}_{r} & i\mathbb{I}_{r} \\ 
\mathbb{I}_{r} & -i\mathbb{I}_{r}%
\end{array}%
\right) \mathcal{R}_{\left( x\right) }\left( T\right) \\
&=&\left( 
\begin{array}{cc}
\left( \boldsymbol{a}^{\dagger }\right) ^{t} & \left( \boldsymbol{a}\right)
^{t}%
\end{array}%
\right) \frac{1}{\sqrt{2}}\left( 
\begin{array}{cc}
\mathbb{I}_{r} & i\mathbb{I}_{r} \\ 
\mathbb{I}_{r} & -i\mathbb{I}_{r}%
\end{array}%
\right) \left( 
\begin{array}{cc}
\mathcal{R}_{\left( x\right) }\left( T\right) _{11} & \mathcal{R}_{\left(
x\right) }\left( T\right) _{12} \\ 
\mathcal{R}_{\left( x\right) }\left( T\right) _{21} & \mathcal{R}_{\left(
x\right) }\left( T\right) _{22}%
\end{array}%
\right) ,
\end{eqnarray}%
which leads to 
\begin{align}
T\left( 
\begin{array}{cc}
\left( \boldsymbol{x}_{1}\right) ^{t} & \left( \boldsymbol{x}_{2}\right) ^{t}%
\end{array}%
\right) \frac{1}{\sqrt{2}}\left( 
\begin{array}{cc}
\mathbb{I}_{r} & i\mathbb{I}_{r} \\ 
\mathbb{I}_{r} & -i\mathbb{I}_{r}%
\end{array}%
\right) ^{\dagger }& =T\left( 
\begin{array}{cc}
\left( \boldsymbol{a}^{\dagger }\right) ^{t} & \left( \boldsymbol{a}\right)
^{t}%
\end{array}%
\right) \frac{1}{\sqrt{2}}\left( 
\begin{array}{cc}
\mathbb{I}_{r} & i\mathbb{I}_{r} \\ 
\mathbb{I}_{r} & -i\mathbb{I}_{r}%
\end{array}%
\right) \frac{1}{\sqrt{2}}\left( 
\begin{array}{cc}
\mathbb{I}_{r} & i\mathbb{I}_{r} \\ 
\mathbb{I}_{r} & -i\mathbb{I}_{r}%
\end{array}%
\right) ^{\dagger } \\
& =T\left( 
\begin{array}{cc}
\left( \boldsymbol{a}^{\dagger }\right) ^{t} & \left( \boldsymbol{a}\right)
^{t}%
\end{array}%
\right) =\left( 
\begin{array}{cc}
\left( \boldsymbol{a}^{\dagger }\right) ^{t} & \left( \boldsymbol{a}\right)
^{t}%
\end{array}%
\right) \mathcal{R}_{\left( a^{\dagger },a\right) }\left( T\right) \\
& =\left( 
\begin{array}{cc}
\left( \boldsymbol{a}^{\dagger }\right) ^{t} & \left( \boldsymbol{a}\right)
^{t}%
\end{array}%
\right) \frac{1}{\sqrt{2}}\left( 
\begin{array}{cc}
\mathbb{I}_{r} & i\mathbb{I}_{r} \\ 
\mathbb{I}_{r} & -i\mathbb{I}_{r}%
\end{array}%
\right) \mathcal{R}_{\left( x\right) }\left( T\right) \frac{1}{\sqrt{2}}%
\left( 
\begin{array}{cc}
\mathbb{I}_{r} & i\mathbb{I}_{r} \\ 
\mathbb{I}_{r} & -i\mathbb{I}_{r}%
\end{array}%
\right) ^{\dagger } \\
& =\left( 
\begin{array}{cc}
\left( \boldsymbol{a}^{\dagger }\right) ^{t} & \left( \boldsymbol{a}\right)
^{t}%
\end{array}%
\right) \frac{1}{\sqrt{2}}\left( 
\begin{array}{cc}
\mathbb{I}_{r} & i\mathbb{I}_{r} \\ 
\mathbb{I}_{r} & -i\mathbb{I}_{r}%
\end{array}%
\right) \left( 
\begin{array}{cc}
\mathcal{R}_{\left( x\right) }\left( T\right) _{11} & \mathcal{R}_{\left(
x\right) }\left( T\right) _{12} \\ 
\mathcal{R}_{\left( x\right) }\left( T\right) _{21} & \mathcal{R}_{\left(
x\right) }\left( T\right) _{22}%
\end{array}%
\right) \frac{1}{\sqrt{2}}\left( 
\begin{array}{cc}
\mathbb{I}_{r} & i\mathbb{I}_{r} \\ 
\mathbb{I}_{r} & -i\mathbb{I}_{r}%
\end{array}%
\right) ^{\dagger }.
\end{align}%
Hence the $\mathcal{R}_{\left( a^{\dagger },a\right) }$ representation is
related to the representation $\mathcal{R}_{\left( x\right) }$ by 
\begin{equation}
\mathcal{R}_{\left( a^{\dagger },a\right) }\left( T\right) =\frac{1}{\sqrt{2}%
}\left( 
\begin{array}{cc}
\mathbb{I}_{r} & i\mathbb{I}_{r} \\ 
\mathbb{I}_{r} & -i\mathbb{I}_{r}%
\end{array}%
\right) \left( 
\begin{array}{cc}
\mathcal{R}_{\left( x\right) }\left( T\right) _{11} & \mathcal{R}_{\left(
x\right) }\left( T\right) _{12} \\ 
\mathcal{R}_{\left( x\right) }\left( T\right) _{21} & \mathcal{R}_{\left(
x\right) }\left( T\right) _{22}%
\end{array}%
\right) \frac{1}{\sqrt{2}}\left( 
\begin{array}{cc}
\mathbb{I}_{r} & i\mathbb{I}_{r} \\ 
\mathbb{I}_{r} & -i\mathbb{I}_{r}%
\end{array}%
\right) ^{\dagger },
\end{equation}%
which gives\ 
\begin{gather}
\mathcal{R}_{\left( a^{\dagger },a\right) }\left( T\right) =\frac{1}{\sqrt{2}%
}\left( 
\begin{array}{cc}
\mathbb{I}_{r} & i\mathbb{I}_{r} \\ 
\mathbb{I}_{r} & -i\mathbb{I}_{r}%
\end{array}%
\right) \left( 
\begin{array}{cc}
\mathcal{R}_{\left( x\right) }\left( T\right) _{11} & \mathcal{R}_{\left(
x\right) }\left( T\right) _{12} \\ 
\mathcal{R}_{\left( x\right) }\left( T\right) _{21} & \mathcal{R}_{\left(
x\right) }\left( T\right) _{22}%
\end{array}%
\right) \frac{1}{\sqrt{2}}\left( 
\begin{array}{cc}
\mathbb{I}_{r} & \mathbb{I}_{r} \\ 
-i\mathbb{I}_{r} & i\mathbb{I}_{r}%
\end{array}%
\right) \\
=\frac{1}{2}\left( 
\begin{array}{cc}
\mathbb{I}_{r} & i\mathbb{I}_{r} \\ 
\mathbb{I}_{r} & -i\mathbb{I}_{r}%
\end{array}%
\right) \left( 
\begin{array}{cc}
\mathcal{R}_{\left( x\right) }\left( T\right) _{11}-i\mathcal{R}_{\left(
x\right) }\left( T\right) _{12} & \mathcal{R}_{\left( x\right) }\left(
T\right) _{11}+i\mathcal{R}_{\left( x\right) }\left( T\right) _{12} \\ 
\mathcal{R}_{\left( x\right) }\left( T\right) _{21}-i\mathcal{R}_{\left(
x\right) }\left( T\right) _{22} & \mathcal{R}_{\left( x\right) }\left(
T\right) _{21}+i\mathcal{R}_{\left( x\right) }\left( T\right) _{22}%
\end{array}%
\right) \\
=\frac{1}{2}\left( 
\begin{array}{cc}
\mathcal{R}_{\left( x\right) }\left( T\right) _{11}-i\mathcal{R}_{\left(
x\right) }\left( T\right) _{12}+i\left( \mathcal{R}_{\left( x\right) }\left(
T\right) _{21}-i\mathcal{R}_{\left( x\right) }\left( T\right) _{22}\right) & 
\mathcal{R}_{\left( x\right) }\left( T\right) _{11}+i\mathcal{R}_{\left(
x\right) }\left( T\right) _{12}+i\left( \mathcal{R}_{\left( x\right) }\left(
T\right) _{21}+i\mathcal{R}_{\left( x\right) }\left( T\right) _{22}\right)
\\ 
\mathcal{R}_{\left( x\right) }\left( T\right) _{11}-i\mathcal{R}_{\left(
x\right) }\left( T\right) _{12}-i\left( \mathcal{R}_{\left( x\right) }\left(
T\right) _{21}-i\mathcal{R}_{\left( x\right) }\left( T\right) _{22}\right) & 
\mathcal{R}_{\left( x\right) }\left( T\right) _{11}+i\mathcal{R}_{\left(
x\right) }\left( T\right) _{12}-i\left( \mathcal{R}_{\left( x\right) }\left(
T\right) _{21}+i\mathcal{R}_{\left( x\right) }\left( T\right) _{22}\right)%
\end{array}%
\right) \\
=\frac{1}{2}\left( 
\begin{array}{cc}
\mathcal{R}_{\left( x\right) }\left( T\right) _{11}+\mathcal{R}_{\left(
x\right) }\left( T\right) _{22}-i\left( \mathcal{R}_{\left( x\right) }\left(
T\right) _{12}-\mathcal{R}_{\left( x\right) }\left( T\right) _{21}\right) & 
\mathcal{R}_{\left( x\right) }\left( T\right) _{11}-\mathcal{R}_{\left(
x\right) }\left( T\right) _{22}+i\left( \mathcal{R}_{\left( x\right) }\left(
T\right) _{12}+\mathcal{R}_{\left( x\right) }\left( T\right) _{21}\right) \\ 
\mathcal{R}_{\left( x\right) }\left( T\right) _{11}-\mathcal{R}_{\left(
x\right) }\left( T\right) _{22}-i\left( \mathcal{R}_{\left( x\right) }\left(
T\right) _{12}+\mathcal{R}_{\left( x\right) }\left( T\right) _{21}\right) & 
\mathcal{R}_{\left( x\right) }\left( T\right) _{11}+\mathcal{R}_{\left(
x\right) }\left( T\right) _{22}+i\left( \mathcal{R}_{\left( x\right) }\left(
T\right) _{12}-\mathcal{R}_{\left( x\right) }\left( T\right) _{21}\right)%
\end{array}%
\right) ,
\end{gather}%
i.e. $\mathcal{R}_{\left( a^{\dagger },a\right) }\left( T\right) $ in terms
of $\mathcal{R}_{\left( x\right) }\left( T\right) $ is 
\begin{gather}
\left( 
\begin{array}{cc}
\mathcal{R}_{\left( a^{\dagger },a\right) }\left( T\right) _{11} & \mathcal{R%
}_{\left( a^{\dagger },a\right) }\left( T\right) _{12} \\ 
\mathcal{R}_{\left( a^{\dagger },a\right) }\left( T\right) _{21} & \mathcal{R%
}_{\left( a^{\dagger },a\right) }\left( T\right) _{22}%
\end{array}%
\right) \\
\parallel \\
\frac{1}{2}\left( 
\begin{array}{cc}
\mathcal{R}_{\left( x\right) }\left( T\right) _{11}+\mathcal{R}_{\left(
x\right) }\left( T\right) _{22}-i\left( \overset{}{\mathcal{R}_{\left(
x\right) }\left( T\right) _{12}-\mathcal{R}_{\left( x\right) }\left(
T\right) _{21}}\right) & \mathcal{R}_{\left( x\right) }\left( T\right) _{11}-%
\mathcal{R}_{\left( x\right) }\left( T\right) _{22}+i\left( \overset{}{%
\mathcal{R}_{\left( x\right) }\left( T\right) _{12}+\mathcal{R}_{\left(
x\right) }\left( T\right) _{21}}\right) \\ 
\mathcal{R}_{\left( x\right) }\left( T\right) _{11}-\mathcal{R}_{\left(
x\right) }\left( T\right) _{22}-i\left( \overset{}{\mathcal{R}_{\left(
x\right) }\left( T\right) _{12}+\mathcal{R}_{\left( x\right) }\left(
T\right) _{21}}\right) & \mathcal{R}_{\left( x\right) }\left( T\right) _{11}+%
\mathcal{R}_{\left( x\right) }\left( T\right) _{22}+i\left( \overset{}{%
\mathcal{R}_{\left( x\right) }\left( T\right) _{12}-\mathcal{R}_{\left(
x\right) }\left( T\right) _{21}}\right)%
\end{array}%
\right) .
\end{gather}%
Thye inverse relationship i.e. $\mathcal{R}_{\left( x\right) }\left(
T\right) $ in terms of is $\mathcal{R}_{\left( a^{\dagger },a\right) }\left(
T\right) $ 
\begin{equation}
\frac{1}{\sqrt{2}}\left( 
\begin{array}{cc}
\mathbb{I}_{r} & i\mathbb{I}_{r} \\ 
\mathbb{I}_{r} & -i\mathbb{I}_{r}%
\end{array}%
\right) ^{\dagger }\left( 
\begin{array}{cc}
\mathcal{R}_{\left( a^{\dagger },a\right) }\left( T\right) _{11} & \mathcal{R%
}_{\left( a^{\dagger },a\right) }\left( T\right) _{12} \\ 
\mathcal{R}_{\left( a^{\dagger },a\right) }\left( T\right) _{21} & \mathcal{R%
}_{\left( a^{\dagger },a\right) }\left( T\right) _{22}%
\end{array}%
\right) \frac{1}{\sqrt{2}}\left( 
\begin{array}{cc}
\mathbb{I}_{r} & i\mathbb{I}_{r} \\ 
\mathbb{I}_{r} & -i\mathbb{I}_{r}%
\end{array}%
\right) =\left( 
\begin{array}{cc}
\mathcal{R}_{\left( x\right) }\left( T\right) _{11} & \mathcal{R}_{\left(
x\right) }\left( T\right) _{12} \\ 
\mathcal{R}_{\left( x\right) }\left( T\right) _{21} & \mathcal{R}_{\left(
x\right) }\left( T\right) _{22}%
\end{array}%
\right)
\end{equation}%
\begin{equation}
=\frac{1}{2}\left( 
\begin{array}{cc}
\mathcal{R}_{\left( a^{\dagger },a\right) }\left( T\right) _{11}+\mathcal{R}%
_{\left( a^{\dagger },a\right) }\left( T\right) _{12}+\mathcal{R}_{\left(
a^{\dagger },a\right) }\left( T\right) _{21}+\mathcal{R}_{\left( a^{\dagger
},a\right) }\left( T\right) _{22} & i\left( \overset{}{\mathcal{R}_{\left(
a^{\dagger },a\right) }\left( T\right) _{11}-\mathcal{R}_{\left( a^{\dagger
},a\right) }\left( T\right) _{12}+\mathcal{R}_{\left( a^{\dagger },a\right)
}\left( T\right) _{21}-\mathcal{R}_{\left( a^{\dagger },a\right) }\left(
T\right) _{22}}\right) \\ 
i\left( \overset{}{\mathcal{R}_{\left( a^{\dagger },a\right) }\left(
T\right) _{11}+\mathcal{R}_{\left( a^{\dagger },a\right) }\left( T\right)
_{12}-\mathcal{R}_{\left( a^{\dagger },a\right) }\left( T\right) _{21}-%
\mathcal{R}_{\left( a^{\dagger },a\right) }\left( T\right) _{22}}\right) & 
\mathcal{R}_{\left( a^{\dagger },a\right) }\left( T\right) _{12}+\mathcal{R}%
_{\left( a^{\dagger },a\right) }\left( T\right) _{21}-\mathcal{R}_{\left(
a^{\dagger },a\right) }\left( T\right) _{11}-\mathcal{R}_{\left( a^{\dagger
},a\right) }\left( T\right) _{22}%
\end{array}%
\right)
\end{equation}%
leading to 
\begin{gather}
\left( 
\begin{array}{cc}
\mathcal{R}_{\left( x\right) }\left( T\right) _{11} & \mathcal{R}_{\left(
x\right) }\left( T\right) _{12} \\ 
\mathcal{R}_{\left( x\right) }\left( T\right) _{21} & \mathcal{R}_{\left(
x\right) }\left( T\right) _{22}%
\end{array}%
\right) \\
\parallel \\
\frac{1}{2}\left( 
\begin{array}{cc}
\mathcal{R}_{\left( a^{\dagger },a\right) }\left( T\right) _{11}+\mathcal{R}%
_{\left( a^{\dagger },a\right) }\left( T\right) _{12}+\mathcal{R}_{\left(
a^{\dagger },a\right) }\left( T\right) _{21}+\mathcal{R}_{\left( a^{\dagger
},a\right) }\left( T\right) _{22} & i\left( \overset{}{\mathcal{R}_{\left(
a^{\dagger },a\right) }\left( T\right) _{11}-\mathcal{R}_{\left( a^{\dagger
},a\right) }\left( T\right) _{12}+\mathcal{R}_{\left( a^{\dagger },a\right)
}\left( T\right) _{21}-\mathcal{R}_{\left( a^{\dagger },a\right) }\left(
T\right) _{22}}\right) \\ 
i\left( \overset{}{\mathcal{R}_{\left( a^{\dagger },a\right) }\left(
T\right) _{11}+\mathcal{R}_{\left( a^{\dagger },a\right) }\left( T\right)
_{12}-\mathcal{R}_{\left( a^{\dagger },a\right) }\left( T\right) _{21}-%
\mathcal{R}_{\left( a^{\dagger },a\right) }\left( T\right) _{22}}\right) & 
\mathcal{R}_{\left( a^{\dagger },a\right) }\left( T\right) _{12}+\mathcal{R}%
_{\left( a^{\dagger },a\right) }\left( T\right) _{21}-\mathcal{R}_{\left(
a^{\dagger },a\right) }\left( T\right) _{11}-\mathcal{R}_{\left( a^{\dagger
},a\right) }\left( T\right) _{22}%
\end{array}%
\right) .
\end{gather}
\end{enumerate}

\item If $\mathcal{R}_{\left( x\right) }\left( T\right) \in GL\left( 2r,%
\mathbb{R}\right) $ then it must be a real matrix as it acts on real linear
combinations of a sel adjoint basis hence 
\begin{equation}
\mathcal{R}_{\left( a^{\dagger },a\right) }\left( T\right) _{11}+\mathcal{R}%
_{\left( a^{\dagger },a\right) }\left( T\right) _{12}+\mathcal{R}_{\left(
a^{\dagger },a\right) }\left( T\right) _{21}+\mathcal{R}_{\left( a^{\dagger
},a\right) }\left( T\right) _{22}=\overline{\mathcal{R}_{\left( a^{\dagger
},a\right) }\left( T\right) _{11}}+\overline{\mathcal{R}_{\left( a^{\dagger
},a\right) }\left( T\right) _{12}}+\overline{\mathcal{R}_{\left( a^{\dagger
},a\right) }\left( T\right) _{21}}+\overline{\mathcal{R}_{\left( a^{\dagger
},a\right) }\left( T\right) _{22}}
\end{equation}%
\begin{equation}
\mathcal{R}_{\left( a^{\dagger },a\right) }\left( T\right) _{11}-\mathcal{R}%
_{\left( a^{\dagger },a\right) }\left( T\right) _{12}+\mathcal{R}_{\left(
a^{\dagger },a\right) }\left( T\right) _{21}-\mathcal{R}_{\left( a^{\dagger
},a\right) }\left( T\right) _{22}=\overline{\mathcal{R}_{\left( a^{\dagger
},a\right) }\left( T\right) _{11}}-\overline{\mathcal{R}_{\left( a^{\dagger
},a\right) }\left( T\right) _{12}}+\overline{\mathcal{R}_{\left( a^{\dagger
},a\right) }\left( T\right) _{21}}-\overline{\mathcal{R}_{\left( a^{\dagger
},a\right) }\left( T\right) _{22}}
\end{equation}%
\begin{equation}
\mathcal{R}_{\left( a^{\dagger },a\right) }\left( T\right) _{11}+\mathcal{R}%
_{\left( a^{\dagger },a\right) }\left( T\right) _{12}-\mathcal{R}_{\left(
a^{\dagger },a\right) }\left( T\right) _{21}-\mathcal{R}_{\left( a^{\dagger
},a\right) }\left( T\right) _{22}=\overline{\mathcal{R}_{\left( a^{\dagger
},a\right) }\left( T\right) _{11}}+\overline{\mathcal{R}_{\left( a^{\dagger
},a\right) }\left( T\right) _{12}}-\overline{\mathcal{R}_{\left( a^{\dagger
},a\right) }\left( T\right) _{21}}-\overline{\mathcal{R}_{\left( a^{\dagger
},a\right) }\left( T\right) _{22}}
\end{equation}%
\begin{equation}
\mathcal{R}_{\left( a^{\dagger },a\right) }\left( T\right) _{12}+\mathcal{R}%
_{\left( a^{\dagger },a\right) }\left( T\right) _{21}-\mathcal{R}_{\left(
a^{\dagger },a\right) }\left( T\right) _{11}-\mathcal{R}_{\left( a^{\dagger
},a\right) }\left( T\right) _{22}=\overline{\mathcal{R}_{\left( a^{\dagger
},a\right) }\left( T\right) _{12}}+\overline{\mathcal{R}_{\left( a^{\dagger
},a\right) }\left( T\right) _{21}}-\overline{\mathcal{R}_{\left( a^{\dagger
},a\right) }\left( T\right) _{11}}-\overline{\mathcal{R}_{\left( a^{\dagger
},a\right) }\left( T\right) _{22}}.
\end{equation}

\begin{example}
Let 
\begin{equation}
T=\left( 
\begin{array}{cc}
u & v \\ 
w & z%
\end{array}%
\right)
\end{equation}%
then 
\begin{eqnarray}
u+z+v+w &=&\overline{u}+\overline{z}+\overline{v}+\overline{w} \\
u-z-v+w &=&\overline{u}-\overline{z}-\overline{v}+\overline{w} \\
u-z+v-w &=&\overline{u}-\overline{z}+\overline{v}-\overline{w} \\
-\left( u+z\right) +v+w &=&-\left( \overline{u}+\overline{z}\right) +%
\overline{v}+\overline{w}.
\end{eqnarray}
\end{example}
\end{enumerate}
\end{remark}

\begin{definition}
A map $T\in \mathcal{L}\left( \mathcal{F}_{1}\right) $ is $\ast $\emph{%
-preserving map} if 
\begin{equation}
T\left( A^{\dagger }\right) =T\left( A\right) ^{\dagger },A\in \mathcal{F}%
_{1}.
\end{equation}
\end{definition}

\begin{definition}
The set of $\ast $-preserving maps $\limfunc{star}\left( \mathcal{F}%
_{1}\right) $ $\subset \mathcal{L}\left( \mathcal{F}_{1}\right) $ is defined
as 
\begin{equation}
\limfunc{star}\left( \mathcal{F}_{1}\right) =\left\{ \left. T\in \mathcal{L}%
\left( \mathcal{F}_{1}\right) \right\vert T\left( A^{\dagger }\right)
=T\left( A\right) ^{\dagger },A\in \mathcal{F}_{1}\right\} .
\end{equation}
\end{definition}

\begin{theorem}
$T\in GL(2r,\mathbb{C)}$ is a nonsingular $\ast $-preserving map $\limfunc{%
Iff}$ 
\begin{equation}
T:\left( \boldsymbol{x}\right) ^{t}\rightarrow \left( \boldsymbol{y}\right)
^{t}\quad \forall \text{ self adjoint bases }\left\{ \left( \boldsymbol{x}%
\right) ,\left( \boldsymbol{y}\right) \right\} \text{ of }\mathcal{F}_{1}.
\end{equation}
\end{theorem}

\begin{theorem}
$T\in \limfunc{star}\left( \mathcal{F}_{1}\right) \subset GL(2r,\mathbb{C)}$
is a nonsingular $\ast $-preserving map $\limfunc{Iff}$ 
\begin{equation}
T\in \limfunc{star}\left( \mathcal{F}_{1}\right) \cong GL(2r,\mathbb{R)}%
\subset GL(2r,\mathbb{C)}.
\end{equation}
\end{theorem}

\begin{corollary}
The reprersentations of $T\in \limfunc{star}\left( \mathcal{F}_{1}\right) $
are real matrices $\mathcal{R}_{\left( x\right) }\left( T\right) $ $\in 
\mathcal{M}(2r,2r,\mathbb{R}^{2r}\mathbb{)}$ $\forall $ self adjoint basis $%
\left( \boldsymbol{x}\right) $ of $\mathcal{F}_{1}.$
\end{corollary}

\begin{lemma}
$\limfunc{star}\left( \mathcal{F}_{1}\right) $ is a group that is isomorphic
to $GL(2r,\mathbb{R)}.$

\begin{proof}
Consider the product 
\begin{equation}
\left( 
\begin{array}{cc}
\mathbb{U}_{1} & \overline{\mathbb{V}}_{1} \\ 
\mathbb{V}_{1} & \overline{\mathbb{U}}_{1}%
\end{array}%
\right) \left( 
\begin{array}{cc}
\mathbb{U}_{2} & \overline{\mathbb{V}}_{2} \\ 
\mathbb{V}_{2} & \overline{\mathbb{U}}_{2}%
\end{array}%
\right) =\left( 
\begin{array}{cc}
\mathbb{U}_{1}\mathbb{U}_{2}+\overline{\mathbb{V}}_{1}\mathbb{V}_{2} & 
\mathbb{U}_{1}\overline{\mathbb{V}}_{2}+\overline{\mathbb{V}}_{1}\overline{%
\mathbb{U}}_{2} \\ 
\mathbb{V}_{1}\mathbb{U}_{2}+\overline{\mathbb{U}}_{1}\mathbb{V}_{2} & 
\mathbb{V}_{1}\overline{\mathbb{V}}_{2}+\overline{\mathbb{U}}_{1}\overline{%
\mathbb{U}}_{2}%
\end{array}%
\right) =\left( 
\begin{array}{cc}
\mathbb{U}_{3} & \overline{\mathbb{V}}_{3} \\ 
\mathbb{V}_{3} & \overline{\mathbb{U}}_{3}%
\end{array}%
\right)
\end{equation}%
so non singular matrices of the form $\left( 
\begin{array}{cc}
\mathbb{U} & \overline{\mathbb{V}} \\ 
\mathbb{V} & \overline{\mathbb{U}}%
\end{array}%
\right) $ produce a group. The unitary matrix $\mathbb{K}=$ $\left( 
\begin{array}{cc}
\frac{1}{\sqrt{2}}\mathbb{I}_{r} & \frac{i}{\sqrt{2}}\mathbb{I}_{r} \\ 
\frac{1}{\sqrt{2}}\mathbb{I}_{r} & -\frac{i}{\sqrt{2}}\mathbb{I}_{r}%
\end{array}%
\right) \in U\left( 2r\right) $ transforms $\left( 
\begin{array}{cc}
\mathbb{U} & \overline{\mathbb{V}} \\ 
\mathbb{V} & \overline{\mathbb{U}}%
\end{array}%
\right) $ to real form $\left( 
\begin{array}{cc}
\func{Re}\left( \mathbb{U+V}\right) & -\func{Im}\left( \mathbb{U+V}\right)
\\ 
\func{Im}\left( \mathbb{U-V}\right) & \func{Re}\left( \mathbb{U-V}\right)%
\end{array}%
\right) .$ Hence $\limfunc{star}\left( \mathcal{F}_{1}\right) \cong GL(2r,%
\mathbb{R)\blacksquare }.$
\end{proof}
\end{lemma}

\begin{remark}
Transformations $T\in \limfunc{star}\left( \mathcal{F}_{1}\right) $ 
\begin{equation}
T\left( 
\begin{array}{cc}
\mathbf{a}^{\dagger } & \mathbf{a}%
\end{array}%
\right) =\left( 
\begin{array}{cc}
\mathbf{a}^{\dagger } & \mathbf{a}%
\end{array}%
\right) \mathbb{T}_{a}=\left( 
\begin{array}{cc}
\mathbf{a}^{\dagger } & \mathbf{a}%
\end{array}%
\right) \mathcal{R}_{a}\left( T\right) =\left( 
\begin{array}{cc}
\mathbf{a}^{\dagger } & \mathbf{a}%
\end{array}%
\right) \left( 
\begin{array}{cc}
\mathbb{U} & \overline{\mathbb{V}} \\ 
\mathbb{V} & \overline{\mathbb{U}}%
\end{array}%
\right) =\left( 
\begin{array}{cc}
\mathbf{b}^{\dagger } & \mathbf{b}%
\end{array}%
\right) ,
\end{equation}%
where 
\begin{equation}
\left( 
\begin{array}{cc}
\mathbb{U} & \overline{\mathbb{V}} \\ 
\mathbb{V} & \overline{\mathbb{U}}%
\end{array}%
\right) =\left( 
\begin{array}{cc}
\mathcal{R}_{\left( a^{\dagger },a\right) }\left( T\right) _{11} & \overline{%
\mathcal{R}_{\left( a^{\dagger },a\right) }\left( T\right) _{21}} \\ 
\mathcal{R}_{\left( a^{\dagger },a\right) }\left( T\right) _{21} & \overline{%
\mathcal{R}_{\left( a^{\dagger },a\right) }\left( T\right) _{11}}%
\end{array}%
\right)
\end{equation}%
then the transformed creation and annihilation operators are given by 
\begin{align}
b_{k}^{\dagger }& =\sum_{1\leq j\leq r}\left( a_{j}^{\dagger
}U_{jk}+a_{j}V_{jk}\right) \\
b_{k}& =\sum_{1\leq j\leq r}\left( a_{j}^{\dagger }\bar{V}_{jk}+a_{j}\bar{U}%
_{jk}\right)
\end{align}
\end{remark}

\begin{theorem}
$T\in \overset{}{\limfunc{star}\left( \mathcal{F}_{1}\right) \subset GL(2r,%
\mathcal{F}_{1}\mathbb{)}\cong GL\left( 2r,\mathbb{C}\right) }$ $\limfunc{Iff%
}$ 
\begin{equation}
\mathcal{R}_{\left( a^{\dagger },a\right) }\left( T\right) =\left( 
\begin{array}{cc}
\mathcal{R}_{\left( a^{\dagger },a\right) }\left( T\right) _{11} & \overline{%
\mathcal{R}_{\left( a^{\dagger },a\right) }\left( T\right) _{21}} \\ 
\mathcal{R}_{\left( a^{\dagger },a\right) }\left( T\right) _{21} & \overline{%
\mathcal{R}_{\left( a^{\dagger },a\right) }\left( T\right) _{11}}%
\end{array}%
\right) =\left( 
\begin{array}{cc}
\mathbb{U} & \overline{\mathbb{V}} \\ 
\mathbb{V} & \overline{\mathbb{U}}%
\end{array}%
\right) \quad \forall \text{ non self adjoint basis }\left( \boldsymbol{a}%
^{\dagger },\boldsymbol{a}\right) \text{ of }\mathcal{F}_{1}.
\end{equation}

\begin{proof}
from the preceding remark.
\end{proof}
\end{theorem}

\begin{remark}
\begin{enumerate}
\item Transformation $K\in \mathcal{L}\left( \mathcal{F}_{1}\right) $ from
Bogolubons to Majoranos 
\begin{equation}
K\left( 
\begin{array}{cc}
\mathbf{a}^{\dagger } & \mathbf{a}%
\end{array}%
\right) =\left( 
\begin{array}{cc}
\mathbf{x}_{1} & \mathbf{x}_{2}%
\end{array}%
\right) =\left( 
\begin{array}{cc}
\mathbf{a}^{\dagger } & \mathbf{a}%
\end{array}%
\right) \mathbb{K}_{a}\mathbb{\ }=\left( 
\begin{array}{cc}
\mathbf{a}^{\dagger } & \mathbf{a}%
\end{array}%
\right) \mathcal{R}_{a}\left( K\right) =\left( 
\begin{array}{cc}
\mathbf{a}^{\dagger } & \mathbf{a}%
\end{array}%
\right) \left( 
\begin{array}{cc}
\frac{1}{\sqrt{2}}\mathbb{I}_{r} & \frac{i}{\sqrt{2}}\mathbb{I}_{r} \\ 
\frac{1}{\sqrt{2}}\mathbb{I}_{r} & -\frac{i}{\sqrt{2}}\mathbb{I}_{r}%
\end{array}%
\right) .
\end{equation}

\item Representations of $T\in \mathcal{L}\left( \mathcal{F}_{1}\right) $
wrt Bogolubons and Majoranos 
\begin{eqnarray}
T\left( 
\begin{array}{cc}
\mathbf{a}^{\dagger } & \mathbf{a}%
\end{array}%
\right) &=&\left( 
\begin{array}{cc}
\mathbf{a}^{\dagger } & \mathbf{a}%
\end{array}%
\right) \mathcal{R}_{a}\left( T\right) \\
T\left( 
\begin{array}{cc}
\mathbf{x}_{1} & \mathbf{x}_{2}%
\end{array}%
\right) &=&\left( 
\begin{array}{cc}
\mathbf{x}_{1} & \mathbf{x}_{2}%
\end{array}%
\right) \mathcal{R}_{x}\left( T\right) \\
T\left( 
\begin{array}{cc}
\mathbf{a}^{\dagger } & \mathbf{a}%
\end{array}%
\right) \mathcal{R}_{a}\left( K\right) &=&\left( 
\begin{array}{cc}
\mathbf{a}^{\dagger } & \mathbf{a}%
\end{array}%
\right) \mathbb{K}_{a}\mathcal{R}_{x}\left( T\right) \\
T\left( 
\begin{array}{cc}
\mathbf{a}^{\dagger } & \mathbf{a}%
\end{array}%
\right) \mathcal{R}_{a}\left( K\right) \mathcal{R}_{a}\left( K\right)
^{\ddagger } &=&\left( 
\begin{array}{cc}
\mathbf{a}^{\dagger } & \mathbf{a}%
\end{array}%
\right) \mathcal{R}_{a}\left( K\right) \mathcal{R}_{x}\left( T\right) 
\mathcal{R}_{a}\left( K\right) ^{\dagger }.
\end{eqnarray}

\item Transformation between representations 
\begin{eqnarray}
\mathcal{R}_{a}\left( T\right) &=&\mathcal{R}_{a}\left( K\right) \mathcal{R}%
_{x}\left( T\right) \mathcal{R}_{a}\left( K\right) ^{\dagger } \\
\left( 
\begin{array}{cc}
\mathcal{R}_{a}\left( T\right) _{11} & \mathcal{R}_{a}\left( T\right) _{12}
\\ 
\mathcal{R}_{a}\left( T\right) _{21} & \mathcal{R}_{a}\left( T\right) _{22}%
\end{array}%
\right) &=&\left( 
\begin{array}{cc}
\frac{1}{\sqrt{2}}\mathbb{I}_{r} & \frac{i}{\sqrt{2}}\mathbb{I}_{r} \\ 
\frac{1}{\sqrt{2}}\mathbb{I}_{r} & -\frac{i}{\sqrt{2}}\mathbb{I}_{r}%
\end{array}%
\right) \left( 
\begin{array}{cc}
\mathcal{R}_{x}\left( T\right) _{11} & \mathcal{R}_{x}\left( T\right) _{12}
\\ 
\mathcal{R}_{x}\left( T\right) _{21} & \mathcal{R}_{x}\left( T\right) _{22}%
\end{array}%
\right) \left( 
\begin{array}{cc}
\frac{1}{\sqrt{2}}\mathbb{I}_{r} & \frac{i}{\sqrt{2}}\mathbb{I}_{r} \\ 
\frac{1}{\sqrt{2}}\mathbb{I}_{r} & -\frac{i}{\sqrt{2}}\mathbb{I}_{r}%
\end{array}%
\right) ^{\dagger }.
\end{eqnarray}
\end{enumerate}
\end{remark}

\begin{remark}
$\widehat{\mathbb{K}}\in \mathcal{L}\left( GL\left( \overset{}{2r,\mathbb{R}}%
\right) ,GL\left( \overset{}{2r,\mathbb{C}}\right) \right) $ maps the
complex unit $i\mathbb{I}_{2r}$ to real form 
\begin{equation}
\widehat{\mathbb{K}}\left( \left( 
\begin{array}{cc}
i\mathbb{I}_{r} & \mathbb{O} \\ 
\mathbb{O} & i\mathbb{I}_{r}%
\end{array}%
\right) \right) =\left( 
\begin{array}{cc}
\mathbb{O} & -\mathbb{I}_{r} \\ 
\mathbb{I}_{r} & \mathbb{O}%
\end{array}%
\right) =\mathbb{J}_{2r}.
\end{equation}%
\textbf{NOTE} that this is \textbf{not} the realification of $i\mathbb{I}%
_{2r}$ which would be $\mathbb{J}_{4r}!$
\end{remark}

\subsection{CAR in Complex and Real Form.}

\begin{remark}
$C,D\in \mathcal{L}\left( \mathcal{F}_{1}\right) $%
\begin{eqnarray}
CD\left( \boldsymbol{A}\right) ^{t} &=&\left( \boldsymbol{A}\right) ^{t}%
\mathcal{R}_{A}\left( C\right) \mathcal{R}_{A}\left( D\right) \\
DC\left( \boldsymbol{A}\right) ^{t} &=&\left( \boldsymbol{A}\right) ^{t}%
\mathcal{R}_{A}\left( D\right) \mathcal{R}_{A}\left( C\right) \\
\left( CD+DC\right) \left( \boldsymbol{A}\right) ^{t} &=&\left\{ \overset{}{%
C,D}\right\} \left( \boldsymbol{A}\right) ^{t}=\left( \boldsymbol{A}\right)
^{t}\left\{ \overset{}{\mathcal{R}_{A}\left( C\right) ,\mathcal{R}_{A}\left(
D\right) }\right\}
\end{eqnarray}
\end{remark}

\begin{remark}
\begin{eqnarray}
T\left( \boldsymbol{A}\right) ^{t} &=&T\left( 
\begin{array}{ccc}
A_{1} & \cdots & A_{2r}%
\end{array}%
\right) =\left( 
\begin{array}{ccc}
TA_{1} & \cdots & TA_{2r}%
\end{array}%
\right) =\left( 
\begin{array}{ccc}
A_{1} & \cdots & A_{2r}%
\end{array}%
\right) \mathcal{R}_{A}\left( T\right) \\
\left( T\left( \boldsymbol{A}\right) ^{t}\right) ^{t} &=&T\left( 
\begin{array}{c}
A_{1} \\ 
\vdots \\ 
A_{2r}%
\end{array}%
\right) =\mathcal{R}_{A}\left( T\right) ^{t}\left( 
\begin{array}{c}
A_{1} \\ 
\vdots \\ 
A_{2r}%
\end{array}%
\right) =\left( 
\begin{array}{c}
TA_{1} \\ 
\vdots \\ 
TA_{2r}%
\end{array}%
\right) .
\end{eqnarray}
\end{remark}

\begin{theorem}
\label{CAR Theorem} Let $T\in \mathcal{L}\left( \mathcal{F}_{1},\mathbb{C}%
\right) $ and 
\begin{equation}
\left( \boldsymbol{B}\right) =T\left( \boldsymbol{A}\right) ^{t}=\left( 
\mathbf{A}\right) ^{t}\mathbb{T}_{a}=\left( \mathbf{A}\right) ^{t}\mathcal{R}%
_{A}\left( T\right) ,
\end{equation}%
where $\left( \mathbf{A}\right) $ and $\left( \mathbf{B}\right) $ are bases
of $\mathcal{F}_{1}$ \textbf{then} 
\begin{equation}
\left\{ \boldsymbol{B},\boldsymbol{B}\right\} =\mathcal{R}_{A}\left(
T\right) ^{t}\left\{ \overset{}{\boldsymbol{A},\boldsymbol{A}}\right\} 
\mathcal{R}_{A}\left( T\right)
\end{equation}%
where 
\begin{equation}
\left\{ \boldsymbol{X},\boldsymbol{X}\right\} =\left( 
\begin{array}{ccc}
\left\{ X_{1},X_{1}\right\} & \cdots & \left\{ X_{1},X_{2r}\right\} \\ 
\vdots & \ddots & \vdots \\ 
\left\{ X_{2r},X_{1}\right\} & \cdots & \left\{ X_{2r},X_{2r}\right\}%
\end{array}%
\right) ,\quad \boldsymbol{X=A},\boldsymbol{B}
\end{equation}
\end{theorem}

\begin{proof}
\begin{equation}
B_{j}=\sum_{1\leq l\leq 2r}A_{l}\mathcal{R}_{A}\left( T\right) _{lj}
\end{equation}%
then 
\begin{eqnarray*}
\left\{ B_{j},B_{k}\right\} &=&\left\{ \sum_{1\leq l\leq 2r}A_{l}\mathcal{R}%
_{A}\left( T\right) _{lj},\sum_{1\leq m\leq 2r}A_{m}\mathcal{R}_{A}\left(
T\right) _{mk}\right\} =\sum_{1\leq l,m\leq 2r}\left\{ A_{l}\mathcal{R}%
_{A}\left( T\right) _{lj},A_{m}\mathcal{R}_{A}\left( T\right) _{mk}\right\}
\\
&=&\sum_{1\leq l,m\leq 2r}\mathcal{R}_{A}\left( T\right) _{jl}^{t}\left\{
A_{l},A_{m}\right\} \mathcal{R}_{A}\left( T\right) _{mk}=\left( \overset{}{%
\mathcal{R}_{A}\left( T\right) ^{t}\left\{ \boldsymbol{A},\boldsymbol{A}%
\right\} \mathcal{R}_{A}\left( T\right) }\right) _{jk}\blacksquare .
\end{eqnarray*}
\end{proof}

\begin{remark}
Noting that $X_{j}=X_{j}^{t}\,\,\,\forall X_{j}\in \mathcal{F}_{1},1\leq
j\leq 2r$ one has the following tensor product identities

\begin{enumerate}
\item 
\begin{enumerate}
\item Column Row/Tensor Product Notation%
\begin{equation}
\left( \boldsymbol{X}\right) \left( \boldsymbol{X}\right) ^{t}=\left( 
\begin{array}{c}
X_{1} \\ 
\vdots \\ 
X_{2r}%
\end{array}%
\right) \otimes \left( 
\begin{array}{ccc}
X_{1} & \cdots & X_{2r}%
\end{array}%
\right) =\left( 
\begin{array}{c}
X_{1} \\ 
\vdots \\ 
X_{2r}%
\end{array}%
\right) \left( 
\begin{array}{ccc}
X_{1} & \cdots & X_{2r}%
\end{array}%
\right) .
\end{equation}

\item Column/Row Tensor Product Notation of Super Vectors Evaluated%
\begin{equation}
\left( \boldsymbol{XX}^{t}\right) =\left( \boldsymbol{X}\right) \left( 
\boldsymbol{X}\right) ^{t}=\left( 
\begin{array}{c}
X_{1} \\ 
\vdots \\ 
X_{2r}%
\end{array}%
\right) \left( 
\begin{array}{ccc}
X_{1} & \cdots & X_{2r}%
\end{array}%
\right) =\left( 
\begin{array}{c}
X_{1}\left( 
\begin{array}{ccc}
X_{1} & \cdots & X_{2r}%
\end{array}%
\right) \\ 
\vdots \\ 
X_{2r}\left( 
\begin{array}{ccc}
X_{1} & \cdots & X_{2r}%
\end{array}%
\right)%
\end{array}%
\right) =\left( 
\begin{array}{ccc}
X_{1}X_{1} & \cdots & X_{1}X_{2r} \\ 
\vdots & \ddots & \vdots \\ 
X_{2r}X_{1} & \cdots & X_{2r}X_{2r}%
\end{array}%
\right) .
\end{equation}
\end{enumerate}

\begin{enumerate}
\item Row/Column/Tensor Product Notation of Super Vectors 
\begin{equation}
\left( \boldsymbol{X}\right) ^{t}\left( \boldsymbol{X}\right) =\left( 
\begin{array}{ccc}
X_{1} & \cdots & X_{2r}%
\end{array}%
\right) \otimes \left( 
\begin{array}{c}
X_{1} \\ 
\vdots \\ 
X_{2r}%
\end{array}%
\right) =\left( 
\begin{array}{ccc}
X_{1} & \cdots & X_{2r}%
\end{array}%
\right) \left( 
\begin{array}{c}
X_{1} \\ 
\vdots \\ 
X_{2r}%
\end{array}%
\right) .
\end{equation}

\item Row/Column Tensor Product Notation of Super Vectors Evaluated 
\begin{equation}
\left( \boldsymbol{X^{t}X}\right) =\left( \boldsymbol{X}\right) ^{t}\left( 
\boldsymbol{X}\right) =\left( 
\begin{array}{ccc}
X_{1} & \cdots & X_{2r}%
\end{array}%
\right) \left( 
\begin{array}{c}
X_{1} \\ 
\vdots \\ 
X_{2r}%
\end{array}%
\right) =\left( 
\begin{array}{ccc}
X_{1}\left( 
\begin{array}{c}
X_{1} \\ 
\vdots \\ 
X_{2r}%
\end{array}%
\right) & \cdots & X_{2r}\left( 
\begin{array}{c}
X_{1} \\ 
\vdots \\ 
X_{2r}%
\end{array}%
\right)%
\end{array}%
\right) =\left( 
\begin{array}{ccc}
X_{1}X_{1} & \cdots & X_{2r}X_{1} \\ 
\vdots & \ddots & \vdots \\ 
X_{1}X_{2r} & \cdots & X_{2r}X_{2r}%
\end{array}%
\right) .
\end{equation}
\end{enumerate}

\item Scalar Product of Super Vectors%
\begin{equation}
\left( 
\begin{array}{ccc}
X_{1} & \cdots & X_{2r}%
\end{array}%
\right) \cdot \left( 
\begin{array}{c}
X_{1} \\ 
\vdots \\ 
X_{2r}%
\end{array}%
\right) =\sum_{1\leq j\leq 2r}X_{j}X_{j}=\limfunc{Tr}\left( \left( 
\begin{array}{ccc}
X_{1} & \cdots & X_{2r}%
\end{array}%
\right) \left( 
\begin{array}{c}
X_{1} \\ 
\vdots \\ 
X_{2r}%
\end{array}%
\right) \right) .
\end{equation}

\item The matrices $\left( \boldsymbol{XX}^{t}\right) $ and $\left( 
\boldsymbol{X^{t}X}\right) $ are symmetric i.e. 
\begin{equation}
\left( \boldsymbol{XX}^{t}\right) =\left( 
\begin{array}{ccc}
X_{1}X_{1} & \cdots & X_{1}X_{2r} \\ 
\vdots & \ddots & \vdots \\ 
X_{2r}X_{1} & \cdots & X_{2r}X_{2r}%
\end{array}%
\right) ^{t}=\left( 
\begin{array}{ccc}
X_{1}X_{1} & \cdots & X_{1}X_{2r} \\ 
\vdots & \ddots & \vdots \\ 
X_{2r}X_{1} & \cdots & X_{2r}X_{2r}%
\end{array}%
\right) =\left( \boldsymbol{XX}^{t}\right) ^{t}
\end{equation}%
and 
\begin{equation}
\left( \boldsymbol{X^{t}X}\right) =\left( 
\begin{array}{ccc}
X_{1}X_{1} & \cdots & X_{2r}X_{1} \\ 
\vdots & \ddots & \vdots \\ 
X_{1}X_{2r} & \cdots & X_{2r}X_{2r}%
\end{array}%
\right) ^{t}=\left( 
\begin{array}{ccc}
X_{1}X_{1} & \cdots & X_{2r}X_{1} \\ 
\vdots & \ddots & \vdots \\ 
X_{1}X_{2r} & \cdots & X_{2r}X_{2r}%
\end{array}%
\right) =\left( \boldsymbol{X^{t}X}\right) ^{t}.
\end{equation}

\item Transformation of $\left( \boldsymbol{X^{t}X}\right) $ 
\begin{eqnarray}
\left( T\boldsymbol{X}\right) \boldsymbol{^{t}}\left( \boldsymbol{TX}\right)
&=&\left( 
\begin{array}{ccc}
TX_{1}TX_{1} & \cdots & TX_{2r}TX_{1} \\ 
\vdots & \ddots & \vdots \\ 
TX_{1}TX_{2r} & \cdots & TX_{2r}TX_{2r}%
\end{array}%
\right) =T\left( 
\begin{array}{ccc}
X_{1} & \cdots & X_{2r}%
\end{array}%
\right) T\left( 
\begin{array}{c}
X_{1} \\ 
\vdots \\ 
X_{2r}%
\end{array}%
\right) =\left( 
\begin{array}{ccc}
TX_{1} & \cdots & TX_{2r}%
\end{array}%
\right) \left( 
\begin{array}{c}
TX_{1} \\ 
\vdots \\ 
TX_{2r}%
\end{array}%
\right) \\
&=&\left( 
\begin{array}{ccc}
X_{1} & \cdots & X_{2r}%
\end{array}%
\right) \mathcal{R}_{X}\left( T\right) \mathcal{R}_{X}\left( T\right)
^{t}\left( 
\begin{array}{c}
X_{1} \\ 
\vdots \\ 
X_{2r}%
\end{array}%
\right)
\end{eqnarray}

\item The details of the action of $T$ on $\left( 
\begin{array}{ccc}
X_{1} & \cdots & X_{2r}%
\end{array}%
\right) $ is given by 
\begin{equation}
TX_{j}=\sum_{1\leq l\leq 2r}X_{l}\mathcal{R}_{X}\left( T\right) _{lj}
\end{equation}%
leading to 
\begin{eqnarray}
TX_{j}TX_{k} &=&\sum_{1\leq l,m\leq 2r}X_{l}\mathcal{R}_{X}\left( T\right)
_{lj}X_{m}\mathcal{R}_{X}\left( T\right) _{mj}=\sum_{1\leq l,m\leq
2r}X_{l}X_{m}\mathcal{R}_{X}\left( T\right) _{lj}\mathcal{R}_{X}\left(
T\right) _{mj}=\sum_{1\leq l,m\leq 2r}X_{l}X_{m}\mathcal{R}_{X}\left(
T\right) _{jl}^{t}\mathcal{R}_{X}\left( T\right) _{mj} \\
&=&\sum_{1\leq l,m\leq 2r}\mathcal{R}_{X}\left( T\right) _{jl}^{t}X_{l}X_{m}%
\mathcal{R}_{X}\left( T\right) _{mj}
\end{eqnarray}

\item This gives 
\begin{equation}
\left( T\boldsymbol{X}\right) \boldsymbol{^{t}}\left( T\boldsymbol{X}\right)
=\mathcal{R}_{X}\left( T\right) ^{t}\left( \boldsymbol{X^{t}X}\right) 
\mathcal{R}_{X}\left( T\right)
\end{equation}%
and 
\begin{equation}
\left( T\boldsymbol{X}\right) \left( T\boldsymbol{X}\right) \boldsymbol{^{t}}%
=\mathcal{R}_{X}\left( T\right) ^{t}\left( \boldsymbol{XX^{t}}\right) 
\mathcal{R}_{X}\left( T\right) .
\end{equation}

\begin{example}
Consider the $r=1$ case and the two matrices 
\begin{equation}
\left( 
\begin{array}{cc}
X_{1}X_{1} & X_{1}X_{2} \\ 
X_{2}X_{1} & X_{2}X_{2}%
\end{array}%
\right)
\end{equation}%
and 
\begin{equation}
\left( 
\begin{array}{cc}
X_{1}X_{1} & X_{2}X_{1} \\ 
X_{1}X_{2} & X_{2}X_{2}%
\end{array}%
\right) .
\end{equation}%
The transformed $\left( 
\begin{array}{cc}
X_{1} & X_{2}%
\end{array}%
\right) $ is given by 
\begin{eqnarray}
TX_{1} &=&X_{1}2+X_{2}3 \\
TX_{2} &=&X_{1}4+X_{2}5
\end{eqnarray}%
and 
\begin{eqnarray}
\left( 
\begin{array}{cc}
X_{1} & X_{2}%
\end{array}%
\right) \left( 
\begin{array}{cc}
2 & 4 \\ 
3 & 5%
\end{array}%
\right) &=&\left( 
\begin{array}{cc}
X_{1}2+X_{2}3 & X_{1}4+X_{2}5%
\end{array}%
\right) \\
\left( 
\begin{array}{cc}
2 & 4 \\ 
3 & 5%
\end{array}%
\right) ^{t}\left( 
\begin{array}{c}
X_{1} \\ 
X_{2}%
\end{array}%
\right) &=&\left( 
\begin{array}{c}
X_{1}2+X_{2}3 \\ 
X_{1}4+X_{2}5%
\end{array}%
\right) ,
\end{eqnarray}%
i.e. 
\begin{equation}
T=\left( 
\begin{array}{cc}
2 & 4 \\ 
3 & 5%
\end{array}%
\right) .
\end{equation}%
The trasformed matrix $\left( T\boldsymbol{X}\right) \left( \boldsymbol{TX}%
\right) ^{t}$ is 
\begin{eqnarray}
\left( T\boldsymbol{X}\right) \left( \boldsymbol{TX}\right) ^{t} &=&\left( 
\begin{array}{cc}
TX_{1}TX_{1} & TX_{1}TX_{2} \\ 
TX_{2}TX_{1} & TX_{2}TX_{2}%
\end{array}%
\right) =\left( 
\begin{array}{cc}
\left( X_{1}2+X_{2}3\right) \left( X_{1}2+X_{2}3\right) & \left(
X_{1}2+X_{2}3\right) \left( X_{1}4+X_{2}5\right) \\ 
\left( X_{1}4+X_{2}5\right) \left( X_{1}2+X_{2}3\right) & \left(
X_{1}4+X_{2}5\right) \left( X_{1}4+X_{2}5\right)%
\end{array}%
\right) \\
&=&\left( 
\begin{array}{cc}
X_{1}2X_{1}2+X_{1}2X_{2}3+X_{2}3X_{1}2+X_{2}3X_{2}3 & 
X_{1}2X_{1}4+X_{1}2X_{2}5+X_{2}3X_{1}4+X_{2}3X_{2}5 \\ 
X_{1}4X_{1}2+X_{1}4X_{2}3+X_{2}5X_{1}2+X_{2}5X_{2}3 & 
X_{1}4X_{1}4+X_{1}4X_{2}5+X_{2}5X_{1}4+X_{2}5X_{2}5%
\end{array}%
\right) \\
&=&\left( 
\begin{array}{cc}
X_{1}X_{1}4+X_{1}X_{2}6+X_{2}X_{1}6+X_{2}X_{2}9 & 
X_{1}X_{1}8+X_{1}X_{2}10+X_{2}X_{1}12+X_{2}X_{2}15 \\ 
X_{1}X_{1}8+X_{1}X_{2}12+X_{2}X_{1}10+X_{2}X15 & 
X_{1}X_{1}16+X_{1}X_{2}20+X_{2}X_{1}20+X_{2}X_{2}25%
\end{array}%
\right)
\end{eqnarray}%
and 
\begin{eqnarray}
&&T^{t}\left( \boldsymbol{X}\right) \left( \boldsymbol{X}\right)
^{t}T=\left( 
\begin{array}{cc}
2 & 4 \\ 
3 & 5%
\end{array}%
\right) ^{t}\left( 
\begin{array}{cc}
X_{1}X_{1} & X_{1}X_{2} \\ 
X_{2}X_{1} & X_{2}X_{2}%
\end{array}%
\right) \left( 
\begin{array}{cc}
2 & 4 \\ 
3 & 5%
\end{array}%
\right) =\left( 
\begin{array}{cc}
2 & 3 \\ 
4 & 5%
\end{array}%
\right) \left( 
\begin{array}{cc}
X_{1}X_{1} & X_{1}X_{2} \\ 
X_{2}X_{1} & X_{2}X_{2}%
\end{array}%
\right) \left( 
\begin{array}{cc}
2 & 4 \\ 
3 & 5%
\end{array}%
\right) \\
&=&\left( 
\begin{array}{cc}
2X_{1}X_{1}+3X_{2}X_{1} & 2X_{1}X_{2}+3X_{2}X_{2} \\ 
4X_{1}X_{1}+5X_{2}X_{1} & 4X_{1}X_{2}+5X_{2}X_{2}%
\end{array}%
\right) \left( 
\begin{array}{cc}
2 & 4 \\ 
3 & 5%
\end{array}%
\right) \\
&=&\left( 
\begin{array}{cc}
4X_{1}X_{1}+6X_{2}X_{1}+6X_{1}X_{1}+9X_{2}X_{1} & 
8X_{1}X_{1}+12X_{2}X_{1}+10X_{1}X_{2}+15X_{2}X_{2} \\ 
8X_{1}X_{1}+10X_{2}X_{1}+12X_{1}X_{2}+15X_{2}X_{2} & 
16X_{1}X_{1}+20X_{2}X_{1}+20X_{1}X_{2}+25X_{2}X_{2}%
\end{array}%
\right)
\end{eqnarray}%
hence 
\begin{equation}
\left( T\boldsymbol{X}\right) \left( \boldsymbol{TX}\right) ^{t}=T^{t}\left( 
\boldsymbol{X}\right) \left( \boldsymbol{X}\right) ^{t}T.
\end{equation}%
The trasformed matrix $\left( T\boldsymbol{X}\right) \boldsymbol{^{t}}\left(
T\boldsymbol{X}\right) $ is%
\begin{eqnarray}
\left( T\boldsymbol{X}\right) \boldsymbol{^{t}}\left( T\boldsymbol{X}\right)
&=&\left( 
\begin{array}{cc}
TX_{1}TX_{1} & TX_{2}TX_{1} \\ 
TX_{1}TX_{2} & TX_{2}TX_{2}%
\end{array}%
\right) =\left( 
\begin{array}{cc}
\left( X_{1}2+X_{2}3\right) \left( X_{1}2+X_{2}3\right) & \left(
X_{1}4+X_{2}5\right) \left( X_{1}2+X_{2}3\right) \\ 
\left( X_{1}2+X_{2}3\right) \left( X_{1}4+X_{2}5\right) & \left(
X_{1}4+X_{2}5\right) \left( X_{1}4+X_{2}5\right)%
\end{array}%
\right) \\
&=&\left( 
\begin{array}{cc}
X_{1}X_{1}4+X_{1}X_{2}6+X_{2}X_{1}6+X_{2}X_{2}9 & 
X_{1}X_{1}8+X_{1}X_{2}12+X_{2}X_{1}10+X_{2}X15 \\ 
X_{1}X_{1}8+X_{1}X_{2}10+X_{2}X_{1}12+X_{2}X_{2}15 & 
X_{1}X_{1}16+X_{1}X_{2}20+X_{2}X_{1}20+X_{2}X_{2}25%
\end{array}%
\right)
\end{eqnarray}%
and 
\begin{eqnarray}
&&T^{t}\left( \boldsymbol{X}\right) ^{t}\left( \boldsymbol{X}\right)
T=\left( 
\begin{array}{cc}
2 & 4 \\ 
3 & 5%
\end{array}%
\right) ^{t}\left( 
\begin{array}{cc}
X_{1}X_{1} & X_{2}X_{1} \\ 
X_{1}X_{2} & X_{2}X_{2}%
\end{array}%
\right) \left( 
\begin{array}{cc}
2 & 4 \\ 
3 & 5%
\end{array}%
\right) =\left( 
\begin{array}{cc}
2 & 3 \\ 
4 & 5%
\end{array}%
\right) \left( 
\begin{array}{cc}
X_{1}X_{1} & X_{2}X_{1} \\ 
X_{1}X_{2} & X_{2}X_{2}%
\end{array}%
\right) \left( 
\begin{array}{cc}
2 & 4 \\ 
3 & 5%
\end{array}%
\right) \\
&=&\left( 
\begin{array}{cc}
2X_{1}X_{1}+3X_{1}X_{2} & 2X_{2}X_{1}+3X_{2}X_{2} \\ 
4X_{1}X_{1}+5X_{1}X_{2} & 4X_{2}X_{1}+5X_{2}X_{2}%
\end{array}%
\right) \left( 
\begin{array}{cc}
2 & 4 \\ 
3 & 5%
\end{array}%
\right) \\
&=&\left( 
\begin{array}{cc}
4X_{1}X_{1}+6X_{1}X_{2}+6X_{2}X_{1}+9X_{2}X_{2} & 
8X_{1}X_{1}+12X_{1}X_{2}+10X_{2}X_{1}+15X_{2}X_{2} \\ 
8X_{1}X_{1}+10X_{1}X_{2}+12X_{2}X_{1}+15X_{2}X_{2} & 
16X_{1}X_{1}+20X_{1}X_{2}+20X_{2}X_{1}+25X_{2}X_{2}%
\end{array}%
\right)
\end{eqnarray}%
hence 
\begin{equation}
\left( T\boldsymbol{X}\right) \boldsymbol{^{t}}\left( T\boldsymbol{X}\right)
=T^{t}\left( \boldsymbol{X}\right) ^{t}\left( \boldsymbol{X}\right) T.
\end{equation}
\end{example}

\item 
\begin{equation}
\left( T\boldsymbol{X}\right) \boldsymbol{^{t}}\left( T\boldsymbol{X}\right)
=T^{t}\left( \boldsymbol{X}\right) ^{t}\left( \boldsymbol{X}\right) T=\left( 
\begin{array}{ccc}
X_{1} & \cdots & X_{2r}%
\end{array}%
\right) \mathcal{R}_{X}\left( T\right) \mathcal{R}_{X}\left( T\right)
^{t}\left( 
\begin{array}{c}
X_{1} \\ 
\vdots \\ 
X_{2r}%
\end{array}%
\right) =\mathcal{R}_{X}\left( T\right) ^{t}\left( 
\begin{array}{ccc}
X_{1} & \cdots & X_{2r}%
\end{array}%
\right) \left( 
\begin{array}{c}
X_{1} \\ 
\vdots \\ 
X_{2r}%
\end{array}%
\right) \mathcal{R}_{X}\left( T\right)
\end{equation}%
and 
\begin{equation}
\left( T\boldsymbol{X}\right) \left( \boldsymbol{TX}\right) ^{t}=T^{t}\left( 
\boldsymbol{X}\right) \left( \boldsymbol{X}\right) ^{t}T=\mathcal{R}%
_{X}\left( T\right) ^{t}\left( 
\begin{array}{c}
X_{1} \\ 
\vdots \\ 
X_{2r}%
\end{array}%
\right) \left( 
\begin{array}{ccc}
X_{1} & \cdots & X_{2r}%
\end{array}%
\right) \mathcal{R}_{X}\left( T\right) .
\end{equation}

\item The anti commutator matrix $\left\{ \boldsymbol{X},\boldsymbol{X}%
\right\} $ can be expressed as 
\begin{eqnarray}
\left\{ \boldsymbol{X},\boldsymbol{X}\right\} &=&\left( 
\begin{array}{ccc}
X_{1}X_{1} & \cdots & X_{1}X_{2r} \\ 
\vdots & \ddots & \vdots \\ 
X_{2r}X_{1} & \cdots & X_{2r}X_{2r}%
\end{array}%
\right) +\left( 
\begin{array}{ccc}
X_{1}X_{1} & \cdots & X_{2r}X_{1} \\ 
\vdots & \ddots & \vdots \\ 
X_{1}X_{2r} & \cdots & X_{2r}X_{2r}%
\end{array}%
\right) \\
&=&\left( 
\begin{array}{c}
X_{1} \\ 
\vdots \\ 
X_{2r}%
\end{array}%
\right) \left( 
\begin{array}{ccc}
X_{1} & \cdots & X_{2r}%
\end{array}%
\right) +\left( 
\begin{array}{ccc}
X_{1} & \cdots & X_{2r}%
\end{array}%
\right) \left( 
\begin{array}{c}
X_{1} \\ 
\vdots \\ 
X_{2r}%
\end{array}%
\right) .
\end{eqnarray}

\item Transformations $T\in \mathcal{L}\left( \mathcal{F}_{1}\right) $ have
the effect 
\begin{eqnarray}
\left\{ T\boldsymbol{X},T\boldsymbol{X}\right\} &=&T^{t}\left( 
\begin{array}{ccc}
X_{1}X_{1} & \cdots & X_{1}X_{2r} \\ 
\vdots & \ddots & \vdots \\ 
X_{2r}X_{1} & \cdots & X_{2r}X_{2r}%
\end{array}%
\right) T+T^{t}\left( 
\begin{array}{ccc}
X_{1}X_{1} & \cdots & X_{2r}X_{1} \\ 
\vdots & \ddots & \vdots \\ 
X_{1}X_{2r} & \cdots & X_{2r}X_{2r}%
\end{array}%
\right) T \\
&=&T^{t}\left( 
\begin{array}{c}
X_{1} \\ 
\vdots \\ 
X_{2r}%
\end{array}%
\right) \left( 
\begin{array}{ccc}
X_{1} & \cdots & X_{2r}%
\end{array}%
\right) T+T^{t}\left( 
\begin{array}{ccc}
X_{1} & \cdots & X_{2r}%
\end{array}%
\right) \left( 
\begin{array}{c}
X_{1} \\ 
\vdots \\ 
X_{2r}%
\end{array}%
\right) T.
\end{eqnarray}
\end{enumerate}
\end{remark}

\begin{theorem}
Transformation of Tensor Products \label{Transformation of Row Column Tensor
Products} 
\begin{equation}
T^{t}\left( 
\begin{array}{ccc}
X_{1} & \cdots & X_{2r}%
\end{array}%
\right) \left( 
\begin{array}{c}
X_{1} \\ 
\vdots \\ 
X_{2r}%
\end{array}%
\right) T=\left( 
\begin{array}{ccc}
X_{1} & \cdots & X_{2r}%
\end{array}%
\right) TT^{t}\left( 
\begin{array}{c}
X_{1} \\ 
\vdots \\ 
X_{2r}%
\end{array}%
\right) .
\end{equation}
\end{theorem}

\subparagraph{Groups Associated with CAR}

\begin{enumerate}
\item $K\in \mathcal{L}\left( \mathcal{F}\right) $%
\begin{equation}
\left( \boldsymbol{Y}\right) =K\left( \boldsymbol{X}\right) =\left( 
\boldsymbol{X}\right) \mathcal{R}_{X}\left( K\right)
\end{equation}%
in particular 
\begin{equation}
\left( 
\begin{array}{cc}
\boldsymbol{x}_{1} & \boldsymbol{x}_{2}%
\end{array}%
\right) =K\left( 
\begin{array}{cc}
\boldsymbol{a}^{\dagger } & \boldsymbol{a}%
\end{array}%
\right) =\left( 
\begin{array}{cc}
\boldsymbol{a}^{\dagger } & \boldsymbol{a}%
\end{array}%
\right) \frac{1}{\sqrt{2}}\left( 
\begin{array}{cc}
\mathbb{I}_{r} & \mathbb{I}_{r} \\ 
i\mathbb{I}_{r} & -i\mathbb{I}_{r}%
\end{array}%
\right) .
\end{equation}%
Note that 
\begin{equation}
K\in U\left( 2r\right) .
\end{equation}

\begin{enumerate}
\item The CAR\ in terms of a creation and annihilation operator basis is 
\begin{gather}
\left( 
\begin{array}{cc}
\left\{ \boldsymbol{a}^{\dagger },\boldsymbol{a}^{\dagger }\right\} & 
\left\{ \boldsymbol{a}^{\dagger },\boldsymbol{a}\right\} \\ 
\left\{ \boldsymbol{a},\boldsymbol{a}^{\dagger }\right\} & \left\{ 
\boldsymbol{a},\boldsymbol{a}\right\}%
\end{array}%
\right) =\left( 
\begin{array}{c}
\boldsymbol{a}^{\dagger } \\ 
\boldsymbol{a}%
\end{array}%
\right) \otimes \left( 
\begin{array}{cc}
\boldsymbol{a}^{\dagger } & \boldsymbol{a}%
\end{array}%
\right) +\left( 
\begin{array}{cc}
\boldsymbol{a}^{\dagger } & \boldsymbol{a}%
\end{array}%
\right) \otimes \left( 
\begin{array}{c}
\boldsymbol{a}^{\dagger } \\ 
\boldsymbol{a}%
\end{array}%
\right) \\
\parallel \\
\left( 
\begin{array}{c}
\boldsymbol{a}^{\dagger } \\ 
\boldsymbol{a}%
\end{array}%
\right) \left( 
\begin{array}{cc}
\boldsymbol{a}^{\dagger } & \boldsymbol{a}%
\end{array}%
\right) +\left( 
\begin{array}{cc}
\boldsymbol{a}^{\dagger } & \boldsymbol{a}%
\end{array}%
\right) \left( 
\begin{array}{c}
\boldsymbol{a}^{\dagger } \\ 
\boldsymbol{a}%
\end{array}%
\right) \\
\parallel \\
\left( 
\begin{array}{cc}
0_{r} & I_{r} \\ 
I_{r} & 0_{r}%
\end{array}%
\right) ,
\end{gather}%
where 
\begin{equation}
I_{r}=\left( 
\begin{array}{cccc}
I_{\mathcal{F}} & \boldsymbol{0} & \cdots & \boldsymbol{0} \\ 
\boldsymbol{0} & \ddots & \vdots & \vdots \\ 
\vdots & \cdots & \ddots & \boldsymbol{0} \\ 
\boldsymbol{0} & \cdots & \boldsymbol{0} & I_{\mathcal{F}}%
\end{array}%
\right) .
\end{equation}

\item The invariance group $G_{CAR}$ of the CAR in this form is 
\begin{equation}
\left( 
\begin{array}{cc}
\left\{ T\boldsymbol{a}^{\dagger },T\boldsymbol{a}^{\dagger }\right\} & 
\left\{ T\boldsymbol{a}^{\dagger },T\boldsymbol{a}\right\} \\ 
\left\{ T\boldsymbol{a},T\boldsymbol{a}^{\dagger }\right\} & \left\{ T%
\boldsymbol{a},T\boldsymbol{a}\right\}%
\end{array}%
\right) =T^{t}\left( 
\begin{array}{cc}
\left\{ \boldsymbol{a}^{\dagger },\boldsymbol{a}^{\dagger }\right\} & 
\left\{ \boldsymbol{a}^{\dagger },\boldsymbol{a}\right\} \\ 
\left\{ \boldsymbol{a},\boldsymbol{a}^{\dagger }\right\} & \left\{ 
\boldsymbol{a},\boldsymbol{a}\right\}%
\end{array}%
\right) T=T^{t}\left( \left( 
\begin{array}{c}
\boldsymbol{a}^{\dagger } \\ 
\boldsymbol{a}%
\end{array}%
\right) \left( 
\begin{array}{cc}
\boldsymbol{a}^{\dagger } & \boldsymbol{a}%
\end{array}%
\right) +\left( 
\begin{array}{cc}
\boldsymbol{a}^{\dagger } & \boldsymbol{a}%
\end{array}%
\right) \left( 
\begin{array}{c}
\boldsymbol{a}^{\dagger } \\ 
\boldsymbol{a}%
\end{array}%
\right) \right) T=T^{t}\left( 
\begin{array}{cc}
0_{r} & I_{r} \\ 
I_{r} & 0_{r}%
\end{array}%
\right) T=\left( 
\begin{array}{cc}
0_{r} & I_{r} \\ 
I_{r} & 0_{r}%
\end{array}%
\right) ,
\end{equation}%
for $T\in GL\left( 2r,\mathbb{C}\right) .$

\item Note that this condition does not ensure 
\begin{equation}
\left( T\left( a\right) \right) ^{\dagger }\neq T\left( a^{\dagger }\right)
\end{equation}%
i.e. in general $T$ is not a $\ast $-map$.$

\item The transformation $K$ from the creation/annihilation form to the
Majorano form gives 
\begin{equation}
K\left( 
\begin{array}{cc}
0_{r} & I_{r} \\ 
I_{r} & 0_{r}%
\end{array}%
\right) K^{t}=\left( 
\begin{array}{cc}
I_{r} & \mathbb{O}_{r} \\ 
\mathbb{O}_{r} & I_{r}%
\end{array}%
\right)
\end{equation}%
in detail 
\begin{equation}
\frac{1}{\sqrt{2}}\left( 
\begin{array}{cc}
\mathbb{I}_{r} & \mathbb{I}_{r} \\ 
i\mathbb{I}_{r} & -i\mathbb{I}_{r}%
\end{array}%
\right) \left( 
\begin{array}{cc}
0_{r} & I_{r} \\ 
I_{r} & 0_{r}%
\end{array}%
\right) \frac{1}{\sqrt{2}}\left( 
\begin{array}{cc}
\mathbb{I}_{r} & i\mathbb{I}_{r} \\ 
\mathbb{I}_{r} & -i\mathbb{I}_{r}%
\end{array}%
\right) =\frac{1}{2}\left( 
\begin{array}{cc}
\mathbb{I}_{r} & \mathbb{I}_{r} \\ 
i\mathbb{I}_{r} & -i\mathbb{I}_{r}%
\end{array}%
\right) \left( 
\begin{array}{cc}
I_{r} & -iI_{r} \\ 
I_{r} & iI_{r}%
\end{array}%
\right) =\left( 
\begin{array}{cc}
I_{r} & \mathbb{O}_{r} \\ 
\mathbb{O}_{r} & I_{r}%
\end{array}%
\right)
\end{equation}%
one has 
\begin{equation}
T\in G_{CAR}\Leftrightarrow T\in KTK^{t}\in O\left( 2r,\mathbb{C}\right) ,
\end{equation}%
thus 
\begin{equation}
G_{CAR}\cong O\left( 2r,\mathbb{C}\right) .
\end{equation}

\item 
\begin{equation}
K\left( 
\begin{array}{cc}
0_{r} & I_{r} \\ 
I_{r} & 0_{r}%
\end{array}%
\right) K^{\dagger }=\left( 
\begin{array}{cc}
O_{r} & iI_{r} \\ 
-iI_{r} & O_{r}%
\end{array}%
\right)
\end{equation}%
in detail%
\begin{equation}
\frac{1}{\sqrt{2}}\left( 
\begin{array}{cc}
\mathbb{I}_{r} & i\mathbb{I}_{r} \\ 
\mathbb{I}_{r} & -i\mathbb{I}_{r}%
\end{array}%
\right) \left( 
\begin{array}{cc}
0_{r} & I_{r} \\ 
I_{r} & 0_{r}%
\end{array}%
\right) \frac{1}{\sqrt{2}}\left( 
\begin{array}{cc}
\mathbb{I}_{r} & \mathbb{I}_{r} \\ 
-i\mathbb{I}_{r} & i\mathbb{I}_{r}%
\end{array}%
\right) =\frac{1}{2}\left( 
\begin{array}{cc}
iI_{r} & I_{r} \\ 
-iI_{r} & I_{r}%
\end{array}%
\right) \left( 
\begin{array}{cc}
\mathbb{I}_{r} & \mathbb{I}_{r} \\ 
-i\mathbb{I}_{r} & i\mathbb{I}_{r}%
\end{array}%
\right) =\left( 
\begin{array}{cc}
O_{r} & iI_{r} \\ 
-iI_{r} & O_{r}%
\end{array}%
\right) .
\end{equation}

\item The CAR\ in terms of a Majorano operator basis is 
\begin{equation}
\left( 
\begin{array}{cc}
\left\{ \boldsymbol{x}_{1},\boldsymbol{x}_{1}\right\} & \left\{ \boldsymbol{x%
}_{1},\boldsymbol{x}_{2}\right\} \\ 
\left\{ \boldsymbol{x}_{2},\boldsymbol{x}_{1}\right\} & \left\{ \boldsymbol{x%
}_{2},\boldsymbol{x}_{2}\right\}%
\end{array}%
\right) =\left( 
\begin{array}{c}
\boldsymbol{x}_{1} \\ 
\boldsymbol{x}_{2}%
\end{array}%
\right) \left( 
\begin{array}{cc}
\boldsymbol{x}_{1} & \boldsymbol{x}_{2}%
\end{array}%
\right) +\left( 
\begin{array}{cc}
\boldsymbol{x}_{1} & \boldsymbol{x}_{2}%
\end{array}%
\right) \left( 
\begin{array}{c}
\boldsymbol{x}_{1} \\ 
\boldsymbol{x}_{2}%
\end{array}%
\right) =\frac{1}{2}\left( 
\begin{array}{cc}
I_{r} & 0_{r} \\ 
0_{r} & I_{r}%
\end{array}%
\right) .
\end{equation}

\item The invariance group $G_{CAR}$ of the CAR in this form is 
\begin{equation}
T^{t}\left( 
\begin{array}{cc}
\left\{ \boldsymbol{x}_{1},\boldsymbol{x}_{1}\right\} & \left\{ \boldsymbol{x%
}_{1},\boldsymbol{x}_{2}\right\} \\ 
\left\{ \boldsymbol{x}_{2},\boldsymbol{x}_{1}\right\} & \left\{ \boldsymbol{x%
}_{2},\boldsymbol{x}_{2}\right\}%
\end{array}%
\right) T=T^{t}\left( \left( 
\begin{array}{c}
\boldsymbol{x}_{1} \\ 
\boldsymbol{x}_{2}%
\end{array}%
\right) \left( 
\begin{array}{cc}
\boldsymbol{x}_{1} & \boldsymbol{x}_{2}%
\end{array}%
\right) +\left( 
\begin{array}{cc}
\boldsymbol{x}_{1} & \boldsymbol{x}_{2}%
\end{array}%
\right) \left( 
\begin{array}{c}
\boldsymbol{x}_{1} \\ 
\boldsymbol{x}_{2}%
\end{array}%
\right) \right) T=T^{t}\frac{1}{2}\left( 
\begin{array}{cc}
I_{r} & 0_{r} \\ 
0_{r} & I_{r}%
\end{array}%
\right) T=\frac{1}{2}\left( 
\begin{array}{cc}
I_{r} & 0_{r} \\ 
0_{r} & I_{r}%
\end{array}%
\right) ,
\end{equation}%
for $T\in GL\left( 2r,\mathbb{C}\right) $ again showing 
\begin{equation}
G_{CAR}\cong O\left( 2r,\mathbb{C}\right) .
\end{equation}
\end{enumerate}

\item The star transformation condition for $\Upsilon \in \mathcal{L}\left( 
\mathcal{F}\right) $ and $A\in \mathcal{F}$%
\begin{equation}
\Upsilon \left( A^{\dagger }\right) =\Upsilon \left( A\right) ^{\dagger }
\end{equation}
\end{enumerate}

\begin{definition}
\begin{equation}
\limfunc{star}\left\{ \mathcal{L}\left( A\right) \right\} =\left\{ \left.
A\right\vert \Upsilon \left( A^{\dagger }\right) =\Upsilon \left( A\right)
^{\dagger }\forall A\in \mathcal{F}_{1}\right\} .
\end{equation}
\end{definition}

\begin{theorem}
\label{embedding o real matrices as complex matrices} The map $\Omega \in 
\mathcal{L}\left( \mathcal{M}\left( 2r,2r,\mathbb{C}\right) ,\mathcal{M}%
\left( 2r,2r,\mathbb{R}\right) \right) $ 
\begin{equation}
\Omega :\mathcal{M}\left( 2r,2r,\mathbb{R}\right) \rightarrow \mathcal{M}%
\left( 2r,2r,\mathbb{C}\right)
\end{equation}%
defined by 
\begin{equation}
\Omega \left( \left( 
\begin{array}{cc}
\mathbb{A} & \mathbb{B} \\ 
\mathbb{C} & \mathbb{D}%
\end{array}%
\right) \right) =\frac{1}{2}\left( 
\begin{array}{cc}
\mathbb{A+D}-i\left( \mathbb{B-C}\right) & \mathbb{A-D}+i\left( \mathbb{B+C}%
\right) \\ 
\mathbb{A-D}-i\left( \mathbb{B+C}\right) & \mathbb{A+D}+i\left( \mathbb{B-C}%
\right)%
\end{array}%
\right)
\end{equation}
is a linear isomorphism.

\begin{proof}
Let 
\begin{equation}
\Omega \left( \left( 
\begin{array}{cc}
\mathbb{A}_{1} & \mathbb{B}_{1} \\ 
\mathbb{C}_{1} & \mathbb{D}_{1}%
\end{array}%
\right) -\left( 
\begin{array}{cc}
\mathbb{A}_{2} & \mathbb{B}_{2} \\ 
\mathbb{C}_{2} & \mathbb{D}_{2}%
\end{array}%
\right) \right) =0
\end{equation}%
then 
\begin{gather}
\frac{1}{2}\left( 
\begin{array}{cc}
\mathbb{A}_{1}\mathbb{+D}_{1}-i\left( \mathbb{B}_{1}\mathbb{-C}_{1}\right) & 
\mathbb{A}_{1}\mathbb{-D}_{1}+i\left( \mathbb{B}_{1}\mathbb{+C}_{1}\right)
\\ 
\mathbb{A}_{1}\mathbb{-D}_{1}-i\left( \mathbb{B}_{1}\mathbb{+C}_{1}\right) & 
\mathbb{A}_{1}\mathbb{+D}_{1}+i\left( \mathbb{B}_{1}\mathbb{-C}_{1}\right)%
\end{array}%
\right) -\frac{1}{2}\left( 
\begin{array}{cc}
\mathbb{A}_{2}\mathbb{+D}_{2}-i\left( \mathbb{B}_{2}\mathbb{-C}_{2}\right) & 
\mathbb{A}_{2}\mathbb{-D}_{2}+i\left( \mathbb{B}_{2}\mathbb{+C}_{2}\right)
\\ 
\mathbb{A}_{2}\mathbb{-D}_{2}-i\left( \mathbb{B}_{2}\mathbb{+C}_{2}\right) & 
\mathbb{A}_{2}\mathbb{+D}_{2}+i\left( \mathbb{B}_{2}\mathbb{-C}_{2}\right)%
\end{array}%
\right) =\mathbb{O}_{2r} \\
\Updownarrow \\
\mathbb{A}_{1}=\mathbb{A}_{2} \\
\mathbb{B}_{1}=\mathbb{B}_{2} \\
\mathbb{C}_{1}=\mathbb{C}_{2} \\
\mathbb{D}_{1}=\mathbb{D}_{2}
\end{gather}%
hence 
\begin{equation}
\ker \left( \Omega \right) =\mathbb{O}_{2r}.\blacksquare
\end{equation}
\end{proof}
\end{theorem}

\begin{corollary}
\begin{equation}
\Omega \left( \left( 
\begin{array}{cc}
\func{Re}\left( \mathbb{U+V}\right) & -\func{Im}\left( \mathbb{U+V}\right)
\\ 
\func{Im}\left( \mathbb{U-V}\right) & \func{Re}\left( \mathbb{U-V}\right)%
\end{array}%
\right) \right) =\left( 
\begin{array}{cc}
\mathbb{U} & \overline{\mathbb{V}} \\ 
\mathbb{V} & \overline{\mathbb{U}}%
\end{array}%
\right) .
\end{equation}

\begin{proof}
define 
\begin{equation}
\left( 
\begin{array}{cc}
\func{Re}\left( \mathbb{U+V}\right) & -\func{Im}\left( \mathbb{U+V}\right)
\\ 
\func{Im}\left( \mathbb{U-V}\right) & \func{Re}\left( \mathbb{U-V}\right)%
\end{array}%
\right) =\left( 
\begin{array}{cc}
\mathbb{A} & \mathbb{B} \\ 
\mathbb{C} & \mathbb{D}%
\end{array}%
\right)
\end{equation}%
then%
\begin{eqnarray}
\mathbb{A} &=&\func{Re}\left( \mathbb{U+V}\right) \\
\mathbb{B} &=&-\func{Im}\left( \mathbb{U+V}\right) \\
\mathbb{C} &=&\func{Im}\left( \mathbb{U-V}\right) \\
\mathbb{D} &=&\func{Re}\left( \mathbb{U-V}\right)
\end{eqnarray}%
and 
\begin{eqnarray}
\mathbb{A+D} &=&2\func{Re}\left( \mathbb{U}\right) \\
\mathbb{A-D} &=&2\func{Re}\left( \mathbb{V}\right) \\
\mathbb{B+C} &=&-2\func{Im}\left( \mathbb{V}\right) \\
\mathbb{B-C} &=&-2\func{Im}\left( \mathbb{U}\right)
\end{eqnarray}%
\begin{equation}
\frac{1}{2}\left( 
\begin{array}{cc}
\mathbb{A+D}-i\left( \mathbb{B-C}\right) & \mathbb{A-D}+i\left( \mathbb{B+C}%
\right) \\ 
\mathbb{A-D}-i\left( \mathbb{B+C}\right) & \mathbb{A+D}+i\left( \mathbb{B-C}%
\right)%
\end{array}%
\right) =\left( 
\begin{array}{cc}
\mathbb{U} & \overline{\mathbb{V}} \\ 
\mathbb{V} & \overline{\mathbb{U}}%
\end{array}%
\right) .\blacksquare
\end{equation}
\end{proof}
\end{corollary}

\begin{theorem}
If $\Upsilon \left( A^{\dagger }\right) =\Upsilon \left( A\right) ^{\dagger
} $ $\forall A\in \mathcal{F}_{1}$ then 
\begin{equation}
\mathcal{R}_{a}\left( A\right) =\left( 
\begin{array}{cc}
\mathcal{R}_{a}\left( \Upsilon \right) _{11} & \overline{\mathcal{R}%
_{a}\left( \Upsilon \right) _{12}} \\ 
\mathcal{R}_{a}\left( \Upsilon \right) _{12} & \overline{\mathcal{R}%
_{a}\left( \Upsilon \right) _{11}}%
\end{array}%
\right)
\end{equation}%
where 
\begin{equation}
\Upsilon \left( 
\begin{array}{cc}
\boldsymbol{a}^{\dagger } & \boldsymbol{a}%
\end{array}%
\right) =\left( 
\begin{array}{cc}
\boldsymbol{a}^{\dagger } & \boldsymbol{a}%
\end{array}%
\right) \mathcal{R}_{a}\left( \Upsilon \right) =\left( 
\begin{array}{cc}
\boldsymbol{a}^{\dagger } & \boldsymbol{a}%
\end{array}%
\right) \left( 
\begin{array}{cc}
\mathcal{R}_{a}\left( \Upsilon \right) _{11} & \overline{\mathcal{R}%
_{a}\left( \Upsilon \right) _{12}} \\ 
\mathcal{R}_{a}\left( \Upsilon \right) _{12} & \overline{\mathcal{R}%
_{a}\left( \Upsilon \right) _{11}}%
\end{array}%
\right) .
\end{equation}

\begin{proof}
\begin{equation}
\end{equation}
\end{proof}
\end{theorem}

\begin{theorem}
Bogolubov Transformation \label{Bogolubov Transformation Theorem} 
\begin{equation}
O\left( 2r,\mathbb{C}\right) \cap \limfunc{star}\left\{ \mathcal{L}\left(
A\right) \right\} =O\left( 2r,\mathbb{R}\right)
\end{equation}

\begin{proof}
\begin{equation}
\end{equation}
\end{proof}
\end{theorem}

\subparagraph{The Tensor Product of Matrices of Any Dimension \label{The
Tensor Product of Matrices of Any Dimension}}

The tensor product of two matrices $A\in \mathcal{M}\left( P\ ,Q\ ,\mathcal{F%
}\right) $ and $B\in \mathcal{M}\left( M\ ,N\ ,\mathcal{F}\right) $ that
have Clifford algebra $\mathcal{F}$ entries 
\begin{gather}
A\otimes B=\left( 
\begin{array}{ccc}
A_{11}B & \cdots & A_{1Q}B \\ 
\vdots & \ddots & \vdots \\ 
A_{P1}B & \cdots & A_{PQ}B%
\end{array}%
\right) =\left( 
\begin{array}{ccc}
A_{11}\left( 
\begin{array}{ccc}
B_{11} & \cdots & B_{1N} \\ 
\vdots & \ddots & \vdots \\ 
B_{M1} & \cdots & B_{MN}%
\end{array}%
\right) & \cdots & A_{1Q}\left( 
\begin{array}{ccc}
B_{11} & \cdots & B_{1N} \\ 
\vdots & \ddots & \vdots \\ 
B_{M1} & \cdots & B_{MN}%
\end{array}%
\right) \\ 
\vdots & \ddots & \vdots \\ 
A_{P1}\left( 
\begin{array}{ccc}
B_{11} & \cdots & B_{1N} \\ 
\vdots & \ddots & \vdots \\ 
B_{M1} & \cdots & B_{MN}%
\end{array}%
\right) & \cdots & A_{PQ}\left( 
\begin{array}{ccc}
B_{11} & \cdots & B_{1N} \\ 
\vdots & \ddots & \vdots \\ 
B_{M1} & \cdots & B_{MN}%
\end{array}%
\right)%
\end{array}%
\right) \\
A\in \mathcal{M}\left( P\ ,Q\ ,\mathcal{F}\right) ,\quad B\in \mathcal{M}%
\left( M\ ,N\ ,\mathcal{F}\right) ,\quad A\otimes B\in \mathcal{M}\left(
PM,QN,\mathcal{F}\right)
\end{gather}

\begin{remark}
\begin{enumerate}
\item In particular consider the tensor product of row super vectors $%
P=1,Q=2r$ with column super vectors with $M=2r$ and $N=1$ i.e. 
\begin{equation}
\left( P,Q,M,N\right) =\left( 1,2r,2r,1\right)
\end{equation}%
then

\begin{enumerate}
\item $\boldsymbol{A}=\left( \boldsymbol{a}\right) ^{t}\in \mathcal{M}\left(
1\ ,2r\ ,\mathcal{F}_{1}\right) $ 
\begin{equation}
\boldsymbol{A}=\left( 
\begin{array}{ccc}
a_{1\ } & \cdots & a_{2r\ }%
\end{array}%
\right) =\left( \boldsymbol{a}\right) ^{t}
\end{equation}%
and

\item $\boldsymbol{B}=\left( \boldsymbol{b}\right) \in \mathcal{M}\left( 2r\
,1\ ,\mathcal{F}_{1}\right) $ 
\begin{equation}
\boldsymbol{B}=\left( 
\begin{array}{c}
b_{1} \\ 
\vdots \\ 
b_{2r}%
\end{array}%
\right) =\left( \boldsymbol{b}\right) ,
\end{equation}%
which gives

\item $\boldsymbol{A}\otimes \boldsymbol{B}=\left( \boldsymbol{a}\right)
^{t}\otimes \left( \boldsymbol{b}\right) \in \mathcal{M}\left( 2r,2r,%
\mathcal{F}_{1}\right) $ that can be evaluated as 
\begin{gather}
\boldsymbol{A}\otimes \boldsymbol{B}=\left( 
\begin{array}{ccc}
A_{11} & \cdots & A_{2r1}%
\end{array}%
\right) \otimes \left( 
\begin{array}{c}
B_{11} \\ 
\vdots \\ 
B_{2r1}%
\end{array}%
\right) =\left( 
\begin{array}{ccc}
A_{11}\left( 
\begin{array}{c}
B_{11} \\ 
\vdots \\ 
B_{2r1}%
\end{array}%
\right) & \cdots & A_{12r}\left( 
\begin{array}{c}
B_{11} \\ 
\vdots \\ 
B_{2r1}%
\end{array}%
\right)%
\end{array}%
\right) =\left( 
\begin{array}{ccc}
A_{11}B_{11} & \cdots & A_{12r}B_{11} \\ 
\vdots & \ddots & \vdots \\ 
A_{11}B_{2r1} & \cdots & A_{12r}B_{2r1}%
\end{array}%
\right) \\
=\left( \boldsymbol{a}\right) ^{t}\otimes \left( \boldsymbol{b}\right)
=\left( 
\begin{array}{ccc}
a_{1\ } & \cdots & a_{2r\ }%
\end{array}%
\right) \otimes \left( 
\begin{array}{c}
b_{1} \\ 
\vdots \\ 
b_{2r}%
\end{array}%
\right) =\left( 
\begin{array}{ccc}
a_{1}b_{1} & \cdots & a_{2r}b_{1} \\ 
\vdots & \ddots & \vdots \\ 
a_{1}b_{2r} & \cdots & a_{2r}b_{2r}%
\end{array}%
\right) .
\end{gather}

\item In particular in the case $r=1$%
\begin{equation}
\left( P,Q,M,N\right) =\left( 2,1,1,2\right)
\end{equation}%
one obtains 
\begin{equation}
\left( \boldsymbol{a}\right) ^{t}\otimes \left( \boldsymbol{b}\right)
=\left( 
\begin{array}{cc}
a_{1} & a_{2}%
\end{array}%
\right) \otimes \left( 
\begin{array}{c}
b_{1} \\ 
b_{2}%
\end{array}%
\right) =\left( 
\begin{array}{cc}
a_{1}\left( 
\begin{array}{c}
b_{1} \\ 
b_{2}%
\end{array}%
\right) & a_{2}\left( 
\begin{array}{c}
b_{1} \\ 
b_{2}%
\end{array}%
\right)%
\end{array}%
\right) =\left( 
\begin{array}{cc}
a_{1}b_{1} & a_{2}b_{1} \\ 
a_{1}b_{2} & a_{2}b_{2}%
\end{array}%
\right)
\end{equation}%
that gives the\emph{\ transpose }expression

\item 
\begin{eqnarray}
\left( 
\begin{array}{cc}
a_{1}b_{1} & a_{2}b_{1} \\ 
a_{1}b_{2} & a_{2}b_{2}%
\end{array}%
\right) ^{t} &=&\left( 
\begin{array}{cc}
\left( a_{1}b_{1}\right) ^{t} & \left( a_{1}b_{2}\right) ^{t} \\ 
\left( a_{2}b_{1}\right) ^{t} & \left( a_{2}b_{2}\right) ^{t}%
\end{array}%
\right) =\left( \overset{}{\left( \boldsymbol{a}\right) ^{t}\otimes \left( 
\boldsymbol{b}\right) }\right) ^{t}=\left( \boldsymbol{b}\right) ^{t}\otimes
\left( \boldsymbol{a}\right) =\left( 
\begin{array}{cc}
b_{1} & b_{2}%
\end{array}%
\right) \otimes \left( 
\begin{array}{c}
a_{1} \\ 
a_{2}%
\end{array}%
\right) \\
&=&\left( 
\begin{array}{cc}
b_{1}\left( 
\begin{array}{c}
a_{1} \\ 
a_{2}%
\end{array}%
\right) & b_{2}\left( 
\begin{array}{c}
a_{1} \\ 
a_{2}%
\end{array}%
\right)%
\end{array}%
\right) =\left( 
\begin{array}{cc}
b_{1}a_{1} & b_{2}a_{1} \\ 
b_{1}a_{2} & b_{2}a_{2}%
\end{array}%
\right) .
\end{eqnarray}
\end{enumerate}

\item In particular in the case $P=2r,Q=1$ and $M=1$ and $N=2r$ i.e. 
\begin{equation}
\left( P,Q,M,N\right) =\left( 2r,1,1,2r\right)
\end{equation}

\begin{enumerate}
\item $\boldsymbol{A}=\left( \boldsymbol{a}\right) \in \mathcal{M}\left( 1\
,2r\ ,\mathcal{F}_{1}\right) $ 
\begin{equation}
\boldsymbol{A}=\left( 
\begin{array}{c}
a_{1} \\ 
\vdots \\ 
a_{2r}%
\end{array}%
\right) =\left( \boldsymbol{a}\right)
\end{equation}%
and

\item $\boldsymbol{B}=\left( \boldsymbol{b}\right) ^{t}\in \mathcal{M}\left(
2r\ ,1\ ,\mathcal{F}_{1}\right) $ 
\begin{equation}
\boldsymbol{B}=\left( 
\begin{array}{ccc}
b_{1\ } & \cdots & b_{2r\ }%
\end{array}%
\right) =\left( \boldsymbol{b}\right) ^{t},
\end{equation}%
which gives

\item the tensor product $A\otimes B\in \mathcal{M}\left( 2r,2r,\mathcal{F}%
_{1}\right) $ 
\begin{eqnarray}
\boldsymbol{A}\otimes \boldsymbol{B} &=&\left( 
\begin{array}{c}
A_{11} \\ 
\vdots \\ 
A_{2r1}%
\end{array}%
\right) \otimes \left( 
\begin{array}{ccc}
B_{11} & \cdots & B_{\ 12r}%
\end{array}%
\right) =\left( 
\begin{array}{c}
A_{11}\left( 
\begin{array}{ccc}
B_{11} & \cdots & B_{\ 12r}%
\end{array}%
\right) \\ 
\vdots \\ 
A_{2r1}\left( 
\begin{array}{ccc}
B_{11} & \cdots & B_{\ 12r}%
\end{array}%
\right)%
\end{array}%
\right) =\left( 
\begin{array}{ccc}
A_{11}B_{11} & \cdots & A_{11}B_{\ 12r} \\ 
\vdots & \ddots & \vdots \\ 
A_{2r1}B_{2r1} & \cdots & A_{2r1}B_{\ 12r}%
\end{array}%
\right) \\
&=&\left( \boldsymbol{a}\right) \otimes \left( \boldsymbol{b}\right)
^{t}=\left( 
\begin{array}{c}
a_{1} \\ 
\vdots \\ 
a_{2r}%
\end{array}%
\right) \otimes \left( 
\begin{array}{ccc}
b_{1\ } & \cdots & b_{2r\ }%
\end{array}%
\right) =\left( 
\begin{array}{ccc}
a_{1}b_{1} & \cdots & a_{2r}b_{1} \\ 
\vdots & \ddots & \vdots \\ 
a_{1}b_{2r} & \cdots & a_{2r}b_{2r}%
\end{array}%
\right)
\end{eqnarray}

\item In particular consider the tensor product of column super vectors
where $P=2,Q=1$ with row super vectors $M=1$ and $N=2$ i.e. 
\begin{equation}
\left( P,Q,M,N\right) =\left( 2,1,1,2\right)
\end{equation}%
then%
\begin{equation}
\left( \boldsymbol{a}\right) \otimes \left( \boldsymbol{b}\right)
^{t}=\left( 
\begin{array}{c}
a_{1} \\ 
a_{2}%
\end{array}%
\right) \otimes \left( 
\begin{array}{cc}
b_{1} & b_{\ 2}%
\end{array}%
\right) =\left( 
\begin{array}{c}
a_{1}\left( 
\begin{array}{cc}
b_{1} & b_{\ 2}%
\end{array}%
\right) \\ 
a_{2}\left( 
\begin{array}{cc}
b_{1} & b_{\ 2}%
\end{array}%
\right)%
\end{array}%
\right) =\left( 
\begin{array}{cc}
a_{1}b_{1} & a_{1}b_{2} \\ 
a_{2}b_{1} & a_{2}b_{2}%
\end{array}%
\right) .
\end{equation}%
that gives the\emph{\ transpose }expression

\item 
\begin{eqnarray}
\left( 
\begin{array}{cc}
a_{1}b_{1} & a_{1}b_{2} \\ 
a_{2}b_{1} & a_{2}b_{2}%
\end{array}%
\right) ^{t} &=&\left( 
\begin{array}{cc}
\left( a_{1}b_{1}\right) ^{t} & \left( a_{2}b_{1}\right) ^{t} \\ 
\left( a_{1}b_{2}\right) ^{t} & \left( a_{2}b_{2}\right) ^{t}%
\end{array}%
\right) =\left( \overset{}{\left( \boldsymbol{a}\right) \otimes \left( 
\boldsymbol{b}\right) ^{t}}\right) ^{t}=\left( \boldsymbol{b}\right) \otimes
\left( \boldsymbol{a}\right) ^{t} \\
&=&\left( 
\begin{array}{c}
b_{1} \\ 
b_{2}%
\end{array}%
\right) \otimes \left( 
\begin{array}{cc}
a_{1} & a_{2}%
\end{array}%
\right) =\left( 
\begin{array}{c}
b_{1}\left( 
\begin{array}{cc}
a_{1} & a_{\ 2}%
\end{array}%
\right) \\ 
b_{2}\left( 
\begin{array}{cc}
a_{1} & a_{\ 2}%
\end{array}%
\right)%
\end{array}%
\right) =\left( 
\begin{array}{cc}
b_{1}a_{1} & b_{1}a_{2} \\ 
b_{2}a_{1} & b_{2}a_{2}%
\end{array}%
\right) .
\end{eqnarray}
\end{enumerate}
\end{enumerate}
\end{remark}

\begin{summary}
For the case $r=1$

\begin{enumerate}
\item 
\begin{equation}
\left( \boldsymbol{a}\right) ^{t}\otimes \left( \boldsymbol{b}\right)
=\left( 
\begin{array}{cc}
a_{1} & a_{2}%
\end{array}%
\right) \otimes \left( 
\begin{array}{c}
b_{1} \\ 
b_{2}%
\end{array}%
\right) =\left( 
\begin{array}{cc}
a_{1}\left( 
\begin{array}{c}
b_{1} \\ 
b_{2}%
\end{array}%
\right) & a_{2}\left( 
\begin{array}{c}
b_{1} \\ 
b_{2}%
\end{array}%
\right)%
\end{array}%
\right) =\left( 
\begin{array}{cc}
a_{1}b_{1} & a_{2}b_{1} \\ 
a_{1}b_{2} & a_{2}b_{2}%
\end{array}%
\right)
\end{equation}

\item 
\begin{equation}
\left( \boldsymbol{b}\right) \otimes \left( \boldsymbol{a}\right)
^{t}=\left( 
\begin{array}{c}
b_{1} \\ 
b_{2}%
\end{array}%
\right) \otimes \left( 
\begin{array}{cc}
a_{1} & a_{2}%
\end{array}%
\right) =\left( 
\begin{array}{c}
b_{1}\left( 
\begin{array}{cc}
a_{1} & a_{\ 2}%
\end{array}%
\right) \\ 
b_{2}\left( 
\begin{array}{cc}
a_{1} & a_{\ 2}%
\end{array}%
\right)%
\end{array}%
\right) =\left( 
\begin{array}{cc}
b_{1}a_{1} & b_{1}a_{2} \\ 
b_{2}a_{1} & b_{2}a_{2}%
\end{array}%
\right)
\end{equation}
\end{enumerate}
\end{summary}

\subparagraph{Transformation of Tensor Products \label{Transformation of
Tensor Products}}

Consider matrices

\begin{enumerate}
\item 
\begin{enumerate}
\item 
\begin{equation}
A\in \mathcal{L}\left( \mathcal{F}^{P},\mathcal{F}^{Q}\right) \cong \mathcal{%
M}\left( P\ ,Q\ ,\mathcal{F}\right)
\end{equation}
and

\item 
\begin{equation}
B\in \mathcal{L}\left( \mathcal{F}^{M},\mathcal{F}^{N}\right) \cong \mathcal{%
M}\left( M\ ,N\ ,\mathcal{F}\right) .
\end{equation}
\end{enumerate}

\item Consider transformations of matrices

\begin{enumerate}
\item 
\begin{equation}
S\in \mathcal{L}\left( \mathcal{L}\left( \mathcal{F}^{P},\mathcal{F}%
^{Q}\right) \ \right)
\end{equation}%
and

\item 
\begin{equation}
T\in \mathcal{L}\left( \mathcal{L}\left( \mathcal{F}^{M},\mathcal{F}%
^{N}\right) \ \right) .
\end{equation}
\end{enumerate}

\item Then
\end{enumerate}

\begin{equation}
\left( SA\right) \otimes \left( TB\right) =\left( S\otimes T\right) \left(
A\otimes B\right)
\end{equation}%
where 
\begin{equation}
\left( S\otimes T\right) \in \mathcal{L}\left( \mathcal{F}^{P}\otimes 
\mathcal{F}^{M},\mathcal{F}^{Q}\otimes \mathcal{F}^{N}\right) \cong \mathcal{%
M}\left( PM\ ,QN\ ,\mathcal{F}\right) .
\end{equation}

\begin{example}
$r=1$%
\begin{eqnarray}
\left( S\boldsymbol{a}\right) ^{t}\otimes \left( T\boldsymbol{b}\right)
&=&S\left( \boldsymbol{a}\right) ^{t}\otimes T\left( \boldsymbol{b}\right)
=\left( S\otimes T\right) \left( \overset{}{\left( \boldsymbol{a}\right)
^{t}\otimes \left( \boldsymbol{b}\right) }\right) =\left( 
\begin{array}{cc}
Sa_{1} & Sa_{2}%
\end{array}%
\right) \otimes \left( 
\begin{array}{c}
Tb_{1} \\ 
Tb_{2}%
\end{array}%
\right) \\
&=&\left( 
\begin{array}{cc}
Sa_{1}\left( 
\begin{array}{c}
Tb_{1} \\ 
Tb_{2}%
\end{array}%
\right) & Sa_{2}\left( 
\begin{array}{c}
Tb_{1} \\ 
Tb_{2}%
\end{array}%
\right)%
\end{array}%
\right) =\left( 
\begin{array}{cc}
S\left( a_{1}\right) T\left( b_{1}\right) & S\left( a_{2}\right) T\left(
b_{1}\right) \\ 
S\left( a_{1}\right) T\left( b_{2}\right) & S\left( a_{2}\right) T\left(
b_{2}\right)%
\end{array}%
\right) .
\end{eqnarray}%
\begin{eqnarray}
&&\left( 
\begin{array}{cc}
Sa_{1} & Sa_{2}%
\end{array}%
\right) \otimes \left( 
\begin{array}{c}
Tb_{1} \\ 
Tb_{2}%
\end{array}%
\right) =\left( 
\begin{array}{cc}
a_{1} & a_{2}%
\end{array}%
\right) \mathcal{R}_{a}\left( S\right) \otimes \mathcal{R}_{b}\left(
T\right) ^{t}\left( 
\begin{array}{c}
b_{1} \\ 
b_{2}%
\end{array}%
\right) \\
&=&\left( 
\begin{array}{cc}
a_{1} & a_{2}%
\end{array}%
\right) \left( 
\begin{array}{cc}
\mathcal{R}_{a}\left( S\right) _{11} & \mathcal{R}_{a}\left( S\right) _{12}
\\ 
\mathcal{R}_{a}\left( S\right) _{21} & \mathcal{R}_{a}\left( S\right) _{22}%
\end{array}%
\right) \otimes \left( 
\begin{array}{cc}
\mathcal{R}_{b}\left( T\right) _{11} & \mathcal{R}_{b}\left( T\right) _{21}
\\ 
\mathcal{R}_{b}\left( T\right) _{12} & \mathcal{R}_{b}\left( T\right) _{22}%
\end{array}%
\right) \left( 
\begin{array}{c}
b_{1} \\ 
b_{2}%
\end{array}%
\right) \\
&=&\left( 
\begin{array}{cc}
a_{1}\mathcal{R}_{a}\left( S\right) _{11}+a_{2}\mathcal{R}_{a}\left(
S\right) _{21} & a_{1}\mathcal{R}_{a}\left( S\right) _{12}+a_{2}\mathcal{R}%
_{a}\left( S\right) _{22}%
\end{array}%
\right) \otimes \left( 
\begin{array}{c}
\mathcal{R}_{b}\left( T\right) _{11}b_{1}+\mathcal{R}_{b}\left( T\right)
_{21}b_{2} \\ 
\mathcal{R}_{b}\left( T\right) _{12}b_{1}+\mathcal{R}_{b}\left( T\right)
_{22}b_{2}%
\end{array}%
\right) \\
&=&\left( 
\begin{array}{cc}
\left( \overset{}{a_{1}\mathcal{R}_{a}\left( S\right) _{11}+a_{2}\mathcal{R}%
_{a}\left( S\right) _{21}}\right) \left( 
\begin{array}{c}
\mathcal{R}_{b}\left( T\right) _{11}b_{1}+\mathcal{R}_{b}\left( T\right)
_{21}b_{2} \\ 
\mathcal{R}_{b}\left( T\right) _{12}b_{1}+\mathcal{R}_{b}\left( T\right)
_{22}b_{2}%
\end{array}%
\right) & \left( \overset{}{a_{1}\mathcal{R}_{a}\left( S\right) _{12}+a_{2}%
\mathcal{R}_{a}\left( S\right) _{22}}\right) \left( 
\begin{array}{c}
\mathcal{R}_{b}\left( T\right) _{11}b_{1}+\mathcal{R}_{b}\left( T\right)
_{21}b_{2} \\ 
\mathcal{R}_{b}\left( T\right) _{12}b_{1}+\mathcal{R}_{b}\left( T\right)
_{22}b_{2}%
\end{array}%
\right)%
\end{array}%
\right) \\
&=&\left( 
\begin{array}{cc}
\left( 
\begin{array}{c}
\left( \overset{}{a_{1}\mathcal{R}_{a}\left( S\right) _{11}+a_{2}\mathcal{R}%
_{a}\left( S\right) _{21}}\right) \left( \overset{}{\mathcal{R}_{b}\left(
T\right) _{11}b_{1}+\mathcal{R}_{b}\left( T\right) _{21}b_{2}}\right) \\ 
\left( \overset{}{a_{1}\mathcal{R}_{a}\left( S\right) _{11}+a_{2}\mathcal{R}%
_{a}\left( S\right) _{21}}\right) \left( \overset{}{\mathcal{R}_{b}\left(
T\right) _{12}b_{1}+\mathcal{R}_{b}\left( T\right) _{22}b_{2}}\right)%
\end{array}%
\right) & \left( 
\begin{array}{c}
\left( \overset{}{a_{1}\mathcal{R}_{a}\left( S\right) _{12}+a_{2}\mathcal{R}%
_{a}\left( S\right) _{22}}\right) \left( \overset{}{\mathcal{R}_{b}\left(
T\right) _{11}b_{1}+\mathcal{R}_{b}\left( T\right) _{21}b_{2}}\right) \\ 
\left( \overset{}{a_{1}\mathcal{R}_{a}\left( S\right) _{12}+a_{2}\mathcal{R}%
_{a}\left( S\right) _{22}}\right) \left( \overset{}{\mathcal{R}_{b}\left(
T\right) _{12}b_{1}+\mathcal{R}_{b}\left( T\right) _{22}b_{2}}\right)%
\end{array}%
\right)%
\end{array}%
\right)
\end{eqnarray}%
\begin{eqnarray}
&=&\left( 
\begin{array}{c}
\overset{}{a_{1}b_{1}\mathcal{R}_{a}\left( S\right) _{11}\overset{}{\mathcal{%
R}_{b}\left( T\right) _{11}+a_{1}b_{2}\mathcal{R}_{a}\left( S\right) _{11}%
\mathcal{R}_{b}\left( T\right) _{21}}+\,a_{2}b_{1}\mathcal{R}_{a}\left(
S\right) _{21}}\overset{}{\mathcal{R}_{b}\left( T\right) _{11}+a_{2}b_{2}%
\mathcal{R}_{a}\left( S\right) _{21}\mathcal{R}_{b}\left( T\right) _{21}} \\ 
\overset{}{a_{1}b_{1}\mathcal{R}_{a}\left( S\right) _{11}\overset{}{\mathcal{%
R}_{b}\left( T\right) _{12}+a_{1}b_{2}\mathcal{R}_{a}\left( S\right) _{11}%
\mathcal{R}_{b}\left( T\right) _{22}}+\,a_{2}b_{1}\mathcal{R}_{a}\left(
S\right) _{21}}\overset{}{\mathcal{R}_{b}\left( T\right) _{12}+a_{2}b_{2}%
\mathcal{R}_{a}\left( S\right) _{21}\mathcal{R}_{b}\left( T\right) _{22}}%
\end{array}%
\right. \\
&&\left. 
\begin{array}{c}
\overset{}{a_{1}b_{1}\mathcal{R}_{a}\left( S\right) _{12}\overset{}{\mathcal{%
R}_{b}\left( T\right) _{11}+a_{1}b_{2}\mathcal{R}_{a}\left( S\right) _{12}%
\mathcal{R}_{b}\left( T\right) _{21}}+\,a_{2}b_{1}\mathcal{R}_{a}\left(
S\right) _{22}}\overset{}{\mathcal{R}_{b}\left( T\right) _{11}+a_{2}b_{2}%
\mathcal{R}_{a}\left( S\right) _{22}\mathcal{R}_{b}\left( T\right) _{21}} \\ 
\overset{}{a_{1}b_{1}\mathcal{R}_{a}\left( S\right) _{12}\overset{}{\mathcal{%
R}_{b}\left( T\right) _{12}+a_{1}b_{2}\mathcal{R}_{a}\left( S\right) _{12}%
\mathcal{R}_{b}\left( T\right) _{22}}+\,a_{2}b_{1}\mathcal{R}_{a}\left(
S\right) _{22}}\overset{}{\mathcal{R}_{b}\left( T\right) _{12}+a_{2}b_{2}%
\mathcal{R}_{a}\left( S\right) _{22}\mathcal{R}_{b}\left( T\right) _{22}}%
\end{array}%
\right)
\end{eqnarray}%
\begin{equation}
\left( 
\begin{array}{cc}
b_{1}a_{1} & b_{1}a_{2} \\ 
b_{2}a_{1} & b_{2}a_{2}%
\end{array}%
\right)
\end{equation}%
\begin{eqnarray}
&&\left( 
\begin{array}{cc}
\left( 
\begin{array}{c}
\left( \mathcal{R}_{a}\left( S\right) _{11}a_{1}+\mathcal{R}_{a}\left(
S\right) _{21}a_{2}\right) \left( b_{1}\mathcal{R}_{b}\left( T\right)
_{11}+b_{2}\mathcal{R}_{b}\left( T\right) _{21}\right) \\ 
\left( \mathcal{R}_{a}\left( S\right) _{11}a_{1}+\mathcal{R}_{a}\left(
S\right) _{21}a_{2}\right) \left( b_{1}\mathcal{R}_{b}\left( T\right)
_{12}+b_{2}\mathcal{R}_{b}\left( T\right) _{22}\right)%
\end{array}%
\right) & \left( 
\begin{array}{c}
\left( \mathcal{R}_{a}\left( S\right) _{12}a_{1}+\mathcal{R}_{a}\left(
S\right) _{22}a_{2}\right) \left( \mathcal{R}_{b}\left( T\right) _{11}b_{1}+%
\mathcal{R}_{b}\left( T\right) _{21}b_{2}\right) \\ 
\left( \mathcal{R}_{a}\left( S\right) _{12}a_{1}+\mathcal{R}_{a}\left(
S\right) _{22}a_{2}\right) \left( b_{1}\mathcal{R}_{b}\left( T\right)
_{12}+b_{2}\mathcal{R}_{b}\left( T\right) _{22}\right)%
\end{array}%
\right)%
\end{array}%
\right) \\
&=&\left( 
\begin{array}{cc}
\left( 
\begin{array}{c}
\left( \mathcal{R}_{a}\left( S\right) _{11}a_{1}+\mathcal{R}_{a}\left(
S\right) _{21}a_{2}\right) \left( b_{1}\mathcal{R}_{b}\left( T\right)
_{11}+b_{2}\mathcal{R}_{b}\left( T\right) _{21}\right) \\ 
\left( \mathcal{R}_{a}\left( S\right) _{11}a_{1}+\mathcal{R}_{a}\left(
S\right) _{21}a_{2}\right) \left( b_{1}\mathcal{R}_{b}\left( T\right)
_{12}+b_{2}\mathcal{R}_{b}\left( T\right) _{22}\right)%
\end{array}%
\right) & \left( 
\begin{array}{c}
\left( \mathcal{R}_{a}\left( S\right) _{12}a_{1}+\mathcal{R}_{a}\left(
S\right) _{22}a_{2}\right) \left( \mathcal{R}_{b}\left( T\right) _{11}b_{1}+%
\mathcal{R}_{b}\left( T\right) _{21}b_{2}\right) \\ 
\left( \mathcal{R}_{a}\left( S\right) _{12}a_{1}+\mathcal{R}_{a}\left(
S\right) _{22}a_{2}\right) \left( b_{1}\mathcal{R}_{b}\left( T\right)
_{12}+b_{2}\mathcal{R}_{b}\left( T\right) _{22}\right)%
\end{array}%
\right)%
\end{array}%
\right)
\end{eqnarray}
\end{example}

\subparagraph{Tensor Product Form of the Canonical Anti Commutator (CAR) 
\label{Tensor Product Form of the Canonical Anti Commutator (CAR)}}

\begin{remark}
Let $r=1$ then 
\begin{eqnarray}
\left( \boldsymbol{a}\right) ^{t}\otimes \left( \boldsymbol{b}\right)
+\left( \boldsymbol{b}\right) \otimes \left( \boldsymbol{a}\right) ^{t}
&=&\left( 
\begin{array}{cc}
a_{1} & a_{2}%
\end{array}%
\right) \otimes \left( 
\begin{array}{c}
b_{1} \\ 
b_{2}%
\end{array}%
\right) +\left( 
\begin{array}{c}
b_{1} \\ 
b_{2}%
\end{array}%
\right) \otimes \left( 
\begin{array}{cc}
a_{1} & a_{2}%
\end{array}%
\right) \\
&=&\left( 
\begin{array}{cc}
a_{1}b_{1} & a_{2}b_{1} \\ 
a_{1}b_{2} & a_{2}b_{2}%
\end{array}%
\right) +\left( 
\begin{array}{cc}
b_{1}a_{1} & b_{1}a_{2} \\ 
b_{2}a_{1} & b_{2}a_{2}%
\end{array}%
\right) =\left( 
\begin{array}{cc}
\left\{ a_{1},b_{1}\right\} & \left\{ a_{2},b_{1}\right\} \\ 
\left\{ a_{1},b_{2}\right\} & \left\{ b_{2},b_{2}\right\}%
\end{array}%
\right)
\end{eqnarray}
\end{remark}

\begin{remark}
Transformations $T\in \mathcal{L}\left( \mathcal{F}_{1}\right) $ have the
following effect on the Anti Commutator of $\left( \boldsymbol{a}\right) $
and $\left( \boldsymbol{b}\right) $

\begin{enumerate}
\item 
\begin{equation}
T\left( \boldsymbol{a}\right) ^{t}\otimes T\left( \boldsymbol{b}\right)
+T\left( \boldsymbol{b}\right) \otimes T\left( \boldsymbol{a}\right)
^{t}=\left( 
\begin{array}{cc}
Ta_{1} & Ta_{2}%
\end{array}%
\right) \otimes \left( 
\begin{array}{c}
Tb_{1} \\ 
Tb_{2}%
\end{array}%
\right) +\left( 
\begin{array}{c}
Tb_{1} \\ 
Tb_{2}%
\end{array}%
\right) \otimes \left( 
\begin{array}{cc}
Ta_{1} & Ta_{\ 2}%
\end{array}%
\right)
\end{equation}

\item In particular if $\left( \boldsymbol{a}\right) =\left( \boldsymbol{b}%
\right) $ 
\begin{eqnarray}
T\left( \boldsymbol{a}\right) ^{t}\otimes T\left( \boldsymbol{a}\right)
+T\left( \boldsymbol{a}\right) \otimes T\left( \boldsymbol{a}\right) ^{t}
&=&\left( 
\begin{array}{cc}
Ta_{1} & Ta_{2}%
\end{array}%
\right) \otimes \left( 
\begin{array}{c}
Ta_{1} \\ 
Ta_{2}%
\end{array}%
\right) +\left( 
\begin{array}{c}
Ta_{1} \\ 
Ta_{2}%
\end{array}%
\right) \otimes \left( 
\begin{array}{cc}
Ta_{1} & Ta_{\ 2}%
\end{array}%
\right) \\
&=&\left( 
\begin{array}{cc}
a_{1} & a_{2}%
\end{array}%
\right) \mathcal{R}_{a}\left( T\right) \otimes \mathcal{R}_{a}\left(
T\right) ^{t}\left( 
\begin{array}{c}
a_{1} \\ 
a_{2}%
\end{array}%
\right) +\mathcal{R}_{a}\left( T\right) ^{t}\left( 
\begin{array}{c}
a_{1} \\ 
a_{2}%
\end{array}%
\right) \otimes \left( 
\begin{array}{cc}
a_{1} & a_{\ 2}%
\end{array}%
\right) \mathcal{R}_{a}\left( T\right) \\
&=&\left( 
\begin{array}{cc}
Ta_{1} & Ta_{2}%
\end{array}%
\right) \otimes \left( 
\begin{array}{c}
Ta_{1} \\ 
Ta_{2}%
\end{array}%
\right) +\left( 
\begin{array}{c}
Ta_{1} \\ 
Ta_{2}%
\end{array}%
\right) \otimes \left( 
\begin{array}{cc}
Ta_{1} & Ta_{\ 2}%
\end{array}%
\right) \\
&=&\left( 
\begin{array}{c}
Ta_{1} \\ 
Ta_{2}%
\end{array}%
\right) ^{t}\otimes \left( 
\begin{array}{cc}
Ta_{1} & Ta_{\ 2}%
\end{array}%
\right) ^{t}+\left( 
\begin{array}{c}
Ta_{1} \\ 
Ta_{2}%
\end{array}%
\right) \otimes \left( 
\begin{array}{cc}
Ta_{1} & Ta_{\ 2}%
\end{array}%
\right) \\
&=&\left( \left( 
\begin{array}{cc}
Ta_{1} & Ta_{\ 2}%
\end{array}%
\right) \otimes \left( 
\begin{array}{c}
Ta_{1} \\ 
Ta_{2}%
\end{array}%
\right) \right) ^{t}+\left( 
\begin{array}{c}
Ta_{1} \\ 
Ta_{2}%
\end{array}%
\right) \otimes \left( 
\begin{array}{cc}
Ta_{1} & Ta_{\ 2}%
\end{array}%
\right)
\end{eqnarray}

\item 
\begin{equation}
\left( 
\begin{array}{cc}
Ta_{1} & Ta_{2}%
\end{array}%
\right) \otimes \left( 
\begin{array}{c}
Ta_{1} \\ 
Ta_{2}%
\end{array}%
\right) +\left( 
\begin{array}{c}
Ta_{1} \\ 
Ta_{2}%
\end{array}%
\right) \otimes \left( 
\begin{array}{cc}
Ta_{1} & Ta_{\ 2}%
\end{array}%
\right) =\left( \left( 
\begin{array}{cc}
Ta_{1} & Ta_{2}%
\end{array}%
\right) \otimes \left( 
\begin{array}{c}
Ta_{1} \\ 
Ta_{2}%
\end{array}%
\right) \right) ^{t}+\left( 
\begin{array}{c}
Ta_{1} \\ 
Ta_{2}%
\end{array}%
\right) \otimes \left( 
\begin{array}{cc}
Ta_{1} & Ta_{\ 2}%
\end{array}%
\right)
\end{equation}%
\begin{eqnarray}
\left( 
\begin{array}{cc}
a_{1} & a_{2}%
\end{array}%
\right) \mathcal{R}_{a}\left( T\right) &=&\left( 
\begin{array}{cc}
Ta_{1} & Ta_{2}%
\end{array}%
\right) \\
\mathcal{R}_{a}\left( T\right) ^{t}\left( 
\begin{array}{c}
a_{1} \\ 
a_{2}%
\end{array}%
\right) &=&\left( 
\begin{array}{c}
Ta_{1} \\ 
Ta_{2}%
\end{array}%
\right)
\end{eqnarray}

\item ?????%
\begin{equation}
\left( 
\begin{array}{cc}
a_{1} & a_{2}%
\end{array}%
\right) \mathcal{R}_{a}\left( T\right) =\left( \mathcal{R}_{a}\left(
T\right) ^{t}\left( 
\begin{array}{c}
a_{1} \\ 
a_{2}%
\end{array}%
\right) \right) ^{t}=\left( 
\begin{array}{cc}
Ta_{1} & Ta_{2}%
\end{array}%
\right)
\end{equation}%
and 
\begin{equation}
\mathcal{R}_{a}\left( T\right) ^{t}\left( 
\begin{array}{c}
a_{1} \\ 
a_{2}%
\end{array}%
\right) =\left( \left( 
\begin{array}{cc}
a_{1} & a_{2}%
\end{array}%
\right) \mathcal{R}_{a}\left( T\right) \right) ^{t}=\left( 
\begin{array}{c}
Ta_{1} \\ 
Ta_{2}%
\end{array}%
\right)
\end{equation}%
which gives 
\begin{eqnarray}
\left( 
\begin{array}{cc}
Ta_{1} & Ta_{2}%
\end{array}%
\right) \otimes \left( 
\begin{array}{c}
Ta_{1} \\ 
Ta_{2}%
\end{array}%
\right) &=&\left( \left( 
\begin{array}{c}
Ta_{1} \\ 
Ta_{2}%
\end{array}%
\right) \otimes \left( 
\begin{array}{cc}
Ta_{1} & Ta_{2}%
\end{array}%
\right) \right) ^{t}=\left( \mathcal{R}_{a}\left( T\right) ^{t}\left( 
\begin{array}{c}
a_{1} \\ 
a_{2}%
\end{array}%
\right) \otimes \left( 
\begin{array}{cc}
a_{1} & a_{2}%
\end{array}%
\right) \mathcal{R}_{a}\left( T\right) \right) ^{t} \\
&=&\left( \mathcal{R}_{a}\left( T\right) ^{t}\left( 
\begin{array}{cc}
a_{1}a_{1} & a_{1}a_{2} \\ 
a_{2}a_{1} & a_{2}a_{2}%
\end{array}%
\right) \mathcal{R}_{a}\left( T\right) \right) ^{t}=\mathcal{R}_{a}\left(
T\right) ^{t}\left( 
\begin{array}{cc}
a_{1}a_{1} & a_{1}a_{2} \\ 
a_{2}a_{1} & a_{2}a_{2}%
\end{array}%
\right) ^{t}\mathcal{R}_{a}\left( T\right) \\
&=&\mathcal{R}_{a}\left( T\right) ^{t}\left( 
\begin{array}{cc}
a_{1}a_{1} & a_{2}a_{1} \\ 
a_{1}a_{2} & a_{2}a_{2}%
\end{array}%
\right) \mathcal{R}_{a}\left( T\right)
\end{eqnarray}
\end{enumerate}
\end{remark}

\begin{corollary}
\begin{eqnarray}
\mathcal{R}_{X}\left( T\right) ^{t}\left\{ \overset{}{\boldsymbol{X},%
\boldsymbol{X}}\right\} \mathcal{R}_{X}\left( T\right) &=&\mathcal{R}%
_{X}\left( T\right) ^{t}\left( \left( 
\begin{array}{ccc}
X_{1}X_{1} & \cdots & X_{1}X_{2r} \\ 
\vdots & \ddots & \vdots \\ 
X_{2r}X_{1} & \cdots & X_{2r}X_{2r}%
\end{array}%
\right) +\left( 
\begin{array}{ccc}
X_{1}X_{1} & \cdots & X_{2r}X_{1} \\ 
\vdots & \ddots & \vdots \\ 
X_{1}X_{2r} & \cdots & X_{2r}X_{2r}%
\end{array}%
\right) \right) \mathcal{R}_{X}\left( T\right) \\
&=&\mathcal{R}_{X}\left( T\right) ^{t}\left( 
\begin{array}{ccc}
X_{1}X_{1} & \cdots & X_{1}X_{2r} \\ 
\vdots & \ddots & \vdots \\ 
X_{2r}X_{1} & \cdots & X_{2r}X_{2r}%
\end{array}%
\right) \mathcal{R}_{X}\left( T\right) +\mathcal{R}_{X}\left( T\right)
^{t}\left( 
\begin{array}{ccc}
X_{1}X_{1} & \cdots & X_{2r}X_{1} \\ 
\vdots & \ddots & \vdots \\ 
X_{1}X_{2r} & \cdots & X_{2r}X_{2r}%
\end{array}%
\right) \mathcal{R}_{X}\left( T\right) \\
&=&\mathcal{R}_{X}\left( T\right) ^{t}\left( 
\begin{array}{c}
X_{1} \\ 
\vdots \\ 
X_{2r}%
\end{array}%
\right) \left( 
\begin{array}{ccc}
X_{1} & \cdots & X_{2r}%
\end{array}%
\right) \mathcal{R}_{X}\left( T\right) +\mathcal{R}_{X}\left( T\right)
^{t}\left( 
\begin{array}{ccc}
X_{1}X_{1} & \cdots & X_{2r}X_{1} \\ 
\vdots & \ddots & \vdots \\ 
X_{1}X_{2r} & \cdots & X_{2r}X_{2r}%
\end{array}%
\right) \mathcal{R}_{X}\left( T\right)
\end{eqnarray}
\end{corollary}

\begin{lemma}
\begin{eqnarray}
\left( 
\begin{array}{cc}
\mathbb{U}^{t} & \overline{\mathbb{V}} \\ 
\mathbb{V}^{t} & \overline{\mathbb{U}}%
\end{array}%
\right) \left( 
\begin{array}{cc}
\left\{ \mathbf{a}^{\dagger }\mathbf{a}^{\dagger }\right\} & \left\{ \mathbf{%
a}^{\dagger }\mathbf{a}\right\} \\ 
\left\{ \mathbf{aa}^{\dagger }\right\} & \left\{ \mathbf{aa}\right\}%
\end{array}%
\right) \left( 
\begin{array}{cc}
\mathbb{U} & \overline{\mathbb{V}} \\ 
\mathbb{V} & \overline{\mathbb{U}}%
\end{array}%
\right) &=&\left( 
\begin{array}{cc}
\mathbb{U}^{t}\left\{ \mathbf{a}^{\dagger }\mathbf{a}^{\dagger }\right\} +%
\overline{\mathbb{V}}\left\{ \mathbf{aa}^{\dagger }\right\} & \mathbb{U}%
^{t}\left\{ \mathbf{a}^{\dagger }\mathbf{a}\right\} +\overline{\mathbb{V}}%
\left\{ \mathbf{aa}\right\} \\ 
\mathbb{V}^{t}\left\{ \mathbf{a}^{\dagger }\mathbf{a}^{\dagger }\right\} +%
\overline{\mathbb{U}}\left\{ \mathbf{aa}^{\dagger }\right\} & \mathbb{V}%
^{t}\left\{ \mathbf{a}^{\dagger }\mathbf{a}\right\} +\overline{\mathbb{U}}%
\left\{ \mathbf{aa}\right\}%
\end{array}%
\right) \left( 
\begin{array}{cc}
\mathbb{U} & \overline{\mathbb{V}} \\ 
\mathbb{V} & \overline{\mathbb{U}}%
\end{array}%
\right) \\
&=&\left( 
\begin{array}{cc}
\mathbb{U}^{t}\left\{ \mathbf{a}^{\dagger }\mathbf{a}^{\dagger }\right\} 
\mathbb{U}+\overline{\mathbb{V}}\left\{ \mathbf{aa}^{\dagger }\right\} 
\mathbb{U+U}^{t}\left\{ \mathbf{a}^{\dagger }\mathbf{a}\right\} \mathbb{V}+%
\overline{\mathbb{V}}\left\{ \mathbf{aa}\right\} \mathbb{V} & \mathbb{U}%
^{t}\left\{ \mathbf{a}^{\dagger }\mathbf{a}\right\} +\overline{\mathbb{V}}%
\left\{ \mathbf{aa}\right\} \\ 
\mathbb{V}^{t}\left\{ \mathbf{a}^{\dagger }\mathbf{a}^{\dagger }\right\} +%
\overline{\mathbb{U}}\left\{ \mathbf{aa}^{\dagger }\right\} & \mathbb{V}%
^{t}\left\{ \mathbf{a}^{\dagger }\mathbf{a}\right\} +\overline{\mathbb{U}}%
\left\{ \mathbf{aa}\right\}%
\end{array}%
\right)
\end{eqnarray}%
as 
\begin{equation}
\left\{ \alpha A,\beta B\right\} =\alpha \beta \left\{ A,B\right\} ,\quad
\alpha ,\beta \in \mathbb{C},\quad A,B\in \mathcal{B}\left( \mathcal{F}%
_{1}\right)
\end{equation}%
This gives 
\begin{equation}
\left( 
\begin{array}{cc}
\mathbf{b}^{\dagger }\mathbf{b}^{\dagger } & \mathbf{b}^{\dagger }\mathbf{b}
\\ 
\mathbf{bb}^{\dagger } & \mathbf{bb}%
\end{array}%
\right) =\left( 
\begin{array}{cc}
\mathbb{U} & \overline{\mathbb{V}} \\ 
\mathbb{V} & \overline{\mathbb{U}}%
\end{array}%
\right) \left( 
\begin{array}{cc}
\mathbf{a}^{\dagger }\mathbf{a}^{\dagger } & \mathbf{a}^{\dagger }\mathbf{a}
\\ 
\mathbf{aa}^{\dagger } & \mathbf{aa}%
\end{array}%
\right) \left( 
\begin{array}{cc}
\mathbb{U} & \overline{\mathbb{V}} \\ 
\mathbb{V} & \overline{\mathbb{U}}%
\end{array}%
\right)
\end{equation}
\end{lemma}

\begin{lemma}
Let 
\begin{equation}
\left( 
\begin{array}{c}
\boldsymbol{a} \\ 
\boldsymbol{b}%
\end{array}%
\right) \left( 
\begin{array}{cc}
\boldsymbol{c} & \boldsymbol{d}%
\end{array}%
\right) =\left( 
\begin{array}{cc}
\boldsymbol{ac} & \boldsymbol{ad} \\ 
\boldsymbol{bc} & \boldsymbol{bd}%
\end{array}%
\right)
\end{equation}%
then the transpose acting on super vectors has the same action on super
matrices and Clifford algebra matrix elements 
\begin{equation}
\left( \left( 
\begin{array}{c}
\boldsymbol{a} \\ 
\boldsymbol{b}%
\end{array}%
\right) \left( 
\begin{array}{cc}
\boldsymbol{c} & \boldsymbol{d}%
\end{array}%
\right) \right) ^{t}=\left( 
\begin{array}{c}
\boldsymbol{c} \\ 
\boldsymbol{d}%
\end{array}%
\right) \left( 
\begin{array}{cc}
\boldsymbol{a} & \boldsymbol{b}%
\end{array}%
\right) =\left( 
\begin{array}{cc}
\boldsymbol{ca} & \boldsymbol{cb} \\ 
\boldsymbol{da} & \boldsymbol{db}%
\end{array}%
\right) =\left( 
\begin{array}{cc}
\boldsymbol{ac} & \boldsymbol{ad} \\ 
\boldsymbol{bc} & \boldsymbol{bd}%
\end{array}%
\right) ^{t}.
\end{equation}

\begin{proof}
from the definitions of $"t"$ in $\mathcal{F}$ and $\mathcal{M}\left( 2r,2r,%
\mathcal{F}_{1}\right) .$
\end{proof}
\end{lemma}

\begin{enumerate}
\item 
\begin{eqnarray}
T\left( 
\begin{array}{cccc}
a_{1}^{\dagger } & a_{2}^{\dagger } & a_{1} & a_{2}%
\end{array}%
\right) &=&\left( 
\begin{array}{cccc}
a_{1}^{\dagger } & a_{2}^{\dagger } & a_{1} & a_{2}%
\end{array}%
\right) \mathcal{R}_{a}\left( T\right) =\left( 
\begin{array}{cccc}
a_{1}^{\dagger } & a_{2}^{\dagger } & a_{1} & a_{2}%
\end{array}%
\right) \left( 
\begin{array}{cc}
\mathcal{R}_{a}\left( T\right) _{11} & \mathcal{R}_{a}\left( T\right) _{12}
\\ 
\mathcal{R}_{a}\left( T\right) _{21} & \mathcal{R}_{a}\left( T\right) _{22}%
\end{array}%
\right) \\
&=&\left( 
\begin{array}{cc}
\boldsymbol{a}^{\dagger } & \boldsymbol{a}%
\end{array}%
\right) \left( 
\begin{array}{cc}
\mathcal{R}_{a}\left( T\right) _{11} & \mathcal{R}_{a}\left( T\right) _{12}
\\ 
\mathcal{R}_{a}\left( T\right) _{21} & \mathcal{R}_{a}\left( T\right) _{22}%
\end{array}%
\right) =\left( 
\begin{array}{cc}
\boldsymbol{b}^{\dagger } & \boldsymbol{b}%
\end{array}%
\right)
\end{eqnarray}%
\begin{eqnarray}
\boldsymbol{b}^{\dagger } &=&\left( 
\begin{array}{cc}
\boldsymbol{a}^{\dagger } & \boldsymbol{a}%
\end{array}%
\right) \left( 
\begin{array}{c}
\mathcal{R}_{a}\left( T\right) _{11} \\ 
\mathcal{R}_{a}\left( T\right) _{21}%
\end{array}%
\right) \\
\boldsymbol{b} &=&\left( 
\begin{array}{cc}
\boldsymbol{a}^{\dagger } & \boldsymbol{a}%
\end{array}%
\right) \left( 
\begin{array}{c}
\mathcal{R}_{a}\left( T\right) _{21} \\ 
\mathcal{R}_{a}\left( T\right) _{22}%
\end{array}%
\right)
\end{eqnarray}

\item 
\begin{gather}
\left( 
\begin{array}{cc}
\left( 
\begin{array}{cc}
\left\{ a_{1}^{\dagger },a_{1}^{\dagger }\right\} & \left\{ a_{1}^{\dagger
},a_{2}^{\dagger }\right\} \\ 
\left\{ a_{2}^{\dagger },a_{1}^{\dagger }\right\} & \left\{ a_{2}^{\dagger
},a_{2}^{\dagger }\right\}%
\end{array}%
\right) & \left( 
\begin{array}{cc}
\left\{ a_{1}^{\dagger },a_{1}\right\} & \left\{ a_{1}^{\dagger
},a_{2}\right\} \\ 
\left\{ a_{2}^{\dagger },a_{1}\right\} & \left\{ a_{2}^{\dagger
},a_{2}\right\}%
\end{array}%
\right) \\ 
\left( 
\begin{array}{cc}
\left\{ a_{1},a_{1}^{\dagger }\right\} & \left\{ a_{1},a_{2}^{\dagger
}\right\} \\ 
\left\{ a_{2},a_{1}^{\dagger }\right\} & \left\{ a_{2},a_{2}^{\dagger
}\right\}%
\end{array}%
\right) & \left( 
\begin{array}{cc}
\left\{ a_{1},a_{1}\right\} & \left\{ a_{1},a_{2}\right\} \\ 
\left\{ a_{2},a_{1}\right\} & \left\{ a_{2},a_{2}\right\}%
\end{array}%
\right)%
\end{array}%
\right) =\left( 
\begin{array}{cc}
\mathbb{O} & \mathbb{I}_{2} \\ 
\mathbb{I}_{2} & \mathbb{O}%
\end{array}%
\right) \\
\parallel
\end{gather}%
\begin{gather}
\left( 
\begin{array}{cc}
\left( 
\begin{array}{cc}
a_{1}^{\dagger }a_{1}^{\dagger } & a_{1}^{\dagger }a_{2}^{\dagger } \\ 
a_{2}^{\dagger }a_{1}^{\dagger } & a_{2}^{\dagger }a_{2}^{\dagger }%
\end{array}%
\right) & \left( 
\begin{array}{cc}
a_{1}^{\dagger }a_{1} & a_{1}^{\dagger }a_{2} \\ 
a_{2}^{\dagger }a_{1} & a_{2}^{\dagger }a_{2}%
\end{array}%
\right) \\ 
\left( 
\begin{array}{cc}
a_{1}a_{1}^{\dagger } & a_{1}a_{2}^{\dagger } \\ 
a_{2}a_{1}^{\dagger } & a_{2}a_{2}^{\dagger }%
\end{array}%
\right) & \left( 
\begin{array}{cc}
a_{1}a_{1} & a_{1}a_{2} \\ 
a_{2}a_{1} & a_{2}a_{2}%
\end{array}%
\right)%
\end{array}%
\right) +\left( 
\begin{array}{cc}
\left( 
\begin{array}{cc}
a_{1}^{\dagger }a_{1}^{\dagger } & a_{2}^{\dagger }a_{1}^{\dagger } \\ 
a_{1}^{\dagger }a_{2}^{\dagger } & a_{2}^{\dagger }a_{2}^{\dagger }%
\end{array}%
\right) & \left( 
\begin{array}{cc}
a_{1}a_{1}^{\dagger } & a_{2}a_{1}^{\dagger } \\ 
a_{1}a_{2}^{\dagger } & a_{2}a_{2}^{\dagger }%
\end{array}%
\right) \\ 
\left( 
\begin{array}{cc}
a_{1}^{\dagger }a_{1} & a_{2}^{\dagger }a_{1} \\ 
a_{1}^{\dagger }a_{2} & a_{2}^{\dagger }a_{2}%
\end{array}%
\right) & \left( 
\begin{array}{cc}
a_{1}a_{1} & a_{2}a_{1} \\ 
a_{1}a_{2} & a_{2}a_{2}%
\end{array}%
\right)%
\end{array}%
\right) \\
\parallel
\end{gather}%
\begin{gather}
\left( 
\begin{array}{cc}
\left( 
\begin{array}{cc}
a_{1}^{\dagger }a_{1}^{\dagger } & a_{1}^{\dagger }a_{2}^{\dagger } \\ 
a_{2}^{\dagger }a_{1}^{\dagger } & a_{2}^{\dagger }a_{2}^{\dagger }%
\end{array}%
\right) & \left( 
\begin{array}{cc}
a_{1}^{\dagger }a_{1} & a_{1}^{\dagger }a_{2} \\ 
a_{2}^{\dagger }a_{1} & a_{2}^{\dagger }a_{2}%
\end{array}%
\right) \\ 
\left( 
\begin{array}{cc}
a_{1}a_{1}^{\dagger } & a_{1}a_{2}^{\dagger } \\ 
a_{2}a_{1}^{\dagger } & a_{2}a_{2}^{\dagger }%
\end{array}%
\right) & \left( 
\begin{array}{cc}
a_{1}a_{1} & a_{1}a_{2} \\ 
a_{2}a_{1} & a_{2}a_{2}%
\end{array}%
\right)%
\end{array}%
\right) +\left( 
\begin{array}{cc}
\left( 
\begin{array}{cc}
\left( a_{1}^{\dagger }a_{1}^{\dagger }\right) ^{t} & \left( a_{1}^{\dagger
}a_{2}^{\dagger }\right) ^{t} \\ 
\left( a_{2}^{\dagger }a_{1}^{\dagger }\right) ^{t} & \left( a_{2}^{\dagger
}a_{2}^{\dagger }\right) ^{t}%
\end{array}%
\right) & \left( 
\begin{array}{cc}
\left( a_{1}^{\dagger }a_{1}\right) ^{t} & \left( a_{1}^{\dagger
}a_{2}\right) ^{t} \\ 
\left( a_{2}^{\dagger }a_{1}\right) ^{t} & \left( a_{2}^{\dagger
}a_{2}\right) ^{t}%
\end{array}%
\right) \\ 
\left( 
\begin{array}{cc}
\left( a_{1}a_{1}^{\dagger }\right) ^{t} & \left( a_{1}a_{2}^{\dagger
}\right) ^{t} \\ 
\left( a_{2}a_{1}^{\dagger }\right) ^{t} & \left( a_{2}a_{2}^{\dagger
}\right) ^{t}%
\end{array}%
\right) & \left( 
\begin{array}{cc}
\left( a_{1}a_{1}\right) ^{t} & \left( a_{1}a_{2}\right) ^{t} \\ 
\left( a_{2}a_{1}\right) ^{t} & \left( a_{2}a_{2}\right) ^{t}%
\end{array}%
\right)%
\end{array}%
\right) \\
\parallel \\
\left( 
\begin{array}{cc}
\left( 
\begin{array}{cc}
a_{1}^{\dagger }a_{1}^{\dagger } & a_{1}^{\dagger }a_{2}^{\dagger } \\ 
a_{2}^{\dagger }a_{1}^{\dagger } & a_{2}^{\dagger }a_{2}^{\dagger }%
\end{array}%
\right) & \left( 
\begin{array}{cc}
a_{1}^{\dagger }a_{1} & a_{1}^{\dagger }a_{2} \\ 
a_{2}^{\dagger }a_{1} & a_{2}^{\dagger }a_{2}%
\end{array}%
\right) \\ 
\left( 
\begin{array}{cc}
a_{1}a_{1}^{\dagger } & a_{1}a_{2}^{\dagger } \\ 
a_{2}a_{1}^{\dagger } & a_{2}a_{2}^{\dagger }%
\end{array}%
\right) & \left( 
\begin{array}{cc}
a_{1}a_{1} & a_{1}a_{2} \\ 
a_{2}a_{1} & a_{2}a_{2}%
\end{array}%
\right)%
\end{array}%
\right) +\left( 
\begin{array}{cc}
\left( 
\begin{array}{cc}
\left( a_{1}^{\dagger }a_{1}^{\dagger }\right) & \left( a_{2}^{\dagger
}a_{1}^{\dagger }\right) \\ 
\left( a_{1}^{\dagger }a_{2}^{\dagger }\right) & \left( a_{2}^{\dagger
}a_{2}^{\dagger }\right)%
\end{array}%
\right) ^{t} & \left( 
\begin{array}{cc}
\left( a_{1}^{\dagger }a_{1}\right) & \left( a_{2}^{\dagger }a_{1}\right) \\ 
\left( a_{1}^{\dagger }a_{2}\right) & \left( a_{2}^{\dagger }a_{2}\right)%
\end{array}%
\right) ^{t} \\ 
\left( 
\begin{array}{cc}
\left( a_{1}a_{1}^{\dagger }\right) & \left( a_{2}a_{1}^{\dagger }\right) \\ 
\left( a_{1}a_{2}^{\dagger }\right) & \left( a_{2}a_{2}^{\dagger }\right)%
\end{array}%
\right) ^{t} & \left( 
\begin{array}{cc}
\left( a_{1}a_{1}\right) & \left( a_{1}a_{2}\right) \\ 
\left( a_{2}a_{1}\right) & \left( a_{2}a_{2}\right)%
\end{array}%
\right) ^{t}%
\end{array}%
\right)
\end{gather}%
\begin{gather}
\parallel \\
\left( 
\begin{array}{cc}
\begin{array}{cc}
a_{1}^{\dagger }a_{1}^{\dagger } & a_{1}^{\dagger }a_{2}^{\dagger } \\ 
a_{2}^{\dagger }a_{1}^{\dagger } & a_{2}^{\dagger }a_{2}^{\dagger }%
\end{array}
& 
\begin{array}{cc}
a_{1}^{\dagger }a_{1} & a_{1}^{\dagger }a_{2} \\ 
a_{2}^{\dagger }a_{1} & a_{2}^{\dagger }a_{2}%
\end{array}
\\ 
\begin{array}{cc}
a_{1}a_{1}^{\dagger } & a_{1}a_{2}^{\dagger } \\ 
a_{2}a_{1}^{\dagger } & a_{2}a_{2}^{\dagger }%
\end{array}
& 
\begin{array}{cc}
a_{1}a_{1} & a_{1}a_{2} \\ 
a_{2}a_{1} & a_{2}a_{2}%
\end{array}%
\end{array}%
\right) +\left( 
\begin{array}{cc}
\begin{array}{cc}
a_{1}^{\dagger }a_{1}^{\dagger } & a_{2}^{\dagger }a_{1}^{\dagger } \\ 
a_{1}^{\dagger }a_{2}^{\dagger } & a_{2}^{\dagger }a_{2}^{\dagger }%
\end{array}
& 
\begin{array}{cc}
a_{1}a_{1}^{\dagger } & a_{2}a_{1}^{\dagger } \\ 
a_{1}a_{2}^{\dagger } & a_{2}a_{2}^{\dagger }%
\end{array}
\\ 
\begin{array}{cc}
a_{1}^{\dagger }a_{1} & a_{2}^{\dagger }a_{1} \\ 
a_{1}^{\dagger }a_{2} & a_{2}^{\dagger }a_{2}%
\end{array}
& 
\begin{array}{cc}
a_{1}a_{1} & a_{1}a_{2} \\ 
a_{2}a_{1} & a_{2}a_{2}%
\end{array}%
\end{array}%
\right) ^{t}
\end{gather}

\begin{enumerate}
\item 
\begin{equation}
\left( 
\begin{array}{c}
\boldsymbol{a}^{\dagger } \\ 
\boldsymbol{a}%
\end{array}%
\right) \left( 
\begin{array}{cc}
\boldsymbol{a}^{\dagger } & \boldsymbol{a}%
\end{array}%
\right) =\left( 
\begin{array}{cc}
\boldsymbol{a}^{\dagger }\boldsymbol{a}^{\dagger } & \boldsymbol{a}^{\dagger
}\boldsymbol{a} \\ 
\boldsymbol{aa}^{\dagger } & \boldsymbol{aa}%
\end{array}%
\right)
\end{equation}

\item 
\begin{equation}
\left( 
\begin{array}{cc}
\left( 
\begin{array}{cc}
a_{1}^{\dagger }a_{1}^{\dagger } & a_{1}^{\dagger }a_{2}^{\dagger } \\ 
a_{2}^{\dagger }a_{1}^{\dagger } & a_{2}^{\dagger }a_{2}^{\dagger }%
\end{array}%
\right) & \left( 
\begin{array}{cc}
a_{1}^{\dagger }a_{1} & a_{1}^{\dagger }a_{2} \\ 
a_{2}^{\dagger }a_{1} & a_{2}^{\dagger }a_{2}%
\end{array}%
\right) \\ 
\left( 
\begin{array}{cc}
a_{1}a_{1}^{\dagger } & a_{1}a_{2}^{\dagger } \\ 
a_{2}a_{1}^{\dagger } & a_{2}a_{2}^{\dagger }%
\end{array}%
\right) & \left( 
\begin{array}{cc}
a_{1}a_{1} & a_{1}a_{2} \\ 
a_{2}a_{1} & a_{2}a_{2}%
\end{array}%
\right)%
\end{array}%
\right) =\left( 
\begin{array}{c}
a_{1}^{\dagger } \\ 
a_{2}^{\dagger } \\ 
a_{1} \\ 
a_{2}%
\end{array}%
\right) \left( 
\begin{array}{cccc}
a_{1}^{\dagger } & a_{2}^{\dagger } & a_{1} & a_{2}%
\end{array}%
\right) =\left( 
\begin{array}{cc}
\left( 
\begin{array}{c}
a_{1}^{\dagger } \\ 
a_{2}^{\dagger }%
\end{array}%
\right) \left( 
\begin{array}{cc}
a_{1}^{\dagger } & a_{2}^{\dagger }%
\end{array}%
\right) & \left( 
\begin{array}{c}
a_{1}^{\dagger } \\ 
a_{2}^{\dagger }%
\end{array}%
\right) \left( 
\begin{array}{cc}
a^{\dagger } & a%
\end{array}%
\right) \\ 
\left( 
\begin{array}{c}
a_{1} \\ 
a_{2}%
\end{array}%
\right) \left( 
\begin{array}{cc}
a_{1}^{\dagger } & a_{2}^{\dagger }%
\end{array}%
\right) & \left( 
\begin{array}{c}
a_{1} \\ 
a_{2}%
\end{array}%
\right) \left( 
\begin{array}{cc}
a_{1} & a_{2}%
\end{array}%
\right)%
\end{array}%
\right)
\end{equation}%
\begin{equation}
\left( 
\begin{array}{cc}
\left( 
\begin{array}{cc}
\left( a_{1}^{\dagger }a_{1}^{\dagger }\right) & \left( a_{2}^{\dagger
}a_{1}^{\dagger }\right) \\ 
\left( a_{1}^{\dagger }a_{2}^{\dagger }\right) & \left( a_{2}^{\dagger
}a_{2}^{\dagger }\right)%
\end{array}%
\right) & \left( 
\begin{array}{cc}
\left( a_{1}a_{1}^{\dagger }\right) & \left( a_{2}a_{1}^{\dagger }\right) \\ 
\left( a_{1}a_{2}^{\dagger }\right) & \left( a_{2}a_{2}^{\dagger }\right)%
\end{array}%
\right) \\ 
\left( 
\begin{array}{cc}
\left( a_{1}^{\dagger }a_{1}\right) & \left( a_{2}^{\dagger }a_{1}\right) \\ 
\left( a_{1}^{\dagger }a_{2}\right) & \left( a_{2}^{\dagger }a_{2}\right)%
\end{array}%
\right) & \left( 
\begin{array}{cc}
\left( a_{1}a_{1}\right) & \left( a_{2}a_{1}\right) \\ 
\left( a_{1}a_{2}\right) & \left( a_{2}a_{2}\right)%
\end{array}%
\right)%
\end{array}%
\right) ^{t}
\end{equation}
\end{enumerate}

\begin{enumerate}
\item 
\begin{equation}
\left( 
\begin{array}{cc}
\left( 
\begin{array}{cc}
a_{1}^{\dagger }a_{1}^{\dagger } & a_{2}^{\dagger }a_{1}^{\dagger } \\ 
a_{1}^{\dagger }a_{2}^{\dagger } & a_{2}^{\dagger }a_{2}^{\dagger }%
\end{array}%
\right) & \left( 
\begin{array}{cc}
a_{1}a_{1}^{\dagger } & a_{2}a_{1}^{\dagger } \\ 
a_{1}a_{2}^{\dagger } & a_{2}a_{2}^{\dagger }%
\end{array}%
\right) \\ 
\left( 
\begin{array}{cc}
a_{1}^{\dagger }a_{1} & a_{2}^{\dagger }a_{1} \\ 
a_{1}^{\dagger }a_{2} & a_{2}^{\dagger }a_{2}%
\end{array}%
\right) & \left( 
\begin{array}{cc}
a_{1}a_{1} & a_{2}a_{1} \\ 
a_{1}a_{2} & a_{2}a_{2}%
\end{array}%
\right)%
\end{array}%
\right) =\left( 
\begin{array}{cc}
\left( 
\begin{array}{cc}
\left( a_{1}^{\dagger }a_{1}^{\dagger }\right) ^{t} & \left( a_{1}^{\dagger
}a_{2}^{\dagger }\right) ^{t} \\ 
\left( a_{2}^{\dagger }a_{1}^{\dagger }\right) ^{t} & \left( a_{2}^{\dagger
}a_{2}^{\dagger }\right) ^{t}%
\end{array}%
\right) & \left( 
\begin{array}{cc}
\left( a_{1}^{\dagger }a_{1}\right) ^{t} & \left( a_{1}^{\dagger
}a_{2}\right) ^{t} \\ 
\left( a_{2}^{\dagger }a_{1}\right) ^{t} & \left( a_{2}^{\dagger
}a_{2}\right) ^{t}%
\end{array}%
\right) \\ 
\left( 
\begin{array}{cc}
\left( a_{1}a_{1}^{\dagger }\right) ^{t} & \left( a_{1}a_{2}^{\dagger
}\right) ^{t} \\ 
\left( a_{2}a_{1}^{\dagger }\right) ^{t} & \left( a_{2}a_{2}^{\dagger
}\right) ^{t}%
\end{array}%
\right) & \left( 
\begin{array}{cc}
\left( a_{1}a_{1}\right) ^{t} & \left( a_{1}a_{2}\right) ^{t} \\ 
\left( a_{2}a_{1}\right) ^{t} & \left( a_{2}a_{2}\right) ^{t}%
\end{array}%
\right)%
\end{array}%
\right)
\end{equation}

\item 
\begin{eqnarray}
\left( 
\begin{array}{cc}
\left( a_{1}^{\dagger }a_{1}^{\dagger }\right) ^{t} & \left( a_{1}^{\dagger
}a_{2}^{\dagger }\right) ^{t} \\ 
\left( a_{2}^{\dagger }a_{1}^{\dagger }\right) ^{t} & \left( a_{2}^{\dagger
}a_{2}^{\dagger }\right) ^{t}%
\end{array}%
\right) &=&\left( 
\begin{array}{cc}
\left( a_{1}^{\dagger }a_{1}^{\dagger }\right) & \left( a_{2}^{\dagger
}a_{1}^{\dagger }\right) \\ 
\left( a_{1}^{\dagger }a_{2}^{\dagger }\right) & \left( a_{2}^{\dagger
}a_{2}^{\dagger }\right)%
\end{array}%
\right) ^{t} \\
\left( 
\begin{array}{cc}
\left( a_{1}^{\dagger }a_{1}\right) ^{t} & \left( a_{1}^{\dagger
}a_{2}\right) ^{t} \\ 
\left( a_{2}^{\dagger }a_{1}\right) ^{t} & \left( a_{2}^{\dagger
}a_{2}\right) ^{t}%
\end{array}%
\right) &=&\left( 
\begin{array}{cc}
\left( a_{1}^{\dagger }a_{1}\right) & \left( a_{2}^{\dagger }a_{1}\right) \\ 
\left( a_{1}^{\dagger }a_{2}\right) & \left( a_{2}^{\dagger }a_{2}\right)%
\end{array}%
\right) ^{t} \\
\left( 
\begin{array}{cc}
\left( a_{1}a_{1}^{\dagger }\right) ^{t} & \left( a_{1}a_{2}^{\dagger
}\right) ^{t} \\ 
\left( a_{2}a_{1}^{\dagger }\right) ^{t} & \left( a_{2}a_{2}^{\dagger
}\right) ^{t}%
\end{array}%
\right) &=&\left( 
\begin{array}{cc}
\left( a_{1}a_{1}^{\dagger }\right) & \left( a_{2}a_{1}^{\dagger }\right) \\ 
\left( a_{1}a_{2}^{\dagger }\right) & \left( a_{2}a_{2}^{\dagger }\right)%
\end{array}%
\right) ^{t} \\
\left( 
\begin{array}{cc}
\left( a_{1}a_{1}\right) ^{t} & \left( a_{1}a_{2}\right) ^{t} \\ 
\left( a_{2}a_{1}\right) ^{t} & \left( a_{2}a_{2}\right) ^{t}%
\end{array}%
\right) &=&\left( 
\begin{array}{cc}
\left( a_{1}a_{1}\right) & \left( a_{2}a_{1}\right) \\ 
\left( a_{1}a_{2}\right) & \left( a_{2}a_{2}\right)%
\end{array}%
\right) ^{t}
\end{eqnarray}

\item 
\begin{eqnarray}
\left( 
\begin{array}{cc}
\left( 
\begin{array}{cc}
a_{1}^{\dagger }a_{1}^{\dagger } & a_{2}^{\dagger }a_{1}^{\dagger } \\ 
a_{1}^{\dagger }a_{2}^{\dagger } & a_{2}^{\dagger }a_{2}^{\dagger }%
\end{array}%
\right) & \left( 
\begin{array}{cc}
a_{1}a_{1}^{\dagger } & a_{2}a_{1}^{\dagger } \\ 
a_{1}a_{2}^{\dagger } & a_{2}a_{2}^{\dagger }%
\end{array}%
\right) \\ 
\left( 
\begin{array}{cc}
a_{1}^{\dagger }a_{1} & a_{2}^{\dagger }a_{1} \\ 
a_{1}^{\dagger }a_{2} & a_{2}^{\dagger }a_{2}%
\end{array}%
\right) & \left( 
\begin{array}{cc}
a_{1}a_{1} & a_{2}a_{1} \\ 
a_{1}a_{2} & a_{2}a_{2}%
\end{array}%
\right)%
\end{array}%
\right) &=&\left( 
\begin{array}{cc}
\left( 
\begin{array}{cc}
\left( a_{1}^{\dagger }a_{1}^{\dagger }\right) & \left( a_{2}^{\dagger
}a_{1}^{\dagger }\right) \\ 
\left( a_{1}^{\dagger }a_{2}^{\dagger }\right) & \left( a_{2}^{\dagger
}a_{2}^{\dagger }\right)%
\end{array}%
\right) ^{t} & \left( 
\begin{array}{cc}
\left( a_{1}^{\dagger }a_{1}\right) & \left( a_{2}^{\dagger }a_{1}\right) \\ 
\left( a_{1}^{\dagger }a_{2}\right) & \left( a_{2}^{\dagger }a_{2}\right)%
\end{array}%
\right) ^{t} \\ 
\left( 
\begin{array}{cc}
\left( a_{1}a_{1}^{\dagger }\right) & \left( a_{2}a_{1}^{\dagger }\right) \\ 
\left( a_{1}a_{2}^{\dagger }\right) & \left( a_{2}a_{2}^{\dagger }\right)%
\end{array}%
\right) ^{t} & \left( 
\begin{array}{cc}
\left( a_{1}a_{1}\right) & \left( a_{2}a_{1}\right) \\ 
\left( a_{1}a_{2}\right) & \left( a_{2}a_{2}\right)%
\end{array}%
\right) ^{t}%
\end{array}%
\right) \\
&=&\left( 
\begin{array}{cc}
\left( 
\begin{array}{cc}
\left( a_{1}^{\dagger }a_{1}^{\dagger }\right) & \left( a_{2}^{\dagger
}a_{1}^{\dagger }\right) \\ 
\left( a_{1}^{\dagger }a_{2}^{\dagger }\right) & \left( a_{2}^{\dagger
}a_{2}^{\dagger }\right)%
\end{array}%
\right) & \left( 
\begin{array}{cc}
\left( a_{1}a_{1}^{\dagger }\right) & \left( a_{2}a_{1}^{\dagger }\right) \\ 
\left( a_{1}a_{2}^{\dagger }\right) & \left( a_{2}a_{2}^{\dagger }\right)%
\end{array}%
\right) \\ 
\left( 
\begin{array}{cc}
\left( a_{1}^{\dagger }a_{1}\right) & \left( a_{2}^{\dagger }a_{1}\right) \\ 
\left( a_{1}^{\dagger }a_{2}\right) & \left( a_{2}^{\dagger }a_{2}\right)%
\end{array}%
\right) & \left( 
\begin{array}{cc}
\left( a_{1}a_{1}\right) & \left( a_{2}a_{1}\right) \\ 
\left( a_{1}a_{2}\right) & \left( a_{2}a_{2}\right)%
\end{array}%
\right)%
\end{array}%
\right) ^{t}
\end{eqnarray}
\end{enumerate}

\item 
\begin{equation}
=\left( 
\begin{array}{cc}
\left( 
\begin{array}{cc}
a_{1}^{\dagger }a_{1}^{\dagger } & a_{2}^{\dagger }a_{1}^{\dagger } \\ 
a_{1}^{\dagger }a_{2}^{\dagger } & a_{2}^{\dagger }a_{2}^{\dagger }%
\end{array}%
\right) ^{t} & \left( 
\begin{array}{cc}
a_{1}^{\dagger }a_{1} & a_{1}^{\dagger }a_{2} \\ 
a_{2}^{\dagger }a_{1} & a_{2}^{\dagger }a_{2}%
\end{array}%
\right) ^{t} \\ 
\left( 
\begin{array}{cc}
a_{1}a_{1}^{\dagger } & a_{2}a_{1}^{\dagger } \\ 
a_{1}a_{2}^{\dagger } & a_{2}a_{2}^{\dagger }%
\end{array}%
\right) ^{t} & \left( 
\begin{array}{cc}
a_{1}a_{1} & a_{2}a_{1} \\ 
a_{1}a_{2} & a_{2}a_{2}%
\end{array}%
\right) ^{t}%
\end{array}%
\right) =\left( 
\begin{array}{cc}
\left( \left( 
\begin{array}{c}
a_{1}^{\dagger } \\ 
a_{2}^{\dagger }%
\end{array}%
\right) \left( 
\begin{array}{cc}
a_{1}^{\dagger } & a_{2}^{\dagger }%
\end{array}%
\right) \right) ^{t} & \left( \left( 
\begin{array}{c}
a_{1}^{\dagger } \\ 
a_{2}^{\dagger }%
\end{array}%
\right) \left( 
\begin{array}{cc}
a_{1} & a_{2}%
\end{array}%
\right) \right) ^{t} \\ 
\left( \left( 
\begin{array}{c}
a_{1} \\ 
a_{2}%
\end{array}%
\right) \left( 
\begin{array}{cc}
a_{1}^{\dagger } & a_{2}^{\dagger }%
\end{array}%
\right) \right) ^{t} & \left( \left( 
\begin{array}{c}
a_{1} \\ 
a_{2}%
\end{array}%
\right) \left( 
\begin{array}{cc}
a_{1} & a_{2}%
\end{array}%
\right) \right) ^{t}%
\end{array}%
\right)
\end{equation}

\item 
\begin{equation}
\left( 
\begin{array}{cc}
\left( \boldsymbol{a}^{\dagger }\boldsymbol{a}^{\dagger }\right) ^{t} & 
\left( \boldsymbol{aa}^{\dagger }\right) ^{t} \\ 
\left( \boldsymbol{a}^{\dagger }\boldsymbol{a}\right) ^{t} & \left( 
\boldsymbol{aa}\right) ^{t}%
\end{array}%
\right) =\left( 
\begin{array}{cc}
\left( 
\begin{array}{cc}
a_{1}^{\dagger }a_{1}^{\dagger } & a_{2}^{\dagger }a_{1}^{\dagger } \\ 
a_{1}^{\dagger }a_{2}^{\dagger } & a_{2}^{\dagger }a_{2}^{\dagger }%
\end{array}%
\right) & \left( 
\begin{array}{cc}
a_{1}^{\dagger }a_{1} & a_{2}^{\dagger }a_{1} \\ 
a_{1}^{\dagger }a_{2} & a_{2}^{\dagger }a_{2}%
\end{array}%
\right) \\ 
\left( 
\begin{array}{cc}
a_{1}a_{1}^{\dagger } & a_{2}a_{1}^{\dagger } \\ 
a_{1}a_{2}^{\dagger } & a_{2}a_{2}^{\dagger }%
\end{array}%
\right) & \left( 
\begin{array}{cc}
a_{1}a_{1} & a_{2}a_{1} \\ 
a_{1}a_{2} & a_{2}a_{2}%
\end{array}%
\right)%
\end{array}%
\right)
\end{equation}

\item 
\begin{equation}
\left( 
\begin{array}{cc}
\boldsymbol{a}^{\dagger }\boldsymbol{a}^{\dagger } & \boldsymbol{a}^{\dagger
}\boldsymbol{a} \\ 
\boldsymbol{aa}^{\dagger } & \boldsymbol{aa}%
\end{array}%
\right) ^{t}=\left( 
\begin{array}{cc}
\left( \boldsymbol{a}^{\dagger }\boldsymbol{a}^{\dagger }\right) ^{t} & 
\left( \boldsymbol{aa}^{\dagger }\right) ^{t} \\ 
\left( \boldsymbol{a}^{\dagger }\boldsymbol{a}\right) ^{t} & \left( 
\boldsymbol{aa}\right) ^{t}%
\end{array}%
\right) =\left( 
\begin{array}{cc}
\left( 
\begin{array}{cc}
a_{1}^{\dagger }a_{1}^{\dagger } & a_{2}^{\dagger }a_{1}^{\dagger } \\ 
a_{1}^{\dagger }a_{2}^{\dagger } & a_{2}^{\dagger }a_{2}^{\dagger }%
\end{array}%
\right) & \left( 
\begin{array}{cc}
a_{1}^{\dagger }a_{1} & a_{2}^{\dagger }a_{1} \\ 
a_{1}^{\dagger }a_{2} & a_{2}^{\dagger }a_{2}%
\end{array}%
\right) \\ 
\left( 
\begin{array}{cc}
a_{1}a_{1}^{\dagger } & a_{2}a_{1}^{\dagger } \\ 
a_{1}a_{2}^{\dagger } & a_{2}a_{2}^{\dagger }%
\end{array}%
\right) & \left( 
\begin{array}{cc}
a_{1}a_{1} & a_{2}a_{1} \\ 
a_{1}a_{2} & a_{2}a_{2}%
\end{array}%
\right)%
\end{array}%
\right)
\end{equation}%
\begin{equation}
\left( 
\begin{array}{cc}
\boldsymbol{0} & \boldsymbol{1} \\ 
\boldsymbol{1} & \boldsymbol{0}%
\end{array}%
\right) \left( 
\begin{array}{cc}
\boldsymbol{a}^{\dagger }\boldsymbol{a}^{\dagger } & \boldsymbol{a}^{\dagger
}\boldsymbol{a} \\ 
\boldsymbol{aa}^{\dagger } & \boldsymbol{aa}%
\end{array}%
\right) \left( 
\begin{array}{cc}
\boldsymbol{0} & \boldsymbol{1} \\ 
\boldsymbol{1} & \boldsymbol{0}%
\end{array}%
\right) =\left( 
\begin{array}{cc}
\boldsymbol{aa}^{\dagger } & \boldsymbol{aa} \\ 
\boldsymbol{a}^{\dagger }\boldsymbol{a}^{\dagger } & \boldsymbol{a}^{\dagger
}\boldsymbol{a}%
\end{array}%
\right) \left( 
\begin{array}{cc}
\boldsymbol{0} & \boldsymbol{1} \\ 
\boldsymbol{1} & \boldsymbol{0}%
\end{array}%
\right) =\left( 
\begin{array}{cc}
\boldsymbol{aa} & \boldsymbol{aa}^{\dagger } \\ 
\boldsymbol{a}^{\dagger }\boldsymbol{a} & \boldsymbol{a}^{\dagger }%
\boldsymbol{a}%
\end{array}%
\right)
\end{equation}

\item 
\begin{equation}
\left( \left( 
\begin{array}{c}
\boldsymbol{a}^{\dagger } \\ 
\boldsymbol{a}%
\end{array}%
\right) \left( 
\begin{array}{cc}
\boldsymbol{a}^{\dagger } & \boldsymbol{a}%
\end{array}%
\right) \right) ^{t}=\left( 
\begin{array}{c}
\boldsymbol{a}^{\dagger } \\ 
\boldsymbol{a}%
\end{array}%
\right) \left( 
\begin{array}{cc}
\boldsymbol{a}^{\dagger } & \boldsymbol{a}%
\end{array}%
\right) =\left( 
\begin{array}{cc}
\boldsymbol{a}^{\dagger }\boldsymbol{a}^{\dagger } & \boldsymbol{a}^{\dagger
}\boldsymbol{a} \\ 
\boldsymbol{aa}^{\dagger } & \boldsymbol{aa}%
\end{array}%
\right) ^{t}=\left( 
\begin{array}{cc}
\left( \boldsymbol{a}^{\dagger }\boldsymbol{a}^{\dagger }\right) ^{t} & 
\left( \boldsymbol{aa}^{\dagger }\right) ^{t} \\ 
\left( \boldsymbol{a}^{\dagger }\boldsymbol{a}\right) ^{t} & \left( 
\boldsymbol{aa}\right) ^{t}%
\end{array}%
\right)
\end{equation}%
\begin{equation}
\left( \left( 
\begin{array}{c}
a_{1}^{\dagger } \\ 
a_{2}^{\dagger } \\ 
a_{1} \\ 
a_{2}%
\end{array}%
\right) \left( 
\begin{array}{cccc}
a_{1}^{\dagger } & a_{2}^{\dagger } & a_{1} & a_{2}%
\end{array}%
\right) \right) ^{t}=\left( 
\begin{array}{c}
a_{1}^{\dagger } \\ 
a_{2}^{\dagger } \\ 
a_{1} \\ 
a_{2}%
\end{array}%
\right) \left( 
\begin{array}{cccc}
a_{1}^{\dagger } & a_{2}^{\dagger } & a_{1} & a_{2}%
\end{array}%
\right)
\end{equation}%
\begin{eqnarray}
\left( 
\begin{array}{cc}
\left( 
\begin{array}{cc}
a_{1}^{\dagger }a_{1}^{\dagger } & a_{1}^{\dagger }a_{2}^{\dagger } \\ 
a_{2}^{\dagger }a_{1}^{\dagger } & a_{2}^{\dagger }a_{2}^{\dagger }%
\end{array}%
\right) & \left( 
\begin{array}{cc}
a_{1}^{\dagger }a_{1} & a_{1}^{\dagger }a_{2} \\ 
a_{2}^{\dagger }a_{1} & a_{2}^{\dagger }a_{2}%
\end{array}%
\right) \\ 
\left( 
\begin{array}{cc}
a_{1}a_{1}^{\dagger } & a_{1}a_{2}^{\dagger } \\ 
a_{2}a_{1}^{\dagger } & a_{2}a_{2}^{\dagger }%
\end{array}%
\right) & \left( 
\begin{array}{cc}
a_{1}a_{1} & a_{1}a_{2} \\ 
a_{2}a_{1} & a_{2}a_{2}%
\end{array}%
\right)%
\end{array}%
\right) ^{t} &=&\left( 
\begin{array}{cc}
\left( 
\begin{array}{cc}
a_{1}^{\dagger }a_{1}^{\dagger } & a_{1}^{\dagger }a_{2}^{\dagger } \\ 
a_{2}^{\dagger }a_{1}^{\dagger } & a_{2}^{\dagger }a_{2}^{\dagger }%
\end{array}%
\right) ^{t} & \left( 
\begin{array}{cc}
a_{1}a_{1}^{\dagger } & a_{1}a_{2}^{\dagger } \\ 
a_{2}a_{1}^{\dagger } & a_{2}a_{2}^{\dagger }%
\end{array}%
\right) ^{t} \\ 
\left( 
\begin{array}{cc}
a_{1}^{\dagger }a_{1} & a_{1}^{\dagger }a_{2} \\ 
a_{2}^{\dagger }a_{1} & a_{2}^{\dagger }a_{2}%
\end{array}%
\right) ^{t} & \left( 
\begin{array}{cc}
a_{1}a_{1} & a_{1}a_{2} \\ 
a_{2}a_{1} & a_{2}a_{2}%
\end{array}%
\right) ^{t}%
\end{array}%
\right) \\
&=&\left( 
\begin{array}{cc}
\left( 
\begin{array}{cc}
a_{1}^{\dagger }a_{1}^{\dagger } & a_{2}^{\dagger }a_{1}^{\dagger } \\ 
a_{1}^{\dagger }a_{2}^{\dagger } & a_{2}^{\dagger }a_{2}^{\dagger }%
\end{array}%
\right) & \left( 
\begin{array}{cc}
a_{1}^{\dagger }a_{1} & a_{1}^{\dagger }a_{2} \\ 
a_{2}^{\dagger }a_{1} & a_{2}^{\dagger }a_{2}%
\end{array}%
\right) \\ 
\left( 
\begin{array}{cc}
a_{1}a_{1}^{\dagger } & a_{1}a_{2}^{\dagger } \\ 
a_{2}a_{1}^{\dagger } & a_{2}a_{2}^{\dagger }%
\end{array}%
\right) & \left( 
\begin{array}{cc}
a_{1}a_{1} & a_{1}a_{2} \\ 
a_{2}a_{1} & a_{2}a_{2}%
\end{array}%
\right)%
\end{array}%
\right)
\end{eqnarray}%
\begin{eqnarray}
\left( 
\begin{array}{cc}
a_{1}^{\dagger }a_{1} & a_{1}^{\dagger }a_{2} \\ 
a_{2}^{\dagger }a_{1} & a_{2}^{\dagger }a_{2}%
\end{array}%
\right) &=&\left( 
\begin{array}{c}
a_{1}^{\dagger } \\ 
a_{2}^{\dagger }%
\end{array}%
\right) \left( 
\begin{array}{cc}
a_{1} & a_{2}%
\end{array}%
\right) \\
\left( \left( 
\begin{array}{c}
a_{1}^{\dagger } \\ 
a_{2}^{\dagger }%
\end{array}%
\right) \left( 
\begin{array}{cc}
a_{1} & a_{2}%
\end{array}%
\right) \right) ^{t} &=&\left( 
\begin{array}{c}
a_{1} \\ 
a_{2}%
\end{array}%
\right) \left( 
\begin{array}{cc}
a_{1}^{\dagger } & a_{2}^{\dagger }%
\end{array}%
\right) =\left( 
\begin{array}{cc}
a_{1}a_{1}^{\dagger } & a_{1}a_{2}^{\dagger } \\ 
a_{2}a_{1}^{\dagger } & a_{2}a_{1}^{\dagger }%
\end{array}%
\right) \\
\left( 
\begin{array}{cc}
a_{1}^{\dagger }a_{1} & a_{1}^{\dagger }a_{2} \\ 
a_{2}^{\dagger }a_{1} & a_{2}^{\dagger }a_{2}%
\end{array}%
\right) ^{t} &=&\left( 
\begin{array}{cc}
a_{1}a_{1}^{\dagger } & a_{1}a_{2}^{\dagger } \\ 
a_{2}a_{1}^{\dagger } & a_{2}a_{2}^{\dagger }%
\end{array}%
\right)
\end{eqnarray}
\end{enumerate}

\subsubsection{$C^{\ast }$-algebra Maps \label{C^ algebra Maps}}

\begin{definition}
A map $T\in \hom \left( \mathcal{B},\mathcal{A}\right) $ is a $c^{\ast }$%
-algebra map if

\begin{enumerate}
\item 
\begin{equation}
T\left( z_{1}a_{1}\right) +T\left( z_{2}a_{2}\right) =T\left(
z_{1}a_{1}+z_{2}a_{2}\right)
\end{equation}

\item 
\begin{equation}
T\left( a_{1}a_{2}\right) =T\left( a_{1}\right) T\left( a_{2}\right)
\end{equation}

\item 
\begin{equation}
T\left( a^{\dagger _{\mathcal{A}}}\right) =T\left( a\right) ^{\dagger _{%
\mathcal{B}}}
\end{equation}

\item 
\begin{equation}
\left\Vert T\left( a\right) \right\Vert _{\mathcal{B}}=\left\Vert
a\right\Vert _{\mathcal{A}}.
\end{equation}
\end{enumerate}
\end{definition}

\begin{remark}
$c^{\ast }$-algebra representations $\mathcal{R}_{c}\in \hom \left( \mathcal{%
B}\left( \mathcal{F}_{1}\right) ,\mathcal{M}\left( 2r,2r,\mathbb{C}\right)
\right) $ have the property 
\begin{equation}
\mathcal{R}_{c}\left( A^{\dagger }\right) =\mathcal{R}_{c}\left( A\right)
^{\dagger }.
\end{equation}
\end{remark}

\begin{theorem}
$T\in GL(2r,\mathcal{F}_{1}\mathbb{)}$ is a $c^{\ast }$-algebra map $%
\limfunc{Iff}$ 
\begin{equation}
T:\left( \boldsymbol{x}\right) ^{t}\rightarrow \left( \boldsymbol{y}\right)
^{t},
\end{equation}%
where $\left( \boldsymbol{x}\right) ^{t}$ and $\left( \boldsymbol{y}\right)
^{t}$ are self adjoint bases of $\mathcal{F}_{1}.$
\end{theorem}

\begin{theorem}
$T\in GL(2r,\mathbb{C)}$ preserves the Clifford Commutation Relationship
form of the CAR $\limfunc{Iff}$ 
\begin{equation}
\mathbb{T}^{t}\mathbb{T=I}_{2r}
\end{equation}%
i.e. iff $T\in O(2r,\mathbb{C)}.$ \emph{However such }$T^{\prime }s$\emph{\
are not in general }$\ast $\emph{-transformations} i.e. 
\begin{equation}
T\in O(2r,\mathbb{C)}\nRightarrow T\left( X^{\dagger }\right) =T\left(
X\right) ^{\dagger }.
\end{equation}

\begin{proof}
\begin{equation}
T\left( 
\begin{array}{cc}
\mathbf{x}_{1} & \mathbf{x}_{2}%
\end{array}%
\right) =\left( \mathbf{x}_{1}\mathbf{x}_{2}\right) \mathbb{T}=\left( 
\mathbf{x}_{1}\mathbf{x}_{2}\right) \left( 
\begin{array}{cc}
\mathbb{T}_{11} & \mathbb{T}_{12} \\ 
\mathbb{T}_{21} & \mathbb{T}_{22}%
\end{array}%
\right) =\left( 
\begin{array}{cc}
\mathbf{y}_{1} & \mathbf{y}_{2}%
\end{array}%
\right)
\end{equation}%
giving 
\begin{gather}
\left( 
\begin{array}{cc}
\left\{ \mathbf{y}_{1},\mathbf{y}_{1}\right\} & \left\{ \mathbf{y}_{1},%
\mathbf{y}_{2}\right\} \\ 
\left\{ \mathbf{y}_{2},\mathbf{y}_{1}\right\} & \left\{ \mathbf{y}_{2},%
\mathbf{y}_{2}\right\}%
\end{array}%
\right) \overset{\triangle }{=}\left( 
\begin{array}{cc}
\mathbf{y}_{1}^{\mathfrak{t}}\mathbf{y}_{1} & \mathbf{y}_{1}^{\mathfrak{t}}%
\mathbf{y}_{2} \\ 
\mathbf{y}_{2}^{\mathfrak{t}}\mathbf{y}_{1} & \mathbf{y}_{2}^{\mathfrak{t}}%
\mathbf{y}_{2}%
\end{array}%
\right) +\left( 
\begin{array}{cc}
\mathbf{y}_{1}^{\mathfrak{t}}\mathbf{y}_{1} & \mathbf{y}_{1}^{\mathfrak{t}}%
\mathbf{y}_{2} \\ 
\mathbf{y}_{2}^{\mathfrak{t}}\mathbf{y}_{1} & \mathbf{y}_{2}^{\mathfrak{t}}%
\mathbf{y}_{2}%
\end{array}%
\right) ^{\mathfrak{t}}  \notag \\
=\mathbb{T}^{t}\left( 
\begin{array}{cc}
\mathbf{x}_{1}^{\mathfrak{t}}\mathbf{x}_{1} & \mathbf{x}_{1}^{\mathfrak{t}}%
\mathbf{x}_{2} \\ 
\mathbf{x}_{2}^{\mathfrak{t}}\mathbf{x}_{1} & \mathbf{x}_{2}^{\mathfrak{t}}%
\mathbf{x}_{2}%
\end{array}%
\right) \mathbb{T}+\left( \mathbb{T}^{t}\left( 
\begin{array}{cc}
\mathbf{x}_{1}^{\mathfrak{t}}\mathbf{x}_{1} & \mathbf{x}_{1}^{\mathfrak{t}}%
\mathbf{x}_{2} \\ 
\mathbf{x}_{2}^{\mathfrak{t}}\mathbf{x}_{1} & \mathbf{x}_{2}^{\mathfrak{t}}%
\mathbf{x}_{2}%
\end{array}%
\right) \mathbb{T}\right) ^{\mathfrak{t}} \\
=\mathbb{T}^{t}\left( 
\begin{array}{cc}
\mathbf{x}_{1}^{\mathfrak{t}}\mathbf{x}_{1} & \mathbf{x}_{1}^{\mathfrak{t}}%
\mathbf{x}_{2} \\ 
\mathbf{x}_{2}^{\mathfrak{t}}\mathbf{x}_{1} & \mathbf{x}_{2}^{\mathfrak{t}}%
\mathbf{x}_{2}%
\end{array}%
\right) \mathbb{T}+\mathbb{T}^{t}\left( 
\begin{array}{cc}
\mathbf{x}_{1}^{\mathfrak{t}}\mathbf{x}_{1} & \mathbf{x}_{1}^{\mathfrak{t}}%
\mathbf{x}_{2} \\ 
\mathbf{x}_{2}^{\mathfrak{t}}\mathbf{x}_{1} & \mathbf{x}_{2}^{\mathfrak{t}}%
\mathbf{x}_{2}%
\end{array}%
\right) ^{\mathfrak{t}}\mathbb{T} \\
=\mathbb{T}^{t}\left( \left( 
\begin{array}{cc}
\left\{ \mathbf{x}_{1},\mathbf{x}_{1}\right\} & \left\{ \mathbf{x}_{1},%
\mathbf{x}_{2}\right\} \\ 
\left\{ \mathbf{x}_{2},\mathbf{x}_{1}\right\} & \left\{ \mathbf{x}_{2},%
\mathbf{x}_{2}\right\}%
\end{array}%
\right) \right) \mathbb{T}.
\end{gather}
\end{proof}
\end{theorem}

The anti commutator $\left\{ ,\right\} $ is a symmetric bilinear form 
\begin{equation}
\left\{ ,\right\} :\mathcal{F}_{1}\times \mathcal{F}_{1}\rightarrow \mathcal{%
F}_{1}
\end{equation}

\begin{notation}
\begin{enumerate}
\item The transpose map $t$ on vectors with $\mathcal{F}_{1}$ elements
interchanges row and column%
\begin{equation}
t\in \left\{ 
\begin{array}{c}
\mathcal{L}\left( \mathcal{M}\left( 2r,1,\mathcal{F}_{1}\right) ,\mathcal{M}%
\left( 1,2r,\mathcal{F}_{1}\right) \right) \\ 
\mathcal{L}\left( \mathcal{M}\left( 1,2r,\mathcal{F}_{1}\right) ,\mathcal{M}%
\left( 2r,1,\mathcal{F}_{1}\right) \right)%
\end{array}%
\right. ,
\end{equation}%
where $\mathcal{F}_{1}$ is the vector space of grade $1$ elements of the
Clifford algebra $\mathcal{F}$ is given by

\item 
\begin{eqnarray}
\left( 
\begin{array}{c}
c_{1} \\ 
\vdots \\ 
c_{2r}%
\end{array}%
\right) ^{t} &=&\left( 
\begin{array}{ccc}
c_{1} & \cdots & c_{2r}%
\end{array}%
\right) \\
\left( 
\begin{array}{ccc}
c_{1} & \cdots & c_{2r}%
\end{array}%
\right) ^{t} &=&\left( 
\begin{array}{c}
c_{1} \\ 
\vdots \\ 
c_{2r}%
\end{array}%
\right)
\end{eqnarray}%
as

\item 
\begin{equation}
c^{t}=c.
\end{equation}
\end{enumerate}
\end{notation}

\begin{notation}
The transpose map $t$ on matrices with $\mathcal{F}_{1}$ elements
interchanges row and column 
\begin{equation}
t\in \mathcal{L}\left( \mathcal{M}\left( \overset{}{2r,2r,\mathcal{F}_{1}}%
\right) ,\mathcal{M}\left( \overset{}{2r,2r,\mathcal{F}_{1}}\right) \right)
\end{equation}%
and is given by 
\begin{equation}
\left( 
\begin{array}{ccc}
c_{11} & \cdots & c_{12r} \\ 
\vdots & \ddots & \vdots \\ 
c_{2r1} & \cdots & c_{2r2r}%
\end{array}%
\right) ^{t}=\left( 
\begin{array}{ccc}
c_{11} & \cdots & c_{12r} \\ 
\vdots & \ddots & \vdots \\ 
c_{12r} & \cdots & c_{2r2r}%
\end{array}%
\right)
\end{equation}
\end{notation}

\begin{notation}
\begin{enumerate}
\item 
\begin{enumerate}
\item 
\begin{equation}
\left[ \left( 
\begin{array}{c}
c_{1} \\ 
\vdots \\ 
c_{2r}%
\end{array}%
\right) \right] \left( 
\begin{array}{ccc}
d_{1} & \cdots & d_{2r}%
\end{array}%
\right) =\left( 
\begin{array}{ccc}
c_{1}\circ d_{1} & \cdots & c_{2r}\circ d_{2r} \\ 
\vdots & \ddots & \vdots \\ 
c_{2r}\circ d_{1} & \cdots & c_{2r}\circ d_{2r}%
\end{array}%
\right) \equiv \left( 
\begin{array}{ccc}
c_{1}d_{1} & \cdots & c_{2r}d_{2r} \\ 
\vdots & \ddots & \vdots \\ 
c_{2r}d_{1} & \cdots & c_{2r}d_{2r}%
\end{array}%
\right)
\end{equation}

\item 
\begin{equation}
\left( \left[ \left( 
\begin{array}{c}
c_{1} \\ 
\vdots \\ 
c_{2r}%
\end{array}%
\right) \right] \left( 
\begin{array}{ccc}
d_{1} & \cdots & d_{2r}%
\end{array}%
\right) \right) ^{t}=\left[ \left( 
\begin{array}{c}
d_{1} \\ 
\vdots \\ 
d_{2r}%
\end{array}%
\right) \right] \left( 
\begin{array}{ccc}
c_{1} & \cdots & c_{2r}%
\end{array}%
\right) =\left( 
\begin{array}{ccc}
d_{1}\circ c_{1} & \cdots & d_{2r}\circ c_{2r} \\ 
\vdots & \ddots & \vdots \\ 
d_{2r}\circ c_{1} & \cdots & d_{2r}\circ c_{2r}%
\end{array}%
\right) \equiv \left( 
\begin{array}{ccc}
d_{1}c_{1} & \cdots & d_{2r}c_{2r} \\ 
\vdots & \ddots & \vdots \\ 
d_{2r}c_{1} & \cdots & d_{2r}c_{2r}%
\end{array}%
\right)
\end{equation}
\end{enumerate}

\begin{enumerate}
\item 
\begin{equation}
\left[ \left( 
\begin{array}{ccc}
c_{1} & \cdots & c_{2r}%
\end{array}%
\right) \right] \left( 
\begin{array}{c}
d_{1} \\ 
\vdots \\ 
d_{2r}%
\end{array}%
\right) =c_{1}\circ d_{1}+\cdots +c_{2r}\circ d_{2r}\equiv c_{1}d_{1}+\cdots
+c_{2r}d_{2r}
\end{equation}

\item 
\begin{equation}
\left( \left[ \left( 
\begin{array}{ccc}
c_{1} & \cdots & c_{2r}%
\end{array}%
\right) \right] \left( 
\begin{array}{c}
d_{1} \\ 
\vdots \\ 
d_{2r}%
\end{array}%
\right) \right) ^{t}=\left[ \left( 
\begin{array}{ccc}
d_{1} & \cdots & d_{2r}%
\end{array}%
\right) \right] \left( 
\begin{array}{c}
c_{1} \\ 
\vdots \\ 
c_{2r}%
\end{array}%
\right) =d_{1}\circ c_{1}+\cdots +d_{2r}\circ c_{2r}\equiv d_{1}c_{1}+\cdots
+d_{2r}c_{2r}.
\end{equation}
\end{enumerate}
\end{enumerate}
\end{notation}

\begin{enumerate}
\item 
\begin{equation}
\mathbb{T}_{a}=\left( 
\begin{array}{cc}
\mathbb{U} & \overline{\mathbb{V}} \\ 
\mathbb{V} & \overline{\mathbb{U}}%
\end{array}%
\right) ,
\end{equation}

\item 
\begin{equation}
\mathbb{T}_{x}=\left( 
\begin{array}{cc}
\func{Re}\left( \mathbb{U+V}\right) & -\func{Im}\left( \mathbb{U+V}\right)
\\ 
\func{Im}\left( \mathbb{U-V}\right) & \func{Re}\left( \mathbb{U-V}\right)%
\end{array}%
\right) .
\end{equation}

\item 
\begin{equation}
\mathbb{K}=\left( 
\begin{array}{cc}
\frac{1}{\sqrt{2}}\mathbb{I}_{r} & \frac{i}{\sqrt{2}}\mathbb{I}_{r} \\ 
\frac{1}{\sqrt{2}}\mathbb{I}_{r} & -\frac{i}{\sqrt{2}}\mathbb{I}_{r}%
\end{array}%
\right)
\end{equation}

\item 
\begin{equation*}
\left( \mathbf{x}_{1},\mathbf{x}_{2}\right) =\left( \mathbf{a}^{\dagger },%
\mathbf{a}\right) \mathbb{K}=\left( \mathbf{a}^{\dagger },\mathbf{a}\right)
\left( 
\begin{array}{cc}
\frac{1}{\sqrt{2}}\mathbb{I}_{r} & \frac{i}{\sqrt{2}}\mathbb{I}_{r} \\ 
\frac{1}{\sqrt{2}}\mathbb{I}_{r} & -\frac{i}{\sqrt{2}}\mathbb{I}_{r}%
\end{array}%
\right)
\end{equation*}%
\begin{align}
\left( \mathbf{c}^{\dagger },\mathbf{c}\right) & =\left( \mathbf{a}^{\dagger
},\mathbf{a}\right) \left( 
\begin{array}{cc}
\mathbb{T} & \mathbb{O} \\ 
\mathbb{O} & \overline{\mathbb{T}}%
\end{array}%
\right) \\
\left( \mathbf{y}_{1},\mathbf{y}_{2}\right) & =\left( \mathbf{x}_{1},\mathbf{%
x}_{2}\right) \left( 
\begin{array}{cc}
\func{Re}\left( \mathbb{T}\right) & -\func{Im}\left( \mathbb{T}\right) \\ 
\func{Im}\left( \mathbb{T}\right) & \func{Re}\left( \mathbb{T}\right)%
\end{array}%
\right) ,
\end{align}

\begin{enumerate}
\item 
\begin{equation}
\left( x_{1},x_{2}\right) =\left( a_{1}^{\dagger },a_{1}\right) \mathbb{K}%
=\left( a_{1}^{\dagger },a_{1}\right) \frac{1}{\sqrt{2}}\left( 
\begin{array}{cc}
1 & i1 \\ 
1 & -i1%
\end{array}%
\right) =\left( \overset{}{\left( \overset{}{\left( a_{1}^{\dagger
},a_{1}\right) \mathbb{K}}\right) _{1},\left( \overset{}{\left(
a_{1}^{\dagger },a_{1}\right) \mathbb{K}}\right) _{2}}\right) =\left(
a_{1}^{\dagger }+a_{1},i\left( a_{1}^{\dagger }-a_{1}\right) \right)
\end{equation}

\item 
\begin{eqnarray}
\left\{ x_{1},x_{1}\right\} &=&\frac{1}{2}\left\{ a_{1}^{\dagger
}+a_{1},a_{1}^{\dagger }+a_{1}\right\} =\frac{1}{2}\left\{ a_{1}^{\dagger
},a_{1}^{\dagger }\right\} +\frac{1}{2}\left\{ a_{1}^{\dagger
},a_{1}\right\} +\frac{1}{2}\left\{ a_{1},a_{1}^{\dagger }\right\} +\frac{1}{%
2}\left\{ a_{1},a_{1}\right\} =1 \\
\left\{ x_{1},x_{2}\right\} &=&\frac{1}{2}\left\{ a_{1}^{\dagger
}+a_{1},i\left( a_{1}^{\dagger }-a_{1}\right) \right\} =\frac{i}{2}\left\{
a_{1}^{\dagger },a_{1}^{\dagger }\right\} -\frac{i}{2}\left\{ a_{1}^{\dagger
},a_{1}\right\} +\frac{i}{2}\left\{ a_{1},a_{1}^{\dagger }\right\} -\frac{i}{%
2}\left\{ a_{1},a_{1}\right\} =0 \\
\left\{ x_{2},x_{2}\right\} &=&\frac{1}{2}\left\{ i\left( a_{1}^{\dagger
}-a_{1}\right) ,i\left( a_{1}^{\dagger }-a_{1}\right) \right\} =-\frac{1}{2}%
\left\{ a_{1}^{\dagger },a_{1}^{\dagger }\right\} +\frac{1}{2}\left\{
a_{1}^{\dagger },a_{1}\right\} +\frac{1}{2}\left\{ a_{1},a_{1}^{\dagger
}\right\} -\frac{1}{2}\left\{ a_{1},a_{1}\right\} =1
\end{eqnarray}
\end{enumerate}

\begin{enumerate}
\item 
\begin{enumerate}
\item 
\begin{equation}
\left( \mathbf{a}^{\dagger }\right) \left( \mathbf{a}\right) ^{t}=\left( 
\begin{array}{c}
a_{1}^{\dagger } \\ 
a_{2}^{\dagger } \\ 
a_{3}^{\dagger }%
\end{array}%
\right) \left( 
\begin{array}{ccc}
a_{1} & a_{2} & a_{3}%
\end{array}%
\right) =\left( 
\begin{array}{ccc}
a_{1}^{\dagger }a_{1} & a_{1}^{\dagger }a_{2} & a_{1}^{\dagger }a_{3} \\ 
a_{2}^{\dagger }a_{1} & a_{2}^{\dagger }a_{2} & a_{2}^{\dagger }a_{3} \\ 
a_{3}^{\dagger }a_{1} & a_{3}^{\dagger }a_{2} & a_{3}^{\dagger }a_{3}%
\end{array}%
\right)
\end{equation}

\item 
\begin{eqnarray}
\left( \overset{}{\left( \mathbf{a}^{\dagger }\right) \left( \mathbf{a}%
\right) ^{t}}\right) ^{t} &=&\left( \mathbf{a}\right) \left( \mathbf{a}%
^{\dagger }\right) ^{t}=\left( \left( 
\begin{array}{c}
a_{1}^{\dagger } \\ 
a_{2}^{\dagger } \\ 
a_{3}^{\dagger }%
\end{array}%
\right) \left( 
\begin{array}{ccc}
a_{1} & a_{2} & a_{3}%
\end{array}%
\right) \right) ^{t}=\left( 
\begin{array}{ccc}
a_{1}^{\dagger }a_{1} & a_{1}^{\dagger }a_{2} & a_{1}^{\dagger }a_{3} \\ 
a_{2}^{\dagger }a_{1} & a_{2}^{\dagger }a_{2} & a_{2}^{\dagger }a_{3} \\ 
a_{3}^{\dagger }a_{1} & a_{3}^{\dagger }a_{2} & a_{3}^{\dagger }a_{3}%
\end{array}%
\right) ^{t} \\
&=&\left( 
\begin{array}{ccc}
\left( a_{1}^{\dagger }a_{1}\right) ^{t} & \left( a_{2}^{\dagger
}a_{1}\right) ^{t} & \left( a_{3}^{\dagger }a_{1}\right) ^{t} \\ 
\left( a_{1}^{\dagger }a_{2}\right) ^{t} & \left( a_{2}^{\dagger
}a_{2}\right) ^{t} & \left( a_{3}^{\dagger }a_{2}\right) ^{t} \\ 
\left( a_{1}^{\dagger }a_{3}\right) ^{t} & \left( a_{2}^{\dagger
}a_{3}\right) ^{t} & \left( a_{3}^{\dagger }a_{3}\right) ^{t}%
\end{array}%
\right) =\left( 
\begin{array}{ccc}
a_{1}a_{1}^{\dagger } & a_{1}a_{2}^{\dagger } & a_{1}a_{3}^{\dagger } \\ 
a_{2}a_{1}^{\dagger } & a_{2}a_{2}^{\dagger } & a_{2}a_{3}^{\dagger } \\ 
a_{3}a_{1}^{\dagger } & a_{3}a_{2}^{\dagger } & a_{3}a_{3}^{\dagger }%
\end{array}%
\right)
\end{eqnarray}

\item 
\begin{gather}
\left( \mathbf{a}^{\dagger }\right) \left( \mathbf{a}\right) ^{t}+\left( 
\overset{}{\left( \mathbf{a}\right) \left( \mathbf{a}^{\dagger }\right) }%
\right) ^{t}=\left( 
\begin{array}{c}
a_{1}^{\dagger } \\ 
a_{2}^{\dagger } \\ 
a_{3}^{\dagger }%
\end{array}%
\right) \left( 
\begin{array}{ccc}
a_{1} & a_{2} & a_{3}%
\end{array}%
\right) +\left( 
\begin{array}{c}
a_{1} \\ 
a_{2} \\ 
a_{3}%
\end{array}%
\right) \left( 
\begin{array}{ccc}
a_{1}^{\dagger } & a_{2}^{\dagger } & a_{3}^{\dagger }%
\end{array}%
\right) \\
=\left( 
\begin{array}{ccc}
a_{1}^{\dagger }a_{1} & a_{1}^{\dagger }a_{2} & a_{1}^{\dagger }a_{3} \\ 
a_{2}^{\dagger }a_{1} & a_{2}^{\dagger }a_{2} & a_{2}^{\dagger }a_{3} \\ 
a_{3}^{\dagger }a_{1} & a_{3}^{\dagger }a_{2} & a_{3}^{\dagger }a_{3}%
\end{array}%
\right) +\left( 
\begin{array}{ccc}
a_{1}a_{1}^{\dagger } & a_{1}a_{2}^{\dagger } & a_{1}a_{3}^{\dagger } \\ 
a_{2}a_{1}^{\dagger } & a_{2}a_{2}^{\dagger } & a_{2}a_{3}^{\dagger } \\ 
a_{3}a_{1}^{\dagger } & a_{3}a_{2}^{\dagger } & a_{3}a_{3}^{\dagger }%
\end{array}%
\right) ^{t}=\left( 
\begin{array}{ccc}
a_{1}^{\dagger }a_{1} & a_{1}^{\dagger }a_{2} & a_{1}^{\dagger }a_{3} \\ 
a_{2}^{\dagger }a_{1} & a_{2}^{\dagger }a_{2} & a_{2}^{\dagger }a_{3} \\ 
a_{3}^{\dagger }a_{1} & a_{3}^{\dagger }a_{2} & a_{3}^{\dagger }a_{3}%
\end{array}%
\right) +\left( 
\begin{array}{ccc}
a_{1}a_{1}^{\dagger } & a_{2}a_{1}^{\dagger } & a_{3}a_{1}^{\dagger } \\ 
a_{1}a_{2}^{\dagger } & a_{2}a_{2}^{\dagger } & a_{3}a_{2}^{\dagger } \\ 
a_{1}a_{3}^{\dagger } & a_{2}a_{3}^{\dagger } & a_{3}a_{3}^{\dagger }%
\end{array}%
\right) \\
=\left( 
\begin{array}{ccc}
a_{1}^{\dagger }a_{1}+a_{1}a_{1}^{\dagger } & a_{1}^{\dagger
}a_{2}+a_{2}a_{1}^{\dagger } & a_{1}^{\dagger }a_{3}+a_{3}a_{1}^{\dagger }
\\ 
a_{2}^{\dagger }a_{1}+a_{1}a_{2}^{\dagger } & a_{2}^{\dagger
}a_{2}+a_{2}a_{2}^{\dagger } & a_{2}^{\dagger }a_{3}+a_{3}a_{2}^{\dagger }
\\ 
a_{3}^{\dagger }a_{1}+a_{1}a_{3}^{\dagger } & a_{3}^{\dagger
}a_{2}+a_{2}a_{3}^{\dagger } & a_{3}^{\dagger }a_{3}+a_{3}a_{3}^{\dagger }%
\end{array}%
\right) =\left( 
\begin{array}{ccc}
\left\{ a_{1}^{\dagger },a_{1}\right\} & \left\{ a_{1}^{\dagger
},a_{2}\right\} & \left\{ a_{1}^{\dagger },a_{3}\right\} \\ 
\left\{ a_{2}^{\dagger },a_{1}\right\} & \left\{ a_{2}^{\dagger
},a_{2}\right\} & \left\{ a_{2}^{\dagger },a_{3}\right\} \\ 
\left\{ a_{3}^{\dagger },a_{1}\right\} & \left\{ a_{3}^{\dagger
},a_{2}\right\} & \left\{ a_{3}^{\dagger },a_{3}\right\}%
\end{array}%
\right)
\end{gather}

\item 
\begin{eqnarray}
\left( 
\begin{array}{c}
\mathbf{a}^{\dagger } \\ 
\mathbf{a}%
\end{array}%
\right) \left( \mathbf{a}^{\dagger },\mathbf{a}\right) &=&\left( 
\begin{array}{cc}
\left( \mathbf{a}^{\dagger }\right) \left( \mathbf{a}^{\dagger }\right) ^{t}
& \left( \mathbf{a}^{\dagger }\right) \left( \mathbf{a}\right) ^{t} \\ 
\left( \mathbf{a}\right) \left( \mathbf{a}^{\dagger }\right) ^{t} & \left( 
\mathbf{a}\right) \left( \mathbf{a}\right) ^{t}%
\end{array}%
\right) =\left( 
\begin{array}{cc}
\left( 
\begin{array}{c}
a_{1}^{\dagger } \\ 
a_{2}^{\dagger } \\ 
a_{3}^{\dagger }%
\end{array}%
\right) \left( 
\begin{array}{ccc}
a_{1}^{\dagger } & a_{2}^{\dagger } & a_{3}^{\dagger }%
\end{array}%
\right) & \left( 
\begin{array}{c}
a_{1}^{\dagger } \\ 
a_{2}^{\dagger } \\ 
a_{3}^{\dagger }%
\end{array}%
\right) \left( 
\begin{array}{ccc}
a_{1} & a_{2} & a_{3}%
\end{array}%
\right) \\ 
\left( 
\begin{array}{c}
a_{1} \\ 
a_{2} \\ 
a_{3}%
\end{array}%
\right) \left( 
\begin{array}{ccc}
a_{1}^{\dagger } & a_{2}^{\dagger } & a_{3}^{\dagger }%
\end{array}%
\right) & \left( 
\begin{array}{c}
a_{1} \\ 
a_{2} \\ 
a_{3}%
\end{array}%
\right) \left( 
\begin{array}{ccc}
a_{1} & a_{2} & a_{3}%
\end{array}%
\right)%
\end{array}%
\right) \\
&=&\left( 
\begin{array}{cc}
\left( 
\begin{array}{ccc}
a_{1}^{\dagger }a_{1}^{\dagger } & a_{1}^{\dagger }a_{2}^{\dagger } & 
a_{1}^{\dagger }a_{3}^{\dagger } \\ 
a_{2}^{\dagger }a_{1}^{\dagger } & a_{2}^{\dagger }a_{2}^{\dagger } & 
a_{2}^{\dagger }a_{3}^{\dagger } \\ 
a_{3}^{\dagger }a_{1}^{\dagger } & a_{3}^{\dagger }a_{2}^{\dagger } & 
a_{3}^{\dagger }a_{3}^{\dagger }%
\end{array}%
\right) & \left( 
\begin{array}{ccc}
a_{1}^{\dagger }a_{1} & a_{1}^{\dagger }a_{2} & a_{1}^{\dagger }a_{3} \\ 
a_{2}^{\dagger }a_{1} & a_{2}^{\dagger }a_{2} & a_{2}^{\dagger }a_{3} \\ 
a_{3}^{\dagger }a_{1} & a_{3}^{\dagger }a_{2} & a_{3}^{\dagger }a_{3}%
\end{array}%
\right) \\ 
\left( 
\begin{array}{ccc}
a_{1}a_{1}^{\dagger } & a_{1}a_{2}^{\dagger } & a_{1}a_{3}^{\dagger } \\ 
a_{2}a_{1}^{\dagger } & a_{2}a_{2}^{\dagger } & a_{2}a_{3}^{\dagger } \\ 
a_{3}a_{1}^{\dagger } & a_{3}a_{2}^{\dagger } & a_{3}a_{3}^{\dagger }%
\end{array}%
\right) & \left( 
\begin{array}{ccc}
a_{1}a_{1} & a_{1}a_{2} & a_{1}a_{3} \\ 
a_{2}a_{1} & a_{2}a_{2} & a_{2}a_{3} \\ 
a_{3}a_{1} & a_{3}a_{2} & a_{3}a_{3}%
\end{array}%
\right)%
\end{array}%
\right)
\end{eqnarray}

\item 
\begin{gather}
\left( 
\begin{array}{cc}
\left( \mathbf{a}^{\dagger }\right) \left( \mathbf{a}^{\dagger }\right) ^{t}
& \left( \mathbf{a}^{\dagger }\right) \left( \mathbf{a}\right) ^{t} \\ 
\left( \mathbf{a}\right) \left( \mathbf{a}^{\dagger }\right) ^{t} & \left( 
\mathbf{a}\right) \left( \mathbf{a}\right) ^{t}%
\end{array}%
\right) ^{t}=\left( 
\begin{array}{cc}
\left( \left( \mathbf{a}^{\dagger }\right) \left( \mathbf{a}^{\dagger
}\right) ^{t}\right) ^{t} & \left( \left( \mathbf{a}\right) \left( \mathbf{a}%
^{\dagger }\right) ^{t}\right) ^{t} \\ 
\left( \left( \mathbf{a}^{\dagger }\right) \left( \mathbf{a}\right)
^{t}\right) ^{t} & \left( \left( \mathbf{a}\right) \left( \mathbf{a}\right)
^{t}\right) ^{t}%
\end{array}%
\right) \\
\parallel \\
\left( 
\begin{array}{cc}
\left( 
\begin{array}{c}
a_{1}^{\dagger } \\ 
a_{2}^{\dagger } \\ 
a_{3}^{\dagger }%
\end{array}%
\right) \left( 
\begin{array}{ccc}
a_{1}^{\dagger } & a_{2}^{\dagger } & a_{3}^{\dagger }%
\end{array}%
\right) & \left( 
\begin{array}{c}
a_{1}^{\dagger } \\ 
a_{2}^{\dagger } \\ 
a_{3}^{\dagger }%
\end{array}%
\right) \left( 
\begin{array}{ccc}
a_{1} & a_{2} & a_{3}%
\end{array}%
\right) \\ 
\left( 
\begin{array}{c}
a_{1} \\ 
a_{2} \\ 
a_{3}%
\end{array}%
\right) \left( 
\begin{array}{ccc}
a_{1}^{\dagger } & a_{2}^{\dagger } & a_{3}^{\dagger }%
\end{array}%
\right) & \left( 
\begin{array}{c}
a_{1} \\ 
a_{2} \\ 
a_{3}%
\end{array}%
\right) \left( 
\begin{array}{ccc}
a_{1} & a_{2} & a_{3}%
\end{array}%
\right)%
\end{array}%
\right) ^{t}=\left( 
\begin{array}{cc}
\left( 
\begin{array}{c}
a_{1}^{\dagger } \\ 
a_{2}^{\dagger } \\ 
a_{3}^{\dagger }%
\end{array}%
\right) \left( 
\begin{array}{ccc}
a_{1}^{\dagger } & a_{2}^{\dagger } & a_{3}^{\dagger }%
\end{array}%
\right) & \left( 
\begin{array}{c}
a_{1} \\ 
a_{2} \\ 
a_{3}%
\end{array}%
\right) \left( 
\begin{array}{ccc}
a_{1}^{\dagger } & a_{2}^{\dagger } & a_{3}^{\dagger }%
\end{array}%
\right) \\ 
\left( 
\begin{array}{c}
a_{1}^{\dagger } \\ 
a_{2}^{\dagger } \\ 
a_{3}^{\dagger }%
\end{array}%
\right) \left( 
\begin{array}{ccc}
a_{1} & a_{2} & a_{3}%
\end{array}%
\right) & \left( 
\begin{array}{c}
a_{1} \\ 
a_{2} \\ 
a_{3}%
\end{array}%
\right) \left( 
\begin{array}{ccc}
a_{1} & a_{2} & a_{3}%
\end{array}%
\right)%
\end{array}%
\right)
\end{gather}

\item 
\begin{equation}
\left( 
\begin{array}{c}
\mathbf{a}^{\dagger } \\ 
\mathbf{a}%
\end{array}%
\right) \left( \mathbf{a}^{\dagger },\mathbf{a}\right) =\left( 
\begin{array}{cc}
\mathbf{a}^{\dagger }\mathbf{a}^{\dagger } & \mathbf{a}^{\dagger }\mathbf{a}
\\ 
\mathbf{aa}^{\dagger } & \mathbf{aa}%
\end{array}%
\right) =\left( 
\begin{array}{cc}
\left( 
\begin{array}{c}
a_{1}^{\dagger } \\ 
\vdots \\ 
a_{r}^{\dagger }%
\end{array}%
\right) \left( 
\begin{array}{ccc}
a_{1}^{\dagger } & \cdots & a_{r}^{\dagger }%
\end{array}%
\right) & \left( 
\begin{array}{c}
a_{1}^{\dagger } \\ 
\vdots \\ 
a_{r}^{\dagger }%
\end{array}%
\right) \left( 
\begin{array}{ccc}
a_{1} & \cdots & a_{r}%
\end{array}%
\right) \\ 
\left( 
\begin{array}{c}
a_{1} \\ 
\vdots \\ 
a_{r}%
\end{array}%
\right) \left( 
\begin{array}{ccc}
a_{1}^{\dagger } & \cdots & a_{r}^{\dagger }%
\end{array}%
\right) & \left( 
\begin{array}{c}
a_{1} \\ 
\vdots \\ 
a_{r}%
\end{array}%
\right) \left( 
\begin{array}{ccc}
a_{1} & \cdots & a_{r}%
\end{array}%
\right)%
\end{array}%
\right)
\end{equation}

\item 
\begin{eqnarray}
\left( \boldsymbol{a}\right) \left( \boldsymbol{b}\right) ^{t} &=&\left( 
\begin{array}{c}
a_{1} \\ 
\vdots \\ 
a_{r}%
\end{array}%
\right) \left( 
\begin{array}{ccc}
b_{1} & \cdots & b_{r}%
\end{array}%
\right) =\left( 
\begin{array}{ccc}
a_{1}b_{1} & \cdots & a_{r}b_{1} \\ 
\vdots & \ddots & \vdots \\ 
a_{1}b_{r} & \cdots & a_{r}b_{r}%
\end{array}%
\right) \\
\left( \boldsymbol{b}\right) \left( \boldsymbol{a}\right) ^{t} &=&\left(
\left( 
\begin{array}{c}
a_{1} \\ 
\vdots \\ 
a_{r}%
\end{array}%
\right) \left( 
\begin{array}{ccc}
b_{1} & \cdots & b_{r}%
\end{array}%
\right) \right) ^{t}=\left( 
\begin{array}{c}
b_{1} \\ 
\vdots \\ 
b_{r}%
\end{array}%
\right) \left( 
\begin{array}{ccc}
a_{1} & \cdots & a_{r}%
\end{array}%
\right) =\left( 
\begin{array}{ccc}
b_{1}a_{1} & \cdots & b_{1}a_{r} \\ 
\vdots & \ddots & \vdots \\ 
b_{r}a_{1} & \cdots & b_{r}a_{r}%
\end{array}%
\right)
\end{eqnarray}%
and 
\begin{equation}
\left( 
\begin{array}{ccc}
a_{1}b_{1} & \cdots & a_{r}b_{1} \\ 
\vdots & \ddots & \vdots \\ 
a_{1}b_{r} & \cdots & a_{r}b_{r}%
\end{array}%
\right) ^{t}=\left( 
\begin{array}{ccc}
\left( a_{1}b_{1}\right) ^{t} & \cdots & \left( a_{1}b_{r}\right) ^{t} \\ 
\vdots & \ddots & \vdots \\ 
\left( a_{r}b_{1}\right) ^{t} & \cdots & \left( a_{r}b_{r}\right) ^{t}%
\end{array}%
\right) =\left( 
\begin{array}{ccc}
b_{1}a_{1} & \cdots & b_{1}a_{r} \\ 
\vdots & \ddots & \vdots \\ 
b_{r}a_{1} & \cdots & b_{r}a_{r}%
\end{array}%
\right) .
\end{equation}

\item and 
\begin{equation}
\left( \left[ \left( 
\begin{array}{c}
a_{1} \\ 
\vdots \\ 
a_{r}%
\end{array}%
\right) \right] \left( 
\begin{array}{ccc}
a_{1}^{\dagger } & \cdots & a_{r}^{\dagger }%
\end{array}%
\right) \right) ^{t}=\left[ \left( 
\begin{array}{c}
a_{1}^{\dagger } \\ 
\vdots \\ 
a_{r}^{\dagger }%
\end{array}%
\right) \right] \left( 
\begin{array}{ccc}
a_{1} & \cdots & a_{r}%
\end{array}%
\right)
\end{equation}

\item 
\begin{gather}
\left( \left( 
\begin{array}{c}
a_{1}^{\dagger } \\ 
\vdots \\ 
a_{r}^{\dagger } \\ 
a_{1} \\ 
\vdots \\ 
a_{r}%
\end{array}%
\right) \left( 
\begin{array}{cccccc}
a_{1}^{\dagger } & \cdots & a_{r}^{\dagger } & a_{1} & \cdots & a_{r}%
\end{array}%
\right) \right) ^{t} \\
\parallel \\
\left( 
\begin{array}{cc}
\left( 
\begin{array}{c}
a_{1}^{\dagger } \\ 
\vdots \\ 
a_{r}^{\dagger }%
\end{array}%
\right) \left( 
\begin{array}{ccc}
a_{1}^{\dagger } & \cdots & a_{r}^{\dagger }%
\end{array}%
\right) & \left( 
\begin{array}{c}
a_{1}^{\dagger } \\ 
\vdots \\ 
a_{r}^{\dagger }%
\end{array}%
\right) \left( 
\begin{array}{ccc}
a_{1} & \cdots & a_{r}%
\end{array}%
\right) \\ 
\left( 
\begin{array}{c}
a_{1} \\ 
\vdots \\ 
a_{r}%
\end{array}%
\right) \left( 
\begin{array}{ccc}
a_{1}^{\dagger } & \cdots & a_{r}^{\dagger }%
\end{array}%
\right) & \left( 
\begin{array}{c}
a_{1} \\ 
\vdots \\ 
a_{r}%
\end{array}%
\right) \left( 
\begin{array}{ccc}
a_{1} & \cdots & a_{r}%
\end{array}%
\right)%
\end{array}%
\right) ^{t} \\
=\left( 
\begin{array}{cc}
\left( \left( 
\begin{array}{c}
a_{1}^{\dagger } \\ 
\vdots \\ 
a_{r}^{\dagger }%
\end{array}%
\right) \left( 
\begin{array}{ccc}
a_{1}^{\dagger } & \cdots & a_{r}^{\dagger }%
\end{array}%
\right) \right) ^{t} & \left( \left( 
\begin{array}{c}
a_{1} \\ 
\vdots \\ 
a_{r}%
\end{array}%
\right) \left( 
\begin{array}{ccc}
a_{1}^{\dagger } & \cdots & a_{r}^{\dagger }%
\end{array}%
\right) \right) ^{t} \\ 
\left( \left( 
\begin{array}{c}
a_{1}^{\dagger } \\ 
\vdots \\ 
a_{r}^{\dagger }%
\end{array}%
\right) \left( 
\begin{array}{ccc}
a_{1} & \cdots & a_{r}%
\end{array}%
\right) \right) ^{t} & \left( \left( 
\begin{array}{c}
a_{1} \\ 
\vdots \\ 
a_{r}%
\end{array}%
\right) \left( 
\begin{array}{ccc}
a_{1} & \cdots & a_{r}%
\end{array}%
\right) \right) ^{t}%
\end{array}%
\right)
\end{gather}
\end{enumerate}

\item 
\begin{equation}
\left( 
\begin{array}{cc}
\left\{ a_{1}^{\dagger },a_{1}^{\dagger }\right\} & \left\{ a_{1}^{\dagger
},a_{1}\right\} \\ 
\left\{ a_{1},a_{1}^{\dagger }\right\} & \left\{ a_{1},a_{1}\right\}%
\end{array}%
\right) =\left( 
\begin{array}{cc}
a_{1}^{\dagger }a_{1}^{\dagger } & a_{1}^{\dagger }a_{1} \\ 
a_{1}a_{1}^{\dagger } & a_{1}a_{1}%
\end{array}%
\right) +\left( 
\begin{array}{cc}
a_{1}^{\dagger }a_{1}^{\dagger } & a_{1}a_{1}^{\dagger } \\ 
a_{1}^{\dagger }a_{1} & a_{1}a_{1}%
\end{array}%
\right) =\left( 
\begin{array}{cc}
a_{1}^{\dagger }a_{1}^{\dagger } & a_{1}^{\dagger }a_{1} \\ 
a_{1}a_{1}^{\dagger } & a_{1}a_{1}%
\end{array}%
\right) +\left( 
\begin{array}{cc}
a_{1}^{\dagger }a_{1}^{\dagger } & a_{1}^{\dagger }a_{1} \\ 
a_{1}a_{1}^{\dagger } & a_{1}a_{1}%
\end{array}%
\right) ^{t}=\left( 
\begin{array}{cc}
0 & 1 \\ 
1 & 0%
\end{array}%
\right)
\end{equation}

\item 
\begin{equation}
\left( 
\begin{array}{c}
x_{1} \\ 
x_{2}%
\end{array}%
\right) \left( 
\begin{array}{cc}
x_{1} & x_{2}%
\end{array}%
\right) =\left( 
\begin{array}{cc}
x_{1}x_{1} & x_{1}x_{2} \\ 
x_{2}x_{1} & x_{2}x_{2}%
\end{array}%
\right)
\end{equation}

\item 
\begin{equation}
\left( 
\begin{array}{c}
x_{1} \\ 
x_{2}%
\end{array}%
\right) \left( 
\begin{array}{cc}
x_{1} & x_{2}%
\end{array}%
\right) =\mathbb{K}^{t}\left( 
\begin{array}{c}
a_{1}^{\dagger } \\ 
a_{1}%
\end{array}%
\right) \left( a_{1}^{\dagger },a_{1}\right) \mathbb{K}=\mathbb{K}^{t}\left( 
\begin{array}{cc}
a_{1}^{\dagger }a_{1}^{\dagger } & a_{1}^{\dagger }a_{1} \\ 
a_{1}a_{1}^{\dagger } & a_{1}a_{1}%
\end{array}%
\right) \mathbb{K}
\end{equation}

\item 
\begin{equation}
\left( 
\begin{array}{cc}
\left\{ x_{1},x_{1}\right\} & \left\{ x_{1},x_{2}\right\} \\ 
\left\{ x_{2},x_{1}\right\} & \left\{ x_{2},x_{2}\right\}%
\end{array}%
\right) =\left( 
\begin{array}{cc}
x_{1}x_{1} & x_{1}x_{2} \\ 
x_{2}x_{1} & x_{2}x_{2}%
\end{array}%
\right) +\left( 
\begin{array}{cc}
x_{1}x_{1} & x_{2}x_{1} \\ 
x_{1}x_{2} & x_{2}x_{2}%
\end{array}%
\right) =\left( 
\begin{array}{cc}
x_{1}x_{1} & x_{1}x_{2} \\ 
x_{2}x_{1} & x_{2}x_{2}%
\end{array}%
\right) +\left( 
\begin{array}{cc}
x_{1}x_{1} & x_{1}x_{2} \\ 
x_{2}x_{1} & x_{2}x_{2}%
\end{array}%
\right) ^{t}=\left( 
\begin{array}{cc}
1 & 0 \\ 
0 & 1%
\end{array}%
\right)
\end{equation}

\item $\mathbb{K}^{t}$ or $\mathbb{K}^{\dagger }$%
\begin{gather}
\left( 
\begin{array}{cc}
\left\{ x_{1},x_{1}\right\} & \left\{ x_{1},x_{2}\right\} \\ 
\left\{ x_{2},x_{1}\right\} & \left\{ x_{2},x_{2}\right\}%
\end{array}%
\right) =\mathbb{K}^{t}\left( \left( 
\begin{array}{cc}
a_{1}^{\dagger }a_{1}^{\dagger } & a_{1}^{\dagger }a_{1} \\ 
a_{1}a_{1}^{\dagger } & a_{1}a_{1}%
\end{array}%
\right) +\left( 
\begin{array}{cc}
a_{1}^{\dagger }a_{1}^{\dagger } & a_{1}^{\dagger }a_{1} \\ 
a_{1}a_{1}^{\dagger } & a_{1}a_{1}%
\end{array}%
\right) ^{t}\right) \mathbb{K} \\
\parallel \\
\mathbb{K}^{t}\left( \left( 
\begin{array}{cc}
a_{1}^{\dagger }a_{1}^{\dagger } & a_{1}^{\dagger }a_{1} \\ 
a_{1}a_{1}^{\dagger } & a_{1}a_{1}%
\end{array}%
\right) +\left( 
\begin{array}{cc}
a_{1}^{\dagger }a_{1}^{\dagger } & a_{1}a_{1}^{\dagger } \\ 
a_{1}^{\dagger }a_{1} & a_{1}a_{1}%
\end{array}%
\right) \right) \mathbb{K=K}^{t}\left( \left( 
\begin{array}{cc}
\left\{ a_{1}^{\dagger },a_{1}^{\dagger }\right\} & \left\{
a_{1},a_{1}^{\dagger }\right\} \\ 
\left\{ a_{1}^{\dagger },a_{1}\right\} & \left\{ a_{1},a_{1}\right\}%
\end{array}%
\right) \right) \mathbb{K} \\
\left. \parallel \right\vert \\
\left( 
\begin{array}{cc}
1 & 0 \\ 
0 & 1%
\end{array}%
\right) =\frac{1}{\sqrt{2}}\left( 
\begin{array}{cc}
1 & 1 \\ 
i1 & -i1%
\end{array}%
\right) \left( 
\begin{array}{cc}
0 & 1 \\ 
1 & 0%
\end{array}%
\right) \frac{1}{\sqrt{2}}\left( 
\begin{array}{cc}
1 & i1 \\ 
1 & -i1%
\end{array}%
\right)
\end{gather}%
$\lceil $%
\begin{equation}
\frac{1}{\sqrt{2}}\left( 
\begin{array}{cc}
1 & 1 \\ 
i1 & -i1%
\end{array}%
\right) \left( 
\begin{array}{cc}
0 & 1 \\ 
1 & 0%
\end{array}%
\right) \frac{1}{\sqrt{2}}\left( 
\begin{array}{cc}
1 & i1 \\ 
1 & -i1%
\end{array}%
\right) =\frac{1}{2}\left( 
\begin{array}{cc}
1 & 1 \\ 
i1 & -i1%
\end{array}%
\right) \left( 
\begin{array}{cc}
1 & -i1 \\ 
1 & i1%
\end{array}%
\right) =\frac{1}{2}\left( 
\begin{array}{cc}
2 & 0 \\ 
0 & 2%
\end{array}%
\right) =\left( 
\begin{array}{cc}
1 & 0 \\ 
0 & 1%
\end{array}%
\right)
\end{equation}%
$\rfloor $
\end{enumerate}

\item 
\begin{align}
\mathbb{T}_{r}& =\left( 
\begin{array}{cc}
\frac{1}{\sqrt{2}}\mathbb{I}_{r} & \frac{1}{\sqrt{2}}\mathbb{I}_{r} \\ 
-\frac{i}{\sqrt{2}}\mathbb{I}_{r} & \frac{i}{\sqrt{2}}\mathbb{I}_{r}%
\end{array}%
\right) \left( 
\begin{array}{cc}
\mathbb{U} & \overline{\mathbb{V}} \\ 
\mathbb{V} & \overline{\mathbb{U}}%
\end{array}%
\right) \left( 
\begin{array}{cc}
\frac{1}{\sqrt{2}}\mathbb{I}_{r} & \frac{i}{\sqrt{2}}\mathbb{I}_{r} \\ 
\frac{1}{\sqrt{2}}\mathbb{I}_{r} & \frac{-i}{\sqrt{2}}\mathbb{I}_{r}%
\end{array}%
\right) \\
& =\left( 
\begin{array}{cc}
\frac{1}{\sqrt{2}}\mathbb{U+}\frac{1}{\sqrt{2}}\mathbb{V} & \frac{1}{\sqrt{2}%
}\overline{\mathbb{V}}+\frac{1}{\sqrt{2}}\overline{\mathbb{U}} \\ 
-\frac{i}{\sqrt{2}}\mathbb{U+}\frac{i}{\sqrt{2}}\mathbb{V} & -\frac{i}{\sqrt{%
2}}\overline{\mathbb{V}}+\frac{i}{\sqrt{2}}\overline{\mathbb{U}}%
\end{array}%
\right) \left( 
\begin{array}{cc}
\frac{1}{\sqrt{2}}\mathbb{I}_{r} & \frac{i}{\sqrt{2}}\mathbb{I}_{r} \\ 
\frac{1}{\sqrt{2}}\mathbb{I}_{r} & \frac{-i}{\sqrt{2}}\mathbb{I}_{r}%
\end{array}%
\right) \\
& =\left( 
\begin{array}{cc}
\frac{1}{2}\mathbb{U}+\frac{1}{2}\mathbb{\overline{\mathbb{U}}+}\frac{1}{2}%
\mathbb{V+}\frac{1}{2}\overline{\mathbb{V}} & \frac{i}{2}\mathbb{U-}\frac{i}{%
2}\mathbb{\overline{\mathbb{U}}+}\frac{i}{2}\mathbb{V-}\frac{i}{2}\overline{%
\mathbb{V}} \\ 
\frac{i}{2}\overline{\mathbb{U}}-\frac{i}{2}\mathbb{U+}\frac{i}{2}\mathbb{V}-%
\frac{i}{2}\overline{\mathbb{V}} & \frac{1}{2}\mathbb{U}+\frac{1}{2}\mathbb{%
\overline{\mathbb{U}}-}\frac{1}{2}\mathbb{V}-\frac{1}{2}\overline{\mathbb{V}}%
\end{array}%
\right) \\
& =\frac{1}{2}\left( 
\begin{array}{cc}
2\func{Re}\left( \mathbb{U+V}\right) & -2\func{Im}\left( \mathbb{U+V}\right)
\\ 
2\func{Im}\left( \mathbb{U-V}\right) & 2\func{Re}\left( \mathbb{U-V}\right)%
\end{array}%
\right) =\left( 
\begin{array}{cc}
\func{Re}\left( \mathbb{U+V}\right) & -\func{Im}\left( \mathbb{U+V}\right)
\\ 
\func{Im}\left( \mathbb{U-V}\right) & \func{Re}\left( \mathbb{U-V}\right)%
\end{array}%
\right) .
\end{align}

\item 
\begin{align}
& \left( 
\begin{array}{cc}
\frac{1}{\sqrt{2}} & \frac{1}{\sqrt{2}} \\ 
-\frac{i}{\sqrt{2}} & \frac{i}{\sqrt{2}}%
\end{array}%
\right) \left( 
\begin{array}{cc}
u & \overline{v} \\ 
v & \overline{u}%
\end{array}%
\right) \left( 
\begin{array}{cc}
\frac{1}{\sqrt{2}} & \frac{i}{\sqrt{2}} \\ 
\frac{1}{\sqrt{2}} & \frac{-i}{\sqrt{2}}%
\end{array}%
\right) =\frac{1}{2}\left( 
\begin{array}{cc}
1 & 1 \\ 
-i & i%
\end{array}%
\right) \left( 
\begin{array}{cc}
u & \overline{v} \\ 
v & \overline{u}%
\end{array}%
\right) \left( 
\begin{array}{cc}
1 & i \\ 
1 & -i%
\end{array}%
\right) \\
& =\frac{1}{2}\left( 
\begin{array}{cc}
u+v & \overline{v}+\overline{u} \\ 
iv-iu & i\overline{u}-i\overline{v}%
\end{array}%
\right) \left( 
\begin{array}{cc}
1 & i \\ 
1 & -i%
\end{array}%
\right) =\frac{1}{2}\left( 
\begin{array}{cc}
u+v+\overline{v}+\overline{u} & i\left( u+v-\overline{v}-\overline{u}\right)
\\ 
i\left( v-u+\overline{u}-\overline{v}\right) & -v+u+\overline{u}-\overline{v}%
\end{array}%
\right) \\
& =\left( 
\begin{array}{cc}
\func{Re}\left( u+v\right) & -\func{Im}\left( u+v\right) \\ 
\func{Im}\left( u-v\right) & \func{Re}\left( u-v\right)%
\end{array}%
\right) ,
\end{align}

\item 
\begin{equation}
\mathcal{C}\left( 
\begin{array}{cc}
\boldsymbol{a}^{\dagger } & \boldsymbol{a}%
\end{array}%
\right) =\left( 
\begin{array}{cc}
\boldsymbol{a}^{\dagger } & \boldsymbol{a}%
\end{array}%
\right) \mathcal{R}_{a}\left( \mathcal{C}\right) =\left( 
\begin{array}{cccccc}
a_{1}^{\dagger } & \cdots & a_{r}^{\dagger } & a_{1} & \cdots & a_{r}%
\end{array}%
\right) \left( 
\begin{array}{cc}
\mathbb{U} & \overline{\mathbb{V}} \\ 
\mathbb{V} & \overline{\mathbb{U}}%
\end{array}%
\right) =\left( 
\begin{array}{cc}
\boldsymbol{a}^{\dagger } & \boldsymbol{a}%
\end{array}%
\right) \left( 
\begin{array}{cc}
\mathbb{U} & \overline{\mathbb{V}} \\ 
\mathbb{V} & \overline{\mathbb{U}}%
\end{array}%
\right) =\left( 
\begin{array}{cc}
\boldsymbol{b}^{\dagger } & \boldsymbol{b}%
\end{array}%
\right)
\end{equation}

\item 
\begin{equation}
\mathcal{C}\left( 
\begin{array}{cc}
\boldsymbol{a}^{\dagger } & \boldsymbol{a}%
\end{array}%
\right) =\left( 
\begin{array}{cc}
\boldsymbol{a}^{\dagger } & \boldsymbol{a}%
\end{array}%
\right) \mathcal{R}_{a}\left( \mathcal{C}\right) =\left( 
\begin{array}{cccc}
a_{1}^{\dagger } & a_{2}^{\dagger } & a_{1} & a_{2}%
\end{array}%
\right) \left( 
\begin{array}{cc}
\mathbb{U} & \overline{\mathbb{V}} \\ 
\mathbb{V} & \overline{\mathbb{U}}%
\end{array}%
\right) =\left( 
\begin{array}{cc}
\boldsymbol{a}^{\dagger } & \boldsymbol{a}%
\end{array}%
\right) \left( 
\begin{array}{cc}
\mathbb{U} & \overline{\mathbb{V}} \\ 
\mathbb{V} & \overline{\mathbb{U}}%
\end{array}%
\right) =\left( 
\begin{array}{cc}
\boldsymbol{b}^{\dagger } & \boldsymbol{b}%
\end{array}%
\right)
\end{equation}

\item 
\begin{equation}
\mathcal{C}\left( 
\begin{array}{cc}
\boldsymbol{a}^{\dagger } & \boldsymbol{a}%
\end{array}%
\right) =\left( 
\begin{array}{cc}
\boldsymbol{a}^{\dagger } & \boldsymbol{a}%
\end{array}%
\right) \mathcal{R}_{a}\left( \mathcal{C}\right) =\left( 
\begin{array}{cccccc}
a_{1}^{\dagger } & \cdots & a_{3}^{\dagger } & a_{1} & \cdots & a_{3}%
\end{array}%
\right) \left( 
\begin{array}{cc}
\mathbb{U} & \overline{\mathbb{V}} \\ 
\mathbb{V} & \overline{\mathbb{U}}%
\end{array}%
\right) =\left( 
\begin{array}{cc}
\boldsymbol{a}^{\dagger } & \boldsymbol{a}%
\end{array}%
\right) \left( 
\begin{array}{cc}
\mathbb{U} & \overline{\mathbb{V}} \\ 
\mathbb{V} & \overline{\mathbb{U}}%
\end{array}%
\right) =\left( 
\begin{array}{cc}
\boldsymbol{b}^{\dagger } & \boldsymbol{b}%
\end{array}%
\right)
\end{equation}

\item 
\begin{equation}
\left( 
\begin{array}{c}
a_{1}^{\dagger } \\ 
\vdots \\ 
a_{3}^{\dagger } \\ 
a_{1} \\ 
\vdots \\ 
a_{3}%
\end{array}%
\right) \left( 
\begin{array}{cccccc}
a_{1}^{\dagger } & \cdots & a_{3}^{\dagger } & a_{1} & \cdots & a_{3}%
\end{array}%
\right) =\left( 
\begin{array}{cccccc}
a_{1}^{\dagger }a_{1}^{\dagger } & a_{1}^{\dagger }a_{2}^{\dagger } & 
a_{1}^{\dagger }a_{3}^{\dagger } & a_{1}^{\dagger }a_{1} & a_{1}^{\dagger
}a_{2} & a_{1}^{\dagger }a_{3} \\ 
a_{2}^{\dagger }a_{1}^{\dagger } & a_{2}^{\dagger }a_{2}^{\dagger } & 
a_{2}^{\dagger }a_{3}^{\dagger } & a_{2}^{\dagger }a_{1} & a_{2}^{\dagger
}a_{2} & a_{2}^{\dagger }a_{3} \\ 
a_{3}^{\dagger }a_{1}^{\dagger } & a_{3}^{\dagger }a_{2}^{\dagger } & 
a_{3}^{\dagger }a_{3}^{\dagger } & a_{3}^{\dagger }a_{1} & a_{3}^{\dagger
}a_{2} & a_{3}^{\dagger }a_{3} \\ 
a_{1}a_{1}^{\dagger } & a_{1}a_{2}^{\dagger } & a_{1}a_{3}^{\dagger } & 
a_{1}a_{1} & a_{1}a_{2} & a_{1}a_{3} \\ 
a_{2}a_{1}^{\dagger } & a_{2}a_{2}^{\dagger } & a_{2}a_{3}^{\dagger } & 
a_{2}a_{1} & a_{2}a_{2} & a_{2}a_{3} \\ 
a_{3}a_{1}^{\dagger } & a_{3}a_{2}^{\dagger } & a_{3}a_{3}^{\dagger } & 
a_{3}a_{1} & a_{3}a_{2} & a_{3}a_{3}%
\end{array}%
\right)
\end{equation}

\item 
\begin{equation}
\left( 
\begin{array}{cc}
\boldsymbol{a} & \boldsymbol{b} \\ 
\boldsymbol{c} & \boldsymbol{d}%
\end{array}%
\right) \rightarrow \left( 
\begin{array}{cc}
\boldsymbol{d} & \boldsymbol{b} \\ 
\boldsymbol{c} & \boldsymbol{a}%
\end{array}%
\right)
\end{equation}%
\begin{equation}
\left( 
\begin{array}{cc}
\mathbb{O} & \mathbb{I}_{3} \\ 
\mathbb{I}_{3} & \mathbb{O}%
\end{array}%
\right) \left( 
\begin{array}{cc}
\boldsymbol{a}^{\dagger }\boldsymbol{a}^{\dagger } & \boldsymbol{a}^{\dagger
}\boldsymbol{a} \\ 
\boldsymbol{aa}^{\dagger } & \boldsymbol{aa}%
\end{array}%
\right) \left( 
\begin{array}{cc}
\mathbb{O} & \mathbb{I}_{3} \\ 
\mathbb{I}_{3} & \mathbb{O}%
\end{array}%
\right) =\left( 
\begin{array}{cc}
\boldsymbol{aa} & \boldsymbol{aa}^{\dagger } \\ 
\boldsymbol{a}^{\dagger }\boldsymbol{a} & \boldsymbol{a}^{\dagger }%
\boldsymbol{a}^{\dagger }%
\end{array}%
\right)
\end{equation}

\item 
\begin{eqnarray}
\left( 
\begin{array}{c}
a_{1}^{\dagger } \\ 
\vdots \\ 
a_{r}^{\dagger } \\ 
a_{1} \\ 
\vdots \\ 
a_{r}%
\end{array}%
\right) \left( 
\begin{array}{cccccc}
a_{1}^{\dagger } & \cdots & a_{r}^{\dagger } & a_{1} & \cdots & a_{r}%
\end{array}%
\right) &=&\left( 
\begin{array}{cccccc}
a_{1}^{\dagger }a_{1}^{\dagger } & \cdots & a_{1}^{\dagger }a_{r}^{\dagger }
& a_{1}^{\dagger }a_{1} & \cdots & a_{1}^{\dagger }a_{r} \\ 
\vdots & \ddots & \vdots & \vdots & \ddots & \vdots \\ 
a_{r}^{\dagger }a_{1}^{\dagger } & \cdots & a_{r}^{\dagger }a_{r}^{\dagger }
& a_{r}^{\dagger }a_{1} & \cdots & a_{r}^{\dagger }a_{r} \\ 
a_{1}a_{1}^{\dagger } & \cdots & a_{1}a_{r}^{\dagger } & a_{1}a_{1} & \cdots
& a_{1}a_{r} \\ 
\vdots & \ddots & \vdots & \vdots & \ddots & \vdots \\ 
a_{r}a_{1}^{\dagger } & \cdots & a_{r}a_{r}^{\dagger } & a_{r}a_{1} & \cdots
& a_{r}a_{r}%
\end{array}%
\right) \\
\left( 
\begin{array}{c}
\boldsymbol{a}^{\dagger } \\ 
\boldsymbol{a}%
\end{array}%
\right) \left( 
\begin{array}{cc}
\boldsymbol{a}^{\dagger } & \boldsymbol{a}%
\end{array}%
\right) &=&\left( 
\begin{array}{cc}
\boldsymbol{a}^{\dagger }\boldsymbol{a}^{\dagger } & \boldsymbol{a}^{\dagger
}\boldsymbol{a} \\ 
\boldsymbol{aa}^{\dagger } & \boldsymbol{aa}%
\end{array}%
\right) \\
\left( 
\begin{array}{c}
\boldsymbol{b}^{\dagger } \\ 
\boldsymbol{b}%
\end{array}%
\right) \left( 
\begin{array}{cc}
\boldsymbol{b}^{\dagger } & \boldsymbol{b}%
\end{array}%
\right) &=&\left( 
\begin{array}{cc}
\boldsymbol{b}^{\dagger }\boldsymbol{b}^{\dagger } & \boldsymbol{b}^{\dagger
}\boldsymbol{b} \\ 
\boldsymbol{bb}^{\dagger } & \boldsymbol{bb}%
\end{array}%
\right) =\left( 
\begin{array}{cc}
\mathbb{U} & \overline{\mathbb{V}} \\ 
\mathbb{V} & \overline{\mathbb{U}}%
\end{array}%
\right) ^{t}\left( 
\begin{array}{c}
\boldsymbol{a}^{\dagger } \\ 
\boldsymbol{a}%
\end{array}%
\right) \left( 
\begin{array}{cc}
\boldsymbol{a}^{\dagger } & \boldsymbol{a}%
\end{array}%
\right) \left( 
\begin{array}{cc}
\mathbb{U} & \overline{\mathbb{V}} \\ 
\mathbb{V} & \overline{\mathbb{U}}%
\end{array}%
\right)
\end{eqnarray}%
\begin{eqnarray}
\left( 
\begin{array}{cc}
\mathbb{O} & \mathbb{I} \\ 
\mathbb{I} & \mathbb{O}%
\end{array}%
\right) \left( 
\begin{array}{cc}
\boldsymbol{a}^{\dagger }\boldsymbol{a}^{\dagger } & \boldsymbol{a}^{\dagger
}\boldsymbol{a} \\ 
\boldsymbol{aa}^{\dagger } & \boldsymbol{aa}%
\end{array}%
\right) \left( 
\begin{array}{cc}
\mathbb{O} & \mathbb{I} \\ 
\mathbb{I} & \mathbb{O}%
\end{array}%
\right) &=&\left( 
\begin{array}{cc}
\boldsymbol{aa}^{\dagger } & \boldsymbol{aa} \\ 
\boldsymbol{a}^{\dagger }\boldsymbol{a}^{\dagger } & \boldsymbol{a}^{\dagger
}\boldsymbol{a}%
\end{array}%
\right) \left( 
\begin{array}{cc}
\mathbb{O} & \mathbb{I} \\ 
\mathbb{I} & \mathbb{O}%
\end{array}%
\right) =\left( 
\begin{array}{cc}
\boldsymbol{aa} & \boldsymbol{aa}^{\dagger } \\ 
\boldsymbol{a}^{\dagger }\boldsymbol{a} & \boldsymbol{a}^{\dagger }%
\boldsymbol{a}^{\dagger }%
\end{array}%
\right) \\
\left( 
\begin{array}{cc}
\boldsymbol{a}^{\dagger }\boldsymbol{a}^{\dagger } & \boldsymbol{a}^{\dagger
}\boldsymbol{a} \\ 
\boldsymbol{aa}^{\dagger } & \boldsymbol{aa}%
\end{array}%
\right) +\left( 
\begin{array}{cc}
\mathbb{O} & \mathbb{I} \\ 
\mathbb{I} & \mathbb{O}%
\end{array}%
\right) \left( 
\begin{array}{cc}
\boldsymbol{a}^{\dagger }\boldsymbol{a}^{\dagger } & \boldsymbol{a}^{\dagger
}\boldsymbol{a} \\ 
\boldsymbol{aa}^{\dagger } & \boldsymbol{aa}%
\end{array}%
\right) &=&\left( 
\begin{array}{cc}
\boldsymbol{a}^{\dagger }\boldsymbol{a}^{\dagger } & \boldsymbol{a}^{\dagger
}\boldsymbol{a} \\ 
\boldsymbol{aa}^{\dagger } & \boldsymbol{aa}%
\end{array}%
\right) +\left( 
\begin{array}{cc}
\boldsymbol{aa}^{\dagger } & \boldsymbol{aa} \\ 
\boldsymbol{a}^{\dagger }\boldsymbol{a}^{\dagger } & \boldsymbol{a}^{\dagger
}\boldsymbol{a}%
\end{array}%
\right)
\end{eqnarray}
\end{enumerate}

\begin{remark}
$\mathcal{M}\left( 2,2,\mathbb{R}\right) $ can be decomposed into the vector
space, (not algebraic)$,$ real direct sum 
\begin{equation}
\mathcal{M}\left( 2,2,\mathbb{R}\right) =\rho \left\{ \mathbb{C}\right\}
\oplus _{\mathbb{R}}\rho \left\{ \mathbb{C}\right\} ^{\bot },
\end{equation}%
where $\rho \left\{ \mathbb{C}\right\} ^{\bot }$ is the orthogonal
compliment wrt the HS inner product in $\mathcal{B}\left( \mathbb{R}%
^{2}\right) \cong \mathcal{M}\left( 2,2,\mathbb{R}\right) $ and is given by 
\begin{equation}
\rho \left\{ \mathbb{C}\right\} ^{\bot }=\left\{ \left. \left( 
\begin{array}{cc}
a & b \\ 
b & -a%
\end{array}%
\right) \right\vert a,b\in \mathbb{R}\right\} .
\end{equation}
\end{remark}

\begin{remark}
\begin{enumerate}
\item 
\begin{equation}
\left( 
\begin{array}{cc}
a & b \\ 
b & -a%
\end{array}%
\right) \left( 
\begin{array}{cc}
a & b \\ 
b & -a%
\end{array}%
\right) ^{t}=\left( 
\begin{array}{cc}
a & b \\ 
b & -a%
\end{array}%
\right) \left( 
\begin{array}{cc}
a & b \\ 
b & -a%
\end{array}%
\right) =\left( 
\begin{array}{cc}
a^{2}+b^{2} & 0 \\ 
0 & a^{2}+b^{2}%
\end{array}%
\right) .
\end{equation}

\item 
\begin{equation}
\det \left( \left( 
\begin{array}{cc}
a & b \\ 
b & -a%
\end{array}%
\right) \right) =-\left( a^{2}+b^{2}\right) <0.
\end{equation}

\item 
\begin{equation}
\left( 
\begin{array}{cc}
a & b \\ 
b & -a%
\end{array}%
\right) \left( 
\begin{array}{cc}
\cos \theta & -\sin \theta \\ 
\sin \theta & \cos \theta%
\end{array}%
\right) =\left( 
\begin{array}{cc}
a\cos \theta +b\sin \theta & -a\sin \theta +b\cos \theta \\ 
b\cos \theta -a\sin \theta & -\left( a\cos \theta +b\sin \theta \right)%
\end{array}%
\right) \Rightarrow \limfunc{Tr}\left\{ \left( 
\begin{array}{cc}
a & b \\ 
b & -a%
\end{array}%
\right) \left( 
\begin{array}{cc}
\cos \theta & -\sin \theta \\ 
\sin \theta & \cos \theta%
\end{array}%
\right) \right\} =0.
\end{equation}
\end{enumerate}
\end{remark}

\begin{remark}
The orthogonal group $O\left( 2,\mathbb{R}\right) $ is given by the disjoint
union of rotations and reflections 
\begin{equation}
O\left( 2,\mathbb{R}\right) =\left\{ \left. \left( 
\begin{array}{cc}
\cos \theta & -\sin \theta \\ 
\sin \theta & \cos \theta%
\end{array}%
\right) \right\vert \theta \in \left[ 0,2\pi \right] \right\} \cup \left\{
\left. \left( 
\begin{array}{cc}
-\cos \theta & \sin \theta \\ 
\sin \theta & \cos \theta%
\end{array}%
\right) \right\vert \theta \in \left\vert 0,2\pi \right) \right\} ,
\end{equation}%
as 
\begin{equation}
\det \left\{ \left( 
\begin{array}{cc}
-\cos \theta & \sin \theta \\ 
\sin \theta & \cos \theta%
\end{array}%
\right) \right\} =-\cos ^{2}\theta -\sin ^{2}\theta =-1.
\end{equation}
\end{remark}

\begin{remark}
If $O\in SO\left( 2,\mathbb{R}\right) \subset O\left( 2,\mathbb{R}\right) $
it is a rotation, while if $O\in O\left( 2,\mathbb{R}\right) $ but $\notin
SO\left( 2,\mathbb{R}\right) $ it is a reflection. If its is a reflection it
is not a complex number.
\end{remark}

Obviously $\rho $ is not surjective, as $\mathbb{R}$-$\dim \left\{ \rho
\left( \mathbb{C}\right) \right\} =2$ and $\mathbb{R}$-$\dim \left\{ M\left(
2,\mathbb{R}\right) \right\} =4,$ but it is injective. The map $\rho $ is a $%
c^{\ast }$-algebra homomorphism as 
\begin{align}
\rho \left( z_{1}z_{2}\right) & =\rho \left( z_{1}\right) \rho \left(
z_{2}\right) \\
\rho \left( \overline{z}\right) & =\rho \left( z\right) ^{\dagger }=\rho
\left( z\right) ^{t} \\
\rho \left( \alpha _{1}z_{1}+\alpha _{2}z_{2}\right) & =\alpha _{1}\rho
\left( z_{1}\right) +\alpha _{2}\rho \left( \alpha _{2}z_{2}\right) \\
\alpha _{1},\alpha _{2}& \in \mathbb{R}.
\end{align}%
$\lceil $That $\rho \left( \mathbb{C}\right) $ is closed under
multiplication is shown by%
\begin{equation}
\left( 
\begin{array}{cc}
x_{1} & -y_{1} \\ 
y_{1} & x_{1}%
\end{array}%
\right) \left( 
\begin{array}{cc}
x_{2} & -y_{2} \\ 
y_{2} & x_{2}%
\end{array}%
\right) =\left( 
\begin{array}{cc}
x_{1}x_{2}-y_{1}y_{2} & -\left( x_{1}y_{2}+y_{1}x_{2}\right) \\ 
y_{1}x_{2}+x_{1}y_{2} & x_{1}x_{2}-y_{1}y_{2}%
\end{array}%
\right) =\left( 
\begin{array}{cc}
x_{3} & -y_{3} \\ 
y_{3} & x_{3}%
\end{array}%
\right) .
\end{equation}%
$\rfloor $

The mapping $\rho $ is isometric as 
\begin{equation}
\left\Vert z\right\Vert ^{2}=\left\Vert \overline{z}z\right\Vert =\overline{z%
}z=x^{2}+y^{2}
\end{equation}%
and 
\begin{align}
\left\Vert \rho \left( z\right) ^{\dagger }\rho \left( z\right) \right\Vert
^{2}& =\limfunc{tr}\left\{ \rho \left( z\right) ^{\dagger }\rho \left(
z\right) \right\} =\limfunc{tr}\left\{ \left( 
\begin{array}{cc}
x & y \\ 
-y & x%
\end{array}%
\right) \left( 
\begin{array}{cc}
x & -y \\ 
y & x%
\end{array}%
\right) \right\} \\
& =\frac{1}{2}\func{Tr}\left\{ \left( 
\begin{array}{cc}
x & y \\ 
-y & x%
\end{array}%
\right) \left( 
\begin{array}{cc}
x & -y \\ 
y & x%
\end{array}%
\right) \right\} =\frac{1}{2}\func{Tr}\left\{ \left( 
\begin{array}{cc}
x^{2}+y^{2} & 0 \\ 
0 & x^{2}+y^{2}%
\end{array}%
\right) \right\} =x^{2}+y^{2}.
\end{align}

One can characterize $\rho \left( \mathbb{C}\right) $ in $\mathcal{M}\left(
2,2,\mathbb{R}\right) $ by observing that if 
\begin{equation}
\left[ A,\left( 
\begin{array}{cc}
0 & 1 \\ 
-1 & 0%
\end{array}%
\right) \right] =\left[ \left( 
\begin{array}{cc}
a & b \\ 
c & d%
\end{array}%
\right) ,\left( 
\begin{array}{cc}
0 & 1 \\ 
-1 & 0%
\end{array}%
\right) \right] =0
\end{equation}%
then 
\begin{align}
\left( 
\begin{array}{cc}
a & b \\ 
c & d%
\end{array}%
\right) \left( 
\begin{array}{cc}
0 & 1 \\ 
-1 & 0%
\end{array}%
\right) & =\left( 
\begin{array}{cc}
-b & a \\ 
-d & c%
\end{array}%
\right) \\
\left( 
\begin{array}{cc}
0 & 1 \\ 
-1 & 0%
\end{array}%
\right) \left( 
\begin{array}{cc}
a & b \\ 
c & d%
\end{array}%
\right) & =\left( 
\begin{array}{cc}
c & d \\ 
-a & -b%
\end{array}%
\right)
\end{align}%
and thus 
\begin{equation}
c=-b,a=d\Rightarrow A=\left( 
\begin{array}{cc}
a & -b \\ 
b & a%
\end{array}%
\right) ,
\end{equation}%
i.e. $A\in \rho \left( \mathbb{C}\right) $.

\subparagraph{ $\protect\rho \left( \mathbb{C}\right) $ in Terms of
Commutants of Complex Units \label{Commutants of Complex Units}}

Formally we can express this as the subalgebra $\rho \left( \mathbb{C}%
\right) \subset \mathcal{B}\left( \mathbb{R}^{2}\right) $ is characterized
by 
\begin{equation}
\rho \left( \mathbb{C}\right) =\ker \left\{ \widehat{J}\right\} =\left\{
\left. \overset{}{A}\right\vert \widehat{J}\left( A\right) =0,A\in \mathcal{B%
}\left( \mathbb{R}^{2}\right) \right\} ,
\end{equation}%
where 
\begin{equation}
\widehat{J}\left( A\right) =\left[ J,A\right]
\end{equation}%
and 
\begin{equation}
J=\left( 
\begin{array}{cc}
0 & -1 \\ 
1 & 0%
\end{array}%
\right) .
\end{equation}

\begin{lemma}
The image $\rho \left( \mathbb{C}\right) $ of the complex numbers $\mathbb{C}
$ in $\mathcal{B}\left( \mathbb{R}^{2}\right) $ is the commutatnt $C\left(
J\right) $ i.e. 
\begin{equation}
\rho \left( \mathbb{C}\right) =C\left( J\right) .
\end{equation}

\begin{proof}
\begin{equation}
C\left( J\right) =\left\{ \left. \overset{}{A\in \mathcal{B}\left( \mathbb{R}%
^{2}\right) }\right\vert \left[ J,A\right] =0\right\}
\end{equation}%
and this set was shown to be equal to $\rho \left( \mathbb{C}\right) .$
\end{proof}
\end{lemma}

The subspace $\ker \left\{ \widehat{J}\right\} \subset \mathcal{B}\left( 
\mathbb{R}^{2}\right) $ is invariant to the operator $\widehat{u_{J}}\in 
\mathcal{B}\left( \mathcal{B}\left( \mathbb{R}^{2}\right) \right) $ defined
by 
\begin{equation}
\widehat{u_{J}}\left( A\right) =e^{J}Ae^{-J}=e^{\widehat{J}}\left( A\right) ,
\end{equation}%
i.e. 
\begin{equation}
\widehat{u_{J}}:\ker \left\{ \widehat{J}\right\} \rightarrow \ker \left\{ 
\widehat{J}\right\} ,
\end{equation}%
where 
\begin{equation}
\widehat{u_{J}}^{t}\left( A\right) =e^{-J}Ae^{J}=e^{-\widehat{J}}\left(
A\right) .
\end{equation}

\begin{remark}
\begin{equation}
\rho \left( \mathbb{C}\right) =\ker \left\{ \widehat{J}\right\}
\end{equation}%
and 
\begin{equation}
\mathcal{B}\left( \mathbb{R}^{2}\right) =\rho \left( \mathbb{C}\right)
\oplus _{\mathbb{R}}\rho \left( \mathbb{C}\right) ^{c}=\ker \left\{ \widehat{%
J}\right\} \oplus _{\mathbb{R}}\ker \left\{ \widehat{J}\right\} ^{c}
\end{equation}
\end{remark}

The HS norm of $\rho \left( z\right) $ wrt the normalized trace is given by 
\begin{equation}
\frac{1}{2}\func{Tr}\left\{ \left( 
\begin{array}{cc}
a & -b \\ 
b & a%
\end{array}%
\right) \left( 
\begin{array}{cc}
a & b \\ 
-b & a%
\end{array}%
\right) \right\} =\frac{1}{2}\func{Tr}\left\{ \left( 
\begin{array}{cc}
a^{2}+b^{2} & 0 \\ 
0 & a^{2}+b^{2}%
\end{array}%
\right) \right\} =a^{2}+b^{2},
\end{equation}%
hence 
\begin{equation}
\left\vert z\right\vert ^{2}=\limfunc{tr}\left\{ \rho \left( z\right)
\right\} .
\end{equation}

Unit complex numbers 
\begin{equation}
\left\vert z\right\vert ^{2}=1
\end{equation}%
represent rigid rotations $\rho \left( z\right) \in O\left( 2,\mathbb{R}%
\right) $ of $\mathbb{R}^{2}$, which can be seen by oberving that 
\begin{equation}
\det \left\{ \rho \left( z\right) \right\} =a^{2}+b^{2}=\left\vert
z\right\vert ^{2}=1
\end{equation}%
showing that 
\begin{equation}
\rho \left( z\right) \in SO\left( 2,\mathbb{R}\right) .
\end{equation}

\begin{theorem}
\begin{equation}
\rho \left\{ \mathbb{C}\right\} =CO^{+}\left( 2,\mathbb{R}\right) .
\end{equation}

\begin{proof}
From the preceding.
\end{proof}
\end{theorem}

\begin{remark}
Non rigid rotations i.e. elastic rotations that stretch and compress are
represented by non unit complex numbers, $z=re^{i\theta },$where $r\neq 1$
and are represented by real matrices $\rho \left( z\right) \in CSO^{+}\left(
2,\mathbb{R}\right) $ that represent complex numbers$.$
\end{remark}

\begin{remark}
If $O\in O\left( 2,\mathbb{R}\right) $ and $\det \left\{ O\right\} =-1$ then 
$O\notin \rho \left( \mathbb{C}\right) .$ ????
\end{remark}

\paragraph{Symplectic Group $Sp\left( 2,\mathbb{R}\right) $}

If $S\in Sp\left( 2,\mathbb{R}\right) $ i.e. is symplectic 
\begin{equation}
SJS^{t}=J
\end{equation}%
then 
\begin{align}
\left( 
\begin{array}{cc}
a & b \\ 
c & d%
\end{array}%
\right) \left( 
\begin{array}{cc}
0 & -1 \\ 
1 & 0%
\end{array}%
\right) \left( 
\begin{array}{cc}
a & c \\ 
b & d%
\end{array}%
\right) & =\left( 
\begin{array}{cc}
0 & -1 \\ 
1 & 0%
\end{array}%
\right) \\
& \Rightarrow -\left( 
\begin{array}{cc}
b & -a \\ 
d & -c%
\end{array}%
\right) \left( 
\begin{array}{cc}
a & c \\ 
b & d%
\end{array}%
\right) =-\left( 
\begin{array}{cc}
0 & 1 \\ 
-1 & 0%
\end{array}%
\right) \\
& \Rightarrow -\left( 
\begin{array}{cc}
ab-ab & bc-ad \\ 
da-cb & dc-cd%
\end{array}%
\right) =-\left( 
\begin{array}{cc}
0 & bc-ad \\ 
da-cb & 0%
\end{array}%
\right) =\left( 
\begin{array}{cc}
0 & -1 \\ 
1 & 0%
\end{array}%
\right)
\end{align}%
and 
\begin{equation}
\left( 
\begin{array}{cc}
a & b \\ 
c & d%
\end{array}%
\right) \in Sp\left( 2,\mathbb{R}\right) \Rightarrow \det \left\{ \left( 
\begin{array}{cc}
a & b \\ 
c & d%
\end{array}%
\right) \right\} =da-cb=1.
\end{equation}%
Which shows that 
\begin{equation}
Sp\left( 2,\mathbb{R}\right) \cong SL\left( 2,\mathbb{R}\right) .
\end{equation}

\paragraph{Orthogonal Group $O\left( 2,\mathbb{R}\right) $}

If $O\in O\left( 2,\mathbb{R}\right) $ i.e. 
\begin{equation}
OO^{t}=I_{2}
\end{equation}%
then 
\begin{align}
\left( 
\begin{array}{cc}
a & b \\ 
c & d%
\end{array}%
\right) \left( 
\begin{array}{cc}
1 & 0 \\ 
0 & 1%
\end{array}%
\right) \left( 
\begin{array}{cc}
a & c \\ 
b & d%
\end{array}%
\right) & =\left( 
\begin{array}{cc}
1 & 0 \\ 
0 & 1%
\end{array}%
\right) \\
\left( 
\begin{array}{cc}
a & b \\ 
c & d%
\end{array}%
\right) \left( 
\begin{array}{cc}
a & c \\ 
b & d%
\end{array}%
\right) & =\left( 
\begin{array}{cc}
a^{2}+b^{2} & ac+bd \\ 
ac+bd & c^{2}+d^{2}%
\end{array}%
\right) =\left( 
\begin{array}{cc}
1 & 0 \\ 
0 & 1%
\end{array}%
\right)
\end{align}%
showing 
\begin{align}
a^{2}+b^{2}& =c^{2}+d^{2}=1 \\
ac& =-bd.
\end{align}%
If 
\begin{equation}
\det \left\{ O\right\} =ad-bc=1,
\end{equation}%
i.e. $O\in SO\left( 2,\mathbb{R}\right) $ then 
\begin{align}
ac+bd& =0\,\&\,ad-bc=1 \\
& \Rightarrow a^{2}d-abc=a \\
& \Rightarrow \left( a^{2}+b^{2}\right) d=a \\
& \Rightarrow d=a,ac+ba=0 \\
& \Rightarrow c=-b
\end{align}%
and thus%
\begin{equation}
O=\left( 
\begin{array}{cc}
a & b \\ 
c & d%
\end{array}%
\right) =\left( 
\begin{array}{cc}
a & -b \\ 
b & a%
\end{array}%
\right) ,
\end{equation}%
and hence 
\begin{equation}
O\in \rho \left( \mathbb{C}\right) .
\end{equation}%
As%
\begin{equation}
\det \left\{ \left( 
\begin{array}{cc}
a & -b \\ 
b & a%
\end{array}%
\right) \right\} =a^{2}+b^{2}=1
\end{equation}%
one has 
\begin{equation}
O\in SO\left( 2,\mathbb{R}\right) .
\end{equation}%
One can thus see that \emph{orthogonal rotations} of $\mathbb{R}^{2}$
correspond to unit complex numbers 
\begin{equation}
SO\left( 2,\mathbb{R}\right) \subset \rho \left( \mathbb{C}\right) \subset 
\mathcal{M}\left( 2,2,\mathbb{R}\right) .
\end{equation}

\paragraph{Indefinite Orthogonal Group $O\left( 1,1,\mathbb{R}\right) $}

If $A\in O\left( 1,1,\mathbb{R}\right) $ 
\begin{equation}
A\left( 
\begin{array}{cc}
1 & 0 \\ 
0 & -1%
\end{array}%
\right) A^{t}=\left( 
\begin{array}{cc}
1 & 0 \\ 
0 & -1%
\end{array}%
\right)
\end{equation}%
then%
\begin{align}
\left( 
\begin{array}{cc}
a & b \\ 
c & d%
\end{array}%
\right) \left( 
\begin{array}{cc}
1 & 0 \\ 
0 & -1%
\end{array}%
\right) \left( 
\begin{array}{cc}
a & c \\ 
b & d%
\end{array}%
\right) & =\left( 
\begin{array}{cc}
1 & 0 \\ 
0 & -1%
\end{array}%
\right) \\
\left( 
\begin{array}{cc}
a & -b \\ 
c & -d%
\end{array}%
\right) \left( 
\begin{array}{cc}
a & c \\ 
b & d%
\end{array}%
\right) & =\left( 
\begin{array}{cc}
a^{2}-b^{2} & ac-bd \\ 
ca-db & c^{2}-d^{2}%
\end{array}%
\right) =\left( 
\begin{array}{cc}
1 & 0 \\ 
0 & -1%
\end{array}%
\right) \\
a^{2}-b^{2}& =d^{2}-c^{2}=1;ac=bd.
\end{align}%
Any element $A\in $ $O\left( 1,1,\mathbb{R}\right) $ can be represented by
matrices belonging to $\mathcal{M}\left( 2,2,\mathbb{R}\right) $ of the form 
\begin{equation}
A=\left( 
\begin{array}{cc}
\cosh \theta & \sinh \theta \\ 
\sinh \theta & \cosh \theta%
\end{array}%
\right)
\end{equation}%
as 
\begin{eqnarray}
&&\left( 
\begin{array}{cc}
\cosh \theta & \sinh \theta \\ 
\sinh \theta & \cosh \theta%
\end{array}%
\right) \left( 
\begin{array}{cc}
1 & 0 \\ 
0 & -1%
\end{array}%
\right) \left( 
\begin{array}{cc}
\cosh \theta & \sinh \theta \\ 
\sinh \theta & \cosh \theta%
\end{array}%
\right) =\left( 
\begin{array}{cc}
\cosh \theta & \sinh \theta \\ 
\sinh \theta & \cosh \theta%
\end{array}%
\right) \left( 
\begin{array}{cc}
\cosh \theta & \sinh \theta \\ 
-\sinh \theta & -\cosh \theta%
\end{array}%
\right) \\
&=&\left( 
\begin{array}{cc}
\cosh ^{2}\theta -\sinh ^{2}\theta & \cosh \theta \sinh \theta -\sinh \theta
\cosh \theta \\ 
\cosh \theta \sinh \theta -\sinh \theta \cosh \theta & \sinh ^{2}\theta
-\cosh ^{2}\theta%
\end{array}%
\right) =\left( 
\begin{array}{cc}
1 & 0 \\ 
0 & -1%
\end{array}%
\right) .
\end{eqnarray}

\subparagraph{Complex Numbers}

Complex numbers can be expressed as 
\begin{eqnarray}
z_{1} &=&\left( x_{1},y_{1}\right) =r_{1}e^{i\theta _{1}}=r_{1}\left( \cos
\theta _{1}+i\sin \theta _{1}\right) \\
z_{2} &=&\left( x_{2},y_{2}\right) =r_{2}e^{i\theta _{2}}=r_{2}\left( \cos
\theta _{2}+i\sin \theta _{2}\right)
\end{eqnarray}%
and their sum as 
\begin{eqnarray}
z &=&re^{i\theta }=z_{1}+z_{2}=\left( x_{1}+x_{2}\right) +i\left(
y_{1}+y_{2}\right) =r\left( \cos \theta +i\sin \theta \right) \\
&=&\sqrt{\left( x_{1}+x_{2}\right) ^{2}+\left( y_{1}+y_{2}\right) ^{2}}\exp
i\left\{ \tan ^{-1}\left( \frac{\left( y_{1}+y_{2}\right) }{\left(
x_{1}+x_{2}\right) }\right) \right\} ,
\end{eqnarray}%
where%
\begin{eqnarray}
r &=&\sqrt{\left( x_{1}+x_{2}\right) ^{2}+\left( y_{1}+y_{2}\right) ^{2}} \\
\theta &=&\tan ^{-1}\left( \frac{\left( y_{1}+y_{2}\right) }{\left(
x_{1}+x_{2}\right) }\right) .
\end{eqnarray}

The general formula for product and division of complex numbers in polar
form is%
\begin{equation}
\frac{z_{1}z_{2}}{z_{3}z_{4}}=\frac{(a+ib)(c+id)}{(e+if)(g+ih)}=\sqrt{\frac{%
\left( a^{2}+b^{2}\right) (c^{2}+d^{2})}{(e^{2}+f^{2})(g^{2}+h^{2})}}\exp
i\left\{ \overset{}{\tan ^{-1}\left( \frac{b}{a}\right) +\tan ^{-1}\left( 
\frac{d}{c}\right) -\tan ^{-1}\left( \frac{f}{e}\right) -\tan ^{-1}\left( 
\frac{h}{g}\right) }\right\} .
\end{equation}

\subparagraph{Complex Matrices $\mathcal{M}\left( 2,2,\mathbb{C}\right) $}

If we consider $\mathcal{M}\left( 2,2,\mathbb{C}\right) $ then the group $%
O\left( 2,\mathbb{R}\right) $ is represented by matrices of the form%
\begin{equation}
U\left( u,v\right) =\left( 
\begin{array}{cc}
u & \overline{v} \\ 
-v & \overline{u}%
\end{array}%
\right) \in U\left( 2\right) \subset \mathcal{M}\left( 2,2,\mathbb{C}\right)
,
\end{equation}%
where%
\begin{eqnarray}
u &=&r_{1}\left( \cos \theta _{1}+i\sin \theta _{1}\right) \\
v &=&r_{2}\left( \cos \theta _{2}+i\sin \theta _{2}\right) ,
\end{eqnarray}%
and%
\begin{eqnarray}
r_{1} &=&\left\vert u\right\vert =\left( u\overline{u}\right) ^{\frac{1}{2}}
\\
\func{Re}\left\{ u\right\} &=&r_{1}\cos \theta _{1} \\
\func{Im}\left\{ u\right\} &=&r_{1}\sin \theta _{1} \\
r_{2} &=&\left\vert v\right\vert =\left( v\overline{v}\right) ^{\frac{1}{2}}
\\
\func{Re}\left\{ v\right\} &=&r_{2}\cos \theta _{2} \\
\func{Im}\left\{ v\right\} &=&r_{2}\sin \theta _{2}.
\end{eqnarray}%
The unitary condition 
\begin{eqnarray}
U\left( u,v\right) U\left( u,v\right) ^{\dagger } &=&\left( 
\begin{array}{cc}
u & \overline{v} \\ 
-v & \overline{u}%
\end{array}%
\right) \left( 
\begin{array}{cc}
\overline{u} & -\overline{v} \\ 
v & u%
\end{array}%
\right) =\left( 
\begin{array}{cc}
\left\vert u\right\vert ^{2}+\left\vert v\right\vert ^{2} & -u\overline{v}+%
\overline{v}u \\ 
-v\overline{u}+\overline{u}v & \left\vert u\right\vert ^{2}+\left\vert
v\right\vert ^{2}%
\end{array}%
\right) \\
&=&\left( 
\begin{array}{cc}
1 & 0 \\ 
0 & 1%
\end{array}%
\right) ,
\end{eqnarray}%
gives 
\begin{eqnarray}
\left\vert u\right\vert ^{2}+\left\vert v\right\vert ^{2} &=&1 \\
-u\overline{v}+\overline{v}u &=&0,
\end{eqnarray}%
or equivalently%
\begin{eqnarray}
\left( 
\begin{array}{cc}
r_{1}e^{i\theta _{1}} & r_{2}e^{-i\theta _{2}} \\ 
-r_{2}e^{i\theta _{2}} & r_{1}e^{-i\theta _{1}}%
\end{array}%
\right) \left( 
\begin{array}{cc}
r_{1}e^{-i\theta _{1}} & -r_{2}e^{-i\theta _{2}} \\ 
r_{2}e^{i\theta _{2}} & r_{1}e^{i\theta _{1}}%
\end{array}%
\right) &=&\left( 
\begin{array}{cc}
r_{1}^{2}+r_{2}^{2} & -r_{1}r_{2}e^{i\left( \theta _{1}-\theta _{2}\right)
}+r_{2}r_{1}e^{i\left( \theta _{1}-\theta _{2}\right) } \\ 
-r_{1}r_{2}e^{i\left( \theta _{1}-\theta _{2}\right) }+r_{2}r_{1}e^{i\left(
\theta _{1}-\theta _{2}\right) } & r_{1}^{2}+r_{2}^{2}%
\end{array}%
\right) \\
&=&\left( 
\begin{array}{cc}
1 & 0 \\ 
0 & 1%
\end{array}%
\right) ,
\end{eqnarray}%
\newline
so 
\begin{equation}
U\left( u,v\right) =\left( 
\begin{array}{cc}
u & \overline{v} \\ 
-v & \overline{u}%
\end{array}%
\right)
\end{equation}%
depends on two \emph{real} parameters $\left( r_{1},\theta _{1}\right) .$

These matrices can be transformed to real form by $\frac{1}{\sqrt{2}}\left( 
\begin{array}{cc}
1 & 1 \\ 
-i & i%
\end{array}%
\right) \in U\left( 2\right) \subset \mathcal{M}\left( 2,2,\mathbb{C}\right) 
$

$\lceil $which is shown by 
\begin{equation}
\frac{1}{\sqrt{2}}\left( 
\begin{array}{cc}
1 & 1 \\ 
-i & i%
\end{array}%
\right) \frac{1}{\sqrt{2}}\left( 
\begin{array}{cc}
1 & i \\ 
1 & -i%
\end{array}%
\right) =\frac{1}{2}\left( 
\begin{array}{cc}
2 & 0 \\ 
0 & 2%
\end{array}%
\right) =\left( 
\begin{array}{cc}
1 & 0 \\ 
0 & 1%
\end{array}%
\right) \in U\left( 2\right) .
\end{equation}%
$\rfloor $%
\begin{eqnarray}
\mathcal{R}_{e}\left( O\right) &=&\widehat{K}\left( \left( 
\begin{array}{cc}
u & \overline{v} \\ 
-v & \overline{u}%
\end{array}%
\right) \right) =\frac{1}{\sqrt{2}}\left( 
\begin{array}{cc}
1 & 1 \\ 
-i & i%
\end{array}%
\right) \left( \left( 
\begin{array}{cc}
u & \overline{v} \\ 
-v & \overline{u}%
\end{array}%
\right) \frac{1}{\sqrt{2}}\left( 
\begin{array}{cc}
1 & i \\ 
1 & -i%
\end{array}%
\right) \right) =\frac{1}{2}\left( 
\begin{array}{cc}
1 & 1 \\ 
-i & i%
\end{array}%
\right) \left( 
\begin{array}{cc}
\left( u+\overline{v}\right) & i\left( u-\overline{v}\right) \\ 
\left( -v+\overline{u}\right) & -i\left( v+\overline{u}\right)%
\end{array}%
\right) \\
&=&\frac{1}{2}\left( 
\begin{array}{cc}
u+\overline{v}-v+\overline{u} & i\left( u-\overline{v}\right) -i\left( v+%
\overline{u}\right) \\ 
-i\left( u+\overline{v}\right) +i\left( -v+\overline{u}\right) & -ii\left( u-%
\overline{v}\right) -ii\left( v+\overline{u}\right)%
\end{array}%
\right) =\frac{1}{2}\left( 
\begin{array}{cc}
u+\overline{u}\ +\overline{v}-v & i\left( u-\overline{u}\right) -i\left( 
\overline{v}+v\right) \\ 
i\left( \overline{u}-u\right) -i\left( \overline{v}+v\right) & \left( u-%
\overline{v}\right) +\left( v+\overline{u}\right)%
\end{array}%
\right) \\
&=&\frac{1}{2}\left( 
\begin{array}{cc}
u+\overline{u}\  & i\left( u-\overline{v}\right) -i\left( v+\overline{u}%
\right) \\ 
-i\left( u+\overline{v}\right) +i\left( -v+\overline{u}\right) & -u+%
\overline{v}+v+\overline{u}%
\end{array}%
\right) =\left( 
\begin{array}{cc}
\func{Re}\left( u+v\right) & -\func{Im}\left( u+v\right) \\ 
\func{Im}\left( u-v\right) & \func{Re}\left( u-v\right)%
\end{array}%
\right)
\end{eqnarray}%
showing that they are are a representation of $O\left( 2,\mathbb{R}\right) .$

$\lceil $ Recall%
\begin{equation}
z-\overline{z}=x+iy-x+iy=2iy
\end{equation}%
then%
\begin{gather}
\frac{1}{2}\left( 
\begin{array}{cc}
1 & 1 \\ 
-i & i%
\end{array}%
\right) \left( 
\begin{array}{cc}
u & \overline{v} \\ 
v & \overline{u}%
\end{array}%
\right) \left( 
\begin{array}{cc}
1 & i \\ 
1 & -i%
\end{array}%
\right) =\frac{1}{2}\left( 
\begin{array}{cc}
u+v & \overline{u}+\overline{v} \\ 
-iu+iv & i\overline{u}-i\overline{v}%
\end{array}%
\right) \left( 
\begin{array}{cc}
1 & i \\ 
1 & -i%
\end{array}%
\right) \\
=\frac{1}{2}\left( 
\begin{array}{cc}
u+v+\overline{u}+\overline{v} & iu+iv-i\overline{u}-i\overline{v} \\ 
-iu+iv+i\overline{u}-i\overline{v} & u-v+\overline{u}-\overline{v}%
\end{array}%
\right) =\frac{1}{2}\left( 
\begin{array}{cc}
u+\overline{u}+v+\overline{v} & i\left( u-\overline{u}\right) +i\left( v-%
\overline{v}\right) \\ 
-i\left( u-\overline{u}\right) +i\left( v-\overline{v}\right) & u+\overline{u%
}-\left( v+\overline{v}\right)%
\end{array}%
\right) \\
=\left( 
\begin{array}{cc}
\func{Re}\left( u+v\right) & -\func{Im}\left( u+v\right) \\ 
\func{Im}\left( u-v\right) & \func{Re}\left( u-v\right)%
\end{array}%
\right) .
\end{gather}%
The orthogonal group condition follows from the unitary group condition%
\begin{gather}
\left( 
\begin{array}{cc}
\func{Re}\left( u+v\right) & -\func{Im}\left( u+v\right) \\ 
\func{Im}\left( u-v\right) & \func{Re}\left( u-v\right)%
\end{array}%
\right) \left( 
\begin{array}{cc}
\func{Re}\left( u+v\right) & -\func{Im}\left( u+v\right) \\ 
\func{Im}\left( u-v\right) & \func{Re}\left( u-v\right)%
\end{array}%
\right) ^{t}= \\
K\left( 
\begin{array}{cc}
u & \overline{v} \\ 
v & \overline{u}%
\end{array}%
\right) K^{\dagger }K\left( 
\begin{array}{cc}
u & \overline{v} \\ 
v & \overline{u}%
\end{array}%
\right) ^{\dagger }K^{\dagger }=K\left( 
\begin{array}{cc}
u & \overline{v} \\ 
v & \overline{u}%
\end{array}%
\right) K^{\dagger }K\left( 
\begin{array}{cc}
\overline{u} & \overline{v} \\ 
v & u%
\end{array}%
\right) K^{\dagger } \\
=K\left( 
\begin{array}{cc}
u & \overline{v} \\ 
v & \overline{u}%
\end{array}%
\right) \left( 
\begin{array}{cc}
\overline{u} & \overline{v} \\ 
v & u%
\end{array}%
\right) K^{\dagger }=KK^{\dagger }=\mathbb{I}_{2}.
\end{gather}%
$\rfloor $

This real matrix can be expressed as 
\begin{equation}
\left( 
\begin{array}{cc}
\func{Re}\left( u+v\right) & -\func{Im}\left( u+v\right) \\ 
\func{Im}\left( u-v\right) & \func{Re}\left( u-v\right)%
\end{array}%
\right) =\left( 
\begin{array}{cc}
r_{s}\cos \theta _{s} & -r_{s}\sin \theta _{s} \\ 
r_{d}\sin \theta _{d} & r_{d}\cos \theta _{d}%
\end{array}%
\right) ,
\end{equation}%
where 
\begin{eqnarray}
r_{s} &=&\sqrt{\left( x_{1}+x_{2}\right) ^{2}+\left( y_{1}+y_{2}\right) ^{2}}
\\
\theta _{s} &=&\tan ^{-1}\left( \frac{\left( y_{1}+y_{2}\right) }{\left(
x_{1}+x_{2}\right) }\right) .
\end{eqnarray}%
and 
\begin{eqnarray}
r_{d} &=&\sqrt{\left( x_{1}-x_{2}\right) ^{2}+\left( y_{1}-y_{2}\right) ^{2}}
\\
\theta _{d} &=&\tan ^{-1}\left( \frac{\left( y_{1}-y_{2}\right) }{\left(
x_{1}-x_{2}\right) }\right) .
\end{eqnarray}%
The determinant of $U\left( u,v\right) $ is given by 
\begin{eqnarray}
\det \left\{ U\left( u,v\right) \right\} &=&\det \left\{ \left( 
\begin{array}{cc}
u & \overline{v} \\ 
v & \overline{u}%
\end{array}%
\right) \right\} =\det \left\{ \left( 
\begin{array}{cc}
r_{s}\cos \theta _{s} & -r_{s}\sin \theta _{s} \\ 
r_{d}\sin \theta _{d} & r_{d}\cos \theta _{d}%
\end{array}%
\right) \right\} \\
&=&\left\vert u\right\vert ^{2}-\left\vert v\right\vert
^{2}=r_{s}r_{d}\left( \cos \theta _{s}\cos \theta _{d}+\sin \theta _{s}\sin
\theta _{d}\right) .
\end{eqnarray}%
The orthogonal group condition gives%
\begin{gather}
\left( 
\begin{array}{cc}
r_{s}\cos \theta _{s} & -r_{s}\sin \theta _{s} \\ 
r_{d}\sin \theta _{d} & r_{d}\cos \theta _{d}%
\end{array}%
\right) \left( 
\begin{array}{cc}
r_{s}\cos \theta _{s} & r_{d}\sin \theta _{d} \\ 
-r_{s}\sin \theta _{s} & r_{d}\cos \theta _{d}%
\end{array}%
\right) = \\
\left( 
\begin{array}{cc}
r_{s}^{2}\left( \cos ^{2}\theta _{s}+\sin ^{2}\theta _{s}\right) & 
r_{s}r_{d}\left( \cos \theta _{s}\sin \theta _{d}-\sin \theta _{s}\cos
\theta _{d}\right) \\ 
r_{s}r_{d}\left( \cos \theta _{s}\sin \theta _{d}-\sin \theta _{s}\cos
\theta _{d}\right) & r_{d}^{2}\left( \cos ^{2}\theta _{d}+\sin ^{2}\theta
_{d}\right)%
\end{array}%
\right) \\
=\left( 
\begin{array}{cc}
1 & 0 \\ 
0 & 1%
\end{array}%
\right) ,
\end{gather}%
which gives 
\begin{eqnarray}
r_{s}^{2} &=&r_{d}^{2}=1 \\
r_{s} &=&\pm 1 \\
r_{d} &=&\pm 1
\end{eqnarray}%
and%
\begin{eqnarray}
\cos \theta _{s}\sin \theta _{d} &=&\sin \theta _{s}\cos \theta _{d} \\
\tan \theta _{d} &=&\tan \theta _{s}.
\end{eqnarray}%
As%
\begin{eqnarray}
\theta _{s} &=&\tan ^{-1}\left( \frac{\left( y_{1}+y_{2}\right) }{\left(
x_{1}+x_{2}\right) }\right) \\
\theta _{d} &=&\tan ^{-1}\left( \frac{\left( y_{1}-y_{2}\right) }{\left(
x_{1}-x_{2}\right) }\right)
\end{eqnarray}%
one obtains 
\begin{eqnarray}
\tan \left\{ \tan ^{-1}\left( \frac{\left( y_{1}+y_{2}\right) }{\left(
x_{1}+x_{2}\right) }\right) \right\} &=&\frac{\left( y_{1}+y_{2}\right) }{%
\left( x_{1}+x_{2}\right) } \\
\tan \left\{ \tan ^{-1}\left( \frac{\left( y_{1}-y_{2}\right) }{\left(
x_{1}-x_{2}\right) }\right) \right\} &=&\frac{\left( y_{1}-y_{2}\right) }{%
\left( x_{1}-x_{2}\right) }
\end{eqnarray}%
giving 
\begin{eqnarray}
&&0=\frac{\left( y_{1}+y_{2}\right) }{\left( x_{1}+x_{2}\right) }-\frac{%
\left( y_{1}-y_{2}\right) }{\left( x_{1}-x_{2}\right) }=\frac{\left(
y_{1}+y_{2}\right) \left( x_{1}-x_{2}\right) -\left( y_{1}-y_{2}\right)
\left( x_{1}+x_{2}\right) }{\left( x_{1}+x_{2}\right) \left(
x_{1}-x_{2}\right) } \\
&&\left. =\right. \frac{%
y_{1}x_{1}-y_{1}x_{2}+y_{2}x_{1}-y_{2}x_{2}-y_{1}x_{1}-y_{1}x_{2}+y_{2}x_{1}+y_{2}x_{2}%
}{x_{1}^{2}-x_{2}^{2}}=\frac{-2y_{1}x_{2}+2y_{2}x_{1}}{x_{1}^{2}-x_{2}^{2}}=%
\frac{2\left( y_{2}x_{1}-y_{1}x_{2}\right) }{x_{1}^{2}-x_{2}^{2}}
\end{eqnarray}%
then%
\begin{eqnarray}
y_{2}x_{1}-y_{1}x_{2} &=&0\Rightarrow \frac{x_{1}}{y_{1}}=\frac{r_{1}\cos
\theta _{1}}{r_{1}\sin \theta _{1}}=\frac{x_{2}}{y_{2}}=\frac{r_{2}\cos
\theta _{2}}{r_{2}\sin \theta _{2}} \\
&\Rightarrow &\tan \theta _{1}=\tan \theta _{2} \\
&\Rightarrow &\theta _{1}=\theta _{2}+n\pi ;\,n\in \mathbb{Z}.
\end{eqnarray}%
Hence $\left( 
\begin{array}{cc}
r_{s}\cos \theta _{s} & -r_{s}\sin \theta _{s} \\ 
r_{d}\sin \theta _{d} & r_{d}\cos \theta _{d}%
\end{array}%
\right) \in O\left( 2,\mathbb{R}\right) $ if $r_{s}^{2}=r_{d}^{2}=1$ and $%
\theta _{1}=\theta _{2}+n\pi $. A special case of this when $r_{s}=r_{d}=1$
and $\theta _{s}=\theta _{d}$ i.e. $v=0$ then 
\begin{equation}
U\left( u,0\right) =\left( 
\begin{array}{cc}
u & 0 \\ 
0 & \overline{u}%
\end{array}%
\right)
\end{equation}%
and $U\left( u,0\right) $ is an orthogonal rotation i.e. 
\begin{equation}
\widehat{K}\left( U\left( u,0\right) \right) =\left( 
\begin{array}{cc}
\cos \theta & -\sin \theta \\ 
\sin \theta & \cos \theta%
\end{array}%
\right) \in SO\left( 2,\mathbb{R}\right) \subset \rho \left( \mathbb{C}%
\right) .
\end{equation}%
The other component, the reflections, with $\det \left\{ O\right\} =-1$ is
given by matrices of the form 
\begin{equation}
O=\left( 
\begin{array}{cc}
\cos \theta & -\sin \theta \\ 
-\sin \theta & -\cos \theta%
\end{array}%
\right) ,
\end{equation}%
which correspond to $r_{s}=r_{d}=1$ and 
\begin{equation}
\theta _{d}=\theta _{s}-\pi .
\end{equation}%
$\lceil $ This can seen by observing%
\begin{eqnarray}
\sin \left( \theta -\pi \right) &=&\sin \theta \cos \pi -\cos \theta \sin
\pi =-\sin \theta \\
\cos \left( \theta -\pi \right) &=&\cos \theta \cos \pi +\sin \theta \sin
\pi =-\cos \theta .
\end{eqnarray}%
$\rfloor $

$\lceil $%
\begin{equation}
\widehat{K}:\mathcal{R}_{M\left( 2,\mathbb{C}\right) }\left( O\left( 2,%
\mathbb{R}\right) \right) \subset M\left( 2,\mathbb{C}\right) \rightarrow 
\mathcal{R}_{M\left( 2,\mathbb{R}\right) }\left( O\left( 2,\mathbb{R}\right)
\right) \subset M\left( 2,\mathbb{R}\right)
\end{equation}%
$\rfloor $

\subsubsection{Complex Numbers from a Clifford Algebra}

For $r=1$ the \emph{normalized} basis, (wrt the $\limfunc{tr}$ inner
product), of the \emph{real} Clifford algebra $\mathcal{F}=Cl\left( 2,0,%
\mathbb{R}\right) $ is 
\begin{equation}
I_{\mathcal{F}},\widetilde{x}_{1},\widetilde{x}_{2},\widetilde{x}_{1}%
\widetilde{x}_{2},
\end{equation}%
its dimension is%
\begin{equation}
\mathbb{R}\text{-}\dim \left\{ \mathcal{F}\right\} =2^{2}=4
\end{equation}%
and 
\begin{equation}
\mathcal{F}\cong M\left( 2,\mathbb{R}\right) .
\end{equation}

\begin{remark}
$M\left( 2,\mathbb{R}\right) $ has an inner product defined by 
\begin{equation}
\left\langle \left. \mathbb{A}\right\vert \mathbb{B}\right\rangle _{M\left(
2,\mathbb{R}\right) }=\frac{1}{2}\func{Tr}_{M\left( 2,\mathbb{R}\right)
}\left\{ \mathbb{A}^{\dagger }\mathbb{B}\right\} =\limfunc{tr}%
\nolimits_{M\left( 2,\mathbb{R}\right) }\left\{ \mathbb{A}^{\dagger }\mathbb{%
B}\right\}
\end{equation}%
and $\mathcal{F}$ an inner product defined by 
\begin{equation}
\left\langle \left. A\right\vert B\right\rangle _{\mathcal{F}}=\limfunc{tr}%
\nolimits_{\mathcal{F}}\left\{ A^{\dagger }B\right\} =\left\langle
\left\langle A^{\dagger }B\right\rangle \right\rangle _{0}.
\end{equation}
\end{remark}

One can express $\mathcal{F}$ in terms of an exterior algebra giving%
\begin{equation}
\mathcal{F}=\tbigwedge \left( \mathcal{F}_{1}\oplus \mathcal{F}_{1}\right)
\cong \tbigwedge \left( \mathbb{R\oplus R}\right) =\tbigwedge \left( \mathbb{%
R}\right) \otimes \tbigwedge \left( \mathbb{R}\right) \cong \mathcal{B}%
\left( \mathbb{R}^{2}\right)
\end{equation}%
$\lceil \mathcal{F}$ is isomorphic to $\mathcal{B}\left( \mathbb{R}%
^{2}\right) $ and not equal to it as $\mathcal{F}$ can be constituted by any
mathematical objects that can be combined in accordance with a $c^{\ast }$%
-algebraic structure, while $\mathcal{B}\left( \mathbb{R}^{2}\right) $ is
constituted by $2\times 2$ matrices of real numbers.$\rfloor $%
\begin{align}
\mathcal{F}& \cong \tbigwedge \mathbb{R}^{2}=\tbigwedge \left( \mathbb{%
R\oplus R}\right) =\tbigwedge\nolimits^{0}\left( \mathbb{R}^{2}\right)
\oplus \tbigwedge\nolimits^{1}\left( \mathbb{R}^{2}\right) \oplus
\tbigwedge\nolimits^{2}\left( \mathbb{R}^{2}\right)
=\tbigwedge\nolimits^{0}\left( \mathbb{R\oplus R}\right) \oplus
\tbigwedge\nolimits^{1}\left( \mathbb{R\oplus R}\right) \oplus
\tbigwedge\nolimits^{2}\left( \mathbb{R\oplus R}\right) \\
& =\mathbb{R}I_{2}\oplus \mathbb{R}^{2}\oplus \tbigwedge\nolimits^{2}\left( 
\mathbb{R}^{2}\right) \cong \left( \left\{ \left. xI_{2}\right\vert x\in 
\mathbb{R}\right\} \right) \oplus \left( \mathbb{R}-\limfunc{linspan}\left\{ 
\widetilde{x}_{1},\widetilde{x}_{2}\right\} \right) \oplus \left( \mathbb{R}-%
\limfunc{linspan}\left\{ \widetilde{x}_{1}\widetilde{x}_{2}\right\} \right)
\\
& =\mathcal{F}_{0}\oplus \mathcal{F}_{1}\oplus \mathcal{F}_{2}.
\end{align}%
From the CR's one obtains 
\begin{equation}
\left( \widetilde{x}_{1}\widetilde{x}_{2}\right) ^{2}=\widetilde{x}_{1}%
\widetilde{x}_{2}\widetilde{x}_{1}\widetilde{x}_{2}=-\widetilde{x}_{1}^{2}%
\widetilde{x}_{2}^{2}=-I_{2},
\end{equation}%
showing that one can define a complex unit 
\begin{equation}
i_{x}\equiv \widetilde{x}_{1}\widetilde{x}_{2}.
\end{equation}%
Any element of $\mathcal{F}$ can be expressed as 
\begin{equation}
A=\alpha _{0}I_{2}+\alpha _{11}\widetilde{x}_{1}+\alpha _{12}\widetilde{x}%
_{2}+\alpha _{3}i=\alpha _{0}I_{2}+\alpha _{3}i_{\widetilde{x}}+\alpha _{11}%
\widetilde{x}_{1}+\alpha _{12}\widetilde{x}_{2},
\end{equation}%
giving 
\begin{equation}
\left\Vert A\right\Vert _{\mathcal{F}}^{2}=\limfunc{tr}\nolimits_{\mathcal{F}%
}\left\{ A^{\dagger }A\right\} =\left\langle \left\langle A^{\dagger
}A\right\rangle \right\rangle _{0}=\left( \alpha _{0}^{2}+\alpha
_{11}^{2}+\alpha _{12}^{2}+\alpha _{3}^{2}\right) .
\end{equation}%
Consider the representation $\mathcal{\Pi }_{x}\left( \mathcal{F}\right) $
of the operators $\left\{ I_{2},x_{1},x_{2},x_{1}x_{2}\right\} $ in $M\left(
2,\mathbb{R}\right) $ given by 
\begin{align}
\mathcal{\Pi }_{x}\left( I_{2}\right) & =\left( 
\begin{array}{cc}
1 & 0 \\ 
0 & 1%
\end{array}%
\right) =\left( 
\begin{array}{cc}
0 & 0 \\ 
0 & 1%
\end{array}%
\right) +\left( 
\begin{array}{cc}
1 & 0 \\ 
0 & 0%
\end{array}%
\right) \\
\mathcal{\Pi }_{x}\left( \widetilde{x}_{1}\right) & =\left( 
\begin{array}{cc}
0 & 1 \\ 
1 & 0%
\end{array}%
\right) =\left( 
\begin{array}{cc}
0 & 1 \\ 
0 & 0%
\end{array}%
\right) +\left( 
\begin{array}{cc}
0 & 0 \\ 
1 & 0%
\end{array}%
\right) \\
\mathcal{\Pi }_{x}\left( \widetilde{x}_{2}\right) & =\left( 
\begin{array}{cc}
1 & 0 \\ 
0 & -1%
\end{array}%
\right) =\left( 
\begin{array}{cc}
1 & 0 \\ 
0 & 0%
\end{array}%
\right) -\left( 
\begin{array}{cc}
0 & 0 \\ 
0 & 1%
\end{array}%
\right) \\
\mathcal{\Pi }_{x}\left( \widetilde{x}_{1}\widetilde{x}_{2}\right) & =\left( 
\begin{array}{cc}
0 & -1 \\ 
1 & 0%
\end{array}%
\right) =\left( 
\begin{array}{cc}
0 & 0 \\ 
1 & 0%
\end{array}%
\right) -\left( 
\begin{array}{cc}
0 & 1 \\ 
0 & 0%
\end{array}%
\right) .
\end{align}%
The squares of the basis elements of this representation are $\pm \mathbb{I}%
_{2},$ which is confirmed by 
\begin{align}
\mathcal{\Pi }_{x}\left( \widetilde{x}_{1}\right) ^{2}& =\mathcal{\Pi }%
_{x}\left( \widetilde{x}_{1}\right) \mathcal{\Pi }_{x}\left( \widetilde{x}%
_{1}\right) =\left( 
\begin{array}{cc}
0 & 1 \\ 
1 & 0%
\end{array}%
\right) \left( 
\begin{array}{cc}
0 & 1 \\ 
1 & 0%
\end{array}%
\right) =\left( 
\begin{array}{cc}
1 & 0 \\ 
0 & 1%
\end{array}%
\right) =\mathbb{I}_{2} \\
\mathcal{\Pi }_{x}\left( \widetilde{x}_{2}\right) ^{2}& =\mathcal{\Pi }%
_{x}\left( \widetilde{x}_{2}\right) \mathcal{\Pi }_{x}\left( \widetilde{x}%
_{2}\right) =\left( 
\begin{array}{cc}
-1 & 0 \\ 
0 & 1%
\end{array}%
\right) \left( 
\begin{array}{cc}
-1 & 0 \\ 
0 & 1%
\end{array}%
\right) =\left( 
\begin{array}{cc}
1 & 0 \\ 
0 & 1%
\end{array}%
\right) =\mathbb{I}_{2} \\
\mathcal{\Pi }_{x}\left( \widetilde{x}_{1}\widetilde{x}_{2}\right) ^{2}& =%
\mathcal{\Pi }_{x}\left( \widetilde{x}_{1}\right) \mathcal{\Pi }_{x}\left( 
\widetilde{x}_{1}\right) =\left( 
\begin{array}{cc}
0 & -1 \\ 
1 & 0%
\end{array}%
\right) \left( 
\begin{array}{cc}
0 & -1 \\ 
1 & 0%
\end{array}%
\right) =\left( 
\begin{array}{cc}
-1 & 0 \\ 
0 & -1%
\end{array}%
\right) =-\mathbb{I}_{2}.
\end{align}%
The state space $\mathcal{S}\left( \mathcal{\Pi }\left( \mathcal{F}\right)
\right) =\mathcal{S}\left( M\left( 2,\mathbb{R}\right) \right) $ has \emph{%
one} $O\left( 2,\mathbb{R}\right) $ \emph{orbit} of extreme states generated
by the density matrices $\mathcal{\Pi }\left( D_{1}\right) $ and $\mathcal{%
\Pi }\left( D_{2}\right) $ tha are given by 
\begin{eqnarray}
f_{D_{1}}\left( A\right) &=&f_{D_{1}}\left( \mathcal{\Pi }_{x}\left(
A\right) \right) =f_{D_{1}}\left( \mathbb{A}\right) =\func{Tr}\left\{ 
\mathcal{\Pi }_{x}\left( A\right) \mathcal{\Pi }_{x}\left( D_{1}\right)
\right\} =\func{Tr}\left\{ \mathcal{\Pi }_{x}\left( A\right) \left( 
\begin{array}{cc}
0 & 0 \\ 
0 & 1%
\end{array}%
\right) \right\} \\
f_{D_{2}}\left( A\right) &=&f_{D_{1}}\left( \mathcal{\Pi }_{x}\left(
A\right) \right) =f_{D_{2}}\left( \mathbb{A}\right) =\func{Tr}\left\{ 
\mathcal{\Pi }_{x}\left( A\right) \mathcal{\Pi }_{x}\left( D_{2}\right)
\right\} =\func{Tr}\left\{ \mathcal{\Pi }_{x}\left( A\right) \left( 
\begin{array}{cc}
1 & 0 \\ 
0 & 0%
\end{array}%
\right) \right\} ,
\end{eqnarray}%
the states $f_{D_{1}}$ are on the same orbit. The state space, $\mathcal{S}%
\left( M\left( 2,\mathbb{R}\right) \right) ,$ of $M\left( 2,\mathbb{R}%
\right) $ is generated by 
\begin{equation}
\mathcal{S}\left( M\left( 2,\mathbb{R}\right) \right) =\limfunc{Conv}\left\{ 
\widehat{O}\left( f_{D_{1}}\right) \right\} \cong \limfunc{Conv}\left\{ 
\widehat{O}\left( \mathcal{\Pi }_{x}\left( D_{1}\right) \right) \right\}
;O_{1}\in O\left( 2,\mathbb{R}\right) .
\end{equation}

$\mathcal{S}\left( M\left( 2,\mathbb{R}\right) \right) $ also contains the
faithful mixed state $\mathcal{\Pi }_{x}\left( D_{\limfunc{tr}}\right) =%
\frac{1}{2}\mathbb{I}_{2}$ that lies on a different orbit to the pure
states. \emph{Any mixed state lies on a distinct orbit characterized by
their eigenvalues.}

Any element in $\mathcal{\Pi }\left( \mathcal{F}\right) $ can be expressed as

\begin{eqnarray}
\mathcal{\Pi }_{x}\left( A\right) &=&\alpha _{0}\mathcal{\Pi }_{x}\left(
I_{2}\right) +\alpha _{11}\mathcal{\Pi }_{x}\left( \widetilde{x}_{1}\right)
+\alpha _{12}\mathcal{\Pi }_{x}\left( \widetilde{x}_{2}\right) +\alpha _{2}%
\mathcal{\Pi }_{x}\left( \widetilde{x}_{1}\widetilde{x}_{2}\right) \\
&=&\alpha _{0}\left( 
\begin{array}{cc}
1 & 0 \\ 
0 & 1%
\end{array}%
\right) +\alpha _{11}\left( 
\begin{array}{cc}
0 & 1 \\ 
1 & 0%
\end{array}%
\right) +\alpha _{12}\left( 
\begin{array}{cc}
-1 & 0 \\ 
0 & 1%
\end{array}%
\right) +\alpha _{2}\left( 
\begin{array}{cc}
0 & 1 \\ 
-1 & 0%
\end{array}%
\right) \\
&=&\left( 
\begin{array}{cc}
\alpha _{0}-\alpha _{12} & \alpha _{11}+\alpha _{2} \\ 
\alpha _{11}-\alpha _{2} & \alpha _{0}+\alpha _{12}%
\end{array}%
\right)
\end{eqnarray}%
and its positive part $\mathcal{\Pi }_{x}\left( A^{\dagger }A\right) =%
\mathcal{\Pi }_{x}\left( A\right) ^{\dagger }\mathcal{\Pi }_{x}\left(
A\right) $ by%
\begin{eqnarray}
\mathcal{\Pi }_{x}\left( A^{\dagger }A\right) &=&\left( 
\begin{array}{cc}
\alpha _{0}-\alpha _{12} & \alpha _{11}-\alpha _{2} \\ 
\alpha _{11}+\alpha _{2} & \alpha _{0}+\alpha _{12}%
\end{array}%
\right) \left( 
\begin{array}{cc}
\alpha _{0}-\alpha _{12} & \alpha _{11}+\alpha _{2} \\ 
\alpha _{11}-\alpha _{2} & \alpha _{0}+\alpha _{12}%
\end{array}%
\right) \\
&=&\left( 
\begin{array}{cc}
\left\Vert A\right\Vert _{\mathcal{F}}^{2}-2\left( \alpha _{0}\alpha
_{12}+\alpha _{2}\alpha _{11}\right) & 2\left( \alpha _{0}\alpha
_{11}-\alpha _{2}\alpha _{12}\right) \\ 
2\left( \alpha _{0}\alpha _{11}-\alpha _{3}\alpha _{12}\right) & \left\Vert
A\right\Vert _{\mathcal{F}}^{2}+2\left( \alpha _{0}\alpha _{12}+\alpha
_{2}\alpha _{11}\right)%
\end{array}%
\right) ,
\end{eqnarray}%
where 
\begin{equation}
\left\Vert A\right\Vert _{\mathcal{F}}^{2}=\left( \alpha _{0}^{2}+\alpha
_{11}^{2}+\alpha _{12}^{2}+\alpha _{2}^{2}\right) .
\end{equation}%
$\lceil $One can confirm that this matrix is positive by%
\begin{eqnarray}
&&\left( 
\begin{array}{cc}
a & b%
\end{array}%
\right) \left( 
\begin{array}{cc}
A_{11}^{\dagger } & A_{21}^{\dagger } \\ 
A_{12}^{\dagger } & A_{22}^{\dagger }%
\end{array}%
\right) \left( 
\begin{array}{cc}
A_{11} & A_{12} \\ 
A_{12} & A_{22}%
\end{array}%
\right) \left( 
\begin{array}{c}
a \\ 
b%
\end{array}%
\right) =\left( \left( 
\begin{array}{cc}
A_{11} & A_{12} \\ 
A_{12} & A_{22}%
\end{array}%
\right) \left( 
\begin{array}{c}
a \\ 
b%
\end{array}%
\right) \right) ^{\dagger }\left( 
\begin{array}{cc}
A_{11} & A_{12} \\ 
A_{12} & A_{22}%
\end{array}%
\right) \left( 
\begin{array}{c}
a \\ 
b%
\end{array}%
\right) \\
&=&\left( 
\begin{array}{c}
A_{11}a+A_{12}b \\ 
A_{12}a+A_{22}b%
\end{array}%
\right) ^{\dagger }\left( 
\begin{array}{c}
A_{11}a+A_{12}b \\ 
A_{12}a+A_{22}b%
\end{array}%
\right) =\left( 
\begin{array}{cc}
a^{\dagger }A_{11}^{\dagger }+b^{\dagger }A_{12}^{\dagger } & a^{\dagger
}A_{12}^{\dagger }+b^{\dagger }A_{22}^{\dagger }%
\end{array}%
\right) \left( 
\begin{array}{c}
A_{11}a+A_{12}b \\ 
A_{12}a+A_{22}b%
\end{array}%
\right) \\
&=&\left( A_{11}a+A_{12}b\right) ^{\dagger }\left( A_{11}a+A_{12}b\right)
+\left( A_{12}a+A_{22}b\right) ^{\dagger }\left( A_{12}a+A_{22}b\right)
=\left( A_{11}a+A_{12}b\right) ^{2}+\left( A_{12}a+A_{22}b\right) ,
\end{eqnarray}%
which can be compared to 
\begin{eqnarray}
&&\left( 
\begin{array}{cc}
a & b%
\end{array}%
\right) \left( 
\begin{array}{cc}
\left( A^{\dagger }A\right) _{11} & \left( A^{\dagger }A\right) _{12} \\ 
\left( A^{\dagger }A\right) _{12} & \left( A^{\dagger }A\right) _{22}%
\end{array}%
\right) \left( 
\begin{array}{c}
a \\ 
b%
\end{array}%
\right) =\left( 
\begin{array}{cc}
a\left( A^{\dagger }A\right) _{11}+b\left( A^{\dagger }A\right) _{12} & 
a\left( A^{\dagger }A\right) _{12}+b\left( A^{\dagger }A\right) _{22}%
\end{array}%
\right) \left( 
\begin{array}{c}
a \\ 
b%
\end{array}%
\right) \\
&=&a\left( A^{\dagger }A\right) _{11}a+b\left( A^{\dagger }A\right)
_{12}a+a\left( A^{\dagger }A\right) _{12}b+b\left( A^{\dagger }A\right)
_{22}b=\left( A^{\dagger }A\right) _{11}a^{2}+2\left( A^{\dagger }A\right)
_{12}ab+\left( A^{\dagger }A\right) _{22}b^{2}
\end{eqnarray}%
giving%
\begin{eqnarray}
&&\left( 
\begin{array}{cc}
a & b%
\end{array}%
\right) \left( 
\begin{array}{cc}
\left( \alpha _{0}^{2}+\alpha _{3}^{2}+\alpha _{11}^{2}+\alpha
_{12}^{2}\right) -2\left( \alpha _{0}\alpha _{12}+\alpha _{3}\alpha
_{11}\right) & 2\left( \alpha _{0}\alpha _{11}-\alpha _{3}\alpha _{12}\right)
\\ 
2\left( \alpha _{0}\alpha _{11}-\alpha _{3}\alpha _{12}\right) & \left(
\alpha _{0}^{2}+\alpha _{3}^{2}+\alpha _{11}^{2}+\alpha _{12}^{2}\right)
+2\left( \alpha _{0}\alpha _{12}+\alpha _{3}\alpha _{11}\right)%
\end{array}%
\right) \left( 
\begin{array}{c}
a \\ 
b%
\end{array}%
\right) = \\
&&\left\{ \left( \alpha _{0}^{2}+\alpha _{3}^{2}+\alpha _{11}^{2}+\alpha
_{12}^{2}\right) -2\left( \alpha _{0}\alpha _{12}+\alpha _{3}\alpha
_{11}\right) \right\} a^{2}+4\left( \alpha _{0}\alpha _{11}-\alpha
_{3}\alpha _{12}\right) ab+\left\{ \left( \alpha _{0}^{2}+\alpha
_{3}^{2}+\alpha _{11}^{2}+\alpha _{12}^{2}\right) +2\left( \alpha _{0}\alpha
_{12}+\alpha _{3}\alpha _{11}\right) \right\} b^{2} \\
&&\left\{ \left( \alpha _{0}-\alpha _{12}\right) ^{2}+\left( \alpha
_{11}-\alpha _{3}\right) ^{2}\right\} a^{2}+4\left( \alpha _{0}\alpha
_{11}-\alpha _{3}\alpha _{12}\right) ab+\left\{ \left( \alpha _{11}+\alpha
_{3}\right) ^{2}+\left( \alpha _{0}+\alpha _{12}\right) ^{2}\right\} b^{2}.
\end{eqnarray}%
$\rfloor $

$\lceil $ Calculation of $A^{\dagger }A$ 
\begin{eqnarray}
&&\left\{ \alpha _{0}\left( 
\begin{array}{cc}
1 & 0 \\ 
0 & 1%
\end{array}%
\right) +\alpha _{11}\left( 
\begin{array}{cc}
0 & 1 \\ 
1 & 0%
\end{array}%
\right) +\alpha _{12}\left( 
\begin{array}{cc}
-1 & 0 \\ 
0 & 1%
\end{array}%
\right) +\alpha _{3}\left( 
\begin{array}{cc}
0 & -1 \\ 
1 & 0%
\end{array}%
\right) \right\} \\
&&\times \left\{ \alpha _{0}\left( 
\begin{array}{cc}
1 & 0 \\ 
0 & 1%
\end{array}%
\right) +\alpha _{11}\left( 
\begin{array}{cc}
0 & 1 \\ 
1 & 0%
\end{array}%
\right) +\alpha _{12}\left( 
\begin{array}{cc}
-1 & 0 \\ 
0 & 1%
\end{array}%
\right) +\alpha _{3}\left( 
\begin{array}{cc}
0 & 1 \\ 
-1 & 0%
\end{array}%
\right) \right\}
\end{eqnarray}%
\begin{eqnarray}
&=&\alpha _{0}\left( 
\begin{array}{cc}
1 & 0 \\ 
0 & 1%
\end{array}%
\right) \left\{ \alpha _{0}\left( 
\begin{array}{cc}
1 & 0 \\ 
0 & 1%
\end{array}%
\right) +\alpha _{11}\left( 
\begin{array}{cc}
0 & 1 \\ 
1 & 0%
\end{array}%
\right) +\alpha _{12}\left( 
\begin{array}{cc}
-1 & 0 \\ 
0 & 1%
\end{array}%
\right) +\alpha _{3}\left( 
\begin{array}{cc}
0 & 1 \\ 
-1 & 0%
\end{array}%
\right) \right\} \\
&&+\alpha _{11}\left( 
\begin{array}{cc}
0 & 1 \\ 
1 & 0%
\end{array}%
\right) \left\{ \alpha _{0}\left( 
\begin{array}{cc}
1 & 0 \\ 
0 & 1%
\end{array}%
\right) +\alpha _{11}\left( 
\begin{array}{cc}
0 & 1 \\ 
1 & 0%
\end{array}%
\right) +\alpha _{12}\left( 
\begin{array}{cc}
-1 & 0 \\ 
0 & 1%
\end{array}%
\right) +\alpha _{3}\left( 
\begin{array}{cc}
0 & 1 \\ 
-1 & 0%
\end{array}%
\right) \right\} \\
&&+\alpha _{12}\left( 
\begin{array}{cc}
-1 & 0 \\ 
0 & 1%
\end{array}%
\right) \left\{ \alpha _{0}\left( 
\begin{array}{cc}
1 & 0 \\ 
0 & 1%
\end{array}%
\right) +\alpha _{11}\left( 
\begin{array}{cc}
0 & 1 \\ 
1 & 0%
\end{array}%
\right) +\alpha _{12}\left( 
\begin{array}{cc}
-1 & 0 \\ 
0 & 1%
\end{array}%
\right) +\alpha _{3}\left( 
\begin{array}{cc}
0 & 1 \\ 
-1 & 0%
\end{array}%
\right) \right\} \\
&&+\alpha _{3}\left( 
\begin{array}{cc}
0 & -1 \\ 
1 & 0%
\end{array}%
\right) \left\{ \alpha _{0}\left( 
\begin{array}{cc}
1 & 0 \\ 
0 & 1%
\end{array}%
\right) +\alpha _{11}\left( 
\begin{array}{cc}
0 & 1 \\ 
1 & 0%
\end{array}%
\right) +\alpha _{12}\left( 
\begin{array}{cc}
-1 & 0 \\ 
0 & 1%
\end{array}%
\right) +\alpha _{3}\left( 
\begin{array}{cc}
0 & 1 \\ 
-1 & 0%
\end{array}%
\right) \right\}
\end{eqnarray}%
\begin{eqnarray}
&=&\left\{ \alpha _{0}\alpha _{0}\left( 
\begin{array}{cc}
1 & 0 \\ 
0 & 1%
\end{array}%
\right) +\alpha _{0}\alpha _{11}\left( 
\begin{array}{cc}
0 & 1 \\ 
1 & 0%
\end{array}%
\right) +\alpha _{0}\alpha _{12}\left( 
\begin{array}{cc}
-1 & 0 \\ 
0 & 1%
\end{array}%
\right) +\alpha _{0}\alpha _{3}\left( 
\begin{array}{cc}
0 & 1 \\ 
-1 & 0%
\end{array}%
\right) \right\} \\
&&+\left\{ 
\begin{array}{c}
\alpha _{0}\alpha _{11}\left( 
\begin{array}{cc}
0 & 1 \\ 
1 & 0%
\end{array}%
\right) \left( 
\begin{array}{cc}
1 & 0 \\ 
0 & 1%
\end{array}%
\right) +\alpha _{11}\alpha _{11}\left( 
\begin{array}{cc}
0 & 1 \\ 
1 & 0%
\end{array}%
\right) \left( 
\begin{array}{cc}
0 & 1 \\ 
1 & 0%
\end{array}%
\right) +\alpha _{12}\alpha _{11}\left( 
\begin{array}{cc}
0 & 1 \\ 
1 & 0%
\end{array}%
\right) \left( 
\begin{array}{cc}
-1 & 0 \\ 
0 & 1%
\end{array}%
\right) \\ 
+\alpha _{3}\alpha _{11}\left( 
\begin{array}{cc}
0 & 1 \\ 
1 & 0%
\end{array}%
\right) \left( 
\begin{array}{cc}
0 & 1 \\ 
-1 & 0%
\end{array}%
\right)%
\end{array}%
\right\} \\
&&+\left\{ 
\begin{array}{c}
\alpha _{0}\alpha _{12}\left( 
\begin{array}{cc}
-1 & 0 \\ 
0 & 1%
\end{array}%
\right) \left( 
\begin{array}{cc}
1 & 0 \\ 
0 & 1%
\end{array}%
\right) +\alpha _{11}\alpha _{12}\left( 
\begin{array}{cc}
-1 & 0 \\ 
0 & 1%
\end{array}%
\right) \left( 
\begin{array}{cc}
0 & 1 \\ 
1 & 0%
\end{array}%
\right) +\alpha _{12}\alpha _{12}\left( 
\begin{array}{cc}
-1 & 0 \\ 
0 & 1%
\end{array}%
\right) \left( 
\begin{array}{cc}
-1 & 0 \\ 
0 & 1%
\end{array}%
\right) \\ 
+\alpha _{3}\alpha _{12}\left( 
\begin{array}{cc}
-1 & 0 \\ 
0 & 1%
\end{array}%
\right) \left( 
\begin{array}{cc}
0 & 1 \\ 
-1 & 0%
\end{array}%
\right)%
\end{array}%
\right\} \\
&&+\left\{ 
\begin{array}{c}
\alpha _{0}\alpha _{3}\left( 
\begin{array}{cc}
0 & -1 \\ 
1 & 0%
\end{array}%
\right) \left( 
\begin{array}{cc}
1 & 0 \\ 
0 & 1%
\end{array}%
\right) +\alpha _{11}\alpha _{3}\left( 
\begin{array}{cc}
0 & -1 \\ 
1 & 0%
\end{array}%
\right) \left( 
\begin{array}{cc}
0 & 1 \\ 
1 & 0%
\end{array}%
\right) +\alpha _{12}\alpha _{3}\left( 
\begin{array}{cc}
0 & -1 \\ 
1 & 0%
\end{array}%
\right) \left( 
\begin{array}{cc}
-1 & 0 \\ 
0 & 1%
\end{array}%
\right) \\ 
+\alpha _{3}\alpha _{3}\left( 
\begin{array}{cc}
0 & -1 \\ 
1 & 0%
\end{array}%
\right) \left( 
\begin{array}{cc}
0 & 1 \\ 
-1 & 0%
\end{array}%
\right)%
\end{array}%
\right\}
\end{eqnarray}%
\begin{eqnarray}
&=&\left\{ \alpha _{0}\alpha _{0}\left( 
\begin{array}{cc}
1 & 0 \\ 
0 & 1%
\end{array}%
\right) +\alpha _{0}\alpha _{11}\left( 
\begin{array}{cc}
0 & 1 \\ 
1 & 0%
\end{array}%
\right) +\alpha _{0}\alpha _{12}\left( 
\begin{array}{cc}
-1 & 0 \\ 
0 & 1%
\end{array}%
\right) +\alpha _{0}\alpha _{3}\left( 
\begin{array}{cc}
0 & 1 \\ 
-1 & 0%
\end{array}%
\right) \right\} \\
&&+\left\{ \alpha _{0}\alpha _{11}\left( 
\begin{array}{cc}
0 & 1 \\ 
1 & 0%
\end{array}%
\right) +\alpha _{11}\alpha _{11}\left( 
\begin{array}{cc}
1 & 0 \\ 
0 & 1%
\end{array}%
\right) +\alpha _{12}\alpha _{11}\left( 
\begin{array}{cc}
0 & 1 \\ 
-1 & 0%
\end{array}%
\right) +\alpha _{3}\alpha _{11}\left( 
\begin{array}{cc}
-1 & 0 \\ 
0 & 1%
\end{array}%
\right) \right\} \\
&&+\left\{ \alpha _{0}\alpha _{12}\left( 
\begin{array}{cc}
-1 & 0 \\ 
0 & 1%
\end{array}%
\right) +\alpha _{11}\alpha _{12}\left( 
\begin{array}{cc}
0 & -1 \\ 
1 & 0%
\end{array}%
\right) +\alpha _{12}\alpha _{12}\left( 
\begin{array}{cc}
1 & 0 \\ 
0 & 1%
\end{array}%
\right) +\alpha _{3}\alpha _{12}\left( 
\begin{array}{cc}
0 & -1 \\ 
-1 & 0%
\end{array}%
\right) \right\} \\
&&+\left\{ \alpha _{0}\alpha _{3}\left( 
\begin{array}{cc}
0 & -1 \\ 
1 & 0%
\end{array}%
\right) +\alpha _{11}\alpha _{3}\left( 
\begin{array}{cc}
-1 & 0 \\ 
0 & 1%
\end{array}%
\right) +\alpha _{12}\alpha _{3}\left( 
\begin{array}{cc}
0 & -1 \\ 
-1 & 0%
\end{array}%
\right) +\alpha _{3}\alpha _{3}\left( 
\begin{array}{cc}
1 & 0 \\ 
0 & 1%
\end{array}%
\right) \right\}
\end{eqnarray}%
\begin{equation}
=\left( \alpha _{0}^{2}+\alpha _{3}^{2}+\alpha _{11}^{2}+\alpha
_{12}^{2}\right) \left( 
\begin{array}{cc}
1 & 0 \\ 
0 & 1%
\end{array}%
\right) +2\left( \alpha _{0}\alpha _{11}-\alpha _{3}\alpha _{12}\right)
\left( 
\begin{array}{cc}
0 & 1 \\ 
1 & 0%
\end{array}%
\right) +2\left( \alpha _{0}\alpha _{12}+\alpha _{3}\alpha _{11}\right)
\left( 
\begin{array}{cc}
-1 & 0 \\ 
0 & 1%
\end{array}%
\right)
\end{equation}%
\begin{eqnarray}
&&\left( 
\begin{array}{cc}
a & b%
\end{array}%
\right) \left\{ 
\begin{array}{c}
\left( \alpha _{0}^{2}+\alpha _{3}^{2}+\alpha _{11}^{2}+\alpha
_{12}^{2}\right) \left( 
\begin{array}{cc}
1 & 0 \\ 
0 & 1%
\end{array}%
\right) +2\left( \alpha _{0}\alpha _{11}-\alpha _{3}\alpha _{12}\right)
\left( 
\begin{array}{cc}
0 & 1 \\ 
1 & 0%
\end{array}%
\right) \\ 
+2\left( \alpha _{0}\alpha _{12}+\alpha _{3}\alpha _{11}\right) \left( 
\begin{array}{cc}
-1 & 0 \\ 
0 & 1%
\end{array}%
\right)%
\end{array}%
\right\} \left( 
\begin{array}{c}
a \\ 
b%
\end{array}%
\right) \\
&=&\left( \alpha _{0}^{2}+\alpha _{3}^{2}+\alpha _{11}^{2}+\alpha
_{12}^{2}\right) \left( a^{2}+b^{2}\right) +2\left( \alpha _{0}\alpha
_{11}-\alpha _{3}\alpha _{12}\right) ab+2\left( \alpha _{0}\alpha
_{12}+\alpha _{3}\alpha _{11}\right) \left( b^{2}-a^{2}\right)
\end{eqnarray}%
giving%
\begin{eqnarray}
\mathcal{\Pi }_{x}\left( A^{\dagger }A\right) &=&\left( 
\begin{array}{cc}
\alpha _{0}-\alpha _{12} & \alpha _{11}-\alpha _{3} \\ 
\alpha _{11}+\alpha _{3} & \alpha _{0}+\alpha _{12}%
\end{array}%
\right) \left( 
\begin{array}{cc}
\alpha _{0}-\alpha _{12} & \alpha _{11}+\alpha _{3} \\ 
\alpha _{11}-\alpha _{3} & \alpha _{0}+\alpha _{12}%
\end{array}%
\right) \\
&=&\left( 
\begin{array}{cc}
\left( \alpha _{0}-\alpha _{12}\right) \left( \alpha _{0}-\alpha
_{12}\right) +\left( \alpha _{11}-\alpha _{3}\right) \left( \alpha
_{11}-\alpha _{3}\right) & \left( \alpha _{0}-\alpha _{12}\right) \left(
\alpha _{11}+\alpha _{3}\right) +\left( \alpha _{0}-\alpha _{12}\right)
\left( \alpha _{0}+\alpha _{12}\right) \\ 
\left( \alpha _{0}-\alpha _{12}\right) \left( \alpha _{11}+\alpha
_{3}\right) +\left( \alpha _{0}-\alpha _{12}\right) \left( \alpha
_{0}+\alpha _{12}\right) & \left( \alpha _{11}+\alpha _{3}\right) \left(
\alpha _{11}+\alpha _{3}\right) \ +\left( \alpha _{0}+\alpha _{12}\right)
\left( \alpha _{0}+\alpha _{12}\right)%
\end{array}%
\right) \\
&=&\left( 
\begin{array}{cc}
\left( \alpha _{0}^{2}+\alpha _{3}^{2}+\alpha _{11}^{2}+\alpha
_{12}^{2}\right) -2\left( \alpha _{0}\alpha _{12}+\alpha _{3}\alpha
_{11}\right) & 2\left( \alpha _{0}\alpha _{11}-\alpha _{3}\alpha _{12}\right)
\\ 
2\left( \alpha _{0}\alpha _{11}-\alpha _{3}\alpha _{12}\right) & \left(
\alpha _{0}^{2}+\alpha _{3}^{2}+\alpha _{11}^{2}+\alpha _{12}^{2}\right)
+2\left( \alpha _{0}\alpha _{12}+\alpha _{3}\alpha _{11}\right)%
\end{array}%
\right) .
\end{eqnarray}%
$\rfloor $

\subparagraph{Density Operators in $\mathcal{F}\protect\cong M\left( 2,%
\mathbb{R}\right) $}

$\lceil $\emph{Self adjoint linear functionals} on $\mathcal{F}$ can be
given as 
\begin{equation}
f_{D}\left( A\right) =2\limfunc{tr}\nolimits_{\mathcal{F}}\left\{ AD\right\}
=\func{Tr}_{\mathcal{F}}\left\{ AD\right\} ,
\end{equation}%
where 
\begin{equation}
D\geq 0,\,\func{Tr}_{\mathcal{F}}\left\{ D\right\} =1
\end{equation}%
giving 
\begin{equation}
f_{D}\left( I_{2}\right) =2\limfunc{tr}\nolimits_{\mathcal{F}}\left\{
D\right\} =\func{Tr}_{\mathcal{F}}\left\{ AI_{2}\right\} =1.
\end{equation}%
Let 
\begin{equation}
D_{I_{2}}=\frac{1}{2}I_{2}
\end{equation}%
then 
\begin{equation}
f_{D_{I_{2}}}\left( I_{2}\right) =2\limfunc{tr}\left\{ I_{2}\frac{1}{2}%
I_{2}\right\} =\func{Tr}\left\{ \frac{1}{2}I_{2}\right\} =1.
\end{equation}%
$\rfloor $

The amplitudes and density operators can be expressed as 
\begin{eqnarray}
Q &=&q_{0}I_{2}+q_{11}\widetilde{x}_{1}+q_{12}\widetilde{x}_{2}+q_{2}%
\widetilde{x}_{1}\widetilde{x}_{2} \\
D &=&I_{2}+d_{11}\widetilde{x}_{1}+d_{12}\widetilde{x}_{2}+d_{3}\widetilde{x}%
_{1}\widetilde{x}_{2}
\end{eqnarray}%
giving%
\begin{eqnarray}
Q^{\dagger }Q &=&I_{2}+2\left( q_{0}q_{11}-q_{2}q_{12}\right) \widetilde{x}%
_{1}+2\left( q_{0}q_{12}+q_{2}q_{11}\right) \widetilde{x}_{2} \\
QQ^{\dagger } &=&I_{2}+2\left( q_{0}q_{11}+q_{2}q_{12}\right) \widetilde{x}%
_{1}+2\left( q_{0}q_{12}-q_{2}q_{11}\right) \widetilde{x}_{2},
\end{eqnarray}%
where 
\begin{equation}
\left\Vert Q\right\Vert _{\mathcal{F}%
}^{2}=q_{0}^{2}+q_{11}^{2}+q_{12}^{2}+q_{2}^{2}=1.
\end{equation}%
Hence using the $\Pi $ representation one obtains%
\begin{equation}
\Pi _{x}\left( Q\right) =\left( 
\begin{array}{cc}
q_{0}-q_{12} & q_{11}+q_{2} \\ 
q_{11}-q_{2} & q_{0}+q_{12}%
\end{array}%
\right)
\end{equation}%
and%
\begin{eqnarray}
\Pi _{x}\left( Q^{\dagger }Q\right) &=&\left( 
\begin{array}{cc}
1+2\left( q_{0}q_{12}+q_{2}q_{11}\right) & 2\left(
q_{0}q_{11}-q_{2}q_{12}\right) \\ 
2\left( q_{0}q_{11}-q_{2}q_{12}\right) & 1-2\left(
q_{0}q_{12}+q_{2}q_{11}\right)%
\end{array}%
\right) \\
\Pi _{x}\left( QQ^{\dagger }\right) &=&\left( 
\begin{array}{cc}
1+2\left( q_{0}q_{12}-q_{2}q_{11}\right) & 2\left(
q_{0}q_{11}-q_{2}q_{12}\right) \\ 
2\left( q_{0}q_{11}-q_{2}q_{12}\right) & 1-2\left(
q_{0}q_{12}-q_{2}q_{11}\right)%
\end{array}%
\right) .
\end{eqnarray}%
In particular consider

\begin{enumerate}
\item 
\begin{equation}
Q=\frac{1}{2}\left( \widetilde{x}_{1}+\widetilde{x}_{1}\widetilde{x}%
_{2}\right) ,
\end{equation}%
which gives 
\begin{gather}
D_{1}=Q^{\dagger }Q=\frac{1}{4}\left( \widetilde{x}_{1}-\widetilde{x}_{1}%
\widetilde{x}_{2}\right) \left( \widetilde{x}_{1}+\widetilde{x}_{1}%
\widetilde{x}_{2}\right) =\frac{1}{2}\left( I_{2}+\widetilde{x}_{2}\right) \\
D_{2}=QQ^{\dagger }=\frac{1}{4}\left( \widetilde{x}_{1}+\widetilde{x}_{1}%
\widetilde{x}_{2}\right) \left( \widetilde{x}_{1}-\widetilde{x}_{1}%
\widetilde{x}_{2}\right) =\frac{1}{2}\left( I_{2}-\widetilde{x}_{2}\right) \\
Q=\frac{1}{2}\left( \widetilde{x}_{1}+\widetilde{x}_{1}\widetilde{x}%
_{2}\right) =q_{0}I_{2}+q_{11}\widetilde{x}_{1}+q_{12}\widetilde{x}_{2}+q_{2}%
\widetilde{x}_{1}\widetilde{x}_{2} \\
q_{0}=0,q_{11}=\frac{1}{2},q_{12}=0,q_{2}=\frac{1}{2}
\end{gather}%
and 
\begin{eqnarray}
D_{1}^{2} &=&\frac{1}{2}\left( I_{2}+\widetilde{x}_{2}\right) \frac{1}{2}%
\left( I_{2}+\widetilde{x}_{2}\right) =\frac{1}{2}\left( I_{2}+\widetilde{x}%
_{2}\right) \\
D_{1}^{2} &=&\frac{1}{2}\left( I_{2}-\widetilde{x}_{2}\right) \frac{1}{2}%
\left( I_{2}-\widetilde{x}_{2}\right) =\frac{1}{2}\left( I_{2}-\widetilde{x}%
_{2}\right) .
\end{eqnarray}%
Using the $\Pi $ representation one obtains the representations of $\mathcal{%
\Pi }\left( Q\right) ,\mathcal{\Pi }\left( D_{1}\right) $ and $\mathcal{\Pi }%
\left( D_{2}\right) $ as 
\begin{eqnarray*}
\mathcal{\Pi }_{x}\left( Q\right) &=&\left( 
\begin{array}{cc}
q_{0}+q_{12} & q_{11}-q_{2} \\ 
q_{11}+q_{2} & q_{0}-q_{12}%
\end{array}%
\right) =\left( 
\begin{array}{cc}
0 & \frac{1}{2}-\frac{1}{2} \\ 
\frac{1}{2}+\frac{1}{2} & 0%
\end{array}%
\right) =\left( 
\begin{array}{cc}
0 & 0 \\ 
1 & 0%
\end{array}%
\right) ,\text{ letting }q_{2}=\frac{1}{2} \\
\mathcal{\Pi }_{x}\left( D_{1}\right) &=&\left( 
\begin{array}{cc}
0 & 1 \\ 
0 & 0%
\end{array}%
\right) \left( 
\begin{array}{cc}
0 & 0 \\ 
1 & 0%
\end{array}%
\right) =\left( 
\begin{array}{cc}
1 & 0 \\ 
0 & 0%
\end{array}%
\right) \\
\mathcal{\Pi }_{x}\left( D_{2}\right) &=&\left( 
\begin{array}{cc}
0 & 0 \\ 
1 & 0%
\end{array}%
\right) \left( 
\begin{array}{cc}
0 & 1 \\ 
0 & 0%
\end{array}%
\right) =\left( 
\begin{array}{cc}
0 & 0 \\ 
0 & 1%
\end{array}%
\right) .
\end{eqnarray*}%
In particular one can choose 
\begin{equation}
q_{2}=q_{11}=\frac{1}{2}
\end{equation}%
to obtain 
\begin{equation}
\mathcal{\Pi }_{x}\left( Q\right) =\left( 
\begin{array}{cc}
0 & 0 \\ 
1 & 0%
\end{array}%
\right) .\rfloor
\end{equation}%
In the case $\mathcal{F}\cong \Pi _{x}\left( \mathcal{F}\right) =M\left( 2,%
\mathbb{R}\right) $ one has $SO\left( 2,\mathbb{R}\right) $ orbits generated
by these density matrices given by 
\begin{eqnarray}
\Pi _{x}\left( D_{1}\left( \theta \right) \right) &=&\left( 
\begin{array}{cc}
\cos \theta & -\sin \theta \\ 
\sin \theta & \cos \theta%
\end{array}%
\right) \left( 
\begin{array}{cc}
1 & 0 \\ 
0 & 0%
\end{array}%
\right) \left( 
\begin{array}{cc}
\cos \theta & \sin \theta \\ 
-\sin \theta & \cos \theta%
\end{array}%
\right) \\
\Pi _{x}\left( D_{2}\left( \theta \right) \right) &=&\left( 
\begin{array}{cc}
\cos \theta & -\sin \theta \\ 
\sin \theta & \cos \theta%
\end{array}%
\right) \left( 
\begin{array}{cc}
0 & 0 \\ 
0 & 1%
\end{array}%
\right) \left( 
\begin{array}{cc}
\cos \theta & \sin \theta \\ 
-\sin \theta & \cos \theta%
\end{array}%
\right) \\
\Pi _{x}\left( D_{\limfunc{tr}}\left( \theta \right) \right) &=&\left( 
\begin{array}{cc}
\cos \theta & -\sin \theta \\ 
\sin \theta & \cos \theta%
\end{array}%
\right) \left( 
\begin{array}{cc}
\frac{1}{2} & 0 \\ 
0 & \frac{1}{2}%
\end{array}%
\right) \left( 
\begin{array}{cc}
\cos \theta & \sin \theta \\ 
-\sin \theta & \cos \theta%
\end{array}%
\right) =\left( 
\begin{array}{cc}
\frac{1}{2} & 0 \\ 
0 & \frac{1}{2}%
\end{array}%
\right)
\end{eqnarray}%
corresponding to the density operators 
\begin{eqnarray}
D_{1}\left( \theta \right) &=&o\left( \theta \right) \left( I_{2}+\widetilde{%
x}_{2}\right) o\left( \theta \right) ^{\dagger }=I_{2}+o\left( \theta
\right) \left( \widetilde{x}_{2}\right) o\left( \theta \right) ^{\dagger } \\
D_{2}\left( \theta \right) &=&o\left( \theta \right) \left( I_{2}-\widetilde{%
x}_{2}\right) o\left( \theta \right) ^{\dagger }=I_{2}-o\left( \theta
\right) \left( \widetilde{x}_{2}\right) o\left( \theta \right) ^{\dagger } \\
D_{\limfunc{tr}}\left( \theta \right) &=&o\left( \theta \right) I_{2}o\left(
\theta \right) ^{\dagger }=I_{2}.
\end{eqnarray}%
The transformation from the basis $\left\vert \widetilde{x}_{1},\widetilde{x}%
_{2}\right) $ to $\left\vert \widetilde{x}_{2},\widetilde{x}_{1}\right) $ is
given by 
\begin{equation}
\left\vert \widetilde{x}_{1},\widetilde{x}_{2}\right) \left( 
\begin{array}{cc}
0 & 1 \\ 
1 & 0%
\end{array}%
\right) =\left\vert \widetilde{x}_{2},\widetilde{x}_{1}\right) ,
\end{equation}%
which is a reflection as 
\begin{equation}
\det \left\{ \left( 
\begin{array}{cc}
0 & 1 \\ 
1 & 0%
\end{array}%
\right) \right\} =-1.
\end{equation}%
This transformation shows that $\mathcal{\Pi }_{x}\left( D_{1}\right) $ and $%
\mathcal{\Pi }_{x}\left( D_{2}\right) $ are on the same $O\left( 2,\mathbb{R}%
\right) $ orbit as 
\begin{equation}
\left( 
\begin{array}{cc}
0 & 1 \\ 
1 & 0%
\end{array}%
\right) \left( 
\begin{array}{cc}
0 & 0 \\ 
0 & 1%
\end{array}%
\right) \left( 
\begin{array}{cc}
0 & 1 \\ 
1 & 0%
\end{array}%
\right) =\left( 
\begin{array}{cc}
0 & 1 \\ 
0 & 0%
\end{array}%
\right) \left( 
\begin{array}{cc}
0 & 1 \\ 
1 & 0%
\end{array}%
\right) =\left( 
\begin{array}{cc}
1 & 0 \\ 
0 & 0%
\end{array}%
\right) .
\end{equation}

\item We can compare to the complexification $M\left( 2,\mathbb{C}\right) $
where we have one dimensional projectors defined by the partial isometries $%
\left\{ a^{\dagger },a\right\} $ 
\begin{equation}
a=\frac{1}{2}\left( \widetilde{x}_{1}-i\widetilde{x}_{2}\right) ,
\end{equation}%
which give 
\begin{gather}
D_{3}=a^{\dagger }a=\frac{1}{4}\left( \widetilde{x}_{1}+i\widetilde{x}%
_{2}\right) \left( \widetilde{x}_{1}-i\widetilde{x}_{2}\right) =\frac{1}{2}%
\left( I_{2}+i\widetilde{x}_{2}\widetilde{x}_{1}\right) \\
D_{4}=aa^{\dagger }=\frac{1}{4}\left( \widetilde{x}_{1}-i\widetilde{x}%
_{2}\right) \left( \widetilde{x}_{1}+i\widetilde{x}_{2}\right) =\frac{1}{2}%
\left( I_{2}-i\widetilde{x}_{2}\widetilde{x}_{1}\right)
\end{gather}%
where 
\begin{gather}
a=\frac{1}{2}\left( \widetilde{x}_{1}-i\widetilde{x}_{2}\right)
=q_{0}I_{2}+q_{11}\widetilde{x}_{1}+q_{12}\widetilde{x}_{2}+q_{2}\widetilde{x%
}_{1}\widetilde{x}_{2} \\
q_{0}=0,q_{11}=\frac{1}{2},q_{12}=-i,q_{2}=0
\end{gather}%
that are idempotent 
\begin{eqnarray}
D_{3}^{2} &=&\frac{1}{4}\left( I_{2}+i\widetilde{x}_{2}\widetilde{x}%
_{1}\right) \left( I_{2}+i\widetilde{x}_{2}\widetilde{x}_{1}\right) =\frac{1%
}{4}\left( I_{2}+2i\widetilde{x}_{2}\widetilde{x}_{1}+I_{2}\right) =\frac{1}{%
2}\left( I_{2}+i\widetilde{x}_{2}\widetilde{x}_{1}\right) \\
D_{4}^{2} &=&\frac{1}{4}\left( I_{2}-i\widetilde{x}_{2}\widetilde{x}%
_{1}\right) \left( I_{2}-i\widetilde{x}_{2}\widetilde{x}_{1}\right) =\frac{1%
}{4}\left( I_{2}-2i\widetilde{x}_{2}\widetilde{x}_{1}+I_{2}\right) =\frac{1}{%
2}\left( I_{2}-i\widetilde{x}_{2}\widetilde{x}_{1}\right) .
\end{eqnarray}%
Using the $\Pi _{x}$ representation one obtains the representations of $%
\mathcal{\Pi }_{x}\left( Q\right) ,\mathcal{\Pi }_{x}\left( D_{3}\right) $
and $\mathcal{\Pi }_{x}\left( D_{4}\right) $ as 
\begin{eqnarray*}
\mathcal{\Pi }_{x}\left( a\right) &=&\left( 
\begin{array}{cc}
q_{0}+q_{12} & q_{11}-q_{2} \\ 
q_{11}+q_{2} & q_{0}-q_{12}%
\end{array}%
\right) =\left( 
\begin{array}{cc}
0-\frac{i}{2} & \frac{1}{2}-0 \\ 
\frac{1}{2}+0 & 0+\frac{i}{2}%
\end{array}%
\right) =\left( 
\begin{array}{cc}
-\frac{i}{2} & \frac{1}{2} \\ 
\frac{1}{2} & \frac{i}{2}%
\end{array}%
\right) \\
\mathcal{\Pi }_{x}\left( D_{3}\right) &=&\left( 
\begin{array}{cc}
-\frac{i}{2} & \frac{1}{2} \\ 
\frac{1}{2} & \frac{i}{2}%
\end{array}%
\right) \left( 
\begin{array}{cc}
\frac{i}{2} & \frac{1}{2} \\ 
\frac{1}{2} & -\frac{i}{2}%
\end{array}%
\right) =\left( 
\begin{array}{cc}
\frac{1}{4}+\frac{1}{4} & -\frac{i}{4}-\frac{i}{4} \\ 
\frac{i}{4}+\frac{i}{4} & \frac{1}{4}+\frac{1}{4}%
\end{array}%
\right) =\left( 
\begin{array}{cc}
\frac{1}{2} & -\frac{i}{2} \\ 
\frac{i}{2} & \frac{1}{2}%
\end{array}%
\right) \\
\mathcal{\Pi }_{x}\left( D_{4}\right) &=&\left( 
\begin{array}{cc}
\frac{i}{2} & \frac{1}{2} \\ 
\frac{1}{2} & -\frac{i}{2}%
\end{array}%
\right) \left( 
\begin{array}{cc}
-\frac{i}{2} & \frac{1}{2} \\ 
\frac{1}{2} & \frac{i}{2}%
\end{array}%
\right) =\left( 
\begin{array}{cc}
\frac{1}{2} & \frac{i}{2} \\ 
-\frac{i}{2} & \frac{1}{2}%
\end{array}%
\right) .
\end{eqnarray*}%
Density matrices in the same $U\left( 2\right) $ orbit as $\mathcal{\Pi }%
_{x}\left( D_{3}\right) $ and $\mathcal{\Pi }_{x}\left( D_{4}\right) $ are
given by 
\begin{eqnarray}
\mathcal{\Pi }_{x}\left( D_{3}\right) \left( \mathbf{\lambda }\right)
&=&U\left( \mathbf{\lambda }\right) \mathcal{\Pi }_{x}\left( D_{3}\right)
U\left( \mathbf{\lambda }\right) ^{\dagger } \\
\mathcal{\Pi }_{x}\left( D_{4}\right) \left( \mathbf{\lambda }\right)
&=&U\left( \mathbf{\lambda }\right) \mathcal{\Pi }_{x}\left( D_{4}\right)
U\left( \mathbf{\lambda }\right) ^{\dagger }.
\end{eqnarray}%
In particular $K\in U\left( 2\right) $ and 
\begin{eqnarray}
K^{\dagger }\mathcal{\Pi }_{x}\left( D_{3}\right) K &=&\frac{1}{2}\left( 
\begin{array}{cc}
1 & i \\ 
1 & -i%
\end{array}%
\right) \left( 
\begin{array}{cc}
\frac{1}{2} & -\frac{i}{2} \\ 
\frac{i}{2} & \frac{1}{2}%
\end{array}%
\right) \left( 
\begin{array}{cc}
1 & 1 \\ 
-i & i%
\end{array}%
\right) =\frac{1}{2}\left( 
\begin{array}{cc}
\frac{1}{2}-\frac{1}{2} & -\frac{i}{2}+\frac{i}{2} \\ 
\frac{1}{2}+\frac{1}{2} & -\frac{i}{2}-\frac{i}{2}%
\end{array}%
\right) \left( 
\begin{array}{cc}
1 & 1 \\ 
-i & i%
\end{array}%
\right) \\
&=&\frac{1}{2}\left( 
\begin{array}{cc}
0 & 0 \\ 
1 & -i%
\end{array}%
\right) \left( 
\begin{array}{cc}
1 & 1 \\ 
-i & i%
\end{array}%
\right) =\left( 
\begin{array}{cc}
0 & 0 \\ 
0 & 1%
\end{array}%
\right) =\mathcal{\Pi }_{x}\left( D_{2}\right)
\end{eqnarray}%
and 
\begin{eqnarray}
K^{\dagger }\mathcal{\Pi }_{x}\left( D_{4}\right) K &=&\frac{1}{2}\left( 
\begin{array}{cc}
1 & i \\ 
1 & -i%
\end{array}%
\right) \left( 
\begin{array}{cc}
\frac{1}{2} & \frac{i}{2} \\ 
-\frac{i}{2} & \frac{1}{2}%
\end{array}%
\right) \left( 
\begin{array}{cc}
1 & 1 \\ 
-i & i%
\end{array}%
\right) =\frac{1}{2}\left( 
\begin{array}{cc}
\frac{1}{2}+\frac{1}{2} & \frac{i}{2}+\frac{i}{2} \\ 
\frac{1}{2}-\frac{1}{2} & \frac{i}{2}-\frac{i}{2}%
\end{array}%
\right) \left( 
\begin{array}{cc}
1 & 1 \\ 
-i & i%
\end{array}%
\right) \\
&=&\frac{1}{2}\left( 
\begin{array}{cc}
1 & i \\ 
0 & 0%
\end{array}%
\right) \left( 
\begin{array}{cc}
1 & 1 \\ 
-i & i%
\end{array}%
\right) =\left( 
\begin{array}{cc}
1 & 0 \\ 
0 & 0%
\end{array}%
\right) =\mathcal{\Pi }_{x}\left( D_{2}\right) .
\end{eqnarray}
\end{enumerate}

\begin{remark}
\begin{equation}
2\limfunc{tr}\nolimits_{\mathcal{F}}\left\{ D\right\} =2\limfunc{tr}%
\nolimits_{\Pi \left( \mathcal{F}\right) }\left\{ \Pi \left( D\right)
\right\} =\func{Tr}_{\Pi \left( \mathcal{F}\right) }\left\{ \Pi \left(
D\right) \right\} =1.
\end{equation}
\end{remark}

\paragraph{GNS\ Representations of $\mathcal{F}$}

\begin{lemma}
The density operators $D_{1}=\frac{1}{2}\left( I_{2}+\widetilde{x}%
_{2}\right) $ and $D_{2}=\frac{1}{2}\left( I_{2}-\widetilde{x}_{2}\right) $
are extreme states of $\mathcal{F}$.
\end{lemma}

\begin{proof}
Let the amplitude $Q_{1}$ be defined by 
\begin{equation}
Q_{1}=\frac{1}{2}\left( \widetilde{x}_{1}+\widetilde{x}_{1}\widetilde{x}%
_{2}\right)
\end{equation}%
then one can form the positive operator 
\begin{equation}
D_{1}=Q_{1}Q_{1}^{\dagger }=\frac{1}{2}\left( \widetilde{x}_{1}-\widetilde{x}%
_{1}\widetilde{x}_{2}\right) \frac{1}{2}\left( \widetilde{x}_{1}+\widetilde{x%
}_{1}\widetilde{x}_{2}\right) =\frac{1}{4}\left( I_{2}+\widetilde{x}_{2}+%
\widetilde{x}_{2}+I_{2}\right) =\frac{1}{2}\left( I_{2}+\widetilde{x}%
_{2}\right) ,
\end{equation}%
where 
\begin{equation}
D_{1}^{2}=\frac{1}{4}\left( I_{2}+\widetilde{x}_{2}\right) ^{2}=\frac{1}{4}%
\left( I_{2}+2\widetilde{x}_{2}+I_{2}\right) =\frac{1}{2}\left( I_{2}+%
\widetilde{x}_{2}\right) =D_{1}
\end{equation}%
and 
\begin{equation}
\func{Tr}\left\{ D_{1}\right\} =1.
\end{equation}%
This shows that $D_{1}$ is a pure state of $\mathcal{B}\left( \mathcal{H}%
\right) \cong \mathcal{B}\left( \mathbb{R}^{2}\right) $ and $Q_{1}$ is a
partial isometry$.$

Let $Q_{2}$ be defined by 
\begin{equation}
Q_{2}=\frac{1}{2}\left( \widetilde{x}_{1}-\widetilde{x}_{1}\widetilde{x}%
_{2}\right)
\end{equation}%
then one can form the positive operator 
\begin{equation}
D_{2}=Q_{2}Q_{2}^{\dagger }=\frac{1}{2}\left( \widetilde{x}_{1}+\widetilde{x}%
_{1}\widetilde{x}_{2}\right) \frac{1}{2}\left( \widetilde{x}_{1}-\widetilde{x%
}_{1}\widetilde{x}_{2}\right) =\frac{1}{4}\left( I_{2}-\widetilde{x}_{2}-%
\widetilde{x}_{2}+I_{2}\right) =\frac{1}{2}\left( I_{2}-\widetilde{x}%
_{2}\right) ,
\end{equation}%
where 
\begin{equation}
D_{1}^{2}=\frac{1}{4}\left( I_{2}-\widetilde{x}_{2}\right) ^{2}=\frac{1}{4}%
\left( I_{2}-2\widetilde{x}_{2}+I_{2}\right) =\frac{1}{2}\left( I_{2}-%
\widetilde{x}_{2}\right) =D_{1}
\end{equation}%
and 
\begin{equation}
\func{Tr}\left\{ D_{2}\right\} =1.
\end{equation}%
This shows that $D_{2}$ is a pure state of $\mathcal{B}\left( \mathcal{H}%
\right) ,$ where $\mathbb{R}$-$\dim \left\{ \mathcal{H}\right\} =2$ and $%
Q_{2}$ is a partial isometry$.$
\end{proof}

\begin{corollary}
$Q_{1}Q_{1}^{\dagger }=Q_{2}^{\dagger }Q_{2}$ and $Q_{1}^{\dagger
}Q_{1}=Q_{2}Q_{2}^{\dagger }.$
\end{corollary}

\begin{lemma}
The maximal left ideals of $\mathcal{F}$ are $\mathcal{L}_{1}=\mathbb{R-}%
\limfunc{linspan}\left\{ I_{2}+\widetilde{x}_{2},\widetilde{x}_{1}+%
\widetilde{x}_{1}\widetilde{x}_{2},\right\} =\ker \left\{ \Phi
_{D_{1}}\right\} $ and $\mathcal{L}_{2}=\mathbb{R-}\limfunc{linspan}\left\{
I_{2}-\widetilde{x}_{2},\widetilde{x}_{1}-\widetilde{x}_{1}\widetilde{x}%
_{2}\right\} =\ker \left\{ \Phi _{D_{2}}\right\} $

\begin{proof}
$\mathcal{F}$ acting on $\mathcal{L}_{1}$ gives 
\begin{eqnarray}
\widetilde{x}_{1}\left( \widetilde{x}_{1}+\widetilde{x}_{1}\widetilde{x}%
_{2}\right) &=&I_{2}+\widetilde{x}_{2} \\
\widetilde{x}_{2}\left( \widetilde{x}_{1}+\widetilde{x}_{1}\widetilde{x}%
_{2}\right) &=&-\widetilde{x}_{1}\widetilde{x}_{2}-\widetilde{x}_{1}=-\left( 
\widetilde{x}_{1}+\widetilde{x}_{1}\widetilde{x}_{2}\right) \\
\widetilde{x}_{1}\widetilde{x}_{2}\left( \widetilde{x}_{1}+\widetilde{x}_{1}%
\widetilde{x}_{2}\right) &=&-\widetilde{x}_{2}-I_{2}=-\left( I_{2}+%
\widetilde{x}_{2}\right)
\end{eqnarray}%
and 
\begin{eqnarray}
\widetilde{x}_{1}\left( I_{2}+\widetilde{x}_{2}\right) &=&\widetilde{x}_{1}+%
\widetilde{x}_{1}\widetilde{x}_{2} \\
\widetilde{x}_{2}\left( I_{2}+\widetilde{x}_{2}\right) &=&\widetilde{x}%
_{2}+I_{2} \\
\widetilde{x}_{1}\widetilde{x}_{2}\left( I_{2}+\widetilde{x}_{2}\right) &=&%
\widetilde{x}_{1}\widetilde{x}_{2}+\widetilde{x}_{1},
\end{eqnarray}%
hence 
\begin{equation}
\mathcal{F}:\mathcal{L}_{1}\rightarrow \mathcal{L}_{1}.
\end{equation}%
$\mathcal{F}$ acting on $\mathcal{L}_{2}$ gives 
\begin{eqnarray}
\widetilde{x}_{1}\left( \widetilde{x}_{1}-\widetilde{x}_{1}\widetilde{x}%
_{2}\right) &=&I_{2}-\widetilde{x}_{2} \\
\widetilde{x}_{2}\left( \widetilde{x}_{1}-\widetilde{x}_{1}\widetilde{x}%
_{2}\right) &=&-\widetilde{x}_{1}\widetilde{x}_{2}+\widetilde{x}_{1}=%
\widetilde{x}_{1}-\widetilde{x}_{1}\widetilde{x}_{2} \\
\widetilde{x}_{1}\widetilde{x}_{2}\left( \widetilde{x}_{1}-\widetilde{x}_{1}%
\widetilde{x}_{2}\right) &=&-\widetilde{x}_{2}-I_{2}=-\left( I_{2}-%
\widetilde{x}_{2}\right)
\end{eqnarray}%
and%
\begin{eqnarray}
\widetilde{x}_{1}\left( I_{2}-\widetilde{x}_{2}\right) &=&\widetilde{x}_{1}-%
\widetilde{x}_{1}\widetilde{x}_{2} \\
\widetilde{x}_{2}\left( I_{2}-\widetilde{x}_{2}\right) &=&\widetilde{x}%
_{2}-I_{2}=-\left( I_{2}-\widetilde{x}_{2}\right) \\
\widetilde{x}_{1}\widetilde{x}_{2}\left( I_{2}-\widetilde{x}_{2}\right) &=&%
\widetilde{x}_{1}\widetilde{x}_{2}-\widetilde{x}_{1},
\end{eqnarray}%
hence 
\begin{equation}
\mathcal{F}:\mathcal{L}_{2}\rightarrow \mathcal{L}_{2}.
\end{equation}%
The left ideals are linearly independent as 
\begin{eqnarray}
\widetilde{x}_{1}+\widetilde{x}_{1}\widetilde{x}_{2} &=&\alpha \left( 
\widetilde{x}_{1}-\widetilde{x}_{1}\widetilde{x}_{2}\right) +\beta \left(
I_{2}-\widetilde{x}_{2}\right) \\
&\Rightarrow &\alpha =1\text{ and }\alpha =-1,\beta =0\text{ }
\end{eqnarray}%
and 
\begin{eqnarray}
\left( I_{2}+\widetilde{x}_{2}\right) &=&\alpha \left( \widetilde{x}_{1}-%
\widetilde{x}_{1}\widetilde{x}_{2}\right) +\beta \left( I_{2}-\widetilde{x}%
_{2}\right) \\
&\Rightarrow &\alpha =0,\beta =1\text{ and }\beta =-1.
\end{eqnarray}%
This shows that as vector spaces 
\begin{equation}
\mathcal{F}=\mathcal{L}_{2}\oplus \mathcal{L}_{1}=\ker \left\{ \Phi
_{D_{1}}\right\} ^{\bot }\oplus \ker \left\{ \Phi _{D_{2}}\right\} ^{\bot }.
\end{equation}
\end{proof}
\end{lemma}

\begin{lemma}
The GNS\ representation $\mathcal{\Pi }_{f}\left( \mathcal{F}\right) \equiv $
$\mathcal{\Pi }_{D_{f}}\left( \mathcal{F}\right) $ on the Hilbert space $%
\mathcal{\Phi }_{f}\left( \mathcal{F}\right) \equiv $ $\mathcal{\Phi }%
_{D_{f}}\left( \mathcal{F}\right) $ produced by the normal state $D_{f}$ is
identical to the GNS\ representation $\mathcal{\Pi }_{\mathcal{\Pi }%
_{D_{f}}\left( \mathcal{F}\right) }\left( \mathcal{\Pi }_{D_{f}}\left( 
\mathcal{F}\right) \right) $ on the Hilbert space $\mathcal{\Phi }_{\mathcal{%
\Pi }_{D_{f}}\left( \mathcal{F}\right) }\left( \mathcal{\Phi }_{D_{f}}\left( 
\mathcal{F}\right) \right) $ i.e.%
\begin{eqnarray}
\left[ \mathcal{\Pi }_{\mathcal{\Pi }_{D_{f}}\left( D_{f}\right) }\circ 
\mathcal{\Pi }_{D_{f}}\right] \left( \mathcal{F}\right) &=&\mathcal{\Pi }%
_{D_{f}}\left( \mathcal{F}\right) \\
\left[ \Phi _{\mathcal{\Pi }_{D_{f}}\left( D_{f}\right) }\circ \mathcal{\Phi 
}_{D_{f}}\right] \left( \mathcal{F}\right) &=&\mathcal{\Phi }_{D_{f}}\left( 
\mathcal{F}\right) .
\end{eqnarray}%
Also 
\begin{eqnarray}
\ker \left\{ \Phi _{D_{f}}\right\} &\subset &\mathcal{F} \\
\ker \left\{ \left[ \Phi _{\mathcal{\Pi }_{D_{f}}\left( D_{f}\right) }\circ 
\mathcal{\Phi }_{D_{f}}\right] \right\} &\subset &\mathcal{\Phi }%
_{D_{f}}\left( \mathcal{F}\right) .
\end{eqnarray}
\end{lemma}

One can observe that $\mathcal{\Pi }_{D_{1}}\left( D_{2}\right) =\left( 
\begin{array}{cc}
0 & 0 \\ 
0 & 1%
\end{array}%
\right) \in \ker \left\{ \Phi _{\mathcal{\Pi }_{D_{1}}\left( D_{1}\right)
}\right\} $ as 
\begin{gather}
\func{Tr}\left\{ \left( 
\begin{array}{cc}
1 & 0 \\ 
0 & 0%
\end{array}%
\right) ^{2}\left( 
\begin{array}{cc}
0 & 0 \\ 
0 & 1%
\end{array}%
\right) \right\} =0 \\
\equiv \func{Tr}\left\{ \left( 
\begin{array}{cc}
0 & 1 \\ 
0 & 0%
\end{array}%
\right) \left( 
\begin{array}{cc}
0 & 0 \\ 
1 & 0%
\end{array}%
\right) \left( 
\begin{array}{cc}
0 & 0 \\ 
0 & 1%
\end{array}%
\right) \right\} =\func{Tr}\left\{ \left( 
\begin{array}{cc}
1 & 0 \\ 
0 & 0%
\end{array}%
\right) \left( 
\begin{array}{cc}
0 & 0 \\ 
0 & 1%
\end{array}%
\right) \right\} =0
\end{gather}%
and $\mathcal{\Pi }_{D_{2}}\left( D_{1}\right) =\left( 
\begin{array}{cc}
1 & 0 \\ 
0 & 0%
\end{array}%
\right) \in \ker \left\{ \Phi _{\mathcal{\Pi }_{D_{2}}\left( D_{2}\right)
}\right\} $ as 
\begin{gather}
\func{Tr}\left\{ \left( 
\begin{array}{cc}
0 & 0 \\ 
0 & 1%
\end{array}%
\right) ^{2}\left( 
\begin{array}{cc}
1 & 0 \\ 
0 & 0%
\end{array}%
\right) \right\} =0 \\
\equiv \func{Tr}\left\{ \left( 
\begin{array}{cc}
0 & 0 \\ 
1 & 0%
\end{array}%
\right) \left( 
\begin{array}{cc}
0 & 1 \\ 
0 & 0%
\end{array}%
\right) \left( 
\begin{array}{cc}
1 & 0 \\ 
0 & 0%
\end{array}%
\right) \right\} =\func{Tr}\left\{ \left( 
\begin{array}{cc}
0 & 0 \\ 
0 & 1%
\end{array}%
\right) \left( 
\begin{array}{cc}
1 & 0 \\ 
0 & 0%
\end{array}%
\right) \right\} =0.
\end{gather}%
The kernals of $\Phi _{\mathcal{\Pi }_{D_{1}}\left( D_{1}\right) }$ and $%
\Phi _{\mathcal{\Pi }_{D_{2}}\left( D_{2}\right) }$ are given in full by 
\begin{equation}
\ker \left\{ \Phi _{\mathcal{\Pi }_{D_{1}}\left( D_{1}\right) }\right\} =%
\limfunc{linspan}\left\{ \left( 
\begin{array}{cc}
1 & 0 \\ 
0 & 0%
\end{array}%
\right) ,\left( 
\begin{array}{cc}
0 & 0 \\ 
1 & 0%
\end{array}%
\right) \right\} =\ker \left\{ \Phi _{\mathcal{\Pi }_{D_{1}}\left(
D_{1}\right) }\right\} ^{\bot }
\end{equation}%
and 
\begin{equation}
\ker \left\{ \Phi _{\mathcal{\Pi }_{D_{2}}\left( D_{2}\right) }\right\} =%
\limfunc{linspan}\left\{ \left( 
\begin{array}{cc}
0 & 0 \\ 
0 & 1%
\end{array}%
\right) ,\left( 
\begin{array}{cc}
0 & 1 \\ 
0 & 0%
\end{array}%
\right) \right\} =\ker \left\{ \Phi _{\mathcal{\Pi }_{D_{2}}\left(
D_{2}\right) }\right\} ^{\bot }
\end{equation}%
and are minimal left ideals in $M\left( 2,\mathbb{R}\right) ,$ note that
they are constituted by matrices that have the same zero columns.

\begin{notation}
We now use 
\begin{eqnarray}
\Phi _{D_{1}} &\equiv &\Phi _{\mathcal{\Pi }_{D_{1}}\left( D_{1}\right) } \\
\Phi _{D_{2}} &\equiv &\Phi _{\mathcal{\Pi }_{D_{2}}\left( D_{2}\right) }.
\end{eqnarray}
\end{notation}

The maps $\Phi _{D_{1}}$, $\Phi _{D_{2}}$ have the following properties

\begin{enumerate}
\item 
\begin{align}
\Phi _{D_{1}}\left( I_{2}\right) & =\Phi _{D_{1}}\left( \left( 
\begin{array}{cc}
1 & 0 \\ 
0 & 1%
\end{array}%
\right) \right) =\Phi _{D_{1}}\left( \left( 
\begin{array}{cc}
0 & 0 \\ 
0 & 1%
\end{array}%
\right) +\left( 
\begin{array}{cc}
1 & 0 \\ 
0 & 0%
\end{array}%
\right) \right) =\Phi _{D_{1}}\left( \left( 
\begin{array}{cc}
0 & 0 \\ 
0 & 1%
\end{array}%
\right) \right) \\
\Phi _{D_{1}}\left( \widetilde{x}_{1}\right) & =\Phi _{D_{1}}\left( \left( 
\begin{array}{cc}
0 & 1 \\ 
1 & 0%
\end{array}%
\right) \right) =\Phi _{D_{1}}\left( \left( 
\begin{array}{cc}
0 & 1 \\ 
0 & 0%
\end{array}%
\right) +\left( 
\begin{array}{cc}
0 & 0 \\ 
1 & 0%
\end{array}%
\right) \right) =\Phi _{D_{1}}\left( \left( 
\begin{array}{cc}
0 & 1 \\ 
0 & 0%
\end{array}%
\right) \right) \\
\Phi _{D_{1}}\left( \widetilde{x}_{2}\right) & =\Phi _{D_{1}}\left( \left( 
\begin{array}{cc}
1 & 0 \\ 
0 & -1%
\end{array}%
\right) \right) =\Phi _{D_{1}}\left( \left( 
\begin{array}{cc}
1 & 0 \\ 
0 & 0%
\end{array}%
\right) -\left( 
\begin{array}{cc}
0 & 0 \\ 
0 & 1%
\end{array}%
\right) \right) =-\Phi _{D_{1}}\left( \left( 
\begin{array}{cc}
0 & 0 \\ 
0 & 1%
\end{array}%
\right) \right) \\
\Phi _{D_{1}}\left( \widetilde{x}_{1}\widetilde{x}_{2}\right) & =\Phi
_{D_{1}}\left( \left( 
\begin{array}{cc}
0 & -1 \\ 
1 & 0%
\end{array}%
\right) \right) =\Phi _{D_{1}}\left( \left( 
\begin{array}{cc}
0 & 0 \\ 
1 & 0%
\end{array}%
\right) -\left( 
\begin{array}{cc}
0 & 1 \\ 
0 & 0%
\end{array}%
\right) \right) =-\Phi _{D_{1}}\left( \left( 
\begin{array}{cc}
0 & 1 \\ 
0 & 0%
\end{array}%
\right) \right)
\end{align}%
showing that 
\begin{eqnarray}
\Phi _{D_{1}}\left( I_{2}\right) &=&-\Phi _{D_{1}}\left( \widetilde{x}%
_{2}\right) \Leftrightarrow \Phi _{D_{1}}\left( I_{2}+\widetilde{x}%
_{2}\right) =0 \\
\Phi _{D_{1}}\left( \widetilde{x}_{1}\widetilde{x}_{2}\right) &=&-\Phi
_{D_{1}}\left( \widetilde{x}_{1}\right) \Leftrightarrow \Phi _{D_{1}}\left( 
\widetilde{x}_{1}\widetilde{x}_{2}+\widetilde{x}_{1}\right) =0
\end{eqnarray}%
and%
\begin{eqnarray}
\Phi _{D_{1}}\left( M\left( 2,\mathbb{R}\right) \right) &=&\mathbb{R-}%
\limfunc{linspan}\left\{ \Phi _{D_{1}}\left( I_{2}\right) ,\Phi
_{D_{1}}\left( \widetilde{x}_{1}\right) \right\} =\mathbb{R-}\limfunc{linspan%
}\left\{ \Phi _{D_{1}}\left( I_{2}\right) ,\Phi _{D_{1}}\left( \widetilde{x}%
_{1}\widetilde{x}_{2}\right) \right\} \\
&=&\mathbb{R-}\limfunc{linspan}\left\{ \Phi _{D_{1}}\left( \widetilde{x}%
_{2}\right) ,\Phi _{D_{1}}\left( \widetilde{x}_{1}\right) \right\} =\mathbb{%
R-}\limfunc{linspan}\left\{ \Phi _{D_{1}}\left( \widetilde{x}_{2}\right)
,\Phi _{D_{1}}\left( \widetilde{x}_{1}\widetilde{x}_{2}\right) \right\} \\
&=&\mathbb{R-}\limfunc{linspan}\left\{ \Phi _{D_{1}}\left( I_{2}-\widetilde{x%
}_{2}\right) ,\Phi _{D_{1}}\left( \widetilde{x}_{1}-\widetilde{x}_{1}%
\widetilde{x}_{2}\right) \right\} .
\end{eqnarray}

\item 
\begin{align}
\Phi _{D_{2}}\left( I_{2}\right) & =\Phi _{D_{2}}\left( 
\begin{array}{cc}
1 & 0 \\ 
0 & 1%
\end{array}%
\right) =\Phi _{D_{2}}\left( \left( 
\begin{array}{cc}
0 & 0 \\ 
0 & 1%
\end{array}%
\right) +\left( 
\begin{array}{cc}
1 & 0 \\ 
0 & 0%
\end{array}%
\right) \right) =\Phi _{D_{2}}\left( \left( 
\begin{array}{cc}
1 & 0 \\ 
0 & 0%
\end{array}%
\right) \right) \\
\Phi _{D_{2}}\left( \widetilde{x}_{1}\right) & =\Phi _{D_{2}}\left( 
\begin{array}{cc}
0 & 1 \\ 
1 & 0%
\end{array}%
\right) =\Phi _{D_{2}}\left( \left( 
\begin{array}{cc}
0 & 1 \\ 
0 & 0%
\end{array}%
\right) +\left( 
\begin{array}{cc}
0 & 0 \\ 
1 & 0%
\end{array}%
\right) \right) =\Phi _{D_{2}}\left( \left( 
\begin{array}{cc}
0 & 0 \\ 
1 & 0%
\end{array}%
\right) \right) \\
\Phi _{D_{2}}\left( \widetilde{x}_{2}\right) & =\Phi _{D_{2}}\left( 
\begin{array}{cc}
1 & 0 \\ 
0 & -1%
\end{array}%
\right) =\Phi _{D_{2}}\left( \left( 
\begin{array}{cc}
1 & 0 \\ 
0 & 0%
\end{array}%
\right) -\left( 
\begin{array}{cc}
0 & 0 \\ 
0 & 1%
\end{array}%
\right) \right) =\Phi _{D_{2}}\left( \left( 
\begin{array}{cc}
1 & 0 \\ 
0 & 0%
\end{array}%
\right) \right) \\
\Phi _{D_{2}}\left( \widetilde{x}_{1}\widetilde{x}_{2}\right) & =\Phi
_{D_{2}}\left( 
\begin{array}{cc}
0 & -1 \\ 
1 & 0%
\end{array}%
\right) =\Phi _{D_{2}}\left( \left( 
\begin{array}{cc}
0 & 0 \\ 
1 & 0%
\end{array}%
\right) -\left( 
\begin{array}{cc}
0 & 1 \\ 
0 & 0%
\end{array}%
\right) \right) =\Phi _{D_{2}}\left( \left( 
\begin{array}{cc}
0 & 0 \\ 
1 & 0%
\end{array}%
\right) \right)
\end{align}%
showing that 
\begin{eqnarray}
\Phi _{D_{2}}\left( I_{2}\right) &=&\Phi _{D_{2}}\left( \widetilde{x}%
_{2}\right) \Leftrightarrow \Phi _{D_{1}}\left( I_{2}-\widetilde{x}%
_{2}\right) =0 \\
\Phi _{D_{2}}\left( \widetilde{x}_{1}\widetilde{x}_{2}\right) &=&\Phi
_{D_{2}}\left( \widetilde{x}_{1}\right) \Leftrightarrow \Phi _{D_{1}}\left( 
\widetilde{x}_{1}\widetilde{x}_{2}-\widetilde{x}_{1}\right) =0.
\end{eqnarray}%
$\lceil $The ideal structure is confirmed by 
\begin{eqnarray}
\left( I_{2}-\widetilde{x}_{2}\right) \left( I_{2}-\widetilde{x}_{2}\right)
&=&2\left( I_{2}-\widetilde{x}_{2}\right) \\
\left( I_{2}-\widetilde{x}_{2}\right) \left( \widetilde{x}_{1}\widetilde{x}%
_{2}-\widetilde{x}_{1}\right) &=&\widetilde{x}_{1}\widetilde{x}_{2}-%
\widetilde{x}_{1}-\widetilde{x}_{2}\widetilde{x}_{1}\widetilde{x}_{2}+%
\widetilde{x}_{2}\widetilde{x}_{1}=0 \\
\left( \widetilde{x}_{1}\widetilde{x}_{2}-\widetilde{x}_{1}\right) \left(
I_{2}-\widetilde{x}_{2}\right) &=&\widetilde{x}_{1}\widetilde{x}_{2}-%
\widetilde{x}_{1}-\widetilde{x}_{1}+\widetilde{x}_{1}\widetilde{x}%
_{2}=2\left( \widetilde{x}_{1}\widetilde{x}_{2}-\widetilde{x}_{1}\right) \\
\left( \widetilde{x}_{1}\widetilde{x}_{2}-\widetilde{x}_{1}\right) \left( 
\widetilde{x}_{1}\widetilde{x}_{2}-\widetilde{x}_{1}\right) &=&-I_{2}-%
\widetilde{x}_{2}+\widetilde{x}_{2}+I_{2}=0 \\
I_{2}-\widetilde{x}_{2}+\widetilde{x}_{1}\widetilde{x}_{2}-\widetilde{x}_{1}
&=&I_{2}-\widetilde{x}_{1}-\widetilde{x}_{2}+\widetilde{x}_{1}\widetilde{x}%
_{2} \\
I_{2}-\widetilde{x}_{2}-\widetilde{x}_{1}\widetilde{x}_{2}+\widetilde{x}_{1}
&=&I_{2}+\widetilde{x}_{1}-\widetilde{x}_{2}-\widetilde{x}_{1}\widetilde{x}%
_{2}
\end{eqnarray}%
$\rfloor $

and%
\begin{eqnarray}
\Phi _{D_{2}}\left( M\left( 2,\mathbb{R}\right) \right) &=&\mathbb{R-}%
\limfunc{linspan}\left\{ \Phi _{D_{2}}\left( I_{2}\right) ,\Phi
_{D_{2}}\left( \widetilde{x}_{1}\right) \right\} =\mathbb{R-}\limfunc{linspan%
}\left\{ \Phi _{D_{2}}\left( I_{2}\right) ,\Phi _{D_{2}}\left( \widetilde{x}%
_{1}\widetilde{x}_{2}\right) \right\} \\
&=&\mathbb{R-}\limfunc{linspan}\left\{ \Phi _{D_{2}}\left( \widetilde{x}%
_{2}\right) ,\Phi _{D_{2}}\left( \widetilde{x}_{1}\right) \right\} =\mathbb{%
R-}\limfunc{linspan}\left\{ \Phi _{D_{2}}\left( \widetilde{x}_{2}\right)
,\Phi _{D_{2}}\left( \widetilde{x}_{1}\widetilde{x}_{2}\right) \right\} \\
&=&\mathbb{R-}\limfunc{linspan}\left\{ \Phi _{D_{1}}\left( I_{2}+\widetilde{x%
}_{2}\right) ,\Phi _{D_{1}}\left( \widetilde{x}_{1}+\widetilde{x}_{1}%
\widetilde{x}_{2}\right) \right\} .
\end{eqnarray}
\end{enumerate}

This shows that in\emph{\ both cases} we can use the basis $\left\{ I_{2},%
\widetilde{x}_{1}\widetilde{x}_{2}\right\} $ to generate the GNS\ spaces $%
\Phi _{D_{1}}\left( \mathcal{F}\right) \cong \Phi _{D_{1}}\left( M\left( 2,%
\mathbb{R}\right) \right) $ and $\Phi _{D_{2}}\left( \mathcal{F}\right)
\cong \Phi _{D_{2}}\left( M\left( 2,\mathbb{R}\right) \right) .$

The representations $\Pi _{D_{1}}$ and $\Pi _{D_{2}}$ of $\mathcal{F}$ are
given by noting that:

\begin{enumerate}
\item 
\begin{eqnarray}
\Phi _{D_{1}}\left( I_{2}\right) &=&-\Phi _{D_{1}}\left( \widetilde{x}%
_{2}\right) \Leftrightarrow \Phi _{D_{1}}\left( I_{2}+\widetilde{x}%
_{2}\right) =0 \\
\Phi _{D_{1}}\left( \widetilde{x}_{1}\right) &=&-\Phi _{D_{1}}\left( 
\widetilde{x}_{1}\widetilde{x}_{2}\right) \Leftrightarrow \Phi
_{D_{1}}\left( \widetilde{x}_{1}+\widetilde{x}_{1}\widetilde{x}_{2}\right)
=0,
\end{eqnarray}%
one obtains%
\begin{equation}
\Phi _{D_{1}}\left( \mathcal{F}\right) =\mathbb{R-}\limfunc{linspan}\left\{
\Phi _{D_{1}}\left( I_{2}\right) ,\Phi _{D_{1}}\left( \widetilde{x}_{1}%
\widetilde{x}_{2}\right) \right\} =\mathbb{R-}\limfunc{linspan}\left\{ \Phi
_{D_{1}}\left( I_{2}-\widetilde{x}_{2}\right) ,\Phi _{D_{1}}\left( 
\widetilde{x}_{1}-\widetilde{x}_{1}\widetilde{x}_{2}\right) \right\} .
\end{equation}%
The norms of $\Phi _{D_{1}}\left( I_{2}\right) $ and $\Phi _{D_{1}}\left( 
\widetilde{x}_{1}\widetilde{x}_{2}\right) $ are%
\begin{eqnarray}
\left\langle \left. \Phi _{D_{1}}\left( I_{2}\right) \right\vert \Phi
_{D_{1}}\left( I_{2}\right) \right\rangle _{D_{1}} &=&\limfunc{tr}\left\{
D_{1}\right\} =\func{Tr}\left\{ \Pi _{D_{1}}\left( D_{1}\right) \right\} =1
\\
\left\langle \left. \Phi _{D_{1}}\left( I_{2}\right) \right\vert \Phi
_{D_{1}}\left( \widetilde{x}_{2}\right) \right\rangle _{D_{1}} &=&\limfunc{tr%
}\left\{ \widetilde{x}_{2}D_{1}\right\} =\func{Tr}\left\{ \Pi _{D_{1}}\left( 
\widetilde{x}_{2}\right) \left( D_{1}\right) \right\} =\func{Tr}\left\{
\left( 
\begin{array}{cc}
1 & 0 \\ 
0 & -1%
\end{array}%
\right) \left( 
\begin{array}{cc}
1 & 0 \\ 
0 & 0%
\end{array}%
\right) \right\} =1 \\
\left\langle \left. \Phi _{D_{1}}\left( \widetilde{x}_{1}\widetilde{x}%
_{2}\right) \right\vert \Phi _{D_{1}}\left( \widetilde{x}_{1}\widetilde{x}%
_{2}\right) \right\rangle _{D_{1}} &=&\limfunc{tr}\left\{ D_{1}\right\} =%
\func{Tr}\left\{ \Pi _{D_{1}}\left( D_{1}\right) \right\} =1 \\
\left\langle \left. \Phi _{D_{1}}\left( \widetilde{x}_{1}\widetilde{x}%
_{2}\right) \right\vert \Phi _{D_{1}}\left( \widetilde{x}_{1}\right)
\right\rangle _{D_{1}} &=&\limfunc{tr}\left\{ \widetilde{x}_{2}\widetilde{x}%
_{1}\widetilde{x}_{1}D_{1}\right\} =\func{Tr}\left\{ \Pi _{D_{1}}\left( 
\widetilde{x}_{2}\widetilde{x}_{1}\widetilde{x}_{1}\right) \Pi
_{D_{1}}\left( D_{1}\right) \right\} =\func{Tr}\left\{ \left( 
\begin{array}{cc}
1 & 0 \\ 
0 & -1%
\end{array}%
\right) \left( 
\begin{array}{cc}
1 & 0 \\ 
0 & 0%
\end{array}%
\right) \right\} =1
\end{eqnarray}%
and they are orthogonal to each other 
\begin{eqnarray}
\left\langle \left. \Phi _{D_{1}}\left( I_{2}\right) \right\vert \Phi
_{D_{1}}\left( \widetilde{x}_{1}\right) \right\rangle _{D_{1}} &=&\limfunc{tr%
}\left\{ \widetilde{x}_{1}D_{1}\right\} =\func{Tr}\left\{ \Pi _{D_{1}}\left( 
\widetilde{x}_{1}\right) \Pi _{D_{1}}\left( D_{1}\right) \right\} =0 \\
\left\langle \left. \Phi _{D_{1}}\left( I_{2}\right) \right\vert \Phi
_{D_{1}}\left( \widetilde{x}_{2}\right) \right\rangle _{D_{1}} &=&\limfunc{tr%
}\left\{ \widetilde{x}_{2}D_{1}\right\} =\func{Tr}\left\{ \Pi _{D_{1}}\left( 
\widetilde{x}_{2}\right) \Pi _{D_{1}}\left( D_{1}\right) \right\} =0 \\
\left\langle \left. \Phi _{D_{1}}\left( I_{2}\right) \right\vert \Phi
_{D_{1}}\left( \widetilde{x}_{1}\widetilde{x}_{2}\right) \right\rangle
_{D_{1}} &=&\limfunc{tr}\left\{ \widetilde{x}_{1}\widetilde{x}%
_{2}D_{1}\right\} =\func{Tr}\left\{ \Pi _{D_{1}}\left( \widetilde{x}_{1}%
\widetilde{x}_{2}\right) \Pi _{D_{1}}\left( D_{1}\right) \right\} =0.
\end{eqnarray}%
The norms of $\frac{1}{\sqrt{2}}\Phi _{D_{1}}\left( I_{2}-\widetilde{x}%
_{2}\right) $ and $\frac{1}{\sqrt{2}}\Phi _{D_{1}}\left( \widetilde{x}_{1}-%
\widetilde{x}_{1}\widetilde{x}_{2}\right) $ are 
\begin{eqnarray}
\frac{1}{2}\left\langle \left. \Phi _{D_{1}}\left( I_{2}-\widetilde{x}%
_{2}\right) \right\vert \Phi _{D_{1}}\left( I_{2}-\widetilde{x}_{2}\right)
\right\rangle _{D_{1}} &=&\frac{1}{2}\func{Tr}\left\{ \Pi _{D_{1}}\left(
I_{2}-\widetilde{x}_{2}\right) \Pi _{D_{1}}\left( I_{2}-\widetilde{x}%
_{2}\right) \Pi _{D_{1}}\left( D_{1}\right) \right\} =1 \\
\frac{1}{2}\left\langle \left. \Phi _{D_{1}}\left( \widetilde{x}_{1}-%
\widetilde{x}_{1}\widetilde{x}_{2}\right) \right\vert \Phi _{D_{1}}\left( 
\widetilde{x}_{1}-\widetilde{x}_{1}\widetilde{x}_{2}\right) \right\rangle
_{D_{1}} &=&\frac{1}{2}\func{Tr}\left\{ \Pi _{D_{1}}\left( \widetilde{x}_{1}-%
\widetilde{x}_{1}\widetilde{x}_{2}\right) ^{\dagger }\Pi _{D_{1}}\left( 
\widetilde{x}_{1}-\widetilde{x}_{1}\widetilde{x}_{2}\right) \Pi
_{D_{1}}\left( D_{1}\right) \right\} =1
\end{eqnarray}%
and they are orthogonal to each other 
\begin{equation}
\frac{1}{2}\left\langle \left. \Phi _{D_{1}}\left( I_{2}-\widetilde{x}%
_{2}\right) \right\vert \Phi _{D_{1}}\left( \widetilde{x}_{1}-\widetilde{x}%
_{1}\widetilde{x}_{2}\right) \right\rangle _{D_{1}}=\frac{1}{2}\func{Tr}%
\left\{ \Pi _{D_{1}}\left( I_{2}-\widetilde{x}_{2}\right) \Pi _{D_{1}}\left( 
\widetilde{x}_{1}-\widetilde{x}_{1}\widetilde{x}_{2}\right) \Pi
_{D_{1}}\left( D_{1}\right) \right\} =0.
\end{equation}%
The representation $\Pi _{D_{1}}$ is defined by 
\begin{eqnarray}
&&\Pi _{D_{1}}\left( \widetilde{x}_{1}\right) \left( \Phi _{D_{1}}\left(
I_{2}-\widetilde{x}_{2}\right) ,\Phi _{D_{1}}\left( \widetilde{x}_{1}-%
\widetilde{x}_{1}\widetilde{x}_{2}\right) \right) =\left( \Phi
_{D_{1}}\left( \widetilde{x}_{1}-\widetilde{x}_{1}\widetilde{x}_{2}\right)
,\Phi _{D_{1}}\left( I_{2}-\widetilde{x}_{2}\right) ,\right) \\
&=&\Phi _{D_{1}}\left( I_{2}-\widetilde{x}_{2}\right) ,\Phi _{D_{1}}\left( 
\widetilde{x}_{1}-\widetilde{x}_{1}\widetilde{x}_{2}\right) \left( 
\begin{array}{cc}
0 & 1 \\ 
1 & 0%
\end{array}%
\right) \\
&&\Pi _{D_{1}}\left( \widetilde{x}_{2}\right) \left( \Phi _{D_{1}}\left(
I_{2}-\widetilde{x}_{2}\right) ,\Phi _{D_{1}}\left( \widetilde{x}_{1}-%
\widetilde{x}_{1}\widetilde{x}_{2}\right) \right) =\left( \Phi
_{D_{1}}\left( \widetilde{x}_{2}-I_{2}\right) ,\Phi _{D_{1}}\left( 
\widetilde{x}_{2}\widetilde{x}_{1}+\widetilde{x}_{1}\right) \right) \\
&=&\left( -\Phi _{D_{1}}\left( I_{2}-\widetilde{x}_{2}\right) ,\Phi
_{D_{1}}\left( \widetilde{x}_{1}-\widetilde{x}_{1}\widetilde{x}_{2}\right)
\right) =\Phi _{D_{1}}\left( I_{2}-\widetilde{x}_{2}\right) ,\Phi
_{D_{1}}\left( \widetilde{x}_{1}-\widetilde{x}_{1}\widetilde{x}_{2}\right)
\left( 
\begin{array}{cc}
-1 & 0 \\ 
0 & 1%
\end{array}%
\right) \\
&&\Pi _{D_{1}}\left( \widetilde{x}_{1}\widetilde{x}_{2}\right) \left( \Phi
_{D_{1}}\left( I_{2}-\widetilde{x}_{2}\right) ,\Phi _{D_{1}}\left( 
\widetilde{x}_{1}-\widetilde{x}_{1}\widetilde{x}_{2}\right) \right) =\left(
\Phi _{D_{1}}\left( \widetilde{x}_{1}\widetilde{x}_{2}-\widetilde{x}%
_{1}\right) ,\Phi _{D_{1}}\left( I_{2}-\widetilde{x}_{1}\right) \right) \\
&=&\Phi _{D_{1}}\left( I_{2}-\widetilde{x}_{2}\right) ,\Phi _{D_{1}}\left( 
\widetilde{x}_{1}-\widetilde{x}_{1}\widetilde{x}_{2}\right) \left( 
\begin{array}{cc}
0 & 1 \\ 
1 & 0%
\end{array}%
\right) .
\end{eqnarray}%
and

\item 
\begin{eqnarray}
\Phi _{D_{2}}\left( I_{2}\right) &=&\Phi _{D_{2}}\left( \widetilde{x}%
_{2}\right) \Leftrightarrow \Phi _{D_{2}}\left( I_{2}-\widetilde{x}%
_{2}\right) =0 \\
\Phi _{D_{2}}\left( \widetilde{x}_{1}\widetilde{x}_{2}\right) &=&\Phi
_{D_{2}}\left( \widetilde{x}_{1}\right) \Leftrightarrow \Phi _{D_{2}}\left( 
\widetilde{x}_{1}-\widetilde{x}_{1}\widetilde{x}_{2}\right) =0.
\end{eqnarray}%
one obtains%
\begin{equation}
\Phi _{D_{2}}\left( \mathcal{F}\right) =\mathbb{R-}\limfunc{linspan}\left\{
\Phi _{D_{2}}\left( I_{2}\right) ,\Phi _{D_{2}}\left( \widetilde{x}_{1}%
\widetilde{x}_{2}\right) \right\} =\mathbb{R-}\limfunc{linspan}\left\{ \frac{%
1}{\sqrt{2}}\Phi _{D_{1}}\left( I_{2}+\widetilde{x}_{2}\right) ,\frac{1}{%
\sqrt{2}}\Phi _{D_{1}}\left( \widetilde{x}_{1}+\widetilde{x}_{1}\widetilde{x}%
_{2}\right) \right\} .
\end{equation}%
giving%
\begin{eqnarray}
&&\Pi _{D_{2}}\left( \widetilde{x}_{1}\right) \left( \Phi _{D_{1}}\left(
I_{2}+\widetilde{x}_{2}\right) ,\Phi _{D_{1}}\left( \widetilde{x}_{1}+%
\widetilde{x}_{1}\widetilde{x}_{2}\right) \right) =\left( \Phi
_{D_{1}}\left( \widetilde{x}_{1}+\widetilde{x}_{1}\widetilde{x}_{2}\right)
,\Phi _{D_{1}}\left( I_{2}+\widetilde{x}_{2}\right) \right) \\
&=&\left( \Phi _{D_{1}}\left( I_{2}+\widetilde{x}_{2}\right) ,\Phi
_{D_{1}}\left( \widetilde{x}_{1}+\widetilde{x}_{1}\widetilde{x}_{2}\right)
\right) \left( 
\begin{array}{cc}
0 & 1 \\ 
1 & 0%
\end{array}%
\right) \\
&&\Pi _{D_{2}}\left( \widetilde{x}_{2}\right) \left( \Phi _{D_{1}}\left(
I_{2}+\widetilde{x}_{2}\right) ,\Phi _{D_{1}}\left( \widetilde{x}_{1}+%
\widetilde{x}_{1}\widetilde{x}_{2}\right) \right) =\left( \Phi
_{D_{1}}\left( I_{2}+\widetilde{x}_{2}\right) ,-\Phi _{D_{2}}\left( 
\widetilde{x}_{1}+\widetilde{x}_{1}\widetilde{x}_{2}\right) \right) \\
&=&\left( \Phi _{D_{1}}\left( I_{2}+\widetilde{x}_{2}\right) ,\Phi
_{D_{1}}\left( \widetilde{x}_{1}+\widetilde{x}_{1}\widetilde{x}_{2}\right)
\right) \left( 
\begin{array}{cc}
1 & 0 \\ 
0 & -1%
\end{array}%
\right) \\
&&\Pi _{D_{2}}\left( \widetilde{x}_{1}\widetilde{x}_{2}\right) \left( \Phi
_{D_{1}}\left( I_{2}+\widetilde{x}_{2}\right) ,\Phi _{D_{1}}\left( 
\widetilde{x}_{1}+\widetilde{x}_{1}\widetilde{x}_{2}\right) \right) =\left(
\Phi _{D_{1}}\left( \widetilde{x}_{1}+\widetilde{x}_{1}\widetilde{x}%
_{2}\right) ,-\Phi _{D_{1}}\left( I_{2}+\widetilde{x}_{2}\right) \right) \\
&=&\left( \Phi _{D_{1}}\left( I_{2}+\widetilde{x}_{2}\right) ,\Phi
_{D_{1}}\left( \widetilde{x}_{1}+\widetilde{x}_{1}\widetilde{x}_{2}\right)
\right) \left( 
\begin{array}{cc}
0 & -1 \\ 
1 & 0%
\end{array}%
\right) .
\end{eqnarray}
\end{enumerate}

\begin{lemma}
$\mathcal{F}$ acting as linear operators on $\mathcal{F}$ is reduced by the
direct sum left module decomposition 
\begin{equation}
\mathcal{F}:\mathcal{L}_{2}\oplus \mathcal{L}_{1}\rightarrow \mathcal{L}%
_{2}\oplus \mathcal{L}_{1}.
\end{equation}
\end{lemma}

Consider the tracial state $\limfunc{tr}$ on $\mathcal{F}.$ This state is
also a linear functional $\Pi _{\limfunc{tr}}\left( \mathcal{F}\right) $ on
where $\limfunc{tr}=\frac{1}{2}\func{Tr}$ and defines a Hilbert-Schmidt
Hilbert space structure i.e. 
\begin{equation}
\left\langle \left. A\right\vert B\right\rangle _{\mathcal{F}}=\limfunc{tr}%
\left\{ A^{\dagger }B\right\} =\limfunc{tr}\left\{ \Pi _{\limfunc{tr}}\left(
A\right) ^{\dagger }\Pi _{\limfunc{tr}}\left( B\right) \right\} =\frac{1}{2}%
\func{Tr}\left\{ \Pi _{\limfunc{tr}}\left( A\right) ^{\dagger }\Pi _{%
\limfunc{tr}}\left( B\right) \right\} .
\end{equation}%
A \emph{conb} for $\mathcal{F}$ wrt $\left\langle \left. {}\right\vert
\right\rangle _{\mathcal{F}}$ is given by $\left\{ \frac{1}{\sqrt{2}}\left(
I_{2}+\widetilde{x}_{2}\right) ,\frac{1}{\sqrt{2}}\left( \widetilde{x}_{1}+%
\widetilde{x}_{1}\widetilde{x}_{2}\right) ,\frac{1}{\sqrt{2}}\left( I_{2}-%
\widetilde{x}_{2}\right) ,\frac{1}{\sqrt{2}}\left( \widetilde{x}_{1}-%
\widetilde{x}_{1}\widetilde{x}_{2}\right) .\right\} $

$\lceil $ 
\begin{eqnarray}
\frac{1}{2}\left\langle \left. \left( \widetilde{x}_{1}+\widetilde{x}_{1}%
\widetilde{x}_{2}\right) \right\vert \left( \widetilde{x}_{1}+\widetilde{x}%
_{1}\widetilde{x}_{2}\right) \right\rangle _{\mathcal{F}} &=&\frac{1}{2}%
\limfunc{tr}\left\{ \left( \widetilde{x}_{1}+\widetilde{x}_{2}\widetilde{x}%
_{1}\right) \left( \widetilde{x}_{1}+\widetilde{x}_{1}\widetilde{x}%
_{2}\right) \right\} =1 \\
\frac{1}{2}\left\langle \left. \left( \widetilde{x}_{1}+\widetilde{x}_{1}%
\widetilde{x}_{2}\right) \right\vert \left( I_{2}+\widetilde{x}_{2}\right)
\right\rangle _{\mathcal{F}} &=&\frac{1}{2}\limfunc{tr}\left\{ \left( 
\widetilde{x}_{1}+\widetilde{x}_{2}\widetilde{x}_{1}\right) \left( I_{2}+%
\widetilde{x}_{2}\right) \right\} =0 \\
\frac{1}{2}\left\langle \left. \left( \widetilde{x}_{1}+\widetilde{x}_{1}%
\widetilde{x}_{2}\right) \right\vert \left( \widetilde{x}_{1}-\widetilde{x}%
_{1}\widetilde{x}_{2}\right) \right\rangle _{\mathcal{F}} &=&\frac{1}{2}%
\limfunc{tr}\left\{ \left( \widetilde{x}_{1}+\widetilde{x}_{2}\widetilde{x}%
_{1}\right) \left( \widetilde{x}_{1}-\widetilde{x}_{1}\widetilde{x}%
_{2}\right) \right\} =0 \\
\frac{1}{2}\left\langle \left. \left( \widetilde{x}_{1}+\widetilde{x}_{1}%
\widetilde{x}_{2}\right) \right\vert \left( I_{2}-\widetilde{x}_{2}\right)
\right\rangle _{\mathcal{F}} &=&\frac{1}{2}\limfunc{tr}\left\{ \left( 
\widetilde{x}_{1}+\widetilde{x}_{2}\widetilde{x}_{1}\right) \left( I_{2}-%
\widetilde{x}_{2}\right) \right\} =0
\end{eqnarray}%
\begin{eqnarray}
\frac{1}{2}\left\langle \left. \left( I_{2}+\widetilde{x}_{2}\right)
\right\vert \left( I_{2}+\widetilde{x}_{2}\right) \right\rangle _{\mathcal{F}%
} &=&\frac{1}{2}\limfunc{tr}\left\{ \left( I_{2}+\widetilde{x}_{2}\right)
\left( I_{2}+\widetilde{x}_{2}\right) \right\} =1 \\
\frac{1}{2}\left\langle \left. \left( I_{2}+\widetilde{x}_{2}\right)
\right\vert \left( \widetilde{x}_{1}-\widetilde{x}_{1}\widetilde{x}%
_{2}\right) \right\rangle _{\mathcal{F}} &=&\frac{1}{2}\limfunc{tr}\left\{
\left( I_{2}+\widetilde{x}_{2}\right) \left( \widetilde{x}_{1}-\widetilde{x}%
_{1}\widetilde{x}_{2}\right) \right\} =0 \\
\frac{1}{2}\left\langle \left. \left( I_{2}+\widetilde{x}_{2}\right)
\right\vert \left( I_{2}-\widetilde{x}_{2}\right) \right\rangle _{\mathcal{F}%
} &=&\frac{1}{2}\limfunc{tr}\left\{ \left( I_{2}+\widetilde{x}_{2}\right)
\left( I_{2}-\widetilde{x}_{2}\right) \right\} =0
\end{eqnarray}%
\begin{eqnarray}
\frac{1}{2}\left\langle \left. \left( \widetilde{x}_{1}-\widetilde{x}_{1}%
\widetilde{x}_{2}\right) \right\vert \left( \widetilde{x}_{1}-\widetilde{x}%
_{1}\widetilde{x}_{2}\right) \right\rangle _{\mathcal{F}} &=&\frac{1}{2}%
\limfunc{tr}\left\{ \left( \widetilde{x}_{1}-\widetilde{x}_{2}\widetilde{x}%
_{1}\right) \left( \widetilde{x}_{1}-\widetilde{x}_{1}\widetilde{x}%
_{2}\right) \right\} =1 \\
\frac{1}{2}\left\langle \left. \left( \widetilde{x}_{1}-\widetilde{x}_{1}%
\widetilde{x}_{2}\right) \right\vert \left( I_{2}-\widetilde{x}_{2}\right)
\right\rangle _{\mathcal{F}} &=&\frac{1}{2}\limfunc{tr}\left\{ \left( 
\widetilde{x}_{1}-\widetilde{x}_{2}\widetilde{x}_{1}\right) \left( I_{2}-%
\widetilde{x}_{2}\right) \right\} =0
\end{eqnarray}%
\begin{equation}
\frac{1}{2}\left\langle \left. \left( I_{2}-\widetilde{x}_{2}\right)
\right\vert \left( I_{2}-\widetilde{x}_{2}\right) \right\rangle _{\mathcal{F}%
}=\frac{1}{2}\limfunc{tr}\left\{ \left( I_{2}-\widetilde{x}_{2}\right)
\left( I_{2}-\widetilde{x}_{2}\right) \right\} =1.
\end{equation}%
$\rfloor $

The left regular representation $\Pi _{L}$ of $\mathcal{F}$ is given by the
action on $\mathcal{L}_{1}=\ker \left\{ \Phi _{D_{1}}\right\} =\ker \left\{
\Phi _{D_{2}}\right\} ^{\bot }=\mathcal{L}_{2}^{\bot }=\left\{ I_{2}+%
\widetilde{x}_{2},\widetilde{x}_{1}+\widetilde{x}_{1}\widetilde{x}%
_{2}\right\} $

\begin{enumerate}
\item 
\begin{eqnarray}
\Pi _{L}\left( \widetilde{x}_{1}\right) \left( \widetilde{x}_{1}+\widetilde{x%
}_{1}\widetilde{x}_{2}\right) &=&\widetilde{x}_{1}\left( \widetilde{x}_{1}+%
\widetilde{x}_{1}\widetilde{x}_{2}\right) =I_{2}+\widetilde{x}_{2} \\
\Pi _{L}\left( \widetilde{x}_{2}\right) \left( \widetilde{x}_{1}+\widetilde{x%
}_{1}\widetilde{x}_{2}\right) &=&-\widetilde{x}_{1}\widetilde{x}_{2}-%
\widetilde{x}_{1}=-\left( \widetilde{x}_{1}+\widetilde{x}_{1}\widetilde{x}%
_{2}\right) \\
\Pi _{L}\left( \widetilde{x}_{1}\widetilde{x}_{2}\right) \left( \widetilde{x}%
_{1}+\widetilde{x}_{1}\widetilde{x}_{2}\right) &=&-\widetilde{x}%
_{2}-I_{2}=-\left( I_{2}+\widetilde{x}_{2}\right)
\end{eqnarray}

\item 
\begin{eqnarray}
\Pi _{L}\left( \widetilde{x}_{1}\right) \left( I_{2}+\widetilde{x}%
_{2}\right) &=&\widetilde{x}_{1}\left( I_{2}+\widetilde{x}_{2}\right) =%
\widetilde{x}_{1}+\widetilde{x}_{1}\widetilde{x}_{2} \\
\Pi _{L}\left( \widetilde{x}_{2}\right) \left( I_{2}+\widetilde{x}%
_{2}\right) &=&\widetilde{x}_{2}\left( I_{2}+\widetilde{x}_{2}\right) =%
\widetilde{x}_{2}+I_{2} \\
\Pi _{L}\left( \widetilde{x}_{1}\widetilde{x}_{2}\right) \left( I_{2}+%
\widetilde{x}_{2}\right) &=&\widetilde{x}_{1}\widetilde{x}_{2}\left( I_{2}+%
\widetilde{x}_{2}\right) =\widetilde{x}_{1}\widetilde{x}_{2}+\widetilde{x}%
_{1},
\end{eqnarray}%
and by the action on $\mathcal{L}_{2}=\ker \left\{ \Phi _{D_{2}}\right\}
=\ker \left\{ \Phi _{D_{1}}\right\} ^{\bot }=\mathcal{L}_{1}^{\bot }=\left\{
I_{2}-\widetilde{x}_{2},\widetilde{x}_{1}-\widetilde{x}_{1}\widetilde{x}%
_{2}\right\} $

\item 
\begin{eqnarray}
\Pi _{L}\left( \widetilde{x}_{1}\right) \left( \widetilde{x}_{1}-\widetilde{x%
}_{1}\widetilde{x}_{2}\right) &=&\widetilde{x}_{1}\left( \widetilde{x}_{1}-%
\widetilde{x}_{1}\widetilde{x}_{2}\right) =I_{2}-\widetilde{x}_{2} \\
\Pi _{L}\left( \widetilde{x}_{2}\right) \left( \widetilde{x}_{1}-\widetilde{x%
}_{1}\widetilde{x}_{2}\right) &=&\widetilde{x}_{2}\left( \widetilde{x}_{1}-%
\widetilde{x}_{1}\widetilde{x}_{2}\right) =-\widetilde{x}_{1}\widetilde{x}%
_{2}+\widetilde{x}_{1} \\
\Pi _{L}\left( \widetilde{x}_{1}\widetilde{x}_{2}\right) \left( \widetilde{x}%
_{1}-\widetilde{x}_{1}\widetilde{x}_{2}\right) &=&\widetilde{x}_{1}%
\widetilde{x}_{2}\left( \widetilde{x}_{1}-\widetilde{x}_{1}\widetilde{x}%
_{2}\right) =-\widetilde{x}_{2}-I_{2}=-\left( I_{2}-\widetilde{x}_{2}\right)
\end{eqnarray}

\item 
\begin{eqnarray}
\Pi _{L}\left( \widetilde{x}_{1}\right) \left( I_{2}-\widetilde{x}%
_{2}\right) &=&\widetilde{x}_{1}\left( I_{2}-\widetilde{x}_{2}\right) =%
\widetilde{x}_{1}-\widetilde{x}_{1}\widetilde{x}_{2} \\
\Pi _{L}\left( \widetilde{x}_{2}\right) \left( I_{2}-\widetilde{x}%
_{2}\right) &=&\widetilde{x}_{2}\left( I_{2}-\widetilde{x}_{2}\right) =%
\widetilde{x}_{2}-I_{2}=-\left( I_{2}-\widetilde{x}_{2}\right) \\
\Pi _{L}\left( \widetilde{x}_{1}\widetilde{x}_{2}\right) \left( I_{2}-%
\widetilde{x}_{2}\right) &=&\widetilde{x}_{1}\widetilde{x}_{2}\left( I_{2}-%
\widetilde{x}_{2}\right) =\widetilde{x}_{1}\widetilde{x}_{2}-\widetilde{x}%
_{1}.
\end{eqnarray}
\end{enumerate}

This can be summarized as 
\begin{eqnarray}
&&\Pi _{L}\left( \widetilde{x}_{1}\right) \left( \frac{1}{\sqrt{2}}\left(
I_{2}+\widetilde{x}_{2}\right) ,\frac{1}{\sqrt{2}}\left( \widetilde{x}_{1}+%
\widetilde{x}_{1}\widetilde{x}_{2}\right) ,\frac{1}{\sqrt{2}}\left( I_{2}-%
\widetilde{x}_{2}\right) \frac{1}{\sqrt{2}}\left( \widetilde{x}_{1}-%
\widetilde{x}_{1}\widetilde{x}_{2}\right) \right) \\
&=&\left( \frac{1}{\sqrt{2}}\left( \widetilde{x}_{1}+\widetilde{x}_{1}%
\widetilde{x}_{2}\right) ,\frac{1}{\sqrt{2}}\left( I_{2}+\widetilde{x}%
_{2}\right) ,\frac{1}{\sqrt{2}}\left( I_{2}-\widetilde{x}_{2}\right) \frac{1%
}{\sqrt{2}}\left( \widetilde{x}_{1}-\widetilde{x}_{1}\widetilde{x}%
_{2}\right) \right) \left( 
\begin{array}{cc}
\left( 
\begin{array}{cc}
0 & 1 \\ 
1 & 0%
\end{array}%
\right) & \left( 
\begin{array}{cc}
0 & 0 \\ 
0 & 0%
\end{array}%
\right) \\ 
\left( 
\begin{array}{cc}
0 & 0 \\ 
0 & 0%
\end{array}%
\right) & \left( 
\begin{array}{cc}
0 & 1 \\ 
1 & 0%
\end{array}%
\right)%
\end{array}%
\right) ,
\end{eqnarray}%
\begin{eqnarray}
&&\Pi _{L}\left( \widetilde{x}_{2}\right) \left( \frac{1}{\sqrt{2}}\left(
I_{2}+\widetilde{x}_{2}\right) ,\frac{1}{\sqrt{2}}\left( \widetilde{x}_{1}+%
\widetilde{x}_{1}\widetilde{x}_{2}\right) ,\frac{1}{\sqrt{2}}\left( 
\widetilde{x}_{1}-\widetilde{x}_{1}\widetilde{x}_{2}\right) ,\frac{1}{\sqrt{2%
}}\left( I_{2}-\widetilde{x}_{2}\right) \right) \\
&=&\left( \frac{1}{\sqrt{2}}\left( \widetilde{x}_{1}+\widetilde{x}_{1}%
\widetilde{x}_{2}\right) ,\frac{1}{\sqrt{2}}\left( I_{2}+\widetilde{x}%
_{2}\right) ,\frac{1}{\sqrt{2}}\left( \widetilde{x}_{1}-\widetilde{x}_{1}%
\widetilde{x}_{2}\right) ,\frac{1}{\sqrt{2}}\left( I_{2}-\widetilde{x}%
_{2}\right) \right) \left( 
\begin{array}{cc}
\left( 
\begin{array}{cc}
1 & 0 \\ 
0 & -1%
\end{array}%
\right) & \left( 
\begin{array}{cc}
0 & 0 \\ 
0 & 0%
\end{array}%
\right) \\ 
\left( 
\begin{array}{cc}
0 & 0 \\ 
0 & 0%
\end{array}%
\right) & \left( 
\begin{array}{cc}
1 & 0 \\ 
0 & -1%
\end{array}%
\right)%
\end{array}%
\right) ,
\end{eqnarray}%
\begin{eqnarray}
&&\Pi _{L}\left( \widetilde{x}_{1}\widetilde{x}_{2}\right) \left( \frac{1}{%
\sqrt{2}}\left( \widetilde{x}_{1}+\widetilde{x}_{1}\widetilde{x}_{2}\right) ,%
\frac{1}{\sqrt{2}}\left( I_{2}+\widetilde{x}_{2}\right) ,\frac{1}{\sqrt{2}}%
\left( I_{2}-\widetilde{x}_{2}\right) \frac{1}{\sqrt{2}}\left( \widetilde{x}%
_{1}-\widetilde{x}_{1}\widetilde{x}_{2}\right) \right) \\
&=&\left( \frac{1}{\sqrt{2}}\left( \widetilde{x}_{1}+\widetilde{x}_{1}%
\widetilde{x}_{2}\right) ,\frac{1}{\sqrt{2}}\left( I_{2}+\widetilde{x}%
_{2}\right) ,\frac{1}{\sqrt{2}}\left( I_{2}-\widetilde{x}_{2}\right) \frac{1%
}{\sqrt{2}}\left( \widetilde{x}_{1}-\widetilde{x}_{1}\widetilde{x}%
_{2}\right) \right) \left( 
\begin{array}{cc}
\left( 
\begin{array}{cc}
0 & -1 \\ 
1 & 0%
\end{array}%
\right) & \left( 
\begin{array}{cc}
0 & 0 \\ 
0 & 0%
\end{array}%
\right) \\ 
\left( 
\begin{array}{cc}
0 & 0 \\ 
0 & 0%
\end{array}%
\right) & \left( 
\begin{array}{cc}
0 & -1 \\ 
1 & 0%
\end{array}%
\right)%
\end{array}%
\right) .
\end{eqnarray}%
The two orthonormal bases $\left\{ \frac{1}{\sqrt{2}}\Phi _{\limfunc{tr}%
}\left( I_{2}+\widetilde{x}_{2}\right) ,\frac{1}{\sqrt{2}}\Phi _{\limfunc{tr}%
}\left( \widetilde{x}_{1}+\widetilde{x}_{1}\widetilde{x}_{2}\right) ,\frac{1%
}{\sqrt{2}}\Phi _{\limfunc{tr}}\left( I_{2}-\widetilde{x}_{2}\right) \frac{1%
}{\sqrt{2}}\Phi _{\limfunc{tr}}\left( \widetilde{x}_{1}-\widetilde{x}_{1}%
\widetilde{x}_{2}\right) \right\} $ and $\left\{ \Phi _{\limfunc{tr}}\left(
I_{2}\right) ,\Phi _{\limfunc{tr}}\left( \widetilde{x}_{1}\right) ,\Phi _{%
\limfunc{tr}}\left( \widetilde{x}_{2}\right) ,\Phi _{\limfunc{tr}}\left( 
\widetilde{x}_{1}\widetilde{x}_{2}\right) \right\} $ of $\Phi _{\limfunc{tr}%
}\left( \mathcal{F}\right) $ $\equiv \mathcal{F}_{HS}$ have the properties:

\begin{enumerate}
\item They are normalized 
\begin{eqnarray}
\left\langle \left. \Phi _{\limfunc{tr}}\left( I_{2}\right) \right\vert \Phi
_{\limfunc{tr}}\left( I_{2}\right) \right\rangle _{\limfunc{tr}} &=&\frac{1}{%
2}\func{Tr}\left\{ I_{2}\right\} =1 \\
\left\langle \left. \Phi _{\limfunc{tr}}\left( \widetilde{x}_{1}\right)
\right\vert \Phi _{\limfunc{tr}}\left( \widetilde{x}_{1}\right)
\right\rangle _{\limfunc{tr}} &=&\frac{1}{2}\func{Tr}\left\{ \widetilde{x}%
_{1}^{2}\right\} =\frac{1}{2}\func{Tr}\left\{ I_{2}\right\} =1 \\
\left\langle \left. \Phi _{\limfunc{tr}}\left( \widetilde{x}_{2}\right)
\right\vert \Phi _{\limfunc{tr}}\left( \widetilde{x}_{2}\right)
\right\rangle _{\limfunc{tr}} &=&\frac{1}{2}\func{Tr}\left\{ \widetilde{x}%
_{2}^{2}\right\} =\frac{1}{2}\func{Tr}\left\{ I_{2}\right\} =1 \\
\left\langle \left. \Phi _{\limfunc{tr}}\left( \widetilde{x}_{1}\widetilde{x}%
_{2}\right) \right\vert \Phi _{\limfunc{tr}}\left( \widetilde{x}_{1}%
\widetilde{x}_{2}\right) \right\rangle _{\limfunc{tr}} &=&\frac{1}{2}\func{Tr%
}\left\{ \widetilde{x}_{2}\widetilde{x}_{1}\widetilde{x}_{1}\widetilde{x}%
_{2}\right\} =\frac{1}{2}\func{Tr}\left\{ I_{2}\right\} =1
\end{eqnarray}%
\begin{eqnarray}
\frac{1}{2}\left\langle \left. \Phi _{\limfunc{tr}}\left( I_{2}-\widetilde{x}%
_{2}\right) \right\vert \Phi _{\limfunc{tr}}\left( I_{2}-\widetilde{x}%
_{2}\right) \right\rangle _{D_{1}} &=&\frac{1}{2}\func{Tr}\left\{ \left(
I_{2}-\widetilde{x}_{2}\right) \left( I_{2}-\widetilde{x}_{2}\right)
\right\} =1 \\
\frac{1}{2}\left\langle \left. \Phi _{\limfunc{tr}}\left( \widetilde{x}_{1}-%
\widetilde{x}_{1}\widetilde{x}_{2}\right) \right\vert \Phi _{\limfunc{tr}%
}\left( \widetilde{x}_{1}-\widetilde{x}_{1}\widetilde{x}_{2}\right)
\right\rangle _{D_{1}} &=&\frac{1}{2}\func{Tr}\left\{ \left( \widetilde{x}%
_{1}-\widetilde{x}_{1}\widetilde{x}_{2}\right) ^{\dagger }\left( \widetilde{x%
}_{1}-\widetilde{x}_{1}\widetilde{x}_{2}\right) \right\} =1 \\
\frac{1}{2}\left\langle \left. \Phi _{\limfunc{tr}}\left( I_{2}+\widetilde{x}%
_{2}\right) \right\vert \Phi _{\limfunc{tr}}\left( I_{2}+\widetilde{x}%
_{2}\right) \right\rangle _{D_{1}} &=&\frac{1}{2}\func{Tr}\left\{ \left(
I_{2}+\widetilde{x}_{2}\right) \left( I_{2}+\widetilde{x}_{2}\right)
\right\} =1 \\
\frac{1}{2}\left\langle \left. \Phi _{\limfunc{tr}}\left( \widetilde{x}_{1}+%
\widetilde{x}_{1}\widetilde{x}_{2}\right) \right\vert \Phi _{\limfunc{tr}%
}\left( \widetilde{x}_{1}+\widetilde{x}_{1}\widetilde{x}_{2}\right)
\right\rangle _{D_{1}} &=&\frac{1}{2}\func{Tr}\left\{ \left( \widetilde{x}%
_{1}+\widetilde{x}_{1}\widetilde{x}_{2}\right) ^{\dagger }\left( \widetilde{x%
}_{1}+\widetilde{x}_{1}\widetilde{x}_{2}\right) \right\} =1,
\end{eqnarray}%
orthogonal 
\begin{eqnarray}
\left\langle \left. \Phi _{\limfunc{tr}}\left( I_{2}\right) \right\vert \Phi
_{\limfunc{tr}}\left( \widetilde{x}_{1}\right) \right\rangle _{\limfunc{tr}}
&=&\frac{1}{2}\func{Tr}\left\{ \widetilde{x}_{1}\right\} =\left\langle
\left. \Phi _{\limfunc{tr}}\left( I_{2}\right) \right\vert \Phi _{\limfunc{tr%
}}\left( \widetilde{x}_{2}\right) \right\rangle _{\limfunc{tr}}=\frac{1}{2}%
\func{Tr}\left\{ \widetilde{x}_{2}\right\} \\
&=&\left\langle \left. \Phi _{\limfunc{tr}}\left( I_{2}\right) \right\vert
\Phi _{\limfunc{tr}}\left( \widetilde{x}_{1}\widetilde{x}_{2}\right)
\right\rangle _{\limfunc{tr}}=\frac{1}{2}\func{Tr}\left\{ \widetilde{x}_{1}%
\widetilde{x}_{2}\right\} =0 \\
\left\langle \left. \Phi _{\limfunc{tr}}\left( \widetilde{x}_{1}\right)
\right\vert \Phi _{\limfunc{tr}}\left( \widetilde{x}_{2}\right)
\right\rangle _{\limfunc{tr}} &=&\frac{1}{2}\func{Tr}\left\{ \widetilde{x}%
_{1}\widetilde{x}_{2}\right\} =\left\langle \left. \Phi _{\limfunc{tr}%
}\left( \widetilde{x}_{1}\right) \right\vert \Phi _{\limfunc{tr}}\left( 
\widetilde{x}_{1}\widetilde{x}_{2}\right) \right\rangle _{\limfunc{tr}}=%
\frac{1}{2}\func{Tr}\left\{ \widetilde{x}_{2}\right\} =0 \\
\left\langle \left. \Phi _{\limfunc{tr}}\left( \widetilde{x}_{2}\right)
\right\vert \Phi _{\limfunc{tr}}\left( \widetilde{x}_{1}\widetilde{x}%
_{2}\right) \right\rangle _{\limfunc{tr}} &=&-\frac{1}{2}\func{Tr}\left\{ 
\widetilde{x}_{1}\right\} =0
\end{eqnarray}%
\begin{eqnarray}
\left\langle \left. \Phi _{\limfunc{tr}}\left( I_{2}-\widetilde{x}%
_{2}\right) \right\vert \Phi _{\limfunc{tr}}\left( \widetilde{x}_{1}-%
\widetilde{x}_{1}\widetilde{x}_{2}\right) \right\rangle _{\limfunc{tr}} &=&%
\frac{1}{2}\func{Tr}\left\{ \left( I_{2}-\widetilde{x}_{2}\right) \left( 
\widetilde{x}_{1}-\widetilde{x}_{1}\widetilde{x}_{2}\right) \right\} =0 \\
\left\langle \left. \Phi _{\limfunc{tr}}\left( I_{2}-\widetilde{x}%
_{2}\right) \right\vert \Phi _{\limfunc{tr}}\left( I_{2}+\widetilde{x}%
_{2}\right) \right\rangle _{\limfunc{tr}} &=&\frac{1}{2}\func{Tr}\left\{
\left( I_{2}-\widetilde{x}_{2}\right) \left( I_{2}+\widetilde{x}_{2}\right)
\right\} =0 \\
\left\langle \left. \Phi _{\limfunc{tr}}\left( I_{2}-\widetilde{x}%
_{2}\right) \right\vert \Phi _{\limfunc{tr}}\left( I_{2}+\widetilde{x}%
_{2}\right) \right\rangle _{\limfunc{tr}} &=&\frac{1}{2}\func{Tr}\left\{
\left( I_{2}-\widetilde{x}_{2}\right) \left( I_{2}+\widetilde{x}_{2}\right)
\right\} =0 \\
\left\langle \left. \Phi _{\limfunc{tr}}\left( \widetilde{x}_{1}-\widetilde{x%
}_{1}\widetilde{x}_{2}\right) \right\vert \Phi _{\limfunc{tr}}\left( I_{2}+%
\widetilde{x}_{2}\right) \right\rangle _{\limfunc{tr}} &=&\frac{1}{2}\func{Tr%
}\left\{ \left( \widetilde{x}_{1}-\widetilde{x}_{1}\widetilde{x}_{2}\right)
^{t}\left( I_{2}+\widetilde{x}_{2}\right) \right\} =0 \\
\left\langle \left. \Phi _{\limfunc{tr}}\left( \widetilde{x}_{1}-\widetilde{x%
}_{1}\widetilde{x}_{2}\right) \right\vert \Phi _{\limfunc{tr}}\left( 
\widetilde{x}_{1}+\widetilde{x}_{1}\widetilde{x}_{2}\right) \right\rangle _{%
\limfunc{tr}} &=&\frac{1}{2}\func{Tr}\left\{ \left( \widetilde{x}_{1}-%
\widetilde{x}_{1}\widetilde{x}_{2}\right) ^{t}\left( \widetilde{x}_{1}+%
\widetilde{x}_{1}\widetilde{x}_{2}\right) \right\} =0 \\
\left\langle \left. \Phi _{\limfunc{tr}}\left( I_{2}+\widetilde{x}%
_{2}\right) \right\vert \Phi _{\limfunc{tr}}\left( \widetilde{x}_{1}+%
\widetilde{x}_{1}\widetilde{x}_{2}\right) \right\rangle _{\limfunc{tr}} &=&%
\frac{1}{2}\func{Tr}\left\{ \left( I_{2}-\widetilde{x}_{2}\right) \left( 
\widetilde{x}_{1}+\widetilde{x}_{1}\widetilde{x}_{2}\right) \right\} =0
\end{eqnarray}%
and

\item are related to each other by the transformation 
\begin{eqnarray}
&&\frac{1}{\sqrt{2}}\left( \Phi _{\limfunc{tr}}\left( I_{2}+\widetilde{x}%
_{2}\right) ,\Phi _{\limfunc{tr}}\left( \widetilde{x}_{1}+\widetilde{x}_{1}%
\widetilde{x}_{2}\right) ,\Phi _{\limfunc{tr}}\left( I_{2}-\widetilde{x}%
_{2}\right) ,\Phi _{\limfunc{tr}}\left( \widetilde{x}_{1}-\widetilde{x}_{1}%
\widetilde{x}_{2}\right) \right) \\
&=&\left( \Phi _{\limfunc{tr}}\left( I_{2}\right) ,\Phi _{\limfunc{tr}%
}\left( \widetilde{x}_{1}\right) ,\Phi _{\limfunc{tr}}\left( \widetilde{x}%
_{2}\right) ,\Phi _{\limfunc{tr}}\left( \widetilde{x}_{1}\widetilde{x}%
_{2}\right) \right) \frac{1}{\sqrt{2}}\left( 
\begin{array}{cccc}
1 & 0 & 1 & 0 \\ 
0 & 1 & 0 & 1 \\ 
1 & 0 & -1 & 0 \\ 
0 & 1 & 0 & -1%
\end{array}%
\right) \\
&=&\left( \Phi _{\limfunc{tr}}\left( I_{2}\right) ,\Phi _{\limfunc{tr}%
}\left( \widetilde{x}_{1}\right) ,\Phi _{\limfunc{tr}}\left( \widetilde{x}%
_{2}\right) ,\Phi _{\limfunc{tr}}\left( \widetilde{x}_{1}\widetilde{x}%
_{2}\right) \right) \frac{1}{\sqrt{2}}\left( 
\begin{array}{cc}
\mathbb{I}_{2} & \mathbb{I}_{2} \\ 
\mathbb{I}_{2} & -\mathbb{I}_{2}%
\end{array}%
\right) .
\end{eqnarray}
\end{enumerate}

From the preceding one can see that the representation $\Pi _{\limfunc{tr}}$
can be expressed as the direct sum of the two irreducible representations $%
\Pi _{D_{1}}$ and $\Pi _{D_{2}}$ i.e. 
\begin{equation}
\Pi _{\limfunc{tr}}\left( \mathcal{F}\right) =\Pi _{D_{1}}\left( \mathcal{F}%
\right) \oplus \Pi _{D_{2}}\left( \mathcal{F}\right)
\end{equation}%
acting on 
\begin{equation}
\Phi _{\limfunc{tr}}\left( \mathcal{F}\right) =\Phi _{D_{1}}\left( \mathcal{F%
}\right) \oplus \Phi _{D_{2}}\left( \mathcal{F}\right) ,
\end{equation}%
that the representations $\Pi _{\limfunc{tr}}$ and $\Pi _{L}$ are the same
i.e 
\begin{equation}
\Pi _{\limfunc{tr}}\equiv \Pi _{L}
\end{equation}%
and 
\begin{equation}
\Pi \left( \mathcal{F}\right) \equiv \Pi _{D_{1}}\equiv \Pi _{D_{2}.}
\end{equation}

The inner product in $\Phi _{D_{1}}\left( \mathcal{F}\right) $ and $\Phi
_{D_{2}}\left( \mathcal{F}\right) $ are given by 
\begin{eqnarray}
\left\langle \left. \Phi _{D}\left( A\right) \right\vert \Phi _{D}\left(
B\right) \right\rangle _{D} &=&\limfunc{tr}\left\{ A^{\dagger }BD\right\} =%
\func{Tr}\left\{ \left\vert \Phi _{D}\left( B\right) \right\rangle
\left\langle \Phi _{D}\left( A\right) \right\vert \mathcal{\Pi }_{D}\left(
D\right) \right\} =2\limfunc{tr}\left\{ \left\vert \Phi _{D}\left( B\right)
\right\rangle \left\langle \Phi _{D}\left( A\right) \right\vert \mathcal{\Pi 
}_{D}\left( D\right) \right\} \\
&=&2\limfunc{tr}\left\{ \Pi _{D}\left( A\right) ^{\dagger }\Pi _{D}\left(
B\right) \mathcal{\Pi }_{D}\left( D\right) \right\} ,
\end{eqnarray}%
where $D=D_{1},D_{2}.$

The set of elements 
\begin{equation}
\mathcal{F}_{\mathbb{C}}=\ker \left\{ \widehat{i_{x}}\right\} =\left\{
A\left. \widehat{i_{x}}\left( A\right) =0\right\vert A\in \mathcal{F}%
\right\} ,
\end{equation}%
where%
\begin{equation}
\widehat{i_{x}}\left( A\right) =\left[ i_{x},A\right]
\end{equation}%
is a subalgebra of $\mathcal{F}$. It is in fact the \emph{even subalgebra} $%
\mathcal{F}_{even}$ i.e. $\mathcal{F}_{\mathbb{C}}=\mathcal{F}_{even},$
where 
\begin{equation}
\mathcal{F}_{even}=\left\{ \left. \alpha _{0}I_{2}+\alpha _{3}\widetilde{x}%
_{1}\widetilde{x}_{2}\right\vert \alpha _{0},\alpha _{3}\in \mathbb{R}%
\right\} =\left\{ \left. \alpha _{0}I_{2}+\alpha _{3}i_{\widetilde{x}%
}\right\vert \alpha _{0},\alpha _{3}\in \mathbb{R}\right\} ,
\end{equation}%
which is closed under Clifford multiplication$.$ Clearly 
\begin{equation}
\mathbb{R}\text{-}\dim \left\{ \mathcal{F}_{even}\right\} =2.
\end{equation}%
The orthogonal compliment of $\mathcal{F}_{even}$ is the odd space $\mathcal{%
F}_{odd},$ which is given by 
\begin{equation}
\mathcal{F}_{odd}=\left\{ \left. \alpha _{11}\widetilde{x}_{1}+\alpha _{12}%
\widetilde{x}_{2}\right\vert \alpha _{11},\alpha _{12}\in \mathbb{R}\right\}
.
\end{equation}

\begin{theorem}
$\mathcal{F}_{even}$ is a representation of the \emph{complex numbers, }$%
\mathbb{C},$ when%
\begin{equation}
\mathbb{R}\text{-}\dim \left\{ \mathcal{F}\right\} =4,
\end{equation}%
$\mathcal{F}$ is the \emph{real} \emph{Clifford algebra} $\equiv Cl\left(
2,0,V\right) ,$ where $\mathbb{R-}\dim \left\{ V\right\} =2$ and 
\begin{equation}
\mathcal{F}=\mathcal{F}_{even}\oplus _{\mathbb{R}}\mathcal{F}_{odd}=\mathcal{%
F}_{0}\oplus _{\mathbb{R}}\mathcal{F}_{1}\oplus _{\mathbb{R}}\mathcal{F}_{2}.
\end{equation}

\begin{corollary}
Every complex number is a sum of a real scalar part and a real bivector part.
\end{corollary}
\end{theorem}

\begin{theorem}
The even subalgebra, $\mathcal{F}_{even}=Cl^{0}\left( 2,0,V\right) ,$ where $%
\mathbb{R-}\dim \left\{ V\right\} =2$ is isomorphic to the Clifford algebra $%
Cl\left( 0,1,V\right) ,$where $\mathbb{R-}\dim \left\{ V\right\} =1.$
\end{theorem}

\begin{proof}
\begin{equation}
Cl\left( 0,1,V\right) =\tbigwedge V=\tbigwedge\nolimits^{0}V\oplus _{\mathbb{%
R}}\tbigwedge\nolimits^{1}V
\end{equation}%
and is generated by $\left\{ I_{1},e\right\} $ where 
\begin{eqnarray}
e^{2} &=&-1I_{1} \\
e\wedge e &=&0.
\end{eqnarray}%
The even subalgebra 
\begin{equation}
Cl^{0}\left( 2,0,V\right) =\tbigwedge\nolimits^{0}V\oplus _{\mathbb{R}%
}\tbigwedge\nolimits^{2}V
\end{equation}%
is generated by $\left\{ I_{2},\left( \widetilde{x}_{1}\widetilde{x}%
_{2}\right) \right\} ,$ where 
\begin{eqnarray}
\left( \widetilde{x}_{1}\widetilde{x}_{2}\right) ^{2} &=&-1I_{2} \\
\left( \widetilde{x}_{1}\widetilde{x}_{2}\right) \wedge \left( \widetilde{x}%
_{1}\widetilde{x}_{2}\right) &=&0,
\end{eqnarray}%
as 
\begin{equation}
\left( \widetilde{x}_{1}\widetilde{x}_{2}\right) \wedge \left( \widetilde{x}%
_{1}\widetilde{x}_{2}\right) =\widetilde{x}_{1}\wedge \widetilde{x}%
_{2}\wedge \widetilde{x}_{1}\wedge \widetilde{x}_{2}=0.
\end{equation}
\end{proof}

The complex number $z\in \mathcal{F}_{even}$ has the vector space
representation, $\mathcal{R}_{x}\left( z\right) $ of 
\begin{equation}
\mathcal{F}_{even}=\mathcal{F}_{0}\oplus _{\mathbb{R}}\mathcal{F}_{2}
\end{equation}%
wrt the orthonormal basis $\left\{ I_{2},\widetilde{x}_{1}\widetilde{x}%
_{2}\right\} ,$ given by 
\begin{equation}
z\left( I_{2},\widetilde{x}_{1}\widetilde{x}_{2}\right) =\left( I_{2},%
\widetilde{x}_{1}\widetilde{x}_{2}\right) \mathcal{R}_{\widetilde{x}}\left(
z\right) =\left( \func{Re}z\right) I_{2}+\left( \func{Im}z\right) \widetilde{%
x}_{1}\widetilde{x}_{2}=\left( I_{2},\widetilde{x}_{1}\widetilde{x}%
_{2}\right) \left( 
\begin{array}{c}
\func{Re}z \\ 
\func{Im}z%
\end{array}%
\right) .
\end{equation}%
$\lceil $%
\begin{gather}
\limfunc{tr}\left\{ I_{2}^{\dagger }I_{2}\right\} =\limfunc{tr}\left\{
\left( \widetilde{x}_{1}\widetilde{x}_{2}\right) ^{\dagger }\left( 
\widetilde{x}_{1}\widetilde{x}_{2}\right) \right\} =1 \\
\limfunc{tr}\left\{ I_{2}^{\dagger }\left( \widetilde{x}_{1}\widetilde{x}%
_{2}\right) \right\} =0
\end{gather}%
$\rfloor $

The representation in $M\left( 2,\mathbb{R}\right) $ is given by 
\begin{equation}
\mathcal{R}\left( z\right) =\left( \func{Re}z\right) \left( 
\begin{array}{cc}
1 & 0 \\ 
0 & 1%
\end{array}%
\right) +\left( \func{Im}z\right) \left( 
\begin{array}{cc}
0 & 1 \\ 
-1 & 0%
\end{array}%
\right) =\left( 
\begin{array}{cc}
\func{Re}z & \func{Im}z \\ 
-\func{Im}z & \func{Re}z%
\end{array}%
\right) .
\end{equation}

One can form the non self adjoint basis of the complexification $\mathcal{F}%
_{1\mathbb{C}}$ of $\mathcal{F}_{1}$ by the transformation 
\begin{equation}
\left( a^{\dagger },a\right) =\frac{1}{\sqrt{2}}\left( \widetilde{x}_{1},%
\widetilde{x}_{2}\right) \frac{1}{\sqrt{2}}\left( 
\begin{array}{cc}
1 & 1 \\ 
i & -i%
\end{array}%
\right) =\left( \frac{1}{2}\left( \widetilde{x}_{1}-i\widetilde{x}%
_{2}\right) ,\frac{1}{2}\left( \widetilde{x}_{1}+i\widetilde{x}_{2}\right)
\right) .
\end{equation}%
$\lceil $%
\begin{eqnarray}
x_{1}^{2} &=&\frac{1}{2}I_{2} \\
\widetilde{x}_{1} &=&\sqrt{2}x_{1} \\
\frac{1}{\sqrt{2}}\widetilde{x}_{1} &=&x_{1}
\end{eqnarray}%
\begin{gather}
\limfunc{tr}\left\{ aa^{\dagger }\right\} =\frac{1}{2}\limfunc{tr}\left\{
\left( x_{1}-ix_{2}\right) \left( x_{1}+ix_{2}\right) \right\} =\frac{1}{2}%
\limfunc{tr}\left\{ \left(
x_{1}^{2}-ix_{2}x_{1}+ix_{1}x_{2}+x_{2}^{2}\right) \right\} =\frac{1}{2}%
\limfunc{tr}\left\{ I_{2}-2ix_{2}x_{1}\right\} =\frac{1}{2} \\
\limfunc{tr}\left\{ I_{2}^{t}I_{2}\right\} =1 \\
\limfunc{tr}\left\{ I_{2}^{\dagger }a^{\dagger }\right\} =0
\end{gather}%
in terms of normalized basis elements 
\begin{gather}
\limfunc{tr}\left\{ aa^{\dagger }\right\} =\limfunc{tr}\left\{ \frac{1}{2}%
\left( \widetilde{x}_{1}-i\widetilde{x}_{2}\right) \frac{1}{2}\left( 
\widetilde{x}_{1}+i\widetilde{x}_{2}\right) \right\} =\frac{1}{4}\limfunc{tr}%
\left\{ \left( \widetilde{x}_{1}^{2}-i\widetilde{x}_{2}\widetilde{x}_{1}+i%
\widetilde{x}_{1}\widetilde{x}_{2}+\widetilde{x}_{2}^{2}\right) \right\} =%
\frac{1}{4}\limfunc{tr}\left\{ 2I_{2}-2i\widetilde{x}_{2}\widetilde{x}%
_{1}\right\} =\frac{1}{2} \\
\limfunc{tr}\left\{ I_{2}^{t}I_{2}\right\} =1 \\
\limfunc{tr}\left\{ I_{2}^{\dagger }a^{\dagger }\right\} =0
\end{gather}%
$\rfloor $

The complex Clifford algebra $\mathcal{F}_{\mathbb{C}}$ is given by 
\begin{equation}
\mathcal{F}_{\mathbb{C}}=\mathbb{C-}\limfunc{linspan}\left\{ I_{2},%
\widetilde{x}_{1},\widetilde{x}_{2},\widetilde{x}_{1}\widetilde{x}%
_{2}\right\} =\mathbb{C-}\limfunc{linspan}\left\{ I_{2},a^{\dagger
},a,a^{\dagger }a\right\} ,
\end{equation}%
one can note that real algebra $\mathcal{F}$ can also be generated by the
non self adjoint basis in the sense of an isomorphism i.e. 
\begin{equation}
\mathcal{F}=\mathbb{R-}\limfunc{linspan}\left\{ I_{2},\widetilde{x}_{1},%
\widetilde{x}_{2},\widetilde{x}_{1}\widetilde{x}_{2}\right\} \cong \mathbb{R-%
}\limfunc{linspan}\left\{ I_{2},a^{\dagger },a,a^{\dagger }a\right\} .
\end{equation}%
One can form the vacuum GNS\ representation $\Pi _{\phi }\left( \mathcal{F}_{%
\mathbb{C}}\right) $ on $\Phi _{\phi }\left( \mathcal{F}_{\mathbb{C}}\right) 
$ wrt the vacuum state $\left\vert \phi \right\rangle \left\langle \phi
\right\vert ,$ where 
\begin{equation}
\ker \left\{ \Phi _{\phi }\right\} =\mathbb{C-}\limfunc{linspan}\left\{
a,a^{\dagger }a\right\} =\left\{ \left. \limfunc{tr}\left\{ A^{\dagger
}A\right\} =0\right\vert A\in \mathcal{F}_{\mathbb{C}}\right\}
\end{equation}%
and 
\begin{equation}
\ker \left\{ \Phi _{\phi }\right\} ^{\bot }=\mathbb{C-}\limfunc{linspan}%
\left\{ a^{\dagger },aa^{\dagger }\right\} =\left\{ \left. \limfunc{tr}%
\left\{ A^{\dagger }A\right\} \neq 0\right\vert A\in \mathcal{F}_{\mathbb{C}%
}\right\} .
\end{equation}%
The representation space $\Phi _{\phi }\left( \mathcal{F}_{\mathbb{C}%
}\right) $ is given by 
\begin{equation}
\Phi _{\phi }\left( \mathcal{F}_{\mathbb{C}}\right) =\mathbb{C-}\limfunc{%
linspan}\left\{ \Phi _{\phi }\left( I_{2}\right) ,\Phi _{\phi }\left(
a^{\dagger }\right) \right\} .
\end{equation}%
As%
\begin{equation}
\left( a^{\dagger },a\right) =\left( \widetilde{x}_{1},\widetilde{x}%
_{2}\right) \frac{1}{2}\left( 
\begin{array}{cc}
1 & 1 \\ 
i & -i%
\end{array}%
\right) =\left( \frac{1}{2}\left( \widetilde{x}_{1}-i\widetilde{x}%
_{2}\right) ,\frac{1}{2}\left( \widetilde{x}_{1}+i\widetilde{x}_{2}\right)
\right)
\end{equation}%
one obtains that the basis of $\mathcal{F}_{1}$ can be expresed in terms of
the creator and annihilator as 
\begin{equation}
\left( \widetilde{x}_{1},\widetilde{x}_{2}\right) =\left( a^{\dagger
},a\right) \left( 
\begin{array}{cc}
1 & -i \\ 
1 & i%
\end{array}%
\right) =\left( a^{\dagger }+a,i\left( a-a^{\dagger }\right) \right) ,
\end{equation}%
and the basis of $\mathcal{F}_{2}$ as 
\begin{equation}
\widetilde{x}_{1}\widetilde{x}_{2}=\left( a^{\dagger }+a\right) i\left(
a-a^{\dagger }\right) =i\left( a^{\dagger }a-aa^{\dagger }\right) =i\left(
I_{2}-2aa^{\dagger }\right) ,
\end{equation}%
which gives 
\begin{eqnarray}
aa^{\dagger } &=&\frac{1}{2}\left( I_{2}+i\widetilde{x}_{1}\widetilde{x}%
_{2}\right) =h \\
a^{\dagger }a &=&\frac{1}{2}\left( I_{2}-i\widetilde{x}_{1}\widetilde{x}%
_{2}\right) =n
\end{eqnarray}%
$\lceil $%
\begin{equation}
\frac{1}{2}\left( I_{2}+i\widetilde{x}_{1}\widetilde{x}_{2}\right) \frac{1}{2%
}\left( I_{2}+i\widetilde{x}_{1}\widetilde{x}_{2}\right) =\frac{1}{2}\left(
I_{2}+2i\widetilde{x}_{1}\widetilde{x}_{2}+I_{2}\right) =\frac{1}{2}\left(
I_{2}+i\widetilde{x}_{1}\widetilde{x}_{2}\right)
\end{equation}%
$\rfloor $

$\lceil $%
\begin{eqnarray}
\frac{1}{2}\left( 
\begin{array}{cc}
1 & 1 \\ 
i & -i%
\end{array}%
\right) \left( 
\begin{array}{cc}
1 & -i \\ 
1 & i%
\end{array}%
\right) &=&\frac{1}{2}\left( 
\begin{array}{cc}
2 & 0 \\ 
0 & 2%
\end{array}%
\right) =\left( 
\begin{array}{cc}
1 & 0 \\ 
0 & 1%
\end{array}%
\right) \\
\left( 
\begin{array}{cc}
1 & -i \\ 
1 & i%
\end{array}%
\right) &=&\left( \frac{1}{2}\left( 
\begin{array}{cc}
1 & 1 \\ 
i & -i%
\end{array}%
\right) \right) ^{-1}
\end{eqnarray}%
$\rfloor $

and the creator and annihilator basis $\left\{ a^{\dagger },a\right\} $ can
be expresed in terms of the basis $\left\{ \widetilde{x}_{1},\widetilde{x}%
_{2}\right\} $ as 
\begin{equation}
\left( a^{\dagger },a\right) =\left( \widetilde{x}_{1},\widetilde{x}%
_{2}\right) \frac{1}{2}\left( 
\begin{array}{cc}
1 & 1 \\ 
-i & i%
\end{array}%
\right) =\frac{1}{2}\left( \widetilde{x}_{1}-i\widetilde{x}_{2},\widetilde{x}%
_{1}-i\widetilde{x}_{2}\right) .
\end{equation}%
Using this one obtains that 
\begin{equation}
\ker \left\{ \Phi _{\phi }\right\} ^{\bot }=\mathbb{C-}\limfunc{linspan}%
\left\{ a^{\dagger },aa^{\dagger }\right\}
\end{equation}%
which can be expressed as 
\begin{equation}
\ker \left\{ \Phi _{\phi }\right\} ^{\bot }=\mathbb{C-}\limfunc{linspan}%
\left\{ \ \left( \widetilde{x}_{1}-i\widetilde{x}_{2}\right) ,\left( I_{2}-i%
\widetilde{x}_{1}\widetilde{x}_{2}\right) \right\} .
\end{equation}%
From these relationships one observes 
\begin{eqnarray}
\Phi _{\phi }\left( \widetilde{x}_{1}\right) &=&\Phi _{\phi }\left(
a^{\dagger }\right) \\
\Phi _{\phi }\left( \widetilde{x}_{2}\right) &=&-i\Phi _{\phi }\left(
a^{\dagger }\right) =-i\Phi _{\phi }\left( \widetilde{x}_{1}\right)
\Leftrightarrow \Phi _{\phi }\left( a^{\dagger }\right) =i\Phi _{\phi
}\left( \widetilde{x}_{2}\right) =\Phi _{\phi }\left( \widetilde{x}%
_{1}\right) \\
\Phi _{\phi }\left( \widetilde{x}_{1}\widetilde{x}_{2}\right) &=&i\Phi
_{\phi }\left( I_{2}-2a^{\dagger }a\right) =i\Phi _{\phi }\left(
I_{2}\right) -2i\Phi _{\phi }\left( a^{\dagger }a\right) =i\Phi _{\phi
}\left( I_{2}\right) .
\end{eqnarray}%
$\lceil $%
\begin{gather}
\Phi _{\phi }\left( \widetilde{x}_{1}\right) =\Phi _{\phi }\left( a^{\dagger
}+a\right) =\Phi _{\phi }\left( a^{\dagger }\right) \\
\Phi _{\phi }\left( \widetilde{x}_{2}\right) =i\Phi _{\phi }\left(
a^{\dagger }-a\right) =i\Phi _{\phi }\left( a^{\dagger }\right) \\
\Phi _{\phi }\left( \widetilde{x}_{1}\right) =i\Phi _{\phi }\left( 
\widetilde{x}_{2}\right) =\Phi _{\phi }\left( a^{\dagger }\right)
\end{gather}%
\begin{equation}
\Phi _{\phi }\left( \widetilde{x}_{1}\widetilde{x}_{2}\right) =i\Phi _{\phi
}\left( \left( a^{\dagger }+a\right) \left( a^{\dagger }-a\right) \right)
=i\Phi _{\phi }\left( a^{\dagger }a-aa^{\dagger }\right) =i\Phi _{\phi
}\left( I_{2}-2a^{\dagger }a\right)
\end{equation}%
\begin{gather}
i\widetilde{x}_{1}\widetilde{x}_{2}=-\left( I_{2}-2a^{\dagger }a\right)
=2a^{\dagger }a-I_{2} \\
\frac{1}{2}\left( I_{2}+i\widetilde{x}_{1}\widetilde{x}_{2}\right)
=a^{\dagger }a=n
\end{gather}%
$\rfloor $

\emph{one can see that}%
\begin{eqnarray}
\Phi _{\phi }\left( \mathcal{F}\right) &=&\mathbb{R-}\limfunc{linspan}%
\left\{ \Phi _{\phi }\left( I_{2}\right) ,\Phi _{\phi }\left( \widetilde{x}%
_{1}\right) ,\Phi _{\phi }\left( \widetilde{x}_{2}\right) ,\Phi _{\phi
}\left( \widetilde{x}_{1}\widetilde{x}_{2}\right) \right\} \\
&=&\mathbb{R-}\limfunc{linspan}\left\{ \Phi _{\phi }\left( I_{2}\right)
,\Phi _{\phi }\left( \widetilde{x}_{1}\widetilde{x}_{2}\right) ,\Phi _{\phi
}\left( \widetilde{x}_{1}\right) ,\Phi _{\phi }\left( \widetilde{x}%
_{2}\right) \right\} \\
&=&\mathbb{R-}\limfunc{linspan}\left\{ \Phi _{\phi }\left( I_{2}\right)
,\Phi _{\phi }\left( \widetilde{x}_{1}\right) ,i\Phi _{\phi }\left( 
\widetilde{x}_{1}\right) ,i\Phi _{\phi }\left( I_{2}\right) \right\} .
\end{eqnarray}
The representation $\Pi _{\phi }$ is given by

\begin{enumerate}
\item 
\begin{eqnarray}
\Pi _{\phi }\left( I\right) \Phi _{\phi }\left( I_{2}\right) &=&\Phi _{\phi
}\left( I_{2}\right) \\
\Pi _{\phi }\left( I\right) \Phi _{\phi }\left( a^{\dagger }\right) &=&\Phi
_{\phi }\left( a^{\dagger }\right)
\end{eqnarray}

\item 
\begin{eqnarray}
\Pi _{\phi }\left( \widetilde{x}_{1}\widetilde{x}_{2}\right) \Phi _{\phi
}\left( I_{2}\right) &=&\Phi _{\phi }\left( \widetilde{x}_{1}\widetilde{x}%
_{2}\right) \\
\Pi _{\phi }\left( \widetilde{x}_{1}\widetilde{x}_{2}\right) \Phi _{\phi
}\left( a^{\dagger }\right) &=&\Pi _{\phi }\left( \widetilde{x}_{1}%
\widetilde{x}_{2}\right) \Phi _{\phi }\left( \widetilde{x}_{1}\right) =-\Phi
_{\phi }\left( \widetilde{x}_{2}\right) =i\Phi _{\phi }\left( \widetilde{x}%
_{1}\right) =i\Phi _{\phi }\left( a^{\dagger }\right)
\end{eqnarray}

\item 
\begin{eqnarray}
\Pi _{\phi }\left( \widetilde{x}_{1}\right) \Phi _{\phi }\left( I_{2}\right)
&=&\Phi _{\phi }\left( \widetilde{x}_{1}\right) \\
\Pi _{\phi }\left( \widetilde{x}_{1}\right) \Phi _{\phi }\left( a^{\dagger
}\right) &=&\Pi _{\phi }\left( \widetilde{x}_{1}\right) \Phi _{\phi }\left( 
\widetilde{x}_{1}\right) =\Phi _{\phi }\left( I_{2}\right)
\end{eqnarray}

\item 
\begin{eqnarray}
\Pi _{\phi }\left( \widetilde{x}_{2}\right) \Phi _{\phi }\left( I_{2}\right)
&=&\Phi _{\phi }\left( \widetilde{x}_{2}\right) \\
\Pi _{\phi }\left( \widetilde{x}_{2}\right) \Phi _{\phi }\left( a^{\dagger
}\right) &=&\Pi _{\phi }\left( \widetilde{x}_{2}\right) \Phi _{\phi }\left( 
\widetilde{x}_{1}\right) =-\Phi _{\phi }\left( \widetilde{x}_{1}\widetilde{x}%
_{2}\right) =-i\Phi _{\phi }\left( I_{2}\right) .
\end{eqnarray}
\end{enumerate}

The even subspace is generated by 
\begin{equation}
\mathcal{F}_{even}=\mathbb{R-}\limfunc{linspan}\left\{ I_{2},\widetilde{x}%
_{1}\widetilde{x}_{2}\right\} =\mathbb{R-}\limfunc{linspan}\left\{
I_{2},i\left( a^{\dagger }a-aa^{\dagger }\right) \right\}
\end{equation}%
and the odd subspace by 
\begin{equation}
\mathcal{F}_{odd}=\mathbb{R-}\limfunc{linspan}\left\{ \widetilde{x}_{1},%
\widetilde{x}_{2}\right\} =\mathbb{R-}\limfunc{linspan}\left\{ \left(
a^{\dagger }+a\right) ,i\left( a-a^{\dagger }\right) \right\} .
\end{equation}%
Hence 
\begin{equation}
\Phi _{\phi }\left( \mathcal{F}_{even}\right) =\mathbb{R-}\limfunc{linspan}%
\left\{ \Phi _{\phi }\left( I_{2}\right) ,\Phi _{\phi }\left( \widetilde{x}%
_{1}\widetilde{x}_{2}\right) \right\} =\mathbb{R-}\limfunc{linspan}\left\{
\Phi _{\phi }\left( I_{2}\right) ,i\Phi _{\phi }\left( I_{2}\right) \right\}
\end{equation}%
and 
\begin{eqnarray}
\Phi _{\phi }\left( \mathcal{F}_{odd}\right) &=&\mathbb{R-}\limfunc{linspan}%
\left\{ \Phi _{\phi }\left( \widetilde{x}_{1}\right) ,\Phi _{\phi }\left( 
\widetilde{x}_{2}\right) \right\} =\mathbb{R-}\limfunc{linspan}\left\{ \Phi
_{\phi }\left( a^{\dagger }+a\right) ,i\Phi _{\phi }\left( a-a^{\dagger
}\right) \right\} \\
&=&\mathbb{R-}\limfunc{linspan}\left\{ \Phi _{\phi }\left( a^{\dagger
}\right) ,i\Phi _{\phi }\left( a^{\dagger }\right) \right\} =\mathbb{R-}%
\limfunc{linspan}\left\{ \Phi _{\phi }\left( \widetilde{x}_{1}\right) ,i\Phi
_{\phi }\left( \widetilde{x}_{1}\right) \right\} .
\end{eqnarray}

The complex numbers are generated by the restriction of $\Pi _{\phi }$ to $%
\mathcal{F}_{even}$ i.e. $\left. \Pi _{\phi }\right\vert _{\mathcal{F}%
_{even}}=\Pi _{\phi }\left( \mathcal{F}_{even}\right) .$ $\mathcal{F}_{even}$
is a closed Clifford sub algebra of $\mathcal{F}$ and its representation $%
\Pi _{\phi }\left( \mathcal{F}_{even}\right) $ leaves both $\Phi _{\phi
}\left( \mathcal{F}_{even}\right) $ and $\Phi _{\phi }\left( \mathcal{F}%
_{odd}\right) $ invariant i.e. 
\begin{eqnarray}
\Pi _{\phi }\left( \mathcal{F}_{even}\right) &:&\Phi _{\phi }\left( \mathcal{%
F}_{even}\right) \rightarrow \Phi _{\phi }\left( \mathcal{F}_{even}\right) \\
\Pi _{\phi }\left( \mathcal{F}_{even}\right) &:&\Phi _{\phi }\left( \mathcal{%
F}_{odd}\right) \rightarrow \Phi _{\phi }\left( \mathcal{F}_{odd}\right) .
\end{eqnarray}%
Thus 
\begin{equation}
\Phi _{\phi }\left( \mathcal{F}_{even}\right) \oplus \Phi _{\phi }\left( 
\mathcal{F}_{odd}\right)
\end{equation}%
reduces the representation $\Pi _{\phi }\left( \mathcal{F}_{even}\right) $
leading to 
\begin{eqnarray}
\Pi _{\phi }\left( I\right) \Phi _{\phi }\left( I_{2}\right) &=&\Phi _{\phi
}\left( I_{2}\right) \\
\Pi _{\phi }\left( I\right) \Phi _{\phi }\left( \widetilde{x}_{1}\widetilde{x%
}_{2}\right) &=&\Phi _{\phi }\left( \widetilde{x}_{1}\widetilde{x}_{2}\right)
\\
\Pi _{\phi }\left( \widetilde{x}_{1}\widetilde{x}_{2}\right) \Phi _{\phi
}\left( I_{2}\right) &=&\Phi _{\phi }\left( \widetilde{x}_{1}\widetilde{x}%
_{2}\right) \\
\Pi _{\phi }\left( \widetilde{x}_{1}\widetilde{x}_{2}\right) \Phi _{\phi
}\left( \widetilde{x}_{1}\widetilde{x}_{2}\right) &=&-\Phi _{\phi }\left(
I_{2}\right)
\end{eqnarray}%
and 
\begin{eqnarray}
\Pi _{\phi }\left( I\right) \Phi _{\phi }\left( \widetilde{x}_{1}\right)
&=&\Phi _{\phi }\left( \widetilde{x}_{1}\right) \\
\Pi _{\phi }\left( I\right) \Phi _{\phi }\left( \widetilde{x}_{2}\right)
&=&\Phi _{\phi }\left( \widetilde{x}_{2}\right) =-i\Phi _{\phi }\left( 
\widetilde{x}_{1}\right) \\
\Pi _{\phi }\left( \widetilde{x}_{1}\widetilde{x}_{2}\right) \Phi _{\phi
}\left( \widetilde{x}_{1}\right) &=&-\Phi _{\phi }\left( \widetilde{x}%
_{2}\right) =i\Phi _{\phi }\left( \widetilde{x}_{1}\right) \\
\Pi _{\phi }\left( \widetilde{x}_{1}\widetilde{x}_{2}\right) \Phi _{\phi
}\left( \widetilde{x}_{2}\right) &=&\Phi _{\phi }\left( \widetilde{x}%
_{1}\right) .
\end{eqnarray}%
These relationships produce the representations $\Pi _{\phi }\left( \mathcal{%
F}_{even}\right) $ of the even subalgebra $\mathcal{F}_{even}$ on $\Phi
_{\phi }\left( \mathcal{F}_{even}\right) $ as 
\begin{eqnarray}
\Pi _{\phi }\left( I\right) \left( \Phi _{\phi }\left( I_{2}\right) ,\Phi
_{\phi }\left( \widetilde{x}_{1}\widetilde{x}_{2}\right) \right) &=&\left(
\Phi _{\phi }\left( I_{2}\right) ,\Phi _{\phi }\left( \widetilde{x}_{1}%
\widetilde{x}_{2}\right) \right) \left( 
\begin{array}{cc}
1 & 0 \\ 
0 & 1%
\end{array}%
\right) =\left( \Phi _{\phi }\left( I_{2}\right) ,i\Phi _{\phi }\left(
I_{2}\right) \right) \left( 
\begin{array}{cc}
1 & 0 \\ 
0 & 1%
\end{array}%
\right) \\
\Pi _{\phi }\left( \widetilde{x}_{1}\widetilde{x}_{2}\right) \left( \Phi
_{\phi }\left( I_{2}\right) ,\Phi _{\phi }\left( \widetilde{x}_{1}\widetilde{%
x}_{2}\right) \right) &=&\left( \Phi _{\phi }\left( I_{2}\right) ,\Phi
_{\phi }\left( \widetilde{x}_{1}\widetilde{x}_{2}\right) \right) \left( 
\begin{array}{cc}
0 & -1 \\ 
1 & 0%
\end{array}%
\right) =\left( \Phi _{\phi }\left( I_{2}\right) ,i\Phi _{\phi }\left(
I_{2}\right) \right) \left( 
\begin{array}{cc}
0 & -1 \\ 
1 & 0%
\end{array}%
\right)
\end{eqnarray}%
and the representations $\Pi _{\phi }\left( \mathcal{F}_{even}\right) $ on $%
\Phi _{\phi }\left( \mathcal{F}_{odd}\right) $ as%
\begin{eqnarray}
\Pi _{\phi }\left( I\right) \left( \Phi _{\phi }\left( \widetilde{x}%
_{1}\right) ,\Phi _{\phi }\left( \widetilde{x}_{2}\right) \right) &=&\left(
\Phi _{\phi }\left( \widetilde{x}_{1}\right) ,\Phi _{\phi }\left( \widetilde{%
x}_{2}\right) \right) \\
&=&\left( \Phi _{\phi }\left( \widetilde{x}_{1}\right) ,\Phi _{\phi }\left( 
\widetilde{x}_{2}\right) \right) \left( 
\begin{array}{cc}
1 & 0 \\ 
0 & 1%
\end{array}%
\right) =\left( \Phi _{\phi }\left( \widetilde{x}_{1}\right) ,-i\Phi _{\phi
}\left( \widetilde{x}_{1}\right) \right) \left( 
\begin{array}{cc}
1 & 0 \\ 
0 & 1%
\end{array}%
\right) \\
\Pi _{\phi }\left( \widetilde{x}_{1}\widetilde{x}_{2}\right) \left( \Phi
_{\phi }\left( \widetilde{x}_{1}\right) ,\Phi _{\phi }\left( \widetilde{x}%
_{2}\right) \right) &=&\left( -\Phi _{\phi }\left( \widetilde{x}_{2}\right)
,\Phi _{\phi }\left( \widetilde{x}_{1}\right) \right) \\
&=&\left( \Phi _{\phi }\left( \widetilde{x}_{1}\right) ,\Phi _{\phi }\left( 
\widetilde{x}_{2}\right) \right) \left( 
\begin{array}{cc}
0 & 1 \\ 
-1 & 0%
\end{array}%
\right) =\left( \Phi _{\phi }\left( \widetilde{x}_{1}\right) ,-i\Phi _{\phi
}\left( \widetilde{x}_{1}\right) \right) \left( 
\begin{array}{cc}
0 & 1 \\ 
-1 & 0%
\end{array}%
\right) .
\end{eqnarray}%
The representation $\Pi _{\phi }\left( \mathcal{F}_{odd}\right) $ of the
generators of $\mathcal{F}_{odd}$ on $\Phi _{\phi }\left( \mathcal{F}%
_{even}\right) $ is given by 
\begin{eqnarray}
\Pi _{\phi }\left( \widetilde{x}_{1}\right) \left( \Phi _{\phi }\left(
I_{2}\right) ,\Phi _{\phi }\left( \widetilde{x}_{1}\widetilde{x}_{2}\right)
\right) &=&\left( \Phi _{\phi }\left( \widetilde{x}_{1}\right) ,\Phi _{\phi
}\left( \widetilde{x}_{2}\right) \right) =\left( \Phi _{\phi }\left(
I_{2}\right) ,\Phi _{\phi }\left( \widetilde{x}_{1}\widetilde{x}_{2}\right)
\right) \left( 
\begin{array}{cc}
1 & 0 \\ 
0 & 1%
\end{array}%
\right) \\
\Pi _{\phi }\left( \widetilde{x}_{2}\right) \left( \Phi _{\phi }\left(
I_{2}\right) ,\Phi _{\phi }\left( \widetilde{x}_{1}\widetilde{x}_{2}\right)
\right) &=&\left( \Phi _{\phi }\left( \widetilde{x}_{2}\right) ,-\Phi _{\phi
}\left( \widetilde{x}_{1}\right) \right) =\left( \Phi _{\phi }\left(
I_{2}\right) ,\Phi _{\phi }\left( \widetilde{x}_{1}\widetilde{x}_{2}\right)
\right) \left( 
\begin{array}{cc}
0 & -1 \\ 
1 & 0%
\end{array}%
\right)
\end{eqnarray}%
and the representations $\Pi _{\phi }\left( \mathcal{F}_{odd}\right) $ on $%
\Phi _{\phi }\left( \mathcal{F}_{odd}\right) $ as%
\begin{eqnarray}
\Pi _{\phi }\left( \widetilde{x}_{1}\right) \left( \Phi _{\phi }\left( 
\widetilde{x}_{1}\right) ,\Phi _{\phi }\left( \widetilde{x}_{2}\right)
\right) &=&\left( \Phi _{\phi }\left( I_{2}\right) ,\Phi _{\phi }\left( 
\widetilde{x}_{1}\widetilde{x}_{2}\right) \right) \\
&=&\left( \Phi _{\phi }\left( I_{2}\right) ,\Phi _{\phi }\left( \widetilde{x}%
_{1}\widetilde{x}_{2}\right) \right) \left( 
\begin{array}{cc}
1 & 0 \\ 
0 & 1%
\end{array}%
\right) =\left( \Phi _{\phi }\left( I_{2}\right) ,i\Phi _{\phi }\left(
I_{2}\right) \right) \left( 
\begin{array}{cc}
1 & 0 \\ 
0 & 1%
\end{array}%
\right) \\
\Pi _{\phi }\left( \widetilde{x}_{2}\right) \left( \Phi _{\phi }\left( 
\widetilde{x}_{1}\right) ,\Phi _{\phi }\left( \widetilde{x}_{2}\right)
\right) &=&\left( -\Phi _{\phi }\left( \widetilde{x}_{1}\widetilde{x}%
_{2}\right) ,\Phi _{\phi }\left( I_{2}\right) \right) \\
&=&\left( \Phi _{\phi }\left( I_{2}\right) ,\Phi _{\phi }\left( \widetilde{x}%
_{1}\widetilde{x}_{2}\right) \right) \left( 
\begin{array}{cc}
0 & 1 \\ 
-1 & 0%
\end{array}%
\right) =\left( \Phi _{\phi }\left( I_{2}\right) ,i\Phi _{\phi }\left(
I_{2}\right) \right) \left( 
\begin{array}{cc}
0 & 1 \\ 
-1 & 0%
\end{array}%
\right) .
\end{eqnarray}

As $\mathcal{F}_{even}$ is a subalgebra one can consider the GNS\
representation, $\Pi _{\limfunc{tr}},$ induced by the faithful tracial state 
\begin{equation}
\limfunc{tr}:\mathcal{F}_{even}\rightarrow \mathbb{R}
\end{equation}%
on the representation space $\Phi _{\limfunc{tr}}\left( \mathcal{F}%
_{even}\right) .$ The representation is given by 
\begin{eqnarray}
\Pi _{\limfunc{tr}}\left( I_{2}\right) \left( \Phi _{\limfunc{tr}}\left(
I_{2}\right) ,\Phi _{\limfunc{tr}}\left( \widetilde{x}_{1}\widetilde{x}%
_{2}\right) \right) &=&\left( \Phi _{\limfunc{tr}}\left( I_{2}\right) ,\Phi
_{\limfunc{tr}}\left( \widetilde{x}_{1}\widetilde{x}_{2}\right) \right)
=\left( \Phi _{\limfunc{tr}}\left( I_{2}\right) ,\Phi _{\limfunc{tr}}\left( 
\widetilde{x}_{1}\widetilde{x}_{2}\right) \right) \left( 
\begin{array}{cc}
1 & 0 \\ 
0 & 1%
\end{array}%
\right) \\
\Pi _{\limfunc{tr}}\left( \widetilde{x}_{1}\widetilde{x}_{2}\right) \left(
\Phi _{\limfunc{tr}}\left( I_{2}\right) ,\Phi _{\limfunc{tr}}\left( 
\widetilde{x}_{1}\widetilde{x}_{2}\right) \right) &=&\left( \Phi _{\limfunc{%
tr}}\left( \widetilde{x}_{1}\widetilde{x}_{2}\right) ,-\Phi _{\limfunc{tr}%
}\left( I_{2}\right) \right) =\left( \Phi _{\limfunc{tr}}\left( I_{2}\right)
,\Phi _{\limfunc{tr}}\left( \widetilde{x}_{1}\widetilde{x}_{2}\right)
\right) \left( 
\begin{array}{cc}
0 & -1 \\ 
1 & 0%
\end{array}%
\right) ,
\end{eqnarray}

\begin{theorem}
$\mathcal{F}_{even}$ is a representation of the \emph{quaternions, }$\mathbb{%
H},$ when%
\begin{equation}
\mathbb{R}-\dim \left\{ \mathcal{F}\right\} =9
\end{equation}%
and $\mathcal{F}$ is the \emph{real} \emph{Clifford algebra} $\equiv
Cl\left( 0,3,\mathbb{R}\right) $ and 
\begin{equation}
\mathcal{F}=\mathcal{F}_{even}\oplus _{\mathbb{R}}\mathcal{F}_{odd}.
\end{equation}
\end{theorem}

In this case 
\begin{equation}
\mathbb{R}-\dim \left\{ \mathcal{F}_{even}\right\} =4,
\end{equation}%
$r=2$ and the \emph{normalized} basis of the \emph{real} Clifford algebra $%
\mathcal{F}$ is 
\begin{equation}
I_{\mathcal{F}},\widetilde{x}_{1},\widetilde{x}_{2},\widetilde{x}_{3},%
\widetilde{x}_{1}\widetilde{x}_{2},\widetilde{x}_{1}\widetilde{x}_{3},%
\widetilde{x}_{2}\widetilde{x}_{3},\widetilde{x}_{1}\widetilde{x}_{2}%
\widetilde{x}_{3}
\end{equation}%
and 
\begin{equation}
\mathbb{R}-\dim \left\{ \mathcal{F}\right\} =8.
\end{equation}%
The even subalgebra 
\begin{equation}
\mathcal{F}_{even}=\mathbb{R-}\limfunc{linspan}\left\{ I_{3},\widetilde{x}%
_{1}\widetilde{x}_{2},\widetilde{x}_{1}\widetilde{x}_{3},\widetilde{x}_{2}%
\widetilde{x}_{3}\right\}
\end{equation}%
is a sub Clifford algebra and 
\begin{eqnarray}
\mathcal{F}_{even} &\cong &Cl^{0}\left( 0,3,\mathbb{R}\right) =\wedge 
\mathcal{F}_{1}=\wedge ^{0}\mathcal{F}_{1}\oplus \wedge ^{2}\mathcal{F}_{1}=%
\mathcal{F}_{0}\oplus \mathcal{F}_{2} \\
\mathcal{F}_{1} &=&\mathbb{R-}\limfunc{linspan}\left\{ \widetilde{x}_{1},%
\widetilde{x}_{2},\widetilde{x}_{3}\right\} .
\end{eqnarray}

$\lceil $This shown by 
\begin{eqnarray}
\widetilde{x}_{1}\widetilde{x}_{2}\widetilde{x}_{1}\widetilde{x}_{3} &=&-%
\widetilde{x}_{2}\widetilde{x}_{3} \\
\widetilde{x}_{1}\widetilde{x}_{3}\widetilde{x}_{1}\widetilde{x}_{2} &=&%
\widetilde{x}_{2}\widetilde{x}_{3} \\
\left\{ \widetilde{x}_{1}\widetilde{x}_{3},\widetilde{x}_{1}\widetilde{x}%
_{2}\right\} &=&0 \\
\widetilde{x}_{1}\widetilde{x}_{2}\widetilde{x}_{2}\widetilde{x}_{3} &=&%
\widetilde{x}_{1}\widetilde{x}_{3} \\
\widetilde{x}_{2}\widetilde{x}_{3}\widetilde{x}_{1}\widetilde{x}_{2} &=&-%
\widetilde{x}_{1}\widetilde{x}_{3} \\
\left\{ \widetilde{x}_{1}\widetilde{x}_{2},\widetilde{x}_{2}\widetilde{x}%
_{3}\right\} &=&0 \\
\widetilde{x}_{1}\widetilde{x}_{3}\widetilde{x}_{2}\widetilde{x}_{3} &=&-%
\widetilde{x}_{1}\widetilde{x}_{2} \\
\widetilde{x}_{2}\widetilde{x}_{3}\widetilde{x}_{1}\widetilde{x}_{3} &=&%
\widetilde{x}_{1}\widetilde{x}_{2} \\
\left\{ \widetilde{x}_{1}\widetilde{x}_{3},\widetilde{x}_{2}\widetilde{x}%
_{3}\right\} &=&0
\end{eqnarray}%
\begin{eqnarray}
\widetilde{x}_{1}\widetilde{x}_{2}\widetilde{x}_{1}\widetilde{x}_{2}
&=&-I_{3} \\
\widetilde{x}_{1}\widetilde{x}_{3}\widetilde{x}_{1}\widetilde{x}_{3}
&=&-I_{3} \\
\widetilde{x}_{2}\widetilde{x}_{3}\widetilde{x}_{2}\widetilde{x}_{3}
&=&-I_{3}.
\end{eqnarray}%
$\rfloor $

$\lceil $ We can rewrite these products in $Cl^{0}\left( 0,3,\mathbb{R}%
\right) $ by 
\begin{eqnarray}
e_{1} &=&\widetilde{x}_{1}\widetilde{x}_{2} \\
e_{2} &=&\widetilde{x}_{1}\widetilde{x}_{3} \\
e_{3} &=&\widetilde{x}_{2}\widetilde{x}_{3} \\
e_{1}e_{2} &=&\widetilde{x}_{1}\widetilde{x}_{2}\widetilde{x}_{1}\widetilde{x%
}_{3}=-\widetilde{x}_{2}\widetilde{x}_{3}=-e_{3} \\
e_{1}e_{3} &=&\widetilde{x}_{1}\widetilde{x}_{2}\widetilde{x}_{2}\widetilde{x%
}_{3}=\widetilde{x}_{1}\widetilde{x}_{3}=e_{2} \\
e_{2}e_{3} &=&\widetilde{x}_{1}\widetilde{x}_{3}\widetilde{x}_{2}\widetilde{x%
}_{3}=-\widetilde{x}_{1}\widetilde{x}_{2}=-e_{1} \\
e_{1}e_{2}e_{3} &=&-\widetilde{x}_{2}\widetilde{x}_{3}\widetilde{x}_{2}%
\widetilde{x}_{3}=I_{3}.
\end{eqnarray}%
$\rfloor $

The Clifford subalgebra $Cl^{0}\left( 0,3,\mathbb{R}\right) $ is isomorphic
to the Clifford algebra $Cl\left( 0,2,\mathbb{R}\right) $, where 
\begin{eqnarray}
\mathcal{F} &\mathcal{=}&Cl\left( 0,2,\mathbb{R}\right) =\mathbb{R-}\limfunc{%
linspan}\left\{ I_{2},E_{1},E_{2},E_{1}E_{2}\right\} =\wedge \mathcal{F}%
_{1}=\wedge ^{0}\mathcal{F}_{1}\oplus \wedge ^{1}\mathcal{F}_{1}\oplus
\wedge ^{2}\mathcal{F}_{1} \\
\mathcal{F}_{1} &=&\mathbb{R-}\limfunc{linspan}\left\{ E_{1},E_{2}\right\}
\end{eqnarray}%
and the isomorphism 
\begin{equation}
T:Cl^{0}\left( 0,3,\mathbb{R}\right) \rightarrow Cl\left( 0,2,\mathbb{R}%
\right)
\end{equation}%
is given by 
\begin{eqnarray}
T &:&I_{3}\rightarrow I_{2} \\
T &:&\widetilde{x}_{1}\widetilde{x}_{2}\rightarrow E_{1} \\
T &:&\widetilde{x}_{1}\widetilde{x}_{3}\rightarrow -E_{2} \\
T &:&\widetilde{x}_{2}\widetilde{x}_{3}\rightarrow E_{1}E_{2}
\end{eqnarray}%
and noting 
\begin{equation}
\mathbb{R}-\dim \left\{ Cl^{0}\left( 0,3,\mathbb{R}\right) \right\} =\mathbb{%
R}-\dim \left\{ Cl\left( 0,2,\mathbb{R}\right) \right\} =2.
\end{equation}%
\emph{The isomorphism }$T$\emph{\ is a Clifford algebra isomorphism but not
grade preserving.}

$\lceil $ The linear properties are given by 
\begin{equation}
T\left( \alpha _{0}I_{3}+\alpha _{1}\widetilde{x}_{1}\widetilde{x}%
_{2}+\alpha _{2}\widetilde{x}_{1}\widetilde{x}_{3}+\alpha _{3}\widetilde{x}%
_{2}\widetilde{x}_{3}\right) =\alpha _{0}T\left( I_{3}\right) +\alpha
_{1}T\left( \widetilde{x}_{1}\widetilde{x}_{2}\right) +\alpha _{2}T\left( 
\widetilde{x}_{1}\widetilde{x}_{3}\right) +\alpha _{3}T\left( \widetilde{x}%
_{2}\widetilde{x}_{3}\right) =\alpha _{1}E_{1}-\alpha _{2}E_{2}+\alpha
_{3}E_{1}E_{2}
\end{equation}%
and the multipicative ones by 
\begin{eqnarray}
T\left( \widetilde{x}_{1}\widetilde{x}_{2}\right) T\left( \widetilde{x}_{1}%
\widetilde{x}_{3}\right) &=&T\left( \widetilde{x}_{1}\widetilde{x}_{2}%
\widetilde{x}_{1}\widetilde{x}_{3}\right) =-E_{1}E_{2}=-T\left( \widetilde{x}%
_{2}\widetilde{x}_{3}\right) =-E_{1}E_{2} \\
T\left( \widetilde{x}_{1}\widetilde{x}_{2}\right) ^{2} &=&E_{1}^{2}=T\left(
-I_{3}\right) =-I_{2} \\
T\left( \widetilde{x}_{1}\widetilde{x}_{3}\right) ^{2} &=&E_{2}^{2}=T\left(
-I_{3}\right) =-I_{2} \\
T\left( \widetilde{x}_{2}\widetilde{x}_{3}\right) ^{2} &=&\left(
E_{1}E_{2}\right) ^{2}=T\left( -I_{3}\right) =-I_{2}.
\end{eqnarray}%
$\rfloor $

\begin{theorem}
$Cl^{0}\left( p,q,\mathbb{R}\right) \cong Cl\left( q,p-1,\mathbb{R}\right)
\cong Cl\left( p,q-1,\mathbb{R}\right) \cong Cl^{0}\left( q,p,\mathbb{R}%
\right) $

\begin{proof}
p92, "An Introduction to Clifford Algebras and Spinors", Jayme Vax, Jr; Rold%
\~{a}o da Rocha, Jr, Oxford Press 2016.
\end{proof}
\end{theorem}

In particular 
\begin{equation}
Cl^{0}\left( 0,3,\mathbb{R}\right) \cong Cl\left( 0,2,\mathbb{R}\right) ,
\end{equation}%
which is given by the map $T.$

\begin{corollary}
$Cl^{0}\left( r,\mathbb{C}\right) \cong Cl\left( r-1,\mathbb{C}\right) .$
\end{corollary}

\subsection{Real Linear Operators}

$A$ is a \emph{real linear operator} defined on the \emph{complex seperable}
Hilbert space $\mathfrak{H}$ 
\begin{equation}
A:\mathfrak{H\rightarrow H}
\end{equation}%
if it has the properties%
\begin{equation}
A\left( x_{1}v_{1}+x_{2}v_{2}\right) =x_{1}A\left( v_{1}\right)
+x_{2}A\left( v_{2}\right) \quad \forall x_{1},x_{2}\in \mathbb{R},\forall
v_{1},v_{2}\in \mathfrak{H.}
\end{equation}%
The Banach space of bounded \emph{real linear} operators on $\mathfrak{H}$
is denoted by $\mathcal{B}_{\mathbb{R}}\left( \mathfrak{H}\right) ,$ where
the norm is defined in the usual way$.$

\begin{lemma}
Any complex Hilbert space, $\mathfrak{H,}$ can expressed as the the \emph{%
real direct sum} of two \emph{real} subspaces $V_{1}$,$iV_{1}$ $\subset 
\mathfrak{H,}$ i.e. closed under \emph{real linear} combinations denoted by%
\begin{equation}
\mathfrak{H}\cong V_{1}\oplus _{\mathbb{R}}iV_{1},
\end{equation}%
where%
\begin{equation}
V_{1}=\mathbb{R}-\limfunc{linspan}\left\{ e_{1},\cdots ,e_{r}\right\} ,
\end{equation}%
and%
\begin{equation}
\mathbb{R}-\dim \left\{ V_{1}\right\} =\mathbb{R}-\dim \left\{
iV_{1}\right\} =r.
\end{equation}
\end{lemma}

\begin{proof}
Let 
\begin{equation}
V_{1}=\mathbb{R}-\limfunc{linspan}\left\{ e_{1},\cdots,e_{r}\right\} ,
\end{equation}
where 
\begin{equation}
\left\{ e_{1},\cdots,e_{r}\right\} \text{ }
\end{equation}
are conb in $\mathfrak{H.}$ Then consider 
\begin{align}
V_{1}\oplus_{\mathbb{R}}iV_{1} & =\mathbb{R}-\limfunc{linspan}\left\{
e_{1},\cdots,e_{r}\right\} \oplus_{\mathbb{R}}\mathbb{R}-\limfunc{linspan}%
\left\{ ie_{1},\cdots,ie_{r}\right\} \\
& =\mathbb{R}-\limfunc{linspan}\left\{ \left( e_{1},0\right) ,\cdots,\left(
e_{r},0\right) ,\left( 0,ie_{1}\right) ,\cdots,\left( 0,ie_{r}\right)
\right\} \\
& =\mathbb{R}-\limfunc{linspan}\left\{ \varepsilon_{1},\cdots
,\varepsilon_{2r}\right\} ,
\end{align}
where 
\begin{align*}
\varepsilon_{j} & =\left( e_{j},0\right) ;1\leq j\leq r \\
\varepsilon_{j} & =\left( 0,ie_{j}\right) ;r+1\leq j\leq2r.
\end{align*}
As $V_{1}$ and $iV_{1}$ are disjoint \emph{real spaces} 
\begin{equation}
V_{1}\oplus_{\mathbb{R}}iV_{1}=V_{1}+_{\mathbb{R}}iV_{1}
\end{equation}%
\begin{equation}
x_{j}\left( e_{j},0\right) +y_{j}\left( 0,ie_{j}\right) =\left(
x_{j}e_{j},0\right) +\left( 0,iy_{j}e_{j}\right) =\left( x_{j}+iy_{j}\right)
e_{j}\quad\forall x_{j},y_{j}\in\mathbb{R}.
\end{equation}
Thus every element of $\mathfrak{H}$ as linear combinations of $\left\{
\left. z_{j}e_{j}\right\vert 1\leq j\leq r\right\} $ and 
\begin{equation}
\mathfrak{H}=V_{1}\oplus_{\mathbb{R}}iV_{1}.
\end{equation}
$.$
\end{proof}

\begin{lemma}
Any complex Hilbert space, $\mathfrak{H,}$ is real isometrically isomorphic
to the real Hilbert space 
\begin{equation}
\mathfrak{H}_{\mathbb{R}}\cong V_{1}\oplus _{\mathbb{R}}V_{1},
\end{equation}%
where a real symmetric inner product in $\mathfrak{H}$ is defined as%
\begin{equation}
\left\langle \left. \eta _{1}\right\vert \eta _{2}\right\rangle _{\mathfrak{H%
}_{\mathbb{R}}}=\frac{1}{2}\func{Re}\left\{ \left\langle \left. \eta
_{1}\right\vert \eta _{2}\right\rangle _{\mathfrak{H}}\right\} =\left\langle
\left. x_{1}\right\vert x_{2}\right\rangle _{V_{1}}+\left\langle \left.
y_{1}\right\vert y_{2}\right\rangle _{V_{1}};x_{1},x_{2}\in
V_{1},y_{1},y_{2}\in V_{1},
\end{equation}%
and $\left\langle \left. \cdot \right\vert \cdot \right\rangle _{V_{1}}$ is
the real symmetric inner product in $V_{1}.$
\end{lemma}

\begin{definition}
A real operator $A\in \mathcal{B}_{\mathbb{R}}\left( \mathfrak{H}\right) $
is \emph{complex} \emph{antilinear} if 
\begin{equation}
\left\{ A,iI_{\mathcal{A}}\right\} =0
\end{equation}%
and \emph{complex} \emph{linear} if 
\begin{equation}
\left[ A,iI_{\mathcal{A}}\right] =0.
\end{equation}
\end{definition}

\begin{remark}
The descriptor "complex" is usually removed when it is understood from the
context.
\end{remark}

The space of \emph{bounded} anti linear operators is denoted $\mathcal{B}%
_{anti}\left( \mathfrak{H}\right) ,$ while the space of linear operators is
denoted, in the standard way, by $\mathcal{B}\left( \mathfrak{H}\right) ,$
both wrt the usual operator norm induced by $\mathfrak{H}.$

\emph{Both} anti linear and linear operators belong to the space of \emph{%
real linear operators} $\mathcal{B}_{\mathbb{R}}\left( \mathfrak{H}\right) $
and any real operator, $A_{\mathbb{R}}\in \mathcal{B}_{\mathbb{R}}\left( 
\mathfrak{H}\right) $ can be decomposed as a sum of a linear operator $A_{%
\mathbb{C}}\in \mathcal{B}\left( \mathfrak{H}\right) $ and an antilinear one 
$A_{a}\in \mathcal{B}_{anti}\left( \mathfrak{H}\right) $ as 
\begin{equation}
A_{\mathbb{R}}=A_{\mathbb{C}}+A_{a},
\end{equation}%
where 
\begin{align}
A_{\mathbb{C}}& =\frac{1}{2}\left( A_{\mathbb{R}}-iA_{\mathbb{R}}i\right) \\
A_{a}& =\frac{1}{2}\left( A_{\mathbb{R}}+iA_{\mathbb{R}}i\right) .
\end{align}%
Complex linear operators, $A^{\mathbb{C}}\in \mathcal{B}\left( \mathfrak{H}_{%
\mathbb{C}}\right) ,$ $\left( \text{called the complexification of the
operator }A_{\mathbb{R}}\right) $ on $\mathfrak{H}_{\mathbb{C}}$ are of the
form 
\begin{equation}
A^{\mathbb{C}}=\left( 
\begin{array}{cc}
A_{\mathbb{C}} & A_{a}\sigma \\ 
\sigma A_{a} & \sigma A_{\mathbb{C}}\sigma%
\end{array}%
\right) ,
\end{equation}%
where 
\begin{equation}
\mathfrak{H}_{\mathbb{C}}\cong V_{1}\oplus _{\mathbb{C}}iV_{2}\cong 
\mathfrak{H}\oplus _{\mathbb{C}}\mathfrak{H\cong \mathfrak{H}\oplus _{%
\mathbb{C}}}\sigma \left( \mathfrak{\mathfrak{H}}\right) \mathfrak{,}
\end{equation}%
$\sigma $ is a complex conjugation $\in \mathcal{B}_{anti}\left( \mathfrak{H}%
\right) $ (defined in the following) and 
\begin{equation}
\mathbb{C}-\dim \left\{ \mathfrak{H}_{\mathbb{C}}\right\} =2r.
\end{equation}

\begin{remark}
The subspaces $V_{1}$,$V_{2}$ can be identical (thus isomorphic) as Hilbert
subspaces or just isomorphic as Hilbert subspaces. If the subspaces are not
equal they can be constituted by elements that have different properties in
the sense of the way they act on other mathematical objects.
\end{remark}

\begin{remark}
The real Banach space $\mathcal{B}_{\mathbb{R}}\left( \mathfrak{H}\right) $
contains both $\mathcal{B}\left( \mathfrak{H}\right) $ and $\mathcal{B}%
_{anti}\left( \mathfrak{H}\right) $ as real subspaces and 
\begin{equation}
\mathcal{B}_{\mathbb{R}}\left( \mathfrak{H}\right) =\mathcal{B}\left( 
\mathfrak{H}\right) \oplus _{\mathbb{R}}\mathcal{B}_{anti}\left( \mathfrak{H}%
\right) .
\end{equation}%
As $\mathcal{B}\left( \mathfrak{H}\right) $ contains $\mathcal{B}_{sa}\left( 
\mathfrak{H}\right) $ as a real subspace one has 
\begin{equation}
\mathcal{B}_{\mathbb{R}}\left( \mathfrak{H}\right) \supset \mathcal{B}\left( 
\mathfrak{H}\right) \supset \mathcal{B}_{sa}\left( \mathfrak{H}\right)
\end{equation}%
as real spaces.
\end{remark}

The spectrum $\mathfrak{s}\left( A\right) $ of both linear and anti linear
operators are defined in the same way i.e. $\lambda\in\mathfrak{s}\left(
A\right) $ $\limfunc{Iff}$ $\left( A-\lambda I_{\mathcal{A}}\right) $ has no
bounded inverse. The approximate point spectrum, $\mathfrak{s}_{ap}(A),$
consists of those $\U{3bb} \in\mathbb{C}$ for which there exists a sequence
of vectors $\left\{ \left. \left\vert e_{j}\right\rangle \right\vert 1\leq
j\leq\infty\right\} $ of unit length for which 
\begin{equation}
\lim\limits_{j\rightarrow\infty}\left\{ \left( A-\lambda I_{\mathcal{A}%
}\right) \left\vert e_{j}\right\rangle \right\} \rightarrow0.
\end{equation}
The point spectrum, $\mathfrak{s}_{p}(A),$ consists of eigenvalues of $A$.
The residual spectrum, $\mathfrak{s}_{res}(A),$ consists of $\U{3bb} \in 
\mathbb{C}$ for which the range of $A-\lambda I_{\mathcal{A}}$ is not dense
in $\mathfrak{H}$.

\subsection{Complex Conjugations}

A \emph{complex conjugation, }$\sigma _{e},$ 
\begin{equation}
\sigma _{e}:\mathfrak{H}\rightarrow \sigma _{e}\left( \mathfrak{H}\right)
\cong \mathfrak{H}
\end{equation}%
on a complex Hilbert space $\mathfrak{H}$ is defined as a \emph{complex
antilinear} 
\begin{equation}
\sigma _{e}\left( z_{1}v_{1}+z_{2}v_{2}\right) =\overline{z_{1}}\sigma
_{e}\left( v_{1}\right) +\overline{z_{2}}\sigma _{e}\left( v_{2}\right)
;z_{1},z_{2}\in \mathbb{C}
\end{equation}%
\emph{isomorphism} that is also an \emph{antilinear involution} 
\begin{equation}
\sigma _{e}^{2}=I_{\mathfrak{H}}.
\end{equation}%
It can be constructed by choosing an equivalence class, $\left[ \mathbf{e}%
\right] ,$ of complete orthonormal bases of $\mathfrak{H}$%
\begin{equation}
\left[ \mathbf{e}\right] =\left\{ \left. \left\{ \left. S\left\vert
e_{j}\right\rangle \right\vert 1\leq j\leq r\right\} \right\vert S\in
O\left( \mathfrak{H,}\mathbb{R}\right) \equiv O\left( r,\mathbb{R}\right)
\right\}
\end{equation}%
with representative basis $\left\{ \left. \left\vert e_{j}\right\rangle
\right\vert 1\leq j\leq r\right\} $ leading to 
\begin{equation}
\lambda =\sum_{1\leq j\leq r}\lambda _{j}\left\vert e_{j}\right\rangle
;\lambda _{j}\in \mathbb{C},
\end{equation}%
such that 
\begin{equation}
\sigma _{e}\left( \lambda \right) =\sum_{1\leq j\leq r}\bar{\lambda}%
_{j}\left\vert e_{j}\right\rangle .
\end{equation}%
This shows that 
\begin{equation}
\sigma _{e}\left( \sigma _{e}\left( \lambda \right) \right) =\sigma
_{e}\left( \sum_{1\leq j\leq r}\bar{\lambda}_{j}\left\vert
e_{j}\right\rangle \right) =\sum_{1\leq j\leq r}\lambda _{j}\left\vert
e_{j}\right\rangle =\lambda ,
\end{equation}%
i.e. 
\begin{equation}
\sigma _{e}^{2}=I_{\mathfrak{H}},
\end{equation}%
so $\sigma _{e}$ is an involution (but not a \emph{linear} involution i.e. $%
\notin \mathcal{B}\left( \mathfrak{H}\right) )$. It is easy to see that $%
\sigma _{e}$ is \emph{an antiunitary operator} i.e. 
\begin{equation}
\left\Vert \sigma _{e}\left( v\right) \right\Vert _{\mathfrak{H}%
}^{2}=\left\langle \left. \sigma _{e}\left( v\right) \right\vert \sigma
_{e}\left( v\right) \right\rangle =\left\langle \left. v\right\vert
v\right\rangle .
\end{equation}%
The involution, $\sigma _{e},$ is a \emph{real linear mapping} on the
complex Hlbert space $\mathfrak{H}$ i.e.%
\begin{equation}
\sigma _{e}\left( x_{1}v_{1}+x_{2}v_{2}\right) =x_{1}\sigma _{e}\left(
v_{1}\right) +x_{2}\sigma _{e}\left( v_{2}\right) ;x_{1},x_{2}\in \mathbb{R}
\end{equation}%
and in particular on the real $2r$-dimensional Hilbert space 
\begin{equation}
\mathfrak{H}_{\mathbb{R}}=V_{\sigma _{e}}^{r}\oplus _{\mathbb{R}}V_{\sigma
_{e}}^{i},
\end{equation}%
where 
\begin{equation}
\left. \sigma _{e}\right\vert _{V_{\sigma _{e}}^{r}}=I_{r}
\end{equation}%
and 
\begin{equation}
\left. \sigma _{e}\right\vert _{V_{\sigma _{e}}^{i}}=-I_{r}.
\end{equation}%
It is easy to see that 
\begin{equation}
V_{\sigma _{e}}^{r}\oplus _{\mathbb{R}}V_{\sigma _{e}}^{i}\cong V_{\sigma
_{f}}^{r}\oplus _{\mathbb{R}}V_{\sigma _{f}}^{i}
\end{equation}%
for all complex conjugations $\sigma _{e}$ and $\sigma _{f}$ even when $%
e\nsim f,$ i.e. $\sigma _{e}$, $\sigma _{f}$ do not belong to the same
equivalence class viz 
\begin{equation}
\left\vert \mathbf{e}\right\rangle \neq o\left\vert \mathbf{f}\right\rangle
;o\in O\left( \mathfrak{H,}\mathbb{R}\right) \equiv O\left( r\mathfrak{,}%
\mathbb{R}\right) .
\end{equation}

The complex conjugation $\sigma_{e}\in\mathcal{B}_{anti}\left( \mathfrak{H}%
\right) $ induces a "complex conjugation", $\sigma_{e}\in\mathcal{B}_{%
\mathbb{R}}\left( \mathfrak{H}\right) ,$ as $\mathcal{B}_{\mathbb{R}}\left( 
\mathfrak{H}\right) \supset\mathcal{B}_{anti}\left( \mathfrak{H}\right) .$
Note that $\mathcal{B}_{\mathbb{R}}\left( \mathfrak{H}\right) \cong\mathcal{B%
}\left( \mathfrak{H}_{\mathbb{R}}\right) \cong\mathcal{B}\left( \mathbb{R}%
^{2r}\right) .$

\begin{notation}
\begin{align}
V_{\sigma _{e}}^{r}& \equiv V_{+\sigma _{e}}\equiv V_{1} \\
V_{\sigma _{e}}^{i}& \equiv V_{-\sigma _{e}}\equiv V_{2}.
\end{align}
\end{notation}

\begin{theorem}
For $\func{Re}_{e},\func{Im}_{e}\in \mathcal{B}_{\mathbb{R}}\left( \mathfrak{%
H}\right) ,\sigma _{e}\in \mathcal{B}_{anti}\left( \mathfrak{H}\right)
\subset \mathcal{B}_{\mathbb{R}}\left( \mathfrak{H}\right) $ defined
previously 
\begin{eqnarray}
\func{Re}_{e} &=&\frac{1}{2}\left\{ I_{\mathfrak{H}}+\sigma _{e}\right\} \\
\func{Im}_{e} &=&-\frac{i}{2}\left\{ I_{\mathfrak{H}}-\sigma _{e}\right\} ,
\end{eqnarray}%
which is equivalent to 
\begin{eqnarray}
\func{Re}_{e}+i\func{Im}_{e} &=&I_{\mathfrak{H}} \\
\func{Re}_{e}-i\func{Im}_{e} &=&\sigma _{e}.
\end{eqnarray}%
$\lceil $%
\begin{eqnarray}
2\func{Re}_{e} &=&I_{\mathfrak{H}}+\sigma _{e} \\
2i\func{Im}_{e} &=&I_{\mathfrak{H}}-\sigma _{e}.
\end{eqnarray}%
$\rfloor $
\end{theorem}

\subparagraph{Complex Symmetric Form}

One can define a complex symmetric inner product in 
\begin{equation}
\left( \left. v_{2}\right\vert v_{1}\right) _{e\mathfrak{H}}=\left( \left.
v_{1}\right\vert v_{2}\right) _{\mathfrak{H}}=\sum_{1\leq j\leq r}\mathcal{R}%
_{e}\left( v_{1}\right) _{j}\mathcal{R}_{e}\left( v_{2}\right) _{j},
\end{equation}%
\emph{that can have positive, zero or negative values, }

$\lceil $In particular \emph{\ } 
\begin{equation}
\sum_{1\leq j\leq r}\mathcal{R}_{e}\left( v_{1}\right) _{j}\mathcal{R}%
_{e}\left( v_{1}\right) _{j}=\sum_{1\leq j\leq r}z_{j}^{2}=\sum_{1\leq j\leq
r}\left( x_{j}+iy_{j}\right) ^{2}=\sum_{1\leq j\leq r}\left(
x_{j}^{2}-y_{j}^{2}+2ix_{j}y_{j}\right)
\end{equation}%
and consider 
\begin{align}
x_{1}y_{1}+x_{2}y_{2}& =0 \\
1\left( 1\right) +1\left( -1\right) & =0.
\end{align}%
$\rfloor $

in $\mathfrak{H}$ with invariance group $O\left( r,\mathbb{C}\right) .$ The
inner product $\left( \left. v_{2}\right\vert v_{1}\right) _{e\mathfrak{H}}$
is not unitarily invariant, i.e. not $U\left( r\right) $-invariant, but is
invariant to the subgroup 
\begin{equation}
O\left( r,\mathbb{R}\right) =O\left( r,\mathbb{C}\right) \cap U\left(
r\right) .
\end{equation}%
Hence the value of the sum 
\begin{equation}
\sum_{1\leq j\leq r}\mathcal{R}_{e}\left( v_{2}\right) _{j}\mathcal{R}%
_{e}\left( v_{1}\right) _{j}
\end{equation}%
in general depends on the representation $\mathcal{R}_{e}$ even if they are
unitarily equivalent i.e. 
\begin{equation}
\sum_{1\leq j\leq r}\mathcal{R}_{e}\left( v_{2}\right) _{j}\mathcal{R}%
_{e}\left( v_{1}\right) _{j}\neq \sum_{1\leq j\leq r}\mathcal{R}_{f}\left(
v_{2}\right) _{j}\mathcal{R}_{f}\left( v_{1}\right) _{j},
\end{equation}%
where 
\begin{equation}
\left\vert \mathbf{e}\right\rangle =u\left\vert \mathbf{f}\right\rangle
;u\in U\left( \mathfrak{H}\right) .
\end{equation}%
They are only equal if they are\emph{\ real orthogonal equivalent} i.e. 
\begin{equation}
\left\vert \mathbf{e}\right\rangle \sim \left\vert \mathbf{f}\right\rangle
\Leftrightarrow \left\vert \mathbf{e}\right\rangle =o\left\vert \mathbf{f}%
\right\rangle ;o\in O\left( \mathfrak{H,}\mathbb{R}\right) \subset U\left( 
\mathfrak{H}\right) .
\end{equation}

\subsubsection{Real Subspaces of $\mathfrak{H}$}

As previosly discussed real subspace, (not closed under complex linear
combinations, but closed under real ones), $V_{+\sigma _{e}}\subset 
\mathfrak{H,}$ can be defined as 
\begin{equation}
V_{+\sigma _{e}}=\mathbb{R-}\limfunc{linspan}\left\{ \left. v+\sigma
_{e}\left( v\right) \right\vert v\in \mathfrak{H}\right\} =\mathbb{R-}%
\limfunc{linspan}\left\{ \left. \sigma _{e}\left( \varsigma \right)
=\varsigma \right\vert \varsigma \in \mathfrak{H}\right\} ,
\end{equation}%
as%
\begin{gather}
v+\sigma _{e}\left( v\right) =\sigma _{e}\left( v\right) +v=\sum_{1\leq
j\leq r}\left( \lambda _{j}+\bar{\lambda}_{j}\right) \left\vert
e_{j}\right\rangle =2\sum_{1\leq j\leq r}\func{Re}\left( \lambda _{j}\right)
\left\vert e_{j}\right\rangle \\
\sigma _{e}\left( v+\sigma _{e}\left( v\right) \right) =\sigma _{e}\left(
v\right) +v=v+\sigma _{e}\left( v\right) .
\end{gather}%
Similarly one can define the real complementary real space $V_{-\sigma _{e}}$
\begin{equation}
V_{-\sigma _{e}}=\mathbb{R-}\limfunc{linspan}\left\{ \left. -i\left(
v-\sigma _{e}\left( v\right) \right) \right\vert v\in \mathfrak{H}\right\} =%
\mathbb{R-}\limfunc{linspan}\left\{ \left. \sigma _{e}\left( \varsigma
\right) =-\varsigma \right\vert \varsigma \in \mathfrak{H}\right\} ,
\end{equation}%
as%
\begin{gather}
-i\left( v-\sigma _{e}\left( v\right) \right) =-i\left( \sigma _{e}\left(
v\right) -v\right) =-i\sum_{1\leq j\leq r}\left( \lambda _{j}-\bar{\lambda}%
_{j}\right) \left\vert e_{j}\right\rangle =2\sum_{1\leq j\leq r}\func{Im}%
\left( \lambda _{j}\right) \left\vert e_{j}\right\rangle \\
\sigma _{e}\left( v-\sigma _{e}\left( v\right) \right) =\sigma _{e}\left(
v\right) -v=-\left( v-\sigma _{e}\left( v\right) \right) .
\end{gather}%
The subspaces $V_{+\sigma _{e}}$ $\subset \mathfrak{H}$ and $V_{-\sigma
_{e}}\subset \mathfrak{H}$ are isomorphic as real Hilbert subspaces. The
vectors $v+\sigma _{e}\left( v\right) $ are eigenvectors of $\sigma _{e}$
with eigenvalue $+1$ and $v-\sigma _{e}\left( v\right) $ and $-i\left(
v-\sigma _{e}\left( v\right) \right) $ with eigenvalues $-1.$ Clearly $%
V_{+\sigma _{e}}$ describes the $e$-real part of a vector, while $V_{-\sigma
_{e}}$ the $e$-imaginary part.

$\lceil$If $v$ is an eigenvector with eigen value $\lambda$ then $av$ is
still an eigenvector with eigen value $\lambda.\rfloor$

Remember $\sigma_{e}\notin\mathcal{B}\left( \mathfrak{H}\right) $ so the
usual properties of eigenspaces of operators do not apply$.$

\begin{notation}
$\sigma_{e}$ denotes a complex conjugation associated with a specific basis $%
\left\vert \mathbf{e}\right\rangle ,$ while $\sigma_{\left[ e\right] }$
denotes any complex conjugation assciated with any basis in the equivalence
class $\left[ e\right] .$ Obviously 
\begin{equation}
\sigma_{e}=\sigma_{\left[ e\right] }
\end{equation}
but the basis might not be the same, though they are related by a real
orthogonal transformation.
\end{notation}

\subsection{Complexification}

\begin{definition}
A \emph{complexification operator}$\,\equiv\,$\emph{complexification, }$J,$%
\emph{\ }of a real even dimensional Hilbert space, $V$, is an antisymmetric
operator $J\in\mathcal{B}_{\mathbb{R}}\left( V\right) \equiv\mathcal{B}_{%
\mathbb{R}}\left( \mathfrak{H}\right) \cong\mathcal{B}\left( \mathbb{R}%
^{2r}\right) $\emph{\ }that has the properties\emph{\ }%
\begin{align}
J & =-J^{t} \\
J^{2} & =-I_{V}.
\end{align}
\end{definition}

\begin{lemma}
If $J$ is a complexification so is $oJo^{t}$ for all $o\in O\left( 2r,%
\mathbb{R}\right) \cong O\left( V,\mathbb{R}\right) $
\end{lemma}

\begin{definition}
A complex conjugation $\sigma\in\mathcal{B}_{\mathbb{R}}\left( V\right) $ is 
\emph{compatible} with a complexification $J\in\mathcal{B}_{\mathbb{R}%
}\left( V\right) $ if 
\begin{equation}
\left\{ \sigma,J\right\} =0.
\end{equation}
\end{definition}

$\lceil$%
\begin{equation}
o\sigma o^{t}oJo^{t}+oJo^{t}o\sigma o^{t}=o\left\{ \sigma,J\right\} o^{t}.
\end{equation}
$\rfloor$

\begin{lemma}
If $\sigma $ $\in \mathcal{B}\left( V\right) $ and $J$ $\in \mathcal{B}%
\left( V\right) $ are \emph{compatible} then the complex conjugations $%
o\sigma o^{t}$ and complexifications $oJo^{t}$ are also \emph{compatible},
for each $o\in O\left( 2r,\mathbb{R}\right) \equiv O\left( V,\mathbb{R}%
\right) .$ Moreover \emph{all} complex conjugations of $V$ and \emph{all }%
complexifications can be generated in this way i.e.%
\begin{equation}
\mathfrak{C}\left( V\right) =\left\{ \left. o\sigma o^{t}\right\vert o\in
O\left( 2r,\mathbb{R}\right) ,\quad \sigma \text{ a complex conjugation}%
\right\}
\end{equation}%
and 
\begin{equation}
\mathfrak{C}\left( V\right) =\left\{ \left. oJo^{t}\right\vert o\in O\left(
2r,\mathbb{R}\right) ,\quad J\text{ a complex unit}\right\} ,
\end{equation}%
$\left( J\text{ is not necessarily compatible with }\sigma \right) ,$ thus
producing \emph{all} compatible complex conjugations and complexifications
when $J$ is compatible with $\sigma $ by 
\begin{align}
o\sigma o^{t}& =o\left\vert \mathbf{e}_{1}\right\rangle \left\langle \mathbf{%
e}_{1}\right\vert o^{t}-o\left\vert \mathbf{e}_{2}\right\rangle \left\langle 
\mathbf{e}_{2}\right\vert o^{t} \\
oJo^{t}& =o\left\vert \mathbf{e}_{2}\right\rangle \left\langle \mathbf{e}%
_{1}\right\vert o^{t}-o\left\vert \mathbf{e}_{1}\right\rangle \left\langle 
\mathbf{e}_{2}\right\vert o^{t}
\end{align}%
as $o$ ranges over all of $O\left( 2r,\mathbb{R}\right) .$ This generation
can be explicitly expressed as%
\begin{align}
o\sigma o^{t}& =o\left\vert \mathbf{e}_{1},\mathbf{e}_{2}\right\rangle
\left( 
\begin{array}{cc}
\mathbb{I} & \mathbb{O} \\ 
\mathbb{O} & -\mathbb{I}%
\end{array}%
\right) \left\langle \mathbf{e}_{1},\mathbf{e}_{2}\right\vert
o^{t}=\left\vert \mathbf{e}_{1},\mathbf{e}_{2}\right\rangle \mathbb{Q}\left( 
\begin{array}{cc}
\mathbb{I} & \mathbb{O} \\ 
\mathbb{O} & -\mathbb{I}%
\end{array}%
\right) \mathbb{Q}^{t}\left\langle \mathbf{e}_{1},\mathbf{e}_{2}\right\vert
\\
oJo^{t}& =o\left\vert \mathbf{e}_{1},\mathbf{e}_{2}\right\rangle \left( 
\begin{array}{cc}
\mathbb{O} & -\mathbb{I} \\ 
\mathbb{I} & \mathbb{O}%
\end{array}%
\right) \left\langle \mathbf{e}_{1},\mathbf{e}_{2}\right\vert
o^{t}=\left\vert \mathbf{e}_{1},\mathbf{e}_{2}\right\rangle \mathbb{Q}\left( 
\begin{array}{cc}
\mathbb{O} & -\mathbb{I} \\ 
\mathbb{I} & \mathbb{O}%
\end{array}%
\right) \mathbb{Q}^{t}\left\langle \mathbf{e}_{1},\mathbf{e}_{2}\right\vert .
\end{align}
\end{lemma}

\begin{lemma}
Any even dimensional real Hilbert space can expanded as the real direct sum
of two orthogonal subspaces as 
\begin{equation}
V=V_{1}\oplus _{\mathbb{R}}V_{2},
\end{equation}%
where $V_{1}$ has the conb $\left\vert \mathbf{e}_{1}\right\rangle $ and $%
V_{2}$ has the conb $\left\vert \mathbf{e}_{2}\right\rangle .$ One can
define a complex conjugation and a compatible complexification associated
with this direct sum decomposition by 
\begin{align}
\sigma & =\left\vert \mathbf{e}_{1}\right\rangle \left\langle \mathbf{e}%
_{1}\right\vert -\left\vert \mathbf{e}_{2}\right\rangle \left\langle \mathbf{%
e}_{2}\right\vert =\left\vert \mathbf{e}_{1},\mathbf{e}_{2}\right\rangle 
\mathcal{R}_{e}\left( \sigma \right) \left\langle \mathbf{e}_{1},\mathbf{e}%
_{2}\right\vert =\left\vert \mathbf{e}_{1},\mathbf{e}_{2}\right\rangle
\left( 
\begin{array}{cc}
\mathbb{I} & \mathbb{O} \\ 
\mathbb{O} & -\mathbb{I}%
\end{array}%
\right) \left\langle \mathbf{e}_{1},\mathbf{e}_{2}\right\vert
=P_{v_{1}}-P_{v_{2}} \\
J& =\left\vert \mathbf{e}_{2}\right\rangle \left\langle \mathbf{e}%
_{1}\right\vert -\left\vert \mathbf{e}_{1}\right\rangle \left\langle \mathbf{%
e}_{2}\right\vert =\left\vert \mathbf{e}_{1},\mathbf{e}_{2}\right\rangle 
\mathcal{R}_{e}\left( J\right) \left\langle \mathbf{e}_{1},\mathbf{e}%
_{2}\right\vert =\left\vert \mathbf{e}_{1},\mathbf{e}_{2}\right\rangle
\left( 
\begin{array}{cc}
\mathbb{O} & -\mathbb{I} \\ 
\mathbb{I} & \mathbb{O}%
\end{array}%
\right) \left\langle \mathbf{e}_{1},\mathbf{e}_{2}\right\vert .
\end{align}%
$\lceil $The projectors $\left\{ P_{v_{1}},P_{v_{2}}\right\} $ and the
complexification have the properties 
\begin{align}
\left\vert \mathbf{e}_{1}\right\rangle \left\langle \mathbf{e}%
_{1}\right\vert \left( \left\vert \mathbf{e}_{2}\right\rangle \left\langle 
\mathbf{e}_{1}\right\vert -\left\vert \mathbf{e}_{1}\right\rangle
\left\langle \mathbf{e}_{2}\right\vert \right) & =-\left\vert \mathbf{e}%
_{1}\right\rangle \left\langle \mathbf{e}_{2}\right\vert \\
\left( \left\vert \mathbf{e}_{2}\right\rangle \left\langle \mathbf{e}%
_{1}\right\vert -\left\vert \mathbf{e}_{1}\right\rangle \left\langle \mathbf{%
e}_{2}\right\vert \right) \left\vert \mathbf{e}_{1}\right\rangle
\left\langle \mathbf{e}_{1}\right\vert & =\left\vert \mathbf{e}%
_{2}\right\rangle \left\langle \mathbf{e}_{1}\right\vert \\
\left\vert \mathbf{e}_{2}\right\rangle \left\langle \mathbf{e}%
_{2}\right\vert \left( \left\vert \mathbf{e}_{2}\right\rangle \left\langle 
\mathbf{e}_{1}\right\vert -\left\vert \mathbf{e}_{1}\right\rangle
\left\langle \mathbf{e}_{2}\right\vert \right) & =\left\vert \mathbf{e}%
_{2}\right\rangle \left\langle \mathbf{e}_{1}\right\vert \\
\left( \left\vert \mathbf{e}_{2}\right\rangle \left\langle \mathbf{e}%
_{1}\right\vert -\left\vert \mathbf{e}_{1}\right\rangle \left\langle \mathbf{%
e}_{2}\right\vert \right) \left\vert \mathbf{e}_{2}\right\rangle
\left\langle \mathbf{e}_{2}\right\vert & =-\left\vert \mathbf{e}%
_{1}\right\rangle \left\langle \mathbf{e}_{2}\right\vert
\end{align}%
giving 
\begin{align}
P_{v_{1}}J& =JP_{v_{2}} \\
J^{-1}P_{v_{1}}J& =-J^{t}P_{v_{1}}J=-JP_{v_{1}}J=P_{v_{2}}
\end{align}%
and%
\begin{align}
P_{v_{2}}J& =JP_{v_{1}} \\
J^{-1}P_{v_{2}}J& =-J^{t}P_{v_{2}}J=-JP_{v_{1}}J=P_{v_{1}}.
\end{align}%
$\rfloor $
\end{lemma}

\begin{lemma}
If $P_{v_{1}},P_{v_{2}},J\in \mathcal{B}\left( V\right) \cong \mathcal{B}%
\left( \mathbb{R}^{2r}\right) $ and 
\begin{equation}
P_{v_{1}}J=JP_{v_{2}},
\end{equation}%
where $\left\{ P_{v_{1}},P_{v_{2}}\right\} $ are complimentary orthogonal
projectors spanning $V$ and $J$ is a complexification i.e.%
\begin{align}
J& =-J^{t} \\
J^{2}& =-I_{2r}
\end{align}%
that can be expanded as 
\begin{equation}
J=\left\vert \mathbf{e}_{1},\mathbf{e}_{2}\right\rangle \left( 
\begin{array}{cc}
\mathbb{O} & -\mathbb{I} \\ 
\mathbb{I} & \mathbb{O}%
\end{array}%
\right) \left\langle \mathbf{e}_{1},\mathbf{e}_{2}\right\vert
\end{equation}%
then 
\begin{equation}
\sigma =P_{V_{1}}-P_{V_{2}}=\left\vert \mathbf{e}_{1},\mathbf{e}%
_{2}\right\rangle \left( 
\begin{array}{cc}
\mathbb{I} & \mathbb{O} \\ 
\mathbb{O} & -\mathbb{I}%
\end{array}%
\right) \left\langle \mathbf{e}_{1},\mathbf{e}_{2}\right\vert =\left\vert 
\mathbf{f}_{1},\mathbf{f}_{2}\right\rangle \left( 
\begin{array}{cc}
\mathbb{I} & \mathbb{O} \\ 
\mathbb{O} & -\mathbb{I}%
\end{array}%
\right) \left\langle \mathbf{f}_{1},\mathbf{f}_{2}\right\vert
\end{equation}%
\emph{are} \emph{compatible conjugation} for all conb's, $\left\{ \left\vert 
\mathbf{e}_{1}\right\rangle \right\} $ and $\left\{ \left\vert \mathbf{f}%
_{1}\right\rangle \right\} ,$ of $V_{1}$. and $\left\{ \left\vert \mathbf{e}%
_{2}\right\rangle \right\} $ and $\left\{ \left\vert \mathbf{f}%
_{2}\right\rangle \right\} ,$ of $V_{2}$ i.e. 
\begin{align}
\left\vert \mathbf{e}_{1}\right\rangle & =o_{1}\left\vert \mathbf{f}%
_{1}\right\rangle \\
\left\vert \mathbf{e}_{2}\right\rangle & =o_{2}\left\vert \mathbf{f}%
_{2}\right\rangle \\
o_{1},o_{2}& \in O\left( r,\mathbb{R}\right) .
\end{align}
\end{lemma}

\begin{corollary}
Let $Q_{V_{12}}\in \mathcal{B}\left( V_{2},V_{1}\right) $ then 
\begin{eqnarray}
Q_{V_{12}} &=&\left\vert \mathbf{e}_{1}\right\rangle \left\langle \mathbf{e}%
_{2}\right\vert \\
Q_{V_{12}}^{\dagger }Q_{V_{12}} &=&\left\vert \mathbf{e}_{2}\right\rangle
\left\langle \mathbf{e}_{1}\right\vert \left\vert \mathbf{e}%
_{1}\right\rangle \left\langle \mathbf{e}_{2}\right\vert =\left\vert \mathbf{%
e}_{2}\right\rangle \left\langle \mathbf{e}_{2}\right\vert =P_{V_{2}} \\
Q_{V_{12}}Q_{V_{12}}^{\dagger } &=&\left\vert \mathbf{e}_{1}\right\rangle
\left\langle \mathbf{e}_{2}\right\vert \left\vert \mathbf{e}%
_{2}\right\rangle \left\langle \mathbf{e}_{1}\right\vert =\left\vert \mathbf{%
e}_{1}\right\rangle \left\langle \mathbf{e}_{1}\right\vert =P_{V_{1}}
\end{eqnarray}
\end{corollary}

We can extend $P_{V_{1}},P_{V_{2}}$ and $Q_{V_{12}}$ to $\mathcal{B}\left(
V\right) $ by 
\begin{eqnarray}
\widetilde{P_{V_{1}}} &=&\left\vert \left( \mathbf{e}_{1},\mathbf{0}\right)
\right\rangle \left\langle \left( \mathbf{e}_{1}\mathbf{,0}\right)
\right\vert \\
\widetilde{P_{V_{2}}} &=&\left\vert \left( \mathbf{0},\mathbf{e}_{2}\right)
\right\rangle \left\langle \left( \mathbf{0,e}_{2}\right) \right\vert \\
\widetilde{Q_{V_{12}}} &=&\left\vert \left( \mathbf{e}_{1},\mathbf{0}\right)
\right\rangle \left\langle \left( \mathbf{0,e}_{2}\right) \right\vert \\
\widetilde{Q_{V_{21}}} &=&\left\vert \left( \mathbf{0,e}_{2}\right)
\right\rangle \left\langle \left( \mathbf{e}_{1},\mathbf{0}\right)
\right\vert \\
\widetilde{Q_{V_{21}}} &=&\widetilde{Q_{V_{21}}}^{\dagger }
\end{eqnarray}%
such that the restrictions 
\begin{eqnarray}
\left. \widetilde{Q_{V_{12}}}\right\vert _{V_{2}} &=&\left\vert \left( 
\mathbf{e}_{1},\mathbf{0}\right) \right\rangle \left\langle \mathbf{e}%
_{2}\right\vert \\
\left. \widetilde{Q_{V_{21}}}\right\vert _{V_{1}} &=&\left\vert \left( 
\mathbf{0,e}_{2}\right) \right\rangle \left\langle \mathbf{e}_{1}\right\vert
\end{eqnarray}
give%
\begin{eqnarray}
P_{V_{2}} &=&\left( \left\vert \left( \mathbf{e}_{1},\mathbf{0}\right)
\right\rangle \left\langle \mathbf{e}_{2}\right\vert \right) ^{\dagger
}\left\vert \left( \mathbf{e}_{1},\mathbf{0}\right) \right\rangle
\left\langle \mathbf{e}_{2}\right\vert =\left\vert \mathbf{e}%
_{2}\right\rangle \left\langle \left( \mathbf{e}_{1},\mathbf{0}\right)
\right\vert \left\vert \left( \mathbf{e}_{1},\mathbf{0}\right) \right\rangle
\left\langle \mathbf{e}_{2}\right\vert =\left( \left. \widetilde{Q_{V_{12}}}%
\right\vert _{V_{2}}\right) ^{\dagger }\left. \widetilde{Q_{V_{12}}}%
\right\vert _{V_{2}}=\left\vert \mathbf{e}_{2}\right\rangle \left\langle 
\mathbf{e}_{2}\right\vert \\
P_{V_{1}} &=&\left( \left\vert \left( \mathbf{0,e}_{1}\right) \right\rangle
\left\langle \mathbf{e}_{1}\right\vert \right) ^{\dagger }\left\vert \left( 
\mathbf{0,e}_{1}\right) \right\rangle \left\langle \mathbf{e}_{1}\right\vert
=\left\vert \mathbf{e}_{1}\right\rangle \left\langle \left( \mathbf{0},%
\mathbf{e}_{2}\right) \right\vert \left\vert \left( \mathbf{0},\mathbf{e}%
_{2}\right) \right\rangle \left\langle \mathbf{e}_{1}\right\vert =\left(
\left. \widetilde{Q_{V_{21}}}\right\vert _{V_{1}}\right) ^{\dagger }\left. 
\widetilde{Q_{V_{21}}}\right\vert _{V_{1}}=\left\vert \mathbf{e}%
_{1}\right\rangle \left\langle \mathbf{e}_{1}\right\vert .
\end{eqnarray}

\begin{proposition}
Compatible conjugations 
\begin{gather}
\sigma _{o}=\left\vert o_{1}\mathbf{e}_{1},o_{2}\mathbf{e}_{2}\right\rangle
\left( 
\begin{array}{cc}
\mathbb{I} & \mathbb{O} \\ 
\mathbb{O} & -\mathbb{I}%
\end{array}%
\right) \left\langle o_{1}\mathbf{e}_{1},o_{2}\mathbf{e}_{2}\right\vert
=\left\vert \mathbf{f}_{1},\mathbf{f}_{2}\right\rangle \left( 
\begin{array}{cc}
\mathbb{I} & \mathbb{O} \\ 
\mathbb{O} & -\mathbb{I}%
\end{array}%
\right) \left\langle \mathbf{f}_{1},\mathbf{f}_{2}\right\vert \\
=\left\vert \mathbf{e}_{1},\mathbf{e}_{2}\right\rangle \left( 
\begin{array}{cc}
\mathbb{O}_{1} & \mathbb{O} \\ 
\mathbb{O} & \mathbb{O}_{2}%
\end{array}%
\right) \left( 
\begin{array}{cc}
\mathbb{I} & \mathbb{O} \\ 
\mathbb{O} & -\mathbb{I}%
\end{array}%
\right) \left( 
\begin{array}{cc}
\mathbb{O}_{1}^{t} & \mathbb{O} \\ 
\mathbb{O} & \mathbb{O}_{2}^{t}%
\end{array}%
\right) \left\langle \mathbf{e}_{1},\mathbf{e}_{2}\right\vert =\left\vert 
\mathbf{e}_{1},\mathbf{e}_{2}\right\rangle \left( 
\begin{array}{cc}
\mathbb{O}_{1}\mathbb{O}_{1}^{t} & \mathbb{O} \\ 
\mathbb{O} & -\mathbb{O}_{2}\mathbb{O}_{2}^{t}%
\end{array}%
\right) \left\langle \mathbf{e}_{1},\mathbf{e}_{2}\right\vert \\
=\left\vert \mathbf{e}_{1},\mathbf{e}_{2}\right\rangle \left( 
\begin{array}{cc}
\mathbb{I} & \mathbb{O} \\ 
\mathbb{O} & -\mathbb{I}%
\end{array}%
\right) \left\langle \mathbf{e}_{1},\mathbf{e}_{2}\right\vert \quad
;o_{1},o_{2}\in O\left( r,\mathbb{R}\right)
\end{gather}%
and complexifications 
\begin{equation}
J_{o\prime }=\left\vert o\prime _{1}\mathbf{e}_{1},o\prime _{2}\mathbf{e}%
_{2}\right\rangle \left( 
\begin{array}{cc}
\mathbb{O} & -\mathbb{I} \\ 
\mathbb{I} & \mathbb{O}%
\end{array}%
\right) \left\langle o\prime _{1}\mathbf{e}_{1},o\prime _{2}\mathbf{e}%
_{2}\right\vert ;o\prime _{1},o\prime _{2}\in O\left( r,\mathbb{R}\right) ,
\end{equation}%
can always be expressed in a common conb of $V_{1}$. and $V_{2}$ as 
\begin{align}
J& =\left\vert o\prime _{1}\mathbf{e}_{1},o\prime _{2}\mathbf{e}%
_{2}\right\rangle \left( 
\begin{array}{cc}
\mathbb{O} & -\mathbb{I} \\ 
\mathbb{I} & \mathbb{O}%
\end{array}%
\right) \left\langle o\prime _{1}\mathbf{e}_{1},o\prime _{2}\mathbf{e}%
_{2}\right\vert =\left\vert \mathbf{h}_{1},\mathbf{h}_{2}\right\rangle
\left( 
\begin{array}{cc}
\mathbb{O} & -\mathbb{I} \\ 
\mathbb{I} & \mathbb{O}%
\end{array}%
\right) \left\langle \mathbf{h}_{1},\mathbf{h}_{2}\right\vert =\left\vert 
\mathbf{h}_{2}\right\rangle \left\langle \mathbf{h}_{1}\right\vert
-\left\vert \mathbf{h}_{1}\right\rangle \left\langle \mathbf{h}%
_{2}\right\vert \\
& =\left\vert \mathbf{e}_{1},\mathbf{e}_{2}\right\rangle \left( 
\begin{array}{cc}
\mathbb{O}_{1} & \mathbb{O} \\ 
\mathbb{O} & \mathbb{O}_{2}%
\end{array}%
\right) \left( 
\begin{array}{cc}
\mathbb{O} & -\mathbb{I} \\ 
\mathbb{I} & \mathbb{O}%
\end{array}%
\right) \left( 
\begin{array}{cc}
\mathbb{O}_{1}^{t} & \mathbb{O} \\ 
\mathbb{O} & \mathbb{O}_{2}^{t}%
\end{array}%
\right) \left\langle \mathbf{e}_{1},\mathbf{e}_{2}\right\vert =\left\vert 
\mathbf{e}_{1},\mathbf{e}_{2}\right\rangle \left( 
\begin{array}{cc}
\mathbb{O} & -\mathbb{O}_{1}\mathbb{O}_{2}^{t} \\ 
\mathbb{O}_{2}\mathbb{O}_{1}^{t} & \mathbb{O}%
\end{array}%
\right) \left\langle \mathbf{e}_{1},\mathbf{e}_{2}\right\vert \\
& =\left\vert \mathbf{e}_{1},\mathbf{e}_{2}\right\rangle \left( 
\begin{array}{cc}
\mathbb{O} & -\mathbb{O}_{3}^{t} \\ 
\mathbb{O}_{3}^{t} & \mathbb{O}%
\end{array}%
\right) \left\langle \mathbf{e}_{1},\mathbf{e}_{2}\right\vert =\left\vert 
\mathbf{e}_{1},\mathbf{e}_{2}\right\rangle \left( 
\begin{array}{cc}
\mathbb{O} & -\mathbb{O}_{3}^{t} \\ 
\mathbb{O}_{3}^{t} & \mathbb{O}%
\end{array}%
\right) \left\langle \mathbf{e}_{1},\mathbf{e}_{2}\right\vert =\left\vert 
\mathbf{e}_{2}\right\rangle \mathbb{O}_{3}\left\langle \mathbf{e}%
_{1}\right\vert -\left\vert \mathbf{e}_{1}\right\rangle \mathbb{O}%
_{3}^{t}\left\langle \mathbf{e}_{2}\right\vert \\
& =\left\vert \mathbf{c}_{2}\right\rangle \left\langle \mathbf{c}%
_{1}\right\vert -\left\vert \mathbf{c}_{1}\right\rangle \left\langle \mathbf{%
c}_{2}\right\vert =\left\vert \mathbf{c}_{1},\mathbf{c}_{2}\right\rangle
\left( 
\begin{array}{cc}
\mathbb{O} & -\mathbb{I} \\ 
\mathbb{I} & \mathbb{O}%
\end{array}%
\right) \left\langle \mathbf{c}_{1},\mathbf{c}_{2}\right\vert ,
\end{align}
\end{proposition}

i.e. as 
\begin{equation}
\sigma=\left\vert \mathbf{c}_{1}\right\rangle \left\langle \mathbf{c}%
_{1}\right\vert -\left\vert \mathbf{c}_{2}\right\rangle \left\langle \mathbf{%
c}_{2}\right\vert =\left\vert \mathbf{c}_{1},\mathbf{c}_{2}\right\rangle
\left( 
\begin{array}{cc}
\mathbb{I} & \mathbb{O} \\ 
\mathbb{O} & -\mathbb{I}%
\end{array}
\right) \left\langle \mathbf{c}_{1},\mathbf{c}_{2}\right\vert
\end{equation}
and 
\begin{equation}
J=\left\vert \mathbf{c}_{2}\right\rangle \left\langle \mathbf{c}%
_{1}\right\vert -\left\vert \mathbf{c}_{1}\right\rangle \left\langle \mathbf{%
c}_{2}\right\vert =\left\vert \mathbf{c}_{1},\mathbf{c}_{2}\right\rangle
\left( 
\begin{array}{cc}
\mathbb{O} & -\mathbb{I} \\ 
\mathbb{I} & \mathbb{O}%
\end{array}
\right) \left\langle \mathbf{c}_{1},\mathbf{c}_{2}\right\vert .
\end{equation}
\emph{Thus wlog one can always assume such forms for \underline{\emph{%
equivalent}} complex conjugations and complexifications. }

The following complex conjugations and complexifications are \emph{compatible%
} but do not in general maintain the same subspace decomposition of $V,$
i.e. $V_{1}$ and $V_{2},$

\begin{enumerate}
\item 
\begin{align}
\left\vert \mathbf{e}_{1},\mathbf{e}_{2}\right\rangle \mathbb{Q}\left( 
\begin{array}{cc}
\mathbb{I} & \mathbb{O} \\ 
\mathbb{O} & -\mathbb{I}%
\end{array}
\right) \mathbb{Q}^{t}\left\langle \mathbf{e}_{1},\mathbf{e}_{2}\right\vert
& =\left\vert \mathbf{f}_{1},\mathbf{f}_{2}\right\rangle \left( 
\begin{array}{cc}
\mathbb{I} & \mathbb{O} \\ 
\mathbb{O} & -\mathbb{I}%
\end{array}
\right) \left\langle \mathbf{f}_{1},\mathbf{f}_{2}\right\vert ;\quad \mathbb{%
Q}\left( 
\begin{array}{cc}
\mathbb{I} & \mathbb{O} \\ 
\mathbb{O} & -\mathbb{I}%
\end{array}
\right) \mathbb{Q}^{t}=\left( 
\begin{array}{cc}
\mathbb{I} & \mathbb{O} \\ 
\mathbb{O} & -\mathbb{I}%
\end{array}
\right) \\
\left\vert \mathbf{e}_{1},\mathbf{e}_{2}\right\rangle \mathbb{Q}\left( 
\begin{array}{cc}
\mathbb{O} & -\mathbb{I} \\ 
\mathbb{I} & \mathbb{O}%
\end{array}
\right) \mathbb{Q}^{t}\left\langle \mathbf{e}_{1},\mathbf{e}_{2}\right\vert
& =\left\vert \mathbf{f}_{1},\mathbf{f}_{2}\right\rangle \left( 
\begin{array}{cc}
\mathbb{O} & -\mathbb{I} \\ 
\mathbb{I} & \mathbb{O}%
\end{array}
\right) \left\langle \mathbf{f}_{1},\mathbf{f}_{2}\right\vert ;\quad \mathbb{%
Q}\left( 
\begin{array}{cc}
\mathbb{O} & -\mathbb{I} \\ 
\mathbb{I} & \mathbb{O}%
\end{array}
\right) \mathbb{Q}^{t}\neq\left( 
\begin{array}{cc}
\mathbb{O} & -\mathbb{I} \\ 
\mathbb{I} & \mathbb{O}%
\end{array}
\right)
\end{align}
i.e. 
\begin{equation}
\mathbb{Q}\in O\left( r,r,\mathbb{R}\right) ,\mathbb{Q}\notin Sp\left( 2r,%
\mathbb{R}\right) ,
\end{equation}

\item 
\begin{align}
\left\vert \mathbf{e}_{1},\mathbf{e}_{2}\right\rangle \mathbb{Q}\left( 
\begin{array}{cc}
\mathbb{I} & \mathbb{O} \\ 
\mathbb{O} & -\mathbb{I}%
\end{array}
\right) \mathbb{Q}^{t}\left\langle \mathbf{e}_{1},\mathbf{e}_{2}\right\vert
& =\left\vert \mathbf{f}_{1},\mathbf{f}_{2}\right\rangle \left( 
\begin{array}{cc}
\mathbb{I} & \mathbb{O} \\ 
\mathbb{O} & -\mathbb{I}%
\end{array}
\right) \left\langle \mathbf{f}_{1},\mathbf{f}_{2}\right\vert ;\quad \mathbb{%
Q}\left( 
\begin{array}{cc}
\mathbb{I} & \mathbb{O} \\ 
\mathbb{O} & -\mathbb{I}%
\end{array}
\right) \mathbb{Q}^{t}\neq\left( 
\begin{array}{cc}
\mathbb{I} & \mathbb{O} \\ 
\mathbb{O} & -\mathbb{I}%
\end{array}
\right) \\
\left\vert \mathbf{e}_{1},\mathbf{e}_{2}\right\rangle \mathbb{Q}\left( 
\begin{array}{cc}
\mathbb{O} & -\mathbb{I} \\ 
\mathbb{I} & \mathbb{O}%
\end{array}
\right) \mathbb{Q}^{t}\left\langle \mathbf{e}_{1},\mathbf{e}_{2}\right\vert
& =\left\vert \mathbf{f}_{1},\mathbf{f}_{2}\right\rangle \left( 
\begin{array}{cc}
\mathbb{O} & -\mathbb{I} \\ 
\mathbb{I} & \mathbb{O}%
\end{array}
\right) \left\langle \mathbf{f}_{1},\mathbf{f}_{2}\right\vert ;\quad \mathbb{%
Q}\left( 
\begin{array}{cc}
\mathbb{O} & -\mathbb{I} \\ 
\mathbb{I} & \mathbb{O}%
\end{array}
\right) \mathbb{Q}^{t}=\left( 
\begin{array}{cc}
\mathbb{O} & -\mathbb{I} \\ 
\mathbb{I} & \mathbb{O}%
\end{array}
\right)
\end{align}
i.e. 
\begin{equation}
\mathbb{Q}\notin O\left( r,r,\mathbb{R}\right) ,\mathbb{Q}\in Sp\left( 2r,%
\mathbb{R}\right) .
\end{equation}
\end{enumerate}

Transformations $\mathbb{Q}$ that \emph{do maintain} the same subspace
structure and produce equivalent and compatible conjugations and
complexifications have the properties 
\begin{equation}
\mathbb{Q}\left( 
\begin{array}{cc}
\mathbb{I} & \mathbb{O} \\ 
\mathbb{O} & -\mathbb{I}%
\end{array}%
\right) \mathbb{Q}^{t}=\left( 
\begin{array}{cc}
\mathbb{I} & \mathbb{O} \\ 
\mathbb{O} & -\mathbb{I}%
\end{array}%
\right)
\end{equation}%
and 
\begin{equation}
\mathbb{Q}\left( 
\begin{array}{cc}
\mathbb{O} & -\mathbb{I} \\ 
\mathbb{I} & \mathbb{O}%
\end{array}%
\right) \mathbb{Q}^{t}=\left( 
\begin{array}{cc}
\mathbb{O} & -\mathbb{I} \\ 
\mathbb{I} & \mathbb{O}%
\end{array}%
\right)
\end{equation}%
i.e. 
\begin{equation}
\mathbb{Q}\in O\left( r,r,\mathbb{R}\right) \cap Sp\left( 2r,\mathbb{R}%
\right) .
\end{equation}%
$\lceil $ Is $O\left( r,r,\mathbb{R}\right) $ similiar to $Sp\left( 2r,%
\mathbb{R}\right) ?$ If it was then 
\begin{align}
\left( 
\begin{array}{cc}
\mathbb{T}_{1} & \mathbb{T}_{2} \\ 
\mathbb{T}_{3} & \mathbb{T}_{4}%
\end{array}%
\right) \left( 
\begin{array}{cc}
\mathbb{O} & -\mathbb{I} \\ 
\mathbb{I} & \mathbb{O}%
\end{array}%
\right) \left( 
\begin{array}{cc}
\mathbb{T}_{1}^{-1} & \mathbb{T}_{3}^{-1} \\ 
\mathbb{T}_{2}^{-1} & \mathbb{T}_{4}^{-1}%
\end{array}%
\right) & =\left( 
\begin{array}{cc}
\mathbb{I} & \mathbb{O} \\ 
\mathbb{O} & -\mathbb{I}%
\end{array}%
\right) \\
\left( 
\begin{array}{cc}
\mathbb{T}_{1} & \mathbb{T}_{2} \\ 
\mathbb{T}_{3} & \mathbb{T}_{4}%
\end{array}%
\right) \left( 
\begin{array}{cc}
\mathbb{O} & -\mathbb{I} \\ 
\mathbb{I} & \mathbb{O}%
\end{array}%
\right) \left( 
\begin{array}{cc}
\mathbb{T}_{1}^{-1} & \mathbb{T}_{3}^{-1} \\ 
\mathbb{T}_{2}^{-1} & \mathbb{T}_{4}^{-1}%
\end{array}%
\right) & =\left( 
\begin{array}{cc}
\mathbb{T}_{2} & -\mathbb{T}_{1} \\ 
\mathbb{T}_{4} & -\mathbb{T}_{3}%
\end{array}%
\right) \left( 
\begin{array}{cc}
\mathbb{T}_{1}^{-1} & \mathbb{T}_{3}^{-1} \\ 
\mathbb{T}_{2}^{-1} & \mathbb{T}_{4}^{-1}%
\end{array}%
\right) =\left( 
\begin{array}{cc}
\mathbb{T}_{2}\mathbb{T}_{1}^{-1}-\mathbb{T}_{1}\mathbb{T}_{2}^{-1} & 
\mathbb{T}_{2}\mathbb{T}_{3}^{-1}-\mathbb{T}_{1}\mathbb{T}_{4}^{-1} \\ 
\mathbb{T}_{4}\mathbb{T}_{1}^{-1}-\mathbb{T}_{3}\mathbb{T}_{2}^{-1} & 
\mathbb{T}_{4}\mathbb{T}_{3}^{-1}-\mathbb{T}_{3}\mathbb{T}_{4}^{-1}%
\end{array}%
\right)
\end{align}%
\begin{align}
\mathbb{T}_{2}\mathbb{T}_{1}^{-1}-\mathbb{T}_{1}\mathbb{T}_{2}^{-1}& =%
\mathbb{I} \\
\mathbb{T}_{4}\mathbb{T}_{3}^{-1}-\mathbb{T}_{3}\mathbb{T}_{4}^{-1}& =-%
\mathbb{I} \\
\mathbb{T}_{2}\mathbb{T}_{3}^{-1}-\mathbb{T}_{1}\mathbb{T}_{4}^{-1}& =%
\mathbb{O} \\
\mathbb{T}_{4}\mathbb{T}_{1}^{-1}-\mathbb{T}_{3}\mathbb{T}_{2}^{-1}& =%
\mathbb{O}
\end{align}%
so there is no transformation $\mathbb{T}$ and the answer is \emph{no}$%
.\rfloor $

\begin{lemma}
If $o\in O\left( r,\mathbb{R}\right) +O\left( r,\mathbb{R}\right) \subset
O\left( 2r,\mathbb{R}\right) $ then $o\sigma o^{t}$ is \emph{equivalent} to
the complex conjugation $\sigma$ and $oJo^{t}$ is \emph{equivalent} to the
complexification $J$. ($o=o_{1}+o_{2}=\widetilde{o_{1}}\widetilde{o_{2}},$
where $o_{1},o_{2}\in O\left( r,\mathbb{R}\right) $)`
\end{lemma}

$\lceil$%
\begin{align}
o_{1} & :V_{1}\rightarrow V_{1} \\
o_{2} & :V_{2}\rightarrow V_{2}
\end{align}%
\begin{align}
o\sigma o^{t} & =o_{1}\left\vert \mathbf{e}_{1}\right\rangle \left\langle 
\mathbf{e}_{1}\right\vert o_{1}^{t}-o_{2}\left\vert \mathbf{e}%
_{2}\right\rangle \left\langle \mathbf{e}_{2}\right\vert
o_{2}^{t}=\left\vert \mathbf{f}_{1}\right\rangle \left\langle \mathbf{f}%
_{1}\right\vert -\left\vert \mathbf{f}_{2}\right\rangle \left\langle \mathbf{%
f}_{2}\right\vert =\left\vert \mathbf{e}_{1}\right\rangle \mathcal{R}%
_{e_{1}}\left( o_{1}o_{1}^{t}\right) \left\langle \mathbf{e}_{1}\right\vert
-\left\vert \mathbf{e}_{2}\right\rangle \mathcal{R}_{e_{2}}\left(
o_{2}o_{2}^{t}\right) \left\langle \mathbf{e}_{2}\right\vert =\left\vert 
\mathbf{e}_{1}\right\rangle \left\langle \mathbf{e}_{1}\right\vert
-\left\vert \mathbf{e}_{2}\right\rangle \left\langle \mathbf{e}%
_{2}\right\vert \\
oJo^{t} & =o_{2}\left\vert \mathbf{e}_{2}\right\rangle \left\langle \mathbf{e%
}_{1}\right\vert o_{1}^{t}-o_{1}\left\vert \mathbf{e}_{1}\right\rangle
\left\langle \mathbf{e}_{2}\right\vert o_{2}^{t}=\left\vert \mathbf{f}%
_{2}\right\rangle \left\langle \mathbf{f}_{1}\right\vert -\left\vert \mathbf{%
f}_{1}\right\rangle \left\langle \mathbf{f}_{2}\right\vert =\left\vert 
\mathbf{e}_{2}\right\rangle \mathcal{R}_{e_{2}}\left( o_{2}\right) \mathcal{R%
}_{e_{1}}\left( o_{1}^{t}\right) \left\langle \mathbf{e}_{1}\right\vert
-\left\vert \mathbf{e}_{1}\right\rangle \mathcal{R}_{e_{1}}\left(
o_{1}\right) \mathcal{R}_{e_{2}}\left( o_{2}^{t}\right) \left\langle \mathbf{%
e}_{2}\right\vert
\end{align}
$\rfloor$

\begin{lemma}
If $o\in O\left( r,\mathbb{R}\right) $ inducing $\widetilde{o}=o+o\in
O\left( 2r,\mathbb{R}\right) $ then $\widetilde{o}\sigma\widetilde{o}^{t}$
and $\sigma$ and $\widetilde{o}J\widetilde{o}^{t}$ and $J$ are the same
complex conjugations and complexifications.
\end{lemma}

\begin{remark}
$O\left( r,\mathbb{R}\right) +O\left( r,\mathbb{R}\right) $ and $\widetilde{%
O\left( r,\mathbb{R}\right) }=\left\{ \left. \widetilde{o}\right\vert o\in
O\left( r,\mathbb{R}\right) \right\} $ are proper subgroups of $O\left( 2r,%
\mathbb{R}\right) $ and $O\left( r,\mathbb{R}\right) +O\left( r,\mathbb{R}%
\right) \supset\widetilde{O\left( r,\mathbb{R}\right) }.$
\end{remark}

\begin{theorem}
Every complex conjugation, $\sigma_{\left[ e\right] }\in\mathcal{B}%
_{anti}\left( \mathfrak{H}\right) ,$ of $\mathfrak{H}$ defines a real direct
sum decomposition of $\mathfrak{H}$, a real Hilbert space $\mathfrak{H}_{%
\mathbb{R}}$ and a complexification $J_{\left[ e\right] }$ of $\mathfrak{H}_{%
\mathbb{R}}$ that anti commutes with $\sigma_{\left[ e\right] }$ i.e. 
\begin{equation}
\left\{ \sigma_{\left[ e\right] },J_{\left[ e\right] }\right\} =0
\end{equation}
and whose complexification via $J_{\left[ e\right] }$ produces $\mathfrak{H.}
$
\end{theorem}

\begin{proof}
$\sigma_{\left[ e\right] }$ produces the \emph{real direct} sum
decomposition of $\mathfrak{H}$ into real eigen spaces associated with the
eigenvalues $\left\{ 1,-1\right\} $ of $\sigma_{\left[ e\right] }$ i.e.%
\begin{equation}
\mathfrak{H}=V_{1e}\oplus_{\mathbb{R}}iV_{2e}\cong V_{+\sigma_{\left[ e%
\right] }}\oplus_{\mathbb{R}}iV_{-\sigma_{\left[ e\right] }},
\end{equation}

where 
\begin{equation}
V_{1e}\equiv V_{+\sigma_{\left[ e\right] }}\text{ and }V_{2e}\equiv
V_{-\sigma_{\left[ e\right] }}
\end{equation}

and one has the following \emph{real} isomorphisms 
\begin{equation}
V_{1e}\cong V_{+\sigma_{\left[ e\right] }}\cong iV_{-\sigma_{\left[ e\right]
}}\cong V_{-\sigma_{\left[ e\right] }},
\end{equation}

that define $2r$ dimensional real Hilbert spaces $\mathfrak{H}_{\mathbb{R}}$
given by 
\begin{equation}
\mathfrak{H}_{\mathbb{R}}=V_{1e}\oplus_{\mathbb{R}}V_{2e}=V_{+\sigma_{\left[
e\right] }}\oplus_{\mathbb{R}}V_{-\sigma_{\left[ e\right] }}\cong
V_{+\sigma_{\left[ e\right] }}\oplus_{\mathbb{R}}V_{+\sigma_{\left[ e\right]
}}\cong V_{-\sigma_{\left[ e\right] }}\oplus_{\mathbb{R}}V_{-\sigma_{\left[ e%
\right] }}.
\end{equation}

The basis of $V_{1e}$ $\subset\mathfrak{H}_{\mathbb{R}}$ inherited from $%
\mathfrak{H}$ is 
\begin{equation}
\left\vert \mathbf{\varepsilon}_{1}\right\rangle \equiv\left\vert \left( 
\mathbf{e},\mathbf{0}\right) \right\rangle ,
\end{equation}

where $\left\vert \mathbf{e}\right\rangle \in\left[ \mathbf{e}\right] ,$
while the corresponding basis of $V_{2e}\subset\mathfrak{H}_{\mathbb{R}}$
inherited

from $\mathfrak{H}$ is 
\begin{equation}
\left\vert \mathbf{\varepsilon}_{2}\right\rangle \equiv\left\vert \left( 
\mathbf{0,e}\right) \right\rangle .
\end{equation}

As $V_{1e}$ equals $V_{2e}$ i.e. 
\begin{equation}
V_{1e}=V_{2e}
\end{equation}

not just isomorphic i.e. 
\begin{equation}
V_{1e}\cong V_{2e}
\end{equation}

one observes that 
\begin{equation}
V_{1e}\oplus_{\mathbb{R}}V_{2e}=V_{1e}\oplus_{\mathbb{R}}V_{1e}\neq V_{1e}+_{%
\mathbb{R}}V_{2e}=V_{1e}+_{\mathbb{R}}V_{1e}.
\end{equation}

The real symmetric inner product, $\left\langle \left. \cdot\right\vert
\cdot\right\rangle _{\mathfrak{H}_{\mathbb{R}}}$ in $\mathfrak{H}_{\mathbb{R}%
}$ is defined by 
\begin{equation}
\left\langle \left. \eta_{1}\right\vert \eta_{2}\right\rangle _{\mathfrak{H}%
_{\mathbb{R}}}=\left\langle \left. x_{1}\right\vert x_{2}\right\rangle
_{V_{1e}}+\left\langle \left. y_{1}\right\vert y_{2}\right\rangle
_{V_{2e}};x_{1},x_{2}\in V_{1e},y_{1},y_{2}\in V_{2e},
\end{equation}

where $\left\langle \left. {}\right\vert \right\rangle _{V_{1e}}$ and $%
\left\langle \left. {}\right\vert \right\rangle _{V_{2e}}$ are real
symmetric inner products in $V_{1e},V_{2e}$. The invariance group of the
inner product $\left\langle \left. \cdot \right\vert \cdot \right\rangle _{%
\mathfrak{H}_{\mathbb{R}}}$ is

$O\left( \mathfrak{H}_{\mathbb{R}}\right) \equiv O\left( 2r,\mathbb{R}%
\right) ,$which is defined by 
\begin{equation}
O\left( 2r,\mathbb{R}\right) =\left\{ \left. g\in GL\left( 2r,\mathbb{R}%
\right) \right\vert \left\{ \left. g\eta_{1}\right\vert g\eta_{2}\right\} _{%
\mathfrak{H}_{\mathbb{R}}}=\left\{ \left. \eta_{1}\right\vert \eta
_{2}\right\} _{\mathfrak{H}_{\mathbb{R}}}\forall\eta_{1},\eta_{2}\in%
\mathfrak{H}_{\mathbb{R}}\right\} \equiv g^{t}g=I_{2r}\equiv I_{\mathfrak{H}%
}.
\end{equation}

The representation of a vector $\eta \in \mathfrak{H}_{\mathbb{R}}$ wrt the
basis $\left\vert \left( \mathbf{e},\mathbf{o}\right) ,\left( \mathbf{o},%
\mathbf{e}\right) \right\rangle $ $\equiv \left\vert \mathbf{\varepsilon }%
_{1},\mathbf{\varepsilon }_{2}\right\rangle $ is 
\begin{equation}
\left\vert \eta \right\rangle =\left\vert \mathbf{\varepsilon }_{1},\mathbf{%
\varepsilon }_{2}\right\rangle \mathcal{R}_{e}\left( \eta \right)
=\left\vert \mathbf{\varepsilon }_{1},\mathbf{\varepsilon }_{2}\right\rangle
\left( 
\begin{array}{c}
\mathbf{x} \\ 
\mathbf{y}%
\end{array}%
\right) .
\end{equation}

$\lceil$ Note that at one extreme of the direct sum 
\begin{equation}
\mathbf{e}_{1}=\mathbf{e}_{2}
\end{equation}

and at the other 
\begin{equation}
\mathbf{e}_{1}\cap\mathbf{e}_{2}=\phi.
\end{equation}

$\rfloor$

The \emph{complex conjugation} $\sigma_{\left[ e\right] }\in\mathcal{B}%
_{anti}\left( \mathfrak{H}\right) $ can be represented by the \emph{complex
conjugation} $\sigma_{\left[ e\right] }$ in $\mathcal{B}\left( \mathfrak{H}_{%
\mathbb{R}}\right) $ as 
\begin{equation}
\sigma_{\left[ e\right] }=\left\vert \mathbf{\varepsilon}_{1}\right\rangle
\left\langle \mathbf{\varepsilon}_{1}\right\vert -\left\vert \mathbf{%
\varepsilon}_{2}\right\rangle \left\langle \mathbf{\varepsilon}%
_{2}\right\vert ,
\end{equation}

where $\left\vert \mathbf{e}\right\rangle $ is the basis of both $V_{1e}$
and $V_{2e},$ while $\left\vert \mathbf{\varepsilon }_{1}\right\rangle $ and 
$\left\vert \mathbf{\varepsilon }_{2}\right\rangle $ are the bases of $%
V_{1e} $ and $V_{2e}$ respectively in the direct sum and we

have the maintained the same symbol for the complex conjugation in both
cases as $\mathcal{B}\left( \mathfrak{H}_{\mathbb{R}}\right) \supset 
\mathcal{B}_{anti}\left( \mathfrak{H}\right) $ as real Banach spaces of

operators acting in $\mathfrak{H}_{\mathbb{R}}$ and $\mathfrak{H}$. $\sigma
_{\left[ e\right] }$ has the representation%
\begin{equation}
\sigma _{\left[ e\right] }\left\vert \mathbf{\varepsilon }_{1},\mathbf{%
\varepsilon }_{2}\right\rangle =\left\vert \mathbf{\varepsilon }_{1},\mathbf{%
\varepsilon }_{2}\right\rangle \mathcal{R}_{e}\left( \sigma _{\left[ e\right]
}\right) =\left\vert \mathbf{\varepsilon }_{1},\mathbf{\varepsilon }%
_{2}\right\rangle \left( 
\begin{array}{cc}
\mathbb{I} & \mathbb{O} \\ 
\mathbb{O} & -\mathbb{I}%
\end{array}%
\right)
\end{equation}

and has invariance group $\sigma_{\left[ e\right] }$ is $O\left( r,r,\mathbb{%
R}\right) .$

One can define a \emph{complexification} $J_{\left[ e\right] }\in \mathcal{B}%
\left( \mathfrak{H}_{\mathbb{R}}\right) $ 
\begin{align}
J_{\left[ e\right] } & =-J_{\left[ e\right] }^{t} \\
J_{\left[ e\right] }^{2} & =-J_{\left[ e\right] }J_{\left[ e\right]
}^{t}=-J_{\left[ e\right] }^{t}J_{\left[ e\right] }=J_{\left[ e\right]
}^{2t}=-I_{\mathfrak{H}_{\mathbb{R}}},
\end{align}

in $\mathfrak{H}_{\mathbb{R}}$ by 
\begin{equation}
J_{\left[ e\right] }=\left\vert \mathbf{\varepsilon}_{2}\right\rangle
\left\langle \mathbf{\varepsilon}_{1}\right\vert -\left\vert \mathbf{%
\varepsilon}_{1}\right\rangle \left\langle \mathbf{\varepsilon}%
_{2}\right\vert ,
\end{equation}

that has the representation%
\begin{equation}
J_{\left[ e\right] }\left\vert \mathbf{\varepsilon}_{1},\mathbf{\varepsilon }%
_{2}\right\rangle =\left\vert \mathbf{\varepsilon}_{1},\mathbf{\varepsilon }%
_{2}\right\rangle \mathcal{R}_{e}\left( J_{\left[ e\right] }\right)
=\left\vert \mathbf{\varepsilon}_{1},\mathbf{\varepsilon}_{2}\right\rangle
\left( 
\begin{array}{cc}
\mathbb{O} & -\mathbb{I} \\ 
\mathbb{I} & \mathbb{O}%
\end{array}
\right) ,
\end{equation}

and invariance group $Sp\left( 2r,\mathbb{R}\right) .$

The complexification $J_{\left[ e\right] }$ and conugation $\sigma_{\left[ e%
\right] }$ are \emph{compatible,} i.e. 
\begin{equation}
\left\{ \sigma_{\left[ e\right] },J_{\left[ e\right] }\right\} =0.
\end{equation}

One can thus express $\sigma_{\left[ e\right] }$ and $J_{\left[ e\right] }$
as%
\begin{align}
\sigma_{\left[ e\right] } & =\left\vert \mathbf{\varepsilon}_{1},\mathbf{%
\varepsilon}_{2}\right\rangle \mathcal{R}_{e}\left( \sigma_{\left[ e\right]
}\right) \left\langle \mathbf{\varepsilon}_{1},\mathbf{\varepsilon}%
_{2}\right\vert =\left\vert \mathbf{\varepsilon}_{1},\mathbf{\varepsilon}%
_{2}\right\rangle \left( 
\begin{array}{cc}
\mathbb{I} & \mathbb{O} \\ 
\mathbb{O} & -\mathbb{I}%
\end{array}
\right) \left\langle \mathbf{\varepsilon}_{1},\mathbf{\varepsilon}%
_{2}\right\vert \\
& =\left( 
\begin{array}{cc}
\left\vert \mathbf{\varepsilon}_{1}\right\rangle & -\left\vert \mathbf{%
\varepsilon}_{2}\right\rangle%
\end{array}
\right) \left( 
\begin{array}{c}
\left\langle \mathbf{\varepsilon}_{1}\right\vert \\ 
\left\langle \mathbf{\varepsilon}_{2}\right\vert%
\end{array}
\right) =\left\vert \mathbf{\varepsilon}_{1}\right\rangle \left\langle 
\mathbf{\varepsilon}_{1}\right\vert -\left\vert \mathbf{\varepsilon}%
_{2}\right\rangle \left\langle \mathbf{\varepsilon}_{2}\right\vert
\end{align}

and%
\begin{align}
J_{\left[ e\right] } & =\left\vert \mathbf{\varepsilon}_{1},\mathbf{%
\varepsilon}_{2}\right\rangle \mathcal{R}_{e}\left( J_{\left[ e\right]
}\right) \left\langle \mathbf{\varepsilon}_{1},\mathbf{\varepsilon }%
_{2}\right\vert =\left\vert \mathbf{\varepsilon}_{1},\mathbf{\varepsilon}%
_{2}\right\rangle \left( 
\begin{array}{cc}
\mathbb{O} & -\mathbb{I} \\ 
\mathbb{I} & \mathbb{O}%
\end{array}
\right) \left\langle \mathbf{\varepsilon}_{1},\mathbf{\varepsilon}%
_{2}\right\vert \\
& =\left( 
\begin{array}{cc}
\mathbb{O} & -\left\vert \mathbf{\varepsilon}_{1}\right\rangle \\ 
\left\vert \mathbf{\varepsilon}_{2}\right\rangle & \mathbb{O}%
\end{array}
\right) \left( 
\begin{array}{c}
\left\langle \mathbf{\varepsilon}_{1}\right\vert \\ 
\left\langle \mathbf{\varepsilon}_{2}\right\vert%
\end{array}
\right) =\left\vert \mathbf{\varepsilon}_{2}\right\rangle \left\langle 
\mathbf{\varepsilon}_{1}\right\vert -\left\vert \mathbf{\varepsilon}%
_{1}\right\rangle \left\langle \mathbf{\varepsilon}_{2}\right\vert .
\end{align}

One can see that $\sigma_{\left[ e\right] }$ and $J_{\left[ e\right] }$ are
invariant to transformations $o\in O\left( r,\mathbb{R}\right) =$ $O\left(
r,r,\mathbb{R}\right) \cap Sp\left( 2r,\mathbb{R}\right) $ if 
\begin{equation}
\left\vert \mathbf{\varepsilon}_{1},\mathbf{\varepsilon}_{2}\right\rangle
=\left\vert \left( \mathbf{e},\mathbf{o}\right) ,\left( \mathbf{o},\mathbf{e}%
\right) \right\rangle ,
\end{equation}

i.e. 
\begin{equation}
o\left( \left\vert \mathbf{\varepsilon}_{1}\right\rangle \left\langle 
\mathbf{\varepsilon}_{1}\right\vert -\left\vert \mathbf{\varepsilon}%
_{2}\right\rangle \left\langle \mathbf{\varepsilon}_{2}\right\vert \right)
o^{t}=\left\vert \mathbf{\varepsilon}_{1}\right\rangle \left\langle \mathbf{%
\varepsilon}_{1}\right\vert -\left\vert \mathbf{\varepsilon}%
_{2}\right\rangle \left\langle \mathbf{\varepsilon}_{2}\right\vert
\end{equation}

and%
\begin{equation}
o\left( \left\vert \mathbf{\varepsilon}_{2}\right\rangle \left\langle 
\mathbf{\varepsilon}_{1}\right\vert -\left\vert \mathbf{\varepsilon}%
_{1}\right\rangle \left\langle \mathbf{\varepsilon}_{2}\right\vert \right)
o^{t}=\left\vert \mathbf{\varepsilon}_{2}\right\rangle \left\langle \mathbf{%
\varepsilon}_{1}\right\vert -\left\vert \mathbf{\varepsilon}%
_{1}\right\rangle \left\langle \mathbf{\varepsilon}_{2}\right\vert ,
\end{equation}

where 
\begin{align}
o\left\vert \mathbf{\varepsilon}_{1},\mathbf{\varepsilon}_{2}\right\rangle &
=\left\vert \mathbf{\varepsilon}_{1},\mathbf{\varepsilon}_{2}\right\rangle
\left( \mathcal{R}_{e}\left( o\right) \oplus\mathcal{R}_{e}\left( o\right)
\right) \\
\left\langle \mathbf{\varepsilon}_{1},\mathbf{\varepsilon}_{2}\right\vert
o^{t} & =\left( \mathcal{R}_{e}\left( o\right) \oplus\mathcal{R}_{e}\left(
o\right) \right) ^{t}\left\langle \mathbf{\varepsilon}_{1},\mathbf{%
\varepsilon}_{2}\right\vert .
\end{align}

As 
\begin{align}
oJ_{\left[ e\right] }o^{t} & =o\left\vert \mathbf{\varepsilon}_{1},\mathbf{%
\varepsilon}_{2}\right\rangle \left( 
\begin{array}{cc}
\mathbb{O} & -\mathbb{I} \\ 
\mathbb{I} & \mathbb{O}%
\end{array}
\right) \left\langle \mathbf{\varepsilon}_{1},\mathbf{\varepsilon}%
_{2}\right\vert o^{t}=\left\vert \mathbf{\varepsilon}_{1},\mathbf{%
\varepsilon }_{2}\right\rangle \left( \mathcal{R}_{e}\left( o\right) \oplus 
\mathcal{R}_{e}\left( o\right) \right) \left( 
\begin{array}{cc}
\mathbb{O} & -\mathbb{I} \\ 
\mathbb{I} & \mathbb{O}%
\end{array}
\right) \left( \mathcal{R}_{e}\left( o\right) \oplus\mathcal{R}_{e}\left(
o\right) \right) ^{t}\left\langle \mathbf{\varepsilon}_{1},\mathbf{%
\varepsilon}_{2}\right\vert \\
& =\left\vert \mathbf{\varepsilon}_{1},\mathbf{\varepsilon}_{2}\right\rangle
\left( 
\begin{array}{cc}
\mathcal{R}_{e}\left( o\right) & \mathbb{O} \\ 
\mathbb{O} & \mathcal{R}_{e}\left( o\right)%
\end{array}
\right) \left( 
\begin{array}{cc}
\mathbb{O} & -\mathbb{I} \\ 
\mathbb{I} & \mathbb{O}%
\end{array}
\right) \left( 
\begin{array}{cc}
\mathcal{R}_{e}\left( o\right) ^{t} & \mathbb{O} \\ 
\mathbb{O} & \mathcal{R}_{e}\left( o\right) ^{t}%
\end{array}
\right) \left\langle \mathbf{\varepsilon}_{1},\mathbf{\varepsilon}%
_{2}\right\vert \\
& =\left\vert \mathbf{\varepsilon}_{1},\mathbf{\varepsilon}_{2}\right\rangle
\left( 
\begin{array}{cc}
\mathbb{O} & -\mathcal{R}_{e}\left( o\right) \\ 
\mathcal{R}_{e}\left( o\right) & \mathbb{O}%
\end{array}
\right) \left( 
\begin{array}{cc}
\mathcal{R}_{e}\left( o\right) ^{t} & \mathbb{O} \\ 
\mathbb{O} & \mathcal{R}_{e}\left( o\right) ^{t}%
\end{array}
\right) \left\langle \mathbf{\varepsilon}_{1},\mathbf{\varepsilon}%
_{2}\right\vert \\
& =\left\vert \mathbf{\varepsilon}_{1},\mathbf{\varepsilon}_{2}\right\rangle
\left( 
\begin{array}{cc}
\mathbb{O} & -\mathcal{R}_{e}\left( o\right) \mathcal{R}_{e}\left( o\right)
^{t} \\ 
\mathcal{R}_{e}\left( o\right) \mathcal{R}_{e}\left( o\right) ^{t} & \mathbb{%
O}%
\end{array}
\right) \left\langle \mathbf{\varepsilon}_{1},\mathbf{\varepsilon}%
_{2}\right\vert =\left\vert \mathbf{\varepsilon}_{1},\mathbf{\varepsilon}%
_{2}\right\rangle \left( 
\begin{array}{cc}
\mathbb{O} & -\mathbb{I} \\ 
\mathbb{I} & \mathbb{O}%
\end{array}
\right) \left\langle \mathbf{\varepsilon}_{1},\mathbf{\varepsilon}%
_{2}\right\vert .
\end{align}

and 
\begin{align}
o\sigma_{\left[ e\right] }o^{t} & =o\left\vert \mathbf{\varepsilon}_{1},%
\mathbf{\varepsilon}_{2}\right\rangle \left( 
\begin{array}{cc}
\mathbb{I} & \mathbb{O} \\ 
\mathbb{O} & -\mathbb{I}%
\end{array}
\right) \left\langle \mathbf{\varepsilon}_{1},\mathbf{\varepsilon}%
_{2}\right\vert o^{t}=\left\vert \mathbf{\varepsilon}_{1},\mathbf{%
\varepsilon }_{2}\right\rangle \left( \mathcal{R}_{e}\left( o\right) \oplus 
\mathcal{R}_{e}\left( o\right) \right) \left( 
\begin{array}{cc}
\mathbb{I} & \mathbb{O} \\ 
\mathbb{O} & -\mathbb{I}%
\end{array}
\right) \left( \mathcal{R}_{e}\left( o\right) \oplus\mathcal{R}_{e}\left(
o\right) \right) ^{t}\left\langle \mathbf{\varepsilon}_{1},\mathbf{%
\varepsilon}_{2}\right\vert \\
& =\left\vert \mathbf{\varepsilon}_{1},\mathbf{\varepsilon}_{2}\right\rangle
\left( 
\begin{array}{cc}
\mathbb{I} & \mathbb{O} \\ 
\mathbb{O} & -\mathbb{I}%
\end{array}
\right) \left\langle \mathbf{\varepsilon}_{1},\mathbf{\varepsilon}%
_{2}\right\vert .
\end{align}

$\lceil$The transformation%
\begin{equation}
\text{ }\left( 
\begin{array}{cc}
\mathbb{T}_{1} & \mathbb{O} \\ 
\mathbb{O} & \mathbb{T}_{2}%
\end{array}
\right) \in O\left( 2r,\mathbb{R}\right)
\end{equation}

but \emph{does not belong to} $O\left( r,\mathbb{R}\right) $ when 
\begin{equation}
\mathbb{T}_{1}\neq\mathbb{T}_{2}
\end{equation}

and 
\begin{equation}
\mathbb{T}_{1},\mathbb{T}_{2}\in O\left( r,\mathbb{R}\right) .
\end{equation}

$\rfloor$

The complexification $J_{\left[ e\right] }$ defines a \emph{symplectic form} 
\begin{equation}
\left\{ \left. \eta_{1}\right\vert \eta_{2}\right\} _{\mathfrak{H}_{\mathbb{R%
}}}=\left\langle \left. \eta_{1}\right\vert J_{\left[ e\right]
}\eta_{2}\right\rangle _{\mathfrak{H}_{\mathbb{R}}}
\end{equation}

on $\mathfrak{H}_{\mathbb{R}}$, which has the $\mathcal{R}_{e}$%
-representation 
\begin{align}
\left\{ \left. \eta_{1}\right\vert \eta_{2}\right\} _{\mathfrak{H}_{\mathbb{R%
}}} & =\left\langle \eta_{1}\left\vert \left\{ \left\vert \mathbf{\varepsilon%
}_{1},\mathbf{\varepsilon}_{2}\right\rangle \left( 
\begin{array}{cc}
\mathbb{O} & -\mathbb{I} \\ 
\mathbb{I} & \mathbb{O}%
\end{array}
\right) \left\langle \mathbf{\varepsilon}_{1},\mathbf{\varepsilon}%
_{2}\right\vert \right\} \right. \eta_{2}\right\rangle \\
& =\mathcal{R}_{e}\left( \eta_{1}\right) ^{t}\left( 
\begin{array}{cc}
\mathbb{O} & -\mathbb{I} \\ 
\mathbb{I} & \mathbb{O}%
\end{array}
\right) \mathcal{R}_{e}\left( \eta_{2}\right) .
\end{align}

The invariance group of the symplectic form $\left\{ \left. \cdot\right\vert
\cdot\right\} _{\mathfrak{H}_{\mathbb{R}}}$ is 
\begin{equation}
Sp\left( 2r,\mathbb{R}\right) =\left\{ \left. g\in GL\left( 2r,\mathbb{R}%
\right) \right\vert \left\{ \left. g\eta_{1}\right\vert g\eta_{2}\right\} _{%
\mathfrak{H}_{\mathbb{R}}}=\left\{ \left. \eta _{1}\right\vert
\eta_{2}\right\} _{\mathfrak{H}_{\mathbb{R}}}\forall\eta _{1},\eta_{2}\in%
\mathfrak{H}_{\mathbb{R}}\right\} ,
\end{equation}

which can be seen to be the same as%
\begin{equation}
Sp\left( 2r,\mathbb{R}\right) =\left\{ \left. g\in GL\left( 2r,\mathbb{R}%
\right) \right\vert \left\langle \left. g\eta_{1}\right\vert J_{\left[ e%
\right] }g\eta_{2}\right\rangle _{\mathfrak{H}_{\mathbb{R}}}=\left\langle
\left. \eta_{1}\right\vert g^{t}J_{\left[ e\right] }g\eta_{2}\right\rangle _{%
\mathfrak{H}_{\mathbb{R}}}=\left\langle \left\{ \left. \eta_{1}\right\vert
\eta_{2}\right\} \right\rangle _{\mathfrak{H}_{\mathbb{R}}}\forall\eta_{1},%
\eta_{2}\in\mathfrak{H}_{\mathbb{R}}\right\} \equiv g^{t}J_{\left[ e\right]
}g=I_{2r}.
\end{equation}

The complexification $J_{\left[ e\right] }\in\mathcal{B}\left( \mathfrak{H}_{%
\mathbb{R}}\right) $ has the effect 
\begin{equation}
J_{\left[ e\right] }\left\vert \mathbf{\varepsilon}_{1}\right\rangle =\left(
\left\vert \mathbf{\varepsilon}_{2}\right\rangle \left\langle \mathbf{%
\varepsilon}_{1}\right\vert -\left\vert \mathbf{\varepsilon}%
_{1}\right\rangle \left\langle \mathbf{\varepsilon}_{2}\right\vert \right)
\left\vert \mathbf{e}_{1}\right\rangle =\left\vert \mathbf{\varepsilon}%
_{2}\right\rangle
\end{equation}

and the basis of $\mathcal{B}\left( \mathfrak{H}_{\mathbb{R}}\right) $ can
be expressed as 
\begin{equation}
\left\vert \mathbf{\varepsilon }_{1},\mathbf{\varepsilon }_{2}\right\rangle
=\left\vert \mathbf{\varepsilon }_{1},J_{\left[ e\right] }\mathbf{%
\varepsilon }_{1}\right\rangle .
\end{equation}

$\lceil$Of course one also has 
\begin{equation}
-J_{\left[ e\right] }\left\vert \mathbf{\varepsilon}_{2}\right\rangle
=\left\vert \mathbf{\varepsilon}_{1}\right\rangle
\end{equation}

so the basis could also be expressed as 
\begin{equation}
\left\vert \mathbf{\varepsilon}_{1},\mathbf{\varepsilon}_{2}\right\rangle
=\left\vert -J_{\left[ e\right] }\mathbf{\varepsilon}_{2},\mathbf{\varepsilon%
}_{2}\right\rangle .
\end{equation}

$\rfloor$

Thus any anti-unitary complex conjugation $\sigma_{\left[ e\right]
}\in\left\{ 
\begin{array}{c}
\mathcal{B}_{anti}\left( \mathfrak{H}\right) \\ 
\mathcal{B}\left( \mathfrak{H}_{\mathbb{R}}\right)%
\end{array}
\right. $ allows a complex Hilbert space $\mathfrak{H}_{e}$ to be expressed
as the

complexification, $\mathfrak{c}_{e},$ of $V_{1}$ by the real linear map 
\begin{gather}
\mathfrak{c}_{e}\mathfrak{:}V_{1}\oplus _{\mathbb{R}}V_{2}=V_{1}\oplus _{%
\mathbb{R}}J_{\left[ e\right] }V_{1}\rightarrow \mathfrak{H} \\
\mathfrak{c}_{e}\left( \eta \right) =v,
\end{gather}

i.e.%
\begin{equation}
\mathfrak{H}=V_{1}\oplus_{\mathbb{R}}J_{\left[ e\right] }V_{1}\equiv
V_{1}\oplus_{\mathbb{R}}iV_{1}=\mathbb{C\otimes}V_{1}=V_{1\mathbb{C}},
\end{equation}

where 
\begin{gather}
\eta=\left( x,y\right) \\
\nu=\mathfrak{c}_{e}\left( \eta\right) =\mathfrak{c}_{e}\left( \left\vert 
\mathbf{\varepsilon}_{1}\right\rangle \mathbf{x}+\left\vert \mathbf{%
\varepsilon}_{2}\right\rangle \mathbf{y}\right) =\mathfrak{c}_{e}\left(
\left\vert \mathbf{\varepsilon}_{1}\right\rangle \mathbf{x}+J_{\left[ e%
\right] }\left\vert \mathbf{\varepsilon}_{1}\right\rangle \mathbf{y}\right)
\rightarrow\left\vert \mathbf{\varepsilon}\right\rangle \func{Re}\mathbf{\nu}%
+i\left\vert \mathbf{\varepsilon}\right\rangle \func{Im}\mathbf{\nu}.
\end{gather}

The inner product in $\mathfrak{H}$ induced by the complexification, $%
\mathfrak{c}_{e},$ of $V_{\left[ \sigma_{e}\right] }$ is defined by%
\begin{align}
\left\langle \left. \nu_{1}\right\vert \nu_{2}\right\rangle _{\mathfrak{H}}
& =\frac{1}{2}\left( \left\langle \left. \eta_{1}\right\vert \eta
_{2}\right\rangle _{\mathfrak{H}_{\mathbb{R}}}+i\left\{ \left. \eta
_{1}\right\vert \eta_{2}\right\} _{\mathfrak{H}_{\mathbb{R}}}\right) \\
& =\frac{1}{2}\left( \left\langle \left. x_{1}\right\vert x_{2}\right\rangle
_{V_{1}}+\left\langle \left. y_{1}\right\vert y_{2}\right\rangle
_{V_{2}}+i\left( \left\langle \left. x_{1}\right\vert y_{2}\right\rangle
_{V_{1}}-\left\langle \left. x_{2}\right\vert y_{1}\right\rangle
_{V_{2}}\right) \right) .
\end{align}

The invariance group of $\left\langle \left. \nu_{1}\right\vert \nu
_{2}\right\rangle _{\mathfrak{H}}$ is thus seen to 
\begin{equation}
U\left( r\right) =Sp\left( 2r,\mathbb{R}\right) \cap O\left( 2r,\mathbb{R}%
\right) .
\end{equation}

The group $U\left( r\right) $ has a realization in both $\mathcal{B}\left( 
\mathfrak{H}_{\mathbb{R}}\right) $ and $\mathcal{B}\left( \mathfrak{H}%
\right) $ 
\begin{gather}
\widehat{\mathfrak{c}_{e}}\left( \left\vert \mathbf{\varepsilon }_{1},%
\mathbf{\varepsilon }_{2}\right\rangle \left( 
\begin{array}{cc}
\mathbb{A} & \mathbb{B} \\ 
\mathbb{C} & \mathbb{D}%
\end{array}%
\right) \left\langle \mathbf{\varepsilon }_{1},\mathbf{\varepsilon }%
_{2}\right\vert \right) =\mathfrak{c}_{e}\left\vert \mathbf{\varepsilon }%
_{1},\mathbf{\varepsilon }_{2}\right\rangle \left( 
\begin{array}{cc}
\mathbb{A} & \mathbb{B} \\ 
\mathbb{C} & \mathbb{D}%
\end{array}%
\right) \left\langle \mathbf{\varepsilon }_{1},\mathbf{\varepsilon }%
_{2}\right\vert \mathfrak{c}_{e}^{t}=\left\vert \mathbf{e},i\mathbf{e}%
\right\rangle \left( 
\begin{array}{cc}
\mathbb{A} & \mathbb{B} \\ 
\mathbb{C} & \mathbb{D}%
\end{array}%
\right) \left\langle \mathbf{e},i\mathbf{e}\right\vert \\
=\left\vert \mathbf{e}\right\rangle \mathbb{A}\left\langle \mathbf{e}%
\right\vert +\left\vert \mathbf{e}\right\rangle \mathbb{B}\left\langle i%
\mathbf{e}\right\vert +\left\vert i\mathbf{e}\right\rangle \mathbb{C}%
\left\langle \mathbf{e}\right\vert +\left\vert i\mathbf{e}\right\rangle 
\mathbb{D}\left\langle i\mathbf{e}\right\vert =\left\vert \mathbf{e}%
\right\rangle \mathbb{A}\left\langle \mathbf{e}\right\vert +\left\vert 
\mathbf{e}\right\rangle \mathbb{D}\left\langle \mathbf{e}\right\vert
+i\left( \left\vert \mathbf{e}\right\rangle \mathbb{C}\left\langle \mathbf{e}%
\right\vert -\left\vert \mathbf{e}\right\rangle \mathbb{D}\left\langle 
\mathbf{e}\right\vert \right) =\left\vert \mathbf{e}\right\rangle \left( 
\mathbb{A+B}\right) \left\langle \mathbf{e}\right\vert +i\left\vert \mathbf{e%
}\right\rangle \left( \mathbb{C-D}\right) \left\langle \mathbf{e}\right\vert
\\
=\left\vert \mathbf{e}\right\rangle \left\{ \left( \mathbb{A+B}\right)
+i\left( \mathbb{C-D}\right) \right\} \left\langle \mathbf{e}\right\vert
=\left\vert \mathbf{e}\right\rangle \left\{ \func{Re}\mathbb{T}+i\func{Im}%
\mathbb{T}\right\} \left\langle \mathbf{e}\right\vert =\left\vert \mathbf{e}%
\right\rangle \mathbb{T}\left\langle \mathbf{e}\right\vert
\end{gather}%
$\lceil $%
\begin{align}
\mathbb{T}^{\dagger }\mathbb{T}& =\left( \func{Re}\mathbb{T}^{t}-i\func{Im}%
\mathbb{T}^{t}\right) \left( \func{Re}\mathbb{T}+i\func{Im}\mathbb{T}\right)
=\func{Re}\mathbb{T}^{t}\func{Re}\mathbb{T+}\func{Im}\mathbb{T}^{t}\func{Im}%
\mathbb{T}+i\left( \func{Re}\mathbb{T}^{t}\func{Im}\mathbb{T}-\func{Im}%
\mathbb{T}^{t}\func{Re}\mathbb{T}\right) =\mathbb{I}_{r} \\
& =\left( \mathbb{A+B}\right) ^{t}\left( \mathbb{A+B}\right) +\left( \mathbb{%
C-D}\right) ^{t}\left( \mathbb{C-D}\right) +i\left( \left( \mathbb{A+B}%
\right) ^{t}\left( \mathbb{C-D}\right) -\left( \mathbb{C-D}\right)
^{t}\left( \mathbb{A+B}\right) \right) \\
& =\mathbb{A}^{t}\mathbb{A+A}^{t}\mathbb{B+B}^{t}\mathbb{A+B}^{t}\mathbb{B}+%
\mathbb{C}^{t}\mathbb{C-C}^{t}\mathbb{D-D}^{t}\mathbb{C+D}^{t}\mathbb{B}+%
\mathbb{D}^{t}\mathbb{D} \\
& +i\left( \mathbb{A}^{t}\mathbb{C-A}^{t}\mathbb{D+B}^{t}\mathbb{C-B}^{t}%
\mathbb{D}-\mathbb{C}^{t}\mathbb{A}-\mathbb{C}^{t}\mathbb{B}+\mathbb{D}^{t}%
\mathbb{A}+\mathbb{D}^{t}\mathbb{B}\right)
\end{align}%
$\rfloor $

The $\mathcal{B}\left( \mathfrak{H}_{\mathbb{R}}\right) $ operators $J_{%
\left[ e\right] }$ and $\sigma_{e}$ anti commute and are thus compatible
i.e. 
\begin{align}
\left\{ J_{\left[ e\right] },\sigma_{\left[ e\right] }\right\} & =\left(
\left\vert \mathbf{\varepsilon}_{2}\right\rangle \left\langle \mathbf{%
\varepsilon}_{1}\right\vert -\left\vert \mathbf{\varepsilon}%
_{1}\right\rangle \left\langle \mathbf{\varepsilon}_{2}\right\vert \right)
\left( \left\vert \mathbf{\varepsilon}_{1}\right\rangle \left\langle \mathbf{%
\varepsilon}_{1}\right\vert -\left\vert \mathbf{\varepsilon}%
_{2}\right\rangle \left\langle \mathbf{\varepsilon}_{2}\right\vert \right)
+\left( \left\vert \mathbf{\varepsilon}_{1}\right\rangle \left\langle 
\mathbf{\varepsilon}_{1}\right\vert -\left\vert \mathbf{\varepsilon}%
_{2}\right\rangle \left\langle \mathbf{\varepsilon}_{2}\right\vert \right)
\left( \left\vert \mathbf{\varepsilon}_{2}\right\rangle \left\langle \mathbf{%
\varepsilon}_{1}\right\vert -\left\vert \mathbf{\varepsilon}%
_{1}\right\rangle \left\langle \mathbf{\varepsilon}_{2}\right\vert \right) \\
& =\left\vert \mathbf{\varepsilon}_{2}\right\rangle \left\langle \mathbf{%
\varepsilon}_{1}\right\vert +\left\vert \mathbf{\varepsilon}%
_{1}\right\rangle \left\langle \mathbf{\varepsilon}_{2}\right\vert
-\left\vert \mathbf{\varepsilon}_{1}\right\rangle \left\langle \mathbf{%
\varepsilon}_{2}\right\vert -\left\vert \mathbf{\varepsilon}%
_{2}\right\rangle \left\langle \mathbf{\varepsilon}_{1}\right\vert =0.
\end{align}
\end{proof}

\begin{remark}
\textbf{Note that} $\mathfrak{H}$ is \emph{NOT} the complexification of $%
\mathfrak{H}_{\mathbb{R}}.$ The complexification, $\mathfrak{H}_{\mathbb{C}%
}, $ of $\mathfrak{H}_{\mathbb{R}}$ has \emph{complex dimension} $2r$ and is
given by 
\begin{equation}
\mathfrak{H}_{\mathbb{C}}=\mathbb{C\otimes}\mathfrak{H}_{\mathbb{R}}=\mathbb{%
C\otimes}\left( V_{\pm\sigma_{e}}\oplus_{\mathbb{R}}V_{\pm
\sigma_{e}}\right) =\mathbb{C\otimes}V_{\pm\sigma_{e}}\oplus_{\mathbb{R}}%
\mathbb{C\otimes}V_{\pm\sigma_{e}}.
\end{equation}
\end{remark}

The map $\mathfrak{c}_{e}$ is isometric if 
\begin{equation}
\left\langle \left. \mathfrak{c}_{e}\left( \nu\right) \right\vert \mathfrak{c%
}_{e}\left( \nu\right) \right\rangle _{\mathfrak{H}}=\left\langle \left.
\nu\right\vert \nu\right\rangle _{\mathfrak{H}_{\mathbb{R}}}\forall \nu\in%
\mathfrak{H}_{\mathbb{R}}
\end{equation}
and a real Hilbert space isomorphism if 
\begin{equation}
\left\langle \left. \mathfrak{c}_{e}\left( \nu_{1}\right) \right\vert 
\mathfrak{c}_{e}\left( \nu_{2}\right) \right\rangle _{\mathfrak{H}%
}=\left\langle \left. \nu_{1}\right\vert \nu_{2}\right\rangle _{\mathfrak{H}%
_{\mathbb{R}}}\forall\nu_{1},\nu_{2}\in\mathfrak{H}_{\mathbb{R}}.
\end{equation}
Two complexifications $\mathfrak{c}_{e}$ and $\mathfrak{c}_{f}$ are
unitarily equivalent if 
\begin{equation}
\mathfrak{c}_{e}=u\mathfrak{c}_{f}u^{\dagger};\text{for a }u\in U\left(
r\right) .
\end{equation}

\paragraph{Example of Complex Conjugations}

Let $\dim\left\{ \mathfrak{H}\right\} =2$ then the following antilinear
operators

\begin{align}
\tau_{0}\left( c_{1}e_{1}+c_{2}e_{2}\right) & =-\overline{c_{2}}e_{1}+%
\overline{c_{1}}e_{2} \\
\tau_{1}\left( c_{1}e_{1}+c_{2}e_{2}\right) & =-\overline{c_{1}}e_{1}+%
\overline{c_{2}}e_{2} \\
\tau_{2}\left( c_{1}e_{1}+c_{2}e_{2}\right) & =i\overline{c_{1}}e_{1}+i%
\overline{c_{2}}e_{2} \\
\tau_{3}\left( c_{1}e_{1}+c_{2}e_{2}\right) & =\overline{c_{1}}e_{1}+%
\overline{c_{2}}e_{2}
\end{align}
are complex conjugations. These definitions can also be expressed in terms
of the conjugation $\sigma_{e}$ as 
\begin{align}
\tau_{0}\left( \mathbf{c}\right) & =\tau_{0}\left\vert \mathbf{e}%
\right\rangle \left( 
\begin{array}{c}
c_{1} \\ 
c_{2}%
\end{array}
\right) =\left\vert \mathbf{e}\right\rangle \sigma_{e}\left( \left( 
\begin{array}{cc}
0 & -1 \\ 
1 & 0%
\end{array}
\right) \left( 
\begin{array}{c}
c_{1} \\ 
c_{2}%
\end{array}
\right) \right) =\left\vert \mathbf{e}\right\rangle \left( 
\begin{array}{c}
-\overline{c_{2}} \\ 
\overline{c_{1}}%
\end{array}
\right) \\
\tau_{1}\left( \mathbf{c}\right) & =\tau_{1}\left\vert \mathbf{e}%
\right\rangle \left( 
\begin{array}{c}
c_{1} \\ 
c_{2}%
\end{array}
\right) =\left\vert \mathbf{e}\right\rangle \sigma_{e}\left( \left( 
\begin{array}{cc}
-1 & 0 \\ 
0 & 1%
\end{array}
\right) \left( 
\begin{array}{c}
c_{1} \\ 
c_{2}%
\end{array}
\right) \right) =\left\vert \mathbf{e}\right\rangle \left( 
\begin{array}{c}
-\overline{c_{1}} \\ 
\overline{c_{2}}%
\end{array}
\right) \\
\tau_{2}\left( \mathbf{c}\right) & =\tau_{2}\left\vert \mathbf{e}%
\right\rangle \left( 
\begin{array}{c}
c_{1} \\ 
c_{2}%
\end{array}
\right) =\left\vert \mathbf{e}\right\rangle \sigma_{e}\left( \left( 
\begin{array}{cc}
i & 0 \\ 
0 & i%
\end{array}
\right) \left( 
\begin{array}{c}
c_{1} \\ 
c_{2}%
\end{array}
\right) \right) \,\,\,=\left\vert \mathbf{e}\right\rangle \left( 
\begin{array}{c}
i\overline{c_{1}} \\ 
i\overline{c_{2}}%
\end{array}
\right) \\
\tau_{3}\left( \mathbf{c}\right) & =\tau_{3}\left\vert \mathbf{e}%
\right\rangle \left( 
\begin{array}{c}
c_{1} \\ 
c_{2}%
\end{array}
\right) =\left\vert \mathbf{e}\right\rangle \sigma_{e}\left( \left( 
\begin{array}{cc}
1 & 0 \\ 
0 & 1%
\end{array}
\right) \left( 
\begin{array}{c}
c_{1} \\ 
c_{2}%
\end{array}
\right) \right) \,\,\,=\left\vert \mathbf{e}\right\rangle \left( 
\begin{array}{c}
\overline{c_{1}} \\ 
\overline{c_{2}}%
\end{array}
\right) .
\end{align}
These conjugations are expressed in a common basis $\left\{
e_{1},e_{2}\right\} $, but for each of these conjugations there exists bases
where they act in standard form i.e. like $\tau_{3}$ above.

The conjugations $\left\{ \left. \tau_{l}\right\vert 1\leq l\leq3\right\} $
are \emph{self antilinear adjoint} i.e. 
\begin{equation}
\left\langle \left. e_{j}\right\vert \tau_{l}e_{k}\right\rangle =\overline{%
\left\langle \left. e_{k}\right\vert \tau_{l}e_{j}\right\rangle };\,1\leq
j,k\leq2;1\leq l\leq3
\end{equation}
and the conjugation $\tau_{0}$ is \emph{skew antilinear adjoint i.e.}%
\begin{equation}
\left\langle \left. e_{j}\right\vert \tau_{0}e_{k}\right\rangle =-\overline{%
\left\langle \left. e_{k}\right\vert \tau_{0}e_{j}\right\rangle };\,1\leq
j,k\leq2.
\end{equation}
\emph{\ }The conjugations $\left\{ \left. \tau_{l}\right\vert 0\leq
l\leq3\right\} $ are orthogonal wrt the $\func{Tr}$ inner product i.e. 
\begin{equation}
\func{Tr}\left\{ \tau_{k}\tau_{l}\right\} =2\delta_{kl};\,0\leq l,k\leq3
\end{equation}
which is well defined as the products $\left\{ \left. \tau_{k}\tau
_{l}\right\vert 0\leq k,l\leq3\right\} $ are linear operators.

One can relate $\left\{ \left. \tau_{l}\right\vert 0\leq l\leq3\right\} $ to
the Pauli operators by 
\begin{align}
\tau_{1}\tau_{2}\tau_{3} & =-i\tau_{0} \\
\tau_{1}\tau_{2} & =i\sigma_{3} \\
\tau_{3}\tau_{1} & =i\sigma_{2} \\
\tau_{2}\tau_{3} & =i\sigma_{1} \\
\tau_{2}\tau_{0} & =\sigma_{2} \\
\tau_{1}\tau_{0} & =\sigma_{1} \\
\tau_{3}\tau_{0} & =\sigma_{3}.
\end{align}

\subsection{The Banach Space of Anti-Linear Operators}

The complex Banach space of antilinear operators is denoted by $\mathcal{B}%
_{anti}\left( \mathfrak{H}\right) ,$ the norm definition in $\mathcal{B}%
_{anti}\left( \mathfrak{H}\right) $ is the same as in $\mathcal{B}\left( 
\mathfrak{H}\right) $ i.e. 
\begin{equation}
\left\Vert A\right\Vert _{\mathcal{B}_{anti}\left( \mathfrak{H}\right)
}=\sup\limits_{\left\Vert \left\langle \left. v\right\vert v\right\rangle
\right\Vert =1}\left\{ \left\langle \left. Av\right\vert Av\right\rangle
\right\} .
\end{equation}%
The \emph{antilinear adjoint, }$A^{\dagger _{a}},$ of an operator $A$ is
defined as 
\begin{equation}
\overline{\left\langle \left. v_{1}\right\vert A^{\dagger
_{a}}v_{2}\right\rangle }=\left\langle \left. A^{\dagger
_{a}}v_{2}\right\vert v_{1}\right\rangle =\overline{\left\langle \left.
v_{2}\right\vert Av_{1}\right\rangle },\,\forall v_{1},v_{2}\in \mathfrak{H,}
\end{equation}%
where $A$ is linear or antilinear. The antilinear adjoint of a complex
conjugation $\sigma _{e}$ is given by 
\begin{equation}
\left\langle \left. \sigma _{e}^{\dagger _{a}}v_{2}\right\vert
v_{1}\right\rangle =\overline{\left\langle \left. v_{2}\right\vert \sigma
_{e}v_{1}\right\rangle }=\sum_{1\leq j\leq r}\mathcal{R}_{e}\left(
v_{2}\right) _{j}\mathcal{R}_{e}\left( v_{1}\right) _{j};\,\,\forall
v_{1},v_{2}\in \mathfrak{H,}
\end{equation}%
which shows that $\sigma _{e}$ \emph{is antilinear self adjoint} as 
\begin{equation}
\left\langle \left. \sigma _{e}v_{2}\right\vert v_{1}\right\rangle
=\sum_{1\leq j\leq r}\mathcal{R}_{e}\left( v_{2}\right) _{j}\mathcal{R}%
_{e}\left( v_{1}\right) _{j}=\overline{\left\langle \left. v_{2}\right\vert
\sigma _{e}v_{1}\right\rangle };\,\,\forall v_{1},v_{2}\in \mathfrak{H}
\end{equation}%
i.e. 
\begin{equation}
\sigma _{e}=\sigma _{e}^{\dagger _{a}}
\end{equation}%
and thus 
\begin{equation}
\sigma _{e}\sigma _{e}^{\dagger _{a}}=\sigma _{e}^{\dagger _{a}}\sigma
_{e}=\sigma _{e}^{2}=I_{\mathcal{B}\left( \mathfrak{H}\right) },
\end{equation}%
which is consistent with $\sigma _{e}$ being anti-unitary. In particular 
\begin{equation}
\left\langle \left. \left( \left\vert v_{2}\right\rangle \left\langle
v_{1}\right\vert \right) ^{\dagger _{a}}v_{2}\right\vert v_{1}\right\rangle =%
\overline{\left\langle \left. v_{2}\right\vert \left\vert v_{2}\right\rangle
\left\langle v_{1}\right\vert v_{1}\right\rangle }=\left\langle \left.
v_{2}\right\vert \left\vert v_{2}\right\rangle \left\langle v_{1}\right\vert
v_{1}\right\rangle .
\end{equation}

Normal adjoint is defined as 
\begin{equation}
\left\langle \left. A^{\dagger}v_{2}\right\vert v_{1}\right\rangle
=\left\langle \left. v_{2}\right\vert Av_{1}\right\rangle
\end{equation}
and normal self adjointness as 
\begin{equation}
A^{\dagger}=A\Rightarrow\left\langle \left. Ae_{j}\right\vert
e_{k}\right\rangle =\overline{\left\langle \left. e_{k}\right\vert
Ae_{j}\right\rangle }=\left\langle \left. e_{j}\right\vert
Ae_{k}\right\rangle .
\end{equation}

Using the antilinear adjoint the norm can be re-expressed as 
\begin{equation}
\left\Vert A\right\Vert _{\mathcal{B}_{anti}\left( \mathfrak{H}\right)
}=\sup\limits_{\left\Vert \left\langle \left. v\right\vert v\right\rangle
\right\Vert =1}\left\{ \left\vert \overline{\left\langle \left. v\right\vert
A^{\dagger_{a}}Av\right\rangle }\right\vert \right\}
=\sup\limits_{\left\Vert \left\langle \left. v\right\vert v\right\rangle
\right\Vert =1}\left\{ \left\vert \left\langle \left. v\right\vert
A^{\dagger_{a}}Av\right\rangle \right\vert \right\} ,
\end{equation}
as the product of two antilinear operators is linear the operator $%
A^{\dagger_{a}}A$ is linear$.$ One can define a hermitian inner product in $%
\mathcal{B}_{2anti}\left( \mathfrak{H}\right) \subseteq\mathcal{B}%
_{anti}\left( \mathfrak{H}\right) $ by 
\begin{equation}
\left\langle \left. A\right\vert B\right\rangle =\func{Tr}\left\{ BA\right\}
<\infty;\,A,B\in\mathcal{B}_{2anti}\left( \mathfrak{H}\right) ,
\end{equation}
which makes sense if $BA\in\mathcal{B}_{1}\left( \mathfrak{H}\right) .$

One should observe that $\mathcal{B}_{anti}\left( \mathfrak{H}\right) $ is
not a $c^{\ast}$-algebra as "$\dagger_{a}"$ is an antilinear operation not a
linear one.

If 
\begin{equation}
A^{\dagger_{a}}=A
\end{equation}
then $A$ is \emph{self antilinear adjoint}, while if 
\begin{equation}
A^{\dagger_{a}}=-A
\end{equation}
then $A$ is \emph{skew antilinear adjoint. }The inner product\emph{\ }$%
\left\langle \left. A\right\vert B\right\rangle $ is \emph{positive definite}
on the \emph{real} space $\mathcal{B}_{anti}^{+}\left( \mathfrak{H}\right) $
of self antilinear adjoint i.e. 
\begin{equation}
\left\langle \left. A\right\vert A\right\rangle =\func{Tr}\left\{
A^{2}\right\} >0;\,A^{\dagger_{a}}=A
\end{equation}
and \emph{negative definite} on the \emph{real} space $\mathcal{B}%
_{anti}^{-}\left( \mathfrak{H}\right) $ of skew antilinear adjoint i.e. 
\begin{equation}
\left\langle \left. A\right\vert A\right\rangle =\func{Tr}\left\{
A^{2}\right\} <0;\,A^{\dagger_{a}}=-A.
\end{equation}

One can introduce a ket-bra operators for $\mathcal{B}_{anti}\left( 
\mathfrak{H}\right) $ that are defined by 
\begin{equation}
\left( \left\vert v_{2}\right\rangle \left\langle v_{1}\right\vert \right)
_{a}\left\vert w\right\rangle =\left\langle \left. w\right\vert
v_{1}\right\rangle \left\vert v_{2}\right\rangle =\overline{\left\langle
\left. v_{1}\right\vert w\right\rangle }\left\vert v_{2}\right\rangle
\end{equation}
giving 
\begin{equation}
\left( \left\vert v_{2}\right\rangle \left\langle v_{1}\right\vert \right)
_{a}\left\vert e_{j}\right\rangle =\left\langle \left. e_{j}\right\vert
v_{1}\right\rangle \left\vert v_{2}\right\rangle =\overline{\left\langle
\left. v_{1}\right\vert e_{j}\right\rangle }\left\vert v_{2}\right\rangle
\end{equation}
and 
\begin{equation}
\left( \left\vert e_{j}\right\rangle \left\langle e_{k}\right\vert \right)
_{a}\left\vert w\right\rangle =\left\langle \left. w\right\vert
e_{k}\right\rangle \left\vert e_{j}\right\rangle =\overline{\left\langle
\left. e_{k}\right\vert w\right\rangle }\left\vert e_{j}\right\rangle =%
\overline{w_{k}}\left\vert e_{j}\right\rangle .
\end{equation}
In particular 
\begin{equation}
\sum_{1\leq j\leq r}\left( \left\vert e_{j}\right\rangle \left\langle
e_{j}\right\vert \right) _{a}\left\vert w\right\rangle =\sum_{1\leq j\leq
r}\left\langle \left. w\right\vert e_{j}\right\rangle \left\vert
e_{j}\right\rangle =\sum_{1\leq j\leq r}\overline{\left\langle \left.
e_{j}\right\vert w\right\rangle }\left\vert e_{j}\right\rangle =\sum_{1\leq
j\leq r}\overline{w}_{j}\left\vert e_{j}\right\rangle ,
\end{equation}
which is the complex conjugation $\sigma_{e}$ i.e. 
\begin{equation}
\sigma_{e}=\sum_{1\leq j\leq r}\left( \left\vert e_{j}\right\rangle
\left\langle e_{j}\right\vert \right) _{a}.
\end{equation}

One can define a basis of self antilinear adjoint ket-bras by \emph{\ } 
\begin{align}
\mathfrak{e}_{jk}^{+} & =\frac{1}{\sqrt{2}}\left( \left( \left\vert
e_{j}\right\rangle \left\langle e_{k}\right\vert \right) _{a}+\left(
\left\vert e_{k}\right\rangle \left\langle e_{j}\right\vert \right)
_{a}\right) ;1\leq j<k\leq r \\
\mathfrak{e}_{jj}^{+} & =\left( \left\vert e_{j}\right\rangle \left\langle
e_{j}\right\vert \right) _{a};1\leq j\leq r,
\end{align}
that form a basis for the real Banach space $\mathcal{B}_{anti}^{e+}\left( 
\mathfrak{H}\right) $ that has dimension $\left. r\left( r+1\right) \right/
2 $ and skew antilinear adjoint ket-bras by%
\begin{equation}
\mathfrak{e}_{jk}^{-}=\frac{1}{\sqrt{2}}\left( \left( \left\vert
e_{j}\right\rangle \left\langle e_{k}\right\vert \right) _{a}-\left(
\left\vert e_{k}\right\rangle \left\langle e_{j}\right\vert \right)
_{a}\right) ;1\leq j<k\leq r,
\end{equation}
that form a basis for the real Banach space $\mathcal{B}_{anti}^{e-}\left( 
\mathfrak{H}\right) $ that has dimension $\left. r\left( r-1\right) \right/
2.$

The action of a general antilinear operator on $\mathfrak{H}$ is given by 
\begin{equation}
\sum_{1\leq j,k,\leq r}\alpha_{jk}\left( \left\vert e_{j}\right\rangle
\left\langle e_{k}\right\vert \right) _{a}\left\vert w\right\rangle
=\sum_{1\leq j,k,\leq r}\alpha_{jk}\overline{\left\langle \left.
e_{k}\right\vert w\right\rangle }\left\vert e_{j}\right\rangle =\sum_{1\leq
j,k,\leq r}\alpha_{jk}\overline{w}_{k}\left\vert e_{j}\right\rangle
=\sum_{1\leq j,\leq r}\left( \mathbf{\alpha}\overline{w}\right)
_{j}\left\vert e_{j}\right\rangle
\end{equation}
and the product of two antilinear operators by 
\begin{equation}
\sum_{1\leq j,k,\leq r}\beta_{jk}\left( \left\vert e_{j}\right\rangle
\left\langle e_{k}\right\vert \right) _{a}\sum_{1\leq j,\leq r}\left( 
\mathbf{\alpha}\overline{w}\right) _{j}\left\vert e_{j}\right\rangle
=\sum_{1\leq j,k,\leq r}\beta_{jk}\overline{\left( \mathbf{\alpha}\overline{w%
}\right) }_{k}\left\vert e_{j}\right\rangle =\sum_{1\leq j,k,\leq
r}\beta_{jk}\left( \overline{\mathbf{\alpha}}w\right) _{k}\left\vert
e_{j}\right\rangle =\sum_{1\leq j,\leq r}\left( \mathbf{\beta}\overline {%
\mathbf{\alpha}}w\right) _{j}\left\vert e_{j}\right\rangle .
\end{equation}
The product of two antilinear operators gives a linear operator 
\begin{equation}
\sum_{1\leq l,m,\leq r}\beta_{lm}\left( \left\vert e_{j}\right\rangle
\left\langle e_{k}\right\vert \right) _{a}\sum_{1\leq j,k,\leq r}\alpha
_{jk}\left( \left\vert e_{j}\right\rangle \left\langle e_{k}\right\vert
\right) _{a}=\sum_{1\leq j,k,\leq r}\lambda_{jk}\left\vert
e_{j}\right\rangle \left\langle e_{k}\right\vert ,
\end{equation}
where 
\begin{equation}
\lambda_{jk}=\left( \mathbf{\beta}\overline{\mathbf{\alpha}}\right) _{jk}.
\end{equation}

This allows one to define a hermitian inner product on $\mathcal{B}%
_{anti}\left( \mathfrak{H}\right) $ by 
\begin{equation}
\left\langle \left. A\right\vert B\right\rangle _{\mathcal{B}_{anti}\left( 
\mathfrak{H}\right) }=\func{Tr}\left\{ \mathbb{B}\overline {\mathbb{A}}%
\right\} =\func{Tr}\left\{ \overline{\mathbb{A}}\mathbb{B}\right\}
=\sum_{1\leq j,\leq r}\left( \mathbb{B}\overline {\mathbb{A}}\right)
_{jj}=\sum_{1\leq j,\leq r}\overline{\left( \mathbb{A}\overline{\mathbb{B}}%
\right) }_{jj}=\overline{\left\langle \left. B\right\vert A\right\rangle }_{%
\mathcal{B}_{anti}\left( \mathfrak{H}\right) }.
\end{equation}

Complex conjugations, $\sigma_{e},$define \emph{complex conjugate operators} 
$\Omega_{\mathbf{e}}\left( A\right) $ belonging to $\mathcal{B}\left( 
\mathfrak{H}\right) $, thus $\Omega_{\mathbf{e}}$ is a \emph{complex
conjugation of the Banach space }$\mathcal{B}\left( \mathfrak{H}\right) $ by
the inner transformation 
\begin{align}
\Omega_{\mathbf{e}} & :\mathcal{B}\left( \mathfrak{H}\right) \rightarrow%
\mathcal{B}\left( \mathfrak{H}\right) \text{ and is antilinear } \\
\Omega_{\mathbf{e}}\left( A\right) & =A^{\ast_{e}}=\sigma_{e}A\sigma
_{e};\,\left( A,\sigma_{e}A\sigma_{e}=A^{\ast_{e}}\right) \in\mathcal{B}%
\left( \mathfrak{H}\right)
\end{align}
$\lceil$%
\begin{align}
\Omega_{\mathbf{e}}\left( \alpha A\right) \left( av_{1}+bv_{2}\right) &
=\sigma_{e}\alpha A\sigma_{e}\left( av_{1}+bv_{2}\right) =\sigma_{e}\alpha
A\sigma_{e}\left( av_{1}\right) +\sigma_{e}\alpha A\sigma_{e}\left(
bv_{2}\right) =\sigma_{e}\alpha A\overline{a}\sigma_{e}\left( v_{1}\right)
+\sigma_{e}\alpha A\overline{b}\sigma_{e}\left( v_{2}\right) \\
& =\sigma_{e}\alpha\overline{a}A\sigma_{e}\left( v_{1}\right) +\sigma
_{e}\alpha\overline{b}A\sigma_{e}\left( v_{2}\right) =\overline{\alpha }%
a\left( \sigma_{e}A\sigma_{e}\right) v_{1}+\overline{\alpha}b\left(
\sigma_{e}A\sigma_{e}\right) v_{2}=\overline{\alpha}a\Omega_{\mathbf{e}%
}\left( A\right) v_{1}+\overline{\alpha}b\Omega_{\mathbf{e}}\left( A\right)
v_{2}
\end{align}%
\begin{equation}
\Omega_{\mathbf{e}}\left( \alpha A+\beta B\right) =\overline{\alpha}%
\sigma_{e}A\sigma_{e}+\overline{\beta}\sigma_{e}B\sigma_{e}=\overline{\alpha 
}\Omega_{\mathbf{e}}\left( A\right) +\overline{\beta}\Omega_{\mathbf{e}%
}\left( B\right)
\end{equation}
$\rfloor$

and \emph{transpose operators} $\Upsilon_{\mathbf{e}}\left( A\right) $ by by
the inner transformation 
\begin{align}
\Upsilon_{\mathbf{e}} & :\mathcal{B}\left( \mathfrak{H}\right) \rightarrow%
\mathcal{B}\left( \mathfrak{H}\right) \text{ and is }\mathbb{C}\text{-linear}
\\
\Upsilon_{\mathbf{e}}\left( A\right) & =A^{t_{e}}=\sigma_{e}A^{\dagger
}\sigma_{e};\,\,A,\sigma_{e}A\sigma_{e}\in\mathcal{B}\left( \mathfrak{H}%
\right)
\end{align}
for $A\in\mathcal{B}\left( \mathfrak{H}\right) .$

$\lceil$%
\begin{equation}
\Upsilon_{\mathbf{e}}\left( \alpha A+\beta B\right) =\sigma_{e}\left( 
\overline{\alpha}A^{\dagger}+\overline{\beta}B^{\dagger}\right) \sigma
_{e}=\alpha\sigma_{e}A^{\dagger}\sigma_{e}+\beta\sigma_{e}B^{\dagger}%
\sigma_{e}=\alpha\Upsilon_{\mathbf{e}}\left( A\right) +\beta\Upsilon _{%
\mathbf{e}}\left( B\right)
\end{equation}
$\rfloor$

The representation $\mathcal{R}_{e}\left( \sigma_{e}A\sigma_{e}\right) $ wrt
the basis $\left\{ \left. \left\vert e_{j}\right\rangle \right\vert 1\leq
j\leq r\right\} $ is given by 
\begin{align}
A^{\ast_{e}}\left\vert \lambda\right\rangle & =\sigma_{e}A\sigma
_{e}\left\vert \lambda\right\rangle =\sum_{1\leq j,k\leq
r}\sigma_{e}A\sigma_{e}\left( \left\vert e_{j}\right\rangle \mathcal{R}%
_{e}\left( \lambda\right) _{j}\right) =\sigma_{e}A\sum_{1\leq j\leq
r}\left\vert e_{j}\right\rangle \overline{\mathcal{R}_{e}\left(
\lambda\right) }_{j}=\sum_{1\leq j\leq r}\sigma_{e}A\left\vert
e_{j}\right\rangle \overline {\mathcal{R}_{e}\left( \lambda\right) }_{j} \\
& =\sum_{1\leq j,k\leq r}\sigma_{e}\left\vert e_{k}\right\rangle \mathcal{R}%
_{e}\left( A\right) _{kj}\overline{\mathcal{R}_{e}\left( \lambda\right) }%
_{j}=\sum_{1\leq j,k\leq r}\sigma_{e}\left( \left\vert e_{k}\right\rangle 
\mathcal{R}_{e}\left( A\right) _{kj}\overline {\mathcal{R}_{e}\left(
\lambda\right) }_{j}\right) =\sum_{1\leq j,k\leq r}\left\vert
e_{k}\right\rangle \overline{\mathcal{R}_{e}\left( A\right) }_{kj}\mathcal{R}%
_{e}\left( \lambda\right) _{j}
\end{align}
$\lceil$%
\begin{equation}
\mathcal{R}_{e}\left( \sigma_{e}A\sigma_{e}\right) _{jk}=\left\langle \left.
e_{j}\right\vert \sigma_{e}A\sigma_{e}\left. e_{k}\right\vert \right\rangle
=\left\langle \left. e_{j}\right\vert \sigma_{e}A\left. e_{k}\right\vert
\right\rangle =\left\langle \overset{}{e_{j}}\left\vert \sigma_{e}\left(
\sum_{1\leq l\leq r}\left\vert e_{l}\right\rangle \mathcal{R}_{e}\left(
A_{lk}\right) \right) \right. \right\rangle =\left\langle \overset{}{e_{j}}%
\left\vert \sum_{1\leq l\leq r}\left\vert e_{l}\right\rangle \overline{%
\mathcal{R}_{e}\left( A_{lk}\right) }\right. \right\rangle =\overline{%
\mathcal{R}_{e}\left( A\right) }_{jk},
\end{equation}
$\rfloor$

i.e. the $e$-representation of $A$ the $e$-complex conjugation $\Omega
_{e}\left( A\right) $ of $A$ is the complex conjugation of the $e$%
-representation of $A$ i.e. 
\begin{equation}
\mathcal{R}_{e}\left( \Omega_{e}\left( A\right) \right) =\overline {\mathcal{%
R}_{e}\left( A\right) },
\end{equation}
but in general 
\begin{equation}
\mathcal{R}_{e}\left( \Omega_{f}\left( A\right) \right) =\mathcal{R}%
_{e}\left( \sigma_{f}A\sigma_{f}\right) \neq\overline{\mathcal{R}_{e}\left(
A\right) }.
\end{equation}

$\lceil $Change of basis $\left\vert \mathbf{e}\right\rangle \leftrightarrow
\left\vert \mathbf{f}\right\rangle $ can be expressed by the unitary
operators $U_{fe}\equiv \left\vert \mathbf{f}\right\rangle \left\langle 
\mathbf{e}\right\vert $ and $U_{ef}\equiv \left\vert \mathbf{e}\right\rangle
\left\langle \mathbf{f}\right\vert $ as 
\begin{align}
\left\vert \mathbf{f}\right\rangle & =U_{fe}\left\vert \mathbf{e}%
\right\rangle =\left( \left\vert \mathbf{f}\right\rangle \left\langle 
\mathbf{e}\right\vert \right) \left\vert \mathbf{e}\right\rangle =\left\vert 
\mathbf{e}\right\rangle \mathcal{R}_{e}\left( \left\vert \mathbf{f}%
\right\rangle \left\langle \mathbf{e}\right\vert \right) =\left\vert \mathbf{%
e}\right\rangle \left\langle \left. \mathbf{e}\right\vert \mathbf{f}%
\right\rangle =\left\vert \mathbf{e}\right\rangle \mathbb{U}_{fe} \\
\left\vert \mathbf{e}\right\rangle & =U_{ef}\left\vert \mathbf{f}%
\right\rangle \,=\left( \left\vert \mathbf{e}\right\rangle \left\langle 
\mathbf{f}\right\vert \right) \left\vert \mathbf{f}\right\rangle
\,\,=\,\,\left\vert \mathbf{f}\right\rangle \mathcal{R}_{f}\left( \left\vert 
\mathbf{e}\right\rangle \left\langle \mathbf{f}\right\vert \right)
=\left\vert \mathbf{f}\right\rangle \left\langle \left. \mathbf{f}%
\right\vert \mathbf{e}\right\rangle \,=\left\vert \mathbf{f}\right\rangle 
\mathbb{U}_{ef}
\end{align}%
and are adjoints of each other i.e. 
\begin{equation}
U_{fe}^{\dagger }=\left( \left\vert \mathbf{f}\right\rangle \left\langle 
\mathbf{e}\right\vert \right) ^{\dagger }=\left\vert \mathbf{e}\right\rangle
\left\langle \mathbf{f}\right\vert =U_{fe},
\end{equation}%
where the unitary operators are the sum of partial isometries i.e. 
\begin{align}
\left\vert \mathbf{f}\right\rangle \left\langle \mathbf{e}\right\vert &
=\sum_{1\leq j\leq r}\left\vert f_{j}\right\rangle \left\langle
e_{j}\right\vert \\
\left\vert \mathbf{e}\right\rangle \left\langle \mathbf{f}\right\vert &
=\sum_{1\leq j\leq r}\left\vert e_{j}\right\rangle \left\langle
f_{j}\right\vert .
\end{align}%
$\lceil $The representations of these unitary operators \emph{are the same
in the two basis} i.e. 
\begin{equation}
\mathcal{R}_{e}\left( \left\vert \mathbf{f}\right\rangle \left\langle 
\mathbf{e}\right\vert \right) =\left\langle \left. \mathbf{e}\right\vert
\left\vert \mathbf{f}\right\rangle \left\langle \mathbf{e}\right\vert 
\mathbf{e}\right\rangle =\left\langle \left. \mathbf{e}\right\vert \mathbf{f}%
\right\rangle
\end{equation}%
\begin{equation}
\mathcal{R}_{f}\left( \left\vert \mathbf{f}\right\rangle \left\langle 
\mathbf{e}\right\vert \right) =\left\langle \left. \mathbf{f}\right\vert
\left\vert \mathbf{f}\right\rangle \left\langle \mathbf{e}\right\vert 
\mathbf{f}\right\rangle =\left\langle \left. \mathbf{e}\right\vert \mathbf{f}%
\right\rangle
\end{equation}%
$\rfloor $

The $\left\vert \mathbf{e}\right\rangle $ and $\left\vert \mathbf{f}%
\right\rangle $ bases produce the representations, $\mathcal{R}_{e}\left(
v\right) $ and $\mathcal{R}_{f}\left( v\right) $ of a vector $\left\vert
v\right\rangle $ 
\begin{align}
\left\vert v\right\rangle & =\left\vert \mathbf{e}\right\rangle \left\langle 
\mathbf{e}\right\vert \left\vert v\right\rangle =\sum_{1\leq j\leq
r}\left\vert e_{j}\right\rangle \left\langle \left. e_{j}\right\vert
v\right\rangle =\sum_{1\leq j\leq r}\left\vert e_{j}\right\rangle \mathcal{R}%
_{e}\left( v\right) _{j} \\
\left\vert v\right\rangle & =\left\vert \mathbf{f}\right\rangle \left\langle 
\mathbf{f}\right\vert \left\vert v\right\rangle =\sum_{1\leq j\leq
r}\left\vert f_{j}\right\rangle \left\langle \left. f_{j}\right\vert
v\right\rangle =\sum_{1\leq j\leq r}\left\vert f_{j}\right\rangle \mathcal{R}%
_{f}\left( v\right) _{j}.
\end{align}%
As 
\begin{equation}
\left\vert v\right\rangle =\left\vert \mathbf{e}\right\rangle \left\langle 
\mathbf{e}\right\vert \left\vert \mathbf{f}\right\rangle \left\langle 
\mathbf{f}\right\vert \left\vert v\right\rangle =\left\vert \mathbf{e}%
\right\rangle \mathbb{U}_{fe}\mathcal{R}_{f}\left( v\right)
\end{equation}%
they are related to each other by 
\begin{equation}
\mathcal{R}_{e}\left( v\right) =\mathbb{U}_{fe}\mathcal{R}_{f}\left(
v\right) \equiv \left\langle \left. \mathbf{e}\right\vert \mathbf{f}%
\right\rangle \mathcal{R}_{f}\left( v\right) .
\end{equation}%
The $\left\vert \mathbf{e}\right\rangle $ and $\left\vert \mathbf{f}%
\right\rangle $ bases produce the representations, $\mathcal{R}_{e}\left(
A\right) $ and $\mathcal{R}_{f}\left( A\right) $ of an operator $A$ by 
\begin{equation}
A=\left\vert \mathbf{f}\right\rangle \left\langle \mathbf{f}\right\vert
A\left\vert \mathbf{f}\right\rangle \left\langle \mathbf{f}\right\vert
=\left\vert \mathbf{f}\right\rangle \mathcal{R}_{f}\left( A\right)
\left\langle \mathbf{f}\right\vert =\left\vert \mathbf{e}\right\rangle
\left\langle \mathbf{e}\right\vert A\left\vert \mathbf{e}\right\rangle
\left\langle \mathbf{e}\right\vert =\left\vert \mathbf{e}\right\rangle 
\mathcal{R}_{e}\left( A\right) \left\langle \mathbf{e}\right\vert
\end{equation}%
and%
\begin{equation}
A=\left\vert \mathbf{e}\right\rangle \left( \left\langle \mathbf{e}%
\right\vert \left\vert \mathbf{f}\right\rangle \right) \mathcal{R}_{f}\left(
A\right) \left( \left\langle \mathbf{f}\right\vert \left\vert \mathbf{e}%
\right\rangle \right) \left\langle \mathbf{e}\right\vert =\left\vert \mathbf{%
e}\right\rangle \fbox{$\mathbb{U}_{f}e\mathcal{R}_{f}\left( A\right) \mathbb{%
U}_{f}e^{\dagger }$}\left\langle \mathbf{e}\right\vert =\left\vert \mathbf{e}%
\right\rangle \fbox{$\mathbb{U}_{e}f^{\dagger }\mathcal{R}_{f}\left(
A\right) \mathbb{U}_{e}f$}\left\langle \mathbf{e}\right\vert =\left\vert 
\mathbf{e}\right\rangle \fbox{$\mathcal{R}_{e}\left( A\right) $}\left\langle 
\mathbf{e}\right\vert ,
\end{equation}%
which are related to each other by 
\begin{gather}
\mathcal{R}_{e}\left( A\right) =\left\langle \mathbf{e}\right\vert
\left\vert \mathbf{f}\right\rangle \mathcal{R}_{f}\left( A\right)
\left\langle \mathbf{f}\right\vert \left\vert \mathbf{e}\right\rangle \equiv 
\mathbb{U}_{fe}\mathcal{R}_{f}\left( A\right) \mathbb{U}_{fe}^{\dagger } \\
\Updownarrow \\
\mathcal{R}_{f}\left( A\right) =\left\langle \mathbf{f}\right\vert
\left\vert \mathbf{e}\right\rangle \mathcal{R}_{e}\left( A\right)
\left\langle \mathbf{e}\right\vert \left\vert \mathbf{f}\right\rangle \equiv 
\mathbb{U}_{fe}^{\dagger }\mathcal{R}_{f}\left( A\right) \mathbb{U}_{fe}.
\end{gather}%
The $f$-complex conjugate, $\left\vert \sigma _{f}\left( \nu \right)
\right\rangle ,$in the $\left\vert \mathbf{e}\right\rangle $ basis can be
expandeds in the $\left\vert \sigma _{f}\left( \mathbf{e}\right)
\right\rangle $ basis as 
\begin{equation}
\left\vert v\right\rangle =\left\vert \sigma _{f}\left( \mathbf{e}\right)
\right\rangle \left\langle \sigma _{f}\left( \mathbf{e}\right) \right\vert
\left\vert v\right\rangle =\sum_{1\leq j\leq r}\left\vert \sigma _{f}\left(
e_{j}\right) \right\rangle \left\langle \left. \sigma _{f}\left(
e_{j}\right) \right\vert v\right\rangle =\sum_{1\leq j\leq r}\left\vert
\sigma _{f}\left( e_{j}\right) \right\rangle \mathcal{R}_{\sigma _{f}\left(
e\right) }\left( v\right) _{j}.
\end{equation}%
The $f$-complex conjugate, $\left\vert \sigma _{f}\left( \mathbf{e}\right)
\right\rangle ,$ of the $\left\vert \mathbf{e}\right\rangle $ basis can be
expanded in terms of the $\left\vert \mathbf{f}\right\rangle $ bases as 
\begin{equation}
\sigma _{f}\left( \left\vert \mathbf{e}\right\rangle \right) =\sigma
_{f}\left( \left( \left\vert \mathbf{e}\right\rangle \left\langle \mathbf{f}%
\right\vert \right) \left\vert \mathbf{f}\right\rangle \right) =\sigma
_{f}\left( \left\vert \mathbf{f}\right\rangle \mathcal{R}_{f}\left(
\left\vert \mathbf{e}\right\rangle \left\langle \mathbf{f}\right\vert
\right) \right) =\left\vert \mathbf{f}\right\rangle \overline{\mathcal{R}%
_{f}\left( \left\vert \mathbf{e}\right\rangle \left\langle \mathbf{f}%
\right\vert \right) }=\left\vert \mathbf{f}\right\rangle \overline{\mathbb{U}%
_{fe}}.
\end{equation}%
In particular each basis element has the expansion 
\begin{equation}
\sigma _{f}\left( \left\vert e_{k}\right\rangle \right) =\sigma _{f}\left(
\sum_{1\leq j\leq r}\left\vert f_{j}\right\rangle \left\langle
f_{j}\right\vert \left\vert e_{k}\right\rangle \right) =\sigma _{f}\left(
\sum_{1\leq j\leq r}\left\vert f_{j}\right\rangle \mathcal{R}_{f}\left(
\left\vert \mathbf{e}\right\rangle \left\langle \mathbf{f}\right\vert
\right) _{jk}\right) =\sum_{1\leq j\leq r}\left\vert f_{j}\right\rangle 
\overline{\mathcal{R}_{f}\left( \left\vert \mathbf{e}\right\rangle
\left\langle \mathbf{f}\right\vert \right) }_{jk}=\sum_{1\leq j\leq
r}\left\vert f_{j}\right\rangle \overline{\left\langle \left. \mathbf{f}%
\right\vert \mathbf{e}\right\rangle }_{jk}=\sum_{1\leq j\leq r}\left\vert
f_{k}\right\rangle \overline{\left( u_{fe}\right) }_{jk}.
\end{equation}

\begin{proposition}
The representation $\mathcal{R}_{f}\left( \Omega _{e}\left( A\right) \right) 
$ of $\Omega _{e}\left( A\right) $ in a basis $\left\{ \left. \left\vert
f_{j}\right\rangle \right\vert 1\leq j\leq r\right\} $ that is not in the
same equivalence class as $\left\{ \left. \left\vert e_{j}\right\rangle
\right\vert 1\leq j\leq r\right\} $ is \textbf{not }the complex conjugate
matrix $\overline{\mathcal{R}_{f}\left( A\right) }$ nor an orthogonal
transformation of $\overline{\mathcal{R}_{f}\left( A\right) }$ i.e.%
\begin{equation}
\mathcal{R}_{f}\left( \Omega _{e}\left( A\right) \right) \equiv \mathcal{R}%
_{f}\left( \sigma _{e}A\sigma _{e}\right) \neq \overline{\mathcal{R}%
_{f}\left( A\right) }.
\end{equation}
\end{proposition}

\begin{proof}
\begin{align}
\Omega_{e}\left( A\right) & =\sigma_{e}A\sigma_{e}=\left\vert \mathbf{f}%
\right\rangle \left\langle \mathbf{f}\right\vert \Omega_{e}\left( A\right)
\left\vert \mathbf{f}\right\rangle \left\langle \mathbf{f}\right\vert
=\left\vert \mathbf{f}\right\rangle \mathcal{R}_{f}\left( \Omega_{e}\left(
A\right) \right) \left\langle \mathbf{f}\right\vert \left\vert \mathbf{e}%
\right\rangle =\left\vert \mathbf{e}\right\rangle \left\langle \mathbf{e}%
\right\vert \Omega_{e}\left( A\right) \left\vert \mathbf{e}\right\rangle
\left\langle \mathbf{e}\right\vert \\
& =\left\vert \mathbf{e}\right\rangle \mathcal{R}_{e}\left( \Omega
_{e}\left( A\right) \right) \left\langle \mathbf{e}\right\vert =\left\vert 
\mathbf{e}\right\rangle \overline{\mathcal{R}_{e}\left( A\right) }%
\left\langle \mathbf{e}\right\vert =\left\vert \mathbf{f}\right\rangle 
\mathbb{U}_{fe}\overline{\mathcal{R}_{e}\left( A\right) }\mathbb{U}%
_{fe}^{\dagger}\left\langle \mathbf{f}\right\vert
\end{align}
thus 
\begin{equation}
\mathcal{R}_{f}\left( \Omega_{e}\left( A\right) \right) =\mathbb{U}_{ef}%
\overline{\mathcal{R}_{e}\left( A\right) }\mathbb{U}_{ef}^{\dagger }=\mathbb{%
U}_{ef}\overline{\mathbb{U}_{ef}^{\dagger}}\overline{\mathcal{R}_{f}\left(
A\right) }\overline{\mathbb{U}_{ef}}\mathbb{U}_{ef}^{\dagger }=\mathbb{U}%
_{ef}\mathbb{U}_{ef}^{t}\overline{\mathcal{R}_{f}\left( A\right) }\overline{%
\mathbb{U}_{ef}}\mathbb{U}_{ef}^{\dagger}=\mathbb{U}_{ef}\mathbb{U}_{ef}^{t}%
\overline{\mathcal{R}_{f}\left( A\right) }\left( \mathbb{U}_{ef}\mathbb{U}%
_{ef}^{t}\right) ^{\dagger}
\end{equation}
i.e. 
\begin{equation}
\mathcal{R}_{f}\left( \Omega_{e}\left( A\right) \right) =\mathbb{U}_{ef}%
\overline{\mathcal{R}_{e}\left( A\right) }\mathbb{U}_{ef}^{\dagger }=\mathbb{%
U}_{ef}\mathbb{U}_{ef}^{t}\overline{\mathcal{R}_{f}\left( A\right) }\left( 
\mathbb{U}_{ef}\mathbb{U}_{ef}^{t}\right) ^{\dagger},
\end{equation}
which can also be written as 
\begin{equation}
\mathcal{R}_{f}\left( \Omega_{e}\left( A\right) \right) =\left\langle \left. 
\mathbf{f}\right\vert \mathbf{e}\right\rangle \left\langle \left. \mathbf{f}%
\right\vert \mathbf{e}\right\rangle ^{t}\overline{\mathcal{R}_{f}\left(
A\right) }\left( \left\langle \left. \mathbf{f}\right\vert \mathbf{e}%
\right\rangle \left\langle \left. \mathbf{f}\right\vert \mathbf{e}%
\right\rangle ^{t}\right) ^{\dagger},
\end{equation}
as would be expected by the change of basis relationship 
\begin{equation}
\mathcal{R}_{f}\left( A\right) =\mathbb{U}_{fe}\mathcal{R}_{e}\left(
A\right) \mathbb{U}_{fe}^{\dagger}.
\end{equation}
However 
\begin{equation}
\overline{\mathcal{R}_{f}\left( A\right) }=\overline{\mathbb{U}_{fe}}%
\overline{\mathcal{R}_{e}\left( A\right) }\overline{\mathbb{U}_{fe}^{\dagger}%
}\neq\mathbb{U}_{fe}\overline{\mathcal{R}_{e}\left( A\right) }\mathbb{U}%
_{fe}^{\dagger}\text{ }=\mathcal{R}_{f}\left( \Omega_{e}\left( A\right)
\right) \text{ unless }\overline{\mathbb{U}_{fe}}=\mathbb{U}_{fe}.
\end{equation}
$\lceil$%
\begin{equation}
\mathcal{R}_{f}\left( \Omega_{e}\left( A\right) \right) =\mathbb{U}_{fe}%
\overline{\mathcal{R}_{e}\left( A\right) }\mathbb{U}_{fe}^{\dagger }=%
\overline{\mathbb{U}_{fe}^{\dagger}}\mathbb{U}_{fe}\overline{\mathcal{R}%
_{e}\left( A\right) }\mathbb{U}_{fe}^{\dagger}\overline{\mathbb{U}_{fe}}=%
\mathbb{U}_{fe}^{t}\mathbb{U}_{fe}\overline{\mathcal{R}_{e}\left( A\right) }%
\overline{\mathbb{U}_{fe}^{t}}\overline{\mathbb{U}_{fe}}=\mathbb{U}_{fe}^{t}%
\mathbb{U}_{fe}\overline{\mathcal{R}_{e}\left( A\right) }\left( \mathbb{U}%
_{fe}^{t}\mathbb{U}_{fe}\right) ^{\dagger}
\end{equation}
$\rfloor$
\end{proof}

Each equivalence class $\left[ \mathbf{e}\right] $ gives a unique complex
conjugate and transpose. An operator $A$ is said to be $\left[ \mathbf{e}%
\right] -$\emph{complex symmetric} if 
\begin{equation}
A^{t_{e}}=A
\end{equation}
and $\left[ \mathbf{e}\right] -$\emph{real }if 
\begin{equation}
A^{\ast_{e}}=A.
\end{equation}
The representations of $\left[ \mathbf{e}\right] -$\emph{complex symmetric}
and $\left[ \mathbf{e}\right] -$\emph{real }operators have complex symmetric
and real representations wrt bases belonging $\left[ \mathbf{e}\right] $
respectively.

One can show that in \emph{finite dimensions \underline{\emph{all}} operators%
} $X\in \mathcal{B}\left( \mathfrak{H}\right) $ \emph{are similar to }a
corresponding $\left[ \mathbf{e}\right] -$\emph{complex symmetric operator}
i.e. 
\begin{equation}
X=TA_{X}T^{-1},
\end{equation}%
where 
\begin{equation}
A_{X}^{t_{e}}=A_{X}
\end{equation}%
for some $\left[ \mathbf{e}\right] .$

\begin{remark}
Note that $A^{\ast _{e}}$and $A^{t_{e}}$ are bases dependent but their
composition (which is commutative) is \emph{NOT }i.e 
\begin{equation}
A^{\dagger }=A^{\ast _{e}t_{e}}.
\end{equation}
\end{remark}

We denote the set of all equivalence classes, $\left[ \mathbf{e}\right] ,$
as $\mathfrak{C}\left( \mathfrak{H}\right) $ 
\begin{equation}
\mathfrak{C}\left( \mathfrak{H}\right) =\left\{ \left. \overset{}{u\left[ 
\mathbf{e}\right] }\right\vert u\in U\left( \mathfrak{H}\right) /O\left( 
\mathfrak{H,}\mathbb{R}\right) \equiv U\left( r\right) /O\left( r\mathfrak{,}%
\mathbb{R}\right) =\frac{O\left( 2r,\mathbb{R}\right) \cap Sp\left( 2r,%
\mathbb{R}\right) }{O\left( r\mathfrak{,}\mathbb{R}\right) }\right\} .
\end{equation}

\subsection{Complex Conjugations in $C^{\ast}$-algebras}

\begin{definition}
Let $\Omega\in\mathcal{B}_{anti}\left( \mathcal{A}\right) $ be
\end{definition}

\begin{enumerate}
\item a complex conjugation acting in $\mathcal{A}$ i.e. 
\begin{align}
\Omega & :\mathcal{A}\longrightarrow\mathcal{A} \\
\Omega\left( \alpha A+\beta B\right) & =\overline{\alpha}\Omega\left(
A\right) +\overline{\beta}\Omega\left( B\right) \\
\Omega^{2} & =I_{\mathcal{A}},
\end{align}

\item a $\ast$ anti-isomorphism i.e. 
\begin{equation}
\Omega\left( A\right) ^{\dagger}=\Omega\left( A^{\dagger}\right) ,
\end{equation}

\item an anti-isometry or an anti-partial isometry if $\mathcal{A}$ is odd
dimensional i.e. 
\begin{equation}
\left\Vert \Omega\left( A\right) \right\Vert _{\mathcal{A}}=\left\Vert
A\right\Vert _{\mathcal{A}},
\end{equation}

\item a homomorphism i.e. 
\begin{equation}
\Omega\left( A\right) \Omega\left( B\right) =\Omega\left( AB\right) ,
\end{equation}
then $\Omega$ is a $c^{\ast}$\emph{-algebra complex conjugation}.

\item $\Omega$ is inner if 
\begin{equation}
\Omega\left( A\right) =\sigma A\sigma,
\end{equation}
where $\sigma$ is anti-unitary.
\end{enumerate}

As $\Omega$ $\in\mathcal{B}_{anti}\left( \mathcal{A}\right) $ one can
consider its spectrum $\mathfrak{s}\left( \Omega\right) $, which is $\left\{
1,-1\right\} $ if $\mathcal{A}$ is even dimensional and leads to the direct
sum decomposition

\begin{equation}
\mathcal{A}=\mathcal{A}_{\Omega+}\oplus_{\mathbb{R}}\mathcal{A}_{\Omega-}\neq%
\mathcal{A}_{\Omega+}\oplus_{\mathbb{C}}\mathcal{A}_{\Omega-},
\end{equation}
and a $1-1$ correspondence to the orthogonal projectors $\left\{ P_{\Omega
+},P_{\Omega-}\right\} \in\mathcal{B}\left( \mathcal{A}_{\Omega+}\oplus_{%
\mathbb{R}}\mathcal{A}_{\Omega-}\right) $ onto these subspacess. If $%
\mathcal{A}$ has a basis this correspondance can be extended to equivalent
classes of bases of $\mathcal{B}\left( \mathcal{A}\right) .$ (The complex
dimension of $\mathcal{A}_{\Omega+}\oplus_{\mathbb{C}}\mathcal{A}_{\Omega-}$
is $2r,$while the complex dimension of $\mathcal{A}_{\Omega+}\oplus _{%
\mathbb{R}}\mathcal{A}_{\Omega-}$ is $r.)$

Multiplication by " $i$ $"$ in $\mathcal{A}$ maps these real subspaces to
each other i.e. 
\begin{align}
i & :\mathcal{A}_{\Omega+}\rightarrow\mathcal{A}_{\Omega-} \\
i & :\mathcal{A}_{\Omega-}\rightarrow-\mathcal{A}_{\Omega+}
\end{align}
\emph{for every complex conjugation }$\Omega.$

For a particular conjugation $\Omega \in \mathcal{B}\left( \mathcal{A}%
_{\Omega +}\oplus _{\mathbb{R}}\mathcal{A}_{\Omega -}\right) \cong \mathcal{B%
}_{anti}\left( \mathcal{A}\right) $ 
\begin{align}
\Omega & :\mathcal{A}_{\Omega +}\rightarrow \mathcal{A}_{\Omega +} \\
\Omega & :\mathcal{A}_{\Omega -}\rightarrow \mathcal{A}_{\Omega -}
\end{align}%
one can construct a real linear map $J_{\Omega }$ 
\begin{equation}
J_{\Omega }:\mathcal{A}_{\Omega +}\oplus _{\mathbb{R}}\mathcal{A}_{\Omega -}%
\mathcal{\rightarrow A}_{\Omega +}\oplus _{\mathbb{R}}\mathcal{A}_{\Omega -}
\end{equation}%
defined by 
\begin{align}
J_{\Omega }& :\mathcal{A}_{\Omega +}\rightarrow \mathcal{A}_{\Omega -} \\
J_{\Omega }& :\mathcal{A}_{\Omega -}\rightarrow -\mathcal{A}_{\Omega +}
\end{align}%
with the properties 
\begin{equation}
J_{\Omega }^{2}=I_{\mathcal{A}}
\end{equation}%
and%
\begin{equation}
J_{\Omega }\Omega =-\Omega J_{\Omega }\equiv \left\{ J_{\Omega },\Omega
\right\} .
\end{equation}%
$\lceil $%
\begin{align}
J_{\Omega }\Omega \left( i\right) & =J_{\Omega }\left( -i\right) =1 \\
\Omega J_{\Omega }\left( i\right) & =J_{\Omega }\left( -1\right) =-1
\end{align}%
$\rfloor $

Hence $J_{\Omega }$\emph{\ is a complexification.} In general 
\begin{align}
J_{\Omega ^{\prime }}& :\mathcal{A}_{\Omega +}\nrightarrow \mathcal{A}%
_{\Omega -} \\
J_{\Omega ^{\prime }}& :\mathcal{A}_{\Omega -}\nrightarrow -\mathcal{A}%
_{\Omega +}
\end{align}%
unless 
\begin{equation}
J_{\Omega ^{\prime }}=uJ_{\Omega }u^{\dagger };u\in O\left( 2r,\mathbb{R}%
\right) .
\end{equation}

\subsubsection{GNS Representation of Complex Conjugations and Complex
Conjugate GNS\ Representations}

\paragraph{States and Complex Conjugations}

\begin{definition}
A state $f\in \mathcal{S}\left( \mathcal{A}\right) $ is compatible with a
complex conjugation $\Omega $ if 
\begin{equation}
f\left( \Omega \left( A\right) \right) =f\left( A\right) \,\forall A\in 
\mathcal{F}
\end{equation}%
i.e. the state $f$ is $\Omega $ invariant.
\end{definition}

A concrete example is given by normal states $f\in \mathcal{S}\left( 
\mathcal{B}\left( \mathfrak{H}\right) \right) \cong \mathcal{S}\left(
M\left( r,\mathbb{C}\right) \right) $

\begin{enumerate}
\item The kernal is given by 
\begin{align}
\ker\Phi_{f_{D}} & =\left\{ A\left\vert f_{D}\left( \Omega_{e}\left(
A^{\dagger}A\right) \right) =f_{D}\left( \Omega_{e}\left( A\right)
^{\dagger}\Omega_{e}\left( A\right) \right) =\func{Tr}\left\{
\Omega_{e}\left( A^{\dagger}A\right) D\right\} \right. \right\} \\
& =\left\{ A\left\vert f_{D}\left( A^{\dagger}A\right) =\func{Tr}\left\{
A^{\dagger}AD\right\} \right. \right\} =0.
\end{align}

\item 
\begin{equation}
f_{D}\left( A^{\dagger}A\right) =\sum_{1\leq j,k\leq r}\mathcal{R}_{e}\left(
A^{\dagger}A\right) _{jk}\mathcal{R}_{e}\left( D\right) _{kj}=\sum_{1\leq
j,k\leq r}\overline{\mathcal{R}_{e}\left( A^{\dagger }A\right) }_{kj}%
\mathcal{R}_{e}\left( D\right) _{kj}
\end{equation}
and%
\begin{align}
f\left( \Omega_{e}\left( A\right) ^{\dagger}\Omega_{e}\left( A\right)
\right) & =f\left( \Omega_{e}\left( A^{\dagger}\right) \Omega_{e}\left(
A\right) \right) =f\left( \Omega_{e}\left( A^{\dagger}A\right) \right) =%
\func{Tr}\left\{ \mathcal{R}_{e}\left( \Omega_{e}\left( A^{\dagger}A\right)
\right) \mathcal{R}_{e}\left( D\right) \right\} \\
& =\sum_{1\leq j,k\leq r}\overline{\mathcal{R}_{e}\left( A^{\dagger
}A\right) }_{jk}\mathcal{R}_{e}\left( D\right) _{kj}=\sum_{1\leq j,k\leq r}%
\mathcal{R}_{e}\left( A^{\dagger}A\right) _{kj}\mathcal{R}_{e}\left(
D\right) _{kj}.
\end{align}

\item The expectation values of the complex conjugate operators are 
\begin{equation}
f_{D}\left( \Omega_{e}\left( A\right) \right) =\func{Tr}\left\{ \mathcal{R}%
_{e}\left( \Omega_{e}\left( A\right) \right) ^{\dagger }\mathcal{R}%
_{e}\left( D\right) \right\} =\func{Tr}\left\{ \left( \overline{\mathcal{R}%
_{e}\left( A\right) }\right) ^{\dagger}\mathcal{R}_{e}\left( D\right)
\right\} =\func{Tr}\left\{ \mathcal{R}_{e}\left( A\right) ^{t}\mathcal{R}%
_{e}\left( D\right) \right\} =\sum_{1\leq j,k\leq r}\mathcal{R}_{e}\left(
A\right) _{kj}\mathcal{R}_{e}\left( D\right) _{kj}.
\end{equation}

\item The representation of the $e$-conjugate of $A$ in the $\mathcal{R}_{f}$
representation is related to the $\mathcal{R}_{e}$ representation of $A$ by 
\begin{equation}
\overline{\mathcal{R}_{f}\left( A\right) }=\overline{\mathbb{U}_{ef}}%
\overline{\mathcal{R}_{e}\left( A\right) }\overline{\mathbb{U}_{ef}^{\dagger}%
}\neq\mathbb{U}_{ef}\overline{\mathcal{R}_{e}\left( A\right) }\mathbb{U}%
_{ef}^{\dagger}\text{ }=\mathcal{R}_{f}\left( \Omega_{e}\left( A\right)
\right) \text{ unless }\overline{\mathbb{U}_{ef}}=\mathbb{U}_{ef},
\end{equation}
where%
\begin{equation}
\mathcal{R}_{f}\left( \Omega_{e}\left( A\right) \right) =\mathbb{U}_{ef}%
\overline{\mathcal{R}_{e}\left( A\right) }\mathbb{U}_{ef}^{\dagger}.
\end{equation}

\item 
\begin{align}
f_{D}\left( \Omega_{e}\left( A\right) \right) & =\func{Tr}\left\{ \mathcal{R}%
_{f}\left( \Omega_{e}\left( A\right) \right) ^{\dagger}\mathcal{R}_{f}\left(
D\right) \right\} =\func{Tr}\left\{ \left( \mathbb{U}_{ef}\overline{\mathcal{%
R}_{e}\left( A\right) }\mathbb{U}_{ef}^{\dagger}\right) ^{\dagger}\mathcal{R}%
_{f}\left( D\right) \right\} \\
& =\func{Tr}\left\{ \mathbb{U}_{ef}\overline{\mathcal{R}_{e}\left( A\right) }%
^{\dagger}\mathbb{U}_{ef}^{\dagger}\mathcal{R}_{f}\left( D\right) \right\} =%
\func{Tr}\left\{ \mathbb{U}_{ef}\mathcal{R}_{e}\left( A\right) ^{t}\mathbb{U}%
_{ef}^{\dagger}\mathcal{R}_{f}\left( D\right) \right\} .
\end{align}
\end{enumerate}

Let $f\in\mathcal{S}\left( \mathcal{A}\right) $ be a state of $\mathcal{A}$
and consider the GNS\ representation $\Pi_{f}\left( \mathcal{A}\right) $
carried by Hilbert space $\Phi_{f}\left( \mathcal{A}\right) .$ The
conjugation, $\Omega,$ is \emph{campatible} with the representation $\Pi_{f}$
if 
\begin{equation}
\Omega:\ker\left\{ \Phi_{f}\right\} \rightarrow\ker\left\{ \Phi _{f}\right\}
\end{equation}
i.e. 
\begin{equation}
A\in\ker\left\{ \Phi_{f}\right\} \Leftrightarrow\Omega\left( A\right)
\in\ker\left\{ \Phi_{f}\right\}
\end{equation}
and $f$ is $\Omega$ complex conjugate invariant.

As both $\Pi_{f}$ and $\Omega$ are homomorphisms so is $\Pi_{f}\circ\Omega$,
which belongs to $\mathcal{B}_{anti}\left( \mathcal{B}\left( \Phi_{f}\left( 
\mathcal{A}\right) \right) \right) $ and produces the \emph{complex
conjugate representation}%
\begin{equation}
\left[ \Omega\left( \Pi_{f}\right) \right] \left( \mathcal{A}\right) \equiv%
\left[ \Pi_{f}\circ\Omega\right] \left( \mathcal{A}\right) \subseteq\mathcal{%
B}\left( \Phi_{f}\left( \mathcal{A}\right) \right) ,
\end{equation}
i.e. 
\begin{equation}
\Omega\circ\Pi_{f}=\Pi_{f}\circ\Omega.
\end{equation}
One can observe that $\Omega$ also produces a complex conjugation of $\Phi
_{f}\left( \mathcal{A}\right) $ by 
\begin{equation}
\sigma\left( \Phi_{f}\left( \mathcal{A}\right) \right) =\Phi_{f}\left(
\Omega\left( \mathcal{A}\right) \right) ,
\end{equation}
i.e.%
\begin{equation}
\sigma\circ\Phi_{f}=\Phi_{f}\circ\Omega.
\end{equation}

\begin{proposition}
\begin{equation}
\Omega\left( \Pi_{f}\right) =\Pi_{f}\left( \Omega\right) =\sigma\Pi
_{f}\sigma;\,\sigma\in\mathcal{B}_{anti}\left( \Phi_{f}\left( \mathcal{A}%
\right) \right) ,\Pi_{f}:\mathcal{A}\rightarrow\mathcal{B}\left( \Phi
_{f}\left( \mathcal{A}\right) \right)
\end{equation}
\end{proposition}

\begin{proof}
The conjugation $\Omega $ defines a conjugation on $\Phi _{f}\left( \mathcal{%
A}\right) $ by the following 
\begin{equation}
\Omega \left( \alpha A+\beta B\right) \Phi _{f}\left( I_{\mathcal{A}}\right)
=\Phi _{f}\left( \Omega \left( \alpha A+\beta B\right) \right) =\overline{%
\alpha }\Phi _{f}\left( \Omega \left( A\right) \right) +\overline{\beta }%
\Phi _{f}\left( \Omega \left( B\right) \right) =\Pi _{f}\left( \Omega
\right) \Phi _{f}\left( \alpha A+\beta B\right) \,\forall A,B\in \mathcal{A},
\end{equation}%
which shows that $\Pi _{f}\left( \Omega \right) $ is antilinear. It is
trivially an involution as 
\begin{equation}
\Pi _{f}\left( \Omega ^{2}\right) =\Pi _{f}\left( \Omega \right) ^{2}.
\end{equation}%
A complex conjugation $\mathfrak{c}$ is defined in $\mathcal{B}\left( \Phi
_{f}\left( \mathcal{A}\right) \right) $ by 
\begin{align}
\mathfrak{c}\left( \Phi _{f}\left( \alpha A+\beta B\right) \right) & =%
\mathfrak{c}\left( \Phi _{f}\left( \alpha A\right) \right) +\mathfrak{c}%
\left( \Phi _{f}\left( \beta B\right) \right) =\overline{\alpha }\mathfrak{c}%
\left( \Phi _{f}\left( A\right) \right) +\overline{\beta }\mathfrak{c}\left(
\Phi _{f}\left( B\right) \right) \,\forall A,B\in \mathcal{A} \\
\mathfrak{c}^{2}& =I_{\mathcal{A}}.
\end{align}%
Thus $\Pi _{f}\left( \Omega \right) $ is a complex conjugation $\mathfrak{c}$
on $\mathcal{B}\left( \Phi _{f}\left( \mathcal{A}\right) \right) .$
\end{proof}

\subsection{Complex Structure}

As discussed in the preceeding a \emph{complex structure} $\mathcal{J}_{%
\mathcal{V}}$ on a real vector space $V$ is a real linear map on the direct
sum of $V$ with a real space $W$ of equal dimension with the following
properties 
\begin{align}
\mathcal{V}_{\mathbb{R}} & =V\oplus_{\mathbb{R}}W=\mathcal{V}+\mathcal{W}=%
\mathcal{V}+\mathcal{V}^{c} \\
\mathcal{J}_{\mathcal{V}} & :V\oplus_{\mathbb{R}}W\rightarrow V\oplus _{%
\mathbb{R}}W \\
\mathcal{J}_{\mathcal{V}}^{2} & =-I_{V\oplus_{\mathbb{R}}W} \\
\mathcal{J}_{\mathcal{V}} & :\mathcal{V}\rightarrow\mathcal{V}^{c} \\
\mathcal{J}_{\mathcal{V}} & :\mathcal{V}^{c}\rightarrow\mathcal{V},
\end{align}
where 
\begin{gather}
\mathcal{J}_{\mathcal{V}}\left( v_{w},0\right) =\left( 0,w_{v}\right)
\leftrightarrow w_{v} \\
\mathcal{J}_{\mathcal{V}}\left( 0,w_{v}\right) =\left( -v_{w},0\right)
\leftrightarrow-v_{w} \\
\left( v_{w},0\right) \in\mathcal{V},\left( 0,w_{v}\right) \in \mathcal{V}%
^{c},v_{w}\in V,w_{v}\in W.
\end{gather}
\emph{\ }$V$ and\emph{\ }$W$\emph{\ do not have to be disjoint} in fact $V$%
\emph{\ }can equal\emph{\ }$W$ and therefore not complementary$.$ If there
is an inner product in $V\oplus_{\mathbb{R}}W$ then we restrict to $\mathcal{%
J}_{\mathcal{V}}$ that are antisymmetric i.e. 
\begin{equation}
\mathcal{J}_{\mathcal{V}}^{t}=-\mathcal{J}_{\mathcal{V}}
\end{equation}
and thus 
\begin{equation}
\mathcal{J}_{\mathcal{V}}\in O\left( V\oplus_{\mathbb{R}}W,\mathbb{R}\right)
.
\end{equation}
Multiplication by the imaginary unit $i$ and thus the complexification of
the real space $\mathcal{V},$ which is ismorphic to $V$ is defined by 
\begin{equation}
i\left( v,w\right) \overset{def}{=}\mathcal{J}_{\mathcal{V}}\left(
v,w\right) =\left( -v_{w},w_{v}\right) ,
\end{equation}
which gives 
\begin{align}
i\left( v_{w},0\right) & =\mathcal{J}_{\mathcal{V}}\left( v_{w},0\right)
=\left( 0,w_{v}\right) =\left( 0,iv_{w}\right) \\
i\left( 0,w_{v}\right) & =\mathcal{J}_{\mathcal{V}}\left( 0,w_{v}\right)
=\left( -v_{w},0\right) =i\left( 0,iv_{w}\right) \\
w_{v} & =iv_{w}.
\end{align}
This shows that $\mathcal{J}_{\mathcal{V}}$ induces a map $J_{V}$ 
\begin{align}
J_{V} & :V\rightarrow W \\
J_{V} & :J_{V}W\rightarrow V
\end{align}
given by 
\begin{align}
\left( 0,J_{V}v_{w}\right) & =\mathcal{J}_{\mathcal{V}}\left( v_{w},0\right)
=\left( 0,w_{v}\right) \\
\mathcal{J}_{\mathcal{V}}\left( 0,J_{V}v_{w}\right) & =\mathcal{J}_{\mathcal{%
V}}\left( 0,w_{v}\right) =\left( -v_{w},0\right)
\end{align}
and thus 
\begin{equation}
J_{V}v_{w}=w_{v};v_{w}\in V,w_{v}\in W.
\end{equation}

Introducing a basis for $\mathcal{V}$ one obtains

\begin{equation*}
\mathcal{J}_{V}\left\vert \left( \mathbf{e,0}\right) \mathbf{,}\left( 
\mathbf{0},\mathbf{f}\right) \right) =\left\vert \left( \mathbf{e,0}\right) 
\mathbf{,}\left( \mathbf{0},\mathbf{f}\right) \right) \mathcal{R}_{e}\left( 
\mathcal{J}_{V}\right) =\left\vert \left( \mathbf{e,0}\right) \mathbf{,}%
\left( \mathbf{0},\mathbf{f}\right) \right) \left( 
\begin{array}{cc}
0 & -\mathbb{I}_{r} \\ 
\mathbb{I}_{r} & 0%
\end{array}
\right) =\left\vert \left( \mathbf{0},\mathbf{f}\right) ,\left( -\mathbf{e,0}%
\right) \right)
\end{equation*}
and 
\begin{equation*}
i\left\vert \left( \mathbf{e,0}\right) \mathbf{,}\left( \mathbf{0},\mathbf{f}%
\right) \right) =\mathcal{J}_{V}\left\vert \left( \mathbf{e,0}\right) 
\mathbf{,}\left( \mathbf{0},\mathbf{f}\right) \right) =\left\vert \mathcal{J}%
_{V}\left( \mathbf{e,0}\right) \mathbf{,}\mathcal{J}_{V}\left( \mathbf{0},%
\mathbf{f}\right) \right) =\left\vert i\left( \mathbf{e,0}\right) \mathbf{,}%
i\left( \mathbf{0},\mathbf{f}\right) \right) =\left\vert \left( \mathbf{0},%
\mathbf{f}\right) ,\left( -\mathbf{e,0}\right) \right) .
\end{equation*}

Once we have made the identification $W\equiv J_{V}V$ one can write 
\begin{equation}
\mathcal{V}_{\mathbb{C}}=V\oplus_{\mathbb{R}}W=V\oplus_{\mathbb{R}}J_{V}V.
\end{equation}
The map $J_{V}$%
\begin{align}
J_{V} & :V\rightarrow W=J_{V}V \\
J_{V} & :W=J_{V}V\rightarrow V
\end{align}
has the following action on basis elements 
\begin{align}
J_{V}\left( e_{1}\right) & \equiv J_{V}e_{1}=ie_{1} \\
J_{V}\left( ie_{1}\right) & \equiv J_{V}J_{V}e_{1}=-e_{1}.
\end{align}

The following isomorphisms $\mathcal{T}_{V\mathcal{V}},\mathcal{T}_{\mathcal{%
V}V},\mathcal{T}_{W\mathcal{W}},\mathcal{T}_{\mathcal{W}W}$ formalize the
relationship between the direct sums and the spaces participating in the
direct sum 
\begin{align}
\mathcal{T}_{V\mathcal{V}} & :V\rightarrow\mathcal{V} \\
\mathcal{T}_{V\mathcal{V}} & :\mathcal{V}\rightarrow V \\
\mathcal{T}_{\mathcal{W}W} & :W\rightarrow\mathcal{W} \\
\mathcal{T}_{W\mathcal{W}} & :\mathcal{W}\rightarrow W,
\end{align}
where%
\begin{align}
V & =\left\{ v\in V\right\} \cong\left\{ \left( v,0\right) \in\mathcal{V}%
\right\} =\mathcal{V} \\
W & =\left\{ w\in W\right\} \cong\left\{ \left( 0,w\right) \in\mathcal{V}%
^{c}\right\} =\mathcal{V}^{c}\equiv\mathcal{W}.
\end{align}
The map $\mathcal{J}_{\mathcal{V}}$ can then be extended by isomorphism to $%
J_{V}$ that is defined on $\mathcal{V}$ and $\mathcal{W}$ i.e. 
\begin{align}
\left. J_{V}\right\vert _{V} & =\mathcal{T}_{V\mathcal{V}}\mathcal{J}_{%
\mathcal{V}}\mathcal{T}_{V\mathcal{V}}^{-1} \\
\left. J_{V}\right\vert _{W} & =\mathcal{T}_{W\mathcal{W}}\mathcal{J}_{%
\mathcal{V}}\mathcal{T}_{W\mathcal{W}}^{-1}
\end{align}
to obtain 
\begin{align}
J_{V} & :V\rightarrow W \\
J_{V} & :W\rightarrow V,
\end{align}
where 
\begin{gather}
J_{V}v_{w}=w_{v} \\
J_{V}w_{v}=-v_{w} \\
v_{w}\in V,w_{v}\in W
\end{gather}
and thus 
\begin{equation}
W=J_{V}V.
\end{equation}
\emph{The complexifiction of }$V$\emph{\ induced by }$W$\emph{\ is then
defined by} 
\begin{equation}
iv\overset{def}{=}J_{V}v\in W,
\end{equation}
so that 
\begin{equation}
iv\in iV\leftrightarrow J_{V}v\in W,
\end{equation}
which gives in particular%
\begin{gather}
i\left( e_{1},0\right) =\left( 0,ie_{1}\right) =\mathcal{J}_{V}\left(
e_{1},0\right) =\left( 0,f_{1}\right) =\left( 0,J_{V}e_{1}\right) \\
i\left( 0,f_{1}\right) =\left( -e_{1},0\right) =\mathcal{J}_{V}\left(
0,f_{1}\right) =\left( -e_{1},0\right) =\left( J_{V}f_{1},0\right) .
\end{gather}
The complexification of $V$ can thus be expressed as 
\begin{equation}
\mathcal{V}_{\mathbb{C}}=V\oplus J_{V}V\equiv V\oplus iV\cong V+J_{V}V\equiv
V+iV
\end{equation}
as $V$ and $J_{V}V$ are disjoint as real spaces i.e. 
\begin{equation}
V\cap J_{V}V=\left\{ 0\right\} .
\end{equation}
In terms of bases of $V+iV$ 
\begin{equation}
\left\vert \mathbf{e,f}\right) =\left\vert \mathbf{e,}J_{V}\mathbf{e}\right)
\leftrightarrow\left\vert \mathbf{e,}i\mathbf{e}\right)
\end{equation}
and 
\begin{equation}
J_{V}\left\vert \mathbf{e,f}\right) =J_{V}\left\vert \mathbf{e,}i\mathbf{e}%
\right) =i\left\vert \mathbf{e,}i\mathbf{e}\right) =\left\vert \mathbf{e,}i%
\mathbf{e}\right) \mathcal{R}_{e}\left( J_{V}\right) =\left\vert \mathbf{e,}i%
\mathbf{e}\right) \left( 
\begin{array}{cc}
0 & -\mathbb{I}_{r} \\ 
\mathbb{I}_{r} & 0%
\end{array}
\right) ,
\end{equation}
\emph{note that} 
\begin{equation}
\mathcal{R}_{e}\left( J_{V}\right) =\mathcal{R}_{e}\left( \mathcal{J}_{%
\mathcal{V}}\right) .
\end{equation}
$\lceil$Using this definition one obtains from the definition of
multiplication by $i$ 
\begin{gather}
\mathcal{J}_{\mathcal{V}}\left\vert \left( e_{1},0\right) ,\left(
0,ie_{1}\right) \right) =\left\vert \left( e_{1},0\right) ,\left(
0,ie_{1}\right) \right) \mathcal{R}_{e}\left( \mathcal{J}_{\mathcal{V}%
}\right) \\
=\left\vert \left( e_{1},0\right) ,\left( 0,ie_{1}\right) \right) \left( 
\begin{array}{cc}
0 & -1 \\ 
1 & 0%
\end{array}
\right) =\left\vert \left( 0,ie_{1}\right) ,\left( -e_{1},0\right) \right) ,
\end{gather}
which one can compare to%
\begin{gather}
J_{V}\left\vert e_{1},ie_{1}\right) =\left\vert e_{1},ie_{1}\right) \mathcal{%
R}_{e}\left( J_{V}\right) \\
=\left\vert e_{1},ie_{1}\right) \left( 
\begin{array}{cc}
0 & -1 \\ 
1 & 0%
\end{array}
\right) =\left\vert ie_{1},e_{1}\right) .
\end{gather}
$\rfloor$

The crucial element in complexification is that $W=J_{V}V$ such that $V$ and 
$J_{V}V$ are disjoint real subspaces together with the identity 
\begin{equation}
J_{V}^{2}=-I_{V\oplus J_{V}V}=-I_{2r}.
\end{equation}

Hence to summarize: complexification of $V$ induced by a space $W$ of the
same dimension is produced by defining multiplication by $i$ by 
\begin{equation}
iv\triangleq J_{V}v=w;v\in V,w\in W.
\end{equation}
$\lceil$%
\begin{align}
J_{V}v & =w;v\in V,w\in W \\
J_{V}w & =-v;v\in V,w\in W.
\end{align}%
\begin{align}
\mathcal{J}_{V} & :\mathcal{V}\rightarrow\mathcal{V}^{c} \\
\mathcal{J}_{V} & :\mathcal{V}^{c}\rightarrow\mathcal{V}
\end{align}%
\begin{align}
J_{V} & :V\rightarrow W \\
J_{V} & :W\rightarrow V
\end{align}%
\begin{equation}
\mathcal{V}_{\mathbb{C}}=\mathcal{V}+\mathcal{J}_{V}\mathcal{V}=V\oplus
J_{V}V\cong V+J_{V}V\equiv V+iV
\end{equation}
\emph{By construction }%
\begin{equation}
V\cap J_{V}V=\left\{ 0\right\}
\end{equation}
\emph{so that} 
\begin{equation}
V\oplus J_{V}V\cong V+J_{V}V=\left\{ z=x+iy;x\in V,iy\in iV\cong J_{V}V,y\in
V\right\} .
\end{equation}
$\rfloor$ 
\begin{equation}
z=x+J_{V}y\leftrightarrow x+iy;x,y\in V
\end{equation}
and 
\begin{equation}
\mathcal{V}_{\mathbb{C}}=V\oplus J_{V}V\cong V_{\mathbb{C}}=V+J_{V}V.
\end{equation}

\subparagraph{Example}

\begin{align}
V & =\mathbb{R-}\limfunc{linspan}\left\{ e_{1},e_{2},e_{3},e_{4}\right\} \\
W & =\mathbb{R-}\limfunc{linspan}\left\{ f_{1},f_{2},f_{3},f_{4}\right\} \\
\mathcal{V} & =\mathbb{R-}\limfunc{linspan}\left\{ \left( e_{1},0\right)
,\left( e_{2},0\right) ,\left( e_{3},0\right) ,\left( e_{4},0\right) \right\}
\\
\mathcal{V}^{c} & =\mathbb{R-}\limfunc{linspan}\left\{ \left( 0,f_{1}\right)
,\left( 0,f_{2}\right) ,\left( 0,f_{3}\right) ,\left( 0,f_{4}\right) \right\}
\\
V\oplus W & =\mathbb{R-}\limfunc{linspan}\left\{ \left( e_{1},0\right)
,\left( e_{2},0\right) ,\left( e_{3},0\right) ,\left( e_{4},0\right) ,\left(
0,f_{1}\right) ,\left( 0,f_{2}\right) ,\left( 0,f_{3}\right) ,\left(
0,f_{4}\right) \right\} \\
& =\mathcal{V}_{1}+\mathcal{V}_{1}^{c}\,\,\mathcal{\,\,\,\,\cong\,\,\,\,}%
\,V+V^{c} \\
V+W & =\mathbb{R-}\limfunc{linspan}\left\{
e_{1},e_{2},e_{3},e_{4},f_{1},f_{2},f_{3},f_{4}\right\} =V\cup W
\end{align}%
\begin{align}
\mathcal{J}_{V}\left( e_{1},0\right) & =\left( 0,f_{1}\right) \Rightarrow
J_{V}e_{1}=f_{5}\overset{\text{def}}{=}ie_{1} \\
\mathcal{J}_{V}\left( e_{2},0\right) & =\left( 0,f_{2}\right) \Rightarrow
J_{V}e_{2}=f_{6}\overset{\text{def}}{=}ie_{2} \\
\mathcal{J}_{V}\left( e_{3},0\right) & =\left( 0,f_{3}\right) \Rightarrow
J_{V}e_{3}=f_{3}\overset{\text{def}}{=}ie_{3} \\
\mathcal{J}_{V}\left( e_{4},0\right) & =\left( 0,f_{4}\right) \Rightarrow
J_{V}e_{4}=f_{4}\overset{\text{def}}{=}ie_{4}
\end{align}%
\begin{align}
\mathcal{J}_{V}\left( 0,f_{1}\right) & =\left( -e_{1},0\right) \Rightarrow
J_{V}f_{1}=-e_{1}\overset{\text{def}}{=}if_{1}=iie_{1} \\
\mathcal{J}_{V}\left( 0,f_{2}\right) & =\left( -e_{2},0\right) \Rightarrow
J_{V}f_{2}=-e_{2}\overset{\text{def}}{=}if_{2}=iie_{2} \\
\mathcal{J}_{V}\left( 0,f_{3}\right) & =\left( -e_{3},0\right) \Rightarrow
J_{V}f_{3}=-e_{3}\overset{\text{def}}{=}if_{3}=iie_{3} \\
\mathcal{J}_{V}\left( 0,f_{4}\right) & =\left( -e_{4},0\right) \Rightarrow
J_{V}f_{4}=-e_{4}\overset{\text{def}}{=}if_{4}=iie_{4}
\end{align}%
\begin{equation}
V_{\mathbb{C}}=\left\{ \left. \left( v_{1},iv_{2}\right) \right\vert
v_{1},v_{2}\in V\right\} \cong\left\{ \left. \left( v_{1},v_{2}\right)
\right\vert v_{1},v_{2}\in V\right\} .
\end{equation}

\emph{External complexification }$\mathbb{C}\otimes V$\emph{\ can be
identified as internal complexification when }$W=V.$

$\mathcal{J}_{\mathcal{V}}$ is defined as above on $V\oplus W=V\oplus V,$
giving $\mathcal{V}_{\mathbb{C}}=V\oplus J_{V}V$ $\cong V_{\mathbb{C}%
}=V\oplus iV$ \emph{and on} 
\begin{gather}
\mathcal{J}_{\mathcal{V}}\left\vert \left( \mathbf{e,0}\right) \mathbf{,}%
\left( \mathbf{0},\mathbf{f}\right) \right) =\left\vert \left( \mathbf{e,0}%
\right) \mathbf{,}\left( \mathbf{0},\mathbf{f}\right) \right) \mathcal{R}%
_{e}\left( \mathcal{J}_{\mathcal{V}}\right) =\left\vert \left( \mathbf{e,0}%
\right) \mathbf{,}\left( \mathbf{0},\mathbf{f}\right) \right) \left( 
\begin{array}{cc}
0 & -\mathbb{I}_{r} \\ 
\mathbb{I}_{r} & 0%
\end{array}
\right) =\left\vert \left( \mathbf{0},\mathbf{f}\right) ,\left( \mathbf{e,0}%
\right) \right) \\
\left( 0,ie_{1}\right) \overset{def}{=}\mathcal{J}_{\mathcal{V}}\left(
e_{1},0\right) \Rightarrow\left( -e_{1},0\right) =\mathcal{J}_{\mathcal{V}%
}\left( 0,ie_{1}\right) \\
\mathcal{J}_{\mathcal{V}}\left\vert \left( \mathbf{e,0}\right) \mathbf{,}%
\left( \mathbf{0},i\mathbf{e}\right) \right) =\left\vert \left( \mathbf{e,0}%
\right) \mathbf{,}\left( \mathbf{0},i\mathbf{e}\right) \right) \mathcal{R}%
_{e}\left( \mathcal{J}_{\mathcal{V}}\right) =\left\vert \left( \mathbf{e,0}%
\right) \mathbf{,}\left( \mathbf{0},i\mathbf{e}\right) \right) \left( 
\begin{array}{cc}
0 & -\mathbb{I}_{r} \\ 
\mathbb{I}_{r} & 0%
\end{array}
\right) =\left\vert -\left( \mathbf{0},i\mathbf{e}\right) ,\left( \mathbf{e,0%
}\right) \right) \\
J_{V}\left\vert \mathbf{e,}i\mathbf{e}\right) =\left\vert \mathbf{e,}i%
\mathbf{e}\right) \mathcal{R}_{e}\left( J_{V}\right) =\left\vert \mathbf{e,}i%
\mathbf{e}\right) \left( 
\begin{array}{cc}
0 & -\mathbb{I}_{r} \\ 
\mathbb{I}_{r} & 0%
\end{array}
\right) =\left\vert i\mathbf{e,}-\mathbf{e}\right) .
\end{gather}

\subsubsection{Group Properties}

A complex structure is an orthogonal transformation $\in O\left( V\mathbf{%
\oplus}W\right) $ as 
\begin{align}
J_{V}^{t}J_{V} & =J_{V}J_{V}^{t}=I_{V\oplus W} \\
\left( 
\begin{array}{cc}
0 & \mathbb{I}_{r} \\ 
-\mathbb{I}_{r} & 0%
\end{array}
\right) \left( 
\begin{array}{cc}
0 & -\mathbb{I}_{r} \\ 
\mathbb{I}_{r} & 0%
\end{array}
\right) & =\left( 
\begin{array}{cc}
0 & -\mathbb{I}_{r} \\ 
\mathbb{I}_{r} & 0%
\end{array}
\right) \left( 
\begin{array}{cc}
0 & \mathbb{I}_{r} \\ 
-\mathbb{I}_{r} & 0%
\end{array}
\right) =\left( 
\begin{array}{cc}
\mathbb{I}_{r} & 0 \\ 
0 & \mathbb{I}_{r}%
\end{array}
\right) .
\end{align}
As all conb's of $V$ are related by orthogonal transformations belonging to $%
O\left( V\right) $ and all conb's of $W$ are related by orthogonal
transformations belonging to $O\left( W\right) $ it is easy to see that the
family of orthogonal complex structures on $V$ and $W$ are given by 
\begin{align}
OJ_{V}O^{t};O & \in O\left( V\right) \subset O\left( V\mathbf{\oplus }%
W\right) \\
OJ_{V^{c}}O^{t};O & \in O\left( V^{c}\right) \subset O\left( V\mathbf{\oplus}%
W\right) ,
\end{align}
where 
\begin{equation}
O=\left( 
\begin{array}{cc}
O_{V} & 0 \\ 
0 & O_{W}%
\end{array}
\right)
\end{equation}
and 
\begin{equation}
O\left( V\right) \cong O\left( W\right) .
\end{equation}

On the other hand one can start with a complex vector space, $V_{\mathbb{C}%
}, $ and consider real subspaces $V_{e},V_{e}^{c}$ generated by a conb $%
\left\{ \mathbf{e}\right\} $ and define a complexification map 
\begin{align}
J_{V_{e}} & :V_{e}\rightarrow V_{e}^{c} \\
J_{V_{e}} & :V_{e}^{c}\rightarrow V_{e} \\
J_{V_{e}} & :V_{e}+V_{e}^{c}\rightarrow V_{e}+V_{e}^{c}
\end{align}
by 
\begin{equation}
J_{V_{e}}e_{j}=ie_{j};1\leq j\leq r.
\end{equation}

$\lceil$%
\begin{align}
\mathcal{J}_{V_{e}} & :V\oplus J_{V_{e}}V\rightarrow V\oplus J_{V_{e}}V \\
J_{V_{e}}\left( c+J_{V_{e}}d\right) & =J_{V_{e}}c-d \\
& \equiv\left( 
\begin{array}{cc}
0 & -\mathbb{I}_{2r} \\ 
\mathbb{I}_{2r} & 0%
\end{array}
\right) \left( 
\begin{array}{c}
c \\ 
d%
\end{array}
\right) =\left( 
\begin{array}{c}
-d \\ 
c%
\end{array}
\right)
\end{align}
$\rfloor$

$J_{V_{e}}$ is the complex unit that complexifies the underlying real vector
space $V_{\sigma_{e}},$ where 
\begin{equation}
V_{\sigma_{e}}=\left\{ \left. \Sigma_{er}\left( x\right) =\frac{1}{2}\left\{
x+\sigma_{e}\left( x\right) \right\} \right\vert x\in V_{\mathbb{C}}\right\}
\end{equation}
and 
\begin{equation}
J_{V_{e}}V_{\sigma_{e}}=\left\{ \left. \Sigma_{ei}\left( x\right) =\frac{1}{2%
}\left\{ x-\sigma_{e}\left( x\right) \right\} \right\vert x\in V_{\mathbb{C}%
}\right\} .
\end{equation}
$\lceil$ Real and imaginary parts are defined by $\sigma$

\begin{equation*}
x+\sigma_{e}\left( x\right) =\left( 
\begin{array}{c}
c \\ 
d%
\end{array}
\right) +\left( 
\begin{array}{c}
c \\ 
-d%
\end{array}
\right) =2\left( 
\begin{array}{c}
c \\ 
0%
\end{array}
\right) =2\func{Re}_{\sigma_{e}}x
\end{equation*}%
\begin{equation*}
x-\sigma_{e}\left( x\right) =\left( 
\begin{array}{c}
c \\ 
d%
\end{array}
\right) -\left( 
\begin{array}{c}
c \\ 
-d%
\end{array}
\right) =2\left( 
\begin{array}{c}
0 \\ 
d%
\end{array}
\right) =2i\func{Im}_{\sigma_{e}}x,
\end{equation*}
showing that the representation of $\mathcal{R}_{e}\left( \sigma\right) $
wrt a basis, $\left\vert \mathbf{e}\right\rangle ,$ of $V_{\sigma}$ is given
by 
\begin{equation}
\mathcal{R}_{e}\left( \sigma_{e}\right) \mathcal{R}_{e}\left( x\right) =%
\mathcal{R}_{e}\left( \sigma_{e}\right) \left( \left( 
\begin{array}{c}
c \\ 
d%
\end{array}
\right) \right) =\left( 
\begin{array}{c}
c \\ 
-d%
\end{array}
\right) ,
\end{equation}
leading to%
\begin{equation}
\sigma_{e}\left( \left\vert \mathbf{e}\right\rangle ,\left\vert i\mathbf{e}%
\right\rangle \right) =\left( \left\vert \mathbf{e}\right\rangle ,\left\vert
i\mathbf{e}\right\rangle \right) \mathcal{R}_{e}\left( \sigma_{e}\right)
=\left( \left\vert \mathbf{e}\right\rangle ,\left\vert i\mathbf{e}%
\right\rangle \right) \left( 
\begin{array}{cc}
\mathbb{I}_{2r} & \mathbb{O} \\ 
\mathbb{O} & -\mathbb{I}_{2r}%
\end{array}
\right) .
\end{equation}
Each $\sigma_{e}$ corresponds to the class of basis generated by 
\begin{equation}
\left\{ \left. O\left\vert \mathbf{e}\right\rangle ,O^{c}\left\vert i\mathbf{%
e}\right\rangle \right\vert O\in O\left( V\right) ,O^{c}\in O\left(
V^{c}\right) \right\}
\end{equation}
that has canonical representative $\left\vert \mathbf{e}\right\rangle .$
Thus we can write $V_{\sigma_{e}}\equiv V_{e}.$ The representation of $%
J_{V_{e}}$ wrt a basis of $V_{e}$ is given by%
\begin{equation}
J_{\mathbb{C}}\left( \left\vert \mathbf{e}\right\rangle ,\left\vert i\mathbf{%
e}\right\rangle \right) =\left( \left\vert \mathbf{e}\right\rangle
,\left\vert i\mathbf{e}\right\rangle \right) R_{e}\left( J_{\mathbb{C}%
}\right) =\left( \left\vert \mathbf{e}\right\rangle ,\left\vert i\mathbf{e}%
\right\rangle \right) \left( 
\begin{array}{cc}
\mathbb{O} & -\mathbb{I}_{2r} \\ 
\mathbb{I}_{2r} & \mathbb{O}%
\end{array}
\right) .
\end{equation}

$\lceil$Consistency in show by 
\begin{align}
\sigma_{e}\left( J_{V_{e}}x\right) & =\left( 
\begin{array}{cc}
\mathbb{I}_{2r} & 0 \\ 
0 & -\mathbb{I}_{2r}%
\end{array}
\right) \left( 
\begin{array}{cc}
0 & -\mathbb{I}_{2r} \\ 
\mathbb{I}_{2r} & 0%
\end{array}
\right) \left( 
\begin{array}{c}
c \\ 
d%
\end{array}
\right) =\left( 
\begin{array}{cc}
\mathbb{I}_{2r} & 0 \\ 
0 & -\mathbb{I}_{2r}%
\end{array}
\right) \left( 
\begin{array}{c}
-d \\ 
c%
\end{array}
\right) =\left( 
\begin{array}{c}
-d \\ 
-c%
\end{array}
\right) \\
& \equiv\sigma_{e}\left( ix\right) =\overline{ix}=-i\overline{x}=-i\left( 
\begin{array}{c}
c \\ 
-d%
\end{array}
\right) =i\left( 
\begin{array}{c}
-c \\ 
d%
\end{array}
\right) =\left( 
\begin{array}{c}
-d \\ 
-c%
\end{array}
\right) ,
\end{align}%
\begin{equation}
\sigma_{e}J_{V_{e}}=-J_{V_{e}}\sigma_{e}
\end{equation}%
\begin{align}
\left( 
\begin{array}{cc}
\mathbb{I}_{2r} & 0 \\ 
0 & -\mathbb{I}_{2r}%
\end{array}
\right) \left( 
\begin{array}{cc}
0 & -\mathbb{I}_{2r} \\ 
\mathbb{I}_{2r} & 0%
\end{array}
\right) & =\left( 
\begin{array}{cc}
0 & -\mathbb{I}_{2r} \\ 
-\mathbb{I}_{2r} & 0%
\end{array}
\right) \\
\left( 
\begin{array}{cc}
0 & -\mathbb{I}_{2r} \\ 
\mathbb{I}_{2r} & 0%
\end{array}
\right) \left( 
\begin{array}{cc}
\mathbb{I}_{2r} & 0 \\ 
0 & -\mathbb{I}_{2r}%
\end{array}
\right) & =\left( 
\begin{array}{cc}
0 & \mathbb{I}_{2r} \\ 
\mathbb{I}_{2r} & 0%
\end{array}
\right)
\end{align}
which is expressed abstractly as 
\begin{equation}
\sigma_{e}\left( a+J_{V_{e}}b)\right) =\sigma_{e}\left( a\right)
+\sigma_{e}\left( J_{V_{e}}b\right) =\sigma_{e}\left( a\right)
-J_{V_{e}}\sigma_{e}\left( b\right) .
\end{equation}
$\rfloor$

Complex structures and complex conjugations are indexed by $\mathfrak{C}%
\left( \mathfrak{H}\right) $. $\sigma_{e}$ and $J_{V_{e}}$ are in $1-1$
correspondence and there \emph{are many} different pairs of complex
structures and associated complex conjugations.

$J_{V_{e}}$ defines a Hermitian scalar product on $V_{e}$ by 
\begin{equation}
\left\langle \left. \left( a,b\right) \right\vert \left( c,d\right)
\right\rangle _{J}=\left( \left. \left( a,b\right) \right\vert \left(
c,d\right) \right) _{V_{e}}+i\left( \left. \left( a,b\right) \right\vert J_{%
\mathbb{C}}\left( c,d\right) \right) _{V_{e}};a,b,c,d\in V_{e}
\end{equation}
$\lceil$%
\begin{align}
& \left\langle \left. \left( a,b\right) \right\vert \left( c,d\right)
\right\rangle _{J}=\left( a,b\right) \left( 
\begin{array}{cc}
\mathbb{I}_{2s} & 0 \\ 
0 & \mathbb{I}_{2s}%
\end{array}
\right) \left( 
\begin{array}{c}
c \\ 
d%
\end{array}
\right) +i\left( a,b\right) \left( 
\begin{array}{cc}
0 & \mathbb{I}_{2s} \\ 
-\mathbb{I}_{2s} & 0%
\end{array}
\right) \left( 
\begin{array}{c}
c \\ 
d%
\end{array}
\right) \\
R_{e}\left( J_{\mathbb{C}}\right) R_{e}\left( J_{\mathbb{C}}^{t}\right) &
=\left( 
\begin{array}{cc}
0 & \mathbb{I}_{2s} \\ 
-\mathbb{I}_{2s} & 0%
\end{array}
\right) \left( 
\begin{array}{cc}
0 & -\mathbb{I}_{2s} \\ 
\mathbb{I}_{2s} & 0%
\end{array}
\right) =\left( 
\begin{array}{cc}
\mathbb{I}_{2s} & 0 \\ 
0 & \mathbb{I}_{2s}%
\end{array}
\right) =\left( 
\begin{array}{cc}
0 & -\mathbb{I}_{2s} \\ 
\mathbb{I}_{2s} & 0%
\end{array}
\right) \left( 
\begin{array}{cc}
0 & \mathbb{I}_{2s} \\ 
-\mathbb{I}_{2s} & 0%
\end{array}
\right) =R_{e}\left( J_{\mathbb{C}}^{t}\right) R_{e}\left( J_{\mathbb{C}%
}\right)
\end{align}
$\rfloor$

to give a complex Hilbert space $V_{e\mathbb{C}}.$

In fact 
\begin{gather}
J=-J^{t} \\
J^{2}=-I=-JJ^{t} \\
\Downarrow \\
JJ^{t}=J^{t}J=I,
\end{gather}%
thus $J\in O\left( V\right) $ if an inner product exists that defines the
transpose map.

\section{Summary}

In general we start with a complex Hilbert space $\mathfrak{H}$ and need to
find real spaces $V_{\sigma}$ such that 
\begin{equation}
\mathfrak{H}=V_{\sigma_{e}\mathbb{C}}.
\end{equation}
In order to define such a real subspace $V_{\sigma_{e}}$ one must first
define a conjugation $\sigma_{e}$ on $\mathfrak{H}$%
\begin{align}
\mathcal{\sigma}_{e} & :\mathfrak{H}\rightarrow\mathfrak{H} \\
\mathcal{\sigma}\left( u+\sigma_{e}\left( v\right) \right) &
=v+\sigma_{e}\left( u\right) .
\end{align}

$\lceil$a \emph{conjugation} 
\begin{equation}
\mathcal{\sigma}_{e}:H\rightarrow\mathcal{\sigma}_{e}\left( H\right)
\end{equation}
is an antiunitary map with 
\begin{align}
\mathcal{\sigma}_{e}^{2} & =I_{H} \\
\left\langle \left. \mathcal{\sigma}_{e}\left( x\right) \right\vert \mathcal{%
\sigma}_{e}\left( y\right) \right\rangle _{H} & =\overline {\left\langle
\left. x\right\vert y\right\rangle }_{H}.
\end{align}
$\rfloor$

Note that 
\begin{equation}
\mathfrak{H}=\mathfrak{H}+\mathcal{\sigma}_{e}\left( \mathfrak{H}\right) ,
\end{equation}
not the direct sum 
\begin{equation}
\mathfrak{H}\oplus\mathcal{\sigma}_{e}\left( \mathfrak{H}\right)
\end{equation}
as $\mathfrak{H\,}\mathcal{\ }$and $\mathcal{\sigma}_{e}\left( \mathfrak{H}%
\right) $ are not disjoint in fact 
\begin{equation}
\mathfrak{H}\mathcal{\cap\sigma}_{e}\left( \mathfrak{H}\right) =\mathfrak{H.}
\end{equation}

One can then define a real subspace $V_{\sigma_{e}}$ of $\mathfrak{H}$ by 
\begin{equation}
V_{\sigma_{e}}=\left\{ \left. x\in\mathfrak{H}\right\vert \mathcal{\sigma }%
_{e}\left( x\right) =x\right\}
\end{equation}
and then form a complexification of $V_{\sigma_{e}}$ wrt a complex structure 
$J_{V_{\sigma_{e}}}$

$\lceil$A \emph{complex structure} on a real vector space $V$ is a real
linear map 
\begin{align}
J & :V\rightarrow W \\
J^{2} & =-I_{V}.
\end{align}
$\rfloor$

that acts on $V_{\sigma_{e}}$ giving a complexification 
\begin{equation}
V_{\sigma_{e}J}=V_{\sigma_{e}}\oplus J_{V_{\sigma_{e}}}V_{\sigma_{e}}\equiv%
\mathfrak{H.}
\end{equation}
The complex structure $J_{V_{\sigma_{e}}}$ induces a real map $J_{V_{\sigma
_{e}}}$ on $\mathfrak{H}$ given by 
\begin{equation}
J_{V_{\sigma_{e}}}:V_{\sigma_{e}}\oplus
J_{V_{\sigma_{e}}}V_{\sigma_{e}}\rightarrow
J_{V_{\sigma_{e}}}V_{\sigma_{e}}\ominus V_{\sigma_{e}}.
\end{equation}
The conjugation $\sigma$

$\lceil$%
\begin{align}
\mathfrak{H} & =V_{\sigma e_{{}}}\mathcal{\oplus}J_{V_{\sigma_{e}}}V_{%
\sigma_{e}} \\
\mathcal{\sigma}_{e}\left( \mathfrak{H}\right) & =V_{\sigma_{e}}\mathcal{%
\ominus}J_{V_{\sigma_{e}}}V_{\sigma_{e}}
\end{align}

$\rfloor$

and complex structure $J$ are compatible if 
\begin{equation}
\mathcal{\sigma}_{e}J_{V_{\sigma}}=-J_{V_{\sigma}}\mathcal{\sigma}_{e}.
\end{equation}
If there exists a basis $\left\vert \mathbf{e}\right\rangle $ of $\mathfrak{H%
}$ such that 
\begin{align}
R_{e}\left( \mathcal{\sigma}_{e}\right) & =\left( 
\begin{array}{cc}
\mathbb{I}_{r} & 0 \\ 
0 & -\mathbb{I}_{r}%
\end{array}
\right) \\
R_{e}\left( J_{V_{\sigma_{e}}}\right) & =\left( 
\begin{array}{cc}
0 & -\mathbb{I}_{r} \\ 
\mathbb{I}_{r} & 0%
\end{array}
\right)
\end{align}
then $\sigma_{e}$ and $J_{V_{\sigma_{e}}}$ are compatible as%
\begin{align}
\left( 
\begin{array}{cc}
\mathbb{I}_{r} & 0 \\ 
0 & -\mathbb{I}_{r}%
\end{array}
\right) \left( 
\begin{array}{cc}
0 & -\mathbb{I}_{r} \\ 
\mathbb{I}_{r} & 0%
\end{array}
\right) & =\left( 
\begin{array}{cc}
0 & -\mathbb{I}_{r} \\ 
-\mathbb{I}_{r} & 0%
\end{array}
\right) \\
\left( 
\begin{array}{cc}
0 & -\mathbb{I}_{r} \\ 
\mathbb{I}_{r} & 0%
\end{array}
\right) \left( 
\begin{array}{cc}
\mathbb{I}_{r} & 0 \\ 
0 & -\mathbb{I}_{r}%
\end{array}
\right) & =\left( 
\begin{array}{cc}
0 & \mathbb{I}_{r} \\ 
\mathbb{I}_{r} & 0%
\end{array}
\right) .
\end{align}

\subsubsection{Generalization}

One can easily generalize this by considering a complexification 
\begin{equation}
\left\vert e_{r+j}\right) =J_{e}\left\vert e_{j}\right) ;1\leq j\leq r,
\end{equation}
where

\begin{equation}
R_{e}\left( J_{e}\right) =\left( 
\begin{array}{cc}
0 & -\mathbb{I}_{r} \\ 
\mathbb{I}_{r} & 0%
\end{array}
\right) ,
\end{equation}
to obtain 
\begin{equation}
\left\vert a\right) =\left\vert z\right\rangle =\left\vert e_{1},\cdots
,e_{r}\right\rangle R_{e}\left( z\right) =\left\vert I_{r}e_{1},\cdots
,I_{r}e_{r},J_{e}e_{1},\cdots,J_{e}e_{r}\right) \left( 
\begin{array}{c}
a_{1} \\ 
\vdots \\ 
\vdots \\ 
a_{2r}%
\end{array}
\right) =\left\vert e_{1},\cdots,e_{2r}\right\rangle R_{e}\left( a\right) ,
\end{equation}
i.e. 
\begin{equation}
V_{r\mathbb{C}}=V_{er}\oplus J_{e}V_{er}.
\end{equation}
Note that 
\begin{equation}
V_{2r}=V_{er}\oplus W_{er}
\end{equation}
where 
\begin{align}
V_{er} & =\mathbb{R}-\limfunc{linspan}\left\{ e_{1},\cdots ,e_{r}\right\} \\
W_{er} & =\mathbb{R}-\limfunc{linspan}\left\{ e_{r+1},\cdots ,e_{2r}\right\}
.
\end{align}
The difference comes from 
\begin{align}
V_{2r} & =\mathbb{R}-\limfunc{linspan}\left\{ e_{1},\cdots
,e_{r},e_{r+1},\cdots,e_{2r}\right\} \\
V_{r\mathbb{C}} & =\mathbb{R}-\limfunc{linspan}\left\{
e_{1}+J_{e}e_{r+1},\cdots,e_{r}+J_{e}e_{2r}\right\} \equiv\mathbb{R}-%
\limfunc{linspan}\left\{ e_{1}+ie_{1},\cdots,e_{r}+ie_{r}\right\} .
\end{align}

One can define a conjugation $\sigma_{e}$ 
\begin{equation}
\sigma_{e}:V_{2r}\rightarrow V_{2r}
\end{equation}
that is compatible with $J_{e}$ by 
\begin{equation}
R_{e}\left( \sigma_{e}\right) =\left( 
\begin{array}{cc}
\mathbb{I}_{r} & 0 \\ 
0 & -\mathbb{I}_{r}%
\end{array}
\right) ,
\end{equation}
i.e. 
\begin{equation}
\sigma_{e}\left( \left\vert \mathbf{e}\right\rangle ,\left\vert \mathbf{%
\tilde{e}}\right\rangle \right) =\left( \left\vert \mathbf{e}\right\rangle
,\left\vert \mathbf{\tilde{e}}\right\rangle \right) R_{e}\left(
\sigma_{e}\right) =\left( \left\vert \mathbf{e}\right\rangle ,\left\vert 
\mathbf{\tilde{e}}\right\rangle \right) \left( 
\begin{array}{cc}
\mathbb{I}_{r} & 0 \\ 
0 & -\mathbb{I}_{r}%
\end{array}
\right) ,
\end{equation}
while the representation of $J_{e}$ is given by%
\begin{equation}
J_{e}\left( \left\vert \mathbf{e}\right\rangle ,\left\vert \mathbf{\tilde{e}}%
\right\rangle \right) =\left( \left\vert \mathbf{e}\right\rangle ,\left\vert 
\mathbf{\tilde{e}}\right\rangle \right) R_{e}\left( J_{e}\right) =\left(
\left\vert \mathbf{e}\right\rangle ,\left\vert \mathbf{\tilde{e}}%
\right\rangle \right) \left( 
\begin{array}{cc}
0 & -\mathbb{I}_{r} \\ 
\mathbb{I}_{r} & 0%
\end{array}
\right) .
\end{equation}

\section{Invariance Groups}

The invariance group of $J_{e}$ $\in Sp\left( 2r,\mathbb{R}\right) $ that is
restricted to be in $O\left( 2r,\mathbb{R}\right) $ i.e. $Sp\left( 2r,%
\mathbb{R}\right) \cap O\left( 2r,\mathbb{R}\right) =U\left( r\right) $ is
given by

$\lceil Sp\left( 2r,\mathbb{R}\right) $ is defined by 
\begin{gather}
\left( 
\begin{array}{cc}
\mathbb{A}^{t} & \mathbb{C}^{t} \\ 
\mathbb{B}^{t} & \mathbb{D}^{t}%
\end{array}
\right) \left( 
\begin{array}{cc}
0 & -\mathbb{I}_{r} \\ 
\mathbb{I}_{r} & 0%
\end{array}
\right) \left( 
\begin{array}{cc}
\mathbb{A} & \mathbb{B} \\ 
\mathbb{C} & \mathbb{D}%
\end{array}
\right) \\
=\left( 
\begin{array}{cc}
\mathbb{C}^{t} & -\mathbb{A}^{t} \\ 
\mathbb{D}^{t} & -\mathbb{B}^{t}%
\end{array}
\right) \left( 
\begin{array}{cc}
\mathbb{A} & \mathbb{B} \\ 
\mathbb{C} & \mathbb{D}%
\end{array}
\right) \\
=\left( 
\begin{array}{cc}
\mathbb{C}^{t}\mathbb{A-A}^{t}\mathbb{C} & \mathbb{C}^{t}\mathbb{B-A}^{t}%
\mathbb{D} \\ 
\mathbb{D}^{t}\mathbb{A-B}^{t}\mathbb{C} & \mathbb{D}^{t}\mathbb{B-B}^{t}%
\mathbb{D}%
\end{array}
\right) =\left( 
\begin{array}{cc}
\mathbb{O} & -\mathbb{I}_{r} \\ 
\mathbb{I}_{r} & \mathbb{O}%
\end{array}
\right) \\
\Updownarrow \\
\mathbb{C}^{t}\mathbb{A-A}^{t}\mathbb{C}=\mathbb{O} \\
\mathbb{C}^{t}\mathbb{B-A}^{t}\mathbb{D}=-\mathbb{I}_{r} \\
\mathbb{D}^{t}\mathbb{A-B}^{t}\mathbb{C}=\mathbb{I}_{r} \\
\mathbb{D}^{t}\mathbb{B-B}^{t}\mathbb{D}=\mathbb{O},
\end{gather}
while $O\left( 2r,\mathbb{R}\right) $ by 
\begin{gather}
\left( 
\begin{array}{cc}
\mathbb{A}^{t} & \mathbb{C}^{t} \\ 
\mathbb{B}^{t} & \mathbb{D}^{t}%
\end{array}
\right) \left( 
\begin{array}{cc}
\mathbb{A} & \mathbb{B} \\ 
\mathbb{C} & \mathbb{D}%
\end{array}
\right) =\left( 
\begin{array}{cc}
\mathbb{A}^{t}\mathbb{A+C}^{t}\mathbb{C} & \mathbb{A}^{t}\mathbb{B+C}^{t}%
\mathbb{D} \\ 
\mathbb{B}^{t}\mathbb{A+D}^{t}\mathbb{C} & \mathbb{B}^{t}\mathbb{B+D}^{t}%
\mathbb{D}%
\end{array}
\right) =\left( 
\begin{array}{cc}
\mathbb{I}_{r} & \mathbb{O} \\ 
\mathbb{O} & \mathbb{I}_{r}%
\end{array}
\right) \\
\Updownarrow \\
\mathbb{A}^{t}\mathbb{A+C}^{t}\mathbb{C=I}_{r} \\
\mathbb{A}^{t}\mathbb{B+C}^{t}\mathbb{D=O} \\
\mathbb{B}^{t}\mathbb{B+D}^{t}\mathbb{D=I}_{r}
\end{gather}
and $O\left( r,r,\mathbb{R}\right) $ by 
\begin{gather}
\left( 
\begin{array}{cc}
\mathbb{A}^{t} & \mathbb{C}^{t} \\ 
\mathbb{B}^{t} & \mathbb{D}^{t}%
\end{array}
\right) \left( 
\begin{array}{cc}
\mathbb{I}_{r} & \mathbb{O} \\ 
\mathbb{O} & -\mathbb{I}_{r}%
\end{array}
\right) \left( 
\begin{array}{cc}
\mathbb{A} & \mathbb{B} \\ 
\mathbb{C} & \mathbb{D}%
\end{array}
\right) =\left( 
\begin{array}{cc}
\mathbb{A}^{t} & -\mathbb{C}^{t} \\ 
\mathbb{B}^{t} & -\mathbb{D}^{t}%
\end{array}
\right) \left( 
\begin{array}{cc}
\mathbb{A} & \mathbb{B} \\ 
\mathbb{C} & \mathbb{D}%
\end{array}
\right) \\
=\left( 
\begin{array}{cc}
\mathbb{A}^{t}\mathbb{A-C}^{t}\mathbb{C} & \mathbb{A}^{t}\mathbb{B-C}^{t}%
\mathbb{D} \\ 
\mathbb{B}^{t}\mathbb{A-D}^{t}\mathbb{C} & \mathbb{B}^{t}\mathbb{B-D}^{t}%
\mathbb{D}%
\end{array}
\right) =\left( 
\begin{array}{cc}
\mathbb{I}_{r} & \mathbb{O} \\ 
\mathbb{O} & -\mathbb{I}_{r}%
\end{array}
\right) \\
\Updownarrow \\
\mathbb{A}^{t}\mathbb{A-C}^{t}\mathbb{C=I}_{r} \\
\mathbb{A}^{t}\mathbb{B-C}^{t}\mathbb{D=O} \\
\mathbb{B}^{t}\mathbb{B-D}^{t}\mathbb{D=-I}_{r}.
\end{gather}
$\rfloor$%
\begin{align}
J_{e}\left( \left\vert \mathbf{e}\right\rangle ,\left\vert \mathbf{\tilde{e}}%
\right\rangle \right) & =\left( \left\vert \mathbf{e}\right\rangle
,\left\vert \mathbf{\tilde{e}}\right\rangle \right) \left( 
\begin{array}{cc}
\mathbb{T}_{1} & 0 \\ 
0 & \mathbb{T}_{2}%
\end{array}
\right) ^{t}\left( 
\begin{array}{cc}
0 & -\mathbb{I}_{r} \\ 
\mathbb{I}_{r} & 0%
\end{array}
\right) \left( 
\begin{array}{cc}
\mathbb{T}_{1} & 0 \\ 
0 & \mathbb{T}_{2}%
\end{array}
\right) \\
& =\left( \left\vert \mathbf{e}\right\rangle ,\left\vert \mathbf{\tilde{e}}%
\right\rangle \right) \left( 
\begin{array}{cc}
\mathbb{T}_{1} & 0 \\ 
0 & \mathbb{T}_{2}%
\end{array}
\right) ^{t}\left( 
\begin{array}{cc}
0 & -\mathbb{T}_{2} \\ 
\mathbb{T}_{1} & 0%
\end{array}
\right) \\
& =\left( \left\vert \mathbf{e}\right\rangle ,\left\vert \mathbf{\tilde{e}}%
\right\rangle \right) \left( 
\begin{array}{cc}
0 & -\mathbb{T}_{1}^{t}\mathbb{T}_{2} \\ 
\mathbb{T}_{2}^{t}\mathbb{T}_{1} & 0%
\end{array}
\right) =\left( \left\vert \mathbf{e}\right\rangle ,\left\vert \mathbf{%
\tilde{e}}\right\rangle \right) \left( 
\begin{array}{cc}
0 & -\mathbb{I}_{r} \\ 
\mathbb{I}_{r} & 0%
\end{array}
\right) \\
& \Rightarrow\text{ }\mathbb{T}_{1}=\mathbb{T}_{2}.
\end{align}
The invariance group of $\sigma_{e}$ that is restricted to be in $O\left( 2r,%
\mathbb{R}\right) $ i.e. $O\left( r,r,\mathbb{R}\right) \cap O\left( 2r,%
\mathbb{R}\right) $ is given by 
\begin{equation}
\sigma_{e}\left( \left\vert \mathbf{e}\right\rangle ,\left\vert \mathbf{%
\tilde{e}}\right\rangle \right) =\left( \left\vert \mathbf{e}\right\rangle
,\left\vert \mathbf{\tilde{e}}\right\rangle \right) \left( 
\begin{array}{cc}
\mathbb{T} & 0 \\ 
0 & \mathbb{T}%
\end{array}
\right) ^{t}\left( 
\begin{array}{cc}
\mathbb{I}_{r} & 0 \\ 
0 & -\mathbb{I}_{r}%
\end{array}
\right) \left( 
\begin{array}{cc}
\mathbb{T} & 0 \\ 
0 & \mathbb{T}%
\end{array}
\right) =\left( \left\vert \mathbf{e}\right\rangle ,\left\vert \mathbf{%
\tilde{e}}\right\rangle \right) \left( 
\begin{array}{cc}
\mathbb{I}_{r} & 0 \\ 
0 & -\mathbb{I}_{r}%
\end{array}
\right) ;\mathbb{T}\in O\left( r,\mathbb{R}\right) .
\end{equation}
Hence%
\begin{gather}
\mathbb{T}\in O\left( r,\mathbb{R}\right) ;\left( 
\begin{array}{cc}
\mathbb{T} & 0 \\ 
0 & \mathbb{T}%
\end{array}
\right) \in O\left( 2r,\mathbb{R}\right) ;\left( 
\begin{array}{cc}
\mathbb{T} & 0 \\ 
0 & \mathbb{T}%
\end{array}
\right) \in Sp\left( 2r,\mathbb{R}\right) ;\left( 
\begin{array}{cc}
\mathbb{T} & 0 \\ 
0 & \mathbb{T}%
\end{array}
\right) \in O\left( r,r,\mathbb{R}\right) \\
Sp\left( 2r,\mathbb{R}\right) \cap O\left( r,r,\mathbb{R}\right) =O\left( r,%
\mathbb{R}\right) .
\end{gather}
One can define a Hermitian inner product in $V_{\mathbb{C}r},$ by 
\begin{equation}
\left\langle \left. z_{1}\right\vert z_{2}\right\rangle =\left( \left.
a_{1}\right\vert a_{2}\right) +i\left( \left. a_{1}\right\vert
J_{e}a_{2}\right) ,
\end{equation}
where 
\begin{equation}
z_{j}=a_{j1}+ia_{j2}.
\end{equation}
The invariance group, $U\left( r\right) $ of $\left\langle \left.
z_{1}\right\vert z_{2}\right\rangle $ is given by 
\begin{align}
\left\langle \left. Tz_{1}\right\vert Tz_{2}\right\rangle & =\left( \left.
Ta_{1}\right\vert Ta_{2}\right) +i\left( \left. Ta_{1}\right\vert
J_{e}Ta_{2}\right) \\
& =\left( \left. a_{1}\right\vert T^{t}Ta_{2}\right) +i\left( \left.
a_{1}\right\vert T^{t}J_{e}Ta_{2}\right) =\left( \left. a_{1}\right\vert
a_{2}\right) +i\left( \left. a_{1}\right\vert J_{e}a_{2}\right)
\end{align}
i.e. 
\begin{equation}
U\left( r\right) =O\left( 2r,\mathbb{R}\right) \cap Sp\left( 2r,\mathbb{R}%
\right) .
\end{equation}

Complexifications belong to the orthogonal group 
\begin{equation}
J^{t}J=JJ^{t}=I,
\end{equation}
(which they must be if $J=-J^{t}$ and $J^{2}=-I$), so if $\left\vert \mathbf{%
e}\right) $ is a conb basis and the representations are as in the preceding 
\begin{equation}
\left( 
\begin{array}{cc}
0 & -\mathbb{I}_{r} \\ 
\mathbb{I}_{r} & 0%
\end{array}
\right) \left( 
\begin{array}{cc}
0 & \mathbb{I}_{r} \\ 
-\mathbb{I}_{r} & 0%
\end{array}
\right) =\left( 
\begin{array}{cc}
\mathbb{I}_{r} & 0 \\ 
0 & \mathbb{I}_{r}%
\end{array}
\right)
\end{equation}
the complexification is an orthogonal opertor i.e. $J\in O\left(
V_{2r}\right) ,$ more over as 
\begin{equation}
\left( 
\begin{array}{cc}
\mathbb{I}_{r} & 0 \\ 
0 & -\mathbb{I}_{r}%
\end{array}
\right) ^{t}\left( 
\begin{array}{cc}
\mathbb{I}_{r} & 0 \\ 
0 & -\mathbb{I}_{r}%
\end{array}
\right) =\left( 
\begin{array}{cc}
\mathbb{I}_{r} & 0 \\ 
0 & \mathbb{I}_{r}%
\end{array}
\right)
\end{equation}
the complex conjugation $\sigma_{e}\in O\left( V_{2r}\right) .$

\subsection{Geminals}

In order to define a complex structure on a real space one just needs a real
antisymmetric form 
\begin{align}
g& :V_{2r}\times V_{2r}\rightarrow \mathbb{R} \\
g& \leftrightarrow \left( g\right\vert \in \left(
\bigwedge\nolimits^{2}V_{2r}\right) ^{d} \\
g\left( v_{1},v_{2}\right) & =\left( g\left\vert v_{1}\wedge v_{2}\right.
\right)
\end{align}%
that can be brought to cannonical form 
\begin{equation}
g=\sum\limits_{1\leq j<k\leq 2r}g_{jk}\left\vert e_{j}e_{k}\right)
=\sum\limits_{1\leq j\leq r}c_{gj}\left\vert \widetilde{e}_{gj}\widetilde{e}%
_{gj+r}\right) =\sum\limits_{1\leq j\leq r}\left\vert e_{gj}e_{gj+r}\right)
;\,c_{j}>0,\,\left( e_{gj}\left\vert e_{gk}\right. \right) =c_{j}\delta
_{jk},
\end{equation}%
where 
\begin{eqnarray}
\left\vert e_{gj}\right) &=&c_{gj}^{\frac{1}{2}}\left\vert \widetilde{e}%
_{gj}\right) \\
\left\vert e_{gj+r}\right) &=&c_{gj}^{\frac{1}{2}}\left\vert \widetilde{e}%
_{gj+r}\right) ,
\end{eqnarray}%
which defines a subspace decomposition (if $g$ no zero canonical
coefficients) 
\begin{equation}
V_{2r}=V_{g}\oplus V_{g}^{\#}
\end{equation}%
and a basis 
\begin{equation}
\left\vert \mathbf{e}\right) =\left\vert \mathbf{e}_{g},\mathbf{e}%
_{g}^{\#}\right)
\end{equation}%
and thus a complexification 
\begin{equation}
J_{g}=\tsum\limits_{1\leq j\leq r}\left\{ \left\vert e_{gj+r}\right) \left(
e_{gj}\right\vert -\left\vert e_{gj}\right) \left( e_{gj+r}\right\vert
\right\} .
\end{equation}

\begin{remark}
Any \emph{real Hilbert space} of bounded operator belonging to the real
operator space $\mathcal{B}\left( \mathcal{H}\right) $ can be expressed as
the sum of symmetric and antisymmetric operators 
\begin{equation}
A=A_{s}+A_{as},
\end{equation}%
where 
\begin{eqnarray}
A_{s} &=&\frac{1}{2}\left( A+A^{t}\right) \\
A_{as} &=&\frac{1}{2}\left( A-A^{t}\right) .
\end{eqnarray}%
Moreover 
\begin{equation}
\mathcal{B}\left( \mathcal{H}\right) =\mathcal{B}\left( \mathcal{H}\right)
_{s}\oplus \mathcal{B}\left( \mathcal{H}\right) _{as}
\end{equation}%
as the direct sum real of subspaces. The subspaces $\mathcal{B}\left( 
\mathcal{H}\right) _{s}$ and $\mathcal{B}\left( \mathcal{H}\right) _{as}$
are respectively closed under the symmetic product 
\begin{eqnarray}
\left\{ ,\right\} &:&\mathcal{B}\left( \mathcal{H}\right) _{s}\times 
\mathcal{B}\left( \mathcal{H}\right) _{s}\rightarrow \mathcal{B}\left( 
\mathcal{H}\right) _{s} \\
\left\{ A,B\right\} &=&AB+BA
\end{eqnarray}%
(and thus a Jordan algebra) and the antisymmetric product 
\begin{eqnarray}
\left[ ,\right] &:&\mathcal{B}\left( \mathcal{H}\right) _{as}\times \mathcal{%
B}\left( \mathcal{H}\right) _{as}\rightarrow \mathcal{B}\left( \mathcal{H}%
\right) _{as} \\
\left[ A,B\right] &=&AB-BA
\end{eqnarray}%
(and thus a Lie algebra).
\end{remark}

\begin{theorem}
As real vector spaces 
\begin{equation}
\mathcal{B}\left( \mathcal{H}\right) _{as}\cong \tbigwedge\nolimits^{2}%
\mathcal{H}
\end{equation}%
and 
\begin{equation}
\mathcal{B}\left( \mathcal{H}\right) _{s}\cong \tbigvee\nolimits^{2}\mathcal{%
H}.
\end{equation}%
$\lceil $%
\begin{eqnarray}
v\wedge w &=&\frac{1}{2}\left( v\otimes w-w\otimes v\right) \\
v\vee w &=&\frac{1}{2}\left( v\otimes w+w\otimes v\right) .
\end{eqnarray}%
$\rfloor $
\end{theorem}

\begin{proof}
One can construct the isomorphisms by 
\begin{eqnarray}
\mathcal{A} &:&\mathcal{B}\left( \mathcal{H}\right) _{as}\rightarrow
\tbigwedge\nolimits^{2}\mathcal{H} \\
\mathcal{A}\left( \left\vert e_{j}\right) \left( e_{k}\right\vert
-\left\vert e_{k}\right) \left( e_{j}\right\vert \right) &=&e_{j}\wedge
e_{k};1\leq j<k\leq 2r
\end{eqnarray}%
and 
\begin{eqnarray}
\mathcal{S} &:&\mathcal{B}\left( \mathcal{H}\right) _{s}\rightarrow
\tbigvee\nolimits^{2}\mathcal{H} \\
\mathcal{S}\left( \left\vert e_{j}\right) \left( e_{k}\right\vert
+\left\vert e_{k}\right) \left( e_{j}\right\vert \right) &=&e_{j}\vee
e_{k};1\leq j\leq k\leq 2r.
\end{eqnarray}
\end{proof}

$\lceil $%
\begin{equation}
\mathcal{B}\left( \mathcal{H}\right) \cong \mathcal{H\otimes H}^{d}
\end{equation}%
$\rfloor $

The canonical forms of $G\in \mathcal{B}\left( \mathcal{H}\right) _{as}$ and
the corresponding $g=\mathcal{A}\left( G\right) \in \tbigwedge\nolimits^{2}%
\mathcal{H}$ are 
\begin{equation}
G=TJ_{e}T^{t}=T\left\{ \tsum\limits_{1\leq j<k\leq 2r}\left( \left\vert
e_{j}\right) \left( e_{k}\right\vert -\left\vert e_{k}\right) \left(
e_{j}\right\vert \right) \right\} T^{t}
\end{equation}
and 
\begin{equation}
g=\tsum\limits_{1\leq j\leq r}\left\vert Te_{j}\wedge Te_{j+r}\right) ,
\end{equation}%
where $T\in GL\left( r,\mathbb{R}\right) .$

The canonical forms of $A\in \mathcal{B}\left( \mathcal{H}\right) _{s}$ and
the corresponding $s=\mathcal{S}\left( A\right) \in \tbigvee\nolimits^{2}%
\mathcal{H}$ are 
\begin{equation}
A=TI_{p,r-q}T^{t}=T\left\{ \tsum\limits_{1\leq j\leq p}\left\vert
e_{j}\right) \left( e_{j}\right\vert -\tsum\limits_{p+1\leq j\leq
2r}\left\vert e_{j}\right) \left( e_{j}\right\vert \right\} T^{t}
\end{equation}%
and 
\begin{equation}
s=\tsum\limits_{1\leq j\leq r}\left\vert Te_{j}\vee Te_{j}\right) ,
\end{equation}%
where $T\in GL\left( 2r,\mathbb{R}\right) .$ Actually only transformation
belonging to $O\left( 2r,\mathbb{R}\right) \ltimes \mathcal{Z}\left(
GL\left( 2r,\mathbb{R}\right) \right) _{+}$ are needed.

$\lceil $Noting 
\begin{equation}
\left( T\wedge T\right) g=\left( T\wedge T\right) \mathcal{A}\left( G\right)
=\mathcal{A}\left( TGT^{t}\right) =\left\vert
e_{1}e_{2},e_{1}e_{3},e_{1}e_{4},e_{2}e_{3},e_{2}e_{4},e_{3}e_{4}\right)
\left( \mathbb{T}\wedge \mathbb{T}\right) \mathbf{g,}
\end{equation}%
where%
\begin{equation}
T_{\left( jk\right) \left( lm\right) }=\left( T\wedge T\right) _{\left(
jk\right) \left( lm\right) }=\frac{1}{2}\left(
T_{jk}T_{lm}-T_{jl}T_{km}-T_{mk}T_{lj}+T_{lm}T_{jk}\right) .
\end{equation}%
$\rfloor $

There are three fundemental antisymmetic forms in $\tbigwedge\nolimits^{2}%
\mathcal{H},$ where $\mathbb{R}$-$\dim \left\{ \mathcal{H}\right\} =4$%
\begin{gather}
G_{1}=\left\vert e_{1},e_{2},e_{3},e_{4}\right) \left( 
\begin{array}{cc}
\left( 
\begin{array}{cc}
0 & 1 \\ 
-1 & 0%
\end{array}%
\right) & \mathbb{O} \\ 
\mathbb{O} & \left( 
\begin{array}{cc}
0 & 1 \\ 
-1 & 0%
\end{array}%
\right)%
\end{array}%
\right) \left( e_{1},e_{2},e_{3},e_{4}\right\vert =\left\vert
-e_{2},e_{1},-e_{4},e_{3}\right) \left( e_{1},e_{2},e_{3},e_{4}\right\vert \\
\rightarrow \left\vert e_{1}\right) \left( e_{2}\right\vert -\left\vert
e_{2}\right) \left( e_{1}\right\vert +\left\vert e_{3}\right) \left(
e_{4}\right\vert -\left\vert e_{4}\right) \left( e_{3}\right\vert
=\left\vert e_{1}e_{2}\right) +\left\vert e_{3}e_{4}\right)
\end{gather}%
\begin{gather}
G_{2}=\left\vert e_{1},e_{2},e_{3},e_{4}\right) \left( 
\begin{array}{cc}
\mathbb{O} & \left( 
\begin{array}{cc}
-1 & 0 \\ 
0 & -1%
\end{array}%
\right) \\ 
\left( 
\begin{array}{cc}
1 & 0 \\ 
0 & 1%
\end{array}%
\right) & \mathbb{O}%
\end{array}%
\right) \left( e_{1},e_{2},e_{3},e_{4}\right\vert =\left\vert
-e_{3},-e_{4},e_{1},e_{2}\right) \left( e_{1},e_{2},e_{3},e_{4}\right\vert \\
\rightarrow \left\vert e_{1}\right) \left( e_{3}\right\vert -\left\vert
e_{3}\right) \left( e_{1}\right\vert +\left\vert e_{2}\right) \left(
e_{4}\right\vert -\left\vert e_{4}\right) \left( e_{2}\right\vert
=\left\vert e_{1}e_{3}\right) +\left\vert e_{2}e_{4}\right)
\end{gather}%
\begin{gather}
G_{3}=\left\vert e_{1},e_{2},e_{3},e_{4}\right) \left( 
\begin{array}{cc}
\mathbb{O} & \left( 
\begin{array}{cc}
0 & -1 \\ 
-1 & 0%
\end{array}%
\right) \\ 
\left( 
\begin{array}{cc}
0 & 1 \\ 
1 & 0%
\end{array}%
\right) & \mathbb{O}%
\end{array}%
\right) \left( e_{1},e_{2},e_{3},e_{4}\right\vert =\left\vert
-e_{4},-e_{3},e_{2},e_{1}\right) \left( e_{1},e_{2},e_{3},e_{4}\right\vert \\
\rightarrow \left\vert e_{1}\right) \left( e_{4}\right\vert -\left\vert
e_{4}\right) \left( e_{1}\right\vert +\left\vert e_{3}\right) \left(
e_{2}\right\vert -\left\vert e_{2}\right) \left( e_{3}\right\vert
=\left\vert e_{1}e_{4}\right) +\left\vert e_{2}e_{3}\right)
\end{gather}%
$\rfloor $

can be generated by by transformations $\mathcal{T}\in GL\left( 2r,\mathbb{R}%
\right) $ 
\begin{equation}
\mathcal{T}\left\vert \mathbf{e}_{g},\mathbf{e}_{g}^{\#}\right) \left( 
\begin{array}{cc}
\mathcal{T}_{11} & \mathcal{T}_{12} \\ 
\mathcal{T}_{21} & \mathcal{T}_{22}%
\end{array}%
\right) =\left\vert \left( \mathbf{e}_{g}\mathcal{T}_{11}+\mathbf{e}_{g}^{\#}%
\mathcal{T}_{21}\right) ,\left( \mathbf{e}_{g}\mathcal{T}_{12}+\mathbf{e}%
_{g}^{\#}\mathcal{T}_{22}\right) \right) =\left\vert \mathbf{f}_{g},\mathbf{f%
}_{g}^{\#}\right)
\end{equation}%
The space $V_{2r}$ can then be made into a complex space by the
decomposition 
\begin{equation}
V_{\mathbb{C}2r}=V_{g}\oplus J_{g}V_{g}^{\#},
\end{equation}%
where the complexification corresponds to the geminal 
\begin{equation}
g=\sum\limits_{1\leq j<k\leq 2r}g_{jk}\left\vert e_{j}e_{k}\right)
=\sum\limits_{1\leq j\leq r}c_{gj}\left\vert \widetilde{e}_{gj}\widetilde{e}%
_{gj+r}\right) =\sum\limits_{1\leq j\leq r}\left\vert e_{gj}e_{gj+r}\right) .
\end{equation}%
Each $g$ defines a symplectic group

\begin{equation}
Sp\left( 2r,\mathbb{R}\right) =\left\{ \left. T\right\vert g\left(
Tv_{1},Tv_{2}\right) =g\left( v_{1},v_{2}\right) \right\}
\end{equation}%
the transformations $T$ can effect both $\left\{ e_{gj}\right\} $ and $%
\left\{ c_{gj}\right\} $ i.e. by changing the normalization of $\left\{
e_{gj}\right\} .$

The real space $V_{2r}$ has a natural inner product 
\begin{equation}
\left( \left. {}\right\vert \right) :V_{2r}\times V_{2r}\rightarrow \mathbb{R%
}
\end{equation}%
such that 
\begin{equation}
\left( \left. \widetilde{e}_{gj}\right\vert \widetilde{e}_{gk}\right)
=\delta _{jk}.
\end{equation}%
One can define a new inner product 
\begin{equation}
\left( \left. {}\right\vert \right) _{g}:V_{2r}\times V_{2r}\rightarrow 
\mathbb{R}
\end{equation}%
by 
\begin{equation}
\left( \left. x\right\vert y\right) _{g}=\left( \left. x\right\vert
C_{g}y\right) ,
\end{equation}%
where 
\begin{eqnarray}
\left\vert e_{gj}\right) &=&c_{gj}^{\frac{1}{2}}\left\vert \widetilde{e}%
_{gj}\right) \\
\left\vert e_{gj+r}\right) &=&c_{gj}^{\frac{1}{2}}\left\vert \widetilde{e}%
_{gj+r}\right) ,
\end{eqnarray}%
\begin{equation}
C_{g}^{-1}=\sum\limits_{1\leq j\leq r}c_{gj}^{-1}\left\{ \left\vert 
\widetilde{e}_{gj}\right) \left( \widetilde{e}_{gj}\right\vert +\left\vert 
\widetilde{e}_{gj+r}\right) \left( \widetilde{e}_{gj+r}\right\vert \right\} .
\end{equation}%
This gives 
\begin{equation}
\left( \left. e_{gj}\right\vert e_{gk}\right) _{g}=\left( \left. c_{gj}^{%
\frac{1}{2}}\widetilde{e}_{gj}\right\vert C_{g}^{-1}c_{g}^{\frac{1}{2}}%
\widetilde{e}_{gk}\right) =\delta _{jk}.
\end{equation}%
This defines the orthogonal group 
\begin{equation}
O\left( 2r,\mathbb{R}\right) =\left\{ \left. \overset{}{T}\right\vert \left(
\left. Tv_{1}\right\vert Tv_{2}\right) _{g}=\left( \left. v_{1}\right\vert
v_{2}\right) _{g}\right\} =\left\{ \left. \overset{}{T}\right\vert \left(
\left. Tv_{1}\right\vert Tv_{2}\right) =\left( \left. v_{1}\right\vert
v_{2}\right) \right\} .
\end{equation}%
One can see that these two definitions are equivalent as 
\begin{eqnarray}
\left( \left. Tv_{1}\right\vert Tv_{2}\right) _{g} &=&\left( \left.
Tv_{1}\right\vert C_{g}^{-1}Tv_{2}\right) =\left( \left. v_{1}\right\vert
T^{t}C_{g}^{-1}Tv_{2}\right) \\
T^{t}C_{g}^{-1}T &=&C_{g}^{-1}\Leftrightarrow C_{g}^{\frac{1}{2}%
}T^{t}C_{g}^{-\frac{1}{2}}C_{g}^{-\frac{1}{2}}TC_{g}^{\frac{1}{2}}=I_{2r}=%
\widetilde{T}^{t}\widetilde{T},
\end{eqnarray}%
where 
\begin{equation}
\widetilde{T}=C_{g}^{-\frac{1}{2}}TC_{g}^{\frac{1}{2}}\Leftrightarrow
T=C_{g}^{\frac{1}{2}}\widetilde{T}C_{g}^{-\frac{1}{2}}.
\end{equation}%
Hence the group of operators $\left\{ \left. \overset{}{\widetilde{T}}%
\right\vert \left( \left. \widetilde{T}v_{1}\right\vert \widetilde{T}%
v_{2}\right) =\left( \left. v_{1}\right\vert v_{2}\right) \right\} $ are
isomorphic to the group of operators $\left\{ \left. \overset{}{T}%
\right\vert \left( \left. Tv_{1}\right\vert Tv_{2}\right) _{g}=\left( \left.
v_{1}\right\vert v_{2}\right) _{g}\right\} $ and both define the same group $%
O\left( 2r,\mathbb{R}\right) .$

\begin{remark}
This shows that one can always define a real symmetric non degenerate inner
product in $V_{2r}$ such that any antisymmetric form $g$ can be expressed in
standard form 
\begin{equation}
g=\sum\limits_{1\leq j\leq r}\left\vert e_{gj}e_{gj+r}\right)
\end{equation}%
in terms of orthonormal vectors wrt this inner product.
\end{remark}

\paragraph{General Inner Products}

If the inner product $\left( \left. {}\right\vert \right) $ is symmetric and
positive definite one can define new inner products by 
\begin{equation}
\left( \left. v_{1}\right\vert v_{2}\right) _{K}=\left( \left.
v_{1}\right\vert Kv_{2}\right) .
\end{equation}%
if they are symmetric then 
\begin{equation}
K=K^{t}
\end{equation}%
if they are symmetric and positive definite then 
\begin{equation}
K>0\Rightarrow K=K^{t}
\end{equation}%
if they are antisymmetric then 
\begin{equation}
K=-K^{t}
\end{equation}%
if they are non degenerate then 
\begin{equation}
K\in GL\left( 2r,\mathbb{R}\right) .
\end{equation}%
One can define generalized adjoints by 
\begin{equation}
\left( \left. A^{\tau }v_{1}\right\vert v_{2}\right) _{K}=\left( \left.
v_{1}\right\vert Av_{2}\right) _{K}
\end{equation}%
thus 
\begin{equation}
\left( \left. KA^{\tau }v_{1}\right\vert v_{2}\right) =\left( \left.
A^{t}Kv_{1}\right\vert v_{2}\right)
\end{equation}%
hence 
\begin{align}
KA^{\tau }& =A^{t}K \\
& \Downarrow \\
A^{\tau }& =K^{-1}A^{t}K.
\end{align}

Let 
\begin{equation}
g\left( v_{1},v_{2}\right) =\left( \left. v_{1}\right\vert Gv_{2}\right)
\end{equation}%
then 
\begin{align}
\left( \left. v_{1}\right\vert Gv_{2}\right) & =-\left( \left.
v_{2}\right\vert Gv_{1}\right) \Longleftrightarrow -\left( \left.
G^{t}v_{2}\right\vert v_{1}\right) =-\left( \left. v_{1}\right\vert
G^{t}v_{2}\right) \\
& \Updownarrow \\
G^{t}& =-G.
\end{align}%
$\lceil $All other antisymmetric bilinear forms on $V_{2r}$ can be defined
by 
\begin{equation}
g_{K}\left( v_{1},v_{2}\right) =-g_{K}\left( v_{2},v_{1}\right) =\left(
\left. v_{1}\right\vert Gv_{2}\right) _{K},
\end{equation}%
where 
\begin{equation}
G^{\tau }=-G.
\end{equation}%
Thus 
\begin{equation}
K^{-1}G^{t}K=-G.
\end{equation}%
$\rfloor $

\emph{Transformations that belong to }%
\begin{equation}
Sp\left( 2r,\mathbb{R}\right) \cap O\left( 2r,\mathbb{R}\right)
\end{equation}
\emph{leave }$\left\{ c_{gj}\right\} $\emph{\ invariant, while changing the
conb }$\left\{ e_{gj}\right\} $\emph{\ to another conb }$\left\{
f_{gj}\right\} .$

\section{Group Decompositions}

\subsection{Unitary transformations $U\left( 2r\right) $}

Unitary transformations $U\left( 2r\right) $ are produced by the Lie algebra
of anti self adjoint operators, which can be expanded as in terms of real
and imaginary parts (wrt a given basis). The imaginary parts generate the
subgroup $O\left( 2r,\mathbb{R}\right) $ 
\begin{align}
u & =\exp i\left\{ \sum_{1\leq j,k\leq2r}\lambda_{jk}\left( \left\vert
e_{j}\right\rangle \left\langle e_{k}\right\vert \right) \right\} \\
& =\exp i\left\{ \sum_{1\leq j<k\leq2r}\left( \func{Re}\left(
\lambda_{jk}\right) +i\func{Im}\left( \lambda_{jk}\right) \right) \left\vert
e_{j}\right\rangle \left\langle e_{k}\right\vert +\sum_{1\leq
j>k\leq2r}\left( \func{Re}\left( \lambda_{jk}\right) -i\func{Im}\left(
\lambda_{jk}\right) \right) \left\vert e_{k}\right\rangle \left\langle
e_{j}\right\vert \right. \\
& \left. +\sum_{1\leq j\leq2r}\func{Re}\left( \lambda_{jj}\right) \left\vert
e_{j}\right\rangle \left\langle e_{j}\right\vert \right\} \\
& =\exp i\left\{ \sum_{1\leq j<k\leq2r}\func{Re}\left( \lambda _{jk}\right)
\left( \left\vert e_{j}\right\rangle \left\langle e_{k}\right\vert
+\left\vert e_{k}\right\rangle \left\langle e_{j}\right\vert \right) +i\func{%
Im}\left( \lambda_{jk}\right) \left( \left\vert e_{j}\right\rangle
\left\langle e_{k}\right\vert -\left\vert e_{k}\right\rangle \left\langle
e_{j}\right\vert \right) \right\} \exp i\left\{ \sum_{1\leq j\leq2r}\func{Re}%
\left( \lambda_{jj}\right) \left\vert e_{j}\right\rangle \left\langle
e_{j}\right\vert \right\}
\end{align}%
\begin{align}
& =\exp i\left\{ \sum_{1\leq j<k\leq2r}\func{Re}\left( \lambda _{jk}\right)
\left( \left\vert e_{j}\right\rangle \left\langle e_{k}\right\vert
+\left\vert e_{k}\right\rangle \left\langle e_{j}\right\vert \right) +i\func{%
Im}\left( \lambda_{jk}\right) \left( \left\vert e_{j}\right\rangle
\left\langle e_{k}\right\vert -\left\vert e_{k}\right\rangle \left\langle
e_{j}\right\vert \right) \right\} \exp i\left\{ \sum_{1\leq j\leq2r}\func{Re}%
\left( \lambda_{jj}\right) \left\vert e_{j}\right\rangle \left\langle
e_{j}\right\vert \right\} \\
& =\exp\left\{ \sum_{1\leq j<k\leq2r}\func{Re}\left( \lambda _{jk}\right)
i\left( \left\vert e_{j}\right\rangle \left\langle e_{k}\right\vert
+\left\vert e_{k}\right\rangle \left\langle e_{j}\right\vert \right) -\func{%
Im}\left( \lambda_{jk}\right) \left( \left\vert e_{j}\right\rangle
\left\langle e_{k}\right\vert -\left\vert e_{k}\right\rangle \left\langle
e_{j}\right\vert \right) \right\} \exp i\left\{ \sum_{1\leq j\leq2r}\func{Re}%
\left( \lambda_{jj}\right) \left\vert e_{j}\right\rangle \left\langle
e_{j}\right\vert \right\} \\
& =\exp\left\{ \sum_{1\leq j<k\leq2r}\alpha_{jk}i\left( \left\vert
e_{j}\right\rangle \left\langle e_{k}\right\vert +\left\vert
e_{k}\right\rangle \left\langle e_{j}\right\vert \right) -\beta_{jk}\left(
\left\vert e_{j}\right\rangle \left\langle e_{k}\right\vert -\left\vert
e_{k}\right\rangle \left\langle e_{j}\right\vert \right) \right\} \exp
i\left\{ \sum_{1\leq j\leq2r}\gamma_{j}i\left\vert e_{j}\right\rangle
\left\langle e_{j}\right\vert \right\} ;\alpha_{jk},\beta_{jk},\gamma_{j}\in%
\mathbb{R} \\
& =\exp\left\{ \sum_{1\leq j<k\leq2r}\tilde{\alpha}_{jk}i\left( \left\vert
e_{j}\right\rangle \left\langle e_{k}\right\vert +\left\vert
e_{k}\right\rangle \left\langle e_{j}\right\vert \right) \right\}
\exp\left\{ \sum_{1\leq j<k\leq2r}\tilde{\beta}_{jk}\left( \left\vert
e_{j}\right\rangle \left\langle e_{k}\right\vert -\left\vert
e_{k}\right\rangle \left\langle e_{j}\right\vert \right) \right\} \times \\
& \times\exp\left\{ \sum_{1\leq j\leq2r}\gamma_{j}i\left\vert
e_{j}\right\rangle \left\langle e_{j}\right\vert \right\} ;\tilde{\alpha}%
_{jk},\tilde{\beta}_{jk},\gamma_{j}\in\mathbb{R}
\end{align}
The unitary group $U\left( 2r\right) $ has the coset factorization 
\begin{equation}
U\left( 2r\right) =\frac{U\left( 2r\right) }{O\left( 2r,\mathbb{R}\right) }%
O\left( 2r,\mathbb{R}\right) ,
\end{equation}
leading to the Lie algebra decomposition%
\begin{equation}
\mathcal{U}\left( 2r\right) =\mathcal{P}\left( 2r\right) \mathcal{\oplus O}%
\left( 2r,\mathbb{R}\right) .
\end{equation}

$\left[ \mathcal{P}\left( 2r\right) ,\mathcal{O}\left( 2r,\mathbb{R}\right) %
\right] \subset \mathcal{P}\left( 2r\right) $ is shown by 
\begin{align}
& \left( \left\vert e_{j}\right\rangle \left\langle e_{k}\right\vert
+\left\vert e_{k}\right\rangle \left\langle e_{j}\right\vert \right) \left(
\left\vert e_{p}\right\rangle \left\langle e_{q}\right\vert -\left\vert
e_{q}\right\rangle \left\langle e_{p}\right\vert \right) =\left\vert
e_{j}\right\rangle \left\langle e_{k}\right\vert \left\vert
e_{p}\right\rangle \left\langle e_{q}\right\vert -\left\vert
e_{j}\right\rangle \left\langle e_{k}\right\vert \left\vert
e_{q}\right\rangle \left\langle e_{p}\right\vert +\left\vert
e_{k}\right\rangle \left\langle e_{j}\right\vert \left\vert
e_{p}\right\rangle \left\langle e_{q}\right\vert -\left\vert
e_{k}\right\rangle \left\langle e_{j}\right\vert \left\vert
e_{q}\right\rangle \left\langle e_{p}\right\vert \\
& =\left\vert e_{j}\right\rangle \left\langle e_{q}\right\vert \delta
_{kp}-\left\vert e_{j}\right\rangle \left\langle e_{p}\right\vert \delta
_{kq}-\left\vert e_{k}\right\rangle \left\langle e_{p}\right\vert \delta
_{jp}+\left\vert e_{k}\right\rangle \left\langle e_{p}\right\vert \delta
_{jq} \\
& \left( \left\vert e_{p}\right\rangle \left\langle e_{q}\right\vert
-\left\vert e_{q}\right\rangle \left\langle e_{p}\right\vert \right) \left(
\left\vert e_{j}\right\rangle \left\langle e_{k}\right\vert +\left\vert
e_{k}\right\rangle \left\langle e_{j}\right\vert \right) =\left\vert
e_{p}\right\rangle \left\langle e_{q}\right\vert \left\vert
e_{j}\right\rangle \left\langle e_{k}\right\vert +\left\vert
e_{p}\right\rangle \left\langle e_{q}\right\vert \left\vert
e_{k}\right\rangle \left\langle e_{j}\right\vert -\left\vert
e_{q}\right\rangle \left\langle e_{p}\right\vert \left\vert
e_{j}\right\rangle \left\langle e_{k}\right\vert -\left\vert
e_{q}\right\rangle \left\langle e_{p}\right\vert \left\vert
e_{k}\right\rangle \left\langle e_{j}\right\vert \\
& =\left\vert e_{p}\right\rangle \left\langle e_{k}\right\vert \delta
_{jq}+\left\vert e_{p}\right\rangle \left\langle e_{j}\right\vert \delta
_{kq}-\left\vert e_{q}\right\rangle \left\langle e_{k}\right\vert \delta
_{jp}-\left\vert e_{q}\right\rangle \left\langle e_{j}\right\vert \delta
_{kp}
\end{align}%
\begin{align}
& \left[ \left( \left\vert e_{j}\right\rangle \left\langle e_{k}\right\vert
+\left\vert e_{k}\right\rangle \left\langle e_{j}\right\vert \right) ,\left(
\left\vert e_{p}\right\rangle \left\langle e_{q}\right\vert -\left\vert
e_{q}\right\rangle \left\langle e_{p}\right\vert \right) \right] =\left(
\left\vert e_{j}\right\rangle \left\langle e_{q}\right\vert +\left\vert
e_{q}\right\rangle \left\langle e_{j}\right\vert \right) \delta _{kp}-\left(
\left\vert e_{p}\right\rangle \left\langle e_{j}\right\vert +\left\vert
e_{j}\right\rangle \left\langle e_{p}\right\vert \right) \delta _{kq} \\
& -\left( \left\vert e_{p}\right\rangle \left\langle e_{k}\right\vert
+\left\vert e_{k}\right\rangle \left\langle e_{p}\right\vert \right) \delta
_{jq}+\left( \left\vert e_{k}\right\rangle \left\langle e_{p}\right\vert
+\left\vert e_{q}\right\rangle \left\langle e_{k}\right\vert \right) \delta
_{jp}
\end{align}%
\begin{equation}
\left[ \left( \left\vert e_{j}\right\rangle \left\langle e_{k}\right\vert
+\left\vert e_{k}\right\rangle \left\langle e_{j}\right\vert \right) ,\left(
\left\vert e_{p}\right\rangle \left\langle e_{q}\right\vert -\left\vert
e_{q}\right\rangle \left\langle e_{p}\right\vert \right) \right]
=P_{jq}\delta _{kp}-P_{pj}\delta _{kq}-P_{pk}\delta _{jq}+P_{kq}\delta _{jp}
\end{equation}%
$\left[ \mathcal{P}\left( 2r\right) ,\mathcal{P}\left( 2r\right) \right]
\subset \mathcal{O}\left( 2r,\mathbb{R}\right) $ is shown by 
\begin{align}
\left( \left\vert e_{j}\right\rangle \left\langle e_{k}\right\vert
+\left\vert e_{k}\right\rangle \left\langle e_{j}\right\vert \right) \left(
\left\vert e_{p}\right\rangle \left\langle e_{q}\right\vert +\left\vert
e_{q}\right\rangle \left\langle e_{p}\right\vert \right) & =\left\vert
e_{j}\right\rangle \left\langle e_{k}\right\vert \left\vert
e_{p}\right\rangle \left\langle e_{q}\right\vert +\left\vert
e_{j}\right\rangle \left\langle e_{k}\right\vert \left\vert
e_{q}\right\rangle \left\langle e_{p}\right\vert +\left\vert
e_{k}\right\rangle \left\langle e_{j}\right\vert \left\vert
e_{p}\right\rangle \left\langle e_{q}\right\vert +\left\vert
e_{k}\right\rangle \left\langle e_{j}\right\vert \left\vert
e_{q}\right\rangle \left\langle e_{p}\right\vert \\
& =\left\vert e_{j}\right\rangle \left\langle e_{q}\right\vert \delta
_{kp}+\left\vert e_{j}\right\rangle \left\langle e_{p}\right\vert \delta
_{kq}+\left\vert e_{k}\right\rangle \left\langle e_{p}\right\vert \delta
_{jp}+\left\vert e_{k}\right\rangle \left\langle e_{p}\right\vert \delta
_{jq} \\
\left( \left\vert e_{p}\right\rangle \left\langle e_{q}\right\vert
+\left\vert e_{q}\right\rangle \left\langle e_{p}\right\vert \right) \left(
\left\vert e_{j}\right\rangle \left\langle e_{k}\right\vert +\left\vert
e_{k}\right\rangle \left\langle e_{j}\right\vert \right) & =\left\vert
e_{p}\right\rangle \left\langle e_{q}\right\vert \left\vert
e_{j}\right\rangle \left\langle e_{k}\right\vert +\left\vert
e_{p}\right\rangle \left\langle e_{q}\right\vert \left\vert
e_{k}\right\rangle \left\langle e_{j}\right\vert +\left\vert
e_{q}\right\rangle \left\langle e_{p}\right\vert \left\vert
e_{j}\right\rangle \left\langle e_{k}\right\vert +\left\vert
e_{q}\right\rangle \left\langle e_{p}\right\vert \left\vert
e_{k}\right\rangle \left\langle e_{j}\right\vert \\
& =\left\vert e_{p}\right\rangle \left\langle e_{k}\right\vert \delta
_{jq}+\left\vert e_{p}\right\rangle \left\langle e_{j}\right\vert \delta
_{kq}+\left\vert e_{q}\right\rangle \left\langle e_{k}\right\vert \delta
_{jp}+\left\vert e_{q}\right\rangle \left\langle e_{j}\right\vert \delta
_{kp}
\end{align}%
\begin{align}
\left[ \left( \left\vert e_{j}\right\rangle \left\langle e_{k}\right\vert
+\left\vert e_{k}\right\rangle \left\langle e_{j}\right\vert \right) ,\left(
\left\vert e_{p}\right\rangle \left\langle e_{q}\right\vert +\left\vert
e_{q}\right\rangle \left\langle e_{p}\right\vert \right) \right] & =\left(
\left\vert e_{j}\right\rangle \left\langle e_{q}\right\vert -\left\vert
e_{q}\right\rangle \left\langle e_{j}\right\vert \right) \delta _{kp}+\left(
\left\vert e_{j}\right\rangle \left\langle e_{p}\right\vert -\left\vert
e_{p}\right\rangle \left\langle e_{j}\right\vert \right) \delta _{kq} \\
& +\left( \left\vert e_{k}\right\rangle \left\langle e_{p}\right\vert
-\left\vert e_{p}\right\rangle \left\langle e_{k}\right\vert \right) \delta
_{jq}+\left( \left\vert e_{k}\right\rangle \left\langle e_{q}\right\vert
-\left\vert e_{q}\right\rangle \left\langle e_{k}\right\vert \right) \delta
_{jp}
\end{align}%
\begin{equation}
\left[ \left( \left\vert e_{j}\right\rangle \left\langle e_{k}\right\vert
+\left\vert e_{k}\right\rangle \left\langle e_{j}\right\vert \right) ,\left(
\left\vert e_{p}\right\rangle \left\langle e_{q}\right\vert +\left\vert
e_{q}\right\rangle \left\langle e_{p}\right\vert \right) \right]
=O_{jq}\delta _{kp}+O_{jp}\delta _{kq}+O_{kp}\delta _{jq}+O_{kq}\delta _{jp}
\end{equation}%
$\left[ \mathcal{O}\left( 2r,\mathbb{R}\right) ,\mathcal{O}\left( 2r,\mathbb{%
R}\right) \right] \subset \mathcal{O}\left( 2r,\mathbb{R}\right) $ is shown
by 
\begin{align}
\left( \left\vert e_{j}\right\rangle \left\langle e_{k}\right\vert
-\left\vert e_{k}\right\rangle \left\langle e_{j}\right\vert \right) \left(
\left\vert e_{p}\right\rangle \left\langle e_{q}\right\vert -\left\vert
e_{q}\right\rangle \left\langle e_{p}\right\vert \right) & =\left\vert
e_{j}\right\rangle \left\langle e_{k}\right\vert \left\vert
e_{p}\right\rangle \left\langle e_{q}\right\vert -\left\vert
e_{j}\right\rangle \left\langle e_{k}\right\vert \left\vert
e_{q}\right\rangle \left\langle e_{p}\right\vert \\
-\left\vert e_{k}\right\rangle \left\langle e_{j}\right\vert \left\vert
e_{p}\right\rangle \left\langle e_{q}\right\vert +\left\vert
e_{k}\right\rangle \left\langle e_{j}\right\vert \left\vert
e_{q}\right\rangle \left\langle e_{p}\right\vert & =\left\vert
e_{j}\right\rangle \left\langle e_{q}\right\vert \delta _{kp}-\left\vert
e_{j}\right\rangle \left\langle e_{p}\right\vert \delta _{kq}-\left\vert
e_{k}\right\rangle \left\langle e_{p}\right\vert \delta _{jp}+\left\vert
e_{k}\right\rangle \left\langle e_{p}\right\vert \delta _{jq} \\
-\left( \left\vert e_{p}\right\rangle \left\langle e_{q}\right\vert
-\left\vert e_{q}\right\rangle \left\langle e_{p}\right\vert \right) \left(
\left\vert e_{j}\right\rangle \left\langle e_{k}\right\vert -\left\vert
e_{k}\right\rangle \left\langle e_{j}\right\vert \right) & =-\left\vert
e_{p}\right\rangle \left\langle e_{q}\right\vert \left\vert
e_{j}\right\rangle \left\langle e_{k}\right\vert +\left\vert
e_{p}\right\rangle \left\langle e_{q}\right\vert \left\vert
e_{k}\right\rangle \left\langle e_{j}\right\vert \\
+\left\vert e_{q}\right\rangle \left\langle e_{p}\right\vert \left\vert
e_{j}\right\rangle \left\langle e_{k}\right\vert -\left\vert
e_{q}\right\rangle \left\langle e_{p}\right\vert \left\vert
e_{k}\right\rangle \left\langle e_{j}\right\vert & =-\left\vert
e_{p}\right\rangle \left\langle e_{k}\right\vert \delta _{jq}+\left\vert
e_{p}\right\rangle \left\langle e_{j}\right\vert \delta _{kq}+\left\vert
e_{q}\right\rangle \left\langle e_{k}\right\vert \delta _{jp}-\left\vert
e_{q}\right\rangle \left\langle e_{j}\right\vert \delta _{kp}
\end{align}%
\begin{align}
\left[ \left( \left\vert e_{j}\right\rangle \left\langle e_{k}\right\vert
-\left\vert e_{k}\right\rangle \left\langle e_{j}\right\vert \right) ,\left(
\left\vert e_{p}\right\rangle \left\langle e_{q}\right\vert -\left\vert
e_{q}\right\rangle \left\langle e_{p}\right\vert \right) \right] & =\left(
\left\vert e_{j}\right\rangle \left\langle e_{q}\right\vert -\left\vert
e_{q}\right\rangle \left\langle e_{j}\right\vert \right) \delta _{kp}-\left(
\left\vert e_{j}\right\rangle \left\langle e_{p}\right\vert -\left\vert
e_{p}\right\rangle \left\langle e_{j}\right\vert \right) \delta _{kq} \\
+\left( \left\vert e_{k}\right\rangle \left\langle e_{p}\right\vert
-\left\vert e_{p}\right\rangle \left\langle e_{k}\right\vert \right) \delta
_{jq}& -\left( \left\vert e_{q}\right\rangle \left\langle e_{k}\right\vert
-\left\vert e_{k}\right\rangle \left\langle e_{q}\right\vert \right) \delta
_{jp}
\end{align}

The preceeding set of Lie group generators satisfy the following
relationships 
\begin{align}
\left[ O_{jk},O_{pq}\right] &
=O_{jq}\delta_{kp}-O_{jp}\delta_{kq}+O_{kp}\delta_{jq}-O_{qk}\delta_{jp} \\
\left[ O_{jk},O_{kl}\right] & =O_{jl} \\
O_{jk} & =-O_{kj}.
\end{align}

\subsection{Orthogonal Transformations $O\left( 2r,\mathbb{R}\right) $}

The orthogonal group $O\left( 2r,\mathbb{R}\right) $ has a coset
factorization 
\begin{equation}
O\left( 2r,\mathbb{R}\right) =\frac{O\left( 2r,\mathbb{R}\right) }{U\left(
r\right) }U\left( r\right) =\left( \frac{O\left( 2r,\mathbb{R}\right) }{%
O\left( 2r,\mathbb{R}\right) \cap Sp_{b}\left( 2r,\mathbb{R}\right) }\right)
\left( O\left( 2r,\mathbb{R}\right) \cap Sp_{b}\left( 2r,\mathbb{R}\right)
\right) .
\end{equation}%
Unitary transformations, $U\left( r\right) ,$ maintains the number of
"particles" i.e are number gauge invariant (gauge transformations are
generated by number operators, so gauge invariance is number operator
dependent), classes of unitary transformations are associated with a
specific number operator $\mathcal{N}_{b}$, thus a specific
complexification, $\widehat{J_{b}}$ and hence a specific realization of $%
Sp\left( 2r,\mathbb{R}\right) .$ All complex structures are generated by the
action of $O\left( 2r,\mathbb{R}\right) $ on a standard $\widehat{J_{a}}$ by 
\begin{equation}
\widehat{J_{b}}=O\widehat{J_{a}}O^{t};O\in O\left( 2r,\mathbb{R}\right) .
\end{equation}%
Clearly 
\begin{equation}
\widehat{J_{a}}=O\widehat{J_{a}}O^{t}
\end{equation}%
if 
\begin{equation}
O\in Sp_{a}\left( 2r,\mathbb{R}\right) ,
\end{equation}%
while 
\begin{equation}
\widehat{J_{a}}\neq O\widehat{J_{a}}O^{t}
\end{equation}%
if 
\begin{equation}
O\notin Sp_{a}\left( 2r,\mathbb{R}\right)
\end{equation}%
i.e. 
\begin{equation}
O\in \frac{O\left( 2r,\mathbb{R}\right) }{O\left( 2r,\mathbb{R}\right) \cap
Sp_{b}\left( 2r,\mathbb{R}\right) }.
\end{equation}

$O\left( 2r,\mathbb{R}\right) $ transformations map $\mathcal{F}%
_{1}\rightarrow\mathcal{F}_{1}$ via inner automorphisms based on the Lie
algebra $\mathcal{F}_{2}$, while $U\left( 2r\right) $ are represented by
inner automorphisms based on the Lie algebra $\mathcal{F}_{even}.$

\subsection{Symplectic Transformations $Sp\left( 2r,\mathbb{R}\right) $}

A set of Lie group generators of $sp\left( 2r,\mathbb{R}\right) $ are given
by 
\begin{align}
Z_{jk} & =\left\vert e_{j}\right\rangle \left\langle e_{k}\right\vert
-\sigma\left( j\right) \sigma\left( k\right) \left( \left\vert
e_{-k}\right\rangle \left\langle e_{-j}\right\vert \right) \\
\sigma\left( j\right) & =1;j>0 \\
\sigma\left( j\right) & =-1;j<0 \\
-r & \leq j,k\leq r.
\end{align}
$\lceil$ Examples:%
\begin{align}
Z_{1-1} & =\left\vert e_{1}\right\rangle \left\langle e_{-1}\right\vert
-\sigma\left( 1\right) \sigma\left( -1\right) \left( \left\vert
e_{1}\right\rangle \left\langle e_{-1}\right\vert \right) =2\left\vert
e_{1}\right\rangle \left\langle e_{-1}\right\vert \\
Z_{1-2} & =\left\vert e_{1}\right\rangle \left\langle e_{-2}\right\vert
-\sigma\left( 1\right) \sigma\left( -2\right) \left( \left\vert
e_{2}\right\rangle \left\langle e_{-1}\right\vert \right) =\left\vert
e_{1}\right\rangle \left\langle e_{-2}\right\vert +\left( \left\vert
e_{2}\right\rangle \left\langle e_{-1}\right\vert \right)
\end{align}

$\rfloor$

This leads to the relationships 
\begin{align}
\left[ Z_{jk},Z_{kl}\right] & =Z_{jl} \\
Z_{-k-j} & =-\sigma\left( j\right) \sigma\left( k\right) Z_{jk}.
\end{align}

Not all transformations belonging to $Sp\left( 2r,\mathbb{R}\right) $ are
inner automorphisms based on the Lie algebra $\mathcal{F}_{2}$, those that
are belong to $O\left( 2r,\mathbb{R}\right) $ and hence to $U\left( r\right)
.$ All transformations belonging to $Sp\left( 2r,\mathbb{R}\right) $ are
inner automorphisms based on the Lie algebra $\mathcal{F}_{even}.$

\section{Extra}

\begin{align}
\mathcal{\sigma}\left( u\right) & =\left( \left\vert \mathbf{e}\right\rangle
,J\left\vert \mathbf{e}\right\rangle \right) R_{e}\left( \sigma\right)
R_{e}\left( u\right) =\left( \left\vert \mathbf{e}\right\rangle ,J\left\vert 
\mathbf{e}\right\rangle \right) \left( 
\begin{array}{cc}
\mathbb{I}_{2r} & 0 \\ 
0 & -\mathbb{I}_{2r}%
\end{array}
\right) \left( 
\begin{array}{c}
\mathbf{a}_{1} \\ 
\mathbf{a}_{2}%
\end{array}
\right) =\left( \left\vert \mathbf{e}\right\rangle ,J\left\vert \mathbf{e}%
\right\rangle \right) \left( 
\begin{array}{c}
\mathbf{b}_{1} \\ 
\mathbf{b}_{2}%
\end{array}
\right) \\
\left( 
\begin{array}{c}
\mathbf{b}_{1} \\ 
\mathbf{b}_{2}%
\end{array}
\right) & =\left( 
\begin{array}{c}
\mathbf{a}_{1} \\ 
-\mathbf{a}_{2}%
\end{array}
\right) \\
J_{\mathbb{C}}\left( u\right) & =\left( \left\vert \mathbf{e}\right\rangle
,J\left\vert \mathbf{e}\right\rangle \right) R_{e}\left( J_{\mathbb{C}%
}\sigma\right) R_{e}\left( u\right) =\left( \left\vert \mathbf{e}%
\right\rangle ,J\left\vert \mathbf{e}\right\rangle \right) \left( 
\begin{array}{cc}
0 & -\mathbb{I}_{2r} \\ 
\mathbb{I}_{2r} & 0%
\end{array}
\right) \left( 
\begin{array}{c}
\mathbf{a}_{1} \\ 
\mathbf{a}_{2}%
\end{array}
\right) =\left( \left\vert \mathbf{e}\right\rangle ,J\left\vert \mathbf{e}%
\right\rangle \right) \left( 
\begin{array}{c}
\mathbf{b}_{1} \\ 
\mathbf{b}_{2}%
\end{array}
\right) \\
\left( 
\begin{array}{c}
\mathbf{b}_{1} \\ 
\mathbf{b}_{2}%
\end{array}
\right) & =\left( 
\begin{array}{c}
-\mathbf{a}_{2} \\ 
\mathbf{a}_{1}%
\end{array}
\right) \\
\mathcal{\sigma}\left( u\right) & =u\Longleftrightarrow\left( 
\begin{array}{c}
\mathbf{a}_{1} \\ 
\mathbf{a}_{2}%
\end{array}
\right) =\left( 
\begin{array}{c}
\mathbf{a}_{1} \\ 
-\mathbf{a}_{2}%
\end{array}
\right) \Longleftrightarrow\mathbf{a}_{2}=0\text{ and }u=\left( 
\begin{array}{c}
\mathbf{a}_{1} \\ 
\mathbf{0}%
\end{array}
\right) .
\end{align}

\begin{align}
\mathcal{H} & =\left\{ x=\left( \left\vert \mathbf{e}\right\rangle
,J\left\vert \mathbf{e}\right\rangle ,\left\vert \mathbf{e}\right\rangle
,J\left\vert \mathbf{e}\right\rangle \right) \mathbf{x}=\left\vert \mathbf{e}%
,J\mathbf{e},\mathbf{e},J\mathbf{e}\right\rangle \left( 
\begin{array}{c}
\mathbf{a}_{1} \\ 
\mathbf{a}_{2} \\ 
\mathbf{a}_{3} \\ 
\mathbf{a}_{4}%
\end{array}
\right) \right\} \\
\mathcal{R}_{e}\left( x\right) & =\left( 
\begin{array}{c}
\mathbf{a}_{1} \\ 
\mathbf{a}_{2} \\ 
\mathbf{a}_{3} \\ 
\mathbf{a}_{4}%
\end{array}
\right) \\
x & =\dsum \limits_{1\leq j\leq r}\left\{ \left\vert e_{j}\right\rangle
a_{1j}+J\left\vert e_{j}\right\rangle a_{2j}+\left\vert e_{j}\right\rangle
a_{3j}-J\left\vert e_{j}\right\rangle a_{4j}\right\}
\end{align}%
\begin{align}
\Sigma_{\sigma} & :\mathcal{H}\rightarrow\mathcal{H} \\
\Sigma_{\sigma}\left( u+\sigma\left( v\right) \right) & =v+\sigma\left(
u\right) , \\
u,v & \in\mathfrak{H}
\end{align}%
\begin{align}
\Sigma_{\sigma}\left( u+\sigma\left( v\right) \right) & =\left\vert \mathbf{e%
},J\mathbf{e},\mathbf{e},J\mathbf{e}\right\rangle \mathcal{R}_{e}\left(
\Sigma_{\sigma}\left( u+\sigma\left( v\right) \right) \right) \\
& =\left\vert \mathbf{e},J\mathbf{e},\mathbf{e},J\mathbf{e}\right\rangle
\left( 
\begin{array}{c}
\mathbf{b}_{1} \\ 
\mathbf{b}_{2} \\ 
\mathbf{b}_{3} \\ 
\mathbf{b}_{4}%
\end{array}
\right) =\left\vert \mathbf{e},J\mathbf{e},\mathbf{e},J\mathbf{e}%
\right\rangle \left( 
\begin{array}{c}
\mathbf{a}_{3} \\ 
-\mathbf{a}_{4} \\ 
\mathbf{a}_{1} \\ 
-\mathbf{a}_{2}%
\end{array}
\right)
\end{align}%
\begin{equation}
\Sigma_{\sigma}\left( x\right) =\dsum \limits_{1\leq j\leq r}\left\{
\left\vert e_{j}\right\rangle a_{3j}-J\left\vert e_{j}\right\rangle
a_{4j}+\left\vert e_{j}\right\rangle a_{1j}-J\left\vert e_{j}\right\rangle
a_{2j}\right\}
\end{equation}%
\begin{equation}
\mathcal{R}_{e}\left( \Sigma_{\sigma}\right) =\left( 
\begin{array}{cccc}
0 & 0 & I & 0 \\ 
0 & 0 & 0 & -I \\ 
I & 0 & 0 & 0 \\ 
0 & -I & 0 & 0%
\end{array}
\right)
\end{equation}%
\begin{align}
\Sigma_{\sigma}\left( x\right) & =x \\
\left( 
\begin{array}{c}
\mathbf{b}_{1} \\ 
\mathbf{b}_{2} \\ 
\mathbf{b}_{3} \\ 
\mathbf{b}_{4}%
\end{array}
\right) & =\left( 
\begin{array}{c}
\mathbf{a}_{3} \\ 
-\mathbf{a}_{4} \\ 
\mathbf{a}_{1} \\ 
-\mathbf{a}_{2}%
\end{array}
\right) =\left( 
\begin{array}{c}
\mathbf{a}_{1} \\ 
\mathbf{a}_{2} \\ 
\mathbf{a}_{3} \\ 
\mathbf{a}_{4}%
\end{array}
\right) \Longleftrightarrow x=\left( 
\begin{array}{c}
\mathbf{a}_{1} \\ 
\mathbf{0} \\ 
\mathbf{a}_{1} \\ 
\mathbf{0}%
\end{array}
\right)
\end{align}

$\lceil$The set of vectors $\left\{ \left. v+\sigma\left( v\right)
\right\vert v\in\mathfrak{H}\right\} $ can also be thought of as belonging
to the larger Hilbert space 
\begin{equation}
\mathcal{H}=\mathfrak{H}\mathcal{\oplus}\overline{\mathfrak{H}}\equiv 
\mathfrak{H}\mathcal{\oplus\sigma}\left( \mathfrak{H}\right) ,
\end{equation}
where $\mathcal{H}$ carries the natural conjugations 
\begin{align}
\Sigma_{\sigma} & :\mathcal{H}\rightarrow\mathcal{H} \\
\Sigma_{\sigma}\left( u+\sigma\left( v\right) \right) & =v+\sigma\left(
u\right) ,
\end{align}
leading to%
\begin{equation}
V_{\sigma}=\left\{ \left. x\in\mathcal{H}\right\vert \Sigma_{\sigma}\left(
x\right) =x\right\} \equiv\left\{ \left. u+\sigma\left( u\right) \in\mathcal{%
H}\right\vert u\in\mathfrak{H}\right\}
\end{equation}
being a real subspace such that 
\begin{equation}
\mathfrak{H}\equiv V_{\sigma\mathbb{C}}.
\end{equation}%
\begin{equation}
\mathcal{H}=V_{\sigma}\oplus iV_{\sigma}\oplus V_{\sigma}\ominus iV_{\sigma}
\end{equation}
$\rfloor$

$\lceil$One can define a real linear isomorphism ? 
\begin{equation}
\beta:\mathfrak{H}\rightarrow V_{\sigma}\subset\mathfrak{H\oplus}%
\sigma\left( \mathfrak{H}\right) =\mathcal{H}
\end{equation}
by%
\begin{equation}
\beta\left( x\right) =\frac{1}{2}\left\{ x+\sigma\left( x\right) \right\}
;x\in\mathfrak{H,}
\end{equation}
that induces a complex structure on $V_{\sigma}$.$\rfloor$

\begin{equation}
\left\langle \left. \left( x+\overline{x}\right) \right\vert i\left( x-%
\overline{x}\right) \right\rangle =i\left\langle \left. x\right\vert
x\right\rangle -i\left\langle \left. x\right\vert \overline{x}\right\rangle
+i\left\langle \left. \overline{x}\right\vert x\right\rangle -i\left\langle
\left. \overline{x}\right\vert \overline{x}\right\rangle
\end{equation}

\subsection{Conformal group}

Being isometries, real orthogonal transforms preserve angles, and are thus
conformal maps, though not all conformal linear transforms are orthogonal.
In classical terms this is the difference between \emph{congruence} and 
\emph{similarity}, as exemplified by $SSS$ (side-side-side) \emph{congruence
of triangles} and $AAA$ (angle-angle-angle) \emph{similarity of triangles}.
The group of conformal linear maps of $\mathbb{R}^{n}$ is denoted $CO(n,%
\mathbb{R})$ for the conformal orthogonal group, and consists of the product
of the orthogonal group with the group of dilations. If $n$ is odd, these
two subgroups do not intersect, and they are a direct product: 
\begin{equation}
CO(2k+1)=O(2k+1,\mathbb{R})\times \mathbb{R}_{\ast },
\end{equation}%
where $\mathbb{R}_{\ast }=\mathbb{R}\setminus \{0\}$ is the real
multiplicative group, while if $n$ is even, these subgroups intersect in $%
\pm I$, so this is not a direct product, but it is a direct product with the
subgroup of dilation by a positive scalar: 
\begin{equation}
CO(2k,\mathbb{R})=O(2k,\mathbb{R})\times \mathbb{R}_{+}.
\end{equation}%
Similarly one can define $CSO(n,\mathbb{R})$; note that this is always: 
\begin{equation}
CSO(n,\mathbb{R})=CO(n,\mathbb{R})\cap GL_{+}(n,\mathbb{R})=SO(n,\mathbb{R}%
)\times \mathbb{R}_{+}.
\end{equation}

\bigskip

\section{Misc}

\subsection{Invariance Groups}

The invariance groups of the real linear maps $J_{V}$ and $\sigma _{V}$ are
given by 
\begin{equation}
Sp\left( V_{\mathbb{R}},\mathbb{R}\right) =\left\{ \left. t\in GL\left( V_{%
\mathbb{R}},\mathbb{R}\right) \right\vert tJ_{V}t^{t}=J_{V}\right\}
\end{equation}%
and 
\begin{equation}
O\left( V,V^{\bot },\mathbb{R}\right) =\left\{ \left. t\in GL\left( V_{%
\mathbb{R}},\mathbb{R}\right) \right\vert t\sigma _{V}t^{t}=\sigma
_{V}\right\} .
\end{equation}

\emph{Polar decomposition} of the group elements give 
\begin{equation*}
t=o_{1}\left( t^{t}t\right) ^{\frac{1}{2}}=\left( t^{t}t\right) ^{\frac{1}{2}%
}o_{2}=s_{1}\left( tt^{t}\right) ^{\frac{1}{2}}=\left( tt^{t}\right) ^{\frac{%
1}{2}}s_{2},
\end{equation*}%
where 
\begin{equation}
o_{1},o_{2},s_{1},s_{2}\in O\left( V_{\mathbb{R}},\mathbb{R}\right) ,t\in
GL\left( V_{\mathbb{R}},\mathbb{R}\right) \subset \mathcal{B}\left( V_{%
\mathbb{R}}\right) .
\end{equation}

In general $t\in GL\left( V_{\mathbb{R}},\mathbb{R}\right) $ transforms one 
\emph{linearly independent basis} of $V_{\mathbb{R}}$ to another, while $%
t=\alpha o_{2},o_{2}\in O\left( V_{\mathbb{R}},\mathbb{R}\right) ,\alpha
>0,\alpha \in \mathfrak{A,}$ where $\mathfrak{A}$ is a maximal abelian
subgroup of $GL\left( V_{\mathbb{R}},\mathbb{R}\right) $ (Maximal Torus)
transforms \emph{one orthogonal basis} of $V_{\mathbb{R}}$ to another and $%
t=o_{2}\in O\left( V_{\mathbb{R}},\mathbb{R}\right) $ transforms one \emph{%
orthonormal basis} of $V_{\mathbb{R}}$ to another.

\begin{definition}
Let $\mathfrak{A}$ be defined by operators of the form 
\begin{equation}
\alpha =\sum_{1\leq j\leq 2r}\left\vert e_{j}\right) \alpha _{j}\left(
e_{j}\right\vert ,
\end{equation}%
where $\left\{ \left. e_{j}\right\vert 1\leq j\leq 2r\right\} $ is an
orthogonal basis of $V_{\mathbb{R}}$ wrt $\left( \left. {}\right\vert
\right) $ and $\alpha _{j}>0,1\leq j\leq 2r.$
\end{definition}

Then

\begin{enumerate}
\item 
\begin{equation}
\alpha \beta =\sum_{1\leq j\leq 2r}\left\vert e_{j}\right) \alpha _{j}\left(
e_{j}\right\vert \sum_{1\leq j\leq 2r}\left\vert e_{k}\right) \beta
_{k}\left( e_{k}\right\vert =\sum_{1\leq j\leq 2r}\left\vert e_{j}\right)
\alpha _{j}\beta _{j}\kappa _{j}\left( e_{j}\right\vert \quad \limfunc{Iff}%
\quad \left( \left. e_{j}\right\vert e_{k}\right) =\kappa _{j}\delta
_{jk},\kappa _{j}\,>0,1\leq j,k\leq 2r
\end{equation}

\item The identity element, $P_{e},$ of the group $\mathfrak{A}$ is given by 
\begin{equation}
P_{e}=\sum_{1\leq j\leq 2r}\left\vert e_{j}\right) \kappa _{j}^{-1}\left(
e_{j}\right\vert .
\end{equation}%
$\lceil $%
\begin{eqnarray}
\sum_{1\leq j\leq 2r}\left\vert e_{j}\right) \alpha _{j}\left(
e_{j}\right\vert \sum_{1\leq j\leq 2r}\left\vert e_{j}\right) \alpha
_{j}\left( e_{j}\right\vert &=&\sum_{1\leq j\leq 2r}\left\vert e_{j}\right)
\alpha _{j}^{2}\kappa _{j}\left( e_{j}\right\vert =\sum_{1\leq j\leq
2r}\left\vert e_{j}\right) \alpha _{j}\left( e_{j}\right\vert \\
\kappa _{j}\alpha _{j} &=&1 \\
\sum_{1\leq j\leq 2r}\left\vert e_{j}\right) \kappa _{j}^{-1}\left(
e_{j}\right\vert \sum_{1\leq j\leq 2r}\left\vert e_{j}\right) \kappa
_{j}^{-1}\left( e_{j}\right\vert &=&\sum_{1\leq j\leq 2r}\left\vert
e_{j}\right) \kappa _{j}^{-2}\kappa _{j}\left( e_{j}\right\vert =\sum_{1\leq
j\leq 2r}\left\vert e_{j}\right) \kappa _{j}^{-1}\left( e_{j}\right\vert .
\end{eqnarray}%
$\rfloor $

\item 
\begin{equation}
\left[ \alpha ,\beta \right] =0.
\end{equation}

\item 
\begin{equation}
\alpha ,\beta >0\Rightarrow \alpha \beta >0.
\end{equation}

\item 
\begin{eqnarray}
\text{If }\alpha \beta &=&\sum_{1\leq j\leq 2r}\left\vert e_{j}\right)
\kappa _{j}^{-1}\left( e_{j}\right\vert =\sum_{1\leq j\leq 2r}\left\vert
e_{j}\right) \alpha _{j}\beta _{j}\kappa _{j}\left( e_{j}\right\vert \\
\text{then }\beta &=&\alpha ^{-1}=\sum_{1\leq j\leq 2r}\left\vert
e_{j}\right) \alpha _{j}^{-1}\kappa _{j}^{-2}\left( e_{j}\right\vert .
\end{eqnarray}

\item Isomorphic subgroups are given by the conjugation i.e. $T\mathfrak{A}%
T^{-1}\cong \mathfrak{A;}T\in GL\left( V_{\mathbb{R}},\mathbb{R}\right) $ 
\begin{eqnarray}
T\alpha T^{-1}T\beta T^{-1} &=&T\alpha \beta T^{-1} \\
\left[ T\alpha T^{-1},T\beta T^{-1}\right] &=&T\left[ \alpha ,\beta \right]
T^{-1}=0 \\
T\alpha T^{-1} &=&\sum_{1\leq j\leq 2r}T\left\vert e_{j}\right) \alpha
_{j}\left( e_{j}\right\vert T^{-1} \\
T\alpha \beta T^{-1} &=&\sum_{1\leq j\leq 2r}T\left\vert e_{j}\right) \alpha
_{j}\beta _{j}\kappa _{j}\left( e_{j}\right\vert T^{-1}=\sum_{1\leq j\leq
2r}\left\vert Te_{j}\right) \alpha _{j}\beta _{j}\kappa _{j}]\left(
T^{-t}e_{j}\right\vert \\
\left( \left. T^{-t}e_{j}\right\vert Te_{k}\right) &=&\left( \left.
e_{j}\right\vert e_{k}\right) .
\end{eqnarray}

\item The bases $\left\{ \left. Te_{j}\right\vert 1\leq j\leq 2r\right\} $
and $\left\{ \left. T^{-t}e_{j}\right\vert 1\leq j\leq 2r\right\} $ are
biorthogonal, so to define $\mathfrak{A}$ one only need biorthogonal bases
and then consider the operators 
\begin{equation}
\alpha =\sum_{1\leq j\leq 2r}\left\vert Te_{j}\right) \alpha _{j}\left(
T^{-t}e_{j}\right\vert =\sum_{1\leq j\leq 2r}\left\vert f_{j}\right) \alpha
_{j}\left( h_{j}\right\vert ,
\end{equation}%
where $\left\{ \left. f_{j}\right\vert 1\leq j\leq 2r\right\} $ and $\left\{
\left. h_{j}\right\vert 1\leq j\leq 2r\right\} $ are biorthogonal wrt $%
\left( \left. {}\right\vert \right) .$

\item $T$ can belong to $O\left( V_{\mathbb{R}},\mathbb{R}\right) \subset
GL\left( V_{\mathbb{R}},\mathbb{R}\right) $ then $T^{-t}=T$ and then the
basis $\left\{ \left. Te_{j}\right\vert 1\leq j\leq 2r\right\} $ is
orthogonal.

\item \textbf{The product of positive operator is not in general positive as
the product of self adjoint opeators is not in general self adjoint, }however%
\textbf{\ }the set of positive elements form 
\begin{equation}
\mathfrak{A}_{+}=\left\{ \left. \alpha \in \mathfrak{A}\right\vert \alpha
>0\right\}
\end{equation}%
\emph{a normal subgroup} of $GL\left( V_{\mathbb{R}},\mathbb{R}\right) $
i.e. 
\begin{equation}
T_{1}\sum_{1\leq j\leq 2r}\left\vert f_{j}\right) \alpha _{j}\left(
h_{j}\right\vert T_{1}^{-1}\in \mathfrak{A}_{+}\,\forall T_{1}\in GL\left(
V_{\mathbb{R}},\mathbb{R}\right) .
\end{equation}%
$\lceil $%
\begin{equation}
\exp \left\{ \zeta _{j}\left( \left\vert e_{j}\right) \left(
e_{j}\right\vert \right) \right\} =I+\zeta _{j}\left( \left\vert
e_{j}\right) \left( e_{j}\right\vert \right) +\frac{1}{2}\zeta
_{j}^{2}\left( \left\vert e_{j}\right) \left( e_{j}\right\vert \right)
+\cdots =I+\left( e^{\zeta _{j}}-1\right) \left\vert e_{j}\right) \left(
e_{j}\right\vert
\end{equation}%
\begin{eqnarray}
O\left( \mathbb{\lambda }\right) &=&\exp \left\{ \sum_{1\leq j<k\leq
2r}\lambda _{jk}\left( \left\vert e_{j}\right) \left( e_{k}\right\vert
-\left\vert e_{k}\right) \left( e_{j}\right\vert \right) \right\} \\
A\left( \mathbb{\zeta }\right) &=&\exp \left\{ \sum_{1\leq j<k\leq 2r}\zeta
_{j}\left( \left\vert e_{j}\right) \left( e_{j}\right\vert \right) \right\}
=\tprod\limits_{1\leq j\leq 2r}\exp \left\{ \zeta _{j}\left( \left\vert
e_{j}\right) \left( e_{j}\right\vert \right) \right\} =\tprod\limits_{1\leq
j\leq 2r}\left( I_{2r}+\left( e^{\zeta _{j}}-1\right) \right) \left\vert
e_{j}\right) \left( e_{j}\right\vert \\
&=&\left( I_{2r}+\left( e^{\zeta _{1}}-1\right) \right) \left\vert
e_{1}\right) \left( e_{1}\right\vert \left( I_{2r}+\left( e^{\zeta
_{2}}-1\right) \right) \left\vert e_{2}\right) \left( e_{2}\right\vert
\cdots \left( I_{2r}+\left( e^{\zeta _{2r}}-1\right) \right) \left\vert
e_{2r}\right) \left( e_{2r}\right\vert \\
&=&\left( I_{2r}+\left( e^{\zeta _{1}}-1\right) \left\vert e_{1}\right)
\left( e_{1}\right\vert \right) \left( I_{2r}+\left( e^{\zeta _{2}}-1\right)
\left\vert e_{2}\right) \left( e_{2}\right\vert \right) \cdots \left(
I_{2r}+\left( e^{\zeta _{2r}}-1\right) \left\vert e_{2r}\right) \left(
e_{2r}\right\vert \right) \\
&=&I_{2r}+\left( e^{\zeta _{1}}-1\right) \left\vert e_{1}\right) \left(
e_{1}\right\vert +\left( e^{\zeta _{2}}-1\right) \left\vert e_{2}\right)
\left( e_{2}\right\vert +\cdots +\left( e^{\zeta _{2r}}-1\right) \left\vert
e_{2r}\right) \left( e_{2r}\right\vert \\
&=&e^{\zeta _{1}}\left\vert e_{1}\right) \left( e_{1}\right\vert +e^{\zeta
_{2}}\left\vert e_{2}\right) \left( e_{2}\right\vert +\cdots +e^{\zeta
_{2r}}\left\vert e_{2r}\right) \left( e_{2r}\right\vert
\end{eqnarray}%
In particular 
\begin{equation}
A\left( \mathbb{\zeta }\right) \left\vert e_{j}\right) =e^{\zeta
_{j}}\left\vert e_{j}\right) =\eta _{j}\left\vert e_{j}\right) ;\,\eta
_{j}\in \left( 0,\infty \right) ,1\leq j\leq 2r
\end{equation}%
so that if $\left\{ \left. \left\vert e_{j}\right) \right\vert 1\leq j\leq
2r\right\} $ then $\left\{ \left. A\left( \mathbb{\zeta }\right) \left\vert
e_{j}\right) \right\vert 1\leq j\leq 2r\right\} $ are also orthogonal, in
general however $A\left( \mathbb{\zeta }\right) $ does \emph{not maintain
orthogonality} as 
\begin{eqnarray}
v_{1} &=&\sum_{1\leq j\leq 2r}v_{1j}\left\vert e_{j}\right) \\
v_{2} &=&\sum_{1\leq j\leq 2r}v_{2j}\left\vert e_{j}\right) \\
\left( \left. v_{1}\right\vert v_{2}\right) &=&\sum_{1\leq j\leq
2r}v_{1j}v_{2j} \\
\left( \left. A\left( \mathbb{\zeta }\right) v_{1}\right\vert A\left( 
\mathbb{\zeta }\right) v_{2}\right) &=&\sum_{1\leq j\leq 2r}e^{2\zeta
_{j}}v_{1j}v_{2j}
\end{eqnarray}%
and 
\begin{equation}
\left( \left. v_{1}\right\vert v_{2}\right) =0\nRightarrow \left( \left.
A\left( \mathbb{\zeta }\right) v_{1}\right\vert A\left( \mathbb{\zeta }%
\right) v_{2}\right) =0.
\end{equation}%
$\rfloor $
\end{enumerate}

\begin{theorem}
All orthogonal bases $\left\{ \left. \left\vert f_{j}\right) \right\vert
1\leq j\leq 2r\right\} $ can be generated from a given orthogonal basis $%
\left\{ \left. \left\vert e_{j}\right) \right\vert 1\leq j\leq 2r\right\} $
by the action of $O\left( \mathbb{\lambda }\right) A\left( \mathbb{\zeta }%
\right) $ on $\left\{ \left. \left\vert e_{j}\right) \right\vert 1\leq j\leq
2r\right\} ,$ where%
\begin{equation}
A\left( \mathbb{\zeta }\right) =\exp \left\{ \sum_{1\leq j\leq 2r}\zeta
_{j}\left( \left\vert e_{j}\right) \left( e_{j}\right\vert \right) \right\}
\in \mathfrak{A}_{+}
\end{equation}%
and 
\begin{equation}
O\left( \mathbb{\lambda }\right) =\exp \left\{ \sum_{1\leq j<k\leq
2r}\lambda _{jk}\left( \left\vert e_{j}\right) \left( e_{k}\right\vert
-\left\vert e_{k}\right) \left( e_{j}\right\vert \right) \right\} \in
O\left( V_{\mathbb{R}},\mathbb{R}\right) .
\end{equation}

\begin{proof}
\begin{equation}
\left( \left. O\left( \mathbb{\lambda }\right) A\left( \mathbb{\zeta }%
\right) e_{j}\right\vert O\left( \mathbb{\lambda }\right) A\left( \mathbb{%
\zeta }\right) e_{k}\right) =\left( \left. e_{j}\right\vert A\left( \mathbb{%
\zeta }\right) ^{2}e_{k}\right) =\alpha _{j}^{2}\left( \left.
e_{j}\right\vert e_{k}\right) =\alpha _{j}^{2}\delta _{jk}.
\end{equation}
\end{proof}
\end{theorem}

The invariance tranformations of $J_{V}$ and $\sigma _{V}$ then produce 
\begin{align}
t\sigma _{V}t^{t}& =\dsum\limits_{1\leq j\leq r}\left\{ \left\vert
te_{j}\right) \left( te_{j}\right\vert -\left\vert te_{j+r}\right) \left(
te_{j+r}\right\vert \right\} =\dsum\limits_{1\leq j\leq r}\left\{ \left\vert
e_{j}\right) \left( e_{j}\right\vert -\left\vert e_{j+r}\right) \left(
e_{j+r}\right\vert \right\} \\
tJ_{V}t^{t}& =\dsum\limits_{1\leq j\leq r}\left\{ \left\vert te_{j}\right)
\left( te_{j+r}\right\vert -\left\vert te_{j+r}\right) \left(
te_{j}\right\vert \right\} =\dsum\limits_{1\leq j\leq r}\left\{ \left\vert
e_{j}\right) \left( e_{j+r}\right\vert -\left\vert e_{j+r}\right) \left(
e_{j}\right\vert \right\} ,
\end{align}%
where%
\begin{align}
t\sigma _{V}t^{t}& =\dsum\limits_{1\leq j\leq r}\left\{ \left\vert
o_{1}\left( t^{t}t\right) ^{\frac{1}{2}}e_{j}\right) \left( o_{1}\left(
t^{t}t\right) ^{\frac{1}{2}}e_{j}\right\vert -\left\vert o_{1}\left(
t^{t}t\right) ^{\frac{1}{2}}e_{j+r}\right) \left( o_{1}\left( t^{t}t\right)
^{\frac{1}{2}}e_{j+r}\right\vert \right\} \\
tJ_{V}t^{t}& =\dsum\limits_{1\leq j\leq r}\left\{ \left\vert o_{1}\left(
t^{t}t\right) ^{\frac{1}{2}}e_{j}\right) \left( o_{1}\left( t^{t}t\right) ^{%
\frac{1}{2}}e_{j+r}\right\vert -\left\vert o_{1}\left( t^{t}t\right) ^{\frac{%
1}{2}}e_{j+r}\right) \left( o_{1}\left( t^{t}t\right) ^{\frac{1}{2}%
}e_{j}\right\vert \right\} ,
\end{align}%
giving%
\begin{align}
t\sigma _{V}t^{t}& =\dsum\limits_{1\leq j\leq r}\left\{ o_{1}\left(
t^{t}t\right) ^{\frac{1}{2}}\left\vert e_{j}\right) \left( e_{j}\right\vert
\left( t^{t}t\right) ^{\frac{1}{2}}o_{1}^{t}-o_{1}\left( t^{t}t\right) ^{%
\frac{1}{2}}\left\vert e_{j+r}\right) \left( e_{j+r}\right\vert \left(
t^{t}t\right) ^{\frac{1}{2}}o_{1}^{t}\right\} \\
tJ_{V}t^{t}& =\dsum\limits_{1\leq j\leq r}\left\{ o_{1}\left( t^{t}t\right)
^{\frac{1}{2}}\left\vert e_{j}\right) \left( e_{j+r}\right\vert \left(
t^{t}t\right) ^{\frac{1}{2}}o_{1}^{t}-o_{1}\left( t^{t}t\right) ^{\frac{1}{2}%
}\left\vert e_{j+r}\right) \left( e_{j}\right\vert \left( t^{t}t\right) ^{%
\frac{1}{2}}o_{1}^{t}\right\} .
\end{align}%
$\lceil $The \emph{polar decomposition is the same as the SVD} i.e.

The polar decomposition gives%
\begin{equation}
t=\left( t^{t}t\right) ^{\frac{1}{2}}o_{2}
\end{equation}%
and leads directlt to the SVD%
\begin{equation}
\left( t^{t}t\right) ^{\frac{1}{2}}o_{2}=\sum_{1\leq j\leq 2r}\left\vert
f_{j}\right) \alpha _{j}\left( f_{j}\right\vert o_{2}=\sum_{1\leq j\leq
2r}\left\vert f_{j}\right) \alpha _{j}\left( o_{2}^{t}f_{j}\right\vert
=\sum_{1\leq j\leq 2r}\left\vert f_{j}\right) \alpha _{j}\left(
h_{j}\right\vert ,
\end{equation}%
where we have used the \emph{spectral theorem} for the real symmetric matrix 
$\left( t^{t}t\right) ^{\frac{1}{2}}$ i.e.%
\begin{equation}
\left( t^{t}t\right) ^{\frac{1}{2}}=\sum_{1\leq j\leq 2r}\left\vert
f_{j}\right) \alpha _{j}\left( f_{j}\right\vert .
\end{equation}%
We then consider the matrix representations given by 
\begin{eqnarray}
\left( \mathbf{h}\right\vert &=&\left( \mathbf{h}\right\vert \left\vert 
\mathbf{e}\right) \left( \mathbf{e}\right\vert =\mathbb{O}_{eh}^{t}\left( 
\mathbf{e}\right\vert \\
\left\vert \mathbf{h}\right) &=&\left\vert \mathbf{e}\right) \left( \mathbf{e%
}\right\vert \left\vert \mathbf{h}\right) =\left\vert \mathbf{e}\right) 
\mathbb{O}_{eh}.
\end{eqnarray}%
that lead to 
\begin{equation}
\sum_{1\leq j\leq 2r}\left\vert f_{j}\right) \alpha _{j}\left(
o_{eh}^{t}f_{j}\right\vert \left\vert \mathbf{e}\right) =\sum_{1\leq j\leq
2r}\left\vert f_{j}\right) \alpha _{j}\left( h_{j}\right\vert \left\vert 
\mathbf{e}\right) =\left\vert \mathbf{f}\right) \mathbb{\alpha }\left( 
\mathbf{h}\right\vert \left\vert \mathbf{e}\right) =\left\vert \mathbf{f}%
\right) \mathbb{\alpha O}_{eh}^{t}\left( \mathbf{e}\right\vert \left\vert 
\mathbf{e}\right) =\left\vert \mathbf{f}\right) \mathbb{\alpha O}%
_{eh}^{t}\left( \mathbf{e}\right\vert =\left\vert \mathbf{f}\right) \mathbb{%
\alpha }\left( \mathbf{h}\right\vert ,
\end{equation}%
which is the SVD i.e.%
\begin{eqnarray}
\left\vert \mathbf{e}\right) \left( \overset{}{\left( \mathbf{e}\right\vert
\left\vert \mathbf{f}\right) \mathbb{\alpha }\left( \mathbf{h}\right\vert
\left\vert \mathbf{e}\right) }\right) \left( \mathbf{e}\right\vert
&=&\left\vert \mathbf{e}\right) \left( \overset{}{\mathcal{R}_{e}\left(
\left\vert \mathbf{f}\right) \mathbb{\alpha }\left( \mathbf{h}\right\vert
\right) }\right) \left( \mathbf{e}\right\vert \\
&\Downarrow & \\
\mathcal{R}_{e}\left( \left\vert \mathbf{f}\right) \mathbb{\alpha }\left( 
\mathbf{h}\right\vert \right) &=&\left( \mathbf{e}\right\vert \left\vert 
\mathbf{f}\right) \mathbb{\alpha }\left( \mathbf{h}\right\vert \left\vert 
\mathbf{e}\right) =\mathbb{O}_{ef}\mathbb{\alpha O}_{he}^{t}.
\end{eqnarray}

\end{document}
